% generated by GAPDoc2LaTeX from XML source (Frank Luebeck)
\documentclass[a4paper,11pt]{report}

            \usepackage{a4wide}
            \newcommand{\bbZ}{\mathbb{Z}}
        
\usepackage[top=37mm,bottom=37mm,left=27mm,right=27mm]{geometry}
\sloppy
\pagestyle{myheadings}
\usepackage{amssymb}
\usepackage[utf8]{inputenc}
\usepackage{makeidx}
\makeindex
\usepackage{color}
\definecolor{FireBrick}{rgb}{0.5812,0.0074,0.0083}
\definecolor{RoyalBlue}{rgb}{0.0236,0.0894,0.6179}
\definecolor{RoyalGreen}{rgb}{0.0236,0.6179,0.0894}
\definecolor{RoyalRed}{rgb}{0.6179,0.0236,0.0894}
\definecolor{LightBlue}{rgb}{0.8544,0.9511,1.0000}
\definecolor{Black}{rgb}{0.0,0.0,0.0}

\definecolor{linkColor}{rgb}{0.0,0.0,0.554}
\definecolor{citeColor}{rgb}{0.0,0.0,0.554}
\definecolor{fileColor}{rgb}{0.0,0.0,0.554}
\definecolor{urlColor}{rgb}{0.0,0.0,0.554}
\definecolor{promptColor}{rgb}{0.0,0.0,0.589}
\definecolor{brkpromptColor}{rgb}{0.589,0.0,0.0}
\definecolor{gapinputColor}{rgb}{0.589,0.0,0.0}
\definecolor{gapoutputColor}{rgb}{0.0,0.0,0.0}

%%  for a long time these were red and blue by default,
%%  now black, but keep variables to overwrite
\definecolor{FuncColor}{rgb}{0.0,0.0,0.0}
%% strange name because of pdflatex bug:
\definecolor{Chapter }{rgb}{0.0,0.0,0.0}
\definecolor{DarkOlive}{rgb}{0.1047,0.2412,0.0064}


\usepackage{fancyvrb}

\usepackage{mathptmx,helvet}
\usepackage[T1]{fontenc}
\usepackage{textcomp}


\usepackage[
            pdftex=true,
            bookmarks=true,        
            a4paper=true,
            pdftitle={Written with GAPDoc},
            pdfcreator={LaTeX with hyperref package / GAPDoc},
            colorlinks=true,
            backref=page,
            breaklinks=true,
            linkcolor=linkColor,
            citecolor=citeColor,
            filecolor=fileColor,
            urlcolor=urlColor,
            pdfpagemode={UseNone}, 
           ]{hyperref}

\newcommand{\maintitlesize}{\fontsize{50}{55}\selectfont}

% write page numbers to a .pnr log file for online help
\newwrite\pagenrlog
\immediate\openout\pagenrlog =\jobname.pnr
\immediate\write\pagenrlog{PAGENRS := [}
\newcommand{\logpage}[1]{\protect\write\pagenrlog{#1, \thepage,}}
%% were never documented, give conflicts with some additional packages

\newcommand{\GAP}{\textsf{GAP}}

%% nicer description environments, allows long labels
\usepackage{enumitem}
\setdescription{style=nextline}

%% depth of toc
\setcounter{tocdepth}{1}





%% command for ColorPrompt style examples
\newcommand{\gapprompt}[1]{\color{promptColor}{\bfseries #1}}
\newcommand{\gapbrkprompt}[1]{\color{brkpromptColor}{\bfseries #1}}
\newcommand{\gapinput}[1]{\color{gapinputColor}{#1}}


\begin{document}

\logpage{[ 0, 0, 0 ]}
\begin{titlepage}
\mbox{}\vfill

\begin{center}{\maintitlesize \textbf{ Semigroups \mbox{}}}\\
\vfill

\hypersetup{pdftitle= Semigroups }
\markright{\scriptsize \mbox{}\hfill  Semigroups  \hfill\mbox{}}
{\Huge \textbf{ A package for semigroups and monoids \mbox{}}}\\
\vfill

{\Huge  5.6.1 \mbox{}}\\[1cm]
{ 18 March 2026 \mbox{}}\\[1cm]
\mbox{}\\[2cm]
{\Large \textbf{ James Mitchell\\
    \mbox{}}}\\
{\Large \textbf{ Marina Anagnostopoulou\texttt{\symbol{45}}Merkouri\\
   \mbox{}}}\\
{\Large \textbf{ Thomas Breuer\\
   \mbox{}}}\\
{\Large \textbf{ Stuart Burrell\\
   \mbox{}}}\\
{\Large \textbf{ Reinis Cirpons\\
    \mbox{}}}\\
{\Large \textbf{ Tom Conti\texttt{\symbol{45}}Leslie\\
   \mbox{}}}\\
{\Large \textbf{ Joseph Edwards\\
    \mbox{}}}\\
{\Large \textbf{ Attila Egri\texttt{\symbol{45}}Nagy\\
   \mbox{}}}\\
{\Large \textbf{ Luke Elliott\\
    \mbox{}}}\\
{\Large \textbf{ Fernando Flores Brito\\
  \mbox{}}}\\
{\Large \textbf{ Tillman Froehlich\\
  \mbox{}}}\\
{\Large \textbf{ Nick Ham\\
   \mbox{}}}\\
{\Large \textbf{ Robert Hancock\\
   \mbox{}}}\\
{\Large \textbf{ Max Horn\\
    \mbox{}}}\\
{\Large \textbf{ Christopher Jefferson\\
    \mbox{}}}\\
{\Large \textbf{ Julius Jonusas\\
   \mbox{}}}\\
{\Large \textbf{ Chinmaya Nagpal\\
 \mbox{}}}\\
{\Large \textbf{ Olexandr Konovalov\\
    \mbox{}}}\\
{\Large \textbf{ Artemis Konstantinidi\\
 \mbox{}}}\\
{\Large \textbf{ Hyeokjun Kwon\\
 \mbox{}}}\\
{\Large \textbf{ Dima V. Pasechnik\\
    \mbox{}}}\\
{\Large \textbf{ Yann Peresse\\
  \mbox{}}}\\
{\Large \textbf{ Markus Pfeiffer\\
   \mbox{}}}\\
{\Large \textbf{ Pramoth Ragavan\\
  \mbox{}}}\\
{\Large \textbf{ Joesph Daynger Ruiz\\
  \mbox{}}}\\
{\Large \textbf{ Christopher Russell\\
 \mbox{}}}\\
{\Large \textbf{ Jack Schmidt\\
   \mbox{}}}\\
{\Large \textbf{ Sergio Siccha\\
  \mbox{}}}\\
{\Large \textbf{ Finn Smith\\
    \mbox{}}}\\
{\Large \textbf{ Ben Spiers\\
 \mbox{}}}\\
{\Large \textbf{ Nicolas Thi{\a'e}ry\\
   \mbox{}}}\\
{\Large \textbf{ Maria Tsalakou\\
    \mbox{}}}\\
{\Large \textbf{ Chris Wensley\\
  \mbox{}}}\\
{\Large \textbf{ Murray Whyte\\
   \mbox{}}}\\
{\Large \textbf{ Wilf A. Wilson\\
   \mbox{}}}\\
{\Large \textbf{ Tianrun Yang\\
 \mbox{}}}\\
{\Large \textbf{ Michael Young\\
    \mbox{}}}\\
{\Large \textbf{ Fabian Zickgraf\\
  \mbox{}}}\\
\hypersetup{pdfauthor= James Mitchell\\
    ;  Marina Anagnostopoulou\texttt{\symbol{45}}Merkouri\\
   ;  Thomas Breuer\\
   ;  Stuart Burrell\\
   ;  Reinis Cirpons\\
    ;  Tom Conti\texttt{\symbol{45}}Leslie\\
   ;  Joseph Edwards\\
    ;  Attila Egri\texttt{\symbol{45}}Nagy\\
   ;  Luke Elliott\\
    ;  Fernando Flores Brito\\
  ;  Tillman Froehlich\\
  ;  Nick Ham\\
   ;  Robert Hancock\\
   ;  Max Horn\\
    ;  Christopher Jefferson\\
    ;  Julius Jonusas\\
   ;  Chinmaya Nagpal\\
 ;  Olexandr Konovalov\\
    ;  Artemis Konstantinidi\\
 ;  Hyeokjun Kwon\\
 ;  Dima V. Pasechnik\\
    ;  Yann Peresse\\
  ;  Markus Pfeiffer\\
   ;  Pramoth Ragavan\\
  ;  Joesph Daynger Ruiz\\
  ;  Christopher Russell\\
 ;  Jack Schmidt\\
   ;  Sergio Siccha\\
  ;  Finn Smith\\
    ;  Ben Spiers\\
 ;  Nicolas Thi{\a'e}ry\\
   ;  Maria Tsalakou\\
    ;  Chris Wensley\\
  ;  Murray Whyte\\
   ;  Wilf A. Wilson\\
   ;  Tianrun Yang\\
 ;  Michael Young\\
    ;  Fabian Zickgraf\\
  }
\end{center}\vfill

\mbox{}\\
{\mbox{}\\
\small \noindent \textbf{ James Mitchell\\
    }  Email: \href{mailto://jdm3@st-andrews.ac.uk} {\texttt{jdm3@st\texttt{\symbol{45}}andrews.ac.uk}}\\
  Homepage: \href{https://jdbm.me} {\texttt{https://jdbm.me}}\\
  Address: \begin{minipage}[t]{8cm}\noindent
 Mathematical Institute, North Haugh, St Andrews, Fife, KY16 9SS, Scotland\\
 \end{minipage}
}\\
{\mbox{}\\
\small \noindent \textbf{ Marina Anagnostopoulou\texttt{\symbol{45}}Merkouri\\
   }  Email: \href{mailto://marina.anagnostopoulou-merkouri@bristol.ac.uk} {\texttt{marina.anagnostopoulou\texttt{\symbol{45}}merkouri@bristol.ac.uk}}\\
  Homepage: \href{https://marinaanagno.github.io} {\texttt{https://marinaanagno.github.io}}}\\
{\mbox{}\\
\small \noindent \textbf{ Thomas Breuer\\
   }  Email: \href{mailto://sam@math.rwth-aachen.de} {\texttt{sam@math.rwth\texttt{\symbol{45}}aachen.de}}\\
  Homepage: \href{https://www.math.rwth-aachen.de/~Thomas.Breuer/} {\texttt{https://www.math.rwth\texttt{\symbol{45}}aachen.de/\texttt{\symbol{126}}Thomas.Breuer/}}}\\
{\mbox{}\\
\small \noindent \textbf{ Stuart Burrell\\
   }  Email: \href{mailto://stuartburrell1994@gmail.com} {\texttt{stuartburrell1994@gmail.com}}\\
  Homepage: \href{https://stuartburrell.github.io} {\texttt{https://stuartburrell.github.io}}}\\
{\mbox{}\\
\small \noindent \textbf{ Reinis Cirpons\\
    }  Email: \href{mailto://rc234@st-andrews.ac.uk} {\texttt{rc234@st\texttt{\symbol{45}}andrews.ac.uk}}\\
  Homepage: \href{https://reinisc.id.lv/} {\texttt{https://reinisc.id.lv/}}\\
  Address: \begin{minipage}[t]{8cm}\noindent
 Mathematical Institute, North Haugh, St Andrews, Fife, KY16 9SS, Scotland\\
 \end{minipage}
}\\
{\mbox{}\\
\small \noindent \textbf{ Tom Conti\texttt{\symbol{45}}Leslie\\
   }  Email: \href{mailto://tom.contileslie@gmail.com} {\texttt{tom.contileslie@gmail.com}}\\
  Homepage: \href{https://tomcontileslie.com/} {\texttt{https://tomcontileslie.com/}}}\\
{\mbox{}\\
\small \noindent \textbf{ Joseph Edwards\\
    }  Email: \href{mailto://jde1@st-andrews.ac.uk} {\texttt{jde1@st\texttt{\symbol{45}}andrews.ac.uk}}\\
  Homepage: \href{https://github.com/Joseph-Edwards} {\texttt{https://github.com/Joseph\texttt{\symbol{45}}Edwards}}\\
  Address: \begin{minipage}[t]{8cm}\noindent
 Mathematical Institute, North Haugh, St Andrews, Fife, KY16 9SS, Scotland\\
 \end{minipage}
}\\
{\mbox{}\\
\small \noindent \textbf{ Attila Egri\texttt{\symbol{45}}Nagy\\
   }  Email: \href{mailto://attila@egri-nagy.hu} {\texttt{attila@egri\texttt{\symbol{45}}nagy.hu}}\\
  Homepage: \href{http://www.egri-nagy.hu} {\texttt{http://www.egri\texttt{\symbol{45}}nagy.hu}}}\\
{\mbox{}\\
\small \noindent \textbf{ Luke Elliott\\
    }  Email: \href{mailto://le27@st-andrews.ac.uk} {\texttt{le27@st\texttt{\symbol{45}}andrews.ac.uk}}\\
  Homepage: \href{https://le27.github.io/Luke-Elliott/} {\texttt{https://le27.github.io/Luke\texttt{\symbol{45}}Elliott/}}\\
  Address: \begin{minipage}[t]{8cm}\noindent
 Mathematical Institute, North Haugh, St Andrews, Fife, KY16 9SS, Scotland\\
 \end{minipage}
}\\
{\mbox{}\\
\small \noindent \textbf{ Fernando Flores Brito\\
  }  Email: \href{mailto://ffloresbrito@gmail.com} {\texttt{ffloresbrito@gmail.com}}}\\
{\mbox{}\\
\small \noindent \textbf{ Tillman Froehlich\\
  }  Email: \href{mailto://trf1@st-andrews.ac.uk} {\texttt{trf1@st\texttt{\symbol{45}}andrews.ac.uk}}}\\
{\mbox{}\\
\small \noindent \textbf{ Nick Ham\\
   }  Email: \href{mailto://nicholas.charles.ham@gmail.com} {\texttt{nicholas.charles.ham@gmail.com}}\\
  Homepage: \href{https://n-ham.github.io} {\texttt{https://n\texttt{\symbol{45}}ham.github.io}}}\\
{\mbox{}\\
\small \noindent \textbf{ Robert Hancock\\
   }  Email: \href{mailto://robert.hancock@maths.ox.ac.uk} {\texttt{robert.hancock@maths.ox.ac.uk}}\\
  Homepage: \href{https://sites.google.com/view/robert-hancock/} {\texttt{https://sites.google.com/view/robert\texttt{\symbol{45}}hancock/}}}\\
{\mbox{}\\
\small \noindent \textbf{ Max Horn\\
    }  Email: \href{mailto://mhorn@rptu.de} {\texttt{mhorn@rptu.de}}\\
  Homepage: \href{https://www.quendi.de/math} {\texttt{https://www.quendi.de/math}}\\
  Address: \begin{minipage}[t]{8cm}\noindent
 Fachbereich Mathematik, RPTU Kaiserslautern\texttt{\symbol{45}}Landau,
Gottlieb\texttt{\symbol{45}}Daimler\texttt{\symbol{45}}Stra{\ss}e 48, 67663
Kaiserslautern, Germany\\
 \end{minipage}
}\\
{\mbox{}\\
\small \noindent \textbf{ Christopher Jefferson\\
    }  Email: \href{mailto://caj21@st-andrews.ac.uk} {\texttt{caj21@st\texttt{\symbol{45}}andrews.ac.uk}}\\
  Homepage: \href{https://heather.cafe/} {\texttt{https://heather.cafe/}}\\
  Address: \begin{minipage}[t]{8cm}\noindent
 Jack Cole Building, North Haugh, St Andrews, Fife, KY16 9SX, Scotland\\
 \end{minipage}
}\\
{\mbox{}\\
\small \noindent \textbf{ Julius Jonusas\\
   }  Email: \href{mailto://j.jonusas@gmail.com} {\texttt{j.jonusas@gmail.com}}\\
  Homepage: \href{http://julius.jonusas.work} {\texttt{http://julius.jonusas.work}}}\\
{\mbox{}\\
\small \noindent \textbf{ Olexandr Konovalov\\
    }  Email: \href{mailto://obk1@st-andrews.ac.uk} {\texttt{obk1@st\texttt{\symbol{45}}andrews.ac.uk}}\\
  Homepage: \href{https://olexandr-konovalov.github.io/} {\texttt{https://olexandr\texttt{\symbol{45}}konovalov.github.io/}}\\
  Address: \begin{minipage}[t]{8cm}\noindent
 Jack Cole Building, North Haugh, St Andrews, Fife, KY16 9SX, Scotland\\
 \end{minipage}
}\\
{\mbox{}\\
\small \noindent \textbf{ Dima V. Pasechnik\\
    }  Email: \href{mailto://dmitrii.pasechnik@cs.ox.ac.uk} {\texttt{dmitrii.pasechnik@cs.ox.ac.uk}}\\
  Homepage: \href{http://users.ox.ac.uk/~coml0531/} {\texttt{http://users.ox.ac.uk/\texttt{\symbol{126}}coml0531/}}\\
  Address: \begin{minipage}[t]{8cm}\noindent
 Pembroke College, St. Aldates, Oxford OX1 1DW, England\\
 \end{minipage}
}\\
{\mbox{}\\
\small \noindent \textbf{ Yann Peresse\\
  }  Email: \href{mailto://y.peresse@herts.ac.uk} {\texttt{y.peresse@herts.ac.uk}}}\\
{\mbox{}\\
\small \noindent \textbf{ Markus Pfeiffer\\
   }  Email: \href{mailto://markus.pfeiffer@morphism.de} {\texttt{markus.pfeiffer@morphism.de}}\\
  Homepage: \href{https://markusp.morphism.de/} {\texttt{https://markusp.morphism.de/}}}\\
{\mbox{}\\
\small \noindent \textbf{ Pramoth Ragavan\\
  }  Email: \href{mailto://pramoth.ragavan@gmail.com} {\texttt{pramoth.ragavan@gmail.com}}}\\
{\mbox{}\\
\small \noindent \textbf{ Joesph Daynger Ruiz\\
  }  Email: \href{mailto://jdruiz@arizona.edu} {\texttt{jdruiz@arizona.edu}}}\\
{\mbox{}\\
\small \noindent \textbf{ Jack Schmidt\\
   }  Email: \href{mailto://jack.schmidt@uky.edu} {\texttt{jack.schmidt@uky.edu}}\\
  Homepage: \href{https://www.ms.uky.edu/~jack/} {\texttt{https://www.ms.uky.edu/\texttt{\symbol{126}}jack/}}}\\
{\mbox{}\\
\small \noindent \textbf{ Sergio Siccha\\
  }  Email: \href{mailto://sergio.siccha@gmail.com} {\texttt{sergio.siccha@gmail.com}}}\\
{\mbox{}\\
\small \noindent \textbf{ Finn Smith\\
    }  Email: \href{mailto://fls3@st-andrews.ac.uk} {\texttt{fls3@st\texttt{\symbol{45}}andrews.ac.uk}}\\
  Homepage: \href{https://flsmith.github.io/} {\texttt{https://flsmith.github.io/}}\\
  Address: \begin{minipage}[t]{8cm}\noindent
 Mathematical Institute, North Haugh, St Andrews, Fife, KY16 9SS, Scotland\\
 \end{minipage}
}\\
{\mbox{}\\
\small \noindent \textbf{ Nicolas Thi{\a'e}ry\\
   }  Email: \href{mailto://nthiery@users.sf.net} {\texttt{nthiery@users.sf.net}}\\
  Homepage: \href{https://nicolas.thiery.name/} {\texttt{https://nicolas.thiery.name/}}}\\
{\mbox{}\\
\small \noindent \textbf{ Maria Tsalakou\\
    }  Email: \href{mailto://mt200@st-andrews.ac.uk} {\texttt{mt200@st\texttt{\symbol{45}}andrews.ac.uk}}\\
  Homepage: \href{https://mariatsalakou.github.io/} {\texttt{https://mariatsalakou.github.io/}}\\
  Address: \begin{minipage}[t]{8cm}\noindent
 Mathematical Institute, North Haugh, St Andrews, Fife, KY16 9SS, Scotland\\
 \end{minipage}
}\\
{\mbox{}\\
\small \noindent \textbf{ Chris Wensley\\
  }  Email: \href{mailto://cdwensley.maths@btinternet.com} {\texttt{cdwensley.maths@btinternet.com}}}\\
{\mbox{}\\
\small \noindent \textbf{ Murray Whyte\\
   }  Email: \href{mailto://mw231@st-andrews.ac.uk} {\texttt{mw231@st\texttt{\symbol{45}}andrews.ac.uk}}\\
  Address: \begin{minipage}[t]{8cm}\noindent
 Mathematical Institute, North Haugh, St Andrews, Fife, KY16 9SS, Scotland\\
 \end{minipage}
}\\
{\mbox{}\\
\small \noindent \textbf{ Wilf A. Wilson\\
   }  Email: \href{mailto://gap@wilf-wilson.net} {\texttt{gap@wilf\texttt{\symbol{45}}wilson.net}}\\
  Homepage: \href{https://wilf.me} {\texttt{https://wilf.me}}}\\
{\mbox{}\\
\small \noindent \textbf{ Michael Young\\
    }  Email: \href{mailto://mct25@st-andrews.ac.uk} {\texttt{mct25@st\texttt{\symbol{45}}andrews.ac.uk}}\\
  Homepage: \href{https://mtorpey.github.io/} {\texttt{https://mtorpey.github.io/}}\\
  Address: \begin{minipage}[t]{8cm}\noindent
 Jack Cole Building, North Haugh, St Andrews, Fife, KY16 9SX, Scotland\\
 \end{minipage}
}\\
{\mbox{}\\
\small \noindent \textbf{ Fabian Zickgraf\\
  }  Email: \href{mailto://f.zickgraf@dashdos.com} {\texttt{f.zickgraf@dashdos.com}}}\\
\end{titlepage}

\newpage\setcounter{page}{2}
{\small 
\section*{Abstract}
\logpage{[ 0, 0, 1 ]}
 The Semigroups package is a GAP package for semigroups, and monoids. There are
particularly efficient methods for finitely presented semigroups and monoids,
and for semigroups and monoids consisting of transformations, partial
permutations, bipartitions, partitioned binary relations, subsemigroups of
regular Rees 0\texttt{\symbol{45}}matrix semigroups, and matrices of various
semirings including boolean matrices, matrices over finite fields, and certain
tropical matrices. Semigroups contains efficient methods for creating
semigroups, monoids, and inverse semigroups and monoids, calculating their
Green's structure, ideals, size, elements, group of units, small generating
sets, testing membership, finding the inverses of a regular element,
factorizing elements over the generators, and so on. It is possible to test if
a semigroup satisfies a particular property, such as if it is regular, simple,
inverse, completely regular, and a large number of further properties. There
are methods for finding presentations for a semigroup, the congruences of a
semigroup, the maximal subsemigroups of a finite semigroup, smaller degree
partial permutation representations, and the character tables of inverse
semigroups. There are functions for producing pictures of the Green's
structure of a semigroup, and for drawing graphical representations of certain
types of elements. \mbox{}}\\[1cm]
{\small 
\section*{Copyright}
\logpage{[ 0, 0, 2 ]}
 {\copyright} by J. D. Mitchell et al.

 \textsf{Semigroups} is free software; you can redistribute it and/or modify it, under the terms of
the GNU General Public License, version 3 of the License, or (at your option)
any later, version. \mbox{}}\\[1cm]
{\small 
\section*{Acknowledgements}
\logpage{[ 0, 0, 3 ]}
 The authors of the \textsf{Semigroups} package would like to thank: 
\begin{description}
\item[{ Manuel Delgado }]  who contributed to the function \texttt{DotString} (\ref{DotString}). 
\item[{ Casey Donoven and Rhiannon Dougall }]  for their contribution to the development of the algorithms for maximal
subsemigroups and smaller degree partial permutation representations. 
\item[{ James East }]  who contributed to the part of the package relating to bipartitions. We also
thank the University of Western Sydney for their support of the development of
this part of the package. 
\item[{ Zak Mesyan }]  who contributed to the code for graph inverse semigroups; see Section \ref{Graph inverse semigroups}. 
\item[{ Yann P{\a'e}resse and Yanhui Wang }]  who contributed to the attribute \texttt{MunnSemigroup} (\ref{MunnSemigroup}). 
\item[{ Jhevon Smith and Ben Steinberg }]  who contributed the function \texttt{CharacterTableOfInverseSemigroup} (\ref{CharacterTableOfInverseSemigroup}). 
\end{description}
 We would also like to acknowledge the support of: EPSRC grant number
GR/S/56085/01; the Carnegie Trust for the Universities of Scotland for funding
the PhD scholarships of Julius Jonu{\v s}as and Wilf A. Wilson when they
worked on this project; the Engineering and Physical Sciences Research Council
(EPSRC) for funding the PhD scholarships of F. Smith (EP/N509759/1) and M.
Young (EP/M506631/1) when they worked on this project. \mbox{}}\\[1cm]
\newpage

\def\contentsname{Contents\logpage{[ 0, 0, 4 ]}}

\tableofcontents
\newpage

 
\chapter{\textcolor{Chapter }{ The \textsf{Semigroups} package }}\label{The Semigroups package}
\logpage{[ 1, 0, 0 ]}
\hyperdef{L}{X7D8D6DB37A0326BE}{}
{
  \index{Semigroups@\textsf{Semigroups} package overview} 
\section{\textcolor{Chapter }{ Introduction }}\label{Introduction}
\logpage{[ 1, 1, 0 ]}
\hyperdef{L}{X7DFB63A97E67C0A1}{}
{
  This is the manual for the \textsf{Semigroups} package for \textsf{GAP} version 5.6.1. \textsf{Semigroups} 5.6.1 is a distant descendant of the \href{ http://schmidt.nuigalway.ie/monoid/index.html} {Monoid package for GAP 3} by Goetz Pfeiffer, Steve A. Linton, Edmund F. Robertson, and Nik Ruskuc.

 From Version 3.0.0, \textsf{Semigroups} includes a copy of the \href{https://libsemigroups.readthedocs.io/en/latest/} {libsemigroups}  C++ library which contains implementations of the
Froidure\texttt{\symbol{45}}Pin, Todd\texttt{\symbol{45}}Coxeter, and
Knuth\texttt{\symbol{45}}Bendix algorithms (among others) that \textsf{Semigroups} utilises. }

 If you find a bug or an issue with the package, please visit the \href{ https://github.com/semigroups/Semigroups/issues } {issue tracker}. 
\section{\textcolor{Chapter }{ Overview }}\label{Overview}
\logpage{[ 1, 2, 0 ]}
\hyperdef{L}{X8389AD927B74BA4A}{}
{
  This manual is organised as follows: 
\begin{description}
\item[{Part I: elements}]  the different types of elements that are introduced in \textsf{Semigroups} are described in Chapters \ref{Bipartitions and blocks}, \ref{Partitioned binary relations (PBRs)}, and \ref{Matrices over semirings}. These include \texttt{Bipartition} (\ref{Bipartition}), \texttt{PBR} (\ref{PBR}), and \texttt{Matrix} (\ref{Matrix:for a filter and a matrix}), which supplement those already defined in the \textsf{GAP} library, such as \texttt{Transformation} (\textbf{Reference: Transformation for an image list}) or \texttt{PartialPerm} (\textbf{Reference: PartialPerm for a domain and image}). 
\item[{Part II: semigroups and monoids defined by generating sets}]  functions and operations for creating semigroups and monoids defined by
generating sets (of the type described in Part I) are described in Chapter \ref{Semigroups and monoids defined by generating sets}. 
\item[{Part III: standard examples and constructions}]  standard examples of semigroups, such as \texttt{FullBooleanMatMonoid} (\ref{FullBooleanMatMonoid}) or \texttt{UniformBlockBijectionMonoid} (\ref{UniformBlockBijectionMonoid}), are described in Chapter \ref{Standard examples}, and standard constructions, such as \texttt{DirectProduct} (\ref{DirectProduct}) are given in Chapter \ref{Standard constructions}. 
\item[{Part IV: the structure of a semigroup or monoid}]  the functionality for determining various structural properties of a given
semigroup or monoid are described in Chapters \ref{Ideals}, \ref{Green's relations}, \ref{Attributes and operations for semigroups}, and \ref{Properties of semigroups}. 
\item[{Part V: congruences, quotients, and homomorphisms}]  methods for creating and manipulating congruences and homomorphisms are
described by Chapters \ref{Congruences} and \ref{Semigroup homomorphisms}. 
\item[{Part VI: finitely presented semigroups and monoids}]  methods for finitely presented semigroups and monoids, in particular, for
Tietze transformations can be found in Chapters \ref{Finitely presented semigroups and Tietze transformations}. 
\item[{Part VII: utilities and helper functions}]  functions for reading and writing semigroups and their elements, and for
visualising semigroups, and some of their elements, can be found in Chapters \ref{Visualising semigroups and elements} and \ref{IO}. 
\end{description}
 }

 }

 
\chapter{\textcolor{Chapter }{Installing \textsf{Semigroups}}}\label{Installing Semigroups}
\logpage{[ 2, 0, 0 ]}
\hyperdef{L}{X82398F3785F63754}{}
{
  
\section{\textcolor{Chapter }{For those in a hurry}}\label{For those in a hurry}
\logpage{[ 2, 1, 0 ]}
\hyperdef{L}{X7DA3059C79842BF3}{}
{
  In this section we give a brief description of how to start using \textsf{Semigroups}.

 It is assumed that you have a working copy of \textsf{GAP} with version number 4.12.1 or higher. The most
up\texttt{\symbol{45}}to\texttt{\symbol{45}}date version of \textsf{GAP} and instructions on how to install it can be obtained from the main \textsf{GAP} webpage \href{https://www.gap-system.org} {\texttt{https://www.gap\texttt{\symbol{45}}system.org}}.

 The following is a summary of the steps that should lead to a successful
installation of \textsf{Semigroups}: 
\begin{itemize}
\item  ensure that the \textsf{datastructures} package version 0.3.0 or higher is available. \textsf{datastructures} must be compiled before \textsf{Semigroups} can be loaded. 
\item  ensure that the \textsf{digraphs} package version 1.6.2 or higher is available. \textsf{digraphs} must be compiled before \textsf{Semigroups} can be loaded. 
\item  ensure that the \textsf{genss} package version 1.6.5 or higher is available. 
\item  ensure that the \href{https://gap-packages.github.io/images/} {images}  package version 1.3.1 or higher is available. 
\item  ensure that the \href{https://gap-packages.github.io/io} {IO}  package version 4.5.1 or higher is available. \href{https://gap-packages.github.io/io} {IO}  must be compiled before \textsf{Semigroups} can be loaded. 
\item  ensure that the \textsf{orb} package version 4.8.2 or higher is available. \textsf{orb} and \textsf{Semigroups} both perform better if \textsf{orb} is compiled. 
\item  download the package archive \texttt{semigroups\texttt{\symbol{45}}5.6.1.tar.gz} from \href{https://semigroups.github.io/Semigroups} {the Semigroups package webpage}. 
\item  unzip and untar the file, this should create a directory called \texttt{semigroups\texttt{\symbol{45}}5.6.1}. 
\item  locate your \textsf{GAP} directory, the one containing the directories \texttt{lib}, \texttt{doc} and so on. Move the directory \texttt{semigroups\texttt{\symbol{45}}5.6.1} into the \texttt{pkg} subdirectory of your \textsf{GAP} directory. 
\item  from version 3.0.0, it is necessary to compile the \textsf{Semigroups} package. \textsf{Semigroups} uses the \href{ https://libsemigroups.github.io/libsemigroups/ } {libsemigroups} C++ library, which requires a compiler implementing the C++14 standard. Inside
the \texttt{pkg/semigroups\texttt{\symbol{45}}5.6.1} directory, in your terminal type 
\begin{Verbatim}[commandchars=!@|,fontsize=\small,frame=single,label=]
  ./configure && make
          
\end{Verbatim}
 Further information about this step can be found in Section \ref{Compiling the kernel module}. 
\item  start \textsf{GAP} in the usual way (i.e. type \texttt{gap} at the command line). 
\item  type \texttt{LoadPackage("semigroups");} 
\end{itemize}
 \emph{\textsc{Please note that} from version 3.0.0: \textsf{Semigroups} can only be loaded if it has been compiled.}

 If you want to check that the package is working correctly, you should run
some of the tests described in Section \ref{Testing your installation}. }

   
\section{\textcolor{Chapter }{Compiling the kernel module}}\label{Compiling the kernel module}
\logpage{[ 2, 2, 0 ]}
\hyperdef{L}{X849F6196875A6DF5}{}
{
  As of version 3.0.0, the \textsf{Semigroups} package has a kernel module written in C++ and this must be compiled. The
kernel module contains the interface to the C++ library \href{https://libsemigroups.readthedocs.io/en/latest/} {libsemigroups} . It is not possible to use the \textsf{Semigroups} package without compiling it. 

 To compile the kernel component inside the \texttt{pkg/semigroups\texttt{\symbol{45}}5.6.1} directory, type 
\begin{Verbatim}[commandchars=!@|,fontsize=\small,frame=single,label=]
  ./configure && make
      
\end{Verbatim}
 

 If you are using GCC to compile \textsf{Semigroups}, then version 5.0 or higher is required. Trying to compile \textsf{Semigroups} with an earlier version of GCC will result in an error at compile time. \textsf{Semigroups} only supports GCC version 5.0 or higher, and clang version 5.0 or higher.

 If you installed the package in a \texttt{pkg} directory other than the standard \texttt{pkg} directory in your \textsf{GAP} installation, then you have to do two things. Firstly during compilation you
have to use the option \texttt{\texttt{\symbol{45}}\texttt{\symbol{45}}with\texttt{\symbol{45}}gaproot=PATH} of the \texttt{configure} script where \texttt{PATH} is a path to the main \textsf{GAP} root directory (if not given the default \texttt{../..} is assumed).

 If you installed \textsf{GAP} on several architectures, you must execute the configure/make step for each of
the architectures. You can either do this immediately after configuring and
compiling \textsf{GAP} itself on this architecture, or alternatively set the environment variable \texttt{CONFIGNAME} to the name of the configuration you used when compiling \textsf{GAP} before running \texttt{./configure}. Note however that your compiler choice and flags (environment variables \texttt{CC} and \texttt{CFLAGS}) need to be chosen to match the setup of the original \textsf{GAP} compilation. For example you have to specify 32\texttt{\symbol{45}}bit or
64\texttt{\symbol{45}}bit mode correctly! }

   
\section{\textcolor{Chapter }{Rebuilding the documentation}}\label{Rebuilding the documentation}
\logpage{[ 2, 3, 0 ]}
\hyperdef{L}{X857CBE5484CF703A}{}
{
  The \textsf{Semigroups} package comes complete with pdf, html, and text versions of the documentation.
However, you might find it necessary, at some point, to rebuild the
documentation. To rebuild the documentation the \textsf{GAPDoc} and \href{https://gap-packages.github.io/AutoDoc} {AutoDoc}  packages are required. To rebuild the documentation type: 
\begin{Verbatim}[commandchars=!@|,fontsize=\small,frame=single,label=]
  gap makedoc.g
      
\end{Verbatim}
 when you're inside the \texttt{pkg/semigroups\texttt{\symbol{45}}5.6.1} directory. }

   
\section{\textcolor{Chapter }{Testing your installation}}\label{Testing your installation}
\logpage{[ 2, 4, 0 ]}
\hyperdef{L}{X7862D3F37C5BBDEF}{}
{
  In this section we describe how to test that \textsf{Semigroups} is working as intended. To quickly test that \textsf{Semigroups} is installed correctly use \texttt{SemigroupsTestInstall} (\ref{SemigroupsTestInstall}). For more extensive tests use \texttt{SemigroupsTestStandard} (\ref{SemigroupsTestStandard}). Finally, for lengthy benchmarking tests use \texttt{SemigroupsTestExtreme} (\ref{SemigroupsTestExtreme}). 

 If something goes wrong, then please review the instructions in Section \ref{For those in a hurry} and ensure that \textsf{Semigroups} has been properly installed. If you continue having problems, please use the \href{https://github.com/semigroups/Semigroups/issues} {issue tracker} to report the issues you are having. 

\subsection{\textcolor{Chapter }{SemigroupsTestInstall}}
\logpage{[ 2, 4, 1 ]}\nobreak
\hyperdef{L}{X80F85B577A3DFCF9}{}
{\noindent\textcolor{FuncColor}{$\triangleright$\enspace\texttt{SemigroupsTestInstall({\mdseries\slshape })\index{SemigroupsTestInstall@\texttt{SemigroupsTestInstall}}
\label{SemigroupsTestInstall}
}\hfill{\scriptsize (function)}}\\
\textbf{\indent Returns:\ }
\texttt{true} or \texttt{false}.



 This function should be called with no argument to test your installation of \textsf{Semigroups} is working correctly. These tests should take no more than a few seconds to
complete. To more comprehensively test that \textsf{Semigroups} is installed correctly use \texttt{SemigroupsTestStandard} (\ref{SemigroupsTestStandard}). }

 

\subsection{\textcolor{Chapter }{SemigroupsTestStandard}}
\logpage{[ 2, 4, 2 ]}\nobreak
\hyperdef{L}{X7C2D57708006AB63}{}
{\noindent\textcolor{FuncColor}{$\triangleright$\enspace\texttt{SemigroupsTestStandard({\mdseries\slshape })\index{SemigroupsTestStandard@\texttt{SemigroupsTestStandard}}
\label{SemigroupsTestStandard}
}\hfill{\scriptsize (function)}}\\
\textbf{\indent Returns:\ }
A list indicating which tests passed and failed and the time take to run each
file.



 This function should be called with no argument to comprehensively test that \textsf{Semigroups} is working correctly. These tests should take no more than a few minutes to
complete. To quickly test that \textsf{Semigroups} is installed correctly use \texttt{SemigroupsTestInstall} (\ref{SemigroupsTestInstall}).

 Each test file is run twice, once when the methods for \texttt{IsActingSemigroup} (\ref{IsActingSemigroup}) are enabled and once when they are disabled. }

 

\subsection{\textcolor{Chapter }{SemigroupsTestExtreme}}
\logpage{[ 2, 4, 3 ]}\nobreak
\hyperdef{L}{X7ED2F9C784B554D8}{}
{\noindent\textcolor{FuncColor}{$\triangleright$\enspace\texttt{SemigroupsTestExtreme({\mdseries\slshape })\index{SemigroupsTestExtreme@\texttt{SemigroupsTestExtreme}}
\label{SemigroupsTestExtreme}
}\hfill{\scriptsize (function)}}\\
\textbf{\indent Returns:\ }
A list indicating which tests passed and failed and the time take to run each
file.



 This function should be called with no argument to run some
long\texttt{\symbol{45}}running tests, which could be used to benchmark \textsf{Semigroups} or test your hardware. These tests should take no more than around half an
hour to complete. To quickly test that \textsf{Semigroups} is installed correctly use \texttt{SemigroupsTestInstall} (\ref{SemigroupsTestInstall}), or to test all aspects of the package use \texttt{SemigroupsTestStandard} (\ref{SemigroupsTestStandard}).

 Each test file is run twice, once when the methods for semigroups satisfying \texttt{IsActingSemigroup} (\ref{IsActingSemigroup}) are enabled and once when they are disabled. }

 

\subsection{\textcolor{Chapter }{SemigroupsTestAll}}
\logpage{[ 2, 4, 4 ]}\nobreak
\hyperdef{L}{X8544F4BD79F0BF3C}{}
{\noindent\textcolor{FuncColor}{$\triangleright$\enspace\texttt{SemigroupsTestAll({\mdseries\slshape })\index{SemigroupsTestAll@\texttt{SemigroupsTestAll}}
\label{SemigroupsTestAll}
}\hfill{\scriptsize (function)}}\\
\textbf{\indent Returns:\ }
\texttt{true} or \texttt{false}.



 This function should be called with no argument to compile the \textsf{Semigroups} package's documentation, run the standard suite of tests, and run all the
examples from the documentation to ensure that their output is correct. The
value returned is \texttt{true} if all the tests succeed, and \texttt{false} otherwise. The whole process should take no more than a few minutes. 

 See \texttt{SemigroupsTestStandard} (\ref{SemigroupsTestStandard}). }

 }

   
\section{\textcolor{Chapter }{More information during a computation}}\logpage{[ 2, 5, 0 ]}
\hyperdef{L}{X798CBC46800AB80F}{}
{
  

\subsection{\textcolor{Chapter }{InfoSemigroups}}
\logpage{[ 2, 5, 1 ]}\nobreak
\hyperdef{L}{X85CD4E6C82BECAF3}{}
{\noindent\textcolor{FuncColor}{$\triangleright$\enspace\texttt{InfoSemigroups\index{InfoSemigroups@\texttt{InfoSemigroups}}
\label{InfoSemigroups}
}\hfill{\scriptsize (info class)}}\\


 \texttt{InfoSemigroups} is the info class of the \textsf{Semigroups} package. The info level is initially set to 0 and no info messages are
displayed. To increase the amount of information displayed during a
computation increase the info level to 2 or 3. To stop all info messages from
being displayed, set the info level to 0. See also  (\textbf{Reference: Info Functions}) and \texttt{SetInfoLevel} (\textbf{Reference: InfoLevel}). }

 }

   }

 
\chapter{\textcolor{Chapter }{ Bipartitions and blocks }}\label{Bipartitions and blocks}
\logpage{[ 3, 0, 0 ]}
\hyperdef{L}{X7C18DB427C9C0917}{}
{
  In this chapter we describe the functions in \textsf{Semigroups} for creating and manipulating bipartitions and semigroups of bipartitions. We
begin by describing what these objects are. 

 A \emph{partition} of a set $X$  is a set of pairwise disjoint non\texttt{\symbol{45}}empty subsets of $X$  whose union is $X$ . A partition of $X$  is the collection of equivalence classes of an equivalence relation on $X$ , and vice versa. 

 Let $n\in$$\mathbb{N}$ , let $\mathbf{n}$ $ = \{1, 2, \ldots, n\}$, and let $-$$\mathbf{n}$ $ = \{-1, -2, \ldots, -n\}$. 

 The \emph{partition monoid} of degree $n$  is the set of all partitions of $\mathbf{n}$ $\cup$\texttt{\symbol{45}}$\mathbf{n}$  with a multiplication we describe below. To avoid conflict with other uses of
the word "partition" in \textsf{GAP}, and to reflect their definition, we have opted to refer to the elements of
the partition monoid as \emph{bipartitions} of degree $n$ ; we will do so from this point on. 

 Let $x$  be any bipartition of degree $n$ . Then $x$  is a set of pairwise disjoint non\texttt{\symbol{45}}empty subsets of $\mathbf{n}$ $\cup$\texttt{\symbol{45}}$\mathbf{n}$  whose union is $\mathbf{n}$ $\cup$\texttt{\symbol{45}}$\mathbf{n}$ ; these subsets are called the \emph{blocks} of $x$ . A block containing elements of both $\mathbf{n}$  and \texttt{\symbol{45}}$\mathbf{n}$  is called a \emph{transverse block}. If $i$ , $j$ $\in$$\mathbf{n}$ $\cup$\texttt{\symbol{45}}$\mathbf{n}$  belong to the same block of a bipartition $x$ , then we write ($i$ , $j$ )$\in$$x$ . 

 Let $x$  and $y$  be bipartitions of degree $n$ . Their product $x$ $y$  can be described as follows. Define $\mathbf{n}$ '$ = \{1', 2', \ldots, n'\}$. From $x$ , create a partition $x$ ' of the set $\mathbf{n}$ $\cup$$\mathbf{n}$ ' by replacing each negative point \texttt{\symbol{45}}$i$  in a block of $x$  by the point $i$ ', and create from $y$  a partition $y$ ' of the set $\mathbf{n}$ '$\cup$\texttt{\symbol{45}}$\mathbf{n}$  by replacing each positive point $i$  in a block of $y$  by the point $i$ '. Then define a relation on the set $\mathbf{n}$ $\cup$$\mathbf{n}$ '$\cup$\texttt{\symbol{45}}$\mathbf{n}$ , where $i$  and $j$  are related if they are related in either $x$ ' or $y$ ', and let $p$  be the transitive closure of this relation. Finally, define $x$ $y$  to be the bipartition of degree $n$  defined by the restriction of the equivalence relation $p$  to the set $\mathbf{n}$ $\cup$\texttt{\symbol{45}}$\mathbf{n}$ . 

 Equivalently, the product $x$ $y$  is defined to be the bipartition where $i$ ,$j$ $\in$$\mathbf{n}$ $\cup$\texttt{\symbol{45}}$\mathbf{n}$  (we assume without loss of generality that $i$ $\geq$$j$ ) belong to the same block of $x$ $y$  if either: 
\begin{itemize}
\item  $i$ \texttt{=}$j$ , 
\item  $i$ , $j$  $\in$ $\mathbf{n}$  and $($$i$ ,$j$ $)$$\in$ $x$ , or 
\item  $i$ , $j$  $\in$ \texttt{\symbol{45}}$\mathbf{n}$  and $($$i$ ,$j$ $)$$\in$ $y$ ; 
\end{itemize}
 or there exists $r\in$$\mathbb{N}$  and  $k(1), k(2),\ldots, k(r)\in \mathbf{n}$  , and one of the following holds: 
\begin{itemize}
\item   $r=2s-1$ for some $s\geq 1$  , $i$ $\in$$\mathbf{n}$ , $j$ $\in$ \texttt{\symbol{45}}$\mathbf{n}$  and  
\[(i,-k(1))\in x,\ (k(1),k(2))\in y,\ (-k(2),-k(3))\in x,\ \ldots,\qquad\]
 
\[\qquad\ldots,\ (-k(2s-2),-k(2s-1))\in x,\ (k(2s-1),j)\in y;\]
   
\item   $r=2s$ for some $s\geq 1$  , and either $i$ ,$j$ $\in$$\mathbf{n}$ , and  
\[(i,-k(1))\in x,\ (k(1),k(2))\in y,\ (-k(2),-k(3))\in x,\ \ldots, (k(2s-1),
k(2s))\in y,\ (-k(2s), j)\in x,\]
   or $i$ ,$j$ $\in$\texttt{\symbol{45}}$\mathbf{n}$ , and  
\[(i,k(1))\in y,\ (-k(1),-k(2))\in x,\ (k(2),k(3))\in y,\ \ldots, (-k(2s-1),
-k(2s))\in x,\ (k(2s), j)\in y.\]
   
\end{itemize}
 This multiplication can be shown to be associative, and so the collection of
all bipartitions of any particular degree is a monoid; the identity element of
the partition monoid of degree $n$ is the bipartition $\left\{\{i,-i\}:i\in\mathbf{n}\right\}$.  A bipartition is a unit if and only if each block is of the form $\{$$i$ ,\texttt{\symbol{45}}$j$ $\}$ for some $i$ , $j$ $\in$$\mathbf{n}$ . Hence the group of units is isomorphic to the symmetric group on $\mathbf{n}$ . 

 Let $x$  be a bipartition of degree $n$ . Then we define $x$ $^*$ to be the bipartition obtained from $x$  by replacing $i$  by \texttt{\symbol{45}}$i$  and \texttt{\symbol{45}}$i$  by $i$  in every block of $x$  for all $i$ $\in$$\mathbf{n}$ . It is routine to verify that if $x$  and $y$  are arbitrary bipartitions of equal degree, then  
\[ (x^*)^*=x,\quad xx^*x=x,\quad x^*xx^*=x^*,\quad (xy)^*=y^*x^*. \]
   In this way, the partition monoid is a \emph{regular *\texttt{\symbol{45}}semigroup}. 

 A bipartition $x$  of degree $n$  is called \emph{planar} if there do not exist distinct blocks $A, U \in$ $x$ , along with $a, b \in A$ and $u, v \in U$, such that $a < u < b < v$. Define $p$  to be the bipartition of degree $n$  with blocks  $\left\{\{i, -(i+1)\}:i\in\{1,\ldots,n-1\right\}\}$ and $\{n,-1\}$  . Note that $p$  is a unit. A bipartition $x$  of degree $n$  is called \emph{annular} if  $x = p^{i} y p^{j}$   for some planar bipartition $y$  of degree $n$ , and some integers $i$  and $j$ . 

 From a graphical perspective, as on Page 873 in \cite{Halverson2005PartitionAlgebras}, a bipartition of degree $n$  is planar if it can be represented as a graph without edges crossing inside of
the rectangle formed by its vertices $\mathbf{n}$ $\cup$\texttt{\symbol{45}}$\mathbf{n}$ . Similarly, as shown in Figure 2 in \cite{auinger2012krohn}, a bipartition of degree $n$  is annular if it can be represented as a graph without edges crossing inside
an annulus. 
\section{\textcolor{Chapter }{The family and categories of bipartitions}}\logpage{[ 3, 1, 0 ]}
\hyperdef{L}{X7850845886902FBF}{}
{
 

\subsection{\textcolor{Chapter }{IsBipartition}}
\logpage{[ 3, 1, 1 ]}\nobreak
\hyperdef{L}{X80F11BEF856E7902}{}
{\noindent\textcolor{FuncColor}{$\triangleright$\enspace\texttt{IsBipartition({\mdseries\slshape obj})\index{IsBipartition@\texttt{IsBipartition}}
\label{IsBipartition}
}\hfill{\scriptsize (Category)}}\\
\textbf{\indent Returns:\ }
\texttt{true} or \texttt{false}.



 Every bipartition in \textsf{GAP} belongs to the category \texttt{IsBipartition}. Basic operations for bipartitions are \texttt{RightBlocks} (\ref{RightBlocks}), \texttt{LeftBlocks} (\ref{LeftBlocks}), \texttt{ExtRepOfObj} (\ref{ExtRepOfObj:for a bipartition}), \texttt{LeftProjection} (\ref{LeftProjection}), \texttt{RightProjection} (\ref{RightProjection}), \texttt{StarOp} (\ref{StarOp:for a bipartition}), \texttt{DegreeOfBipartition} (\ref{DegreeOfBipartition}), \texttt{RankOfBipartition} (\ref{RankOfBipartition}), multiplication of two bipartitions of equal degree is via \texttt{*}. }

 

\subsection{\textcolor{Chapter }{IsBipartitionCollection}}
\logpage{[ 3, 1, 2 ]}\nobreak
\hyperdef{L}{X82F5D10C85489832}{}
{\noindent\textcolor{FuncColor}{$\triangleright$\enspace\texttt{IsBipartitionCollection({\mdseries\slshape obj})\index{IsBipartitionCollection@\texttt{IsBipartitionCollection}}
\label{IsBipartitionCollection}
}\hfill{\scriptsize (Category)}}\\
\noindent\textcolor{FuncColor}{$\triangleright$\enspace\texttt{IsBipartitionCollColl({\mdseries\slshape obj})\index{IsBipartitionCollColl@\texttt{IsBipartitionCollColl}}
\label{IsBipartitionCollColl}
}\hfill{\scriptsize (Category)}}\\
\textbf{\indent Returns:\ }
\texttt{true} or \texttt{false}.



 Every collection of bipartitions belongs to the category \texttt{IsBipartitionCollection}. For example, bipartition semigroups belong to \texttt{IsBipartitionCollection}.

 Every collection of collections of bipartitions belongs to \texttt{IsBipartitionCollColl}. For example, a list of bipartition semigroups belongs to \texttt{IsBipartitionCollColl}. }

 }

   
\section{\textcolor{Chapter }{Creating bipartitions}}\label{creating-bipartitions}
\logpage{[ 3, 2, 0 ]}
\hyperdef{L}{X85D77073820C7E72}{}
{
  There are several ways of creating bipartitions in \textsf{GAP}, which are described in this section. The maximum degree of a bipartition is
set as $2 ^ 29 - 1$. In reality, it is unlikely to be possible to create bipartitions of degrees
as small as $2 ^ 24$ because they require too much memory. 

\subsection{\textcolor{Chapter }{Bipartition}}
\logpage{[ 3, 2, 1 ]}\nobreak
\hyperdef{L}{X7E052E6378A5B758}{}
{\noindent\textcolor{FuncColor}{$\triangleright$\enspace\texttt{Bipartition({\mdseries\slshape blocks})\index{Bipartition@\texttt{Bipartition}}
\label{Bipartition}
}\hfill{\scriptsize (function)}}\\
\textbf{\indent Returns:\ }
A bipartition.



 \texttt{Bipartition} returns the bipartition \texttt{x} with equivalence classes \mbox{\texttt{\mdseries\slshape blocks}}, which should be a list of duplicate\texttt{\symbol{45}}free lists whose
union is \texttt{[\texttt{\symbol{45}}n .. \texttt{\symbol{45}}1]} union \texttt{[1 .. n]} for some positive integer \texttt{n}. 

 \texttt{Bipartition} returns an error if the argument does not define a bipartition. 
\begin{Verbatim}[commandchars=!@|,fontsize=\small,frame=single,label=Example]
  !gapprompt@gap>| !gapinput@x := Bipartition([[1, -1], [2, 3, -3], [-2]]);|
  <bipartition: [ 1, -1 ], [ 2, 3, -3 ], [ -2 ]>
\end{Verbatim}
 }

 

\subsection{\textcolor{Chapter }{BipartitionByIntRep}}
\logpage{[ 3, 2, 2 ]}\nobreak
\hyperdef{L}{X846AA7568435D2CE}{}
{\noindent\textcolor{FuncColor}{$\triangleright$\enspace\texttt{BipartitionByIntRep({\mdseries\slshape list})\index{BipartitionByIntRep@\texttt{BipartitionByIntRep}}
\label{BipartitionByIntRep}
}\hfill{\scriptsize (operation)}}\\
\textbf{\indent Returns:\ }
A bipartition.



 It is possible to create a bipartition using its internal representation. The
argument \mbox{\texttt{\mdseries\slshape list}} must be a list of positive integers not greater than \texttt{n}, of length \texttt{2 * n}, and where \texttt{i} appears in the list only if \texttt{i\texttt{\symbol{45}}1} occurs earlier in the list. 

 For example, the internal representation of the bipartition with blocks 
\begin{Verbatim}[commandchars=!@|,fontsize=\small,frame=single,label=Example]
  [1, -1], [2, 3, -2], [-3]
\end{Verbatim}
 has internal representation 
\begin{Verbatim}[commandchars=!@|,fontsize=\small,frame=single,label=Example]
  [1, 2, 2, 1, 2, 3]
\end{Verbatim}
 The internal representation indicates that the number \texttt{1} is in class \texttt{1}, the number \texttt{2} is in class \texttt{2}, the number \texttt{3} is in class \texttt{2}, the number \texttt{\texttt{\symbol{45}}1} is in class \texttt{1}, the number \texttt{\texttt{\symbol{45}}2} is in class \texttt{2}, and \texttt{\texttt{\symbol{45}}3} is in class \texttt{3}. As another example, \texttt{[1, 3, 2, 1]} is not the internal representation of any bipartition since there is no \texttt{2} before the \texttt{3} in the second position.

 In its first form \texttt{BipartitionByIntRep} verifies that the argument \mbox{\texttt{\mdseries\slshape list}} is the internal representation of a bipartition. 

 See also \texttt{IntRepOfBipartition} (\ref{IntRepOfBipartition}). 
\begin{Verbatim}[commandchars=!@|,fontsize=\small,frame=single,label=Example]
  !gapprompt@gap>| !gapinput@BipartitionByIntRep([1, 2, 2, 1, 3, 4]);|
  <bipartition: [ 1, -1 ], [ 2, 3 ], [ -2 ], [ -3 ]>
\end{Verbatim}
 }

 

\subsection{\textcolor{Chapter }{IdentityBipartition}}
\logpage{[ 3, 2, 3 ]}\nobreak
\hyperdef{L}{X8379B0538101FBC8}{}
{\noindent\textcolor{FuncColor}{$\triangleright$\enspace\texttt{IdentityBipartition({\mdseries\slshape n})\index{IdentityBipartition@\texttt{IdentityBipartition}}
\label{IdentityBipartition}
}\hfill{\scriptsize (operation)}}\\
\textbf{\indent Returns:\ }
The identity bipartition.



 Returns the identity bipartition with degree \mbox{\texttt{\mdseries\slshape n}}. 
\begin{Verbatim}[commandchars=!@|,fontsize=\small,frame=single,label=Example]
  !gapprompt@gap>| !gapinput@IdentityBipartition(10);|
  <block bijection: [ 1, -1 ], [ 2, -2 ], [ 3, -3 ], [ 4, -4 ],
   [ 5, -5 ], [ 6, -6 ], [ 7, -7 ], [ 8, -8 ], [ 9, -9 ], [ 10, -10 ]>
\end{Verbatim}
 }

 

\subsection{\textcolor{Chapter }{LeftOne (for a bipartition)}}
\logpage{[ 3, 2, 4 ]}\nobreak
\hyperdef{L}{X824EDD4582AAA8C7}{}
{\noindent\textcolor{FuncColor}{$\triangleright$\enspace\texttt{LeftOne({\mdseries\slshape x})\index{LeftOne@\texttt{LeftOne}!for a bipartition}
\label{LeftOne:for a bipartition}
}\hfill{\scriptsize (attribute)}}\\
\noindent\textcolor{FuncColor}{$\triangleright$\enspace\texttt{LeftProjection({\mdseries\slshape x})\index{LeftProjection@\texttt{LeftProjection}}
\label{LeftProjection}
}\hfill{\scriptsize (attribute)}}\\
\textbf{\indent Returns:\ }
A bipartition.



 The \texttt{LeftProjection} of a bipartition \mbox{\texttt{\mdseries\slshape x}} is the bipartition \texttt{\mbox{\texttt{\mdseries\slshape x}} * Star(\mbox{\texttt{\mdseries\slshape x}})}. It is so\texttt{\symbol{45}}named, since the left and right blocks of the
left projection equal the left blocks of \mbox{\texttt{\mdseries\slshape x}}. 

 The left projection \texttt{e} of \mbox{\texttt{\mdseries\slshape x}} is also a bipartition with the property that \texttt{e * \mbox{\texttt{\mdseries\slshape x}} = \mbox{\texttt{\mdseries\slshape x}}}. \texttt{LeftOne} and \texttt{LeftProjection} are synonymous. 
\begin{Verbatim}[commandchars=!@|,fontsize=\small,frame=single,label=Example]
  !gapprompt@gap>| !gapinput@x := Bipartition([|
  !gapprompt@>| !gapinput@  [1, 4, -1, -2, -6], [2, 3, 5, -4], [6, -3], [-5]]);;|
  !gapprompt@gap>| !gapinput@LeftOne(x);|
  <block bijection: [ 1, 4, -1, -4 ], [ 2, 3, 5, -2, -3, -5 ],
   [ 6, -6 ]>
  !gapprompt@gap>| !gapinput@LeftBlocks(x);|
  <blocks: [ 1*, 4* ], [ 2*, 3*, 5* ], [ 6* ]>
  !gapprompt@gap>| !gapinput@RightBlocks(LeftOne(x));|
  <blocks: [ 1*, 4* ], [ 2*, 3*, 5* ], [ 6* ]>
  !gapprompt@gap>| !gapinput@LeftBlocks(LeftOne(x));|
  <blocks: [ 1*, 4* ], [ 2*, 3*, 5* ], [ 6* ]>
  !gapprompt@gap>| !gapinput@LeftOne(x) * x = x;|
  true
\end{Verbatim}
 }

 

\subsection{\textcolor{Chapter }{RightOne (for a bipartition)}}
\logpage{[ 3, 2, 5 ]}\nobreak
\hyperdef{L}{X790B71108070FAC2}{}
{\noindent\textcolor{FuncColor}{$\triangleright$\enspace\texttt{RightOne({\mdseries\slshape x})\index{RightOne@\texttt{RightOne}!for a bipartition}
\label{RightOne:for a bipartition}
}\hfill{\scriptsize (attribute)}}\\
\noindent\textcolor{FuncColor}{$\triangleright$\enspace\texttt{RightProjection({\mdseries\slshape x})\index{RightProjection@\texttt{RightProjection}}
\label{RightProjection}
}\hfill{\scriptsize (attribute)}}\\
\textbf{\indent Returns:\ }
A bipartition.



 The \texttt{RightProjection} of a bipartition \mbox{\texttt{\mdseries\slshape x}} is the bipartition \texttt{Star(\mbox{\texttt{\mdseries\slshape x}}) * \mbox{\texttt{\mdseries\slshape x}}}. It is so\texttt{\symbol{45}}named, since the left and right blocks of the
right projection equal the right blocks of \mbox{\texttt{\mdseries\slshape x}}. 

 The right projection \texttt{e} of \mbox{\texttt{\mdseries\slshape x}} is also a bipartition with the property that \texttt{\mbox{\texttt{\mdseries\slshape x}} * e = \mbox{\texttt{\mdseries\slshape x}}}. \texttt{RightOne} and \texttt{RightProjection} are synonymous. 
\begin{Verbatim}[commandchars=!@|,fontsize=\small,frame=single,label=Example]
  !gapprompt@gap>| !gapinput@x := Bipartition([[1, -1, -4], [2, -2, -3], [3, 4], [5, -5]]);;|
  !gapprompt@gap>| !gapinput@RightOne(x);|
  <block bijection: [ 1, 4, -1, -4 ], [ 2, 3, -2, -3 ], [ 5, -5 ]>
  !gapprompt@gap>| !gapinput@RightBlocks(RightOne(x));|
  <blocks: [ 1*, 4* ], [ 2*, 3* ], [ 5* ]>
  !gapprompt@gap>| !gapinput@LeftBlocks(RightOne(x));|
  <blocks: [ 1*, 4* ], [ 2*, 3* ], [ 5* ]>
  !gapprompt@gap>| !gapinput@RightBlocks(x);|
  <blocks: [ 1*, 4* ], [ 2*, 3* ], [ 5* ]>
  !gapprompt@gap>| !gapinput@x * RightOne(x) = x;|
  true
\end{Verbatim}
 }

 

\subsection{\textcolor{Chapter }{StarOp (for a bipartition)}}
\logpage{[ 3, 2, 6 ]}\nobreak
\hyperdef{L}{X7CE00E0C79F62745}{}
{\noindent\textcolor{FuncColor}{$\triangleright$\enspace\texttt{StarOp({\mdseries\slshape x})\index{StarOp@\texttt{StarOp}!for a bipartition}
\label{StarOp:for a bipartition}
}\hfill{\scriptsize (operation)}}\\
\noindent\textcolor{FuncColor}{$\triangleright$\enspace\texttt{Star({\mdseries\slshape x})\index{Star@\texttt{Star}!for a bipartition}
\label{Star:for a bipartition}
}\hfill{\scriptsize (attribute)}}\\
\textbf{\indent Returns:\ }
A bipartition.



 \texttt{StarOp} returns the unique bipartition \texttt{g} with the property that: \texttt{\mbox{\texttt{\mdseries\slshape x}} * g * \mbox{\texttt{\mdseries\slshape x}} = \mbox{\texttt{\mdseries\slshape x}}}, \texttt{RightBlocks(\mbox{\texttt{\mdseries\slshape x}}) = LeftBlocks(g)}, and \texttt{LeftBlocks(\mbox{\texttt{\mdseries\slshape x}}) = RightBlocks(g)}. The star \texttt{g} can be obtained from \mbox{\texttt{\mdseries\slshape x}} by changing the sign of every integer in the external representation of \mbox{\texttt{\mdseries\slshape x}}. 
\begin{Verbatim}[commandchars=!@|,fontsize=\small,frame=single,label=Example]
  !gapprompt@gap>| !gapinput@x := Bipartition([[1, -4], [2, 3, 4], [5], [-1], [-2, -3], [-5]]);|
  <bipartition: [ 1, -4 ], [ 2, 3, 4 ], [ 5 ], [ -1 ], [ -2, -3 ],
   [ -5 ]>
  !gapprompt@gap>| !gapinput@y := Star(x);|
  <bipartition: [ 1 ], [ 2, 3 ], [ 4, -1 ], [ 5 ], [ -2, -3, -4 ],
   [ -5 ]>
  !gapprompt@gap>| !gapinput@x * y * x = x;|
  true
  !gapprompt@gap>| !gapinput@LeftBlocks(x) = RightBlocks(y);|
  true
  !gapprompt@gap>| !gapinput@RightBlocks(x) = LeftBlocks(y);|
  true
\end{Verbatim}
 }

 

\subsection{\textcolor{Chapter }{TensorBipartitions (for a pair of bipartitions)}}
\logpage{[ 3, 2, 7 ]}\nobreak
\hyperdef{L}{X7F7D823084F79C7D}{}
{\noindent\textcolor{FuncColor}{$\triangleright$\enspace\texttt{TensorBipartitions({\mdseries\slshape x, y})\index{TensorBipartitions@\texttt{TensorBipartitions}!for a pair of bipartitions}
\label{TensorBipartitions:for a pair of bipartitions}
}\hfill{\scriptsize (operation)}}\\
\noindent\textcolor{FuncColor}{$\triangleright$\enspace\texttt{TensorBipartitions({\mdseries\slshape coll})\index{TensorBipartitions@\texttt{TensorBipartitions}!for a list of bipartitions}
\label{TensorBipartitions:for a list of bipartitions}
}\hfill{\scriptsize (operation)}}\\
\textbf{\indent Returns:\ }
A bipartition.



 This operation returns the tensor product of the bipartitions \mbox{\texttt{\mdseries\slshape x}} and \mbox{\texttt{\mdseries\slshape y}}, or of all of the bipartitions in \mbox{\texttt{\mdseries\slshape coll}}. Roughly speaking, the \emph{tensor product} of bipartitions \mbox{\texttt{\mdseries\slshape x}} and \mbox{\texttt{\mdseries\slshape y}} is obtained by writing the bipartitions next to each other and reindexing the
second one \mbox{\texttt{\mdseries\slshape y}}. More precisely, if \mbox{\texttt{\mdseries\slshape x}} has degree \texttt{m} and \mbox{\texttt{\mdseries\slshape y}} has degree \texttt{n}, then the tensor product is the partition of \texttt{[\texttt{\symbol{45}}m \texttt{\symbol{45}} n .. \texttt{\symbol{45}}1, 1 .. m
+ n]} given by: 
\begin{itemize}
\item  the blocks containing values in \texttt{[\texttt{\symbol{45}}m .. \texttt{\symbol{45}}1, 1 .. m]} are precisely the blocks of \mbox{\texttt{\mdseries\slshape x}}; 
\item  the blocks containing values in \texttt{[\texttt{\symbol{45}}m \texttt{\symbol{45}} n .. \texttt{\symbol{45}}m
\texttt{\symbol{45}} 1, m + 1 ... m + n]} are obtained from the blocks of \mbox{\texttt{\mdseries\slshape y}} by adding \texttt{m} to the positive, and subtracting \texttt{m} from the negative values. 
\end{itemize}
 If the argument is a list of bipartitions \mbox{\texttt{\mdseries\slshape coll}}, then the tensor product of the bipartitions in the list (from left to right)
is returned. An error is given if \mbox{\texttt{\mdseries\slshape coll}} is empty.

 See also: \texttt{IrreducibleComponentsOfBipartition} (\ref{IrreducibleComponentsOfBipartition}) and \texttt{IsIrreducibleBipartition} (\ref{IsIrreducibleBipartition}).  
\begin{Verbatim}[commandchars=!@|,fontsize=\small,frame=single,label=Example]
  !gapprompt@gap>| !gapinput@TensorBipartitions(|
  !gapprompt@>| !gapinput@Bipartition([[1, 2, -2], [-1]]),|
  !gapprompt@>| !gapinput@Bipartition([[1, 2, 3, -2], [-1], [-3]]));|
  <bipartition: [ 1, 2, -2 ], [ 3, 4, 5, -4 ], [ -1 ], [ -3 ], [ -5 ]>
\end{Verbatim}
 }

 

\subsection{\textcolor{Chapter }{RandomBipartition}}
\logpage{[ 3, 2, 8 ]}\nobreak
\hyperdef{L}{X8077265981409CCB}{}
{\noindent\textcolor{FuncColor}{$\triangleright$\enspace\texttt{RandomBipartition({\mdseries\slshape [rs, ]n})\index{RandomBipartition@\texttt{RandomBipartition}}
\label{RandomBipartition}
}\hfill{\scriptsize (operation)}}\\
\noindent\textcolor{FuncColor}{$\triangleright$\enspace\texttt{RandomBlockBijection({\mdseries\slshape [rs, ]n})\index{RandomBlockBijection@\texttt{RandomBlockBijection}}
\label{RandomBlockBijection}
}\hfill{\scriptsize (operation)}}\\
\textbf{\indent Returns:\ }
A bipartition.



 If \mbox{\texttt{\mdseries\slshape n}} is a positive integer, then \texttt{RandomBipartition} returns a random bipartition of degree \mbox{\texttt{\mdseries\slshape n}}, and \texttt{RandomBlockBijection} returns a random block bijection of degree \mbox{\texttt{\mdseries\slshape n}}. 

 If the optional first argument \mbox{\texttt{\mdseries\slshape rs}} is a random source, then this is used to generate the bipartition returned by \texttt{RandomBipartition} and \texttt{RandomBlockBijection}. 

 Note that neither of these functions has a uniform distribution. 
\begin{Verbatim}[commandchars=!@|,fontsize=\small,frame=single,label=Example]
  !gapprompt@gap>| !gapinput@x := RandomBipartition(6);|
  <bipartition: [ 1, 2, 3, 4 ], [ 5 ], [ 6, -2, -3, -4 ], [ -1, -5 ], [ -6 ]>
  !gapprompt@gap>| !gapinput@x := RandomBlockBijection(4);|
  <block bijection: [ 1, 4, -2 ], [ 2, -4 ], [ 3, -1, -3 ]>
\end{Verbatim}
 }

 }

   
\section{\textcolor{Chapter }{Changing the representation of a bipartition}}\label{changing-rep-bipartitions}
\logpage{[ 3, 3, 0 ]}
\hyperdef{L}{X7C2C44D281A0D2C9}{}
{
  It is possible that a bipartition can be represented as another type of
object, or that another type of \textsf{GAP} object can be represented as a bipartition. In this section, we describe the
functions in the \textsf{Semigroups} package for changing the representation of bipartition, or for changing the
representation of another type of object to that of a bipartition.

 The operations \texttt{AsPermutation} (\ref{AsPermutation:for a bipartition}), \texttt{AsPartialPerm} (\ref{AsPartialPerm:for a bipartition}), \texttt{AsTransformation} (\ref{AsTransformation:for a bipartition}) can be used to convert bipartitions into permutations, partial permutations,
or transformations where appropriate. 

\subsection{\textcolor{Chapter }{AsBipartition}}
\logpage{[ 3, 3, 1 ]}\nobreak
\hyperdef{L}{X855126D98583C181}{}
{\noindent\textcolor{FuncColor}{$\triangleright$\enspace\texttt{AsBipartition({\mdseries\slshape x[, n]})\index{AsBipartition@\texttt{AsBipartition}}
\label{AsBipartition}
}\hfill{\scriptsize (operation)}}\\
\textbf{\indent Returns:\ }
A bipartition.



 \texttt{AsBipartition} returns the bipartition, permutation, transformation, or partial permutation \mbox{\texttt{\mdseries\slshape x}}, as a bipartition of degree \mbox{\texttt{\mdseries\slshape n}}.

 There are several possible arguments for \texttt{AsBipartition}: 
\begin{description}
\item[{permutations}]  If \mbox{\texttt{\mdseries\slshape x}} is a permutation and \mbox{\texttt{\mdseries\slshape n}} is a positive integer, then \texttt{AsBipartition(\mbox{\texttt{\mdseries\slshape x}}, \mbox{\texttt{\mdseries\slshape n}})} returns the bipartition on \texttt{[1 .. \mbox{\texttt{\mdseries\slshape n}}]} with classes \texttt{[i, i \texttt{\symbol{94}} \mbox{\texttt{\mdseries\slshape x}}]} for all \texttt{i = 1 .. n}.

 If no positive integer \mbox{\texttt{\mdseries\slshape n}} is specified, then the largest moved point of \mbox{\texttt{\mdseries\slshape x}} is used as the value for \mbox{\texttt{\mdseries\slshape n}}; see \texttt{LargestMovedPoint} (\textbf{Reference: LargestMovedPoint for a permutation}). 
\item[{transformations}]  If \mbox{\texttt{\mdseries\slshape x}} is a transformation and \mbox{\texttt{\mdseries\slshape n}} is a positive integer such that \mbox{\texttt{\mdseries\slshape x}} is a transformation of \texttt{[1 .. \mbox{\texttt{\mdseries\slshape n}}]}, then \texttt{AsTransformation} returns the bipartition with classes $(i)f ^ {-1}\cup \{i\}$ for all \texttt{i} in the image of \mbox{\texttt{\mdseries\slshape x}}.

 If the positive integer \mbox{\texttt{\mdseries\slshape n}} is not specified, then the degree of \mbox{\texttt{\mdseries\slshape x}} is used as the value for \mbox{\texttt{\mdseries\slshape n}}. 
\item[{partial permutations}]  If \mbox{\texttt{\mdseries\slshape x}} is a partial permutation and \mbox{\texttt{\mdseries\slshape n}} is a positive integer, then \texttt{AsBipartition} returns the bipartition with classes \texttt{[i, i \texttt{\symbol{94}} \mbox{\texttt{\mdseries\slshape x}}]} for \texttt{i} in \texttt{[1 .. \mbox{\texttt{\mdseries\slshape n}}]}. Thus the degree of the returned bipartition is the maximum of \mbox{\texttt{\mdseries\slshape n}} and the values \texttt{i \texttt{\symbol{94}} \mbox{\texttt{\mdseries\slshape x}}} where \texttt{i} in \texttt{[1 .. \mbox{\texttt{\mdseries\slshape n}}]}.

 If the optional argument \mbox{\texttt{\mdseries\slshape n}} is not present, then the default value of the maximum of the largest moved
point and the largest image of a moved point of \mbox{\texttt{\mdseries\slshape x}} plus \texttt{1} is used. 
\item[{bipartitions}]  If \mbox{\texttt{\mdseries\slshape x}} is a bipartition and \mbox{\texttt{\mdseries\slshape n}} is a non\texttt{\symbol{45}}negative integer, then \texttt{AsBipartition} returns a bipartition corresponding to \mbox{\texttt{\mdseries\slshape x}} with degree \mbox{\texttt{\mdseries\slshape n}}. 

 If \mbox{\texttt{\mdseries\slshape n}} equals the degree of \mbox{\texttt{\mdseries\slshape x}}, then \mbox{\texttt{\mdseries\slshape x}} is returned. If \mbox{\texttt{\mdseries\slshape n}} is less than the degree of \mbox{\texttt{\mdseries\slshape x}}, then this function returns the bipartition obtained from \mbox{\texttt{\mdseries\slshape x}} by removing the values exceeding \mbox{\texttt{\mdseries\slshape n}} or less than \mbox{\texttt{\mdseries\slshape \texttt{\symbol{45}}n}} from the blocks of \mbox{\texttt{\mdseries\slshape x}}. If \mbox{\texttt{\mdseries\slshape n}} is greater than the degree of \mbox{\texttt{\mdseries\slshape x}}, then this function returns the bipartition with the same blocks as \mbox{\texttt{\mdseries\slshape x}} and the singleton blocks \texttt{i} and \texttt{\texttt{\symbol{45}}i} for all \texttt{i} greater than the degree of \mbox{\texttt{\mdseries\slshape x}} 
\item[{pbrs}]  If \mbox{\texttt{\mdseries\slshape x}} is a pbr satisfying \texttt{IsBipartitionPBR} (\ref{IsBipartitionPBR}) and \mbox{\texttt{\mdseries\slshape n}} is a non\texttt{\symbol{45}}negative integer, then \texttt{AsBipartition} returns the bipartition corresponding to \mbox{\texttt{\mdseries\slshape x}} with degree \mbox{\texttt{\mdseries\slshape n}}. 
\end{description}
 
\begin{Verbatim}[commandchars=!@|,fontsize=\small,frame=single,label=Example]
  !gapprompt@gap>| !gapinput@x := Transformation([3, 5, 3, 4, 1, 2]);;|
  !gapprompt@gap>| !gapinput@AsBipartition(x, 5);|
  <bipartition: [ 1, 3, -3 ], [ 2, -5 ], [ 4, -4 ], [ 5, -1 ], [ -2 ]>
  !gapprompt@gap>| !gapinput@AsBipartition(x);|
  <bipartition: [ 1, 3, -3 ], [ 2, -5 ], [ 4, -4 ], [ 5, -1 ],
   [ 6, -2 ], [ -6 ]>
  !gapprompt@gap>| !gapinput@AsBipartition(x, 10);|
  <bipartition: [ 1, 3, -3 ], [ 2, -5 ], [ 4, -4 ], [ 5, -1 ],
   [ 6, -2 ], [ 7, -7 ], [ 8, -8 ], [ 9, -9 ], [ 10, -10 ], [ -6 ]>
  !gapprompt@gap>| !gapinput@AsBipartition((1, 3)(2, 4));|
  <block bijection: [ 1, -3 ], [ 2, -4 ], [ 3, -1 ], [ 4, -2 ]>
  !gapprompt@gap>| !gapinput@AsBipartition((1, 3)(2, 4), 10);|
  <block bijection: [ 1, -3 ], [ 2, -4 ], [ 3, -1 ], [ 4, -2 ],
   [ 5, -5 ], [ 6, -6 ], [ 7, -7 ], [ 8, -8 ], [ 9, -9 ], [ 10, -10 ]>
  !gapprompt@gap>| !gapinput@x := PartialPerm([1, 2, 3, 4, 5, 6], [6, 7, 1, 4, 3, 2]);;|
  !gapprompt@gap>| !gapinput@AsBipartition(x, 11);|
  <bipartition: [ 1, -6 ], [ 2, -7 ], [ 3, -1 ], [ 4, -4 ], [ 5, -3 ],
   [ 6, -2 ], [ 7 ], [ 8 ], [ 9 ], [ 10 ], [ 11 ], [ -5 ], [ -8 ],
   [ -9 ], [ -10 ], [ -11 ]>
  !gapprompt@gap>| !gapinput@AsBipartition(x);|
  <bipartition: [ 1, -6 ], [ 2, -7 ], [ 3, -1 ], [ 4, -4 ], [ 5, -3 ],
   [ 6, -2 ], [ 7 ], [ -5 ]>
  !gapprompt@gap>| !gapinput@AsBipartition(Transformation([1, 1, 2]), 1);|
  <block bijection: [ 1, -1 ]>
  !gapprompt@gap>| !gapinput@x := Bipartition([[1, 2, -2], [3], [4, 5, 6, -1],|
  !gapprompt@>| !gapinput@                     [-3, -4, -5, -6]]);;|
  !gapprompt@gap>| !gapinput@AsBipartition(x, 0);|
  <empty bipartition>
  !gapprompt@gap>| !gapinput@AsBipartition(x, 2);|
  <bipartition: [ 1, 2, -2 ], [ -1 ]>
  !gapprompt@gap>| !gapinput@AsBipartition(x, 8);|
  <bipartition: [ 1, 2, -2 ], [ 3 ], [ 4, 5, 6, -1 ], [ 7 ], [ 8 ],
   [ -3, -4, -5, -6 ], [ -7 ], [ -8 ]>
  !gapprompt@gap>| !gapinput@x := PBR(|
  !gapprompt@>| !gapinput@[[-1, 1, 2, 3, 4], [-1, 1, 2, 3, 4],|
  !gapprompt@>| !gapinput@ [-1, 1, 2, 3, 4], [-1, 1, 2, 3, 4]],|
  !gapprompt@>| !gapinput@[[-1, 1, 2, 3, 4], [-2], [-3], [-4]]);;|
  !gapprompt@gap>| !gapinput@AsBipartition(x);|
  <bipartition: [ 1, 2, 3, 4, -1 ], [ -2 ], [ -3 ], [ -4 ]>
  !gapprompt@gap>| !gapinput@AsBipartition(x, 2);|
  <bipartition: [ 1, 2, -1 ], [ -2 ]>
  !gapprompt@gap>| !gapinput@AsBipartition(x, 4);|
  <bipartition: [ 1, 2, 3, 4, -1 ], [ -2 ], [ -3 ], [ -4 ]>
  !gapprompt@gap>| !gapinput@AsBipartition(x, 5);|
  <bipartition: [ 1, 2, 3, 4, -1 ], [ 5 ], [ -2 ], [ -3 ], [ -4 ],
   [ -5 ]>
  !gapprompt@gap>| !gapinput@AsBipartition(x, 0);|
  <empty bipartition>
\end{Verbatim}
 }

 

\subsection{\textcolor{Chapter }{AsBlockBijection}}
\logpage{[ 3, 3, 2 ]}\nobreak
\hyperdef{L}{X85A5AD2B7F3B776F}{}
{\noindent\textcolor{FuncColor}{$\triangleright$\enspace\texttt{AsBlockBijection({\mdseries\slshape x[, n]})\index{AsBlockBijection@\texttt{AsBlockBijection}}
\label{AsBlockBijection}
}\hfill{\scriptsize (operation)}}\\
\textbf{\indent Returns:\ }
A block bijection.



 When the argument \mbox{\texttt{\mdseries\slshape x}} is a partial perm and \mbox{\texttt{\mdseries\slshape n}} is a positive integer which is greater than the maximum of the degree and
codegree of \mbox{\texttt{\mdseries\slshape x}}, this function returns a block bijection corresponding to \mbox{\texttt{\mdseries\slshape x}}. This block bijection has the same non\texttt{\symbol{45}}singleton classes
as \texttt{g := AsBipartition(\mbox{\texttt{\mdseries\slshape x}}, \mbox{\texttt{\mdseries\slshape n}})} and one additional class which is the union the singleton classes of \texttt{g}. 

 If the optional second argument \mbox{\texttt{\mdseries\slshape n}} is not present, then the maximum of the degree and codegree of \mbox{\texttt{\mdseries\slshape x}} plus 1 is used by default. If the second argument \mbox{\texttt{\mdseries\slshape n}} is not greater than this maximum, then an error is given. 

 This is the value at \mbox{\texttt{\mdseries\slshape x}} of the embedding of the symmetric inverse monoid into the dual symmetric
inverse monoid given in the FitzGerald\texttt{\symbol{45}}Leech Theorem \cite{Fitzgerald1998aa}. 

 When the argument \mbox{\texttt{\mdseries\slshape x}} is a partial perm bipartition (see \texttt{IsPartialPermBipartition} (\ref{IsPartialPermBipartition})) then this operation returns \texttt{AsBlockBijection(AsPartialPerm(\mbox{\texttt{\mdseries\slshape x}})[, \mbox{\texttt{\mdseries\slshape n}}])}. 
\begin{Verbatim}[commandchars=!@|,fontsize=\small,frame=single,label=Example]
  !gapprompt@gap>| !gapinput@x := PartialPerm([1, 2, 3, 6, 7, 10], [9, 5, 6, 1, 7, 8]);|
  [2,5][3,6,1,9][10,8](7)
  !gapprompt@gap>| !gapinput@AsBipartition(x, 11);|
  <bipartition: [ 1, -9 ], [ 2, -5 ], [ 3, -6 ], [ 4 ], [ 5 ],
   [ 6, -1 ], [ 7, -7 ], [ 8 ], [ 9 ], [ 10, -8 ], [ 11 ], [ -2 ],
   [ -3 ], [ -4 ], [ -10 ], [ -11 ]>
  !gapprompt@gap>| !gapinput@AsBlockBijection(x, 10);|
  Error, the 2nd argument (a pos. int.) is less than or equal to the max\
  imum of the degree and codegree of the 1st argument (a partial perm)
  !gapprompt@gap>| !gapinput@AsBlockBijection(x, 11);|
  <block bijection: [ 1, -9 ], [ 2, -5 ], [ 3, -6 ],
   [ 4, 5, 8, 9, 11, -2, -3, -4, -10, -11 ], [ 6, -1 ], [ 7, -7 ],
   [ 10, -8 ]>
  !gapprompt@gap>| !gapinput@x := Bipartition([[1, -3], [2], [3, -2], [-1]]);;|
  !gapprompt@gap>| !gapinput@IsPartialPermBipartition(x);|
  true
  !gapprompt@gap>| !gapinput@AsBlockBijection(x);|
  <block bijection: [ 1, -3 ], [ 2, 4, -1, -4 ], [ 3, -2 ]>
\end{Verbatim}
 }

 

\subsection{\textcolor{Chapter }{AsTransformation (for a bipartition)}}
\logpage{[ 3, 3, 3 ]}\nobreak
\hyperdef{L}{X7CE91D0C83865214}{}
{\noindent\textcolor{FuncColor}{$\triangleright$\enspace\texttt{AsTransformation({\mdseries\slshape x})\index{AsTransformation@\texttt{AsTransformation}!for a bipartition}
\label{AsTransformation:for a bipartition}
}\hfill{\scriptsize (attribute)}}\\
\textbf{\indent Returns:\ }
A transformation.



 When the argument \mbox{\texttt{\mdseries\slshape x}} is a bipartition, that mathematically defines a transformation, this function
returns that transformation. A bipartition \mbox{\texttt{\mdseries\slshape x}} defines a transformation if and only if its right blocks are the image list of
a permutation of \texttt{[1 .. n]} where \texttt{n} is the degree of \mbox{\texttt{\mdseries\slshape x}}. 

 See \texttt{IsTransBipartition} (\ref{IsTransBipartition}). 
\begin{Verbatim}[commandchars=!@|,fontsize=\small,frame=single,label=Example]
  !gapprompt@gap>| !gapinput@x := Bipartition([[1, -3], [2, -2], [3, 5, 10, -7],|
  !gapprompt@>| !gapinput@                     [4, -12], [6, 7, -6], [8, -5], [9, -11],|
  !gapprompt@>| !gapinput@                     [11, 12, -10], [-1], [-4], [-8], [-9]]);;|
  !gapprompt@gap>| !gapinput@AsTransformation(x);|
  Transformation( [ 3, 2, 7, 12, 7, 6, 6, 5, 11, 7, 10, 10 ] )
  !gapprompt@gap>| !gapinput@IsTransBipartition(x);|
  true
  !gapprompt@gap>| !gapinput@x := Bipartition([[1, 5], [2, 4, 8, 10],|
  !gapprompt@>| !gapinput@                     [3, 6, 7, -1, -2], [9, -4, -6, -9],|
  !gapprompt@>| !gapinput@                     [-3, -5], [-7, -8], [-10]]);;|
  !gapprompt@gap>| !gapinput@AsTransformation(x);|
  Error, the argument (a bipartition) does not define a transformation
\end{Verbatim}
 }

 

\subsection{\textcolor{Chapter }{AsPartialPerm (for a bipartition)}}
\logpage{[ 3, 3, 4 ]}\nobreak
\hyperdef{L}{X7C5212EF7A200E63}{}
{\noindent\textcolor{FuncColor}{$\triangleright$\enspace\texttt{AsPartialPerm({\mdseries\slshape x})\index{AsPartialPerm@\texttt{AsPartialPerm}!for a bipartition}
\label{AsPartialPerm:for a bipartition}
}\hfill{\scriptsize (operation)}}\\
\textbf{\indent Returns:\ }
A partial perm.



 When the argument \mbox{\texttt{\mdseries\slshape x}} is a bipartition that mathematically defines a partial perm, this function
returns that partial perm. 

 A bipartition \mbox{\texttt{\mdseries\slshape x}} defines a partial perm if and only if its numbers of left and right blocks
both equal its degree. 

 See \texttt{IsPartialPermBipartition} (\ref{IsPartialPermBipartition}). 
\begin{Verbatim}[commandchars=!@|,fontsize=\small,frame=single,label=Example]
  !gapprompt@gap>| !gapinput@x := Bipartition([[1, -4], [2, -2], [3, -10], [4, -5],|
  !gapprompt@>| !gapinput@                     [5, -9], [6], [7], [8, -6], [9, -3], [10, -8],|
  !gapprompt@>| !gapinput@                     [-1], [-7]]);;|
  !gapprompt@gap>| !gapinput@IsPartialPermBipartition(x);|
  true
  !gapprompt@gap>| !gapinput@AsPartialPerm(x);|
  [1,4,5,9,3,10,8,6](2)
  !gapprompt@gap>| !gapinput@x := Bipartition([[1, -2, -4], [2, 3, 4, -3], [-1]]);;|
  !gapprompt@gap>| !gapinput@IsPartialPermBipartition(x);|
  false
  !gapprompt@gap>| !gapinput@AsPartialPerm(x);|
  Error, the argument (a bipartition) does not define a partial perm
\end{Verbatim}
 }

 

\subsection{\textcolor{Chapter }{AsPermutation (for a bipartition)}}
\logpage{[ 3, 3, 5 ]}\nobreak
\hyperdef{L}{X7C684CD38405DBEF}{}
{\noindent\textcolor{FuncColor}{$\triangleright$\enspace\texttt{AsPermutation({\mdseries\slshape x})\index{AsPermutation@\texttt{AsPermutation}!for a bipartition}
\label{AsPermutation:for a bipartition}
}\hfill{\scriptsize (attribute)}}\\
\textbf{\indent Returns:\ }
A permutation.



 When the argument \mbox{\texttt{\mdseries\slshape x}} is a bipartition that mathematically defines a permutation, this function
returns that permutation. 

 A bipartition \mbox{\texttt{\mdseries\slshape x}} defines a permutation if and only if its numbers of left, right, and
transverse blocks all equal its degree.

 See \texttt{IsPermBipartition} (\ref{IsPermBipartition}). 
\begin{Verbatim}[commandchars=!@|,fontsize=\small,frame=single,label=Example]
  !gapprompt@gap>| !gapinput@x := Bipartition([[1, -6], [2, -4], [3, -2], [4, -5],|
  !gapprompt@>| !gapinput@                     [5, -3], [6, -1]]);;|
  !gapprompt@gap>| !gapinput@IsPermBipartition(x);|
  true
  !gapprompt@gap>| !gapinput@AsPermutation(x);|
  (1,6)(2,4,5,3)
  !gapprompt@gap>| !gapinput@AsBipartition(last) = x;|
  true
\end{Verbatim}
 }

 }

   
\section{\textcolor{Chapter }{Operators for bipartitions}}\label{operators-bipartitions}
\logpage{[ 3, 4, 0 ]}
\hyperdef{L}{X83F2C3C97E8FFA49}{}
{
  
\begin{description}
\item[{\texttt{\mbox{\texttt{\mdseries\slshape f}} * \mbox{\texttt{\mdseries\slshape g}}}}]  \index{*@\texttt{*} (for bipartitions)} returns the composition of \mbox{\texttt{\mdseries\slshape f}} and \mbox{\texttt{\mdseries\slshape g}} when \mbox{\texttt{\mdseries\slshape f}} and \mbox{\texttt{\mdseries\slshape g}} are bipartitions. 
\item[{\texttt{\mbox{\texttt{\mdseries\slshape f}} {\textless} \mbox{\texttt{\mdseries\slshape g}}}}]  \index{<@\texttt{{\textless}} (for bipartitions)} returns \texttt{true} if the internal representation of \mbox{\texttt{\mdseries\slshape f}} is lexicographically less than the internal representation of \mbox{\texttt{\mdseries\slshape g}} and \texttt{false} if it is not. 
\item[{\texttt{\mbox{\texttt{\mdseries\slshape f}} = \mbox{\texttt{\mdseries\slshape g}}}}]  \index{=@\texttt{=} (for bipartitions)} returns \texttt{true} if the bipartition \mbox{\texttt{\mdseries\slshape f}} equals the bipartition \mbox{\texttt{\mdseries\slshape g}} and returns \texttt{false} if it does not. 
\end{description}
 

\subsection{\textcolor{Chapter }{PartialPermLeqBipartition}}
\logpage{[ 3, 4, 1 ]}\nobreak
\hyperdef{L}{X7A39D36086647536}{}
{\noindent\textcolor{FuncColor}{$\triangleright$\enspace\texttt{PartialPermLeqBipartition({\mdseries\slshape x, y})\index{PartialPermLeqBipartition@\texttt{PartialPermLeqBipartition}}
\label{PartialPermLeqBipartition}
}\hfill{\scriptsize (operation)}}\\
\textbf{\indent Returns:\ }
\texttt{true} or \texttt{false}.



 If \mbox{\texttt{\mdseries\slshape x}} and \mbox{\texttt{\mdseries\slshape y}} are partial perm bipartitions, i.e. they satisfy \texttt{IsPartialPermBipartition} (\ref{IsPartialPermBipartition}), then this function returns \texttt{AsPartialPerm(\mbox{\texttt{\mdseries\slshape x}}) {\textless} AsPartialPerm(\mbox{\texttt{\mdseries\slshape y}})}. }

 

\subsection{\textcolor{Chapter }{NaturalLeqPartialPermBipartition}}
\logpage{[ 3, 4, 2 ]}\nobreak
\hyperdef{L}{X8608D78F83D55108}{}
{\noindent\textcolor{FuncColor}{$\triangleright$\enspace\texttt{NaturalLeqPartialPermBipartition({\mdseries\slshape x, y})\index{NaturalLeqPartialPermBipartition@\texttt{NaturalLeqPartialPermBipartition}}
\label{NaturalLeqPartialPermBipartition}
}\hfill{\scriptsize (operation)}}\\
\textbf{\indent Returns:\ }
\texttt{true} or \texttt{false}.



 The \emph{natural partial order} $\leq$ on an inverse semigroup \texttt{S} is defined by \texttt{s} $\leq$ \texttt{t} if there exists an idempotent \texttt{e} in \texttt{S} such that \texttt{s = et}. Hence if \mbox{\texttt{\mdseries\slshape x}} and \mbox{\texttt{\mdseries\slshape y}} are partial perm bipartitions, then \mbox{\texttt{\mdseries\slshape x}} $\leq$ \mbox{\texttt{\mdseries\slshape y}} if and only if \texttt{AsPartialPerm(\mbox{\texttt{\mdseries\slshape x}})} is a restriction of \texttt{AsPartialPerm(\mbox{\texttt{\mdseries\slshape y}})}. 

 \texttt{NaturalLeqPartialPermBipartition} returns \texttt{true} if \texttt{AsPartialPerm(\mbox{\texttt{\mdseries\slshape x}})} is a restriction of \texttt{AsPartialPerm(\mbox{\texttt{\mdseries\slshape y}})} and \texttt{false} if it is not. Note that since this is a partial order and not a total order,
it is possible that \mbox{\texttt{\mdseries\slshape x}} and \mbox{\texttt{\mdseries\slshape y}} are incomparable with respect to the natural partial order. }

 

\subsection{\textcolor{Chapter }{NaturalLeqBlockBijection}}
\logpage{[ 3, 4, 3 ]}\nobreak
\hyperdef{L}{X79E8FA077E24C1F4}{}
{\noindent\textcolor{FuncColor}{$\triangleright$\enspace\texttt{NaturalLeqBlockBijection({\mdseries\slshape x, y})\index{NaturalLeqBlockBijection@\texttt{NaturalLeqBlockBijection}}
\label{NaturalLeqBlockBijection}
}\hfill{\scriptsize (operation)}}\\
\textbf{\indent Returns:\ }
\texttt{true} or \texttt{false}.



 The \emph{natural partial order} $\leq$ on an inverse semigroup \texttt{S} is defined by \texttt{s} $\leq$ \texttt{t} if there exists an idempotent \texttt{e} in \texttt{S} such that \texttt{s = et}. Hence if \mbox{\texttt{\mdseries\slshape x}} and \mbox{\texttt{\mdseries\slshape y}} are block bijections, then \mbox{\texttt{\mdseries\slshape x}} $\leq$ \mbox{\texttt{\mdseries\slshape y}} if and only if \mbox{\texttt{\mdseries\slshape x}} contains \mbox{\texttt{\mdseries\slshape y}}. 

 \texttt{NaturalLeqBlockBijection} returns \texttt{true} if \mbox{\texttt{\mdseries\slshape x}} is contained in \mbox{\texttt{\mdseries\slshape y}} and \texttt{false} if it is not. Note that since this is a partial order and not a total order,
it is possible that \mbox{\texttt{\mdseries\slshape x}} and \mbox{\texttt{\mdseries\slshape y}} are incomparable with respect to the natural partial order. 
\begin{Verbatim}[commandchars=!@|,fontsize=\small,frame=single,label=Example]
  !gapprompt@gap>| !gapinput@x := Bipartition([[1, 2, -3], [3, -1, -2], [4, -4],|
  !gapprompt@>| !gapinput@                     [5, -5], [6, -6], [7, -7],|
  !gapprompt@>| !gapinput@                     [8, -8], [9, -9], [10, -10]]);;|
  !gapprompt@gap>| !gapinput@y := Bipartition([[1, -2], [2, -1], [3, -3],|
  !gapprompt@>| !gapinput@                     [4, -4], [5, -5], [6, -6], [7, -7],|
  !gapprompt@>| !gapinput@                     [8, -8], [9, -9], [10, -10]]);;|
  !gapprompt@gap>| !gapinput@z := Bipartition([Union([1 .. 10], [-10 .. -1])]);;|
  !gapprompt@gap>| !gapinput@NaturalLeqBlockBijection(x, y);|
  false
  !gapprompt@gap>| !gapinput@NaturalLeqBlockBijection(y, x);|
  false
  !gapprompt@gap>| !gapinput@NaturalLeqBlockBijection(z, x);|
  true
  !gapprompt@gap>| !gapinput@NaturalLeqBlockBijection(z, y);|
  true
\end{Verbatim}
 }

 

\subsection{\textcolor{Chapter }{PermLeftQuoBipartition}}
\logpage{[ 3, 4, 4 ]}\nobreak
\hyperdef{L}{X7D9F5A248028FF52}{}
{\noindent\textcolor{FuncColor}{$\triangleright$\enspace\texttt{PermLeftQuoBipartition({\mdseries\slshape x, y})\index{PermLeftQuoBipartition@\texttt{PermLeftQuoBipartition}}
\label{PermLeftQuoBipartition}
}\hfill{\scriptsize (operation)}}\\
\textbf{\indent Returns:\ }
A permutation.



 If \mbox{\texttt{\mdseries\slshape x}} and \mbox{\texttt{\mdseries\slshape y}} are bipartitions with equal left and right blocks, then \texttt{PermLeftQuoBipartition} returns the permutation of the indices of the right blocks of \mbox{\texttt{\mdseries\slshape x}} (and \mbox{\texttt{\mdseries\slshape y}}) induced by \texttt{Star(\mbox{\texttt{\mdseries\slshape x}}) * \mbox{\texttt{\mdseries\slshape y}}}. 

 \texttt{PermLeftQuoBipartition} verifies that \mbox{\texttt{\mdseries\slshape x}} and \mbox{\texttt{\mdseries\slshape y}} have equal left and right blocks, and returns an error if they do not. 
\begin{Verbatim}[commandchars=!@|,fontsize=\small,frame=single,label=Example]
  !gapprompt@gap>| !gapinput@x := Bipartition([[1, 4, 6, 7, 8, 10], [2, 5, -1, -2, -8],|
  !gapprompt@>| !gapinput@                     [3, -3, -6, -7, -9], [9, -4, -5], [-10]]);;|
  !gapprompt@gap>| !gapinput@y := Bipartition([[1, 4, 6, 7, 8, 10], [2, 5, -3, -6, -7, -9],|
  !gapprompt@>| !gapinput@                     [3, -4, -5], [9, -1, -2, -8], [-10]]);;|
  !gapprompt@gap>| !gapinput@PermLeftQuoBipartition(x, y);|
  (1,2,3)
  !gapprompt@gap>| !gapinput@Star(x) * y;|
  <bipartition: [ 1, 2, 8, -3, -6, -7, -9 ], [ 3, 6, 7, 9, -4, -5 ],
   [ 4, 5, -1, -2, -8 ], [ 10 ], [ -10 ]>
\end{Verbatim}
 }

 }

   
\section{\textcolor{Chapter }{Attributes for bipartitons}}\label{attributes-bipartitions}
\logpage{[ 3, 5, 0 ]}
\hyperdef{L}{X87F3A304814797CE}{}
{
  In this section we describe various attributes that a bipartition can possess. 

\subsection{\textcolor{Chapter }{DegreeOfBipartition}}
\logpage{[ 3, 5, 1 ]}\nobreak
\hyperdef{L}{X780F5E00784FE58C}{}
{\noindent\textcolor{FuncColor}{$\triangleright$\enspace\texttt{DegreeOfBipartition({\mdseries\slshape x})\index{DegreeOfBipartition@\texttt{DegreeOfBipartition}}
\label{DegreeOfBipartition}
}\hfill{\scriptsize (attribute)}}\\
\noindent\textcolor{FuncColor}{$\triangleright$\enspace\texttt{DegreeOfBipartitionCollection({\mdseries\slshape x})\index{DegreeOfBipartitionCollection@\texttt{DegreeOfBipartitionCollection}}
\label{DegreeOfBipartitionCollection}
}\hfill{\scriptsize (attribute)}}\\
\textbf{\indent Returns:\ }
A positive integer.



 The degree of a bipartition is, roughly speaking, the number of points where
it is defined. More precisely, if \mbox{\texttt{\mdseries\slshape x}} is a bipartition defined on \texttt{2 * n} points, then the degree of \mbox{\texttt{\mdseries\slshape x}} is \texttt{n}. 

 The degree of a collection \mbox{\texttt{\mdseries\slshape coll}} of bipartitions of equal degree is just the degree of any (and every)
bipartition in \mbox{\texttt{\mdseries\slshape coll}}. The degree of collection of bipartitions of unequal degrees is not defined. 
\begin{Verbatim}[commandchars=!@|,fontsize=\small,frame=single,label=Example]
  !gapprompt@gap>| !gapinput@x := Bipartition([[1, 7, -3, -8], [2, 6],|
  !gapprompt@>| !gapinput@                     [3], [4, -7, -9], [5, 9, -2],|
  !gapprompt@>| !gapinput@                     [8, -1, -4, -6], [-5]]);;|
  !gapprompt@gap>| !gapinput@DegreeOfBipartition(x);|
  9
  !gapprompt@gap>| !gapinput@S := BrauerMonoid(5);|
  <regular bipartition *-monoid of degree 5 with 3 generators>
  !gapprompt@gap>| !gapinput@IsBipartitionCollection(S);|
  true
  !gapprompt@gap>| !gapinput@DegreeOfBipartitionCollection(S);|
  5
\end{Verbatim}
 }

 

\subsection{\textcolor{Chapter }{RankOfBipartition}}
\logpage{[ 3, 5, 2 ]}\nobreak
\hyperdef{L}{X82074756826AD2C2}{}
{\noindent\textcolor{FuncColor}{$\triangleright$\enspace\texttt{RankOfBipartition({\mdseries\slshape x})\index{RankOfBipartition@\texttt{RankOfBipartition}}
\label{RankOfBipartition}
}\hfill{\scriptsize (attribute)}}\\
\noindent\textcolor{FuncColor}{$\triangleright$\enspace\texttt{NrTransverseBlocks({\mdseries\slshape x})\index{NrTransverseBlocks@\texttt{NrTransverseBlocks}!for a bipartition}
\label{NrTransverseBlocks:for a bipartition}
}\hfill{\scriptsize (attribute)}}\\
\textbf{\indent Returns:\ }
The rank of a bipartition.



 When the argument is a bipartition \mbox{\texttt{\mdseries\slshape x}}, \texttt{RankOfBipartition} returns the number of blocks of \mbox{\texttt{\mdseries\slshape x}} containing both positive and negative entries, i.e. the number of transverse
blocks of \mbox{\texttt{\mdseries\slshape x}}. 

 \texttt{NrTransverseBlocks} is just a synonym for \texttt{RankOfBipartition}. 
\begin{Verbatim}[commandchars=!@|,fontsize=\small,frame=single,label=Example]
  !gapprompt@gap>| !gapinput@x := Bipartition([[1, 2, 6, 7, -4, -5, -7], [3, 4, 5, -1, -3],|
  !gapprompt@>| !gapinput@                     [8, -9], [9, -2], [-6], [-8]]);|
  <bipartition: [ 1, 2, 6, 7, -4, -5, -7 ], [ 3, 4, 5, -1, -3 ],
   [ 8, -9 ], [ 9, -2 ], [ -6 ], [ -8 ]>
  !gapprompt@gap>| !gapinput@RankOfBipartition(x);|
  4
\end{Verbatim}
 }

 

\subsection{\textcolor{Chapter }{ExtRepOfObj (for a bipartition)}}
\logpage{[ 3, 5, 3 ]}\nobreak
\hyperdef{L}{X86F6506C780C6E08}{}
{\noindent\textcolor{FuncColor}{$\triangleright$\enspace\texttt{ExtRepOfObj({\mdseries\slshape x})\index{ExtRepOfObj@\texttt{ExtRepOfObj}!for a bipartition}
\label{ExtRepOfObj:for a bipartition}
}\hfill{\scriptsize (operation)}}\\
\textbf{\indent Returns:\ }
A partition of \texttt{[1 .. 2 * n]}.



 If \texttt{n} is the degree of the bipartition \mbox{\texttt{\mdseries\slshape x}}, then \texttt{ExtRepOfObj} returns the partition of \texttt{[\texttt{\symbol{45}}n .. \texttt{\symbol{45}}1]} union \texttt{[1 .. n]} corresponding to \mbox{\texttt{\mdseries\slshape x}} as a sorted list of duplicate\texttt{\symbol{45}}free lists. 
\begin{Verbatim}[commandchars=!@|,fontsize=\small,frame=single,label=Example]
  !gapprompt@gap>| !gapinput@x := Bipartition([[1, 5, -3], [2, 4, -2, -4], [3, -1, -5]]);|
  <block bijection: [ 1, 5, -3 ], [ 2, 4, -2, -4 ], [ 3, -1, -5 ]>
  !gapprompt@gap>| !gapinput@ExtRepOfObj(x);|
  [ [ 1, 5, -3 ], [ 2, 4, -2, -4 ], [ 3, -1, -5 ] ]
\end{Verbatim}
 }

 

\subsection{\textcolor{Chapter }{IntRepOfBipartition}}
\logpage{[ 3, 5, 4 ]}\nobreak
\hyperdef{L}{X7ECD393A854C073B}{}
{\noindent\textcolor{FuncColor}{$\triangleright$\enspace\texttt{IntRepOfBipartition({\mdseries\slshape x})\index{IntRepOfBipartition@\texttt{IntRepOfBipartition}}
\label{IntRepOfBipartition}
}\hfill{\scriptsize (attribute)}}\\
\textbf{\indent Returns:\ }
A list of positive integers.



 If \mbox{\texttt{\mdseries\slshape x}} is a bipartition with degree \texttt{n}, then \texttt{IntRepOfBipartition} returns the \emph{internal representation} of \mbox{\texttt{\mdseries\slshape x}}: a list of length \texttt{2 * n} containing positive integers which correspond to the blocks of \mbox{\texttt{\mdseries\slshape x}}. 

 If \texttt{i} is in \texttt{[1 .. n]}, then \texttt{list[i]} refers to the point \texttt{i}; if \texttt{i} is in \texttt{[n + 1 .. 2 * n]}, then \texttt{list[i]} refers to the point \texttt{n \texttt{\symbol{45}} i} (a negative point). Two points lie in the same block of the bipartition if and
only if their entries in the list are equal. 

 See also \texttt{BipartitionByIntRep} (\ref{BipartitionByIntRep}). 
\begin{Verbatim}[commandchars=!@|,fontsize=\small,frame=single,label=Example]
  !gapprompt@gap>| !gapinput@x := Bipartition([[1, -3], [3, 4], [2, -1, -2], [-4]]);|
  <bipartition: [ 1, -3 ], [ 2, -1, -2 ], [ 3, 4 ], [ -4 ]>
  !gapprompt@gap>| !gapinput@IntRepOfBipartition(x);|
  [ 1, 2, 3, 3, 2, 2, 1, 4 ]
\end{Verbatim}
 }

 

\subsection{\textcolor{Chapter }{RightBlocks}}
\logpage{[ 3, 5, 5 ]}\nobreak
\hyperdef{L}{X86A10B138230C2A4}{}
{\noindent\textcolor{FuncColor}{$\triangleright$\enspace\texttt{RightBlocks({\mdseries\slshape x})\index{RightBlocks@\texttt{RightBlocks}}
\label{RightBlocks}
}\hfill{\scriptsize (attribute)}}\\
\textbf{\indent Returns:\ }
The right blocks of a bipartition.



 \texttt{RightBlocks} returns the right blocks of the bipartition \mbox{\texttt{\mdseries\slshape x}}. 

 The \emph{right blocks} of a bipartition \mbox{\texttt{\mdseries\slshape x}} are just the intersections of the blocks of \mbox{\texttt{\mdseries\slshape x}} with \texttt{[\texttt{\symbol{45}}n .. \texttt{\symbol{45}}1]} where \texttt{n} is the degree of \mbox{\texttt{\mdseries\slshape x}}, the values in transverse blocks are positive, and the values in
non\texttt{\symbol{45}}transverse blocks are negative. 

 The right blocks of a bipartition are \textsf{GAP} objects in their own right, and are not simply a list of blocks of \mbox{\texttt{\mdseries\slshape x}}; see \ref{section-blocks} for more information. 

 The significance of this notion lies in the fact that bipartitions \mbox{\texttt{\mdseries\slshape x}} and \texttt{y} are $\mathcal{L}$\texttt{\symbol{45}}related in the partition monoid if and only if they have
equal right blocks. 
\begin{Verbatim}[commandchars=!@|,fontsize=\small,frame=single,label=Example]
  !gapprompt@gap>| !gapinput@x := Bipartition([[1, 4, 7, 8, -4], [2, 3, 5, -2, -7],|
  !gapprompt@>| !gapinput@                     [6, -1], [-3], [-5, -6, -8]]);;|
  !gapprompt@gap>| !gapinput@RightBlocks(x);|
  <blocks: [ 1* ], [ 2*, 7* ], [ 3 ], [ 4* ], [ 5, 6, 8 ]>
  !gapprompt@gap>| !gapinput@LeftBlocks(x);|
  <blocks: [ 1*, 4*, 7*, 8* ], [ 2*, 3*, 5* ], [ 6* ]>
\end{Verbatim}
 }

 

\subsection{\textcolor{Chapter }{LeftBlocks}}
\logpage{[ 3, 5, 6 ]}\nobreak
\hyperdef{L}{X7B9B364379D8F4E8}{}
{\noindent\textcolor{FuncColor}{$\triangleright$\enspace\texttt{LeftBlocks({\mdseries\slshape x})\index{LeftBlocks@\texttt{LeftBlocks}}
\label{LeftBlocks}
}\hfill{\scriptsize (attribute)}}\\
\textbf{\indent Returns:\ }
The left blocks of a bipartition.



 \texttt{LeftBlocks} returns the left blocks of the bipartition \mbox{\texttt{\mdseries\slshape x}}. 

 The \emph{left blocks} of a bipartition \mbox{\texttt{\mdseries\slshape x}} are just the intersections of the blocks of \mbox{\texttt{\mdseries\slshape x}} with \texttt{[1..n]} where \texttt{n} is the degree of \mbox{\texttt{\mdseries\slshape x}}, the values in transverse blocks are positive, and the values in
non\texttt{\symbol{45}}transverse blocks are negative. 

 The left blocks of a bipartition are \textsf{GAP} objects in their own right, and are not simply a list of blocks of \mbox{\texttt{\mdseries\slshape x}}; see \ref{section-blocks} for more information. 

 The significance of this notion lies in the fact that bipartitions \mbox{\texttt{\mdseries\slshape x}} and \texttt{y} are $\mathcal{R}$\texttt{\symbol{45}}related in the partition monoid if and only if they have
equal left blocks. 
\begin{Verbatim}[commandchars=!@|,fontsize=\small,frame=single,label=Example]
  !gapprompt@gap>| !gapinput@x := Bipartition([[1, 4, 7, 8, -4], [2, 3, 5, -2, -7],|
  !gapprompt@>| !gapinput@                     [6, -1], [-3], [-5, -6, -8]]);;|
  !gapprompt@gap>| !gapinput@RightBlocks(x);|
  <blocks: [ 1* ], [ 2*, 7* ], [ 3 ], [ 4* ], [ 5, 6, 8 ]>
  !gapprompt@gap>| !gapinput@LeftBlocks(x);|
  <blocks: [ 1*, 4*, 7*, 8* ], [ 2*, 3*, 5* ], [ 6* ]>
\end{Verbatim}
 }

 

\subsection{\textcolor{Chapter }{IrreducibleComponentsOfBipartition}}
\logpage{[ 3, 5, 7 ]}\nobreak
\hyperdef{L}{X80979FA37D1559DA}{}
{\noindent\textcolor{FuncColor}{$\triangleright$\enspace\texttt{IrreducibleComponentsOfBipartition({\mdseries\slshape x})\index{IrreducibleComponentsOfBipartition@\texttt{IrreducibleComponentsOfBipartition}}
\label{IrreducibleComponentsOfBipartition}
}\hfill{\scriptsize (operation)}}\\
\textbf{\indent Returns:\ }
A list of bipartitions.



 The irreducible components of the bipartition \mbox{\texttt{\mdseries\slshape x}} are the maximal degree bipartitions \texttt{y1, ..., ym} such that \texttt{TensorBipartitions(y1, ... ym)} equals \mbox{\texttt{\mdseries\slshape x}}. See also: \texttt{TensorBipartitions} (\ref{TensorBipartitions:for a pair of bipartitions}) and \texttt{IsIrreducibleBipartition} (\ref{IsIrreducibleBipartition}).  
\begin{Verbatim}[commandchars=!@|,fontsize=\small,frame=single,label=Example]
  !gapprompt@gap>| !gapinput@x := TensorBipartitions(|
  !gapprompt@>| !gapinput@Bipartition([[1, 2, -2], [-1]]),|
  !gapprompt@>| !gapinput@Bipartition([[1, 2, 3, -2], [-1], [-3]]));|
  <bipartition: [ 1, 2, -2 ], [ 3, 4, 5, -4 ], [ -1 ], [ -3 ], [ -5 ]>
  !gapprompt@gap>| !gapinput@coll := IrreducibleComponentsOfBipartition(x);|
  [ <bipartition: [ 1, 2, -2 ], [ -1 ]>, 
  <bipartition: [ 1, 2, 3, -2 ], [ -1 ], [ -3 ]> ]
  !gapprompt@gap>| !gapinput@TensorBipartitions(coll) = x;|
  true
\end{Verbatim}
 }

 

\subsection{\textcolor{Chapter }{NrLeftBlocks}}
\logpage{[ 3, 5, 8 ]}\nobreak
\hyperdef{L}{X79AEDB5382FD25CF}{}
{\noindent\textcolor{FuncColor}{$\triangleright$\enspace\texttt{NrLeftBlocks({\mdseries\slshape x})\index{NrLeftBlocks@\texttt{NrLeftBlocks}}
\label{NrLeftBlocks}
}\hfill{\scriptsize (attribute)}}\\
\textbf{\indent Returns:\ }
A non\texttt{\symbol{45}}negative integer.



 When the argument is a bipartition \mbox{\texttt{\mdseries\slshape x}}, \texttt{NrLeftBlocks} returns the number of left blocks of \mbox{\texttt{\mdseries\slshape x}}, i.e. the number of blocks of \mbox{\texttt{\mdseries\slshape x}} intersecting \texttt{[1 .. n]} non\texttt{\symbol{45}}trivially. 
\begin{Verbatim}[commandchars=!@|,fontsize=\small,frame=single,label=Example]
  !gapprompt@gap>| !gapinput@x := Bipartition([[1, 2, 3, 4, 5, 6, 8], [7, -2, -3],|
  !gapprompt@>| !gapinput@                     [-1, -4, -7, -8], [-5, -6]]);;|
  !gapprompt@gap>| !gapinput@NrLeftBlocks(x);|
  2
  !gapprompt@gap>| !gapinput@LeftBlocks(x);|
  <blocks: [ 1, 2, 3, 4, 5, 6, 8 ], [ 7* ]>
\end{Verbatim}
 }

 

\subsection{\textcolor{Chapter }{NrRightBlocks}}
\logpage{[ 3, 5, 9 ]}\nobreak
\hyperdef{L}{X86385A3C8662E1A7}{}
{\noindent\textcolor{FuncColor}{$\triangleright$\enspace\texttt{NrRightBlocks({\mdseries\slshape x})\index{NrRightBlocks@\texttt{NrRightBlocks}}
\label{NrRightBlocks}
}\hfill{\scriptsize (attribute)}}\\
\textbf{\indent Returns:\ }
A non\texttt{\symbol{45}}negative integer.



 When the argument is a bipartition \mbox{\texttt{\mdseries\slshape x}}, \texttt{NrRightBlocks} returns the number of right blocks of \mbox{\texttt{\mdseries\slshape x}}, i.e. the number of blocks of \mbox{\texttt{\mdseries\slshape x}} intersecting \texttt{[\texttt{\symbol{45}}n .. \texttt{\symbol{45}}1]} non\texttt{\symbol{45}}trivially. 
\begin{Verbatim}[commandchars=!@|,fontsize=\small,frame=single,label=Example]
  !gapprompt@gap>| !gapinput@x := Bipartition([[1, 2, 3, 4, 6, -2, -7], [5, -1, -3, -8],|
  !gapprompt@>| !gapinput@                     [7, -4, -6], [8], [-5]]);;|
  !gapprompt@gap>| !gapinput@RightBlocks(x);|
  <blocks: [ 1*, 3*, 8* ], [ 2*, 7* ], [ 4*, 6* ], [ 5 ]>
  !gapprompt@gap>| !gapinput@NrRightBlocks(x);|
  4
\end{Verbatim}
 }

 

\subsection{\textcolor{Chapter }{NrBlocks (for blocks)}}
\logpage{[ 3, 5, 10 ]}\nobreak
\hyperdef{L}{X8110B6557A98FB5C}{}
{\noindent\textcolor{FuncColor}{$\triangleright$\enspace\texttt{NrBlocks({\mdseries\slshape blocks})\index{NrBlocks@\texttt{NrBlocks}!for blocks}
\label{NrBlocks:for blocks}
}\hfill{\scriptsize (attribute)}}\\
\noindent\textcolor{FuncColor}{$\triangleright$\enspace\texttt{NrBlocks({\mdseries\slshape f})\index{NrBlocks@\texttt{NrBlocks}!for a bipartition}
\label{NrBlocks:for a bipartition}
}\hfill{\scriptsize (attribute)}}\\
\textbf{\indent Returns:\ }
A positive integer.



 If \mbox{\texttt{\mdseries\slshape blocks}} is some blocks or \mbox{\texttt{\mdseries\slshape f}} is a bipartition, then \texttt{NrBlocks} returns the number of blocks in \mbox{\texttt{\mdseries\slshape blocks}} or \mbox{\texttt{\mdseries\slshape f}}, respectively. 
\begin{Verbatim}[commandchars=!@|,fontsize=\small,frame=single,label=Example]
  !gapprompt@gap>| !gapinput@blocks := BLOCKS_NC([[-1, -2, -3, -4], [-5], [6]]);|
  <blocks: [ 1, 2, 3, 4 ], [ 5 ], [ 6* ]>
  !gapprompt@gap>| !gapinput@NrBlocks(blocks);|
  3
  !gapprompt@gap>| !gapinput@x := Bipartition([|
  !gapprompt@>| !gapinput@  [1, 5], [2, 4, -2, -4], [3, 6, -1, -5, -6], [-3]]);|
  <bipartition: [ 1, 5 ], [ 2, 4, -2, -4 ], [ 3, 6, -1, -5, -6 ],
   [ -3 ]>
  !gapprompt@gap>| !gapinput@NrBlocks(x);|
  4
\end{Verbatim}
 }

 

\subsection{\textcolor{Chapter }{DomainOfBipartition}}
\logpage{[ 3, 5, 11 ]}\nobreak
\hyperdef{L}{X8657EE2B79E1DD02}{}
{\noindent\textcolor{FuncColor}{$\triangleright$\enspace\texttt{DomainOfBipartition({\mdseries\slshape x})\index{DomainOfBipartition@\texttt{DomainOfBipartition}}
\label{DomainOfBipartition}
}\hfill{\scriptsize (attribute)}}\\
\textbf{\indent Returns:\ }
A list of positive integers.



 If \mbox{\texttt{\mdseries\slshape x}} is a bipartition, then \texttt{DomainOfBipartition} returns the domain of \mbox{\texttt{\mdseries\slshape x}}. The \emph{domain} of \mbox{\texttt{\mdseries\slshape x}} consists of those numbers \texttt{i} in \texttt{[1 .. n]} such that \texttt{i} is contained in a transverse block of \mbox{\texttt{\mdseries\slshape x}}, where \texttt{n} is the degree of \mbox{\texttt{\mdseries\slshape x}} (see \texttt{DegreeOfBipartition} (\ref{DegreeOfBipartition})). 
\begin{Verbatim}[commandchars=!@|,fontsize=\small,frame=single,label=Example]
  !gapprompt@gap>| !gapinput@x := Bipartition([[1, 2], [3, 4, 5, -5], [6, -6],|
  !gapprompt@>| !gapinput@                     [-1, -2, -3], [-4]]);|
  <bipartition: [ 1, 2 ], [ 3, 4, 5, -5 ], [ 6, -6 ], [ -1, -2, -3 ],
   [ -4 ]>
  !gapprompt@gap>| !gapinput@DomainOfBipartition(x);|
  [ 3, 4, 5, 6 ]
\end{Verbatim}
 }

 

\subsection{\textcolor{Chapter }{CodomainOfBipartition}}
\logpage{[ 3, 5, 12 ]}\nobreak
\hyperdef{L}{X84569A187A211332}{}
{\noindent\textcolor{FuncColor}{$\triangleright$\enspace\texttt{CodomainOfBipartition({\mdseries\slshape x})\index{CodomainOfBipartition@\texttt{CodomainOfBipartition}}
\label{CodomainOfBipartition}
}\hfill{\scriptsize (attribute)}}\\
\textbf{\indent Returns:\ }
A list of positive integers.



 If \mbox{\texttt{\mdseries\slshape x}} is a bipartition, then \texttt{CodomainOfBipartition} returns the codomain of \mbox{\texttt{\mdseries\slshape x}}. The \emph{codomain} of \mbox{\texttt{\mdseries\slshape x}} consists of those numbers \texttt{i} in \texttt{[\texttt{\symbol{45}}n .. \texttt{\symbol{45}}1]} such that \texttt{i} is contained in a transverse block of \mbox{\texttt{\mdseries\slshape x}}, where \texttt{n} is the degree of \mbox{\texttt{\mdseries\slshape x}} (see \texttt{DegreeOfBipartition} (\ref{DegreeOfBipartition})). 
\begin{Verbatim}[commandchars=!@|,fontsize=\small,frame=single,label=Example]
  !gapprompt@gap>| !gapinput@x := Bipartition([[1, 2], [3, 4, 5, -5], [6, -6],|
  !gapprompt@>| !gapinput@                     [-1, -2, -3], [-4]]);|
  <bipartition: [ 1, 2 ], [ 3, 4, 5, -5 ], [ 6, -6 ], [ -1, -2, -3 ],
   [ -4 ]>
  !gapprompt@gap>| !gapinput@CodomainOfBipartition(x);|
  [ -5, -6 ]
\end{Verbatim}
 }

 

\subsection{\textcolor{Chapter }{IsTransBipartition}}
\logpage{[ 3, 5, 13 ]}\nobreak
\hyperdef{L}{X79C556827A578509}{}
{\noindent\textcolor{FuncColor}{$\triangleright$\enspace\texttt{IsTransBipartition({\mdseries\slshape x})\index{IsTransBipartition@\texttt{IsTransBipartition}}
\label{IsTransBipartition}
}\hfill{\scriptsize (property)}}\\
\textbf{\indent Returns:\ }
\texttt{true} or \texttt{false}.



 If the bipartition \mbox{\texttt{\mdseries\slshape x}} defines a transformation, then \texttt{IsTransBipartition} returns \texttt{true}, and if not, then \texttt{false} is returned.

 A bipartition \mbox{\texttt{\mdseries\slshape x}} defines a transformation if and only if the number of left blocks equals the
number of transverse blocks and the number of right blocks equals the degree. 
\begin{Verbatim}[commandchars=!@|,fontsize=\small,frame=single,label=Example]
  !gapprompt@gap>| !gapinput@x := Bipartition([[1, 4, -2], [2, 5, -6], [3, -7],|
  !gapprompt@>| !gapinput@                     [6, 7, -9], [8, 9, -1], [10, -5],|
  !gapprompt@>| !gapinput@                     [-3], [-4], [-8], [-10]]);;|
  !gapprompt@gap>| !gapinput@IsTransBipartition(x);|
  true
  !gapprompt@gap>| !gapinput@x := Bipartition([[1, 4, -3, -6], [2, 5, -4, -5],|
  !gapprompt@>| !gapinput@                     [3, 6, -1], [-2]]);;|
  !gapprompt@gap>| !gapinput@IsTransBipartition(x);|
  false
  !gapprompt@gap>| !gapinput@Number(PartitionMonoid(3), IsTransBipartition);|
  27
\end{Verbatim}
 }

 

\subsection{\textcolor{Chapter }{IsDualTransBipartition}}
\logpage{[ 3, 5, 14 ]}\nobreak
\hyperdef{L}{X7F0B8ACC7C9A937F}{}
{\noindent\textcolor{FuncColor}{$\triangleright$\enspace\texttt{IsDualTransBipartition({\mdseries\slshape x})\index{IsDualTransBipartition@\texttt{IsDualTransBipartition}}
\label{IsDualTransBipartition}
}\hfill{\scriptsize (property)}}\\
\textbf{\indent Returns:\ }
\texttt{true} or \texttt{false}.



 If the star of the bipartition \mbox{\texttt{\mdseries\slshape x}} defines a transformation, then \texttt{IsDualTransBipartition} returns \texttt{true}, and if not, then \texttt{false} is returned.

 A bipartition is the dual of a transformation if and only if its number of
right blocks equals its number of transverse blocks and its number of left
blocks equals its degree. 
\begin{Verbatim}[commandchars=!@|,fontsize=\small,frame=single,label=Example]
  !gapprompt@gap>| !gapinput@x := Bipartition([[1, -8, -9], [2, -1, -4], [3],|
  !gapprompt@>| !gapinput@                     [4], [5, -10], [6, -2, -5], [7, -3],|
  !gapprompt@>| !gapinput@                     [8], [9, -6, -7], [10]]);;|
  !gapprompt@gap>| !gapinput@IsDualTransBipartition(x);|
  true
  !gapprompt@gap>| !gapinput@x := Bipartition([[1, 4, -3, -6], [2, 5, -4, -5],|
  !gapprompt@>| !gapinput@                     [3, 6, -1], [-2]]);;|
  !gapprompt@gap>| !gapinput@IsTransBipartition(x);|
  false
  !gapprompt@gap>| !gapinput@Number(PartitionMonoid(3), IsDualTransBipartition);|
  27
\end{Verbatim}
 }

 

\subsection{\textcolor{Chapter }{IsPermBipartition}}
\logpage{[ 3, 5, 15 ]}\nobreak
\hyperdef{L}{X8031B53E7D0ECCFA}{}
{\noindent\textcolor{FuncColor}{$\triangleright$\enspace\texttt{IsPermBipartition({\mdseries\slshape x})\index{IsPermBipartition@\texttt{IsPermBipartition}}
\label{IsPermBipartition}
}\hfill{\scriptsize (property)}}\\
\textbf{\indent Returns:\ }
\texttt{true} or \texttt{false}.



 If the bipartition \mbox{\texttt{\mdseries\slshape x}} defines a permutation, then \texttt{IsPermBipartition} returns \texttt{true}, and if not, then \texttt{false} is returned.

 A bipartition is a permutation if its numbers of left, right, and transverse
blocks all equal its degree. 
\begin{Verbatim}[commandchars=!@|,fontsize=\small,frame=single,label=Example]
  !gapprompt@gap>| !gapinput@x := Bipartition([|
  !gapprompt@>| !gapinput@  [1, 4, -1], [2, -3], [3, 6, -5], [5, -2, -4, -6]]);;|
  !gapprompt@gap>| !gapinput@IsPermBipartition(x);|
  false
  !gapprompt@gap>| !gapinput@x := Bipartition([[1, -3], [2, -4], [3, -6], [4, -1],|
  !gapprompt@>| !gapinput@                     [5, -5], [6, -2], [7, -8], [8, -7]]);;|
  !gapprompt@gap>| !gapinput@IsPermBipartition(x);|
  true
\end{Verbatim}
 }

 

\subsection{\textcolor{Chapter }{IsPartialPermBipartition}}
\logpage{[ 3, 5, 16 ]}\nobreak
\hyperdef{L}{X87C771D37B1FE95C}{}
{\noindent\textcolor{FuncColor}{$\triangleright$\enspace\texttt{IsPartialPermBipartition({\mdseries\slshape x})\index{IsPartialPermBipartition@\texttt{IsPartialPermBipartition}}
\label{IsPartialPermBipartition}
}\hfill{\scriptsize (property)}}\\
\textbf{\indent Returns:\ }
\texttt{true} or \texttt{false}.



 If the bipartition \mbox{\texttt{\mdseries\slshape x}} defines a partial permutation, then \texttt{IsPartialPermBipartition} returns \texttt{true}, and if not, then \texttt{false} is returned.

 A bipartition \mbox{\texttt{\mdseries\slshape x}} defines a partial permutation if and only if the numbers of left and right
blocks of \mbox{\texttt{\mdseries\slshape x}} equal the degree of \mbox{\texttt{\mdseries\slshape x}}. 
\begin{Verbatim}[commandchars=!@|,fontsize=\small,frame=single,label=Example]
  !gapprompt@gap>| !gapinput@x := Bipartition([|
  !gapprompt@>| !gapinput@  [1, 4, -1], [2, -3], [3, 6, -5], [5, -2, -4, -6]]);;|
  !gapprompt@gap>| !gapinput@IsPartialPermBipartition(x);|
  false
  !gapprompt@gap>| !gapinput@x := Bipartition([[1, -3], [2], [-4], [3, -6], [4, -1],|
  !gapprompt@>| !gapinput@                     [5, -5], [6, -2], [7, -8], [8, -7]]);;|
  !gapprompt@gap>| !gapinput@IsPermBipartition(x);|
  false
  !gapprompt@gap>| !gapinput@IsPartialPermBipartition(x);|
  true
\end{Verbatim}
 }

 

\subsection{\textcolor{Chapter }{IsBlockBijection}}
\logpage{[ 3, 5, 17 ]}\nobreak
\hyperdef{L}{X829494DF7FD6CFEC}{}
{\noindent\textcolor{FuncColor}{$\triangleright$\enspace\texttt{IsBlockBijection({\mdseries\slshape x})\index{IsBlockBijection@\texttt{IsBlockBijection}}
\label{IsBlockBijection}
}\hfill{\scriptsize (property)}}\\
\textbf{\indent Returns:\ }
\texttt{true} or \texttt{false}.



 If the bipartition \mbox{\texttt{\mdseries\slshape x}} induces a bijection from the quotient of \texttt{[1 .. n]} by the blocks of \mbox{\texttt{\mdseries\slshape f}} to the quotient of \texttt{[\texttt{\symbol{45}}n .. \texttt{\symbol{45}}1]} by the blocks of \mbox{\texttt{\mdseries\slshape f}}, then \texttt{IsBlockBijection} return \texttt{true}, and if not, then it returns \texttt{false}. 

 A bipartition is a block bijection if and only if its number of blocks, left
blocks and right blocks are equal. 
\begin{Verbatim}[commandchars=!@|,fontsize=\small,frame=single,label=Example]
  !gapprompt@gap>| !gapinput@x := Bipartition([[1, 4, 5, -2], [2, 3, -1], [6, -5, -6],|
  !gapprompt@>| !gapinput@                     [-3, -4]]);;|
  !gapprompt@gap>| !gapinput@IsBlockBijection(x);|
  false
  !gapprompt@gap>| !gapinput@x := Bipartition([[1, 2, -3], [3, -1, -2], [4, -4], [5, -5]]);;|
  !gapprompt@gap>| !gapinput@IsBlockBijection(x);|
  true
\end{Verbatim}
 }

 

\subsection{\textcolor{Chapter }{IsUniformBlockBijection}}
\logpage{[ 3, 5, 18 ]}\nobreak
\hyperdef{L}{X79D54AD8833B9551}{}
{\noindent\textcolor{FuncColor}{$\triangleright$\enspace\texttt{IsUniformBlockBijection({\mdseries\slshape x})\index{IsUniformBlockBijection@\texttt{IsUniformBlockBijection}}
\label{IsUniformBlockBijection}
}\hfill{\scriptsize (property)}}\\
\textbf{\indent Returns:\ }
\texttt{true} or \texttt{false}.



 If the bipartition \mbox{\texttt{\mdseries\slshape x}} is a block bijection where every block contains an equal number of positive
and negative entries, then \texttt{IsUniformBlockBijection} returns \texttt{true}, and otherwise it returns \texttt{false}. 
\begin{Verbatim}[commandchars=!@|,fontsize=\small,frame=single,label=Example]
  !gapprompt@gap>| !gapinput@x := Bipartition([[1, 2, -3, -4], [3, -5], [4, -6],|
  !gapprompt@>| !gapinput@[5, -7], [6, -8], [7, -9], [8, -1], [9, -2]]);;|
  !gapprompt@gap>| !gapinput@IsBlockBijection(x);|
  true
  !gapprompt@gap>| !gapinput@x := Bipartition([[1, 2, -3], [3, -1, -2], [4, -4],|
  !gapprompt@>| !gapinput@[5, -5]]);;|
  !gapprompt@gap>| !gapinput@IsUniformBlockBijection(x);|
  false
\end{Verbatim}
 }

 

\subsection{\textcolor{Chapter }{CanonicalBlocks}}
\logpage{[ 3, 5, 19 ]}\nobreak
\hyperdef{L}{X7B87B9B081FF88BB}{}
{\noindent\textcolor{FuncColor}{$\triangleright$\enspace\texttt{CanonicalBlocks({\mdseries\slshape blocks})\index{CanonicalBlocks@\texttt{CanonicalBlocks}}
\label{CanonicalBlocks}
}\hfill{\scriptsize (attribute)}}\\
\textbf{\indent Returns:\ }
Blocks of a bipartition.



 If \mbox{\texttt{\mdseries\slshape blocks}} is the blocks of a bipartition, then the function \texttt{CanonicalBlocks} returns a canonical representative of \mbox{\texttt{\mdseries\slshape blocks}}. 

 In particular, let \texttt{C(n)} be a largest class such that any element of \texttt{C(n)} is blocks of a bipartition of degree \texttt{n} and such that for every pair of elements \texttt{x} and \texttt{y} of \texttt{C(n)} the number of signed, and similarly unsigned, blocks of any given size in both \texttt{x} and \texttt{y} are the same. Then \texttt{CanonicalBlocks} returns a canonical representative of a class \texttt{C(n)} containing \mbox{\texttt{\mdseries\slshape blocks}} where \texttt{n} is the degree of \mbox{\texttt{\mdseries\slshape blocks}}. 
\begin{Verbatim}[commandchars=!@|,fontsize=\small,frame=single,label=Example]
  !gapprompt@gap>| !gapinput@B := BLOCKS_NC([[-1, -3], [2, 4, 7], [5, 6]]);|
  <blocks: [ 1, 3 ], [ 2*, 4*, 7* ], [ 5*, 6* ]>
  !gapprompt@gap>| !gapinput@CanonicalBlocks(B);|
  <blocks: [ 1*, 2* ], [ 3*, 4*, 5* ], [ 6, 7 ]>
\end{Verbatim}
 }

 

\subsection{\textcolor{Chapter }{IsIrreducibleBipartition}}
\logpage{[ 3, 5, 20 ]}\nobreak
\hyperdef{L}{X78999D6F8231AE79}{}
{\noindent\textcolor{FuncColor}{$\triangleright$\enspace\texttt{IsIrreducibleBipartition({\mdseries\slshape x})\index{IsIrreducibleBipartition@\texttt{IsIrreducibleBipartition}}
\label{IsIrreducibleBipartition}
}\hfill{\scriptsize (operation)}}\\
\textbf{\indent Returns:\ }
\texttt{true} or \texttt{false}.



 A bipartition is \emph{irreducible} if it is not the tensor product of bipartitions of strictly smaller degree.
This function returns \texttt{true} if \mbox{\texttt{\mdseries\slshape x}} is irreducible and \texttt{false} if it is reducible. See also: \texttt{TensorBipartitions} (\ref{TensorBipartitions:for a pair of bipartitions}) and \texttt{IrreducibleComponentsOfBipartition} (\ref{IrreducibleComponentsOfBipartition}).  
\begin{Verbatim}[commandchars=!@|,fontsize=\small,frame=single,label=Example]
  !gapprompt@gap>| !gapinput@IsIrreducibleBipartition(TensorBipartitions(|
  !gapprompt@>| !gapinput@Bipartition([[1, 2, -2], [-1]]),|
  !gapprompt@>| !gapinput@Bipartition([[1, 2, 3, -2], [-1], [-3]])));|
  false
  !gapprompt@gap>| !gapinput@IsIrreducibleBipartition(Bipartition([[1, 2, 3, -2], [-1], [-3]]));|
  true
\end{Verbatim}
 }

 }

   
\section{\textcolor{Chapter }{Creating blocks and their attributes}}\label{section-blocks}
\logpage{[ 3, 6, 0 ]}
\hyperdef{L}{X87684C148592F831}{}
{
  As described above the left and right blocks of a bipartition characterise
Green's $\mathcal{R}$\texttt{\symbol{45}} and $\mathcal{L}$\texttt{\symbol{45}}relation of the partition monoid; see \texttt{LeftBlocks} (\ref{LeftBlocks}) and \texttt{RightBlocks} (\ref{RightBlocks}). The left or right blocks of a bipartition are \textsf{GAP} objects in their own right. 

 In this section, we describe the functions in the \textsf{Semigroups} package for creating and manipulating the left or right blocks of a
bipartition. 

\subsection{\textcolor{Chapter }{IsBlocks}}
\logpage{[ 3, 6, 1 ]}\nobreak
\hyperdef{L}{X7D77092078EC860C}{}
{\noindent\textcolor{FuncColor}{$\triangleright$\enspace\texttt{IsBlocks({\mdseries\slshape obj})\index{IsBlocks@\texttt{IsBlocks}}
\label{IsBlocks}
}\hfill{\scriptsize (Category)}}\\
\textbf{\indent Returns:\ }
\texttt{true} or \texttt{false}.



 Every blocks object in \textsf{GAP} belongs to the category \texttt{IsBlocks}. Basic operations for blocks are \texttt{ExtRepOfObj} (\ref{ExtRepOfObj:for a blocks}), \texttt{RankOfBlocks} (\ref{RankOfBlocks}), \texttt{DegreeOfBlocks} (\ref{DegreeOfBlocks}), \texttt{OnRightBlocks} (\ref{OnRightBlocks}), and \texttt{OnLeftBlocks} (\ref{OnLeftBlocks}). }

 

\subsection{\textcolor{Chapter }{BLOCKS{\textunderscore}NC}}
\logpage{[ 3, 6, 2 ]}\nobreak
\hyperdef{L}{X81302B217DCAAE6F}{}
{\noindent\textcolor{FuncColor}{$\triangleright$\enspace\texttt{BLOCKS{\textunderscore}NC({\mdseries\slshape classes})\index{BLOCKS{\textunderscore}NC@\texttt{BLOCKS{\textunderscore}NC}}
\label{BLOCKSuScoreNC}
}\hfill{\scriptsize (function)}}\\
\textbf{\indent Returns:\ }
A blocks.



 This function makes it possible to create a \textsf{GAP} object corresponding to the left or right blocks of a bipartition without
reference to any bipartitions. 

 \texttt{BLOCKS{\textunderscore}NC} returns the blocks with equivalence classes \mbox{\texttt{\mdseries\slshape classes}}, which should be a list of duplicate\texttt{\symbol{45}}free lists consisting
solely of positive or negative integers, where the union of the absolute
values of the lists is \texttt{[1 .. n]} for some \texttt{n}. The blocks with positive entries correspond to transverse blocks and the
classes with negative entries correspond to non\texttt{\symbol{45}}transverse
blocks. 

 This method function does not check that its arguments are valid, and should
be used with caution. 
\begin{Verbatim}[commandchars=!@|,fontsize=\small,frame=single,label=Example]
  !gapprompt@gap>| !gapinput@BLOCKS_NC([[1], [2], [-3, -6], [-4, -5]]);|
  <blocks: [ 1* ], [ 2* ], [ 3, 6 ], [ 4, 5 ]>
\end{Verbatim}
 }

 

\subsection{\textcolor{Chapter }{ExtRepOfObj (for a blocks)}}
\logpage{[ 3, 6, 3 ]}\nobreak
\hyperdef{L}{X7D2CB12279623CE2}{}
{\noindent\textcolor{FuncColor}{$\triangleright$\enspace\texttt{ExtRepOfObj({\mdseries\slshape blocks})\index{ExtRepOfObj@\texttt{ExtRepOfObj}!for a blocks}
\label{ExtRepOfObj:for a blocks}
}\hfill{\scriptsize (operation)}}\\
\textbf{\indent Returns:\ }
A list of integers.



 If \texttt{n} is the degree of a bipartition with left or right blocks \mbox{\texttt{\mdseries\slshape blocks}}, then \texttt{ExtRepOfObj} returns the partition corresponding to \mbox{\texttt{\mdseries\slshape blocks}} as a sorted list of duplicate\texttt{\symbol{45}}free lists. 
\begin{Verbatim}[commandchars=!@|,fontsize=\small,frame=single,label=Example]
  !gapprompt@gap>| !gapinput@blocks := BLOCKS_NC([[1, 6], [2, 3, 7], [4, 5], [-8]]);;|
  !gapprompt@gap>| !gapinput@ExtRepOfObj(blocks);|
  [ [ 1, 6 ], [ 2, 3, 7 ], [ 4, 5 ], [ -8 ] ]
\end{Verbatim}
 }

 

\subsection{\textcolor{Chapter }{RankOfBlocks}}
\logpage{[ 3, 6, 4 ]}\nobreak
\hyperdef{L}{X787D22AE7FA69239}{}
{\noindent\textcolor{FuncColor}{$\triangleright$\enspace\texttt{RankOfBlocks({\mdseries\slshape blocks})\index{RankOfBlocks@\texttt{RankOfBlocks}}
\label{RankOfBlocks}
}\hfill{\scriptsize (attribute)}}\\
\noindent\textcolor{FuncColor}{$\triangleright$\enspace\texttt{NrTransverseBlocks({\mdseries\slshape blocks})\index{NrTransverseBlocks@\texttt{NrTransverseBlocks}!for blocks}
\label{NrTransverseBlocks:for blocks}
}\hfill{\scriptsize (attribute)}}\\
\textbf{\indent Returns:\ }
A non\texttt{\symbol{45}}negative integer.



 When the argument \mbox{\texttt{\mdseries\slshape blocks}} is the left or right blocks of a bipartition, \texttt{RankOfBlocks} returns the number of blocks of \mbox{\texttt{\mdseries\slshape blocks}} containing only positive entries, i.e. the number of transverse blocks in \mbox{\texttt{\mdseries\slshape blocks}}. 

 \texttt{NrTransverseBlocks} is a synonym of \texttt{RankOfBlocks} in this context. 
\begin{Verbatim}[commandchars=!@|,fontsize=\small,frame=single,label=Example]
  !gapprompt@gap>| !gapinput@blocks := BLOCKS_NC([[-1, -2, -4, -6], [3, 10, 12], [5, 7],|
  !gapprompt@>| !gapinput@                       [8], [9], [-11]]);;|
  !gapprompt@gap>| !gapinput@RankOfBlocks(blocks);|
  4
\end{Verbatim}
 }

 

\subsection{\textcolor{Chapter }{DegreeOfBlocks}}
\logpage{[ 3, 6, 5 ]}\nobreak
\hyperdef{L}{X8527DC6A8771C2BE}{}
{\noindent\textcolor{FuncColor}{$\triangleright$\enspace\texttt{DegreeOfBlocks({\mdseries\slshape blocks})\index{DegreeOfBlocks@\texttt{DegreeOfBlocks}}
\label{DegreeOfBlocks}
}\hfill{\scriptsize (attribute)}}\\
\textbf{\indent Returns:\ }
A non\texttt{\symbol{45}}negative integer.



 The degree of \mbox{\texttt{\mdseries\slshape blocks}} is the number of points \texttt{n} where it is defined, i.e. the union of the blocks in \mbox{\texttt{\mdseries\slshape blocks}} will be \texttt{[1 .. n]} after taking the absolute value of every element. 
\begin{Verbatim}[commandchars=!@|,fontsize=\small,frame=single,label=Example]
  !gapprompt@gap>| !gapinput@blocks := BLOCKS_NC([[-1, -11], [2], [3, 5, 6, 7], [4, 8], [9, 10],|
  !gapprompt@>| !gapinput@                       [12]]);;|
  !gapprompt@gap>| !gapinput@DegreeOfBlocks(blocks);|
  12
\end{Verbatim}
 }

 

\subsection{\textcolor{Chapter }{ProjectionFromBlocks}}
\logpage{[ 3, 6, 6 ]}\nobreak
\hyperdef{L}{X815D99A983B2355F}{}
{\noindent\textcolor{FuncColor}{$\triangleright$\enspace\texttt{ProjectionFromBlocks({\mdseries\slshape blocks})\index{ProjectionFromBlocks@\texttt{ProjectionFromBlocks}}
\label{ProjectionFromBlocks}
}\hfill{\scriptsize (attribute)}}\\
\textbf{\indent Returns:\ }
A bipartition.



 When the argument \mbox{\texttt{\mdseries\slshape blocks}} is the left or right blocks of a bipartition, this operation returns the
unique bipartition whose left and right blocks are equal to \mbox{\texttt{\mdseries\slshape blocks}}. 

 If \mbox{\texttt{\mdseries\slshape blocks}} is the left blocks of a bipartition \texttt{x}, then this operation returns a bipartition equal to the left projection of \texttt{x}. The analogous statement holds when \mbox{\texttt{\mdseries\slshape blocks}} is the right blocks of a bipartition. 
\begin{Verbatim}[commandchars=!@|,fontsize=\small,frame=single,label=Example]
  !gapprompt@gap>| !gapinput@x := Bipartition([[1], [2, -2, -3], [3], [-1]]);|
  <bipartition: [ 1 ], [ 2, -2, -3 ], [ 3 ], [ -1 ]>
  !gapprompt@gap>| !gapinput@ProjectionFromBlocks(LeftBlocks(x));|
  <bipartition: [ 1 ], [ 2, -2 ], [ 3 ], [ -1 ], [ -3 ]>
  !gapprompt@gap>| !gapinput@LeftProjection(x);|
  <bipartition: [ 1 ], [ 2, -2 ], [ 3 ], [ -1 ], [ -3 ]>
  !gapprompt@gap>| !gapinput@ProjectionFromBlocks(RightBlocks(x));|
  <bipartition: [ 1 ], [ 2, 3, -2, -3 ], [ -1 ]>
  !gapprompt@gap>| !gapinput@RightProjection(x);|
  <bipartition: [ 1 ], [ 2, 3, -2, -3 ], [ -1 ]>
\end{Verbatim}
 }

 }

 
\section{\textcolor{Chapter }{Actions on blocks}}\logpage{[ 3, 7, 0 ]}
\hyperdef{L}{X7A45E0067F344683}{}
{
  Bipartitions act on left and right blocks in several ways, which are described
in this section. 

\subsection{\textcolor{Chapter }{OnRightBlocks}}
\logpage{[ 3, 7, 1 ]}\nobreak
\hyperdef{L}{X7B701DA37F75E77B}{}
{\noindent\textcolor{FuncColor}{$\triangleright$\enspace\texttt{OnRightBlocks({\mdseries\slshape blocks, x})\index{OnRightBlocks@\texttt{OnRightBlocks}}
\label{OnRightBlocks}
}\hfill{\scriptsize (operation)}}\\
\textbf{\indent Returns:\ }
The blocks of a bipartition.



 \texttt{OnRightBlocks} returns the right blocks of the product \texttt{g * \mbox{\texttt{\mdseries\slshape x}}} where \texttt{g} is any bipartition whose right blocks are equal to \mbox{\texttt{\mdseries\slshape blocks}}. 
\begin{Verbatim}[commandchars=!@|,fontsize=\small,frame=single,label=Example]
  !gapprompt@gap>| !gapinput@x := Bipartition([[1, 4, 5, 8], [2, 3, 7], [6, -3, -4, -5],|
  !gapprompt@>| !gapinput@                     [-1, -2, -6], [-7, -8]]);;|
  !gapprompt@gap>| !gapinput@y := Bipartition([[1, 5], [2, 4, 8, -2], [3, 6, 7, -3, -4],|
  !gapprompt@>| !gapinput@                     [-1, -6, -8], [-5, -7]]);;|
  !gapprompt@gap>| !gapinput@RightBlocks(y * x);|
  <blocks: [ 1, 2, 6 ], [ 3*, 4*, 5* ], [ 7, 8 ]>
  !gapprompt@gap>| !gapinput@OnRightBlocks(RightBlocks(y), x);|
  <blocks: [ 1, 2, 6 ], [ 3*, 4*, 5* ], [ 7, 8 ]>
\end{Verbatim}
 }

 

\subsection{\textcolor{Chapter }{OnLeftBlocks}}
\logpage{[ 3, 7, 2 ]}\nobreak
\hyperdef{L}{X7A5A4AF57BEA2313}{}
{\noindent\textcolor{FuncColor}{$\triangleright$\enspace\texttt{OnLeftBlocks({\mdseries\slshape blocks, x})\index{OnLeftBlocks@\texttt{OnLeftBlocks}}
\label{OnLeftBlocks}
}\hfill{\scriptsize (operation)}}\\
\textbf{\indent Returns:\ }
The blocks of a bipartition.



 \texttt{OnLeftBlocks} returns the left blocks of the product \texttt{\mbox{\texttt{\mdseries\slshape x}} * y} where \texttt{y} is any bipartition whose left blocks are equal to \mbox{\texttt{\mdseries\slshape blocks}}. 
\begin{Verbatim}[commandchars=!@|,fontsize=\small,frame=single,label=Example]
  !gapprompt@gap>| !gapinput@x := Bipartition([[1, 5, 7, -1, -3, -4, -6], [2, 3, 6, 8],|
  !gapprompt@>| !gapinput@                     [4, -2, -5, -8], [-7]]);;|
  !gapprompt@gap>| !gapinput@y := Bipartition([[1, 3, -4, -5], [2, 4, 5, 8], [6, -1, -3],|
  !gapprompt@>| !gapinput@                     [7, -2, -6, -7, -8]]);;|
  !gapprompt@gap>| !gapinput@LeftBlocks(x * y);|
  <blocks: [ 1*, 4*, 5*, 7* ], [ 2, 3, 6, 8 ]>
  !gapprompt@gap>| !gapinput@OnLeftBlocks(LeftBlocks(y), x);|
  <blocks: [ 1*, 4*, 5*, 7* ], [ 2, 3, 6, 8 ]>
\end{Verbatim}
 }

 }

   
\section{\textcolor{Chapter }{ Semigroups of bipartitions }}\logpage{[ 3, 8, 0 ]}
\hyperdef{L}{X876C963F830719E2}{}
{
  Semigroups and monoids of bipartitions can be created in the usual way in \textsf{GAP} using the functions \texttt{Semigroup} (\textbf{Reference: Semigroup}) and \texttt{Monoid} (\textbf{Reference: Monoid}); see Chapter \ref{Semigroups and monoids defined by generating sets} for more details. 

 It is possible to create inverse semigroups and monoids of bipartitions using \texttt{InverseSemigroup} (\textbf{Reference: InverseSemigroup}) and \texttt{InverseMonoid} (\textbf{Reference: InverseMonoid}) when the argument is a collection of block bijections or partial perm
bipartions; see \texttt{IsBlockBijection} (\ref{IsBlockBijection}) and \texttt{IsPartialPermBipartition} (\ref{IsPartialPermBipartition}). Note that every bipartition semigroup in \textsf{Semigroups} is finite. 

\subsection{\textcolor{Chapter }{IsBipartitionSemigroup}}
\logpage{[ 3, 8, 1 ]}\nobreak
\hyperdef{L}{X810BFF647C4E191E}{}
{\noindent\textcolor{FuncColor}{$\triangleright$\enspace\texttt{IsBipartitionSemigroup({\mdseries\slshape S})\index{IsBipartitionSemigroup@\texttt{IsBipartitionSemigroup}}
\label{IsBipartitionSemigroup}
}\hfill{\scriptsize (filter)}}\\
\noindent\textcolor{FuncColor}{$\triangleright$\enspace\texttt{IsBipartitionMonoid({\mdseries\slshape S})\index{IsBipartitionMonoid@\texttt{IsBipartitionMonoid}}
\label{IsBipartitionMonoid}
}\hfill{\scriptsize (filter)}}\\
\textbf{\indent Returns:\ }
\texttt{true} or \texttt{false}.



 A \emph{bipartition semigroup} is simply a semigroup consisting of bipartitions. An object \mbox{\texttt{\mdseries\slshape obj}} is a bipartition semigroup in \textsf{GAP} if it satisfies \texttt{IsSemigroup} (\textbf{Reference: IsSemigroup}) and \texttt{IsBipartitionCollection} (\ref{IsBipartitionCollection}).

 A \emph{bipartition monoid} is a monoid consisting of bipartitions. An object \mbox{\texttt{\mdseries\slshape obj}} is a bipartition monoid in \textsf{GAP} if it satisfies \texttt{IsMonoid} (\textbf{Reference: IsMonoid}) and \texttt{IsBipartitionCollection} (\ref{IsBipartitionCollection}).

 Note that it is possible for a bipartition semigroup to have a multiplicative
neutral element (i.e. an identity element) but not to satisfy \texttt{IsBipartitionMonoid}. For example, 
\begin{Verbatim}[commandchars=!@|,fontsize=\small,frame=single,label=Example]
  !gapprompt@gap>| !gapinput@x := Bipartition([|
  !gapprompt@>| !gapinput@[1, 4, -2], [2, 5, -6], [3, -7], [6, 7, -9], [8, 9, -1],|
  !gapprompt@>| !gapinput@[10, -5], [-3], [-4], [-8], [-10]]);;|
  !gapprompt@gap>| !gapinput@S := Semigroup(x, One(x));|
  <commutative bipartition monoid of degree 10 with 1 generator>
  !gapprompt@gap>| !gapinput@IsMonoid(S);|
  true
  !gapprompt@gap>| !gapinput@IsBipartitionMonoid(S);|
  true
  !gapprompt@gap>| !gapinput@S := Semigroup([|
  !gapprompt@>| !gapinput@Bipartition([|
  !gapprompt@>| !gapinput@  [1, -3], [2, -8], [3, 8, -1], [4, -4], [5, -5], [6, -6],|
  !gapprompt@>| !gapinput@  [7, -7], [9, 10, -10], [-2], [-9]]),|
  !gapprompt@>| !gapinput@Bipartition([|
  !gapprompt@>| !gapinput@  [1, -1], [2, -2], [3, -3], [4, -4], [5, -5], [6, -6],|
  !gapprompt@>| !gapinput@  [7, -7], [8, -8], [9, 10, -10], [-9]])]);;|
  !gapprompt@gap>| !gapinput@One(S);|
  fail
  !gapprompt@gap>| !gapinput@MultiplicativeNeutralElement(S);|
  <bipartition: [ 1, -1 ], [ 2, -2 ], [ 3, -3 ], [ 4, -4 ], [ 5, -5 ],
   [ 6, -6 ], [ 7, -7 ], [ 8, -8 ], [ 9, 10, -10 ], [ -9 ]>
  !gapprompt@gap>| !gapinput@IsMonoid(S);|
  false
\end{Verbatim}
 In this example \texttt{S} cannot be converted into a monoid using \texttt{AsMonoid} (\textbf{Reference: AsMonoid}) since the \texttt{One} (\textbf{Reference: One}) of any element in \texttt{S} differs from the multiplicative neutral element. 

 For more details see \texttt{IsMagmaWithOne} (\textbf{Reference: IsMagmaWithOne}). }

 

\subsection{\textcolor{Chapter }{IsBlockBijectionSemigroup}}
\logpage{[ 3, 8, 2 ]}\nobreak
\hyperdef{L}{X80C37124794636F3}{}
{\noindent\textcolor{FuncColor}{$\triangleright$\enspace\texttt{IsBlockBijectionSemigroup({\mdseries\slshape S})\index{IsBlockBijectionSemigroup@\texttt{IsBlockBijectionSemigroup}}
\label{IsBlockBijectionSemigroup}
}\hfill{\scriptsize (property)}}\\
\noindent\textcolor{FuncColor}{$\triangleright$\enspace\texttt{IsBlockBijectionMonoid({\mdseries\slshape S})\index{IsBlockBijectionMonoid@\texttt{IsBlockBijectionMonoid}}
\label{IsBlockBijectionMonoid}
}\hfill{\scriptsize (filter)}}\\
\textbf{\indent Returns:\ }
\texttt{true} or \texttt{false}.



 A \emph{block bijection semigroup} is simply a semigroup consisting of block bijections. A \emph{block bijection monoid} is a monoid consisting of block bijections.

 An object in \textsf{GAP} is a block bijection monoid if it satisfies \texttt{IsMonoid} (\textbf{Reference: IsMonoid}) and \texttt{IsBlockBijectionSemigroup}.

 See \texttt{IsBlockBijection} (\ref{IsBlockBijection}). }

 

\subsection{\textcolor{Chapter }{IsPartialPermBipartitionSemigroup}}
\logpage{[ 3, 8, 3 ]}\nobreak
\hyperdef{L}{X79A706A582ABE558}{}
{\noindent\textcolor{FuncColor}{$\triangleright$\enspace\texttt{IsPartialPermBipartitionSemigroup({\mdseries\slshape S})\index{IsPartialPermBipartitionSemigroup@\texttt{IsPartialPermBipartitionSemigroup}}
\label{IsPartialPermBipartitionSemigroup}
}\hfill{\scriptsize (property)}}\\
\noindent\textcolor{FuncColor}{$\triangleright$\enspace\texttt{IsPartialPermBipartitionMonoid({\mdseries\slshape S})\index{IsPartialPermBipartitionMonoid@\texttt{IsPartialPermBipartitionMonoid}}
\label{IsPartialPermBipartitionMonoid}
}\hfill{\scriptsize (filter)}}\\
\textbf{\indent Returns:\ }
\texttt{true} or \texttt{false}.



 A \emph{partial perm bipartition semigroup} is simply a semigroup consisting of partial perm bipartitions. A \emph{partial perm bipartition monoid} is a monoid consisting of partial perm bipartitions.

 An object in \textsf{GAP} is a partial perm bipartition monoid if it satisfies \texttt{IsMonoid} (\textbf{Reference: IsMonoid}) and \texttt{IsPartialPermBipartitionSemigroup}.

 See \texttt{IsPartialPermBipartition} (\ref{IsPartialPermBipartition}). }

 

\subsection{\textcolor{Chapter }{IsPermBipartitionGroup}}
\logpage{[ 3, 8, 4 ]}\nobreak
\hyperdef{L}{X7DEE07577D7379AC}{}
{\noindent\textcolor{FuncColor}{$\triangleright$\enspace\texttt{IsPermBipartitionGroup({\mdseries\slshape S})\index{IsPermBipartitionGroup@\texttt{IsPermBipartitionGroup}}
\label{IsPermBipartitionGroup}
}\hfill{\scriptsize (property)}}\\
\textbf{\indent Returns:\ }
\texttt{true} or \texttt{false}.



 A \emph{perm bipartition group} is simply a semigroup consisting of perm bipartitions.

 See \texttt{IsPermBipartition} (\ref{IsPermBipartition}).

 }

 

\subsection{\textcolor{Chapter }{DegreeOfBipartitionSemigroup}}
\logpage{[ 3, 8, 5 ]}\nobreak
\hyperdef{L}{X8162E2BB7CF144F5}{}
{\noindent\textcolor{FuncColor}{$\triangleright$\enspace\texttt{DegreeOfBipartitionSemigroup({\mdseries\slshape S})\index{DegreeOfBipartitionSemigroup@\texttt{DegreeOfBipartitionSemigroup}}
\label{DegreeOfBipartitionSemigroup}
}\hfill{\scriptsize (attribute)}}\\
\textbf{\indent Returns:\ }
A non\texttt{\symbol{45}}negative integer.



 The \emph{degree} of a bipartition semigroup \mbox{\texttt{\mdseries\slshape S}} is just the degree of any (and every) element of \mbox{\texttt{\mdseries\slshape S}}. 
\begin{Verbatim}[commandchars=!@|,fontsize=\small,frame=single,label=Example]
  !gapprompt@gap>| !gapinput@DegreeOfBipartitionSemigroup(JonesMonoid(8));|
  8
\end{Verbatim}
 }

 }

   }

 
\chapter{\textcolor{Chapter }{ Partitioned binary relations (PBRs) }}\label{Partitioned binary relations (PBRs)}
\logpage{[ 4, 0, 0 ]}
\hyperdef{L}{X85A717D1790B7BB5}{}
{
  In this chapter we describe the functions in \textsf{Semigroups} for creating and manipulating partitioned binary relations, henceforth
referred to by their acronym PBRs. We begin by describing what these objects
are. 

 PBRs were introduced in the paper \cite{Martin2011aa} as, roughly speaking, the maximum generalization of bipartitions and related
objects. Although, mathematically, bipartitions are a special type of PBR, in \textsf{Semigroups} bipartitions and PBRs are currently distinct types of objects. It is possible
to change the representation from bipartition to PBR, and from PBR to
bipartition, when appropriate; see Section \ref{Changing the representation of a PBR} for more details. The reason for this distinct is largely historical,
bipartition appeared in the literature, and in the \textsf{Semigroups} package, before PBRs.    
\section{\textcolor{Chapter }{The family and categories of PBRs}}\logpage{[ 4, 1, 0 ]}
\hyperdef{L}{X7C40DA67826FF873}{}
{
  

\subsection{\textcolor{Chapter }{IsPBR}}
\logpage{[ 4, 1, 1 ]}\nobreak
\hyperdef{L}{X82CCBADC80AE2D15}{}
{\noindent\textcolor{FuncColor}{$\triangleright$\enspace\texttt{IsPBR({\mdseries\slshape obj})\index{IsPBR@\texttt{IsPBR}}
\label{IsPBR}
}\hfill{\scriptsize (Category)}}\\
\textbf{\indent Returns:\ }
 \texttt{true} or \texttt{false}. 



 Every PBR in \textsf{GAP} belongs to the category \texttt{IsPBR}. Basic operations for PBRs are \texttt{DegreeOfPBR} (\ref{DegreeOfPBR}), \texttt{ExtRepOfObj} (\ref{ExtRepOfObj:for a PBR}), \texttt{PBRNumber} (\ref{PBRNumber}), \texttt{NumberPBR} (\ref{NumberPBR}), \texttt{StarOp} (\ref{StarOp:for a PBR}), and multiplication of two PBRs of equal degree is via \texttt{*}. }

 

\subsection{\textcolor{Chapter }{IsPBRCollection}}
\logpage{[ 4, 1, 2 ]}\nobreak
\hyperdef{L}{X854A9CEA7AC14C0A}{}
{\noindent\textcolor{FuncColor}{$\triangleright$\enspace\texttt{IsPBRCollection({\mdseries\slshape obj})\index{IsPBRCollection@\texttt{IsPBRCollection}}
\label{IsPBRCollection}
}\hfill{\scriptsize (Category)}}\\
\noindent\textcolor{FuncColor}{$\triangleright$\enspace\texttt{IsPBRCollColl({\mdseries\slshape obj})\index{IsPBRCollColl@\texttt{IsPBRCollColl}}
\label{IsPBRCollColl}
}\hfill{\scriptsize (Category)}}\\
\textbf{\indent Returns:\ }
\texttt{true} or \texttt{false}.



 Every collection of PBRs belongs to the category \texttt{IsPBRCollection}. For example, PBR semigroups belong to \texttt{IsPBRCollection}. 

 Every collection of collections of PBRs belongs to \texttt{IsPBRCollColl}. For example, a list of PBR semigroups belongs to \texttt{IsPBRCollColl}. }

 }

   
\section{\textcolor{Chapter }{Creating PBRs}}\label{Creating PBRs}
\logpage{[ 4, 2, 0 ]}
\hyperdef{L}{X8758C4FB81D2C2A1}{}
{
  There are several ways of creating PBRs in \textsf{GAP}, which are described in this section. 

\subsection{\textcolor{Chapter }{PBR}}
\logpage{[ 4, 2, 1 ]}\nobreak
\hyperdef{L}{X82A8646F7C70CF3B}{}
{\noindent\textcolor{FuncColor}{$\triangleright$\enspace\texttt{PBR({\mdseries\slshape left, right})\index{PBR@\texttt{PBR}}
\label{PBR}
}\hfill{\scriptsize (operation)}}\\
\textbf{\indent Returns:\ }
A PBR.



 The arguments \mbox{\texttt{\mdseries\slshape left}} and \mbox{\texttt{\mdseries\slshape right}} of this function should each be a list of length \texttt{n} whose entries are lists of integers in the ranges \texttt{[\texttt{\symbol{45}}n .. \texttt{\symbol{45}}1]} and \texttt{[1 .. n]} for some \texttt{n} greater than 0. 

 Given such an argument, \texttt{PBR} returns the PBR \texttt{x} where: 
\begin{itemize}
\item  for each \texttt{i} in the range \texttt{[1 .. n]} there is an edge from \texttt{i} to every \texttt{j} in \mbox{\texttt{\mdseries\slshape left[i]}}; 
\item  for each \texttt{i} in the range \texttt{[\texttt{\symbol{45}}n .. \texttt{\symbol{45}}1]} there is an edge from \texttt{i} to every \texttt{j} in \mbox{\texttt{\mdseries\slshape right[\texttt{\symbol{45}}i]}}; 
\end{itemize}
 \texttt{PBR} returns an error if the argument does not define a PBR. 
\begin{Verbatim}[commandchars=!@|,fontsize=\small,frame=single,label=Example]
  !gapprompt@gap>| !gapinput@PBR([[-3, -2, -1, 2, 3], [-1], [-3, -2, 1, 2]],|
  !gapprompt@>| !gapinput@       [[-2, -1, 1, 2, 3], [3], [-3, -2, -1, 1, 3]]);|
  PBR([ [ -3, -2, -1, 2, 3 ], [ -1 ], [ -3, -2, 1, 2 ] ],
    [ [ -2, -1, 1, 2, 3 ], [ 3 ], [ -3, -2, -1, 1, 3 ] ])
\end{Verbatim}
 }

 

\subsection{\textcolor{Chapter }{RandomPBR}}
\logpage{[ 4, 2, 2 ]}\nobreak
\hyperdef{L}{X82FE736F7F11B157}{}
{\noindent\textcolor{FuncColor}{$\triangleright$\enspace\texttt{RandomPBR({\mdseries\slshape n[, p]})\index{RandomPBR@\texttt{RandomPBR}}
\label{RandomPBR}
}\hfill{\scriptsize (operation)}}\\
\textbf{\indent Returns:\ }
A PBR.



 If \mbox{\texttt{\mdseries\slshape n}} is a positive integer and \mbox{\texttt{\mdseries\slshape p}} is an float between \texttt{0} and \texttt{1}, then \texttt{RandomPBR} returns a random PBR of degree \mbox{\texttt{\mdseries\slshape n}} where the probability of there being an edge from \texttt{i} to \texttt{j} is approximately \texttt{p}. 

 If the optional second argument is not present, then a random value \mbox{\texttt{\mdseries\slshape p}} is used (chosen with uniform probability). 
\begin{Verbatim}[commandchars=!@|,fontsize=\small,frame=single,label=Example]
  !gapprompt@gap>| !gapinput@RandomPBR(6);|
  PBR(
    [ [ -5, 1, 2, 3 ], [ -6, -3, -1, 2, 5 ], [ -5, -2, 2, 3, 5 ],
        [ -6, -4, -1, 2, 3, 6 ], [ -4, -1, 2, 4 ],
        [ -5, -3, -1, 1, 2, 3, 5 ] ],
    [ [ -6, -4, -2, 1, 3, 5, 6 ], [ -5, -2, 1, 2, 3, 5 ],
        [ -6, -5, -2, 1, 5 ], [ -6, -5, -3, -2, 1, 3, 4 ],
        [ -6, -5, -4, -2, 3, 5 ], [ -6, -4, -2, -1, 1, 2, 6 ] ])
\end{Verbatim}
 }

 

\subsection{\textcolor{Chapter }{EmptyPBR}}
\logpage{[ 4, 2, 3 ]}\nobreak
\hyperdef{L}{X8646781B7EAE04C0}{}
{\noindent\textcolor{FuncColor}{$\triangleright$\enspace\texttt{EmptyPBR({\mdseries\slshape n})\index{EmptyPBR@\texttt{EmptyPBR}}
\label{EmptyPBR}
}\hfill{\scriptsize (operation)}}\\
\textbf{\indent Returns:\ }
A PBR.



 If \mbox{\texttt{\mdseries\slshape n}} is a positive integer, then \texttt{EmptyPBR} returns the PBR of degree \mbox{\texttt{\mdseries\slshape n}} with no edges. 
\begin{Verbatim}[commandchars=!@|,fontsize=\small,frame=single,label=Example]
  !gapprompt@gap>| !gapinput@x := EmptyPBR(3);|
  PBR([ [  ], [  ], [  ] ], [ [  ], [  ], [  ] ])
  !gapprompt@gap>| !gapinput@IsEmptyPBR(x);|
  true
\end{Verbatim}
 }

 

\subsection{\textcolor{Chapter }{IdentityPBR}}
\logpage{[ 4, 2, 4 ]}\nobreak
\hyperdef{L}{X80D20EA3816DC862}{}
{\noindent\textcolor{FuncColor}{$\triangleright$\enspace\texttt{IdentityPBR({\mdseries\slshape n})\index{IdentityPBR@\texttt{IdentityPBR}}
\label{IdentityPBR}
}\hfill{\scriptsize (operation)}}\\
\textbf{\indent Returns:\ }
A PBR.



 If \mbox{\texttt{\mdseries\slshape n}} is a positive integer, then \texttt{IdentityPBR} returns the identity PBR of degree \mbox{\texttt{\mdseries\slshape n}}. This PBR has \texttt{2}\mbox{\texttt{\mdseries\slshape n}} edges: specifically, for each \texttt{i} in the ranges \texttt{[1 .. n]} and \texttt{[\texttt{\symbol{45}}n .. \texttt{\symbol{45}}1]}, the identity PBR has an edge from \texttt{i} to \texttt{\texttt{\symbol{45}}i}. 
\begin{Verbatim}[commandchars=!@|,fontsize=\small,frame=single,label=Example]
  !gapprompt@gap>| !gapinput@x := IdentityPBR(3);|
  PBR([ [ -1 ], [ -2 ], [ -3 ] ], [ [ 1 ], [ 2 ], [ 3 ] ])
  !gapprompt@gap>| !gapinput@IsIdentityPBR(x);|
  true
\end{Verbatim}
 }

 

\subsection{\textcolor{Chapter }{UniversalPBR}}
\logpage{[ 4, 2, 5 ]}\nobreak
\hyperdef{L}{X847BA0177D90E9D7}{}
{\noindent\textcolor{FuncColor}{$\triangleright$\enspace\texttt{UniversalPBR({\mdseries\slshape n})\index{UniversalPBR@\texttt{UniversalPBR}}
\label{UniversalPBR}
}\hfill{\scriptsize (operation)}}\\
\textbf{\indent Returns:\ }
A PBR.



 If \mbox{\texttt{\mdseries\slshape n}} is a positive integer, then \texttt{UniversalPBR} returns the PBR of degree \mbox{\texttt{\mdseries\slshape n}} with \texttt{4 * n \texttt{\symbol{94}} 2} edges, i.e. every possible edge. 
\begin{Verbatim}[commandchars=!@|,fontsize=\small,frame=single,label=Example]
  !gapprompt@gap>| !gapinput@x := UniversalPBR(2);|
  PBR([ [ -2, -1, 1, 2 ], [ -2, -1, 1, 2 ] ],
    [ [ -2, -1, 1, 2 ], [ -2, -1, 1, 2 ] ])
  !gapprompt@gap>| !gapinput@IsUniversalPBR(x);|
  true
\end{Verbatim}
 }

 }

   
\section{\textcolor{Chapter }{Changing the representation of a PBR}}\label{Changing the representation of a PBR}
\logpage{[ 4, 3, 0 ]}
\hyperdef{L}{X86B714987C01895F}{}
{
  It is possible that a PBR can be represented as another type of object, or
that another type of \textsf{GAP} object can be represented as a PBR. In this section, we describe the functions
in the \textsf{Semigroups} package for changing the representation of PBR, or for changing the
representation of another type of object to that of a PBR. 

 The operations \texttt{AsPermutation} (\ref{AsPermutation:for a PBR}), \texttt{AsPartialPerm} (\ref{AsPartialPerm:for a PBR}), \texttt{AsTransformation} (\ref{AsTransformation:for a PBR}), \texttt{AsBipartition} (\ref{AsBipartition}), \texttt{AsBooleanMat} (\ref{AsBooleanMat}) can be used to convert PBRs into permutations, partial permutations,
transformations, bipartitions, and boolean matrices where appropriate. 

\subsection{\textcolor{Chapter }{AsPBR}}
\logpage{[ 4, 3, 1 ]}\nobreak
\hyperdef{L}{X81CBBE6080439596}{}
{\noindent\textcolor{FuncColor}{$\triangleright$\enspace\texttt{AsPBR({\mdseries\slshape x[, n]})\index{AsPBR@\texttt{AsPBR}}
\label{AsPBR}
}\hfill{\scriptsize (operation)}}\\
\textbf{\indent Returns:\ }
A PBR.



 \texttt{AsPBR} returns the boolean matrix, bipartition, transformation, partial permutation,
or permutation \mbox{\texttt{\mdseries\slshape x}} as a PBR of degree \mbox{\texttt{\mdseries\slshape n}}. 

 There are several possible arguments for \texttt{AsPBR}: 
\begin{description}
\item[{bipartitions}]  If \mbox{\texttt{\mdseries\slshape x}} is a bipartition and \mbox{\texttt{\mdseries\slshape n}} is a positive integer, then \texttt{AsPBR} returns a PBR corresponding to \mbox{\texttt{\mdseries\slshape x}} with degree \mbox{\texttt{\mdseries\slshape n}}. The resulting PBR has an edge from \texttt{i} to \texttt{j} whenever \texttt{i} and \texttt{j} belong to the same block of \mbox{\texttt{\mdseries\slshape x}}.

 If the optional second argument \mbox{\texttt{\mdseries\slshape n}} is not specified, then degree of the bipartition \mbox{\texttt{\mdseries\slshape x}} is used by default. 
\item[{boolean matrices}]  If \mbox{\texttt{\mdseries\slshape x}} is a boolean matrix of even dimension \texttt{2 * m} and \mbox{\texttt{\mdseries\slshape n}} is a positive integer, then \texttt{AsPBR} returns a PBR corresponding to \mbox{\texttt{\mdseries\slshape x}} with degree \mbox{\texttt{\mdseries\slshape n}}. If the optional second argument \mbox{\texttt{\mdseries\slshape n}} is not specified, then dimension of the boolean matrix \mbox{\texttt{\mdseries\slshape x}} is used by default. 
\item[{transformations, partial perms, permutations}]  If \mbox{\texttt{\mdseries\slshape x}} is a transformation, partial perm, or permutation and \mbox{\texttt{\mdseries\slshape n}} is a positive integer, then \texttt{AsPBR} is a synonym for \texttt{AsPBR(AsBipartition(\mbox{\texttt{\mdseries\slshape x}}, \mbox{\texttt{\mdseries\slshape n}}))}. If the optional second argument \mbox{\texttt{\mdseries\slshape n}} is not specified, then \texttt{AsPBR} is a synonym for \texttt{AsPBR(AsBipartition(\mbox{\texttt{\mdseries\slshape x}}))}. See \texttt{AsBipartition} (\ref{AsBipartition}) for more details. 
\end{description}
 
\begin{Verbatim}[commandchars=!@|,fontsize=\small,frame=single,label=Example]
  !gapprompt@gap>| !gapinput@x := Bipartition([[1, 2, -1], [3, -2], [4, -3, -4]]);|
  <block bijection: [ 1, 2, -1 ], [ 3, -2 ], [ 4, -3, -4 ]>
  !gapprompt@gap>| !gapinput@AsPBR(x, 2);|
  PBR([ [ -1, 1, 2 ], [ -1, 1, 2 ] ], [ [ -1, 1, 2 ], [ -2 ] ])
  !gapprompt@gap>| !gapinput@AsPBR(x, 5);|
  PBR([ [ -1, 1, 2 ], [ -1, 1, 2 ], [ -2, 3 ], [ -4, -3, 4 ], [  ] ],
    [ [ -1, 1, 2 ], [ -2, 3 ], [ -4, -3, 4 ], [ -4, -3, 4 ], [  ] ])
  !gapprompt@gap>| !gapinput@AsPBR(x);|
  PBR([ [ -1, 1, 2 ], [ -1, 1, 2 ], [ -2, 3 ], [ -4, -3, 4 ] ],
    [ [ -1, 1, 2 ], [ -2, 3 ], [ -4, -3, 4 ], [ -4, -3, 4 ] ])
  !gapprompt@gap>| !gapinput@mat := Matrix(IsBooleanMat, [[1, 0, 0, 1],|
  !gapprompt@>| !gapinput@                                [0, 1, 1, 0],|
  !gapprompt@>| !gapinput@                                [1, 0, 1, 1],|
  !gapprompt@>| !gapinput@                                [0, 0, 0, 1]]);;|
  !gapprompt@gap>| !gapinput@AsPBR(mat);|
  PBR([ [ -2, 1 ], [ -1, 2 ] ], [ [ -2, -1, 1 ], [ -2 ] ])
  !gapprompt@gap>| !gapinput@AsPBR(mat, 2);|
  PBR([ [ 1 ] ], [ [ -1 ] ])
  !gapprompt@gap>| !gapinput@AsPBR(mat, 6);|
  PBR([ [ -2, 1 ], [ -1, 2 ], [  ] ], [ [ -2, -1, 1 ], [ -2 ], [  ] ])
  !gapprompt@gap>| !gapinput@x := Transformation([2, 2, 1]);;|
  !gapprompt@gap>| !gapinput@AsPBR(x);|
  PBR([ [ -2 ], [ -2 ], [ -1 ] ], [ [ 3 ], [ 1, 2 ], [  ] ])
  !gapprompt@gap>| !gapinput@AsPBR(x, 2);|
  PBR([ [ -2 ], [ -2 ] ], [ [  ], [ 1, 2 ] ])
  !gapprompt@gap>| !gapinput@AsPBR(x, 4);|
  PBR([ [ -2 ], [ -2 ], [ -1 ], [ -4 ] ],
    [ [ 3 ], [ 1, 2 ], [  ], [ 4 ] ])
  !gapprompt@gap>| !gapinput@x := PartialPerm([4, 3]);|
  [1,4][2,3]
  !gapprompt@gap>| !gapinput@AsPBR(x);|
  PBR([ [ -4 ], [ -3 ], [  ], [  ] ], [ [  ], [  ], [ 2 ], [ 1 ] ])
  !gapprompt@gap>| !gapinput@AsPBR(x, 2);|
  PBR([ [  ], [  ] ], [ [  ], [  ] ])
  !gapprompt@gap>| !gapinput@AsPBR(x, 5);|
  PBR([ [ -4 ], [ -3 ], [  ], [  ], [  ] ],
    [ [  ], [  ], [ 2 ], [ 1 ], [  ] ])
  !gapprompt@gap>| !gapinput@x := (1, 3)(2, 4);|
  (1,3)(2,4)
  !gapprompt@gap>| !gapinput@AsPBR(x);|
  PBR([ [ -3, 1 ], [ -4, 2 ], [ -1, 3 ], [ -2, 4 ] ],
    [ [ -1, 3 ], [ -2, 4 ], [ -3, 1 ], [ -4, 2 ] ])
  !gapprompt@gap>| !gapinput@AsPBR(x, 5);|
  PBR([ [ -3, 1 ], [ -4, 2 ], [ -1, 3 ], [ -2, 4 ], [ -5, 5 ] ],
    [ [ -1, 3 ], [ -2, 4 ], [ -3, 1 ], [ -4, 2 ], [ -5, 5 ] ])
\end{Verbatim}
 }

 

\subsection{\textcolor{Chapter }{AsTransformation (for a PBR)}}
\logpage{[ 4, 3, 2 ]}\nobreak
\hyperdef{L}{X8407F516825A514A}{}
{\noindent\textcolor{FuncColor}{$\triangleright$\enspace\texttt{AsTransformation({\mdseries\slshape x})\index{AsTransformation@\texttt{AsTransformation}!for a PBR}
\label{AsTransformation:for a PBR}
}\hfill{\scriptsize (attribute)}}\\
\textbf{\indent Returns:\ }
A transformation.



 When the argument \mbox{\texttt{\mdseries\slshape x}} is a PBR which satisfies \texttt{IsTransformationPBR} (\ref{IsTransformationPBR}), then this attribute returns that transformation. 
\begin{Verbatim}[commandchars=!@|,fontsize=\small,frame=single,label=Example]
  !gapprompt@gap>| !gapinput@x := PBR([[-3], [-3], [-2]], [[], [3], [1, 2]]);;|
  !gapprompt@gap>| !gapinput@IsTransformationPBR(x);|
  true
  !gapprompt@gap>| !gapinput@AsTransformation(x);|
  Transformation( [ 3, 3, 2 ] )
  !gapprompt@gap>| !gapinput@x := PBR([[1], [1, 2]], [[-2, -1], [-2, -1]]);;|
  !gapprompt@gap>| !gapinput@AsTransformation(x);|
  Error, the argument (a pbr) does not define a transformation
\end{Verbatim}
 }

 

\subsection{\textcolor{Chapter }{AsPartialPerm (for a PBR)}}
\logpage{[ 4, 3, 3 ]}\nobreak
\hyperdef{L}{X795B1C16819905E8}{}
{\noindent\textcolor{FuncColor}{$\triangleright$\enspace\texttt{AsPartialPerm({\mdseries\slshape x})\index{AsPartialPerm@\texttt{AsPartialPerm}!for a PBR}
\label{AsPartialPerm:for a PBR}
}\hfill{\scriptsize (operation)}}\\
\textbf{\indent Returns:\ }
A partial perm.



 When the argument \mbox{\texttt{\mdseries\slshape x}} is a PBR which satisfies \texttt{IsPartialPermPBR} (\ref{IsPartialPermPBR}), then this function returns that partial perm. 
\begin{Verbatim}[commandchars=!@|,fontsize=\small,frame=single,label=Example]
  !gapprompt@gap>| !gapinput@x := PBR([[-1, 1], [-3, 2], [-4, 3], [4], [5]],|
  !gapprompt@>| !gapinput@            [[-1, 1], [-2], [-3, 2], [-4, 3], [-5]]);;|
  !gapprompt@gap>| !gapinput@IsPartialPermPBR(x);|
  true
  !gapprompt@gap>| !gapinput@AsPartialPerm(x);|
  [2,3,4](1)
\end{Verbatim}
 }

 

\subsection{\textcolor{Chapter }{AsPermutation (for a PBR)}}
\logpage{[ 4, 3, 4 ]}\nobreak
\hyperdef{L}{X86786B297FBCD064}{}
{\noindent\textcolor{FuncColor}{$\triangleright$\enspace\texttt{AsPermutation({\mdseries\slshape x})\index{AsPermutation@\texttt{AsPermutation}!for a PBR}
\label{AsPermutation:for a PBR}
}\hfill{\scriptsize (attribute)}}\\
\textbf{\indent Returns:\ }
A permutation.



 When the argument \mbox{\texttt{\mdseries\slshape x}} is a PBR which satisfies \texttt{IsPermPBR} (\ref{IsPermPBR}), then this attribute returns that permutation. 
\begin{Verbatim}[commandchars=!@|,fontsize=\small,frame=single,label=Example]
  !gapprompt@gap>| !gapinput@x := PBR([[-1, 1], [-4, 2], [-2, 3], [-3, 4]],|
  !gapprompt@>| !gapinput@            [[-1, 1], [-2, 3], [-3, 4], [-4, 2]]);;|
  !gapprompt@gap>| !gapinput@IsPermPBR(x);|
  true
  !gapprompt@gap>| !gapinput@AsPermutation(x);|
  (2,4,3)
\end{Verbatim}
 }

 }

   
\section{\textcolor{Chapter }{Operators for PBRs}}\label{Operators for PBRs}
\logpage{[ 4, 4, 0 ]}
\hyperdef{L}{X872B5817878660E5}{}
{
  
\begin{description}
\item[{\texttt{\mbox{\texttt{\mdseries\slshape x}} * \mbox{\texttt{\mdseries\slshape y}}}}]  \index{*@\texttt{ * } (for PBRs)} returns the product of \mbox{\texttt{\mdseries\slshape x}} and \mbox{\texttt{\mdseries\slshape y}} when \mbox{\texttt{\mdseries\slshape x}} and \mbox{\texttt{\mdseries\slshape y}} are PBRs. 
\item[{\texttt{\mbox{\texttt{\mdseries\slshape x}} {\textless} \mbox{\texttt{\mdseries\slshape y}}}}]  \index{<@\texttt{{\textless}} (for PBRs)} returns \texttt{true} if the degree of \mbox{\texttt{\mdseries\slshape x}} is less than the degree of \mbox{\texttt{\mdseries\slshape y}}, or the degrees are equal and the out\texttt{\symbol{45}}neighbours of \mbox{\texttt{\mdseries\slshape x}} (as a list of list of positive integers) is lexicographically less than the
out\texttt{\symbol{45}}neighbours of \mbox{\texttt{\mdseries\slshape y}}. 
\item[{\texttt{\mbox{\texttt{\mdseries\slshape x}} = \mbox{\texttt{\mdseries\slshape y}}}}]  \index{=@\texttt{=} (for PBRs)} returns \texttt{true} if the PBR \mbox{\texttt{\mdseries\slshape x}} equals the PBR \mbox{\texttt{\mdseries\slshape y}} and returns \texttt{false} if it does not. 
\end{description}
 }

   
\section{\textcolor{Chapter }{Attributes for PBRs}}\label{Attributes for PBRs}
\logpage{[ 4, 5, 0 ]}
\hyperdef{L}{X78EC8E597EB99730}{}
{
  In this section we describe various attributes that a PBR can possess. 

\subsection{\textcolor{Chapter }{StarOp (for a PBR)}}
\logpage{[ 4, 5, 1 ]}\nobreak
\hyperdef{L}{X7DFC277E80A50C2F}{}
{\noindent\textcolor{FuncColor}{$\triangleright$\enspace\texttt{StarOp({\mdseries\slshape x})\index{StarOp@\texttt{StarOp}!for a PBR}
\label{StarOp:for a PBR}
}\hfill{\scriptsize (operation)}}\\
\noindent\textcolor{FuncColor}{$\triangleright$\enspace\texttt{Star({\mdseries\slshape x})\index{Star@\texttt{Star}!for a PBR}
\label{Star:for a PBR}
}\hfill{\scriptsize (attribute)}}\\
\textbf{\indent Returns:\ }
A PBR.



 \texttt{StarOp} returns the unique PBR \texttt{y} obtained by exchanging the positive and negative numbers in \mbox{\texttt{\mdseries\slshape x}} (i.e. multiplying \texttt{ExtRepOfObj} (\ref{ExtRepOfObj:for a PBR}) by \texttt{\texttt{\symbol{45}}1} and swapping its first and second components). 
\begin{Verbatim}[commandchars=!@|,fontsize=\small,frame=single,label=Example]
  !gapprompt@gap>| !gapinput@x := PBR([[], [-1], []], [[-3, -2, 2, 3], [-2, 1], []]);;|
  !gapprompt@gap>| !gapinput@Star(x);|
  PBR([ [ -3, -2, 2, 3 ], [ -1, 2 ], [  ] ], [ [  ], [ 1 ], [  ] ])
\end{Verbatim}
 }

 

\subsection{\textcolor{Chapter }{DegreeOfPBR}}
\logpage{[ 4, 5, 2 ]}\nobreak
\hyperdef{L}{X785B576B7823D626}{}
{\noindent\textcolor{FuncColor}{$\triangleright$\enspace\texttt{DegreeOfPBR({\mdseries\slshape x})\index{DegreeOfPBR@\texttt{DegreeOfPBR}}
\label{DegreeOfPBR}
}\hfill{\scriptsize (attribute)}}\\
\noindent\textcolor{FuncColor}{$\triangleright$\enspace\texttt{DegreeOfPBRCollection({\mdseries\slshape x})\index{DegreeOfPBRCollection@\texttt{DegreeOfPBRCollection}}
\label{DegreeOfPBRCollection}
}\hfill{\scriptsize (attribute)}}\\
\textbf{\indent Returns:\ }
A positive integer.



 The degree of a PBR is, roughly speaking, the number of points where it is
defined. More precisely, if \mbox{\texttt{\mdseries\slshape x}} is a PBR defined on \texttt{2 * n} points, then the degree of \mbox{\texttt{\mdseries\slshape x}} is \texttt{n}. 

 The degree of a collection \mbox{\texttt{\mdseries\slshape coll}} of PBRs of equal degree is just the degree of any (and every) PBR in \mbox{\texttt{\mdseries\slshape coll}}. The degree of collection of PBRs of unequal degrees is not defined. 
\begin{Verbatim}[commandchars=!@|,fontsize=\small,frame=single,label=Example]
  !gapprompt@gap>| !gapinput@x := PBR([[-2], [-2, -1, 2, 3], [-1, 1, 2, 3]],|
  !gapprompt@>| !gapinput@            [[-1, 1], [2, 3], [-3, 2, 3]]);|
  PBR([ [ -2 ], [ -2, -1, 2, 3 ], [ -1, 1, 2, 3 ] ],
    [ [ -1, 1 ], [ 2, 3 ], [ -3, 2, 3 ] ])
  !gapprompt@gap>| !gapinput@DegreeOfPBR(x);|
  3
  !gapprompt@gap>| !gapinput@S := FullPBRMonoid(2);|
  <pbr monoid of degree 2 with 10 generators>
  !gapprompt@gap>| !gapinput@DegreeOfPBRCollection(S);|
  2
\end{Verbatim}
 }

 

\subsection{\textcolor{Chapter }{ExtRepOfObj (for a PBR)}}
\logpage{[ 4, 5, 3 ]}\nobreak
\hyperdef{L}{X78302D7E81BB1E54}{}
{\noindent\textcolor{FuncColor}{$\triangleright$\enspace\texttt{ExtRepOfObj({\mdseries\slshape x})\index{ExtRepOfObj@\texttt{ExtRepOfObj}!for a PBR}
\label{ExtRepOfObj:for a PBR}
}\hfill{\scriptsize (operation)}}\\
\textbf{\indent Returns:\ }
A pair of lists of lists of integers.



 If \texttt{n} is the degree of the PBR \mbox{\texttt{\mdseries\slshape x}}, then \texttt{ExtRepOfObj} returns the argument required by \texttt{PBR} (\ref{PBR}) to create a PBR equal to \mbox{\texttt{\mdseries\slshape x}}, i.e. \texttt{PBR(ExtRepOfObj(\mbox{\texttt{\mdseries\slshape x}}))} returns a PBR equal to \mbox{\texttt{\mdseries\slshape x}}. 
\begin{Verbatim}[commandchars=!@|,fontsize=\small,frame=single,label=Example]
  !gapprompt@gap>| !gapinput@x := PBR([[-1, 1], [-2, 2]],|
  !gapprompt@>| !gapinput@            [[-2, -1, 1], [-1, 1, 2]]);|
  PBR([ [ -1, 1 ], [ -2, 2 ] ], [ [ -2, -1, 1 ], [ -1, 1, 2 ] ])
  !gapprompt@gap>| !gapinput@ExtRepOfObj(x);|
  [ [ [ -1, 1 ], [ -2, 2 ] ], [ [ -2, -1, 1 ], [ -1, 1, 2 ] ] ]
\end{Verbatim}
 }

 

\subsection{\textcolor{Chapter }{PBRNumber}}
\logpage{[ 4, 5, 4 ]}\nobreak
\hyperdef{L}{X7F4C8A2B79E6D963}{}
{\noindent\textcolor{FuncColor}{$\triangleright$\enspace\texttt{PBRNumber({\mdseries\slshape m, n})\index{PBRNumber@\texttt{PBRNumber}}
\label{PBRNumber}
}\hfill{\scriptsize (operation)}}\\
\noindent\textcolor{FuncColor}{$\triangleright$\enspace\texttt{NumberPBR({\mdseries\slshape mat})\index{NumberPBR@\texttt{NumberPBR}}
\label{NumberPBR}
}\hfill{\scriptsize (operation)}}\\
\textbf{\indent Returns:\ }
A PBR, or a positive integer.



 These functions implement a bijection from the set of all PBRs of degree \mbox{\texttt{\mdseries\slshape n}} and the numbers \texttt{[1 .. 2 \texttt{\symbol{94}} (4 * \mbox{\texttt{\mdseries\slshape n}} \texttt{\symbol{94}} 2)]}. 

 More precisely, if \mbox{\texttt{\mdseries\slshape m}} and \mbox{\texttt{\mdseries\slshape n}} are positive integers such that \mbox{\texttt{\mdseries\slshape m}} is at most \texttt{2 \texttt{\symbol{94}} (4 * \mbox{\texttt{\mdseries\slshape n}} \texttt{\symbol{94}} 2)}, then \texttt{PBRNumber} returns the \mbox{\texttt{\mdseries\slshape m}}th PBR of degree \mbox{\texttt{\mdseries\slshape n}}. 

 If \mbox{\texttt{\mdseries\slshape mat}} is a PBR of degree \mbox{\texttt{\mdseries\slshape n}}, then \texttt{NumberPBR} returns the number in \texttt{[1 .. 2 \texttt{\symbol{94}} (4 * \mbox{\texttt{\mdseries\slshape n}} \texttt{\symbol{94}} 2)]} that corresponds to \mbox{\texttt{\mdseries\slshape mat}}. 
\begin{Verbatim}[commandchars=!@|,fontsize=\small,frame=single,label=Example]
  !gapprompt@gap>| !gapinput@S := FullPBRMonoid(1);|
  <pbr monoid of degree 1 with 4 generators>
  !gapprompt@gap>| !gapinput@List(S, NumberPBR);|
  [ 3, 15, 5, 7, 8, 1, 4, 11, 13, 16, 6, 2, 9, 12, 14, 10 ]
\end{Verbatim}
 }

 

\subsection{\textcolor{Chapter }{IsEmptyPBR}}
\logpage{[ 4, 5, 5 ]}\nobreak
\hyperdef{L}{X82FD0AB179ED4AFD}{}
{\noindent\textcolor{FuncColor}{$\triangleright$\enspace\texttt{IsEmptyPBR({\mdseries\slshape x})\index{IsEmptyPBR@\texttt{IsEmptyPBR}}
\label{IsEmptyPBR}
}\hfill{\scriptsize (property)}}\\
\textbf{\indent Returns:\ }
\texttt{true} or \texttt{false}.



 A PBR is \textsc{empty} if it has no edges. \texttt{IsEmptyPBR} returns \texttt{true} if the PBR \mbox{\texttt{\mdseries\slshape x}} is empty and \texttt{false} if it is not. 
\begin{Verbatim}[commandchars=!@|,fontsize=\small,frame=single,label=Example]
  !gapprompt@gap>| !gapinput@x := PBR([[]], [[]]);;|
  !gapprompt@gap>| !gapinput@IsEmptyPBR(x);|
  true
  !gapprompt@gap>| !gapinput@x := PBR([[-2, 1], [2]], [[-1], [-2, 1]]);|
  PBR([ [ -2, 1 ], [ 2 ] ], [ [ -1 ], [ -2, 1 ] ])
  !gapprompt@gap>| !gapinput@IsEmptyPBR(x);|
  false
\end{Verbatim}
 }

 

\subsection{\textcolor{Chapter }{IsIdentityPBR}}
\logpage{[ 4, 5, 6 ]}\nobreak
\hyperdef{L}{X7E263B2F7B838D6E}{}
{\noindent\textcolor{FuncColor}{$\triangleright$\enspace\texttt{IsIdentityPBR({\mdseries\slshape x})\index{IsIdentityPBR@\texttt{IsIdentityPBR}}
\label{IsIdentityPBR}
}\hfill{\scriptsize (property)}}\\
\textbf{\indent Returns:\ }
\texttt{true} or \texttt{false}.



 A PBR of degree \texttt{n} is the \textsc{identity} PBR of degree \texttt{n} if it is the identity of the full PBR monoid of degree \texttt{n}. The identity PBR of degree \texttt{n} has \texttt{2n} edges. Specifically, for each \texttt{i} in the ranges \texttt{[1 .. n]} and \texttt{[\texttt{\symbol{45}}n .. \texttt{\symbol{45}}1]}, the identity PBR has an edge from \texttt{i} to \texttt{\texttt{\symbol{45}}i}. 

 \texttt{IsIdentityPBR} returns \texttt{true} is the PBR \mbox{\texttt{\mdseries\slshape x}} is an identity PBR and \texttt{false} if it is not. 
\begin{Verbatim}[commandchars=!@|,fontsize=\small,frame=single,label=Example]
  !gapprompt@gap>| !gapinput@x := PBR([[-2], [-1]], [[1], [2]]);|
  PBR([ [ -2 ], [ -1 ] ], [ [ 1 ], [ 2 ] ])
  !gapprompt@gap>| !gapinput@IsIdentityPBR(x);|
  false
  !gapprompt@gap>| !gapinput@x := PBR([[-1]], [[1]]);|
  PBR([ [ -1 ] ], [ [ 1 ] ])
  !gapprompt@gap>| !gapinput@IsIdentityPBR(x);|
  true
\end{Verbatim}
 }

 

\subsection{\textcolor{Chapter }{IsUniversalPBR}}
\logpage{[ 4, 5, 7 ]}\nobreak
\hyperdef{L}{X7A280FC27BAD0EF0}{}
{\noindent\textcolor{FuncColor}{$\triangleright$\enspace\texttt{IsUniversalPBR({\mdseries\slshape x})\index{IsUniversalPBR@\texttt{IsUniversalPBR}}
\label{IsUniversalPBR}
}\hfill{\scriptsize (property)}}\\
\textbf{\indent Returns:\ }
\texttt{true} or \texttt{false}.



 A PBR of degree \texttt{n} is \textsc{universal} if it has \texttt{4 * n \texttt{\symbol{94}} 2} edges, i.e. every possible edge. 
\begin{Verbatim}[commandchars=!@|,fontsize=\small,frame=single,label=Example]
  !gapprompt@gap>| !gapinput@x := PBR([[]], [[]]);|
  PBR([ [  ] ], [ [  ] ])
  !gapprompt@gap>| !gapinput@IsUniversalPBR(x);|
  false
  !gapprompt@gap>| !gapinput@x := PBR([[-2, 1], [2]], [[-1], [-2, 1]]);|
  PBR([ [ -2, 1 ], [ 2 ] ], [ [ -1 ], [ -2, 1 ] ])
  !gapprompt@gap>| !gapinput@IsUniversalPBR(x);|
  false
  !gapprompt@gap>| !gapinput@x := PBR([[-1, 1]], [[-1, 1]]);|
  PBR([ [ -1, 1 ] ], [ [ -1, 1 ] ])
  !gapprompt@gap>| !gapinput@IsUniversalPBR(x);|
  true
\end{Verbatim}
 }

 

\subsection{\textcolor{Chapter }{IsBipartitionPBR}}
\logpage{[ 4, 5, 8 ]}\nobreak
\hyperdef{L}{X81EC86397E098BC8}{}
{\noindent\textcolor{FuncColor}{$\triangleright$\enspace\texttt{IsBipartitionPBR({\mdseries\slshape x})\index{IsBipartitionPBR@\texttt{IsBipartitionPBR}}
\label{IsBipartitionPBR}
}\hfill{\scriptsize (property)}}\\
\noindent\textcolor{FuncColor}{$\triangleright$\enspace\texttt{IsBlockBijectionPBR({\mdseries\slshape x})\index{IsBlockBijectionPBR@\texttt{IsBlockBijectionPBR}}
\label{IsBlockBijectionPBR}
}\hfill{\scriptsize (property)}}\\
\textbf{\indent Returns:\ }
\texttt{true} or \texttt{false}.



 If the PBR \mbox{\texttt{\mdseries\slshape x}} defines a bipartition, then \texttt{IsBipartitionPBR} returns \texttt{true}, and if not, then it returns \texttt{false}. 

 A PBR \mbox{\texttt{\mdseries\slshape x}} defines a bipartition if and only if when considered as a boolean matrix it is
an equivalence.

 If \mbox{\texttt{\mdseries\slshape x}} satisfies \texttt{IsBipartitionPBR} and when considered as a bipartition it is a block bijection, then \texttt{IsBlockBijectionPBR} returns \texttt{true}. 
\begin{Verbatim}[commandchars=!@|,fontsize=\small,frame=single,label=Example]
  !gapprompt@gap>| !gapinput@x := PBR([[-1, 3], [-1, 3], [-2, 1, 2, 3]],|
  !gapprompt@>| !gapinput@            [[-2, -1, 2], [-2, -1, 1, 2, 3],|
  !gapprompt@>| !gapinput@             [-2, -1, 1, 2]]);|
  PBR([ [ -1, 3 ], [ -1, 3 ], [ -2, 1, 2, 3 ] ],
    [ [ -2, -1, 2 ], [ -2, -1, 1, 2, 3 ], [ -2, -1, 1, 2 ] ])
  !gapprompt@gap>| !gapinput@IsBipartitionPBR(x);|
  false
  !gapprompt@gap>| !gapinput@x := PBR([[-2, -1, 1], [2, 3], [2, 3]],|
  !gapprompt@>| !gapinput@            [[-2, -1, 1], [-2, -1, 1], [-3]]);|
  PBR([ [ -2, -1, 1 ], [ 2, 3 ], [ 2, 3 ] ],
    [ [ -2, -1, 1 ], [ -2, -1, 1 ], [ -3 ] ])
  !gapprompt@gap>| !gapinput@IsBipartitionPBR(x);|
  true
  !gapprompt@gap>| !gapinput@IsBlockBijectionPBR(x);|
  false
\end{Verbatim}
 }

 

\subsection{\textcolor{Chapter }{IsTransformationPBR}}
\logpage{[ 4, 5, 9 ]}\nobreak
\hyperdef{L}{X7AF425D17BBE9023}{}
{\noindent\textcolor{FuncColor}{$\triangleright$\enspace\texttt{IsTransformationPBR({\mdseries\slshape x})\index{IsTransformationPBR@\texttt{IsTransformationPBR}}
\label{IsTransformationPBR}
}\hfill{\scriptsize (property)}}\\
\textbf{\indent Returns:\ }
 \texttt{true} or \texttt{false}. 



 If the PBR \mbox{\texttt{\mdseries\slshape x}} defines a transformation, then \texttt{IsTransformationPBR} returns \texttt{true}, and if not, then \texttt{false} is returned. 

 A PBR \mbox{\texttt{\mdseries\slshape x}} defines a transformation if and only if it satisfies \texttt{IsBipartitionPBR} (\ref{IsBipartitionPBR}) and when it is considered as a bipartition it satisfies \texttt{IsTransBipartition} (\ref{IsTransBipartition}). 

 With this definition, \texttt{AsPBR} (\ref{AsPBR}) and \texttt{AsTransformation} (\ref{AsTransformation:for a PBR}) define mutually inverse isomorphisms from the full transformation monoid of
degree \texttt{n} to the submonoid of the full PBR monoid of degree \texttt{n} consisting of all the elements satisfying \texttt{IsTransformationPBR}. 
\begin{Verbatim}[commandchars=!@|,fontsize=\small,frame=single,label=Example]
  !gapprompt@gap>| !gapinput@x := PBR([[-3], [-1], [-3]], [[2], [], [1, 3]]);|
  PBR([ [ -3 ], [ -1 ], [ -3 ] ], [ [ 2 ], [  ], [ 1, 3 ] ])
  !gapprompt@gap>| !gapinput@IsTransformationPBR(x);|
  true
  !gapprompt@gap>| !gapinput@x := AsTransformation(x);|
  Transformation( [ 3, 1, 3 ] )
  !gapprompt@gap>| !gapinput@AsPBR(x) * AsPBR(x) = AsPBR(x ^ 2);|
  true
  !gapprompt@gap>| !gapinput@Number(FullPBRMonoid(1), IsTransformationPBR);|
  1
  !gapprompt@gap>| !gapinput@x := PBR([[-2, -1, 2], [-2, 1, 2]], [[-1, 1], [-2]]);|
  PBR([ [ -2, -1, 2 ], [ -2, 1, 2 ] ], [ [ -1, 1 ], [ -2 ] ])
  !gapprompt@gap>| !gapinput@IsTransformationPBR(x);|
  false
\end{Verbatim}
 }

 

\subsection{\textcolor{Chapter }{IsDualTransformationPBR}}
\logpage{[ 4, 5, 10 ]}\nobreak
\hyperdef{L}{X7962D03186B1AFDF}{}
{\noindent\textcolor{FuncColor}{$\triangleright$\enspace\texttt{IsDualTransformationPBR({\mdseries\slshape x})\index{IsDualTransformationPBR@\texttt{IsDualTransformationPBR}}
\label{IsDualTransformationPBR}
}\hfill{\scriptsize (property)}}\\
\textbf{\indent Returns:\ }
\texttt{true} or \texttt{false}.



 If the PBR \mbox{\texttt{\mdseries\slshape x}} defines a dual transformation, then \texttt{IsDualTransformationPBR} returns \texttt{true}, and if not, then \texttt{false} is returned. 

 A PBR \mbox{\texttt{\mdseries\slshape x}} defines a dual transformation if and only if \texttt{Star(\mbox{\texttt{\mdseries\slshape x}})} satisfies \texttt{IsTransformationPBR} (\ref{IsTransformationPBR}). 

 
\begin{Verbatim}[commandchars=!@|,fontsize=\small,frame=single,label=Example]
  !gapprompt@gap>| !gapinput@x := PBR([[-3, 1, 3], [-1, 2], [-3, 1, 3]],|
  !gapprompt@>| !gapinput@            [[-1, 2], [-2], [-3, 1, 3]]);|
  PBR([ [ -3, 1, 3 ], [ -1, 2 ], [ -3, 1, 3 ] ],
    [ [ -1, 2 ], [ -2 ], [ -3, 1, 3 ] ])
  !gapprompt@gap>| !gapinput@IsDualTransformationPBR(x);|
  false
  !gapprompt@gap>| !gapinput@IsDualTransformationPBR(Star(x));|
  true
  !gapprompt@gap>| !gapinput@Number(FullPBRMonoid(1), IsDualTransformationPBR);|
  1
\end{Verbatim}
 }

 

\subsection{\textcolor{Chapter }{IsPartialPermPBR}}
\logpage{[ 4, 5, 11 ]}\nobreak
\hyperdef{L}{X7883CD5D824CC236}{}
{\noindent\textcolor{FuncColor}{$\triangleright$\enspace\texttt{IsPartialPermPBR({\mdseries\slshape x})\index{IsPartialPermPBR@\texttt{IsPartialPermPBR}}
\label{IsPartialPermPBR}
}\hfill{\scriptsize (property)}}\\
\textbf{\indent Returns:\ }
\texttt{true} or \texttt{false}.



 If the PBR \mbox{\texttt{\mdseries\slshape x}} defines a partial permutation, then \texttt{IsPartialPermPBR} returns \texttt{true}, and if not, then \texttt{false} is returned. 

 A PBR \mbox{\texttt{\mdseries\slshape x}} defines a partial perm if and only if it satisfies \texttt{IsBipartitionPBR} (\ref{IsBipartitionPBR}) and and when it is considered as a bipartition it satisfies \texttt{IsPartialPermBipartition} (\ref{IsPartialPermBipartition}). 

 With this definition, \texttt{AsPBR} (\ref{AsPBR}) and \texttt{AsPartialPerm} (\ref{AsPartialPerm:for a PBR}) define mutually inverse isomorphisms from the symmetric inverse monoid of
degree \texttt{n} to the submonoid of the full PBR monoid of degree \texttt{n} consisting of all the elements satisfying \texttt{IsPartialPermPBR}. 
\begin{Verbatim}[commandchars=!@|,fontsize=\small,frame=single,label=Example]
  !gapprompt@gap>| !gapinput@x := PBR([[-1, 1], [2]], [[-1, 1], [-2]]);|
  PBR([ [ -1, 1 ], [ 2 ] ], [ [ -1, 1 ], [ -2 ] ])
  !gapprompt@gap>| !gapinput@IsPartialPermPBR(x);|
  true
  !gapprompt@gap>| !gapinput@x := PartialPerm([3, 1]);|
  [2,1,3]
  !gapprompt@gap>| !gapinput@AsPBR(x) * AsPBR(x) = AsPBR(x ^ 2);|
  true
  !gapprompt@gap>| !gapinput@Number(FullPBRMonoid(1), IsPartialPermPBR);|
  2
\end{Verbatim}
 }

 

\subsection{\textcolor{Chapter }{IsPermPBR}}
\logpage{[ 4, 5, 12 ]}\nobreak
\hyperdef{L}{X85B21BB0835FE166}{}
{\noindent\textcolor{FuncColor}{$\triangleright$\enspace\texttt{IsPermPBR({\mdseries\slshape x})\index{IsPermPBR@\texttt{IsPermPBR}}
\label{IsPermPBR}
}\hfill{\scriptsize (property)}}\\
\textbf{\indent Returns:\ }
\texttt{true} or \texttt{false}.



 If the PBR \mbox{\texttt{\mdseries\slshape x}} defines a permutation, then \texttt{IsPermPBR} returns \texttt{true}, and if not, then \texttt{false} is returned. 

 A PBR \mbox{\texttt{\mdseries\slshape x}} defines a permutation if and only if it satisfies \texttt{IsBipartitionPBR} (\ref{IsBipartitionPBR}) and and when it is considered as a bipartition it satisfies \texttt{IsPermBipartition} (\ref{IsPermBipartition}). 

 With this definition, \texttt{AsPBR} (\ref{AsPBR}) and \texttt{AsPermutation} (\ref{AsPermutation:for a PBR}) define mutually inverse isomorphisms from the symmetric group of degree \texttt{n} to the subgroup of the full PBR monoid of degree \texttt{n} consisting of all the elements satisfying \texttt{IsPermPBR} (i.e. the \texttt{GroupOfUnits} (\ref{GroupOfUnits}) of the full PBR monoid of degree \texttt{n}). 
\begin{Verbatim}[commandchars=!@|,fontsize=\small,frame=single,label=Example]
  !gapprompt@gap>| !gapinput@x := PBR([[-2, 1], [-4, 2], [-1, 3], [-3, 4]],|
  !gapprompt@>| !gapinput@[[-1, 3], [-2, 1], [-3, 4], [-4, 2]]);;|
  !gapprompt@gap>| !gapinput@IsPermPBR(x);|
  true
  !gapprompt@gap>| !gapinput@x := (1, 5)(2, 4, 3);|
  (1,5)(2,4,3)
  !gapprompt@gap>| !gapinput@y := (1, 4, 3)(2, 5);|
  (1,4,3)(2,5)
  !gapprompt@gap>| !gapinput@AsPBR(x) * AsPBR(y) = AsPBR(x * y);|
  true
  !gapprompt@gap>| !gapinput@Number(FullPBRMonoid(1), IsPermPBR);|
  1
\end{Verbatim}
 }

 }

   
\section{\textcolor{Chapter }{ Semigroups of PBRs }}\logpage{[ 4, 6, 0 ]}
\hyperdef{L}{X7ECD4BBD7A0E834E}{}
{
  Semigroups and monoids of PBRs can be created in the usual way in \textsf{GAP} using the functions \texttt{Semigroup} (\textbf{Reference: Semigroup}) and \texttt{Monoid} (\textbf{Reference: Monoid}); see Chapter \ref{Semigroups and monoids defined by generating sets} for more details. 

 It is possible to create inverse semigroups and monoids of PBRs using \texttt{InverseSemigroup} (\textbf{Reference: InverseSemigroup}) and \texttt{InverseMonoid} (\textbf{Reference: InverseMonoid}) when the argument is a collection of PBRs satisfying \texttt{IsBipartitionPBR} (\ref{IsBipartitionPBR}) and when considered as bipartitions, the collection satisfies \texttt{IsGeneratorsOfInverseSemigroup}. 

 Note that every PBR semigroup in \textsf{Semigroups} is finite. 

\subsection{\textcolor{Chapter }{IsPBRSemigroup}}
\logpage{[ 4, 6, 1 ]}\nobreak
\hyperdef{L}{X8554A3F878A4DC73}{}
{\noindent\textcolor{FuncColor}{$\triangleright$\enspace\texttt{IsPBRSemigroup({\mdseries\slshape S})\index{IsPBRSemigroup@\texttt{IsPBRSemigroup}}
\label{IsPBRSemigroup}
}\hfill{\scriptsize (filter)}}\\
\noindent\textcolor{FuncColor}{$\triangleright$\enspace\texttt{IsPBRMonoid({\mdseries\slshape S})\index{IsPBRMonoid@\texttt{IsPBRMonoid}}
\label{IsPBRMonoid}
}\hfill{\scriptsize (filter)}}\\
\textbf{\indent Returns:\ }
\texttt{true} or \texttt{false}.



 A \emph{PBR semigroup} is simply a semigroup consisting of PBRs. An object \mbox{\texttt{\mdseries\slshape obj}} is a PBR semigroup in \textsf{GAP} if it satisfies \texttt{IsSemigroup} (\textbf{Reference: IsSemigroup}) and \texttt{IsPBRCollection} (\ref{IsPBRCollection}).

 A \emph{PBR monoid} is a monoid consisting of PBRs. An object \mbox{\texttt{\mdseries\slshape obj}} is a PBR monoid in \textsf{GAP} if it satisfies \texttt{IsMonoid} (\textbf{Reference: IsMonoid}) and \texttt{IsPBRCollection} (\ref{IsPBRCollection}).

 Note that it is possible for a PBR semigroup to have a multiplicative neutral
element (i.e. an identity element) but not to satisfy \texttt{IsPBRMonoid}. For example, 
\begin{Verbatim}[commandchars=!@|,fontsize=\small,frame=single,label=Example]
  !gapprompt@gap>| !gapinput@x := PBR([[-2, -1, 3], [-2, 2], [-3, -2, 1, 2, 3]],|
  !gapprompt@>| !gapinput@            [[-3, -2, -1, 2, 3], [-3, -2, -1, 2, 3], [-1]]);;|
  !gapprompt@gap>| !gapinput@S := Semigroup(x, One(x));|
  <commutative pbr monoid of degree 3 with 1 generator>
  !gapprompt@gap>| !gapinput@IsMonoid(S);|
  true
  !gapprompt@gap>| !gapinput@IsPBRMonoid(S);|
  true
  !gapprompt@gap>| !gapinput@S := Semigroup([|
  !gapprompt@>| !gapinput@ PBR([[-2, 1], [-3, 2], [-1, 3], [-4, 4, 5], [-4, 4, 5]],|
  !gapprompt@>| !gapinput@     [[-1, 3], [-2, 1], [-3, 2], [-4, 4, 5], [-5]]),|
  !gapprompt@>| !gapinput@ PBR([[-2, 1], [-1, 2], [-3, 3], [-4, 4, 5], [-4, 4, 5]],|
  !gapprompt@>| !gapinput@     [[-1, 2], [-2, 1], [-3, 3], [-4, 4, 5], [-5]]),|
  !gapprompt@>| !gapinput@ PBR([[-1, 1, 3], [-2, 2], [-1, 1, 3], [-4, 4, 5], [-4, 4, 5]],|
  !gapprompt@>| !gapinput@     [[-1, 1, 3], [-2, 2], [-3], [-4, 4, 5], [-5]])]);|
  <pbr semigroup of degree 5 with 3 generators>
  !gapprompt@gap>| !gapinput@One(S);|
  fail
  !gapprompt@gap>| !gapinput@MultiplicativeNeutralElement(S);|
  PBR([ [ -1, 1 ], [ -2, 2 ], [ -3, 3 ], [ -4, 4, 5 ], [ -4, 4, 5 ] ],
    [ [ -1, 1 ], [ -2, 2 ], [ -3, 3 ], [ -4, 4, 5 ], [ -5 ] ])
  !gapprompt@gap>| !gapinput@IsPBRMonoid(S);|
  false
\end{Verbatim}
 In this example \texttt{S} cannot be converted into a monoid using \texttt{AsMonoid} (\textbf{Reference: AsMonoid}) since the \texttt{One} (\textbf{Reference: One}) of any element in \texttt{S} differs from the multiplicative neutral element. 

 For more details see \texttt{IsMagmaWithOne} (\textbf{Reference: IsMagmaWithOne}). }

 

\subsection{\textcolor{Chapter }{DegreeOfPBRSemigroup}}
\logpage{[ 4, 6, 2 ]}\nobreak
\hyperdef{L}{X80FC004C7B65B4C0}{}
{\noindent\textcolor{FuncColor}{$\triangleright$\enspace\texttt{DegreeOfPBRSemigroup({\mdseries\slshape S})\index{DegreeOfPBRSemigroup@\texttt{DegreeOfPBRSemigroup}}
\label{DegreeOfPBRSemigroup}
}\hfill{\scriptsize (attribute)}}\\
\textbf{\indent Returns:\ }
A non\texttt{\symbol{45}}negative integer.



 The \emph{degree} of a PBR semigroup \mbox{\texttt{\mdseries\slshape S}} is just the degree of any (and every) element of \mbox{\texttt{\mdseries\slshape S}}. 
\begin{Verbatim}[commandchars=!@|,fontsize=\small,frame=single,label=Example]
  !gapprompt@gap>| !gapinput@S := Semigroup(|
  !gapprompt@>| !gapinput@ PBR([[-1, 1], [-2, 2], [-3, 3]],|
  !gapprompt@>| !gapinput@     [[-1, 1], [-2, 2], [-3, 3]]),|
  !gapprompt@>| !gapinput@ PBR([[1, 2], [1, 2], [-3, 3]],|
  !gapprompt@>| !gapinput@     [[-2, -1], [-2, -1], [-3, 3]]),|
  !gapprompt@>| !gapinput@ PBR([[-1, 1], [2, 3], [2, 3]],|
  !gapprompt@>| !gapinput@     [[-1, 1], [-3, -2], [-3, -2]]));|
  <pbr semigroup of degree 3 with 3 generators>
  !gapprompt@gap>| !gapinput@DegreeOfPBRSemigroup(S);|
  3
\end{Verbatim}
 }

 }

 }

 
\chapter{\textcolor{Chapter }{ Matrices over semirings }}\label{Matrices over semirings}
\logpage{[ 5, 0, 0 ]}
\hyperdef{L}{X82D6B7FE7CAC0AFA}{}
{
    In this chapter we describe the functionality in \textsf{Semigroups} for creating matrices over semirings. \textsc{Only square matrices are currently supported.} We use the term \textsc{matrix} to mean \textsc{square matrix} everywhere in this manual. 

 For reference, matrices over the following semirings are currently supported: 
\begin{description}
\item[{the Boolean semiring}]  the set $\{0, 1\}$ where $0 + 0 = 0$, $0 + 1 = 1 + 1 = 1 + 0 = 1$, $1\cdot 0 = 0 \cdot 0 = 0 \cdot 1 = 0$, and $1\cdot 1 = 1$. 
\item[{the max\texttt{\symbol{45}}plus semiring}]  the set of integers and negative infinity $\mathbb{Z}\cup \{-\infty\}$ with operations max and plus. 
\item[{the min\texttt{\symbol{45}}plus semiring}]  the set of integers and infinity $\mathbb{Z}\cup \{\infty\}$ with operations min and plus; 
\item[{tropical max\texttt{\symbol{45}}plus semirings}]  the set $\{-\infty, 0, 1, \ldots, t\}$ for some threshold $t$ with operations max and plus; 
\item[{tropical min\texttt{\symbol{45}}plus semirings}]  the set $\{0, 1, \ldots, t, \infty\}$ for some threshold $t$ with operations min and plus; 
\item[{the semiring $\mathbb{N}_{t,p}$}]  the semiring $\mathbb{N}_{t,p} = \{0, 1, \ldots, t, t + 1, \ldots, t + p - 1\}$ for some threshold $t$ and period $p$ under addition and multiplication modulo the congruence $t = t + p$; 
\item[{the integers}]  the usual ring of integers; 
\item[{finite fields}]  the finite fields \texttt{GF(q\texttt{\symbol{94}}d)} for prime \texttt{q} and some positive integer \texttt{d}. 
\end{description}
 With the exception of matrices of finite fields, semigroups of matrices in \textsf{Semigroups} are of the second type described in Section \ref{Underlying algorithms}. In other words, a version of the Froidure\texttt{\symbol{45}}Pin Algorithm \cite{Froidure1997aa} is used to compute semigroups of these types, i.e it is possible that all of
the elements of such a semigroup are enumerated and stored in the memory of
your computer.   
\section{\textcolor{Chapter }{Creating matrices over semirings}}\label{Matrices over arbitrary semirings}
\logpage{[ 5, 1, 0 ]}
\hyperdef{L}{X7ECF673C7BE2384D}{}
{
  In this section we describe the two main operations for creating matrices over
semirings in \textsf{Semigroups}, and the categories, attributes, and operations which apply to every matrix
over one of the semirings given at the start of this chapter.

 There are several special methods for boolean matrices, which can be found in
Section \ref{Boolean matrices}. There are also several special methods for finite fields, which can be found
in section \ref{Matrices over finite fields}. 

\subsection{\textcolor{Chapter }{IsMatrixOverSemiring}}
\logpage{[ 5, 1, 1 ]}\nobreak
\hyperdef{L}{X8711618C7A8A1B60}{}
{\noindent\textcolor{FuncColor}{$\triangleright$\enspace\texttt{IsMatrixOverSemiring({\mdseries\slshape obj})\index{IsMatrixOverSemiring@\texttt{IsMatrixOverSemiring}}
\label{IsMatrixOverSemiring}
}\hfill{\scriptsize (Category)}}\\
\textbf{\indent Returns:\ }
\texttt{true} or \texttt{false}.



 Every matrix over a semiring in \textsf{Semigroups} is a member of the category \texttt{IsMatrixOverSemiring}, which is a subcategory of \texttt{IsMultiplicativeElementWithOne} (\textbf{Reference: IsMultiplicativeElementWithOne}), \texttt{IsAssociativeElement} (\textbf{Reference: IsAssociativeElement}), and \texttt{IsPositionalObjectRep}; see  (\textbf{Reference: Representation}). 

 Every matrix over a semiring in \textsf{Semigroups} is a square matrix. 

 Basic operations for matrices over semirings are: \texttt{DimensionOfMatrixOverSemiring} (\ref{DimensionOfMatrixOverSemiring}), \texttt{TransposedMat} (\textbf{Reference: TransposedMat}), and \texttt{One} (\textbf{Reference: One}). }

 

\subsection{\textcolor{Chapter }{IsMatrixOverSemiringCollection}}
\logpage{[ 5, 1, 2 ]}\nobreak
\hyperdef{L}{X86F696B883677D6B}{}
{\noindent\textcolor{FuncColor}{$\triangleright$\enspace\texttt{IsMatrixOverSemiringCollection({\mdseries\slshape obj})\index{IsMatrixOverSemiringCollection@\texttt{IsMatrixOverSemiringCollection}}
\label{IsMatrixOverSemiringCollection}
}\hfill{\scriptsize (Category)}}\\
\noindent\textcolor{FuncColor}{$\triangleright$\enspace\texttt{IsMatrixOverSemiringCollColl({\mdseries\slshape obj})\index{IsMatrixOverSemiringCollColl@\texttt{IsMatrixOverSemiringCollColl}}
\label{IsMatrixOverSemiringCollColl}
}\hfill{\scriptsize (Category)}}\\
\textbf{\indent Returns:\ }
\texttt{true} or \texttt{false}.



 Every collection of matrices over the same semiring belongs to the category \texttt{IsMatrixOverSemiringCollection}. For example, semigroups of matrices over a semiring belong to \texttt{IsMatrixOverSemiringCollection}. 

 Every collection of collections of matrices over the same semiring belongs to
the category \texttt{IsMatrixOverSemiringCollColl}. For example, a list of semigroups of matrices over semirings belongs to \texttt{IsMatrixOverSemiringCollColl}. }

 

\subsection{\textcolor{Chapter }{DimensionOfMatrixOverSemiring}}
\logpage{[ 5, 1, 3 ]}\nobreak
\hyperdef{L}{X7C1CDA817CE076FD}{}
{\noindent\textcolor{FuncColor}{$\triangleright$\enspace\texttt{DimensionOfMatrixOverSemiring({\mdseries\slshape mat})\index{DimensionOfMatrixOverSemiring@\texttt{DimensionOfMatrixOverSemiring}}
\label{DimensionOfMatrixOverSemiring}
}\hfill{\scriptsize (attribute)}}\\
\textbf{\indent Returns:\ }
A positive integer.



 If \mbox{\texttt{\mdseries\slshape mat}} is a matrix over a semiring (i.e. belongs to the category \texttt{IsMatrixOverSemiring} (\ref{IsMatrixOverSemiring})), then \mbox{\texttt{\mdseries\slshape mat}} is a square \texttt{n} by \texttt{n} matrix. \texttt{DimensionOfMatrixOverSemiring} returns the dimension \texttt{n} of \mbox{\texttt{\mdseries\slshape mat}}. 
\begin{Verbatim}[commandchars=!@|,fontsize=\small,frame=single,label=Example]
  !gapprompt@gap>| !gapinput@x := BooleanMat([[1, 0, 0, 1],|
  !gapprompt@>| !gapinput@                    [0, 1, 1, 0],|
  !gapprompt@>| !gapinput@                    [1, 0, 1, 1],|
  !gapprompt@>| !gapinput@                    [0, 0, 0, 1]]);|
  Matrix(IsBooleanMat, [[1, 0, 0, 1], [0, 1, 1, 0], [1, 0, 1, 1],
    [0, 0, 0, 1]])
  !gapprompt@gap>| !gapinput@DimensionOfMatrixOverSemiring(x);|
  4
\end{Verbatim}
 }

 

\subsection{\textcolor{Chapter }{DimensionOfMatrixOverSemiringCollection}}
\logpage{[ 5, 1, 4 ]}\nobreak
\hyperdef{L}{X7FF0B2A783BA2D06}{}
{\noindent\textcolor{FuncColor}{$\triangleright$\enspace\texttt{DimensionOfMatrixOverSemiringCollection({\mdseries\slshape coll})\index{DimensionOfMatrixOverSemiringCollection@\texttt{Dimension}\-\texttt{Of}\-\texttt{Matrix}\-\texttt{Over}\-\texttt{Semiring}\-\texttt{Collection}}
\label{DimensionOfMatrixOverSemiringCollection}
}\hfill{\scriptsize (attribute)}}\\
\textbf{\indent Returns:\ }
A positive integer.



 If \mbox{\texttt{\mdseries\slshape coll}} is a collection of matrices over a semiring (i.e. belongs to the category \texttt{IsMatrixOverSemiringCollection} (\ref{IsMatrixOverSemiringCollection})), then the elements of \mbox{\texttt{\mdseries\slshape coll}} are square \texttt{n} by \texttt{n} matrices. \texttt{DimensionOfMatrixOverSemiringCollection} returns the dimension \texttt{n} of these matrices. 
\begin{Verbatim}[commandchars=!@|,fontsize=\small,frame=single,label=Example]
  !gapprompt@gap>| !gapinput@x := BooleanMat([[1, 0, 0, 1],|
  !gapprompt@>| !gapinput@                    [0, 1, 1, 0],|
  !gapprompt@>| !gapinput@                    [1, 0, 1, 1],|
  !gapprompt@>| !gapinput@                    [0, 0, 0, 1]]);|
  Matrix(IsBooleanMat, [[1, 0, 0, 1], [0, 1, 1, 0], [1, 0, 1, 1],
    [0, 0, 0, 1]])
  !gapprompt@gap>| !gapinput@DimensionOfMatrixOverSemiringCollection(Semigroup(x));|
  4
\end{Verbatim}
 }

 

\subsection{\textcolor{Chapter }{Matrix (for a filter and a matrix)}}
\logpage{[ 5, 1, 5 ]}\nobreak
\hyperdef{L}{X7DCA234C86ED8BD3}{}
{\noindent\textcolor{FuncColor}{$\triangleright$\enspace\texttt{Matrix({\mdseries\slshape filt, mat[, threshold[, period]]})\index{Matrix@\texttt{Matrix}!for a filter and a matrix}
\label{Matrix:for a filter and a matrix}
}\hfill{\scriptsize (operation)}}\\
\noindent\textcolor{FuncColor}{$\triangleright$\enspace\texttt{Matrix({\mdseries\slshape semiring, mat})\index{Matrix@\texttt{Matrix}!for a semiring and a matrix}
\label{Matrix:for a semiring and a matrix}
}\hfill{\scriptsize (operation)}}\\
\textbf{\indent Returns:\ }
A matrix over semiring.



 This operation can be used to construct a matrix over a semiring in \textsf{Semigroups}. 

 In its first form, the first argument \mbox{\texttt{\mdseries\slshape filt}} specifies the filter to be used to create the matrix, the second argument \mbox{\texttt{\mdseries\slshape mat}} is a \textsf{GAP} matrix (i.e. a list of lists) compatible with \mbox{\texttt{\mdseries\slshape filt}}, the third and fourth arguments \mbox{\texttt{\mdseries\slshape threshold}} and \mbox{\texttt{\mdseries\slshape period}} (if required) must be positive integers. 

 
\begin{description}
\item[{\mbox{\texttt{\mdseries\slshape filt}}}]  This must be one of the filters given in Section \ref{IsXMatrix}. 
\item[{\mbox{\texttt{\mdseries\slshape mat}}}]  This must be a list of \texttt{n} lists each of length \texttt{n} (i.e. a square matrix), consisting of elements belonging to the underlying
semiring described by \mbox{\texttt{\mdseries\slshape filt}}, and \mbox{\texttt{\mdseries\slshape threshold}} and \mbox{\texttt{\mdseries\slshape period}} if present. An error is given if \mbox{\texttt{\mdseries\slshape mat}} is not compatible with the other arguments. 

 For example, if \mbox{\texttt{\mdseries\slshape filt}} is \texttt{IsMaxPlusMatrix}, then the entries of \mbox{\texttt{\mdseries\slshape mat}} must belong to the max\texttt{\symbol{45}}plus semiring, i.e. they must be
integers or \texttt{\symbol{45}}$\infty$. 

 The supported semirings are fully described at the start of this chapter. 
\item[{\mbox{\texttt{\mdseries\slshape threshold}}}]  If \mbox{\texttt{\mdseries\slshape filt}} is any of \texttt{IsTropicalMaxPlusMatrix} (\ref{IsTropicalMaxPlusMatrix}), \texttt{IsTropicalMinPlusMatrix} (\ref{IsTropicalMinPlusMatrix}), or \texttt{IsNTPMatrix} (\ref{IsNTPMatrix}), then this argument specifies the threshold of the underlying semiring of the
matrix being created. 
\item[{\mbox{\texttt{\mdseries\slshape period}}}]  If \mbox{\texttt{\mdseries\slshape filt}} is \texttt{IsNTPMatrix} (\ref{IsNTPMatrix}), then this argument specifies the period of the underlying semiring of the
matrix being created. 
\end{description}
 In its second form, the arguments should be a semiring \mbox{\texttt{\mdseries\slshape semiring}} and matrix \mbox{\texttt{\mdseries\slshape mat}} with entries in \mbox{\texttt{\mdseries\slshape semiring}}. Currently, the only supported semirings are finite fields of prime order,
and the integers \texttt{Integers} (\textbf{Reference: Integers}). 

 The function \texttt{BooleanMat} (\ref{BooleanMat}) is provided for specifically creating boolean matrices. 
\begin{Verbatim}[commandchars=!@|,fontsize=\small,frame=single,label=Example]
  !gapprompt@gap>| !gapinput@Matrix(IsBooleanMat, [[1, 0, 0, 0],|
  !gapprompt@>| !gapinput@                         [0, 0, 0, 0],|
  !gapprompt@>| !gapinput@                         [1, 1, 1, 1],|
  !gapprompt@>| !gapinput@                         [1, 0, 1, 1]]);|
  Matrix(IsBooleanMat, [[1, 0, 0, 0], [0, 0, 0, 0], [1, 1, 1, 1],
    [1, 0, 1, 1]])
  !gapprompt@gap>| !gapinput@Matrix(IsMaxPlusMatrix, [[4, 0, -2],|
  !gapprompt@>| !gapinput@                            [1, -3, 0],|
  !gapprompt@>| !gapinput@                            [5, -1, -4]]);|
  Matrix(IsMaxPlusMatrix, [[4, 0, -2], [1, -3, 0], [5, -1, -4]])
  !gapprompt@gap>| !gapinput@Matrix(IsMinPlusMatrix, [[-1, infinity],|
  !gapprompt@>| !gapinput@                            [1, -1]]);|
  Matrix(IsMinPlusMatrix, [[-1, infinity], [1, -1]])
  !gapprompt@gap>| !gapinput@Matrix(IsTropicalMaxPlusMatrix, [[3, 2, 4],|
  !gapprompt@>| !gapinput@                                    [3, 1, 1],|
  !gapprompt@>| !gapinput@                                    [-infinity, 1, 1]],|
  !gapprompt@>| !gapinput@          9);|
  Matrix(IsTropicalMaxPlusMatrix, [[3, 2, 4], [3, 1, 1],
    [-infinity, 1, 1]], 9)
  !gapprompt@gap>| !gapinput@Matrix(IsTropicalMinPlusMatrix, [[1, 1, 1],|
  !gapprompt@>| !gapinput@                                    [0, 3, 0],|
  !gapprompt@>| !gapinput@                                    [1, 1, 3]],|
  !gapprompt@>| !gapinput@          9);|
  Matrix(IsTropicalMinPlusMatrix, [[1, 1, 1], [0, 3, 0], [1, 1, 3]], 9)
  !gapprompt@gap>| !gapinput@Matrix(IsNTPMatrix, [[0, 0, 0],|
  !gapprompt@>| !gapinput@                        [2, 0, 1],|
  !gapprompt@>| !gapinput@                        [2, 2, 2]],|
  !gapprompt@>| !gapinput@          2, 1);|
  Matrix(IsNTPMatrix, [[0, 0, 0], [2, 0, 1], [2, 2, 2]], 2, 1)
  !gapprompt@gap>| !gapinput@Matrix(Integers, [[-1, -2, 0],|
  !gapprompt@>| !gapinput@                     [0, 3, -1],|
  !gapprompt@>| !gapinput@                     [1, 0, -3]]);|
  <3x3-matrix over Integers>
  !gapprompt@gap>| !gapinput@Matrix(Integers, [[-1, -2, 0],|
  !gapprompt@>| !gapinput@                     [0, 3, -1],|
  !gapprompt@>| !gapinput@                     [1, 0, -3]]);|
  <3x3-matrix over Integers>
\end{Verbatim}
  }

 

\subsection{\textcolor{Chapter }{AsMatrix (for a filter and a matrix)}}
\logpage{[ 5, 1, 6 ]}\nobreak
\hyperdef{L}{X85426D8885431ECE}{}
{\noindent\textcolor{FuncColor}{$\triangleright$\enspace\texttt{AsMatrix({\mdseries\slshape filt, mat})\index{AsMatrix@\texttt{AsMatrix}!for a filter and a matrix}
\label{AsMatrix:for a filter and a matrix}
}\hfill{\scriptsize (operation)}}\\
\noindent\textcolor{FuncColor}{$\triangleright$\enspace\texttt{AsMatrix({\mdseries\slshape filt, mat, threshold})\index{AsMatrix@\texttt{AsMatrix}!for a filter, matrix, and threshold}
\label{AsMatrix:for a filter, matrix, and threshold}
}\hfill{\scriptsize (operation)}}\\
\noindent\textcolor{FuncColor}{$\triangleright$\enspace\texttt{AsMatrix({\mdseries\slshape filt, mat, threshold, period})\index{AsMatrix@\texttt{AsMatrix}!for a filter, matrix, threshold, and period}
\label{AsMatrix:for a filter, matrix, threshold, and period}
}\hfill{\scriptsize (operation)}}\\
\textbf{\indent Returns:\ }
A matrix.



 This operation can be used to change the representation of certain matrices
over semirings. If \mbox{\texttt{\mdseries\slshape mat}} is a matrix over a semiring (in the category \texttt{IsMatrixOverSemiring} (\ref{IsMatrixOverSemiring})), then \texttt{AsMatrix} returns a new matrix corresponding to \mbox{\texttt{\mdseries\slshape mat}} of the type specified by the filter \mbox{\texttt{\mdseries\slshape filt}}, and if applicable the arguments \mbox{\texttt{\mdseries\slshape threshold}} and \mbox{\texttt{\mdseries\slshape period}}. The dimension of the matrix \mbox{\texttt{\mdseries\slshape mat}} is not changed by this operation. 

 The version of the operation with arguments \mbox{\texttt{\mdseries\slshape filt}} and \mbox{\texttt{\mdseries\slshape mat}} can be applied to: 
\begin{itemize}
\item  \texttt{IsMinPlusMatrix} (\ref{IsMinPlusMatrix}) and a tropical min\texttt{\symbol{45}}plus matrix (i.e. convert a tropical
min\texttt{\symbol{45}}plus matrix to a (non\texttt{\symbol{45}}tropical)
min\texttt{\symbol{45}}plus matrix); 
\item  \texttt{IsMaxPlusMatrix} (\ref{IsMaxPlusMatrix}) and a tropical max\texttt{\symbol{45}}plus matrix;  
\end{itemize}
 The version of the operation with arguments \mbox{\texttt{\mdseries\slshape filt}}, \mbox{\texttt{\mdseries\slshape mat}}, and \mbox{\texttt{\mdseries\slshape threshold}} can be applied to: 
\begin{itemize}
\item  \texttt{IsTropicalMinPlusMatrix} (\ref{IsTropicalMinPlusMatrix}), a tropical min\texttt{\symbol{45}}plus or min\texttt{\symbol{45}}plus
matrix, and a value for the threshold of the resulting matrix. 
\item  \texttt{IsTropicalMaxPlusMatrix} (\ref{IsTropicalMaxPlusMatrix}) and a tropical max\texttt{\symbol{45}}plus,  or max\texttt{\symbol{45}}plus matrix, and a value for the threshold of the
resulting matrix. 
\end{itemize}
 The version of the operation with arguments \mbox{\texttt{\mdseries\slshape filt}}, \mbox{\texttt{\mdseries\slshape mat}}, \mbox{\texttt{\mdseries\slshape threshold}}, and \mbox{\texttt{\mdseries\slshape period}} can be applied to \texttt{IsNTPMatrix} (\ref{IsNTPMatrix}) and an ntp matrix, or integer matrix. 

 When converting matrices with negative entries to an ntp, tropical
max\texttt{\symbol{45}}plus, or tropical min\texttt{\symbol{45}}plus matrix,
the entry is replaced with its absolute value. 

 When converting non\texttt{\symbol{45}}tropical matrices to tropical matrices
entries higher than the specified threshold are reduced to the threshold. 
\begin{Verbatim}[commandchars=!@|,fontsize=\small,frame=single,label=Example]
  !gapprompt@gap>| !gapinput@mat := Matrix(IsTropicalMinPlusMatrix, [[0, 1, 3],|
  !gapprompt@>| !gapinput@                                           [1, 1, 6],|
  !gapprompt@>| !gapinput@                                           [0, 4, 2]], 10);;|
  !gapprompt@gap>| !gapinput@AsMatrix(IsMinPlusMatrix, mat);|
  Matrix(IsMinPlusMatrix, [[0, 1, 3], [1, 1, 6], [0, 4, 2]])
  !gapprompt@gap>| !gapinput@mat := Matrix(IsTropicalMaxPlusMatrix, [[-infinity, -infinity, 3],|
  !gapprompt@>| !gapinput@                                           [0, 1, 3],|
  !gapprompt@>| !gapinput@                                           [4, 1, 0]], 10);;|
  !gapprompt@gap>| !gapinput@AsMatrix(IsMaxPlusMatrix, mat);|
  Matrix(IsMaxPlusMatrix, [[-infinity, -infinity, 3], [0, 1, 3],
    [4, 1, 0]])
  !gapprompt@gap>| !gapinput@mat := Matrix(IsNTPMatrix, [[1, 2, 2],|
  !gapprompt@>| !gapinput@                               [0, 2, 0],|
  !gapprompt@>| !gapinput@                               [1, 3, 0]], 4, 5);;|
  !gapprompt@gap>| !gapinput@Matrix(Integers, mat);|
  <3x3-matrix over Integers>
  !gapprompt@gap>| !gapinput@mat := Matrix(IsMinPlusMatrix, [[0, 1, 3], [1, 1, 6], [0, 4, 2]]);;|
  !gapprompt@gap>| !gapinput@mat := AsMatrix(IsTropicalMinPlusMatrix, mat, 2);|
  Matrix(IsTropicalMinPlusMatrix, [[0, 1, 2], [1, 1, 2], [0, 2, 2]], 2)
  !gapprompt@gap>| !gapinput@mat := AsMatrix(IsTropicalMinPlusMatrix, mat, 1);|
  Matrix(IsTropicalMinPlusMatrix, [[0, 1, 1], [1, 1, 1], [0, 1, 1]], 1)
  !gapprompt@gap>| !gapinput@mat := Matrix(IsTropicalMaxPlusMatrix, [[-infinity, -infinity, 3],|
  !gapprompt@>| !gapinput@                                           [0, 1, 3],|
  !gapprompt@>| !gapinput@                                           [4, 1, 0]], 10);;|
  !gapprompt@gap>| !gapinput@AsMatrix(IsTropicalMaxPlusMatrix, mat, 4);|
  Matrix(IsTropicalMaxPlusMatrix, [[-infinity, -infinity, 3],
    [0, 1, 3], [4, 1, 0]], 4)
  !gapprompt@gap>| !gapinput@mat := Matrix(IsMaxPlusMatrix, [[-infinity, -infinity, 3],|
  !gapprompt@>| !gapinput@                                   [0, 1, 3],|
  !gapprompt@>| !gapinput@                                   [4, 1, 0]]);;|
  !gapprompt@gap>| !gapinput@AsMatrix(IsTropicalMaxPlusMatrix, mat, 10);|
  Matrix(IsTropicalMaxPlusMatrix, [[-infinity, -infinity, 3],
    [0, 1, 3], [4, 1, 0]], 10)
  !gapprompt@gap>| !gapinput@mat := Matrix(IsNTPMatrix, [[0, 1, 0],|
  !gapprompt@>| !gapinput@                               [1, 3, 1],|
  !gapprompt@>| !gapinput@                               [1, 0, 1]], 10, 10);;|
  !gapprompt@gap>| !gapinput@mat := AsMatrix(IsNTPMatrix, mat, 5, 6);|
  Matrix(IsNTPMatrix, [[0, 1, 0], [1, 3, 1], [1, 0, 1]], 5, 6)
  !gapprompt@gap>| !gapinput@mat := AsMatrix(IsNTPMatrix, mat, 2, 6);|
  Matrix(IsNTPMatrix, [[0, 1, 0], [1, 3, 1], [1, 0, 1]], 2, 6)
  !gapprompt@gap>| !gapinput@mat := AsMatrix(IsNTPMatrix, mat, 2, 1);|
  Matrix(IsNTPMatrix, [[0, 1, 0], [1, 2, 1], [1, 0, 1]], 2, 1)
  !gapprompt@gap>| !gapinput@mat := Matrix(Integers, mat);|
  <3x3-matrix over Integers>
  !gapprompt@gap>| !gapinput@AsMatrix(IsNTPMatrix, mat, 1, 2);|
  Matrix(IsNTPMatrix, [[0, 1, 0], [1, 2, 1], [1, 0, 1]], 1, 2)
\end{Verbatim}
 }

 

\subsection{\textcolor{Chapter }{RandomMatrix (for a filter and a matrix)}}
\logpage{[ 5, 1, 7 ]}\nobreak
\hyperdef{L}{X82172D747D66C8CC}{}
{\noindent\textcolor{FuncColor}{$\triangleright$\enspace\texttt{RandomMatrix({\mdseries\slshape filt, dim[, threshold[, period]]})\index{RandomMatrix@\texttt{RandomMatrix}!for a filter and a matrix}
\label{RandomMatrix:for a filter and a matrix}
}\hfill{\scriptsize (function)}}\\
\noindent\textcolor{FuncColor}{$\triangleright$\enspace\texttt{RandomMatrix({\mdseries\slshape semiring, dim})\index{RandomMatrix@\texttt{RandomMatrix}!for a semiring and a matrix}
\label{RandomMatrix:for a semiring and a matrix}
}\hfill{\scriptsize (function)}}\\
\textbf{\indent Returns:\ }
A matrix over semiring.



 This operation can be used to construct a random matrix over a semiring in \textsf{Semigroups}. The usage of \texttt{RandomMatrix} is similar to that of \texttt{Matrix} (\ref{Matrix:for a filter and a matrix}). 

 In its first form, the first argument \mbox{\texttt{\mdseries\slshape filt}} specifies the filter to be used to create the matrix, the second argument \mbox{\texttt{\mdseries\slshape dim}} is dimension of the matrix, the third and fourth arguments \mbox{\texttt{\mdseries\slshape threshold}} and \mbox{\texttt{\mdseries\slshape period}} (if required) must be positive integers. 

 
\begin{description}
\item[{\mbox{\texttt{\mdseries\slshape filt}}}]  This must be one of the filters given in Section \ref{IsXMatrix}. 
\item[{\mbox{\texttt{\mdseries\slshape dim}}}]  This must be a positive integer. 
\item[{\mbox{\texttt{\mdseries\slshape threshold}}}]  If \mbox{\texttt{\mdseries\slshape filt}} is any of \texttt{IsTropicalMaxPlusMatrix} (\ref{IsTropicalMaxPlusMatrix}), \texttt{IsTropicalMinPlusMatrix} (\ref{IsTropicalMinPlusMatrix}), or \texttt{IsNTPMatrix} (\ref{IsNTPMatrix}), then this argument specifies the threshold of the underlying semiring of the
matrix being created. 
\item[{\mbox{\texttt{\mdseries\slshape period}}}]  If \mbox{\texttt{\mdseries\slshape filt}} is \texttt{IsNTPMatrix} (\ref{IsNTPMatrix}), then this argument specifies the period of the underlying semiring of the
matrix being created. 
\end{description}
 In its second form, the arguments should be a semiring \mbox{\texttt{\mdseries\slshape semiring}} and dimension \mbox{\texttt{\mdseries\slshape dim}}. Currently, the only supported semirings are finite fields of prime order and
the integers \texttt{Integers} (\textbf{Reference: Integers}). 
\begin{Verbatim}[commandchars=!@|,fontsize=\small,frame=single,label=Example]
  !gapprompt@gap>| !gapinput@RandomMatrix(IsBooleanMat, 3);|
  Matrix(IsBooleanMat, [[1, 0, 0], [1, 0, 1], [1, 0, 1]])
  !gapprompt@gap>| !gapinput@RandomMatrix(IsMaxPlusMatrix, 2);|
  Matrix(IsMaxPlusMatrix, [[1, -infinity], [1, 0]])
  !gapprompt@gap>| !gapinput@RandomMatrix(IsMinPlusMatrix, 3);|
  Matrix(IsMinPlusMatrix, [[infinity, 2, infinity], [4, 0, -2], [1, -3, 0]])
  !gapprompt@gap>| !gapinput@RandomMatrix(IsTropicalMaxPlusMatrix, 3, 5);|
  Matrix(IsTropicalMaxPlusMatrix, [[5, 1, 4], [1, -infinity, 1], [1, 0, 2]],
    5)
  !gapprompt@gap>| !gapinput@RandomMatrix(IsTropicalMinPlusMatrix, 3, 2);|
  Matrix(IsTropicalMinPlusMatrix, [[1, -infinity, -infinity], [1, 1, 1],
    [2, 2, 1]], 2)
  !gapprompt@gap>| !gapinput@RandomMatrix(IsNTPMatrix, 3, 2, 5);|
  Matrix(IsNTPMatrix, [[1, 1, 1], [1, 1, 0], [3, 0, 1]], 2, 5)
  !gapprompt@gap>| !gapinput@RandomMatrix(Integers, 2);|
  Matrix(Integers, [[1, 3], [0, 0]])
  !gapprompt@gap>| !gapinput@RandomMatrix(Integers, 2);|
  Matrix(Integers, [[-1, 0], [0, -1]])
  !gapprompt@gap>| !gapinput@RandomMatrix(GF(5), 1);|
  Matrix(GF(5), [[Z(5)^0]])
\end{Verbatim}
 }

 
\subsection{\textcolor{Chapter }{Matrix filters}}\label{IsXMatrix}
\logpage{[ 5, 1, 8 ]}
\hyperdef{L}{X782480C686F1A663}{}
{
\noindent\textcolor{FuncColor}{$\triangleright$\enspace\texttt{IsBooleanMat({\mdseries\slshape obj})\index{IsBooleanMat@\texttt{IsBooleanMat}}
\label{IsBooleanMat}
}\hfill{\scriptsize (Category)}}\\
\noindent\textcolor{FuncColor}{$\triangleright$\enspace\texttt{IsMaxPlusMatrix({\mdseries\slshape obj})\index{IsMaxPlusMatrix@\texttt{IsMaxPlusMatrix}}
\label{IsMaxPlusMatrix}
}\hfill{\scriptsize (Category)}}\\
\noindent\textcolor{FuncColor}{$\triangleright$\enspace\texttt{IsMinPlusMatrix({\mdseries\slshape obj})\index{IsMinPlusMatrix@\texttt{IsMinPlusMatrix}}
\label{IsMinPlusMatrix}
}\hfill{\scriptsize (Category)}}\\
\noindent\textcolor{FuncColor}{$\triangleright$\enspace\texttt{IsTropicalMatrix({\mdseries\slshape obj})\index{IsTropicalMatrix@\texttt{IsTropicalMatrix}}
\label{IsTropicalMatrix}
}\hfill{\scriptsize (Category)}}\\
\noindent\textcolor{FuncColor}{$\triangleright$\enspace\texttt{IsTropicalMaxPlusMatrix({\mdseries\slshape obj})\index{IsTropicalMaxPlusMatrix@\texttt{IsTropicalMaxPlusMatrix}}
\label{IsTropicalMaxPlusMatrix}
}\hfill{\scriptsize (Category)}}\\
\noindent\textcolor{FuncColor}{$\triangleright$\enspace\texttt{IsTropicalMinPlusMatrix({\mdseries\slshape obj})\index{IsTropicalMinPlusMatrix@\texttt{IsTropicalMinPlusMatrix}}
\label{IsTropicalMinPlusMatrix}
}\hfill{\scriptsize (Category)}}\\
\noindent\textcolor{FuncColor}{$\triangleright$\enspace\texttt{IsNTPMatrix({\mdseries\slshape obj})\index{IsNTPMatrix@\texttt{IsNTPMatrix}}
\label{IsNTPMatrix}
}\hfill{\scriptsize (Category)}}\\
\noindent\textcolor{FuncColor}{$\triangleright$\enspace\texttt{Integers({\mdseries\slshape obj})\index{Integers@\texttt{Integers}}
\label{Integers}
}\hfill{\scriptsize (Category)}}\\
\textbf{\indent Returns:\ }
\texttt{true} or \texttt{false}.



 Every matrix over a semiring in \textsf{Semigroups} is a member of one of these categories, which are subcategory of \texttt{IsMatrixOverSemiring} (\ref{IsMatrixOverSemiring}). 

 \texttt{IsTropicalMatrix} is a supercategory of \texttt{IsTropicalMaxPlusMatrix} and \texttt{IsTropicalMinPlusMatrix}.  

 Basic operations for matrices over semirings include: multiplication via
\texttt{\symbol{92}}*, \texttt{DimensionOfMatrixOverSemiring} (\ref{DimensionOfMatrixOverSemiring}), \texttt{One} (\textbf{Reference: One}), the underlying list of lists used to create the matrix can be accessed using \texttt{AsList} (\ref{AsList}), the rows of \texttt{mat} can be accessed using \texttt{mat[i]} where \texttt{i} is between \texttt{1} and the dimension of the matrix, it also possible to loop over the rows of a
matrix; for tropical matrices \texttt{ThresholdTropicalMatrix} (\ref{ThresholdTropicalMatrix}); for ntp matrices \texttt{ThresholdNTPMatrix} (\ref{ThresholdNTPMatrix}) and \texttt{PeriodNTPMatrix} (\ref{PeriodNTPMatrix}). 

 For matrices over finite fields see Section \ref{Matrices over finite fields}; for Boolean matrices more details can be found in Section \ref{Boolean matrices}. }

 
\subsection{\textcolor{Chapter }{Matrix collection filters}}\label{IsXMatrixCollection}
\logpage{[ 5, 1, 9 ]}
\hyperdef{L}{X86233A3E86512493}{}
{
\noindent\textcolor{FuncColor}{$\triangleright$\enspace\texttt{IsBooleanMatCollection({\mdseries\slshape obj})\index{IsBooleanMatCollection@\texttt{IsBooleanMatCollection}}
\label{IsBooleanMatCollection}
}\hfill{\scriptsize (Category)}}\\
\noindent\textcolor{FuncColor}{$\triangleright$\enspace\texttt{IsBooleanMatCollColl({\mdseries\slshape obj})\index{IsBooleanMatCollColl@\texttt{IsBooleanMatCollColl}}
\label{IsBooleanMatCollColl}
}\hfill{\scriptsize (Category)}}\\
\noindent\textcolor{FuncColor}{$\triangleright$\enspace\texttt{IsMatrixOverFiniteFieldCollection({\mdseries\slshape obj})\index{IsMatrixOverFiniteFieldCollection@\texttt{IsMatrixOverFiniteFieldCollection}}
\label{IsMatrixOverFiniteFieldCollection}
}\hfill{\scriptsize (Category)}}\\
\noindent\textcolor{FuncColor}{$\triangleright$\enspace\texttt{IsMatrixOverFiniteFieldCollColl({\mdseries\slshape obj})\index{IsMatrixOverFiniteFieldCollColl@\texttt{IsMatrixOverFiniteFieldCollColl}}
\label{IsMatrixOverFiniteFieldCollColl}
}\hfill{\scriptsize (Category)}}\\
\noindent\textcolor{FuncColor}{$\triangleright$\enspace\texttt{IsMaxPlusMatrixCollection({\mdseries\slshape obj})\index{IsMaxPlusMatrixCollection@\texttt{IsMaxPlusMatrixCollection}}
\label{IsMaxPlusMatrixCollection}
}\hfill{\scriptsize (Category)}}\\
\noindent\textcolor{FuncColor}{$\triangleright$\enspace\texttt{IsMaxPlusMatrixCollColl({\mdseries\slshape obj})\index{IsMaxPlusMatrixCollColl@\texttt{IsMaxPlusMatrixCollColl}}
\label{IsMaxPlusMatrixCollColl}
}\hfill{\scriptsize (Category)}}\\
\noindent\textcolor{FuncColor}{$\triangleright$\enspace\texttt{IsMinPlusMatrixCollection({\mdseries\slshape obj})\index{IsMinPlusMatrixCollection@\texttt{IsMinPlusMatrixCollection}}
\label{IsMinPlusMatrixCollection}
}\hfill{\scriptsize (Category)}}\\
\noindent\textcolor{FuncColor}{$\triangleright$\enspace\texttt{IsMinPlusMatrixCollColl({\mdseries\slshape obj})\index{IsMinPlusMatrixCollColl@\texttt{IsMinPlusMatrixCollColl}}
\label{IsMinPlusMatrixCollColl}
}\hfill{\scriptsize (Category)}}\\
\noindent\textcolor{FuncColor}{$\triangleright$\enspace\texttt{IsTropicalMatrixCollection({\mdseries\slshape obj})\index{IsTropicalMatrixCollection@\texttt{IsTropicalMatrixCollection}}
\label{IsTropicalMatrixCollection}
}\hfill{\scriptsize (Category)}}\\
\noindent\textcolor{FuncColor}{$\triangleright$\enspace\texttt{IsTropicalMaxPlusMatrixCollection({\mdseries\slshape obj})\index{IsTropicalMaxPlusMatrixCollection@\texttt{IsTropicalMaxPlusMatrixCollection}}
\label{IsTropicalMaxPlusMatrixCollection}
}\hfill{\scriptsize (Category)}}\\
\noindent\textcolor{FuncColor}{$\triangleright$\enspace\texttt{IsTropicalMaxPlusMatrixCollColl({\mdseries\slshape obj})\index{IsTropicalMaxPlusMatrixCollColl@\texttt{IsTropicalMaxPlusMatrixCollColl}}
\label{IsTropicalMaxPlusMatrixCollColl}
}\hfill{\scriptsize (Category)}}\\
\noindent\textcolor{FuncColor}{$\triangleright$\enspace\texttt{IsTropicalMinPlusMatrixCollection({\mdseries\slshape obj})\index{IsTropicalMinPlusMatrixCollection@\texttt{IsTropicalMinPlusMatrixCollection}}
\label{IsTropicalMinPlusMatrixCollection}
}\hfill{\scriptsize (Category)}}\\
\noindent\textcolor{FuncColor}{$\triangleright$\enspace\texttt{IsTropicalMinPlusMatrixCollColl({\mdseries\slshape obj})\index{IsTropicalMinPlusMatrixCollColl@\texttt{IsTropicalMinPlusMatrixCollColl}}
\label{IsTropicalMinPlusMatrixCollColl}
}\hfill{\scriptsize (Category)}}\\
\noindent\textcolor{FuncColor}{$\triangleright$\enspace\texttt{IsNTPMatrixCollection({\mdseries\slshape obj})\index{IsNTPMatrixCollection@\texttt{IsNTPMatrixCollection}}
\label{IsNTPMatrixCollection}
}\hfill{\scriptsize (Category)}}\\
\noindent\textcolor{FuncColor}{$\triangleright$\enspace\texttt{IsNTPMatrixCollColl({\mdseries\slshape obj})\index{IsNTPMatrixCollColl@\texttt{IsNTPMatrixCollColl}}
\label{IsNTPMatrixCollColl}
}\hfill{\scriptsize (Category)}}\\
\textbf{\indent Returns:\ }
 \texttt{true} or \texttt{false}. 



 Every collection of matrices over the same semiring in \textsf{Semigroups} belongs to one of the categories above. For example, semigroups of boolean
matrices belong to \texttt{IsBooleanMatCollection}. 

 Similarly, every collection of collections of matrices over the same semiring
in \textsf{Semigroups} belongs to one of the categories above. }

 

\subsection{\textcolor{Chapter }{AsList}}
\logpage{[ 5, 1, 10 ]}\nobreak
\hyperdef{L}{X8289FCCC8274C89D}{}
{\noindent\textcolor{FuncColor}{$\triangleright$\enspace\texttt{AsList({\mdseries\slshape mat})\index{AsList@\texttt{AsList}}
\label{AsList}
}\hfill{\scriptsize (attribute)}}\\
\noindent\textcolor{FuncColor}{$\triangleright$\enspace\texttt{AsMutableList({\mdseries\slshape mat})\index{AsMutableList@\texttt{AsMutableList}}
\label{AsMutableList}
}\hfill{\scriptsize (operation)}}\\
\textbf{\indent Returns:\ }
A list of lists.



 If \mbox{\texttt{\mdseries\slshape mat}} is a matrix over a semiring (in the category \texttt{IsMatrixOverSemiring} (\ref{IsMatrixOverSemiring})), then \texttt{AsList} returns the underlying list of lists of semiring elements corresponding to \mbox{\texttt{\mdseries\slshape mat}}. In this case, the returned list and all of its entries are immutable. 

 The operation \texttt{AsMutableList} returns a mutable copy of the underlying list of lists of the matrix over
semiring \mbox{\texttt{\mdseries\slshape mat}}. 
\begin{Verbatim}[commandchars=!@|,fontsize=\small,frame=single,label=Example]
  !gapprompt@gap>| !gapinput@mat := Matrix(IsNTPMatrix,|
  !gapprompt@>| !gapinput@                 [[0, 1, 0], [1, 3, 1], [1, 0, 1]], 5, 6);|
  Matrix(IsNTPMatrix, [[0, 1, 0], [1, 3, 1], [1, 0, 1]], 5, 6)
  !gapprompt@gap>| !gapinput@list := AsList(mat);|
  [ [ 0, 1, 0 ], [ 1, 3, 1 ], [ 1, 0, 1 ] ]
  !gapprompt@gap>| !gapinput@IsMutable(list);|
  false
  !gapprompt@gap>| !gapinput@IsMutable(list[1]);|
  false
  !gapprompt@gap>| !gapinput@list := AsMutableList(mat);|
  [ [ 0, 1, 0 ], [ 1, 3, 1 ], [ 1, 0, 1 ] ]
  !gapprompt@gap>| !gapinput@IsMutable(list);|
  true
  !gapprompt@gap>| !gapinput@IsMutable(list[1]);|
  true
  !gapprompt@gap>| !gapinput@mat = Matrix(IsNTPMatrix, AsList(mat), 5, 6);|
  true
\end{Verbatim}
 }

  

\subsection{\textcolor{Chapter }{ThresholdTropicalMatrix}}
\logpage{[ 5, 1, 11 ]}\nobreak
\hyperdef{L}{X7D21408E845E4648}{}
{\noindent\textcolor{FuncColor}{$\triangleright$\enspace\texttt{ThresholdTropicalMatrix({\mdseries\slshape mat})\index{ThresholdTropicalMatrix@\texttt{ThresholdTropicalMatrix}}
\label{ThresholdTropicalMatrix}
}\hfill{\scriptsize (attribute)}}\\
\textbf{\indent Returns:\ }
A positive integer.



 If \mbox{\texttt{\mdseries\slshape mat}} is a tropical matrix (i.e. belongs to the category \texttt{IsTropicalMatrix} (\ref{IsTropicalMatrix})), then \texttt{ThresholdTropicalMatrix} returns the threshold (i.e. the largest integer) of the underlying semiring. 
\begin{Verbatim}[commandchars=!@|,fontsize=\small,frame=single,label=Example]
  !gapprompt@gap>| !gapinput@mat := Matrix(IsTropicalMaxPlusMatrix,|
  !gapprompt@>| !gapinput@[[0, 3, 0, 2],|
  !gapprompt@>| !gapinput@ [1, 1, 1, 0],|
  !gapprompt@>| !gapinput@ [-infinity, 1, -infinity, 1],|
  !gapprompt@>| !gapinput@ [0, -infinity, 2, -infinity]], 10);|
  Matrix(IsTropicalMaxPlusMatrix, [[0, 3, 0, 2], [1, 1, 1, 0],
    [-infinity, 1, -infinity, 1], [0, -infinity, 2, -infinity]], 10)
  !gapprompt@gap>| !gapinput@ThresholdTropicalMatrix(mat);|
  10
  !gapprompt@gap>| !gapinput@mat := Matrix(IsTropicalMaxPlusMatrix,|
  !gapprompt@>| !gapinput@[[0, 3, 0, 2],|
  !gapprompt@>| !gapinput@ [1, 1, 1, 0],|
  !gapprompt@>| !gapinput@ [-infinity, 1, -infinity, 1],|
  !gapprompt@>| !gapinput@ [0, -infinity, 2, -infinity]], 3);|
  Matrix(IsTropicalMaxPlusMatrix, [[0, 3, 0, 2], [1, 1, 1, 0],
    [-infinity, 1, -infinity, 1], [0, -infinity, 2, -infinity]], 3)
  !gapprompt@gap>| !gapinput@ThresholdTropicalMatrix(mat);|
  3
\end{Verbatim}
 }

 

\subsection{\textcolor{Chapter }{ThresholdNTPMatrix}}
\logpage{[ 5, 1, 12 ]}\nobreak
\hyperdef{L}{X7874559881FE8779}{}
{\noindent\textcolor{FuncColor}{$\triangleright$\enspace\texttt{ThresholdNTPMatrix({\mdseries\slshape mat})\index{ThresholdNTPMatrix@\texttt{ThresholdNTPMatrix}}
\label{ThresholdNTPMatrix}
}\hfill{\scriptsize (attribute)}}\\
\noindent\textcolor{FuncColor}{$\triangleright$\enspace\texttt{PeriodNTPMatrix({\mdseries\slshape mat})\index{PeriodNTPMatrix@\texttt{PeriodNTPMatrix}}
\label{PeriodNTPMatrix}
}\hfill{\scriptsize (attribute)}}\\
\textbf{\indent Returns:\ }
A positive integer.



 An \textsc{ntp matrix} is a matrix with entries in a semiring $\mathbb{N}_{t,p} = \{0, 1, \ldots, t, t + 1, \ldots, t + p - 1\}$ for some threshold $t$ and period $p$ under addition and multiplication modulo the congruence $t = t + p$. 

 If \mbox{\texttt{\mdseries\slshape mat}} is a ntp matrix (i.e. belongs to the category \texttt{IsNTPMatrix} (\ref{IsNTPMatrix})), then \texttt{ThresholdNTPMatrix} and \texttt{PeriodNTPMatrix} return the threshold and period of the underlying semiring, respectively. 
\begin{Verbatim}[commandchars=!@|,fontsize=\small,frame=single,label=Example]
  !gapprompt@gap>| !gapinput@mat := Matrix(IsNTPMatrix, [[1, 1, 0],|
  !gapprompt@>| !gapinput@                               [2, 1, 0],|
  !gapprompt@>| !gapinput@                               [0, 1, 1]],|
  !gapprompt@>| !gapinput@                 1, 2);|
  Matrix(IsNTPMatrix, [[1, 1, 0], [2, 1, 0], [0, 1, 1]], 1, 2)
  !gapprompt@gap>| !gapinput@ThresholdNTPMatrix(mat);|
  1
  !gapprompt@gap>| !gapinput@PeriodNTPMatrix(mat);|
  2
  !gapprompt@gap>| !gapinput@mat := Matrix(IsNTPMatrix, [[2, 1, 3],|
  !gapprompt@>| !gapinput@                               [0, 5, 1],|
  !gapprompt@>| !gapinput@                               [4, 1, 0]],|
  !gapprompt@>| !gapinput@                 3, 4);|
  Matrix(IsNTPMatrix, [[2, 1, 3], [0, 5, 1], [4, 1, 0]], 3, 4)
  !gapprompt@gap>| !gapinput@ThresholdNTPMatrix(mat);|
  3
  !gapprompt@gap>| !gapinput@PeriodNTPMatrix(mat);|
  4
\end{Verbatim}
 }

 }

   
\section{\textcolor{Chapter }{Operators for matrices over semirings}}\label{Operators for matrices over semirings}
\logpage{[ 5, 2, 0 ]}
\hyperdef{L}{X807E402687741CDA}{}
{
  
\begin{description}
\item[{\texttt{\mbox{\texttt{\mdseries\slshape mat1}} * \mbox{\texttt{\mdseries\slshape mat2}}}}]  \index{*@\texttt{*} (for matrices over a semiring)} returns the product of the matrices \mbox{\texttt{\mdseries\slshape mat1}} and \mbox{\texttt{\mdseries\slshape mat2}} of equal dimension over the same semiring using the usual matrix
multiplication with the operations \texttt{+} and \texttt{*} from the underlying semiring. 
\item[{\texttt{\mbox{\texttt{\mdseries\slshape mat1}} {\textless} \mbox{\texttt{\mdseries\slshape mat2}}}}]  \index{<@\texttt{{\textless}} (for matrices over a semiring)} returns \texttt{true} if when considered as a list of rows, the matrix \mbox{\texttt{\mdseries\slshape mat1}} is short\texttt{\symbol{45}}lex less than the matrix \mbox{\texttt{\mdseries\slshape mat2}}, and \texttt{false} if this is not the case. This means that a matrix of lower dimension is less
than a matrix of higher dimension. 
\item[{\texttt{\mbox{\texttt{\mdseries\slshape mat1}} = \mbox{\texttt{\mdseries\slshape mat2}}}}]  \index{=@\texttt{=} (for matrices over a semiring)} returns \texttt{true} if the matrix \mbox{\texttt{\mdseries\slshape mat1}} equals the matrix \mbox{\texttt{\mdseries\slshape mat2}} (i.e. the entries are equal and the underlying semirings are equal) and
returns \texttt{false} if it does not. 
\end{description}
 }

   
\section{\textcolor{Chapter }{ Boolean matrices }}\label{Boolean matrices}
\logpage{[ 5, 3, 0 ]}
\hyperdef{L}{X844A32A184E5EB75}{}
{
  In this section we describe the operations, properties, and attributes in \textsf{Semigroups} specifically for Boolean matrices. These include: 
\begin{itemize}
\item  \texttt{NumberBooleanMat} (\ref{NumberBooleanMat}) 
\item  \texttt{Successors} (\ref{Successors}) 
\item  \texttt{IsRowTrimBooleanMat} (\ref{IsRowTrimBooleanMat}), \texttt{IsColTrimBooleanMat} (\ref{IsColTrimBooleanMat}), and \texttt{IsTrimBooleanMat} (\ref{IsTrimBooleanMat}), 
\item  \texttt{CanonicalBooleanMat} (\ref{CanonicalBooleanMat}) 
\item  \texttt{IsSymmetricBooleanMat} (\ref{IsSymmetricBooleanMat}) 
\item  \texttt{IsAntiSymmetricBooleanMat} (\ref{IsAntiSymmetricBooleanMat}) 
\item  \texttt{IsTransitiveBooleanMat} (\ref{IsTransitiveBooleanMat}) 
\item  \texttt{IsReflexiveBooleanMat} (\ref{IsReflexiveBooleanMat}) 
\item  \texttt{IsTotalBooleanMat} (\ref{IsTotalBooleanMat}) 
\item  \texttt{IsOntoBooleanMat} (\ref{IsOntoBooleanMat}) 
\item  \texttt{IsPartialOrderBooleanMat} (\ref{IsPartialOrderBooleanMat}) 
\item  \texttt{IsEquivalenceBooleanMat} (\ref{IsEquivalenceBooleanMat}) 
\end{itemize}
 

\subsection{\textcolor{Chapter }{BooleanMat}}
\logpage{[ 5, 3, 1 ]}\nobreak
\hyperdef{L}{X84A16D4D7D015885}{}
{\noindent\textcolor{FuncColor}{$\triangleright$\enspace\texttt{BooleanMat({\mdseries\slshape arg})\index{BooleanMat@\texttt{BooleanMat}}
\label{BooleanMat}
}\hfill{\scriptsize (function)}}\\
\textbf{\indent Returns:\ }
A boolean matrix.



 \texttt{BooleanMat} returns the boolean matrix \texttt{mat} defined by its argument. The argument can be any of the following: 
\begin{description}
\item[{a matrix with entries \texttt{0} and/or \texttt{1}}]  the argument \mbox{\texttt{\mdseries\slshape arg}} is list of \texttt{n} lists of length \texttt{n} consisting of the values \texttt{0} and \texttt{1}; 
\item[{a matrix with entries \texttt{true} and/or \texttt{false}}]  the argument \mbox{\texttt{\mdseries\slshape arg}} is list of \texttt{n} lists of length \texttt{n} consisting of the values \texttt{true} and \texttt{false}; 
\item[{successors}]  the argument \mbox{\texttt{\mdseries\slshape arg}} is list of \texttt{n} sublists of consisting of positive integers not greater than \texttt{n}. In this case, the entry \texttt{j} in the sublist in position \texttt{i} of \mbox{\texttt{\mdseries\slshape arg}} indicates that the entry in position \texttt{(i, j)} of the created boolean matrix is \texttt{true}. 
\end{description}
 \texttt{BooleanMat} returns an error if the argument is not one of the above types. 
\begin{Verbatim}[commandchars=!@|,fontsize=\small,frame=single,label=Example]
  !gapprompt@gap>| !gapinput@x := BooleanMat([[true, false], [true, true]]);|
  Matrix(IsBooleanMat, [[1, 0], [1, 1]])
  !gapprompt@gap>| !gapinput@y := BooleanMat([[1, 0], [1, 1]]);|
  Matrix(IsBooleanMat, [[1, 0], [1, 1]])
  !gapprompt@gap>| !gapinput@z := BooleanMat([[1], [1, 2]]);|
  Matrix(IsBooleanMat, [[1, 0], [1, 1]])
  !gapprompt@gap>| !gapinput@x = y;|
  true
  !gapprompt@gap>| !gapinput@y = z;|
  true
  !gapprompt@gap>| !gapinput@Display(x);|
  1 0
  1 1
\end{Verbatim}
 }

 

\subsection{\textcolor{Chapter }{AsBooleanMat}}
\logpage{[ 5, 3, 2 ]}\nobreak
\hyperdef{L}{X7DA524567E0E7E16}{}
{\noindent\textcolor{FuncColor}{$\triangleright$\enspace\texttt{AsBooleanMat({\mdseries\slshape x[, n]})\index{AsBooleanMat@\texttt{AsBooleanMat}}
\label{AsBooleanMat}
}\hfill{\scriptsize (operation)}}\\
\textbf{\indent Returns:\ }
A boolean matrix.



 \texttt{AsBooleanMat} returns the pbr, bipartition, permutation, transformation, or partial
permutation \mbox{\texttt{\mdseries\slshape x}}, as a boolean matrix of dimension \mbox{\texttt{\mdseries\slshape n}}. 

 There are several possible arguments for \texttt{AsBooleanMat}: 
\begin{description}
\item[{permutations}]  If \mbox{\texttt{\mdseries\slshape x}} is a permutation and \mbox{\texttt{\mdseries\slshape n}} is a positive integer, then \texttt{AsBooleanMat(\mbox{\texttt{\mdseries\slshape x}}, \mbox{\texttt{\mdseries\slshape n}})} returns the boolean matrix \texttt{mat} of dimension \mbox{\texttt{\mdseries\slshape n}} such that \texttt{mat[i][j] = true} if and only if \texttt{j = i \texttt{\symbol{94}} x}. 

 If no positive integer \mbox{\texttt{\mdseries\slshape n}} is specified, then the largest moved point of \mbox{\texttt{\mdseries\slshape x}} is used as the value for \mbox{\texttt{\mdseries\slshape n}}; see \texttt{LargestMovedPoint} (\textbf{Reference: LargestMovedPoint for a permutation}). 
\item[{transformations}]  If \mbox{\texttt{\mdseries\slshape x}} is a transformation and \mbox{\texttt{\mdseries\slshape n}} is a positive integer such that \mbox{\texttt{\mdseries\slshape x}} is a transformation of \texttt{[1 .. \mbox{\texttt{\mdseries\slshape n}}]}, then \texttt{AsTransformation} returns the boolean matrix \texttt{mat} of dimension \mbox{\texttt{\mdseries\slshape n}} such that \texttt{mat[i][j] = true} if and only if \texttt{j = i \texttt{\symbol{94}} x}. 

 If the positive integer \mbox{\texttt{\mdseries\slshape n}} is not specified, then the degree of \mbox{\texttt{\mdseries\slshape f}} is used as the value for \mbox{\texttt{\mdseries\slshape n}}. 
\item[{partial permutations}]  If \mbox{\texttt{\mdseries\slshape x}} is a partial permutation and \mbox{\texttt{\mdseries\slshape n}} is a positive integer such that \texttt{i \texttt{\symbol{94}} \mbox{\texttt{\mdseries\slshape x}} {\textless}= n} for all \texttt{i} in \texttt{[1 .. \mbox{\texttt{\mdseries\slshape n}}]}, then \texttt{AsBooleanMat} returns the boolean matrix \texttt{mat} of dimension \mbox{\texttt{\mdseries\slshape n}} such that \texttt{mat[i][j] = true} if and only if \texttt{j = i \texttt{\symbol{94}} x}. 

 If the optional argument \mbox{\texttt{\mdseries\slshape n}} is not present, then the default value of the maximum of degree and the
codegree of \mbox{\texttt{\mdseries\slshape x}} is used. 
\item[{bipartitions}]  If \mbox{\texttt{\mdseries\slshape x}} is a bipartition and \mbox{\texttt{\mdseries\slshape n}} is any non\texttt{\symbol{45}}negative integer, then \texttt{AsBooleanMat} returns the boolean matrix \texttt{mat} of dimension \mbox{\texttt{\mdseries\slshape n}} such that \texttt{mat[i][j] = true} if and only if \texttt{i} and \texttt{j} belong to the same block of \mbox{\texttt{\mdseries\slshape x}}. 

 If the optional argument \mbox{\texttt{\mdseries\slshape n}} is not present, then twice the degree of \mbox{\texttt{\mdseries\slshape x}} is used by default. 
\item[{pbrs}]  If \mbox{\texttt{\mdseries\slshape x}} is a pbr and \mbox{\texttt{\mdseries\slshape n}} is any non\texttt{\symbol{45}}negative integer, then \texttt{AsBooleanMat} returns the boolean matrix \texttt{mat} of dimension \mbox{\texttt{\mdseries\slshape n}} such that \texttt{mat[i][j] = true} if and only if \texttt{i} and \texttt{j} are related in \mbox{\texttt{\mdseries\slshape x}}. 

 If the optional argument \mbox{\texttt{\mdseries\slshape n}} is not present, then twice the degree of \mbox{\texttt{\mdseries\slshape x}} is used by default. 
\end{description}
 
\begin{Verbatim}[commandchars=!@|,fontsize=\small,frame=single,label=Example]
  !gapprompt@gap>| !gapinput@Display(AsBooleanMat((1, 2), 5));|
  0 1 0 0 0
  1 0 0 0 0
  0 0 1 0 0
  0 0 0 1 0
  0 0 0 0 1
  !gapprompt@gap>| !gapinput@Display(AsBooleanMat((1, 2)));|
  0 1
  1 0
  !gapprompt@gap>| !gapinput@x := Transformation([1, 3, 4, 1, 3]);;|
  !gapprompt@gap>| !gapinput@Display(AsBooleanMat(x));|
  1 0 0 0 0
  0 0 1 0 0
  0 0 0 1 0
  1 0 0 0 0
  0 0 1 0 0
  !gapprompt@gap>| !gapinput@Display(AsBooleanMat(x, 4));|
  1 0 0 0
  0 0 1 0
  0 0 0 1
  1 0 0 0
  !gapprompt@gap>| !gapinput@x := PartialPerm([1, 2, 3, 6, 8, 10],|
  !gapprompt@>| !gapinput@                    [2, 6, 7, 9, 1, 5]);|
  [3,7][8,1,2,6,9][10,5]
  !gapprompt@gap>| !gapinput@Display(AsBooleanMat(x));|
  0 1 0 0 0 0 0 0 0 0
  0 0 0 0 0 1 0 0 0 0
  0 0 0 0 0 0 1 0 0 0
  0 0 0 0 0 0 0 0 0 0
  0 0 0 0 0 0 0 0 0 0
  0 0 0 0 0 0 0 0 1 0
  0 0 0 0 0 0 0 0 0 0
  1 0 0 0 0 0 0 0 0 0
  0 0 0 0 0 0 0 0 0 0
  0 0 0 0 1 0 0 0 0 0
  !gapprompt@gap>| !gapinput@x := Bipartition([[1, 4, -2, -3], [2, 3, 5, -5], [-1, -4]]);|
  <bipartition: [ 1, 4, -2, -3 ], [ 2, 3, 5, -5 ], [ -1, -4 ]>
  !gapprompt@gap>| !gapinput@y := AsBooleanMat(x);|
  <10x10 boolean matrix>
  !gapprompt@gap>| !gapinput@Display(y);|
  1 0 0 1 0 0 1 1 0 0
  0 1 1 0 1 0 0 0 0 1
  0 1 1 0 1 0 0 0 0 1
  1 0 0 1 0 0 1 1 0 0
  0 1 1 0 1 0 0 0 0 1
  0 0 0 0 0 1 0 0 1 0
  1 0 0 1 0 0 1 1 0 0
  1 0 0 1 0 0 1 1 0 0
  0 0 0 0 0 1 0 0 1 0
  0 1 1 0 1 0 0 0 0 1
  !gapprompt@gap>| !gapinput@IsEquivalenceBooleanMat(y);|
  true
  !gapprompt@gap>| !gapinput@AsBooleanMat(x, 1);|
  Matrix(IsBooleanMat, [[1]])
  !gapprompt@gap>| !gapinput@Display(AsBooleanMat(x, 1));|
  1
  !gapprompt@gap>| !gapinput@Display(AsBooleanMat(x, 2));|
  1 0
  0 1
  !gapprompt@gap>| !gapinput@Display(AsBooleanMat(x, 3));|
  1 0 0
  0 1 1
  0 1 1
  !gapprompt@gap>| !gapinput@Display(AsBooleanMat(x, 11));|
  1 0 0 1 0 0 1 1 0 0 0
  0 1 1 0 1 0 0 0 0 1 0
  0 1 1 0 1 0 0 0 0 1 0
  1 0 0 1 0 0 1 1 0 0 0
  0 1 1 0 1 0 0 0 0 1 0
  0 0 0 0 0 1 0 0 1 0 0
  1 0 0 1 0 0 1 1 0 0 0
  1 0 0 1 0 0 1 1 0 0 0
  0 0 0 0 0 1 0 0 1 0 0
  0 1 1 0 1 0 0 0 0 1 0
  0 0 0 0 0 0 0 0 0 0 0
  !gapprompt@gap>| !gapinput@x := PBR(|
  !gapprompt@>| !gapinput@[[-1, 1], [2, 3], [-3, 2, 3]],|
  !gapprompt@>| !gapinput@[[-1, 1, 2], [-3, -1, 1, 3], [-3, -1, 1, 2, 3]]);;|
  !gapprompt@gap>| !gapinput@AsBooleanMat(x);|
  Matrix(IsBooleanMat, [[1, 0, 0, 1, 0, 0], [0, 1, 1, 0, 0, 0],
    [0, 1, 1, 0, 0, 1], [1, 1, 0, 1, 0, 0], [1, 0, 1, 1, 0, 1],
    [1, 1, 1, 1, 0, 1]])
  !gapprompt@gap>| !gapinput@Display(AsBooleanMat(x));|
  1 0 0 1 0 0
  0 1 1 0 0 0
  0 1 1 0 0 1
  1 1 0 1 0 0
  1 0 1 1 0 1
  1 1 1 1 0 1
\end{Verbatim}
 }

 

\subsection{\textcolor{Chapter }{\texttt{\symbol{92}}in}}
\logpage{[ 5, 3, 3 ]}\nobreak
\hyperdef{L}{X87BDB89B7AAFE8AD}{}
{\noindent\textcolor{FuncColor}{$\triangleright$\enspace\texttt{\texttt{\symbol{92}}in({\mdseries\slshape mat1, mat2})\index{\texttt{\symbol{92}}in@\texttt{\texttt{\symbol{92}}in}}
\label{bSlashin}
}\hfill{\scriptsize (operation)}}\\
\textbf{\indent Returns:\ }
\texttt{true} or \texttt{false}.



 If \mbox{\texttt{\mdseries\slshape mat1}} and \mbox{\texttt{\mdseries\slshape mat2}} are boolean matrices, then \texttt{\mbox{\texttt{\mdseries\slshape mat1}} in \mbox{\texttt{\mdseries\slshape mat2}}} returns \texttt{true} if the binary relation defined by \mbox{\texttt{\mdseries\slshape mat1}} is a subset of that defined by \mbox{\texttt{\mdseries\slshape mat2}}. 
\begin{Verbatim}[commandchars=!@|,fontsize=\small,frame=single,label=Example]
  !gapprompt@gap>| !gapinput@x := BooleanMat([[1, 0, 0, 1], [0, 0, 0, 0],|
  !gapprompt@>| !gapinput@                    [1, 0, 1, 1], [0, 1, 1, 1]]);;|
  !gapprompt@gap>| !gapinput@y := BooleanMat([[1, 0, 1, 0], [1, 1, 1, 0],|
  !gapprompt@>| !gapinput@                    [0, 1, 1, 0], [1, 1, 1, 1]]);;|
  !gapprompt@gap>| !gapinput@x in y;|
  false
  !gapprompt@gap>| !gapinput@y in y;|
  true
\end{Verbatim}
 }

 

\subsection{\textcolor{Chapter }{OnBlist}}
\logpage{[ 5, 3, 4 ]}\nobreak
\hyperdef{L}{X8629FA5F7B682078}{}
{\noindent\textcolor{FuncColor}{$\triangleright$\enspace\texttt{OnBlist({\mdseries\slshape blist, mat})\index{OnBlist@\texttt{OnBlist}}
\label{OnBlist}
}\hfill{\scriptsize (function)}}\\
\textbf{\indent Returns:\ }
A boolean list.



 If \mbox{\texttt{\mdseries\slshape blist}} is a boolean list of length \texttt{n} and \mbox{\texttt{\mdseries\slshape mat}} is boolean matrices of dimension \texttt{n}, then \texttt{OnBlist} returns the product of \mbox{\texttt{\mdseries\slshape blist}} (thought of as a row vector over the boolean semiring) and \mbox{\texttt{\mdseries\slshape mat}}. 
\begin{Verbatim}[commandchars=!@|,fontsize=\small,frame=single,label=Example]
  !gapprompt@gap>| !gapinput@mat := BooleanMat([[1, 0, 0, 1],|
  !gapprompt@>| !gapinput@                      [0, 0, 0, 0],|
  !gapprompt@>| !gapinput@                      [1, 0, 1, 1],|
  !gapprompt@>| !gapinput@                      [0, 1, 1, 1]]);;|
  !gapprompt@gap>| !gapinput@blist := BlistList([1 .. 4], [1, 2]);|
  [ true, true, false, false ]
  !gapprompt@gap>| !gapinput@OnBlist(blist, mat);|
  [ true, false, false, true ]
\end{Verbatim}
 }

 

\subsection{\textcolor{Chapter }{Successors}}
\logpage{[ 5, 3, 5 ]}\nobreak
\hyperdef{L}{X85E2FD8B82652876}{}
{\noindent\textcolor{FuncColor}{$\triangleright$\enspace\texttt{Successors({\mdseries\slshape mat})\index{Successors@\texttt{Successors}}
\label{Successors}
}\hfill{\scriptsize (attribute)}}\\
\textbf{\indent Returns:\ }
A list of lists of positive integers.



 A row of a boolean matrix of dimension \texttt{n} can be thought of of as the characteristic function of a subset \texttt{S} of \texttt{[1 .. n]}, i.e. \texttt{i in S} if and only if the \texttt{i}th component of the row equals $1$. We refer to the subset \texttt{S} as the \textsc{successors} of the row. 

 If \mbox{\texttt{\mdseries\slshape mat}} is a boolean matrix, then \texttt{Successors} returns the list of successors of the rows of \mbox{\texttt{\mdseries\slshape mat}}. 
\begin{Verbatim}[commandchars=!@|,fontsize=\small,frame=single,label=Example]
  !gapprompt@gap>| !gapinput@mat := BooleanMat([[1, 0, 1, 1],|
  !gapprompt@>| !gapinput@                      [1, 0, 0, 0],|
  !gapprompt@>| !gapinput@                      [0, 0, 1, 0],|
  !gapprompt@>| !gapinput@                      [1, 1, 0, 0]]);;|
  !gapprompt@gap>| !gapinput@Successors(mat);|
  [ [ 1, 3, 4 ], [ 1 ], [ 3 ], [ 1, 2 ] ]
\end{Verbatim}
 }

 

\subsection{\textcolor{Chapter }{BooleanMatNumber}}
\logpage{[ 5, 3, 6 ]}\nobreak
\hyperdef{L}{X7E0FD5878106AB66}{}
{\noindent\textcolor{FuncColor}{$\triangleright$\enspace\texttt{BooleanMatNumber({\mdseries\slshape m, n})\index{BooleanMatNumber@\texttt{BooleanMatNumber}}
\label{BooleanMatNumber}
}\hfill{\scriptsize (operation)}}\\
\noindent\textcolor{FuncColor}{$\triangleright$\enspace\texttt{NumberBooleanMat({\mdseries\slshape mat})\index{NumberBooleanMat@\texttt{NumberBooleanMat}}
\label{NumberBooleanMat}
}\hfill{\scriptsize (operation)}}\\
\textbf{\indent Returns:\ }
A boolean matrix, or a positive integer.



 These functions implement a bijection from the set of all boolean matrices of
dimension \mbox{\texttt{\mdseries\slshape n}} and the numbers \texttt{[1 .. 2 \texttt{\symbol{94}} (\mbox{\texttt{\mdseries\slshape n}} \texttt{\symbol{94}} 2)]}. 

 More precisely, if \mbox{\texttt{\mdseries\slshape m}} and \mbox{\texttt{\mdseries\slshape n}} are positive integers such that \mbox{\texttt{\mdseries\slshape m}} is at most \texttt{2 \texttt{\symbol{94}} (\mbox{\texttt{\mdseries\slshape n}} \texttt{\symbol{94}} 2)}, then \texttt{BooleanMatNumber} returns the \mbox{\texttt{\mdseries\slshape m}}th \mbox{\texttt{\mdseries\slshape n}} by \mbox{\texttt{\mdseries\slshape n}} boolean matrix. 

 If \mbox{\texttt{\mdseries\slshape mat}} is an \mbox{\texttt{\mdseries\slshape n}} by \mbox{\texttt{\mdseries\slshape n}} boolean matrix, then \texttt{NumberBooleanMat} returns the number in \texttt{[1 .. 2 \texttt{\symbol{94}} (\mbox{\texttt{\mdseries\slshape n}} \texttt{\symbol{94}} 2)]} that corresponds to \mbox{\texttt{\mdseries\slshape mat}}. 
\begin{Verbatim}[commandchars=!@|,fontsize=\small,frame=single,label=Example]
  !gapprompt@gap>| !gapinput@mat := BooleanMat([[0, 1, 1, 0],|
  !gapprompt@>| !gapinput@                      [1, 0, 1, 1],|
  !gapprompt@>| !gapinput@                      [1, 1, 0, 1],|
  !gapprompt@>| !gapinput@                      [0, 1, 0, 1]]);;|
  !gapprompt@gap>| !gapinput@NumberBooleanMat(mat);|
  27606
  !gapprompt@gap>| !gapinput@Display(BooleanMatNumber(27606, 4));|
  0 1 1 0
  1 0 1 1
  1 1 0 1
  0 1 0 1
\end{Verbatim}
 }

 

\subsection{\textcolor{Chapter }{BlistNumber}}
\logpage{[ 5, 3, 7 ]}\nobreak
\hyperdef{L}{X793A1C277C1D7D6D}{}
{\noindent\textcolor{FuncColor}{$\triangleright$\enspace\texttt{BlistNumber({\mdseries\slshape m, n})\index{BlistNumber@\texttt{BlistNumber}}
\label{BlistNumber}
}\hfill{\scriptsize (function)}}\\
\noindent\textcolor{FuncColor}{$\triangleright$\enspace\texttt{NumberBlist({\mdseries\slshape blist})\index{NumberBlist@\texttt{NumberBlist}}
\label{NumberBlist}
}\hfill{\scriptsize (function)}}\\
\textbf{\indent Returns:\ }
A boolean list, or a positive integer.



 These functions implement a bijection from the set of all boolean lists of
length \mbox{\texttt{\mdseries\slshape n}} and the numbers \texttt{[1 .. 2 \texttt{\symbol{94}} \mbox{\texttt{\mdseries\slshape n}}]}. 

 More precisely, if \mbox{\texttt{\mdseries\slshape m}} and \mbox{\texttt{\mdseries\slshape n}} are positive integers such that \mbox{\texttt{\mdseries\slshape m}} is at most \texttt{2 \texttt{\symbol{94}} \mbox{\texttt{\mdseries\slshape n}}}, then \texttt{BlistNumber} returns the \mbox{\texttt{\mdseries\slshape m}}th boolean list of length \mbox{\texttt{\mdseries\slshape n}}. 

 If \mbox{\texttt{\mdseries\slshape blist}} is a boolean list of length \mbox{\texttt{\mdseries\slshape n}}, then \texttt{NumberBlist} returns the number in \texttt{[1 .. 2 \texttt{\symbol{94}} \mbox{\texttt{\mdseries\slshape n}}]} that corresponds to \mbox{\texttt{\mdseries\slshape blist}}. 
\begin{Verbatim}[commandchars=!@|,fontsize=\small,frame=single,label=Example]
  !gapprompt@gap>| !gapinput@blist := BlistList([1 .. 10], []);|
  [ false, false, false, false, false, false, false, false, false,
    false ]
  !gapprompt@gap>| !gapinput@NumberBlist(blist);|
  1
  !gapprompt@gap>| !gapinput@blist := BlistList([1 .. 10], [10]);|
  [ false, false, false, false, false, false, false, false, false, true
   ]
  !gapprompt@gap>| !gapinput@NumberBlist(blist);|
  2
  !gapprompt@gap>| !gapinput@BlistNumber(1, 10);|
  [ false, false, false, false, false, false, false, false, false,
    false ]
  !gapprompt@gap>| !gapinput@BlistNumber(2, 10);|
  [ false, false, false, false, false, false, false, false, false, true
   ]
\end{Verbatim}
 }

 

\subsection{\textcolor{Chapter }{CanonicalBooleanMat (for a perm group, perm group and boolean matrix)}}
\logpage{[ 5, 3, 8 ]}\nobreak
\hyperdef{L}{X7EEA5011862E6298}{}
{\noindent\textcolor{FuncColor}{$\triangleright$\enspace\texttt{CanonicalBooleanMat({\mdseries\slshape G, H, mat})\index{CanonicalBooleanMat@\texttt{CanonicalBooleanMat}!for a perm group, perm group and boolean matrix}
\label{CanonicalBooleanMat:for a perm group, perm group and boolean matrix}
}\hfill{\scriptsize (operation)}}\\
\noindent\textcolor{FuncColor}{$\triangleright$\enspace\texttt{CanonicalBooleanMat({\mdseries\slshape G, mat})\index{CanonicalBooleanMat@\texttt{CanonicalBooleanMat}!for a perm group and boolean matrix}
\label{CanonicalBooleanMat:for a perm group and boolean matrix}
}\hfill{\scriptsize (operation)}}\\
\noindent\textcolor{FuncColor}{$\triangleright$\enspace\texttt{CanonicalBooleanMat({\mdseries\slshape mat})\index{CanonicalBooleanMat@\texttt{CanonicalBooleanMat}}
\label{CanonicalBooleanMat}
}\hfill{\scriptsize (attribute)}}\\
\textbf{\indent Returns:\ }
A boolean matrix.



 This operation returns a fixed representative of the orbit of the boolean
matrix \mbox{\texttt{\mdseries\slshape mat}} under the action of the permutation group \mbox{\texttt{\mdseries\slshape G}} on its rows and the permutation group \mbox{\texttt{\mdseries\slshape H}} on its columns. 

 In its second form, when only a single permutation group \mbox{\texttt{\mdseries\slshape G}} is specified, \mbox{\texttt{\mdseries\slshape G}} acts on the rows and columns of \mbox{\texttt{\mdseries\slshape mat}} independently. 

 In its third form, when only a boolean matrix is specified, \texttt{CanonicalBooleanMat} returns a fixed representative of the orbit of \mbox{\texttt{\mdseries\slshape mat}} under the action of the symmetric group on its rows, and, independently, on
its columns. In other words, \texttt{CanonicalBooleanMat} returns a canonical boolean matrix equivalent to \mbox{\texttt{\mdseries\slshape mat}} up to rearranging rows and columns. This version of \texttt{CanonicalBooleanMat} uses \textsf{digraphs} and its interface with the \href{http://www.tcs.tkk.fi/Software/bliss/} {bliss}  library for computing automorphism groups and canonical forms of graphs \cite{JunttilaKaski}. As a consequence, \texttt{CanonicalBooleanMat} with a single argument is significantly faster than the versions with 2 or 3
arguments. 
\begin{Verbatim}[commandchars=!@|,fontsize=\small,frame=single,label=Example]
  !gapprompt@gap>| !gapinput@mat := BooleanMat([[1, 1, 1, 0, 0, 0],|
  !gapprompt@>| !gapinput@                      [0, 0, 0, 1, 0, 1],|
  !gapprompt@>| !gapinput@                      [1, 0, 0, 1, 0, 1],|
  !gapprompt@>| !gapinput@                      [0, 0, 0, 0, 0, 0],|
  !gapprompt@>| !gapinput@                      [0, 1, 1, 1, 1, 1],|
  !gapprompt@>| !gapinput@                      [0, 1, 1, 0, 1, 0]]);|
  Matrix(IsBooleanMat, [[1, 1, 1, 0, 0, 0], [0, 0, 0, 1, 0, 1],
    [1, 0, 0, 1, 0, 1], [0, 0, 0, 0, 0, 0], [0, 1, 1, 1, 1, 1],
    [0, 1, 1, 0, 1, 0]])
  !gapprompt@gap>| !gapinput@CanonicalBooleanMat(mat);|
  Matrix(IsBooleanMat, [[0, 0, 0, 0, 0, 0], [1, 1, 0, 0, 0, 0],
    [0, 0, 1, 1, 1, 0], [1, 1, 0, 0, 1, 0], [0, 0, 1, 1, 0, 1],
    [1, 1, 1, 1, 0, 1]])
  !gapprompt@gap>| !gapinput@Display(CanonicalBooleanMat(mat));|
  0 0 0 0 0 0
  1 1 0 0 0 0
  0 0 1 1 1 0
  1 1 0 0 1 0
  0 0 1 1 0 1
  1 1 1 1 0 1
  !gapprompt@gap>| !gapinput@Display(CanonicalBooleanMat(Group((1, 3)), mat));|
  0 1 1 0 0 1
  0 0 1 0 0 1
  1 1 0 1 0 0
  0 0 0 0 0 0
  1 0 1 1 1 1
  1 0 0 1 1 0
  !gapprompt@gap>| !gapinput@Display(CanonicalBooleanMat(Group((1, 3)), Group(()), mat));|
  1 1 1 0 0 0
  0 0 0 1 0 1
  0 1 0 1 0 1
  0 0 0 0 0 0
  1 0 1 1 1 1
  1 0 1 0 1 0
\end{Verbatim}
 }

 

\subsection{\textcolor{Chapter }{IsRowTrimBooleanMat}}
\logpage{[ 5, 3, 9 ]}\nobreak
\hyperdef{L}{X794C91597CC9F784}{}
{\noindent\textcolor{FuncColor}{$\triangleright$\enspace\texttt{IsRowTrimBooleanMat({\mdseries\slshape mat})\index{IsRowTrimBooleanMat@\texttt{IsRowTrimBooleanMat}}
\label{IsRowTrimBooleanMat}
}\hfill{\scriptsize (property)}}\\
\noindent\textcolor{FuncColor}{$\triangleright$\enspace\texttt{IsColTrimBooleanMat({\mdseries\slshape mat})\index{IsColTrimBooleanMat@\texttt{IsColTrimBooleanMat}}
\label{IsColTrimBooleanMat}
}\hfill{\scriptsize (property)}}\\
\noindent\textcolor{FuncColor}{$\triangleright$\enspace\texttt{IsTrimBooleanMat({\mdseries\slshape mat})\index{IsTrimBooleanMat@\texttt{IsTrimBooleanMat}}
\label{IsTrimBooleanMat}
}\hfill{\scriptsize (property)}}\\
\textbf{\indent Returns:\ }
\texttt{true} or \texttt{false}.



 A row or column of a boolean matrix of dimension \texttt{n} can be thought of of as the characteristic function of a subset \texttt{S} of \texttt{[1 .. n]}, i.e. \texttt{i in S} if and only if the \texttt{i}th component of the row or column equals \texttt{1}. 

 A boolean matrix is \textsc{row trim} if no subset induced by a row of \mbox{\texttt{\mdseries\slshape mat}} is contained in the subset induced by any other row of \mbox{\texttt{\mdseries\slshape mat}}. \textsc{Column trim} is defined analogously. A boolean matrix is \textsc{trim} if it is both row and column trim. 
\begin{Verbatim}[commandchars=!@|,fontsize=\small,frame=single,label=Example]
  !gapprompt@gap>| !gapinput@mat := BooleanMat([[0, 0, 1, 1],|
  !gapprompt@>| !gapinput@                      [1, 0, 1, 0],|
  !gapprompt@>| !gapinput@                      [1, 1, 0, 0],|
  !gapprompt@>| !gapinput@                      [0, 1, 0, 1]]);;|
  !gapprompt@gap>| !gapinput@IsTrimBooleanMat(mat);|
  true
  !gapprompt@gap>| !gapinput@mat := BooleanMat([[0, 1, 1, 0],|
  !gapprompt@>| !gapinput@                      [0, 0, 1, 0],|
  !gapprompt@>| !gapinput@                      [1, 0, 0, 1],|
  !gapprompt@>| !gapinput@                      [1, 0, 1, 0]]);;|
  !gapprompt@gap>| !gapinput@IsRowTrimBooleanMat(mat);|
  false
  !gapprompt@gap>| !gapinput@IsColTrimBooleanMat(mat);|
  false
\end{Verbatim}
 }

 

\subsection{\textcolor{Chapter }{IsSymmetricBooleanMat}}
\logpage{[ 5, 3, 10 ]}\nobreak
\hyperdef{L}{X7D22BA78790EFBC6}{}
{\noindent\textcolor{FuncColor}{$\triangleright$\enspace\texttt{IsSymmetricBooleanMat({\mdseries\slshape mat})\index{IsSymmetricBooleanMat@\texttt{IsSymmetricBooleanMat}}
\label{IsSymmetricBooleanMat}
}\hfill{\scriptsize (property)}}\\
\textbf{\indent Returns:\ }
\texttt{true} or \texttt{false}.



 A boolean matrix is \textsc{symmetric} if it is symmetric about the main diagonal, i.e. \texttt{\mbox{\texttt{\mdseries\slshape mat}}[i][j] = \mbox{\texttt{\mdseries\slshape mat}}[j][i]} for all \texttt{i, j} in the range \texttt{[1 .. n]} where \texttt{n} is the dimension of \mbox{\texttt{\mdseries\slshape mat}}. 
\begin{Verbatim}[commandchars=!@|,fontsize=\small,frame=single,label=Example]
  !gapprompt@gap>| !gapinput@mat := BooleanMat([[0, 1, 1, 0],|
  !gapprompt@>| !gapinput@                      [1, 0, 1, 1],|
  !gapprompt@>| !gapinput@                      [1, 1, 0, 1],|
  !gapprompt@>| !gapinput@                      [0, 1, 0, 1]]);|
  Matrix(IsBooleanMat, [[0, 1, 1, 0], [1, 0, 1, 1], [1, 1, 0, 1],
    [0, 1, 0, 1]])
  !gapprompt@gap>| !gapinput@IsSymmetricBooleanMat(mat);|
  false
  !gapprompt@gap>| !gapinput@mat := BooleanMat([[0, 1, 1, 0],|
  !gapprompt@>| !gapinput@                      [1, 0, 1, 1],|
  !gapprompt@>| !gapinput@                      [1, 1, 0, 1],|
  !gapprompt@>| !gapinput@                      [0, 1, 1, 1]]);|
  Matrix(IsBooleanMat, [[0, 1, 1, 0], [1, 0, 1, 1], [1, 1, 0, 1],
    [0, 1, 1, 1]])
  !gapprompt@gap>| !gapinput@IsSymmetricBooleanMat(mat);|
  true
\end{Verbatim}
 }

 

\subsection{\textcolor{Chapter }{IsReflexiveBooleanMat}}
\logpage{[ 5, 3, 11 ]}\nobreak
\hyperdef{L}{X7C373B7D87044050}{}
{\noindent\textcolor{FuncColor}{$\triangleright$\enspace\texttt{IsReflexiveBooleanMat({\mdseries\slshape mat})\index{IsReflexiveBooleanMat@\texttt{IsReflexiveBooleanMat}}
\label{IsReflexiveBooleanMat}
}\hfill{\scriptsize (property)}}\\
\textbf{\indent Returns:\ }
\texttt{true} or \texttt{false}.



 A boolean matrix is \textsc{reflexive} if every entry in the main diagonal is \texttt{true}, i.e. \texttt{\mbox{\texttt{\mdseries\slshape mat}}[i][i] = true} for all \texttt{i} in the range \texttt{[1 .. n]} where \texttt{n} is the dimension of \mbox{\texttt{\mdseries\slshape mat}}. 
\begin{Verbatim}[commandchars=!@|,fontsize=\small,frame=single,label=Example]
  !gapprompt@gap>| !gapinput@mat := BooleanMat([[1, 0, 0, 0],|
  !gapprompt@>| !gapinput@                      [1, 1, 0, 0],|
  !gapprompt@>| !gapinput@                      [0, 1, 0, 1],|
  !gapprompt@>| !gapinput@                      [1, 1, 1, 1]]);|
  Matrix(IsBooleanMat, [[1, 0, 0, 0], [1, 1, 0, 0], [0, 1, 0, 1],
    [1, 1, 1, 1]])
  !gapprompt@gap>| !gapinput@IsReflexiveBooleanMat(mat);|
  false
  !gapprompt@gap>| !gapinput@mat := BooleanMat([[1, 1, 1, 0],|
  !gapprompt@>| !gapinput@                      [1, 1, 1, 1],|
  !gapprompt@>| !gapinput@                      [1, 1, 1, 1],|
  !gapprompt@>| !gapinput@                      [0, 1, 1, 1]]);|
  Matrix(IsBooleanMat, [[1, 1, 1, 0], [1, 1, 1, 1], [1, 1, 1, 1],
    [0, 1, 1, 1]])
  !gapprompt@gap>| !gapinput@IsReflexiveBooleanMat(mat);|
  true
\end{Verbatim}
 }

 

\subsection{\textcolor{Chapter }{IsTransitiveBooleanMat}}
\logpage{[ 5, 3, 12 ]}\nobreak
\hyperdef{L}{X7CDAD39B856AC3E5}{}
{\noindent\textcolor{FuncColor}{$\triangleright$\enspace\texttt{IsTransitiveBooleanMat({\mdseries\slshape mat})\index{IsTransitiveBooleanMat@\texttt{IsTransitiveBooleanMat}}
\label{IsTransitiveBooleanMat}
}\hfill{\scriptsize (property)}}\\
\textbf{\indent Returns:\ }
\texttt{true} or \texttt{false}.



 A boolean matrix is \textsc{transitive} if whenever \texttt{\mbox{\texttt{\mdseries\slshape mat}}[i][j] = true} and \texttt{\mbox{\texttt{\mdseries\slshape mat}}[j][k] = true} then \texttt{\mbox{\texttt{\mdseries\slshape mat}}[i][k] = true}. 
\begin{Verbatim}[commandchars=!@|,fontsize=\small,frame=single,label=Example]
  !gapprompt@gap>| !gapinput@x := BooleanMat([[1, 0, 0, 1],|
  !gapprompt@>| !gapinput@                    [1, 0, 1, 1],|
  !gapprompt@>| !gapinput@                    [1, 1, 1, 0],|
  !gapprompt@>| !gapinput@                    [0, 1, 1, 0]]);|
  Matrix(IsBooleanMat, [[1, 0, 0, 1], [1, 0, 1, 1], [1, 1, 1, 0],
    [0, 1, 1, 0]])
  !gapprompt@gap>| !gapinput@IsTransitiveBooleanMat(x);|
  false
  !gapprompt@gap>| !gapinput@x := BooleanMat([[1, 1, 1, 1],|
  !gapprompt@>| !gapinput@                    [1, 1, 1, 1],|
  !gapprompt@>| !gapinput@                    [1, 1, 1, 1],|
  !gapprompt@>| !gapinput@                    [1, 1, 1, 1]]);|
  Matrix(IsBooleanMat, [[1, 1, 1, 1], [1, 1, 1, 1], [1, 1, 1, 1],
    [1, 1, 1, 1]])
  !gapprompt@gap>| !gapinput@IsTransitiveBooleanMat(x);|
  true
\end{Verbatim}
 }

 

\subsection{\textcolor{Chapter }{IsAntiSymmetricBooleanMat}}
\logpage{[ 5, 3, 13 ]}\nobreak
\hyperdef{L}{X8570C8A08549383D}{}
{\noindent\textcolor{FuncColor}{$\triangleright$\enspace\texttt{IsAntiSymmetricBooleanMat({\mdseries\slshape mat})\index{IsAntiSymmetricBooleanMat@\texttt{IsAntiSymmetricBooleanMat}}
\label{IsAntiSymmetricBooleanMat}
}\hfill{\scriptsize (property)}}\\
\textbf{\indent Returns:\ }
\texttt{true} or \texttt{false}.



 A boolean matrix is \textsc{anti\texttt{\symbol{45}}symmetric} if whenever \texttt{\mbox{\texttt{\mdseries\slshape mat}}[i][j] = true} and \texttt{\mbox{\texttt{\mdseries\slshape mat}}[j][i] = true} then \texttt{i = j}. 
\begin{Verbatim}[commandchars=!@|,fontsize=\small,frame=single,label=Example]
  !gapprompt@gap>| !gapinput@x := BooleanMat([[1, 0, 0, 1],|
  !gapprompt@>| !gapinput@                    [1, 0, 1, 1],|
  !gapprompt@>| !gapinput@                    [1, 1, 1, 0],|
  !gapprompt@>| !gapinput@                    [0, 1, 1, 0]]);|
  Matrix(IsBooleanMat, [[1, 0, 0, 1], [1, 0, 1, 1], [1, 1, 1, 0],
    [0, 1, 1, 0]])
  !gapprompt@gap>| !gapinput@IsAntiSymmetricBooleanMat(x);|
  false
  !gapprompt@gap>| !gapinput@x := BooleanMat([[1, 0, 0, 1],|
  !gapprompt@>| !gapinput@                    [1, 0, 1, 0],|
  !gapprompt@>| !gapinput@                    [1, 0, 1, 0],|
  !gapprompt@>| !gapinput@                    [0, 1, 1, 0]]);|
  Matrix(IsBooleanMat, [[1, 0, 0, 1], [1, 0, 1, 0], [1, 0, 1, 0],
    [0, 1, 1, 0]])
  !gapprompt@gap>| !gapinput@IsAntiSymmetricBooleanMat(x);|
  true
\end{Verbatim}
 }

 

\subsection{\textcolor{Chapter }{IsTotalBooleanMat}}
\logpage{[ 5, 3, 14 ]}\nobreak
\hyperdef{L}{X7A68D87982A07C6F}{}
{\noindent\textcolor{FuncColor}{$\triangleright$\enspace\texttt{IsTotalBooleanMat({\mdseries\slshape mat})\index{IsTotalBooleanMat@\texttt{IsTotalBooleanMat}}
\label{IsTotalBooleanMat}
}\hfill{\scriptsize (property)}}\\
\noindent\textcolor{FuncColor}{$\triangleright$\enspace\texttt{IsOntoBooleanMat({\mdseries\slshape mat})\index{IsOntoBooleanMat@\texttt{IsOntoBooleanMat}}
\label{IsOntoBooleanMat}
}\hfill{\scriptsize (property)}}\\
\textbf{\indent Returns:\ }
\texttt{true} or \texttt{false}.



 A boolean matrix is \textsc{total} if there is at least one entry in every row is \texttt{true}. Similarly, a boolean matrix is \textsc{onto} if at least one entry in every column is \texttt{true}. 
\begin{Verbatim}[commandchars=!@|,fontsize=\small,frame=single,label=Example]
  !gapprompt@gap>| !gapinput@x := BooleanMat([[1, 0, 0, 1],|
  !gapprompt@>| !gapinput@                    [1, 0, 1, 1],|
  !gapprompt@>| !gapinput@                    [1, 1, 1, 0],|
  !gapprompt@>| !gapinput@                    [0, 1, 1, 0]]);|
  Matrix(IsBooleanMat, [[1, 0, 0, 1], [1, 0, 1, 1], [1, 1, 1, 0],
    [0, 1, 1, 0]])
  !gapprompt@gap>| !gapinput@IsTotalBooleanMat(x);|
  true
  !gapprompt@gap>| !gapinput@IsOntoBooleanMat(x);|
  true
  !gapprompt@gap>| !gapinput@x := BooleanMat([[1, 0, 0, 1],|
  !gapprompt@>| !gapinput@                    [1, 0, 1, 0],|
  !gapprompt@>| !gapinput@                    [0, 0, 0, 0],|
  !gapprompt@>| !gapinput@                    [0, 1, 1, 0]]);|
  Matrix(IsBooleanMat, [[1, 0, 0, 1], [1, 0, 1, 0], [0, 0, 0, 0],
    [0, 1, 1, 0]])
  !gapprompt@gap>| !gapinput@IsTotalBooleanMat(x);|
  false
  !gapprompt@gap>| !gapinput@IsOntoBooleanMat(x);|
  true
\end{Verbatim}
 }

 

\subsection{\textcolor{Chapter }{IsPartialOrderBooleanMat}}
\logpage{[ 5, 3, 15 ]}\nobreak
\hyperdef{L}{X7D9BECEA7E9B72A7}{}
{\noindent\textcolor{FuncColor}{$\triangleright$\enspace\texttt{IsPartialOrderBooleanMat({\mdseries\slshape mat})\index{IsPartialOrderBooleanMat@\texttt{IsPartialOrderBooleanMat}}
\label{IsPartialOrderBooleanMat}
}\hfill{\scriptsize (property)}}\\
\textbf{\indent Returns:\ }
\texttt{true} or \texttt{false}.



 A boolean matrix is a \textsc{partial order} if it is reflexive, transitive, and anti\texttt{\symbol{45}}symmetric. 
\begin{Verbatim}[commandchars=!@|,fontsize=\small,frame=single,label=Example]
  !gapprompt@gap>| !gapinput@Number(FullBooleanMatMonoid(3), IsPartialOrderBooleanMat);|
  19
\end{Verbatim}
 }

 

\subsection{\textcolor{Chapter }{IsEquivalenceBooleanMat}}
\logpage{[ 5, 3, 16 ]}\nobreak
\hyperdef{L}{X82EA957982B79827}{}
{\noindent\textcolor{FuncColor}{$\triangleright$\enspace\texttt{IsEquivalenceBooleanMat({\mdseries\slshape mat})\index{IsEquivalenceBooleanMat@\texttt{IsEquivalenceBooleanMat}}
\label{IsEquivalenceBooleanMat}
}\hfill{\scriptsize (property)}}\\
\textbf{\indent Returns:\ }
\texttt{true} or \texttt{false}.



 A boolean matrix is an \textsc{equivalence} if it is reflexive, transitive, and symmetric. 
\begin{Verbatim}[commandchars=!@|,fontsize=\small,frame=single,label=Example]
  !gapprompt@gap>| !gapinput@Number(FullBooleanMatMonoid(3), IsEquivalenceBooleanMat);|
  5
  !gapprompt@gap>| !gapinput@Bell(3);|
  5
\end{Verbatim}
 }

 

\subsection{\textcolor{Chapter }{IsTransformationBooleanMat}}
\logpage{[ 5, 3, 17 ]}\nobreak
\hyperdef{L}{X7E6B588887D34A0A}{}
{\noindent\textcolor{FuncColor}{$\triangleright$\enspace\texttt{IsTransformationBooleanMat({\mdseries\slshape mat})\index{IsTransformationBooleanMat@\texttt{IsTransformationBooleanMat}}
\label{IsTransformationBooleanMat}
}\hfill{\scriptsize (property)}}\\
\textbf{\indent Returns:\ }
\texttt{true} or \texttt{false}.



 A boolean matrix is a \textsc{transformation} if every row contains precisely one \texttt{true} value. 
\begin{Verbatim}[commandchars=!@|,fontsize=\small,frame=single,label=Example]
  !gapprompt@gap>| !gapinput@Number(FullBooleanMatMonoid(3), IsTransformationBooleanMat);|
  27
\end{Verbatim}
 }

 }

   
\section{\textcolor{Chapter }{ Matrices over finite fields }}\label{Matrices over finite fields}
\logpage{[ 5, 4, 0 ]}
\hyperdef{L}{X873822B6830CE367}{}
{
  In this section we describe some operations in \textsf{Semigroups} for matrices over finite fields that are required for such matrices to form
semigroups satisfying \texttt{IsActingSemigroup} (\ref{IsActingSemigroup}). 

 From v5.0.0, \textsf{Semigroups} uses the \textsf{GAP} library implementation of matrices over finite fields belonging to the
category \texttt{IsMatrixObj} (\textbf{Reference: IsMatrixObj}) rather than the previous implementation in the \textsf{Semigroups} package. This means that from v5.0.0, matrices over a finite field no longer
belong to the category \texttt{IsMatrixOverSemiring} (\ref{IsMatrixOverSemiring}). 

 The following methods are implemented in \textsf{Semigroups} for matrix objects over finite fields. 

\subsection{\textcolor{Chapter }{RowSpaceBasis (for a matrix over finite field)}}
\logpage{[ 5, 4, 1 ]}\nobreak
\hyperdef{L}{X857E626783CCF766}{}
{\noindent\textcolor{FuncColor}{$\triangleright$\enspace\texttt{RowSpaceBasis({\mdseries\slshape m})\index{RowSpaceBasis@\texttt{RowSpaceBasis}!for a matrix over finite field}
\label{RowSpaceBasis:for a matrix over finite field}
}\hfill{\scriptsize (attribute)}}\\
\noindent\textcolor{FuncColor}{$\triangleright$\enspace\texttt{RowSpaceTransformation({\mdseries\slshape m})\index{RowSpaceTransformation@\texttt{RowSpaceTransformation}!for a matrix over finite field}
\label{RowSpaceTransformation:for a matrix over finite field}
}\hfill{\scriptsize (attribute)}}\\
\noindent\textcolor{FuncColor}{$\triangleright$\enspace\texttt{RowSpaceTransformationInv({\mdseries\slshape m})\index{RowSpaceTransformationInv@\texttt{RowSpaceTransformationInv}!for a matrix over finite field}
\label{RowSpaceTransformationInv:for a matrix over finite field}
}\hfill{\scriptsize (attribute)}}\\


 If \mbox{\texttt{\mdseries\slshape m}} is a matrix object over a finite field, then to compute the value of any of
the above attributes, a canonical basis for the row space of \mbox{\texttt{\mdseries\slshape m}} is computed along with an invertible matrix \texttt{RowSpaceTransformation} such that \texttt{m * RowSpaceTransformation(m) = RowSpaceBasis(m)}. \texttt{RowSpaceTransformationInv(m)} is the inverse of \texttt{RowSpaceTransformation(m)}. 
\begin{Verbatim}[commandchars=!@|,fontsize=\small,frame=single,label=Example]
  !gapprompt@gap>| !gapinput@x := Matrix(GF(4), Z(4) ^ 0 * [[1, 1, 0], [0, 1, 1], [1, 1, 1]]);|
  [ [ Z(2)^0, Z(2)^0, 0*Z(2) ], [ 0*Z(2), Z(2)^0, Z(2)^0 ],
    [ Z(2)^0, Z(2)^0, Z(2)^0 ] ]
  !gapprompt@gap>| !gapinput@RowSpaceBasis(x);|
  <rowbasis of rank 3 over GF(2^2)>
  !gapprompt@gap>| !gapinput@RowSpaceTransformation(x);|
  [ [ 0*Z(2), Z(2)^0, Z(2)^0 ], [ Z(2)^0, Z(2)^0, Z(2)^0 ],
    [ Z(2)^0, 0*Z(2), Z(2)^0 ] ]
\end{Verbatim}
 }

 

\subsection{\textcolor{Chapter }{RightInverse (for a matrix over finite field)}}
\logpage{[ 5, 4, 2 ]}\nobreak
\hyperdef{L}{X8733B04781B682E5}{}
{\noindent\textcolor{FuncColor}{$\triangleright$\enspace\texttt{RightInverse({\mdseries\slshape m})\index{RightInverse@\texttt{RightInverse}!for a matrix over finite field}
\label{RightInverse:for a matrix over finite field}
}\hfill{\scriptsize (attribute)}}\\
\noindent\textcolor{FuncColor}{$\triangleright$\enspace\texttt{LeftInverse({\mdseries\slshape m})\index{LeftInverse@\texttt{LeftInverse}!for a matrix over finite field}
\label{LeftInverse:for a matrix over finite field}
}\hfill{\scriptsize (attribute)}}\\
\textbf{\indent Returns:\ }
A matrix over a finite field.



 These attributes contain a semigroup left\texttt{\symbol{45}}inverse, and a
semigroup right\texttt{\symbol{45}}inverse of the matrix \mbox{\texttt{\mdseries\slshape m}} respectively. 
\begin{Verbatim}[commandchars=!@|,fontsize=\small,frame=single,label=Example]
  !gapprompt@gap>| !gapinput@x := Matrix(GF(4), Z(4) ^ 0 * [[1, 1, 0], [0, 0, 0], [1, 1, 1]]);|
  [ [ Z(2)^0, Z(2)^0, 0*Z(2) ], [ 0*Z(2), 0*Z(2), 0*Z(2) ],
    [ Z(2)^0, Z(2)^0, Z(2)^0 ] ]
  !gapprompt@gap>| !gapinput@LeftInverse(x);|
  [ [ Z(2)^0, Z(2)^0, 0*Z(2) ], [ 0*Z(2), 0*Z(2), 0*Z(2) ],
    [ Z(2)^0, 0*Z(2), Z(2)^0 ] ]
  !gapprompt@gap>| !gapinput@Display(LeftInverse(x) * x);|
   1 1 .
   . . .
   . . 1
\end{Verbatim}
 }

 }

   
\section{\textcolor{Chapter }{ Matrices over the integers }}\label{Matrices over the integers}
\logpage{[ 5, 5, 0 ]}
\hyperdef{L}{X8770A88E82AA24B7}{}
{
  In this section we describe operations in \textsf{Semigroups} specifically for integer matrices. 

 From v5.0.0, \textsf{Semigroups} uses the \textsf{GAP} library implementation of matrices over the integers belonging to the category \texttt{IsMatrixObj} (\textbf{Reference: IsMatrixObj}) rather than the previous implementation in the \textsf{Semigroups} package. This means that from v5.0.0, matrices over the integers no longer
belong to the category \texttt{IsMatrixOverSemiring} (\ref{IsMatrixOverSemiring}).

 The following methods are implemented in \textsf{Semigroups} for matrix objects over the integers. 

\subsection{\textcolor{Chapter }{InverseOp (for an integer matrix)}}
\logpage{[ 5, 5, 1 ]}\nobreak
\hyperdef{L}{X7BC66ECE8378068E}{}
{\noindent\textcolor{FuncColor}{$\triangleright$\enspace\texttt{InverseOp({\mdseries\slshape mat})\index{InverseOp@\texttt{InverseOp}!for an integer matrix}
\label{InverseOp:for an integer matrix}
}\hfill{\scriptsize (operation)}}\\
\textbf{\indent Returns:\ }
An integer matrix.



 If \mbox{\texttt{\mdseries\slshape mat}} is an integer matrix (i.e. belongs to the category \texttt{Integers} (\ref{Integers})) whose inverse (if it exists) is also an integer matrix, then \texttt{InverseOp} returns the inverse of \mbox{\texttt{\mdseries\slshape mat}}. 

 An integer matrix has an integer matrix inverse if and only if it has
determinant one. 
\begin{Verbatim}[commandchars=!@|,fontsize=\small,frame=single,label=Example]
  !gapprompt@gap>| !gapinput@mat := Matrix(Integers, [[0, 0, -1],|
  !gapprompt@>| !gapinput@                            [0, 1, 0],|
  !gapprompt@>| !gapinput@                            [1, 0, 0]]);|
  <3x3-matrix over Integers>
  !gapprompt@gap>| !gapinput@InverseOp(mat);|
  <3x3-matrix over Integers>
  !gapprompt@gap>| !gapinput@mat * InverseOp(mat) = One(mat);|
  true
\end{Verbatim}
 }

 

\subsection{\textcolor{Chapter }{IsTorsion (for an integer matrix)}}
\logpage{[ 5, 5, 2 ]}\nobreak
\hyperdef{L}{X7CA636F080777C36}{}
{\noindent\textcolor{FuncColor}{$\triangleright$\enspace\texttt{IsTorsion({\mdseries\slshape mat})\index{IsTorsion@\texttt{IsTorsion}!for an integer matrix}
\label{IsTorsion:for an integer matrix}
}\hfill{\scriptsize (attribute)}}\\
\textbf{\indent Returns:\ }
\texttt{true} or \texttt{false}



 If \mbox{\texttt{\mdseries\slshape mat}} is an integer matrix (i.e. belongs to the category \texttt{Integers} (\ref{Integers})), then \texttt{IsTorsion} returns \texttt{true} if \mbox{\texttt{\mdseries\slshape mat}} is torsion and \texttt{false} otherwise. 

 An integer matrix \mbox{\texttt{\mdseries\slshape mat}} is torsion if and only if there exists an integer \texttt{n} such that \mbox{\texttt{\mdseries\slshape mat}} to the power of \texttt{n} is equal to the identity matrix. 
\begin{Verbatim}[commandchars=!@|,fontsize=\small,frame=single,label=Example]
  !gapprompt@gap>| !gapinput@mat := Matrix(Integers, [[0, 0, -1],|
  !gapprompt@>| !gapinput@                            [0, 1, 0],|
  !gapprompt@>| !gapinput@                            [1, 0, 0]]);|
  <3x3-matrix over Integers>
  !gapprompt@gap>| !gapinput@IsTorsion(mat);|
  true
  !gapprompt@gap>| !gapinput@mat := Matrix(Integers, [[0, 0, -1, 0],|
  !gapprompt@>| !gapinput@                            [0, -1, 0, 0],|
  !gapprompt@>| !gapinput@                            [4, 4, 2, -1],|
  !gapprompt@>| !gapinput@                            [1, 1, 0, 3]]);|
  <4x4-matrix over Integers>
  !gapprompt@gap>| !gapinput@IsTorsion(mat);|
  false
\end{Verbatim}
 }

 

\subsection{\textcolor{Chapter }{Order}}
\logpage{[ 5, 5, 3 ]}\nobreak
\hyperdef{L}{X84F59A2687C62763}{}
{\noindent\textcolor{FuncColor}{$\triangleright$\enspace\texttt{Order({\mdseries\slshape mat})\index{Order@\texttt{Order}}
\label{Order}
}\hfill{\scriptsize (attribute)}}\\
\textbf{\indent Returns:\ }
An integer or \texttt{infinity}.



 If \mbox{\texttt{\mdseries\slshape mat}} is an integer matrix, then \texttt{InverseOp} returns the order of \mbox{\texttt{\mdseries\slshape mat}}. The order of \mbox{\texttt{\mdseries\slshape mat}} is the smallest integer power of \mbox{\texttt{\mdseries\slshape mat}} equal to the identity. If no such integer exists, the order is equal to \texttt{infinity}. 
\begin{Verbatim}[commandchars=!@|,fontsize=\small,frame=single,label=Example]
  !gapprompt@gap>| !gapinput@mat := Matrix(Integers, [[0, 0, -1, 0],|
  !gapprompt@>| !gapinput@                            [0, -1, 0, 0],|
  !gapprompt@>| !gapinput@                            [4, 4, 2, -1],|
  !gapprompt@>| !gapinput@                            [1, 1, 0, 3]]);;|
  !gapprompt@gap>| !gapinput@Order(mat);|
  infinity
  !gapprompt@gap>| !gapinput@mat := Matrix(Integers, [[0, 0, -1],|
  !gapprompt@>| !gapinput@                            [0, 1, 0],|
  !gapprompt@>| !gapinput@                            [1, 0, 0]]);;|
  !gapprompt@gap>| !gapinput@Order(mat);|
  4
\end{Verbatim}
 }

 }

   
\section{\textcolor{Chapter }{ Max\texttt{\symbol{45}}plus and min\texttt{\symbol{45}}plus matrices }}\label{Max-plus and min-plus matrices}
\logpage{[ 5, 6, 0 ]}
\hyperdef{L}{X86BFFFBC87F2AB1E}{}
{
  In this section we describe operations in \textsf{Semigroups} specifically for max\texttt{\symbol{45}}plus and min\texttt{\symbol{45}}plus
matrices. These are in addition to those given elsewhere in this chapter for
arbitrary matrices over semirings. These include: 
\begin{itemize}
\item  \texttt{InverseOp} (\ref{InverseOp}) 
\item  \texttt{RadialEigenvector} (\ref{RadialEigenvector}) 
\item  \texttt{SpectralRadius} (\ref{SpectralRadius}) 
\item  \texttt{UnweightedPrecedenceDigraph} (\ref{UnweightedPrecedenceDigraph}) 
\end{itemize}
 

\subsection{\textcolor{Chapter }{InverseOp}}
\logpage{[ 5, 6, 1 ]}\nobreak
\hyperdef{L}{X82EC4F49877D6EB1}{}
{\noindent\textcolor{FuncColor}{$\triangleright$\enspace\texttt{InverseOp({\mdseries\slshape mat})\index{InverseOp@\texttt{InverseOp}}
\label{InverseOp}
}\hfill{\scriptsize (operation)}}\\
\textbf{\indent Returns:\ }
A max\texttt{\symbol{45}}plus matrix.



 If \mbox{\texttt{\mdseries\slshape mat}} is an invertible max\texttt{\symbol{45}}plus matrix (i.e. belongs to the
category \texttt{IsMaxPlusMatrix} (\ref{IsMaxPlusMatrix}) and is a row permutation applied to the identity), then \texttt{InverseOp} returns the inverse of \mbox{\texttt{\mdseries\slshape mat}}. This method is described in \cite{Kacie2009aa}. 
\begin{Verbatim}[commandchars=!@|,fontsize=\small,frame=single,label=Example]
  !gapprompt@gap>| !gapinput@InverseOp(Matrix(IsMaxPlusMatrix, [[-infinity, -infinity, 0],|
  !gapprompt@>| !gapinput@                                      [0, -infinity, -infinity],|
  !gapprompt@>| !gapinput@                                      [-infinity, 0, -infinity]]));|
  Matrix(IsMaxPlusMatrix, [[-infinity, 0, -infinity],
    [-infinity, -infinity, 0], [0, -infinity, -infinity]])
\end{Verbatim}
 }

 

\subsection{\textcolor{Chapter }{RadialEigenvector}}
\logpage{[ 5, 6, 2 ]}\nobreak
\hyperdef{L}{X83663A5387042B69}{}
{\noindent\textcolor{FuncColor}{$\triangleright$\enspace\texttt{RadialEigenvector({\mdseries\slshape mat})\index{RadialEigenvector@\texttt{RadialEigenvector}}
\label{RadialEigenvector}
}\hfill{\scriptsize (operation)}}\\
\textbf{\indent Returns:\ }
A vector.



 If \mbox{\texttt{\mdseries\slshape mat}} is a \texttt{n} by \texttt{n} max\texttt{\symbol{45}}plus matrix (i.e. belongs to the category \texttt{IsMaxPlusMatrix} (\ref{IsMaxPlusMatrix})), then \texttt{RadialEigenvector} returns an eigenvector corresponding to the eigenvalue equal to the spectral
radius of \mbox{\texttt{\mdseries\slshape mat}}. This method is described in \cite{Kacie2009aa}. 
\begin{Verbatim}[commandchars=!@|,fontsize=\small,frame=single,label=Example]
  !gapprompt@gap>| !gapinput@RadialEigenvector(Matrix(IsMaxPlusMatrix, [[0, -3], [-2, -10]]));|
  [ 0, -2 ]
\end{Verbatim}
 }

 

\subsection{\textcolor{Chapter }{SpectralRadius}}
\logpage{[ 5, 6, 3 ]}\nobreak
\hyperdef{L}{X83FCFB368743E4BA}{}
{\noindent\textcolor{FuncColor}{$\triangleright$\enspace\texttt{SpectralRadius({\mdseries\slshape mat})\index{SpectralRadius@\texttt{SpectralRadius}}
\label{SpectralRadius}
}\hfill{\scriptsize (operation)}}\\
\textbf{\indent Returns:\ }
A rational number.



 If \mbox{\texttt{\mdseries\slshape mat}} is a max\texttt{\symbol{45}}plus matrix (i.e. belongs to the category \texttt{IsMaxPlusMatrix} (\ref{IsMaxPlusMatrix})), then \texttt{SpectralRadius} returns the supremum of the set of eigenvalues of \mbox{\texttt{\mdseries\slshape mat}}. This method is described in \cite{Bacelli1992aa}. 
\begin{Verbatim}[commandchars=!@|,fontsize=\small,frame=single,label=Example]
  !gapprompt@gap>| !gapinput@SpectralRadius(Matrix(IsMaxPlusMatrix, [[-infinity, 1, -4],|
  !gapprompt@>| !gapinput@                                           [2, -infinity, 0],|
  !gapprompt@>| !gapinput@                                           [0, 1, 1]]));|
  3/2
\end{Verbatim}
 }

 

\subsection{\textcolor{Chapter }{UnweightedPrecedenceDigraph}}
\logpage{[ 5, 6, 4 ]}\nobreak
\hyperdef{L}{X869F60527C2B9328}{}
{\noindent\textcolor{FuncColor}{$\triangleright$\enspace\texttt{UnweightedPrecedenceDigraph({\mdseries\slshape mat})\index{UnweightedPrecedenceDigraph@\texttt{UnweightedPrecedenceDigraph}}
\label{UnweightedPrecedenceDigraph}
}\hfill{\scriptsize (operation)}}\\
\textbf{\indent Returns:\ }
A digraph.



 If \mbox{\texttt{\mdseries\slshape mat}} is a max\texttt{\symbol{45}}plus matrix (i.e. belongs to the category \texttt{IsMaxPlusMatrix} (\ref{IsMaxPlusMatrix})), then \texttt{UnweightedPrecedenceDigraph} returns the unweighted precedence digraph corresponding to \mbox{\texttt{\mdseries\slshape mat}}. 

 For an \texttt{n} by \texttt{n} matrix \mbox{\texttt{\mdseries\slshape mat}}, the unweighted precedence digraph is defined as the digraph with \texttt{n} vertices where vertex \texttt{i} is adjacent to vertex \texttt{j} if and only if \mbox{\texttt{\mdseries\slshape mat}}\texttt{[i][j]} is not equal to \texttt{\texttt{\symbol{45}}infinity}. 
\begin{Verbatim}[commandchars=!@|,fontsize=\small,frame=single,label=Example]
  !gapprompt@gap>| !gapinput@UnweightedPrecedenceDigraph(Matrix(IsMaxPlusMatrix, [[2, -2, 0],|
  !gapprompt@>| !gapinput@[-infinity, 10, -2], [-infinity, 2, 1]]));|
  <immutable digraph with 3 vertices, 7 edges>
\end{Verbatim}
 }

 }

   
\section{\textcolor{Chapter }{ Matrix semigroups }}\label{Matrix semigroups}
\logpage{[ 5, 7, 0 ]}
\hyperdef{L}{X79B614AA803BD103}{}
{
  In this section we describe operations and attributes in \textsf{Semigroups} specifically for matrix semigroups. These are in addition to those given
elsewhere in this chapter for arbitrary matrices over semirings. These
include: 
\begin{itemize}
\item  \texttt{IsXMatrixSemigroup} (\ref{IsXMatrixSemigroup}) 
\item  \texttt{IsFinite} (\ref{IsFinite}) 
\item  \texttt{IsTorsion} (\ref{IsTorsion}) 
\item  \texttt{NormalizeSemigroup} (\ref{NormalizeSemigroup}) 
\end{itemize}
 Random matrix semigroups can be created by using the function \texttt{RandomSemigroup} (\ref{RandomSemigroup}). We can also create a matrix semigroup isomorphic to a given semigroup by
using \texttt{IsomorphismSemigroup} (\ref{IsomorphismSemigroup}) and \texttt{AsSemigroup} (\ref{AsSemigroup}). These functions require a filter, and accept any of the filters in Section \ref{IsXMatrixSemigroup}. 

 There are corresponding functions which can be used for matrix monoids: \texttt{RandomMonoid} (\ref{RandomMonoid}), \texttt{IsomorphismMonoid} (\ref{IsomorphismMonoid}), and \texttt{AsMonoid} (\ref{AsMonoid}). These can be used with the filters in Section \ref{IsXMatrixMonoid}. 

 
\subsection{\textcolor{Chapter }{Matrix semigroup filters}}\label{IsXMatrixSemigroup}
\logpage{[ 5, 7, 1 ]}
\hyperdef{L}{X7DC6EB0680B3E4DD}{}
{
\noindent\textcolor{FuncColor}{$\triangleright$\enspace\texttt{IsMatrixOverSemiringSemigroup({\mdseries\slshape obj})\index{IsMatrixOverSemiringSemigroup@\texttt{IsMatrixOverSemiringSemigroup}}
\label{IsMatrixOverSemiringSemigroup}
}\hfill{\scriptsize (Category)}}\\
\noindent\textcolor{FuncColor}{$\triangleright$\enspace\texttt{IsBooleanMatSemigroup({\mdseries\slshape obj})\index{IsBooleanMatSemigroup@\texttt{IsBooleanMatSemigroup}}
\label{IsBooleanMatSemigroup}
}\hfill{\scriptsize (Category)}}\\
\noindent\textcolor{FuncColor}{$\triangleright$\enspace\texttt{IsMatrixOverFiniteFieldSemigroup({\mdseries\slshape obj})\index{IsMatrixOverFiniteFieldSemigroup@\texttt{IsMatrixOverFiniteFieldSemigroup}}
\label{IsMatrixOverFiniteFieldSemigroup}
}\hfill{\scriptsize (Category)}}\\
\noindent\textcolor{FuncColor}{$\triangleright$\enspace\texttt{IsMaxPlusMatrixSemigroup({\mdseries\slshape obj})\index{IsMaxPlusMatrixSemigroup@\texttt{IsMaxPlusMatrixSemigroup}}
\label{IsMaxPlusMatrixSemigroup}
}\hfill{\scriptsize (Category)}}\\
\noindent\textcolor{FuncColor}{$\triangleright$\enspace\texttt{IsMinPlusMatrixSemigroup({\mdseries\slshape obj})\index{IsMinPlusMatrixSemigroup@\texttt{IsMinPlusMatrixSemigroup}}
\label{IsMinPlusMatrixSemigroup}
}\hfill{\scriptsize (Category)}}\\
\noindent\textcolor{FuncColor}{$\triangleright$\enspace\texttt{IsTropicalMatrixSemigroup({\mdseries\slshape obj})\index{IsTropicalMatrixSemigroup@\texttt{IsTropicalMatrixSemigroup}}
\label{IsTropicalMatrixSemigroup}
}\hfill{\scriptsize (Category)}}\\
\noindent\textcolor{FuncColor}{$\triangleright$\enspace\texttt{IsTropicalMaxPlusMatrixSemigroup({\mdseries\slshape obj})\index{IsTropicalMaxPlusMatrixSemigroup@\texttt{IsTropicalMaxPlusMatrixSemigroup}}
\label{IsTropicalMaxPlusMatrixSemigroup}
}\hfill{\scriptsize (Category)}}\\
\noindent\textcolor{FuncColor}{$\triangleright$\enspace\texttt{IsTropicalMinPlusMatrixSemigroup({\mdseries\slshape obj})\index{IsTropicalMinPlusMatrixSemigroup@\texttt{IsTropicalMinPlusMatrixSemigroup}}
\label{IsTropicalMinPlusMatrixSemigroup}
}\hfill{\scriptsize (Category)}}\\
\noindent\textcolor{FuncColor}{$\triangleright$\enspace\texttt{IsNTPMatrixSemigroup({\mdseries\slshape obj})\index{IsNTPMatrixSemigroup@\texttt{IsNTPMatrixSemigroup}}
\label{IsNTPMatrixSemigroup}
}\hfill{\scriptsize (Category)}}\\
\noindent\textcolor{FuncColor}{$\triangleright$\enspace\texttt{IsIntegerMatrixSemigroup({\mdseries\slshape obj})\index{IsIntegerMatrixSemigroup@\texttt{IsIntegerMatrixSemigroup}}
\label{IsIntegerMatrixSemigroup}
}\hfill{\scriptsize (Category)}}\\
\textbf{\indent Returns:\ }
\texttt{true} or \texttt{false}.



 The above are the currently supported types of matrix semigroups. For monoids
see Section \ref{IsXMatrixMonoid}. }

 
\subsection{\textcolor{Chapter }{Matrix monoid filters}}\label{IsXMatrixMonoid}
\logpage{[ 5, 7, 2 ]}
\hyperdef{L}{X8616225581BC7414}{}
{
\noindent\textcolor{FuncColor}{$\triangleright$\enspace\texttt{IsMatrixOverSemiringMonoid({\mdseries\slshape obj})\index{IsMatrixOverSemiringMonoid@\texttt{IsMatrixOverSemiringMonoid}}
\label{IsMatrixOverSemiringMonoid}
}\hfill{\scriptsize (Category)}}\\
\noindent\textcolor{FuncColor}{$\triangleright$\enspace\texttt{IsBooleanMatMonoid({\mdseries\slshape obj})\index{IsBooleanMatMonoid@\texttt{IsBooleanMatMonoid}}
\label{IsBooleanMatMonoid}
}\hfill{\scriptsize (Category)}}\\
\noindent\textcolor{FuncColor}{$\triangleright$\enspace\texttt{IsMatrixOverFiniteFieldMonoid({\mdseries\slshape obj})\index{IsMatrixOverFiniteFieldMonoid@\texttt{IsMatrixOverFiniteFieldMonoid}}
\label{IsMatrixOverFiniteFieldMonoid}
}\hfill{\scriptsize (Category)}}\\
\noindent\textcolor{FuncColor}{$\triangleright$\enspace\texttt{IsMaxPlusMatrixMonoid({\mdseries\slshape obj})\index{IsMaxPlusMatrixMonoid@\texttt{IsMaxPlusMatrixMonoid}}
\label{IsMaxPlusMatrixMonoid}
}\hfill{\scriptsize (Category)}}\\
\noindent\textcolor{FuncColor}{$\triangleright$\enspace\texttt{IsMinPlusMatrixMonoid({\mdseries\slshape obj})\index{IsMinPlusMatrixMonoid@\texttt{IsMinPlusMatrixMonoid}}
\label{IsMinPlusMatrixMonoid}
}\hfill{\scriptsize (Category)}}\\
\noindent\textcolor{FuncColor}{$\triangleright$\enspace\texttt{IsTropicalMatrixMonoid({\mdseries\slshape obj})\index{IsTropicalMatrixMonoid@\texttt{IsTropicalMatrixMonoid}}
\label{IsTropicalMatrixMonoid}
}\hfill{\scriptsize (Category)}}\\
\noindent\textcolor{FuncColor}{$\triangleright$\enspace\texttt{IsTropicalMaxPlusMatrixMonoid({\mdseries\slshape obj})\index{IsTropicalMaxPlusMatrixMonoid@\texttt{IsTropicalMaxPlusMatrixMonoid}}
\label{IsTropicalMaxPlusMatrixMonoid}
}\hfill{\scriptsize (Category)}}\\
\noindent\textcolor{FuncColor}{$\triangleright$\enspace\texttt{IsTropicalMinPlusMatrixMonoid({\mdseries\slshape obj})\index{IsTropicalMinPlusMatrixMonoid@\texttt{IsTropicalMinPlusMatrixMonoid}}
\label{IsTropicalMinPlusMatrixMonoid}
}\hfill{\scriptsize (Category)}}\\
\noindent\textcolor{FuncColor}{$\triangleright$\enspace\texttt{IsNTPMatrixMonoid({\mdseries\slshape obj})\index{IsNTPMatrixMonoid@\texttt{IsNTPMatrixMonoid}}
\label{IsNTPMatrixMonoid}
}\hfill{\scriptsize (Category)}}\\
\noindent\textcolor{FuncColor}{$\triangleright$\enspace\texttt{IsIntegerMatrixMonoid({\mdseries\slshape obj})\index{IsIntegerMatrixMonoid@\texttt{IsIntegerMatrixMonoid}}
\label{IsIntegerMatrixMonoid}
}\hfill{\scriptsize (Category)}}\\
\textbf{\indent Returns:\ }
\texttt{true} or \texttt{false}.



 The above are the currently supported types of matrix monoids. They correspond
to the matrix semigroup types in Section \ref{IsXMatrixSemigroup}. }

 

\subsection{\textcolor{Chapter }{IsFinite}}
\logpage{[ 5, 7, 3 ]}\nobreak
\hyperdef{L}{X808A4061809A6E67}{}
{\noindent\textcolor{FuncColor}{$\triangleright$\enspace\texttt{IsFinite({\mdseries\slshape S})\index{IsFinite@\texttt{IsFinite}}
\label{IsFinite}
}\hfill{\scriptsize (property)}}\\
\textbf{\indent Returns:\ }
\texttt{true} or \texttt{false}. 



 If \mbox{\texttt{\mdseries\slshape S}} is a max\texttt{\symbol{45}}plus or min\texttt{\symbol{45}}plus matrix
semigroup (i.e. belongs to the category \texttt{IsMaxPlusMatrixSemigroup} (\ref{IsMaxPlusMatrixSemigroup})), then \texttt{IsFinite} returns \texttt{true} if \mbox{\texttt{\mdseries\slshape S}} is finite and \texttt{false} otherwise. This method is based on \cite{Gaubert1996aa} (max\texttt{\symbol{45}}plus) and \cite{Simon1978aa} (min\texttt{\symbol{45}}plus). For min\texttt{\symbol{45}}plus matrix
semigroups, this method is only valid if each min\texttt{\symbol{45}}plus
matrix in the semigroup contains only non\texttt{\symbol{45}}negative
integers. Note, this method is terminating and does not involve enumerating
semigroups. 
\begin{Verbatim}[commandchars=!@|,fontsize=\small,frame=single,label=Example]
  !gapprompt@gap>| !gapinput@IsFinite(Semigroup(Matrix(IsMaxPlusMatrix,|
  !gapprompt@>| !gapinput@                             [[0, -3],|
  !gapprompt@>| !gapinput@                              [-2, -10]])));|
  true
  !gapprompt@gap>| !gapinput@IsFinite(Semigroup(Matrix(IsMaxPlusMatrix,|
  !gapprompt@>| !gapinput@                             [[1, -infinity, 2],|
  !gapprompt@>| !gapinput@                              [-2, 4, -infinity],|
  !gapprompt@>| !gapinput@                              [1, 0, 3]])));|
  false
  !gapprompt@gap>| !gapinput@IsFinite(Semigroup(Matrix(IsMinPlusMatrix,|
  !gapprompt@>| !gapinput@                             [[infinity, 0],|
  !gapprompt@>| !gapinput@                              [5, 4]])));|
  false
  !gapprompt@gap>| !gapinput@IsFinite(Semigroup(Matrix(IsMinPlusMatrix,|
  !gapprompt@>| !gapinput@                             [[1, 0],|
  !gapprompt@>| !gapinput@                              [0, infinity]])));|
  true
\end{Verbatim}
 }

 

\subsection{\textcolor{Chapter }{IsTorsion}}
\logpage{[ 5, 7, 4 ]}\nobreak
\hyperdef{L}{X80C6B26284721409}{}
{\noindent\textcolor{FuncColor}{$\triangleright$\enspace\texttt{IsTorsion({\mdseries\slshape S})\index{IsTorsion@\texttt{IsTorsion}}
\label{IsTorsion}
}\hfill{\scriptsize (attribute)}}\\
\textbf{\indent Returns:\ }
\texttt{true} or \texttt{false}.



 If \mbox{\texttt{\mdseries\slshape S}} is a max\texttt{\symbol{45}}plus matrix semigroup (i.e. belongs to the
category \texttt{IsMaxPlusMatrixSemigroup} (\ref{IsMaxPlusMatrixSemigroup})), then \texttt{IsTorsion} returns \texttt{true} if \mbox{\texttt{\mdseries\slshape S}} is torsion and \texttt{false} otherwise. This method is based on \cite{Gaubert1996aa} and draws on \cite{Burrell2016aa}, \cite{Bacelli1992aa} and \cite{Kacie2009aa}. 
\begin{Verbatim}[commandchars=!@|,fontsize=\small,frame=single,label=Example]
  !gapprompt@gap>| !gapinput@IsTorsion(Semigroup(Matrix(IsMaxPlusMatrix,|
  !gapprompt@>| !gapinput@                             [[0, -3],|
  !gapprompt@>| !gapinput@                              [-2, -10]])));|
  true
  !gapprompt@gap>| !gapinput@IsTorsion(Semigroup(Matrix(IsMaxPlusMatrix,|
  !gapprompt@>| !gapinput@                              [[1, -infinity, 2],|
  !gapprompt@>| !gapinput@                               [-2, 4, -infinity],|
  !gapprompt@>| !gapinput@                               [1, 0, 3]])));|
  false
\end{Verbatim}
 }

 

\subsection{\textcolor{Chapter }{NormalizeSemigroup}}
\logpage{[ 5, 7, 5 ]}\nobreak
\hyperdef{L}{X873DE466868DA849}{}
{\noindent\textcolor{FuncColor}{$\triangleright$\enspace\texttt{NormalizeSemigroup({\mdseries\slshape S})\index{NormalizeSemigroup@\texttt{NormalizeSemigroup}}
\label{NormalizeSemigroup}
}\hfill{\scriptsize (operation)}}\\
\textbf{\indent Returns:\ }
A semigroup.



 This method applies to max\texttt{\symbol{45}}plus matrix semigroups (i.e.
those belonging to the category \texttt{IsMaxPlusMatrixSemigroup} (\ref{IsMaxPlusMatrixSemigroup})) that are finitely generated, such that the spectral radius of the matrix
equal to the sum of the generators (with respect to the
max\texttt{\symbol{45}}plus semiring) is zero. \texttt{NormalizeSemigroup} returns a semigroup of matrices all with strictly
non\texttt{\symbol{45}}positive entries. Note that the output is isomorphic to
a min\texttt{\symbol{45}}plus matrix semigroup. This method is based on \cite{Gaubert1996aa} and \cite{Burrell2016aa}. 
\begin{Verbatim}[commandchars=!@|,fontsize=\small,frame=single,label=Example]
  !gapprompt@gap>| !gapinput@NormalizeSemigroup(Semigroup(Matrix(IsMaxPlusMatrix,|
  !gapprompt@>| !gapinput@                                       [[0, -3],|
  !gapprompt@>| !gapinput@                                        [-2, -10]])));|
  <commutative semigroup of 2x2 max-plus matrices with 1 generator>
\end{Verbatim}
 }

 }

   }

 
\chapter{\textcolor{Chapter }{ Semigroups and monoids defined by generating sets }}\label{Semigroups and monoids defined by generating sets}
\logpage{[ 6, 0, 0 ]}
\hyperdef{L}{X7995B4F18672DDB0}{}
{
  In this chapter we describe the various ways that semigroups and monoids
defined by generating sets can be created in \textsf{Semigroups}; where the generators are, for example, those elements described in earlier
chapters of this manual. 
\section{\textcolor{Chapter }{Underlying algorithms}}\label{Underlying algorithms}
\logpage{[ 6, 1, 0 ]}
\hyperdef{L}{X7A19D22B7A05CC2F}{}
{
  Computing the Green's structure of a semigroup or monoid is a fundamental step
in almost every algorithm for semigroups. There are two fundamental algorithms
in the \textsf{Semigroups} package for computing the Green's structure of a semigroup defined by a set of
generators. In the next two subsections we briefly describe these two
algorithms. 
\subsection{\textcolor{Chapter }{ Acting semigroups }}\label{Acting semigroups}
\logpage{[ 6, 1, 1 ]}
\hyperdef{L}{X7A3AC74C7FF85825}{}
{
  The first of the fundamental algorithms for computing a semigroup defined by a
generating set is described in \cite{Mitchell2019aa}. When applied to a semigroup or monoid with relatively large subgroups, or $\mathcal{D}$\texttt{\symbol{45}}classes, these are the most efficient methods in the \textsf{Semigroups} package. For example, the complexity of computing, say, the size of a
transformation semigroup that happens to be a group, is the same as the
complexity of the Schreier\texttt{\symbol{45}}Sims Algorithm (polynomial in
the number of points acted on by the transformations) for a permutation group. 

 In theory, these algorithms can be applied to compute any subsemigroup of a
regular semigroup; but so far in the \textsf{Semigroups} package they are only implemented for semigroups of: transformations (see  (\textbf{Reference: Transformations})), partial permutations (see  (\textbf{Reference: Partial permutations})), bipartitions (see Chapter \ref{Bipartitions and blocks}), matrices over a finite field (see Section \ref{Matrices over finite fields}); subsemigroups of regular Rees 0\texttt{\symbol{45}}matrix semigroups over
permutation groups (see Chapter  (\textbf{Reference: Rees Matrix Semigroups})), and subsemigroups of McAlister triples (see Section \ref{McAlister
        triple semigroups}). 

 We refer to semigroups to which the algorithms in \cite{Mitchell2019aa} can be applied as \emph{acting semigroups}, and such semigroups belong to the category \texttt{IsActingSemigroup} (\ref{IsActingSemigroup}).

 If you know \emph{a priori} that the semigroup you want to compute is large and $\mathcal{J}$\texttt{\symbol{45}}trivial, then you can disable the special methods for
acting semigroups when you create the semigroups; see Section \ref{Options when creating semigroups} for more details.

 It is harder for the acting semigroup algorithms to compute Green's $\mathcal{L}$\texttt{\symbol{45}} and $\mathcal{H}$\texttt{\symbol{45}}classes of a transformation semigroup and the methods used
to compute with Green's $\mathcal{R}$\texttt{\symbol{45}} and $\mathcal{D}$\texttt{\symbol{45}}classes are the most efficient in \textsf{Semigroups}. Thus, if you are computing with a transformation semigroup, wherever
possible, it is advisable to use the commands relating to Green's $\mathcal{R}$\texttt{\symbol{45}} or $\mathcal{D}$\texttt{\symbol{45}}classes rather than those relating to Green's $\mathcal{L}$\texttt{\symbol{45}} or $\mathcal{H}$\texttt{\symbol{45}}classes. No such difficulties are present when computing
with the other types of acting semigroups in \textsf{Semigroups}.

 There are methods in \textsf{Semigroups} for computing individual Green's classes of an acting semigroup without
computing the entire data structure of the underlying semigroup; see \texttt{GreensRClassOfElementNC} (\ref{GreensRClassOfElementNC}). It is also possible to compute the $\mathcal{R}$\texttt{\symbol{45}}classes, the number of elements and test membership in a
semigroup without computing all the elements; see, for example, \texttt{GreensRClasses} (\ref{GreensRClasses}), \texttt{RClassReps} (\ref{RClassReps}), \texttt{IteratorOfRClasses} (\ref{IteratorOfRClasses}), or \texttt{NrRClasses} (\ref{NrRClasses}). This may be useful if you want to study a very large semigroup where
computing all the elements of the semigroup is not feasible.

 }

 

\subsection{\textcolor{Chapter }{IsActingSemigroup}}
\logpage{[ 6, 1, 2 ]}\nobreak
\hyperdef{L}{X7F69D8FC7D578A0C}{}
{\noindent\textcolor{FuncColor}{$\triangleright$\enspace\texttt{IsActingSemigroup({\mdseries\slshape obj})\index{IsActingSemigroup@\texttt{IsActingSemigroup}}
\label{IsActingSemigroup}
}\hfill{\scriptsize (Category)}}\\
\textbf{\indent Returns:\ }
\texttt{true} or \texttt{false}.



 Every acting semigroup, as defined in Section \ref{Acting semigroups}, belongs to this category. 
\begin{Verbatim}[commandchars=!@|,fontsize=\small,frame=single,label=Example]
  !gapprompt@gap>| !gapinput@S := Semigroup(Transformation([1, 3, 2]));;|
  !gapprompt@gap>| !gapinput@IsActingSemigroup(S);|
  true
  !gapprompt@gap>| !gapinput@CanUseFroidurePin(S);|
  true
  !gapprompt@gap>| !gapinput@S := FreeSemigroup(3);;|
  !gapprompt@gap>| !gapinput@IsActingSemigroup(S);|
  false
\end{Verbatim}
 }

 
\subsection{\textcolor{Chapter }{The Froidure\texttt{\symbol{45}}Pin Algorithm}}\label{FroidurePin}
\logpage{[ 6, 1, 3 ]}
\hyperdef{L}{X7E2DE9767D5D82F7}{}
{
  The second fundamental algorithm for computing finite semigroups is the
Froidure\texttt{\symbol{45}}Pin Algorithm \cite{Froidure1997aa}. The \textsf{Semigroups} package contains two implementations of the Froidure\texttt{\symbol{45}}Pin
Algorithm: one in the \href{https://libsemigroups.readthedocs.io/en/latest/} {libsemigroups}  C++ library and the other within the \textsf{Semigroups} package kernel module.

 Both implementations outperform the algorithms for acting semigroups when
applied to semigroups with small (trivial) subgroups. This method is also used
to determine the structure of a semigroup when the algorithms described in \cite{Mitchell2019aa} do not apply. It is possible to specify which methods should be applied to a
given semigroup; see Section \ref{Options when creating semigroups}.

 A semigroup to which the Froidure\texttt{\symbol{45}}Pin Algorithm can be
applied in \textsf{Semigroups} satisfy \texttt{CanUseFroidurePin} (\ref{CanUseFroidurePin}). Every acting semigroup in \textsf{Semigroups} satisfies \texttt{CanUseFroidurePin} (\ref{CanUseFroidurePin}) and the Froidure\texttt{\symbol{45}}Pin Algorithm is used to compute certain
properties or attributes. 

 Currently, the \href{https://libsemigroups.readthedocs.io/en/latest/} {libsemigroups}  implementation of the Froidure\texttt{\symbol{45}}Pin Algorithm can be applied
to semigroups consisting of the following types of elements: transformations
(see  (\textbf{Reference: Transformations})), partial permutations (see  (\textbf{Reference: Partial permutations})), bipartitions (see Chapter \ref{Bipartitions and blocks}), partitioned binary relations (see Chapter \ref{Partitioned binary relations (PBRs)}) as defined in \cite{Martin2011aa}; and matrices over the following semirings: 
\begin{itemize}
\item  the \emph{Boolean semiring} $\{0, 1\}$ where $0 + 0 = 0$, $0 + 1 = 1 + 1 = 1 + 0 = 1$, and $1\cdot 0 = 0 \cdot 0 = 0 \cdot 1 = 0$ 
\item  finite fields; 
\item  the \emph{max\texttt{\symbol{45}}plus semiring} of natural numbers and negative infinity $\mathbb{N}\cup \{-\infty\}$ with operations max and plus; 
\item  the \emph{min\texttt{\symbol{45}}plus semiring} of natural numbers and infinity $\mathbb{N}\cup \{\infty\}$ with operations min and plus; 
\item  the \emph{tropical max\texttt{\symbol{45}}plus semiring} $\{-\infty, 0, 1, \ldots, t\}$ for some threshold $t$ with operations max and plus; 
\item  the \emph{tropical min\texttt{\symbol{45}}plus semiring} $\{0, 1, \ldots, t, \infty\}$ for some threshold $t$ with operations min and plus; 
\item  the semiring $\mathbb{N}_{t, p} = \{0, 1, \ldots, t, t + 1, \ldots, t + p - 1\}$ for some threshold $t$ and period $p$ under addition and multiplication modulo the congruence $t = t + p$. 
\end{itemize}
 (see Chapter \ref{Matrices over semirings}).

 The version of the Froidure\texttt{\symbol{45}}Pin Algorithm \cite{Froidure1997aa} written in C within the \textsf{Semigroups} package kernel module can be used to compute any other semigroup in \textsf{GAP} which satisfies \texttt{CanUseGapFroidurePin} (\ref{CanUseGapFroidurePin}). In theory, any finite semigroup can be computed using this algorithm.
However, the condition that the semigroup has satisfies \texttt{CanUseGapFroidurePin} (\ref{CanUseGapFroidurePin}) is imposed to avoid this method being used when it is inappropriate. If
implementing a new type of semigroup in \textsf{GAP}, then simply do 
\begin{Verbatim}[commandchars=!@|,fontsize=\small,frame=single,label=Example]
  InstallTrueMethod(CanUseGapFroidurePin,
                         MyNewSemigroupType);
\end{Verbatim}
 to make your new semigroup type \texttt{MyNewSemigroupType} use this version of the Froidure\texttt{\symbol{45}}Pin Algorithm. To make
this work efficiently it is necessary that a hash function is implemented for
the elements of \texttt{MyNewSemigroupType}; more details will be included in a future edition of this manual. 

 Mostly due to the way that \textsf{GAP} handles memory, this implementation is approximately 4 times slower than the
implementation in \href{https://libsemigroups.readthedocs.io/en/latest/} {libsemigroups} . This version of the Froidure\texttt{\symbol{45}}Pin Algorithm is included
because it applies to a wider class of semigroups than those currently
implemented in \href{https://libsemigroups.readthedocs.io/en/latest/} {libsemigroups}  and it is more straightforward to extend the classes of semigroup to which it
applies. }

 

\subsection{\textcolor{Chapter }{CanUseFroidurePin}}
\logpage{[ 6, 1, 4 ]}\nobreak
\hyperdef{L}{X7FEE8CFA87E7B872}{}
{\noindent\textcolor{FuncColor}{$\triangleright$\enspace\texttt{CanUseFroidurePin({\mdseries\slshape obj})\index{CanUseFroidurePin@\texttt{CanUseFroidurePin}}
\label{CanUseFroidurePin}
}\hfill{\scriptsize (property)}}\\
\noindent\textcolor{FuncColor}{$\triangleright$\enspace\texttt{CanUseGapFroidurePin({\mdseries\slshape obj})\index{CanUseGapFroidurePin@\texttt{CanUseGapFroidurePin}}
\label{CanUseGapFroidurePin}
}\hfill{\scriptsize (property)}}\\
\noindent\textcolor{FuncColor}{$\triangleright$\enspace\texttt{CanUseLibsemigroupsFroidurePin({\mdseries\slshape obj})\index{CanUseLibsemigroupsFroidurePin@\texttt{CanUseLibsemigroupsFroidurePin}}
\label{CanUseLibsemigroupsFroidurePin}
}\hfill{\scriptsize (property)}}\\
\textbf{\indent Returns:\ }
\texttt{true} or \texttt{false}.



 Every semigroup satisfying \texttt{CanUseFroidurePin} is a valid input for the Froidure\texttt{\symbol{45}}Pin algorithm; see
Section \ref{FroidurePin} for more details.

 Basic operations for semigroups satisfying \texttt{CanUseFroidurePin} are: \texttt{AsListCanonical} (\ref{AsListCanonical}), \texttt{EnumeratorCanonical} (\ref{EnumeratorCanonical}), \texttt{IteratorCanonical} (\ref{IteratorCanonical}), \texttt{PositionCanonical} (\ref{PositionCanonical}), \texttt{Enumerate} (\ref{Enumerate}), and \texttt{IsEnumerated} (\ref{IsEnumerated}).  
\begin{Verbatim}[commandchars=!@|,fontsize=\small,frame=single,label=Example]
  !gapprompt@gap>| !gapinput@S := Semigroup(Transformation([1, 3, 2]));;|
  !gapprompt@gap>| !gapinput@CanUseFroidurePin(S);|
  true
  !gapprompt@gap>| !gapinput@S := FreeSemigroup(3);;|
  !gapprompt@gap>| !gapinput@CanUseFroidurePin(S);|
  false
\end{Verbatim}
 }

 }

 
\section{\textcolor{Chapter }{Semigroups represented by generators}}\label{Semigroups represented by generators}
\logpage{[ 6, 2, 0 ]}
\hyperdef{L}{X79BD00A682BDED7A}{}
{
  

\subsection{\textcolor{Chapter }{InverseMonoidByGenerators}}
\logpage{[ 6, 2, 1 ]}\nobreak
\hyperdef{L}{X79A15C7C83BBA60B}{}
{\noindent\textcolor{FuncColor}{$\triangleright$\enspace\texttt{InverseMonoidByGenerators({\mdseries\slshape coll[, opts]})\index{InverseMonoidByGenerators@\texttt{InverseMonoidByGenerators}}
\label{InverseMonoidByGenerators}
}\hfill{\scriptsize (operation)}}\\
\noindent\textcolor{FuncColor}{$\triangleright$\enspace\texttt{InverseSemigroupByGenerators({\mdseries\slshape coll[, opts]})\index{InverseSemigroupByGenerators@\texttt{InverseSemigroupByGenerators}}
\label{InverseSemigroupByGenerators}
}\hfill{\scriptsize (operation)}}\\
\textbf{\indent Returns:\ }
An inverse monoid or semigroup.



 If \mbox{\texttt{\mdseries\slshape coll}} is a collection satisfying \texttt{IsGeneratorsOfInverseSemigroup}, then \texttt{InverseMonoidByGenerators} and \texttt{InverseSemigroupByGenerators} return the inverse monoid and semigroup generated by \mbox{\texttt{\mdseries\slshape coll}}, respectively. 

 If present, the optional second argument \mbox{\texttt{\mdseries\slshape opts}} should be a record containing the values of the options for the semigroup
being created, as described in Section \ref{Options when creating semigroups}.

 }

 }

 
\section{\textcolor{Chapter }{Options when creating semigroups}}\label{Options when creating semigroups}
\logpage{[ 6, 3, 0 ]}
\hyperdef{L}{X799EBA2F819D8867}{}
{
  When using any of the functions: 
\begin{itemize}
\item \texttt{InverseSemigroup} (\textbf{Reference: InverseSemigroup}), 
\item \texttt{InverseMonoid} (\textbf{Reference: InverseMonoid}), 
\item \texttt{Semigroup} (\textbf{Reference: Semigroup}), 
\item \texttt{Monoid} (\textbf{Reference: Monoid}), 
\item \texttt{SemigroupByGenerators} (\textbf{Reference: SemigroupByGenerators}), 
\item \texttt{MonoidByGenerators} (\textbf{Reference: MonoidByGenerators}), 
\item \texttt{ClosureSemigroup} (\ref{ClosureSemigroup}), 
\item \texttt{ClosureMonoid} (\ref{ClosureMonoid}), 
\item \texttt{ClosureInverseSemigroup} (\ref{ClosureInverseSemigroup}), 
\item \texttt{ClosureInverseMonoid} (\ref{ClosureInverseMonoid}), 
\item \texttt{SemigroupIdeal} (\ref{SemigroupIdeal})
\end{itemize}
 a record can be given as an optional final argument. The components of this
record specify the values of certain options for the semigroup being created.
A list of these options and their default values is given below. 

 Assume that \mbox{\texttt{\mdseries\slshape S}} is the semigroup created by one of the functions given above and that either: \mbox{\texttt{\mdseries\slshape S}} is generated by a collection \mbox{\texttt{\mdseries\slshape gens}}; or \mbox{\texttt{\mdseries\slshape S}} is an ideal of such a semigroup. 
\begin{description}
\item[{\texttt{acting}}]  this component should be \texttt{true} or \texttt{false}. Roughly speaking, there are two types of methods in the \textsf{Semigroups} package: those for semigroups which have to be fully enumerated, and those for
semigroups that do not; see Section \ref{Introduction}. In order for a semigroup to use the latter methods in  \textsf{Semigroups} it must satisfy \texttt{IsActingSemigroup} (\ref{IsActingSemigroup}). By default any semigroup or monoid of transformations, partial permutations,
Rees 0\texttt{\symbol{45}}matrix elements, or bipartitions satisfies \texttt{IsActingSemigroup}.

 There are cases (such as when it is known \emph{a priori} that the semigroup is $\mathcal{D}$\texttt{\symbol{45}}trivial), when it might be preferable to use the methods
that involve fully enumerating a semigroup. In other words, it might be
desirable to disable the more sophisticated methods for acting semigroups. If
this is the case, then the value of this component can be set \texttt{false} when the semigroup is created. Following this none of the special methods for
acting semigroup will be used to compute anything about the semigroup. 
\item[{\texttt{regular}}]  this component should be \texttt{true} or \texttt{false}. If it is known \emph{a priori} that the semigroup \texttt{S} being created is a regular semigroup, then this component can be set to \texttt{true}. In this case, \texttt{S} knows it is a regular semigroup and can take advantage of the methods for
regular semigroups in \textsf{Semigroups}. It is usually much more efficient to compute with a regular semigroup that
to compute with a non\texttt{\symbol{45}}regular semigroup.

 If this option is set to \texttt{true} when the semigroup being defined is \textsc{not} regular, then the results might be unpredictable. 

 The default value for this option is \texttt{false}. 
\item[{\texttt{hashlen}}]  this component should be a positive integer, which roughly specifies the
lengths of the hash tables used internally by \textsf{Semigroups}. \textsf{Semigroups} uses hash tables in several fundamental methods. The lengths of these tables
are a compromise between performance and memory usage; larger tables provide
better performance for large computations but use more memory. Note that it is
unlikely that you will need to specify this option unless you find that \textsf{GAP} runs out of memory unexpectedly or that the performance of \textsf{Semigroups} is poorer than expected. If you find that \textsf{GAP} runs out of memory unexpectedly, or you plan to do a large number of
computations with relatively small semigroups (say with tens of thousands of
elements), then you might consider setting \texttt{hashlen} to be less than the default value of \texttt{12517} for each of these semigroups. If you find that the performance of \textsf{Semigroups} is unexpectedly poor, or you plan to do a computation with a very large
semigroup (say, more than 10 million elements), then you might consider
setting \texttt{hashlen} to be greater than the default value of \texttt{12517}. 

 You might find it useful to set the info level of the info class \texttt{InfoOrb} to 2 or higher since this will indicate when hash tables used by \textsf{Semigroups} are being grown; see \texttt{SetInfoLevel} (\textbf{Reference: InfoLevel}). 
\item[{\texttt{small}}] if this component is set to \texttt{true}, then \textsf{Semigroups} will compute a small subset of \mbox{\texttt{\mdseries\slshape gens}} that generates \mbox{\texttt{\mdseries\slshape S}} at the time that \mbox{\texttt{\mdseries\slshape S}} is created. This will increase the amount of time required to create \mbox{\texttt{\mdseries\slshape S}} substantially, but may decrease the amount of time required for subsequent
calculations with \mbox{\texttt{\mdseries\slshape S}}. If this component is set to \texttt{false}, then \textsf{Semigroups} will return the semigroup generated by \mbox{\texttt{\mdseries\slshape gens}} without modifying \mbox{\texttt{\mdseries\slshape gens}}. The default value for this component is \texttt{false}.

 This option is ignored when passed to \texttt{ClosureSemigroup} (\ref{ClosureSemigroup}) or \texttt{ClosureInverseSemigroup} (\ref{ClosureInverseSemigroup}). 
\item[{\texttt{cong{\textunderscore}by{\textunderscore}ker{\textunderscore}trace{\textunderscore}threshold}}] this should be a positive integer, which specifies a semigroup size. If \mbox{\texttt{\mdseries\slshape S}} is a semigroup with inverse op, and \mbox{\texttt{\mdseries\slshape S}} has a size greater than or equal to this threshold, then any congruence
defined on it may use the "kernel and trace" method to perform calculations.
If its size is less than the threshold, then other methods will be used
instead. The "kernel and trace" method has better complexity than the generic
method, but has large overheads which make it a poor choice for small
semigroups. The default value for this component is \texttt{10 \texttt{\symbol{94}} 5}. See Section \ref{Congruences on inverse semigroups} for more information about the "kernel and trace" method. 
\end{description}
 
\begin{Verbatim}[commandchars=!@|,fontsize=\small,frame=single,label=Example]
  !gapprompt@gap>| !gapinput@S := Semigroup(Transformation([1, 2, 3, 3]),|
  !gapprompt@>| !gapinput@                  rec(hashlen := 100003, small := false));|
  <commutative transformation semigroup of degree 4 with 1 generator>
\end{Verbatim}
 The default values of the options described above are stored in a global
variable named \texttt{SEMIGROUPS.DefaultOptionsRec} (\ref{SEMIGROUPS.DefaultOptionsRec}). If you want to change the default values of these options for a single \textsf{GAP} session, then you can simply redefine the value in \textsf{GAP}. For example, to change the option \texttt{small} from the default value of \mbox{\texttt{\mdseries\slshape false}} use: 
\begin{Verbatim}[commandchars=!@|,fontsize=\small,frame=single,label=Example]
  !gapprompt@gap>| !gapinput@SEMIGROUPS.DefaultOptionsRec.small := true;|
  true
\end{Verbatim}
 If you want to change the default values of the options stored in \texttt{SEMIGROUPS.DefaultOptionsRec} (\ref{SEMIGROUPS.DefaultOptionsRec}) for all \textsf{GAP} sessions, then you can edit these values in the file \texttt{semigroups\texttt{\symbol{45}}5.6.1/gap/options.g}. 

 

\subsection{\textcolor{Chapter }{SEMIGROUPS.DefaultOptionsRec}}
\logpage{[ 6, 3, 1 ]}\nobreak
\hyperdef{L}{X78CF5DCC7C697BB3}{}
{\noindent\textcolor{FuncColor}{$\triangleright$\enspace\texttt{SEMIGROUPS.DefaultOptionsRec\index{SEMIGROUPS.DefaultOptionsRec@\texttt{SEMIGROUPS.DefaultOptionsRec}}
\label{SEMIGROUPS.DefaultOptionsRec}
}\hfill{\scriptsize (global variable)}}\\


 This global variable is a record whose components contain the default values
of certain options for semigroups. A description of these options is given
above in Section \ref{Options when creating semigroups}. 

 The value of \texttt{SEMIGROUPS.DefaultOptionsRec} is defined in the file \texttt{semigroups/gap/options.g}. }

 }

 
\section{\textcolor{Chapter }{Subsemigroups and supersemigroups}}\logpage{[ 6, 4, 0 ]}
\hyperdef{L}{X87AA2EB6810B4631}{}
{
  

\subsection{\textcolor{Chapter }{ClosureSemigroup}}
\logpage{[ 6, 4, 1 ]}\nobreak
\hyperdef{L}{X7BE36790862AE26F}{}
{\noindent\textcolor{FuncColor}{$\triangleright$\enspace\texttt{ClosureSemigroup({\mdseries\slshape S, coll[, opts]})\index{ClosureSemigroup@\texttt{ClosureSemigroup}}
\label{ClosureSemigroup}
}\hfill{\scriptsize (operation)}}\\
\noindent\textcolor{FuncColor}{$\triangleright$\enspace\texttt{ClosureMonoid({\mdseries\slshape S, coll[, opts]})\index{ClosureMonoid@\texttt{ClosureMonoid}}
\label{ClosureMonoid}
}\hfill{\scriptsize (operation)}}\\
\noindent\textcolor{FuncColor}{$\triangleright$\enspace\texttt{ClosureInverseSemigroup({\mdseries\slshape S, coll[, opts]})\index{ClosureInverseSemigroup@\texttt{ClosureInverseSemigroup}}
\label{ClosureInverseSemigroup}
}\hfill{\scriptsize (operation)}}\\
\noindent\textcolor{FuncColor}{$\triangleright$\enspace\texttt{ClosureInverseMonoid({\mdseries\slshape S, coll[, opts]})\index{ClosureInverseMonoid@\texttt{ClosureInverseMonoid}}
\label{ClosureInverseMonoid}
}\hfill{\scriptsize (operation)}}\\
\textbf{\indent Returns:\ }
A semigroup, monoid, inverse semigroup, or inverse monoid.



 These operations return the semigroup, monoid, inverse semigroup or inverse
monoid generated by the argument \mbox{\texttt{\mdseries\slshape S}} and the collection of elements \mbox{\texttt{\mdseries\slshape coll}} after removing duplicates and elements from \mbox{\texttt{\mdseries\slshape coll}} that are already in \mbox{\texttt{\mdseries\slshape S}}. In most cases, the new semigroup knows at least as much information about
its structure as was already known about that of \mbox{\texttt{\mdseries\slshape S}}. 

 When \texttt{X} is any of \texttt{Semigroup} (\textbf{Reference: Semigroup}), \texttt{Monoid} (\textbf{Reference: Monoid}), \texttt{InverseSemigroup} (\textbf{Reference: InverseSemigroup}), or \texttt{InverseMonoid} (\textbf{Reference: InverseMonoid}), the argument \mbox{\texttt{\mdseries\slshape S}} of the operation \texttt{ClosureX} must belong to the category \texttt{IsX}, and \texttt{ClosureX(\mbox{\texttt{\mdseries\slshape S}}, \mbox{\texttt{\mdseries\slshape coll}})} returns an object in the category \texttt{IsX} such that 
\begin{Verbatim}[commandchars=!@|,fontsize=\small,frame=single,label=Example]
          ClosureX(S, coll) = X(S, coll);
\end{Verbatim}
 but may have fewer generators, if for example, \mbox{\texttt{\mdseries\slshape coll}} contains a duplicates or elements already known to belong to \mbox{\texttt{\mdseries\slshape S}}.

 For example, the argument \mbox{\texttt{\mdseries\slshape S}} of \texttt{ClosureInverseSemigroup} must be an inverse semigroup in the category \texttt{IsInverseSemigroup} (\textbf{Reference: IsInverseSemigroup}). \texttt{ClosureInverseSemigroup(\mbox{\texttt{\mdseries\slshape S}}, \mbox{\texttt{\mdseries\slshape coll}})} returns an inverse semigroup which is equal to \texttt{InverseSemigroup(\mbox{\texttt{\mdseries\slshape S}}, \mbox{\texttt{\mdseries\slshape coll}})}. 

 If present, the optional third argument \mbox{\texttt{\mdseries\slshape opts}} should be a record containing the values of the options for the semigroup
being created as described in Section \ref{Options when creating semigroups}.

 
\begin{Verbatim}[commandchars=!@|,fontsize=\small,frame=single,label=Example]
  !gapprompt@gap>| !gapinput@gens := [Transformation([2, 6, 7, 2, 6, 1, 1, 5]),|
  !gapprompt@>| !gapinput@            Transformation([3, 8, 1, 4, 5, 6, 7, 1]),|
  !gapprompt@>| !gapinput@            Transformation([4, 3, 2, 7, 7, 6, 6, 5]),|
  !gapprompt@>| !gapinput@            Transformation([7, 1, 7, 4, 2, 5, 6, 3])];;|
  !gapprompt@gap>| !gapinput@S := Monoid(gens[1]);;|
  !gapprompt@gap>| !gapinput@for x in gens do|
  !gapprompt@>| !gapinput@     S := ClosureSemigroup(S, x);|
  !gapprompt@>| !gapinput@   od;|
  !gapprompt@gap>| !gapinput@S;|
  <transformation monoid of degree 8 with 4 generators>
  !gapprompt@gap>| !gapinput@Size(S);|
  233606
  !gapprompt@gap>| !gapinput@S := Monoid(PartialPerm([1]));|
  <trivial partial perm group of rank 1 with 1 generator>
  !gapprompt@gap>| !gapinput@T := ClosureMonoid(S, [PartialPerm([2 .. 5])]);|
  <partial perm monoid of rank 5 with 2 generators>
  !gapprompt@gap>| !gapinput@One(T);|
  <identity partial perm on [ 1, 2, 3, 4, 5 ]>
  !gapprompt@gap>| !gapinput@T := ClosureSemigroup(S, [PartialPerm([2 .. 5])]);|
  <partial perm semigroup of rank 4 with 2 generators>
  !gapprompt@gap>| !gapinput@One(T);|
  fail
  !gapprompt@gap>| !gapinput@ClosureInverseMonoid(DualSymmetricInverseMonoid(3),|
  !gapprompt@>| !gapinput@                        DClass(DualSymmetricInverseMonoid(3),|
  !gapprompt@>| !gapinput@                               IdentityBipartition(3)));|
  <inverse block bijection monoid of degree 3 with 3 generators>
  !gapprompt@gap>| !gapinput@S := InverseSemigroup(Bipartition([[1, -1, -3], [2, 3, -2]]),|
  !gapprompt@>| !gapinput@                         Bipartition([[1, -3], [2, -2], [3, -1]]));;|
  !gapprompt@gap>| !gapinput@T := ClosureInverseSemigroup(S, DClass(PartitionMonoid(3),|
  !gapprompt@>| !gapinput@IdentityBipartition(3)));|
  <inverse block bijection semigroup of degree 3 with 3 generators>
  !gapprompt@gap>| !gapinput@T := ClosureInverseSemigroup(T, [T.1, T.1, T.1]);|
  <inverse block bijection semigroup of degree 3 with 3 generators>
  !gapprompt@gap>| !gapinput@S := InverseMonoid([|
  !gapprompt@>| !gapinput@PartialPerm([5, 9, 10, 0, 6, 3, 8, 4, 0]),|
  !gapprompt@>| !gapinput@PartialPerm([10, 7, 0, 8, 0, 0, 5, 9, 1])]);;|
  !gapprompt@gap>| !gapinput@x := PartialPerm([|
  !gapprompt@>| !gapinput@5, 1, 7, 3, 10, 0, 2, 12, 0, 14, 11, 0, 16, 0, 0, 0, 0, 6, 9, 15]);|
  [4,3,7,2,1,5,10,14][8,12][13,16][18,6][19,9][20,15](11)
  !gapprompt@gap>| !gapinput@S := ClosureInverseSemigroup(S, x);|
  <inverse partial perm semigroup of rank 19 with 4 generators>
  !gapprompt@gap>| !gapinput@Size(S);|
  9744
  !gapprompt@gap>| !gapinput@T := Idempotents(SymmetricInverseSemigroup(10));;|
  !gapprompt@gap>| !gapinput@S := ClosureInverseSemigroup(S, T);|
  <inverse partial perm semigroup of rank 19 with 14 generators>
\end{Verbatim}
 }

 

\subsection{\textcolor{Chapter }{SubsemigroupByProperty (for a semigroup and function)}}
\logpage{[ 6, 4, 2 ]}\nobreak
\hyperdef{L}{X7E5B4C5A82F9E0E0}{}
{\noindent\textcolor{FuncColor}{$\triangleright$\enspace\texttt{SubsemigroupByProperty({\mdseries\slshape S, func})\index{SubsemigroupByProperty@\texttt{SubsemigroupByProperty}!for a semigroup and function}
\label{SubsemigroupByProperty:for a semigroup and function}
}\hfill{\scriptsize (operation)}}\\
\noindent\textcolor{FuncColor}{$\triangleright$\enspace\texttt{SubsemigroupByProperty({\mdseries\slshape S, func, limit})\index{SubsemigroupByProperty@\texttt{SubsemigroupByProperty}!for a semigroup, function, and limit on the size of the subsemigroup}
\label{SubsemigroupByProperty:for a semigroup, function, and limit on the size of the subsemigroup}
}\hfill{\scriptsize (operation)}}\\
\textbf{\indent Returns:\ }
A semigroup.



 \texttt{SubsemigroupByProperty} returns the subsemigroup of the semigroup \mbox{\texttt{\mdseries\slshape S}} generated by those elements of \mbox{\texttt{\mdseries\slshape S}} fulfilling \mbox{\texttt{\mdseries\slshape func}} (which should be a function returning \texttt{true} or \texttt{false}).

 If no elements of \mbox{\texttt{\mdseries\slshape S}} fulfil \mbox{\texttt{\mdseries\slshape func}}, then \texttt{fail} is returned.

 If the optional third argument \mbox{\texttt{\mdseries\slshape limit}} is present and a positive integer, then once the subsemigroup has at least \mbox{\texttt{\mdseries\slshape limit}} elements the computation stops. 
\begin{Verbatim}[commandchars=!@|,fontsize=\small,frame=single,label=Example]
  !gapprompt@gap>| !gapinput@func := function(x)|
  !gapprompt@>| !gapinput@     local n;|
  !gapprompt@>| !gapinput@     n := DegreeOfTransformation(x);|
  !gapprompt@>| !gapinput@     return 1 ^ x <> 1 and ForAll([1 .. n], y -> y = 1 or y ^ x = y);|
  !gapprompt@>| !gapinput@   end;|
  function( x ) ... end
  !gapprompt@gap>| !gapinput@T := SubsemigroupByProperty(FullTransformationSemigroup(3), func);|
  <transformation semigroup of size 2, degree 3 with 2 generators>
  !gapprompt@gap>| !gapinput@T := SubsemigroupByProperty(FullTransformationSemigroup(4), func);|
  <transformation semigroup of size 3, degree 4 with 3 generators>
  !gapprompt@gap>| !gapinput@T := SubsemigroupByProperty(FullTransformationSemigroup(5), func);|
  <transformation semigroup of size 4, degree 5 with 4 generators>
\end{Verbatim}
 }

 

\subsection{\textcolor{Chapter }{InverseSubsemigroupByProperty}}
\logpage{[ 6, 4, 3 ]}\nobreak
\hyperdef{L}{X832AEDCC7BA9E5F5}{}
{\noindent\textcolor{FuncColor}{$\triangleright$\enspace\texttt{InverseSubsemigroupByProperty({\mdseries\slshape S, func})\index{InverseSubsemigroupByProperty@\texttt{InverseSubsemigroupByProperty}}
\label{InverseSubsemigroupByProperty}
}\hfill{\scriptsize (operation)}}\\
\textbf{\indent Returns:\ }
An inverse semigroup.



 \texttt{InverseSubsemigroupByProperty} returns the inverse subsemigroup of the inverse semigroup \mbox{\texttt{\mdseries\slshape S}} generated by those elements of \mbox{\texttt{\mdseries\slshape S}} fulfilling \mbox{\texttt{\mdseries\slshape func}} (which should be a function returning \texttt{true} or \texttt{false}).

 If no elements of \mbox{\texttt{\mdseries\slshape S}} fulfil \mbox{\texttt{\mdseries\slshape func}}, then \texttt{fail} is returned.

 If the optional third argument \mbox{\texttt{\mdseries\slshape limit}} is present and a positive integer, then once the subsemigroup has at least \mbox{\texttt{\mdseries\slshape limit}} elements the computation stops. 
\begin{Verbatim}[commandchars=!@|,fontsize=\small,frame=single,label=Example]
  !gapprompt@gap>| !gapinput@IsIsometry := function(f)|
  !gapprompt@>| !gapinput@local n, i, j, k, l;|
  !gapprompt@>| !gapinput@ n := RankOfPartialPerm(f);|
  !gapprompt@>| !gapinput@ for i in [1 .. n - 1] do|
  !gapprompt@>| !gapinput@   k := DomainOfPartialPerm(f)[i];|
  !gapprompt@>| !gapinput@   for j in [i + 1 .. n] do|
  !gapprompt@>| !gapinput@     l := DomainOfPartialPerm(f)[j];|
  !gapprompt@>| !gapinput@     if not AbsInt(k ^ f - l ^ f) = AbsInt(k - l) then|
  !gapprompt@>| !gapinput@       return false;|
  !gapprompt@>| !gapinput@     fi;|
  !gapprompt@>| !gapinput@   od;|
  !gapprompt@>| !gapinput@ od;|
  !gapprompt@>| !gapinput@ return true;|
  !gapprompt@>| !gapinput@end;;|
  !gapprompt@gap>| !gapinput@S := InverseSubsemigroupByProperty(SymmetricInverseSemigroup(5),|
  !gapprompt@>| !gapinput@IsIsometry);;|
  !gapprompt@gap>| !gapinput@Size(S);|
  142
\end{Verbatim}
 }

 }

 
\section{\textcolor{Chapter }{Changing the representation of a semigroup}}\label{Changing the representation of a semigroup}
\logpage{[ 6, 5, 0 ]}
\hyperdef{L}{X82CCC1A781650878}{}
{
  The \textsf{Semigroups} package provides two convenient constructors \texttt{IsomorphismSemigroup} (\ref{IsomorphismSemigroup}) and \texttt{IsomorphismMonoid} (\ref{IsomorphismMonoid}) for changing the representation of a given semigroup or monoid. These methods
can be used to find an isomorphism from any semigroup to a semigroup of any
other type, provided such an isomorphism exists. 

 Note that at present neither \texttt{IsomorphismSemigroup} (\ref{IsomorphismSemigroup}) nor \texttt{IsomorphismMonoid} (\ref{IsomorphismMonoid}) can be used to determine whether two given semigroups, or monoids, are
isomorphic.

 Some methods for \texttt{IsomorphismSemigroup} (\ref{IsomorphismSemigroup}) and \texttt{IsomorphismMonoid} (\ref{IsomorphismMonoid}) are based on methods for the \textsf{GAP} library operations: 
\begin{itemize}
\item  \texttt{IsomorphismReesMatrixSemigroup} (\textbf{Reference: IsomorphismReesMatrixSemigroup}), 
\item  \texttt{AntiIsomorphismTransformationSemigroup} (\textbf{Reference: AntiIsomorphismTransformationSemigroup}), 
\item  \texttt{IsomorphismTransformationSemigroup} (\textbf{Reference: IsomorphismTransformationSemigroup}) and \texttt{IsomorphismTransformationMonoid} (\textbf{Reference: IsomorphismTransformationMonoid}), 
\item  \texttt{IsomorphismPartialPermSemigroup} (\textbf{Reference: IsomorphismPartialPermSemigroup}) and \texttt{IsomorphismPartialPermMonoid} (\textbf{Reference: IsomorphismPartialPermMonoid}), 
\item  \texttt{IsomorphismFpSemigroup} (\textbf{Reference: IsomorphismFpSemigroup}) and \texttt{IsomorphismFpMonoid}. 
\end{itemize}
 The operation \texttt{IsomorphismMonoid} (\ref{IsomorphismMonoid}) can be used to return an isomorphism from a semigroup which is mathematically
a monoid (but does not below to the category of monoids in \textsf{GAP} \texttt{IsMonoid} (\textbf{Reference: IsMonoid})) into a monoid. This is the primary purpose of the operation \texttt{IsomorphismMonoid} (\ref{IsomorphismMonoid}). Either \texttt{IsomorphismSemigroup} (\ref{IsomorphismSemigroup}) or \texttt{IsomorphismMonoid} (\ref{IsomorphismMonoid}) can be used to change the representation of a monoid, but only the latter is
guaranteed to return an object in the category of monoids. 

 
\begin{Verbatim}[commandchars=!@|,fontsize=\small,frame=single,label=Example]
  !gapprompt@gap>| !gapinput@S := Monoid(Transformation([1, 4, 6, 2, 5, 3, 7, 8, 9, 9]),|
  !gapprompt@>| !gapinput@               Transformation([6, 3, 2, 7, 5, 1, 8, 8, 9, 9]));;|
  !gapprompt@gap>| !gapinput@AsSemigroup(IsBooleanMatSemigroup, S);|
  <monoid of 10x10 boolean matrices with 2 generators>
  !gapprompt@gap>| !gapinput@AsMonoid(IsBooleanMatMonoid, S);|
  <monoid of 10x10 boolean matrices with 2 generators>
  !gapprompt@gap>| !gapinput@S := Semigroup(Transformation([1, 4, 6, 2, 5, 3, 7, 8, 9, 9]),|
  !gapprompt@>| !gapinput@                  Transformation([6, 3, 2, 7, 5, 1, 8, 8, 9, 9]));;|
  !gapprompt@gap>| !gapinput@AsSemigroup(IsBooleanMatSemigroup, S);|
  <semigroup of 10x10 boolean matrices with 2 generators>
  !gapprompt@gap>| !gapinput@AsMonoid(IsBooleanMatMonoid, S);|
  <monoid of 8x8 boolean matrices with 2 generators>
  !gapprompt@gap>| !gapinput@M := Monoid([|
  !gapprompt@>| !gapinput@Bipartition([[1, -3], [2, 3, 6], [4, 7, -6], [5, -8], [8, -4, -5],|
  !gapprompt@>| !gapinput@             [-1], [-2], [-7]]),|
  !gapprompt@>| !gapinput@Bipartition([[1, 3, -6], [2, -8], [4, 8, -1], [5], [6, -3, -4],|
  !gapprompt@>| !gapinput@             [7], [-2], [-5], [-7]]),|
  !gapprompt@>| !gapinput@Bipartition([[1, 2, 4, -3, -7, -8], [3, 5, 6, 8, -4, -6],|
  !gapprompt@>| !gapinput@             [7, -1, -2, -5]])]);;|
  !gapprompt@gap>| !gapinput@AsMonoid(IsPBRMonoid, M);|
  <pbr monoid of size 163, degree 163 with 3 generators>
  !gapprompt@gap>| !gapinput@AsSemigroup(IsPBRSemigroup, M);|
  <pbr semigroup of size 163, degree 8 with 4 generators>
\end{Verbatim}
 There are some further methods in \textsf{Semigroups} for obtaining an isomorphism from a Rees matrix, or
0\texttt{\symbol{45}}matrix, semigroup to another such semigroup with
particular properties; \texttt{RMSNormalization} (\ref{RMSNormalization}) and \texttt{RZMSNormalization} (\ref{RZMSNormalization}). 

\subsection{\textcolor{Chapter }{IsomorphismSemigroup}}
\logpage{[ 6, 5, 1 ]}\nobreak
\hyperdef{L}{X838F18E87F765697}{}
{\noindent\textcolor{FuncColor}{$\triangleright$\enspace\texttt{IsomorphismSemigroup({\mdseries\slshape filt, S})\index{IsomorphismSemigroup@\texttt{IsomorphismSemigroup}}
\label{IsomorphismSemigroup}
}\hfill{\scriptsize (operation)}}\\
\textbf{\indent Returns:\ }
An isomorphism of semigroups.



 \texttt{IsomorphismSemigroup} can be used to find an isomorphism from a given semigroup to a semigroup of
another type, provided such an isomorphism exists. 

 The first argument \mbox{\texttt{\mdseries\slshape filt}} must be of the form \texttt{IsXSemigroup}, for example, \texttt{IsTransformationSemigroup} (\textbf{Reference: IsTransformationSemigroup}), \texttt{IsFpSemigroup} (\textbf{Reference: IsFpSemigroup}), and \texttt{IsPBRSemigroup} (\ref{IsPBRSemigroup}) are some possible values for \mbox{\texttt{\mdseries\slshape filt}}. Note that \mbox{\texttt{\mdseries\slshape filt}} should not be of the form \texttt{IsXMonoid}; see \texttt{IsomorphismMonoid} (\ref{IsomorphismMonoid}). The second argument \mbox{\texttt{\mdseries\slshape S}} should be a semigroup. 

 \texttt{IsomorphismSemigroup} returns an isomorphism from \mbox{\texttt{\mdseries\slshape S}} to a semigroup \mbox{\texttt{\mdseries\slshape T}} of the type described by \mbox{\texttt{\mdseries\slshape filt}}, if such an isomorphism exists. More precisely, if \texttt{T} is the range of the returned isomorphism, then \texttt{\mbox{\texttt{\mdseries\slshape filt}}(T)} will return \texttt{true}. For example, if \mbox{\texttt{\mdseries\slshape filt}} is \texttt{IsTransformationSemigroup}, then the range of the returned isomorphism will be a transformation
semigroup. 

 An error is returned if there is no isomorphism from \mbox{\texttt{\mdseries\slshape S}} to a semigroup satisfying \mbox{\texttt{\mdseries\slshape filt}}. For example, there is no method for \texttt{IsomorphismSemigroup} when \mbox{\texttt{\mdseries\slshape filt}} is, say, \texttt{IsReesMatrixSemigroup} (\textbf{Reference: IsReesMatrixSemigroup}) and when \mbox{\texttt{\mdseries\slshape S}} is a non\texttt{\symbol{45}}simple semigroup. Similarly, there is no method
when \mbox{\texttt{\mdseries\slshape filt}} is \texttt{IsPartialPermSemigroup} (\textbf{Reference: IsPartialPermSemigroup}) and when \mbox{\texttt{\mdseries\slshape S}} is a non\texttt{\symbol{45}}inverse semigroup. 

 In some cases, if no better method is installed, \texttt{IsomorphismSemigroup} returns an isomorphism found by composing an isomorphism from \mbox{\texttt{\mdseries\slshape S}} to a transformation semigroup \texttt{T}, and an isomorphism from \texttt{T} to a semigroup of type \mbox{\texttt{\mdseries\slshape filt}}. 

 Note that if the argument \mbox{\texttt{\mdseries\slshape S}} belongs to the category of monoids \texttt{IsMonoid} (\textbf{Reference: IsMonoid}), then \texttt{IsomorphismSemigroup} will often, but not always, return a monoid isomorphism. 
\begin{Verbatim}[commandchars=!@|,fontsize=\small,frame=single,label=Example]
  !gapprompt@gap>| !gapinput@S := Semigroup([|
  !gapprompt@>| !gapinput@Bipartition([|
  !gapprompt@>| !gapinput@  [1, 2], [3, 6, -2], [4, 5, -3, -4], [-1, -6], [-5]]),|
  !gapprompt@>| !gapinput@Bipartition([|
  !gapprompt@>| !gapinput@  [1, -4], [2, 3, 4, 5], [6], [-1, -6], [-2, -3], [-5]])]);|
  <bipartition semigroup of degree 6 with 2 generators>
  !gapprompt@gap>| !gapinput@IsomorphismSemigroup(IsTransformationSemigroup, S);|
  <bipartition semigroup of size 11, degree 6 with 2 generators> ->
  <transformation semigroup of size 11, degree 12 with 2 generators>
  !gapprompt@gap>| !gapinput@IsomorphismSemigroup(IsBooleanMatSemigroup, S);|
  <bipartition semigroup of size 11, degree 6 with 2 generators> ->
  <semigroup of size 11, 12x12 boolean matrices with 2 generators>
  !gapprompt@gap>| !gapinput@IsomorphismSemigroup(IsFpSemigroup, S);|
  <bipartition semigroup of size 11, degree 6 with 2 generators> ->
  <fp semigroup with 2 generators and 5 relations of length 27>
  !gapprompt@gap>| !gapinput@S := InverseSemigroup([|
  !gapprompt@>| !gapinput@PartialPerm([1, 2, 3, 6, 8, 10],|
  !gapprompt@>| !gapinput@            [2, 6, 7, 9, 1, 5]),|
  !gapprompt@>| !gapinput@PartialPerm([1, 2, 3, 4, 6, 7, 8, 10],|
  !gapprompt@>| !gapinput@            [3, 8, 1, 9, 4, 10, 5, 6])]);;|
  !gapprompt@gap>| !gapinput@IsomorphismSemigroup(IsBipartitionSemigroup, S);|
  <inverse partial perm semigroup of rank 10 with 2 generators> ->
  <inverse bipartition semigroup of degree 10 with 2 generators>
  !gapprompt@gap>| !gapinput@S := SymmetricInverseMonoid(4);|
  <symmetric inverse monoid of degree 4>
  !gapprompt@gap>| !gapinput@IsomorphismSemigroup(IsBlockBijectionSemigroup, S);|
  <symmetric inverse monoid of degree 4> ->
  <inverse block bijection monoid of degree 5 with 3 generators>
  !gapprompt@gap>| !gapinput@Size(Range(last));|
  209
  !gapprompt@gap>| !gapinput@S := Semigroup([|
  !gapprompt@>| !gapinput@PartialPerm([3, 1]), PartialPerm([1, 3, 4])]);;|
  !gapprompt@gap>| !gapinput@IsomorphismSemigroup(IsBlockBijectionSemigroup, S);|
  <partial perm semigroup of rank 3 with 2 generators> ->
  <block bijection semigroup of degree 5 with 2 generators>
\end{Verbatim}
 }

 

\subsection{\textcolor{Chapter }{IsomorphismMonoid}}
\logpage{[ 6, 5, 2 ]}\nobreak
\hyperdef{L}{X83D03BE678C9974F}{}
{\noindent\textcolor{FuncColor}{$\triangleright$\enspace\texttt{IsomorphismMonoid({\mdseries\slshape filt, S})\index{IsomorphismMonoid@\texttt{IsomorphismMonoid}}
\label{IsomorphismMonoid}
}\hfill{\scriptsize (operation)}}\\
\textbf{\indent Returns:\ }
An isomorphism of monoids.



 \texttt{IsomorphismMonoid} can be used to find an isomorphism from a given semigroup which is
mathematically a monoid (but might not belong to the category of monoids in \textsf{GAP}) to a monoid, provided such an isomorphism exists. 

 The first argument \mbox{\texttt{\mdseries\slshape filt}} must be of the form \texttt{IsXMonoid}, for example, \texttt{IsTransformationMonoid} (\textbf{Reference: IsTransformationMonoid}), \texttt{IsFpMonoid} (\textbf{Reference: IsFpMonoid}), and \texttt{IsBipartitionMonoid} (\ref{IsBipartitionMonoid}) are some possible values for \mbox{\texttt{\mdseries\slshape filt}}. Note that \mbox{\texttt{\mdseries\slshape filt}} should not be of the form \texttt{IsXSemigroup}; see \texttt{IsomorphismSemigroup} (\ref{IsomorphismSemigroup}). The second argument \mbox{\texttt{\mdseries\slshape S}} should be a semigroup which is mathematically a monoid but which may or may
not belong to the category \texttt{IsMonoid} (\textbf{Reference: IsMonoid}) of monoids in \textsf{GAP}, i.e. \mbox{\texttt{\mdseries\slshape S}} must satisfy \texttt{IsMonoidAsSemigroup} (\ref{IsMonoidAsSemigroup}). 

 \texttt{IsomorphismMonoid} returns a monoid isomorphism from \mbox{\texttt{\mdseries\slshape S}} to a semigroup \mbox{\texttt{\mdseries\slshape T}} of the type described by \mbox{\texttt{\mdseries\slshape filt}}, if such an isomorphism exists. In this context, a \emph{monoid isomorphism} is a semigroup isomorphism that maps the \texttt{MultiplicativeNeutralElement} (\textbf{Reference: MultiplicativeNeutralElement}) of \mbox{\texttt{\mdseries\slshape S}} to the \texttt{One} (\textbf{Reference: One}) of \mbox{\texttt{\mdseries\slshape T}}. If \texttt{T} is the range of the returned isomorphism, then \texttt{\mbox{\texttt{\mdseries\slshape filt}}(T)} will return \texttt{true}. For example, if \mbox{\texttt{\mdseries\slshape filt}} is \texttt{IsTransformationMonoid}, then the range of the returned isomorphism will be a transformation monoid. 

 An error is returned if there is no isomorphism from \mbox{\texttt{\mdseries\slshape S}} to a monoid satisfying \mbox{\texttt{\mdseries\slshape filt}}. For example, there is no method for \texttt{IsomorphismMonoid} when \mbox{\texttt{\mdseries\slshape filt}} is, say, \texttt{IsReesZeroMatrixSemigroup} (\textbf{Reference: IsReesZeroMatrixSemigroup}) and when \mbox{\texttt{\mdseries\slshape S}} is a not 0\texttt{\symbol{45}}simple. Similarly, there is no method when \mbox{\texttt{\mdseries\slshape filt}} is \texttt{IsPartialPermMonoid} (\textbf{Reference: IsPartialPermMonoid}) and when \mbox{\texttt{\mdseries\slshape S}} is a non\texttt{\symbol{45}}inverse monoid. 

 In some cases, if no better method is installed, \texttt{IsomorphismMonoid} returns an isomorphism found by composing an isomorphism from \mbox{\texttt{\mdseries\slshape S}} to a transformation monoid \texttt{T}, and an isomorphism from \texttt{T} to a monoid of type \mbox{\texttt{\mdseries\slshape filt}}. 
\begin{Verbatim}[commandchars=!@|,fontsize=\small,frame=single,label=Example]
  !gapprompt@gap>| !gapinput@S := Semigroup(Transformation([1, 4, 6, 2, 5, 3, 7, 8, 9, 9]),|
  !gapprompt@>| !gapinput@                  Transformation([6, 3, 2, 7, 5, 1, 8, 8, 9, 9]));|
  <transformation semigroup of degree 10 with 2 generators>
  !gapprompt@gap>| !gapinput@IsomorphismMonoid(IsTransformationMonoid, S);|
  <transformation semigroup of degree 10 with 2 generators> ->
  <transformation monoid of degree 8 with 2 generators>
  !gapprompt@gap>| !gapinput@IsomorphismMonoid(IsBooleanMatMonoid, S);|
  <transformation semigroup of degree 10 with 2 generators> ->
  <monoid of 8x8 boolean matrices with 2 generators>
  !gapprompt@gap>| !gapinput@IsomorphismMonoid(IsFpMonoid, S);|
  <transformation semigroup of degree 10 with 2 generators> ->
  <fp monoid with 2 generators and 17 relations of length 278>
\end{Verbatim}
 }

 

\subsection{\textcolor{Chapter }{AsSemigroup}}
\logpage{[ 6, 5, 3 ]}\nobreak
\hyperdef{L}{X80ED104F85AE5134}{}
{\noindent\textcolor{FuncColor}{$\triangleright$\enspace\texttt{AsSemigroup({\mdseries\slshape filt, S})\index{AsSemigroup@\texttt{AsSemigroup}}
\label{AsSemigroup}
}\hfill{\scriptsize (operation)}}\\
\textbf{\indent Returns:\ }
A semigroup.



 \texttt{AsSemigroup(\mbox{\texttt{\mdseries\slshape filt}}, \mbox{\texttt{\mdseries\slshape S}})} is just shorthand for \texttt{Range(IsomorphismSemigroup(\mbox{\texttt{\mdseries\slshape filt}}, \mbox{\texttt{\mdseries\slshape S}}))}, when \mbox{\texttt{\mdseries\slshape S}} is a semigroup; see \texttt{IsomorphismSemigroup} (\ref{IsomorphismSemigroup}) for more details. 

 Note that if the argument \mbox{\texttt{\mdseries\slshape S}} belongs to the category of monoids \texttt{IsMonoid} (\textbf{Reference: IsMonoid}), then \texttt{AsSemigroup} will often, but not always, return a monoid. A monoid is not returned if there
is not a good monoid isomorphism from \mbox{\texttt{\mdseries\slshape S}} to a monoid of the required type, but there is a good semigroup isomorphism. 

 If it is not possible to convert the semigroup \mbox{\texttt{\mdseries\slshape S}} to a semigroup of type \mbox{\texttt{\mdseries\slshape filt}}, then an error is given. 
\begin{Verbatim}[commandchars=!@|,fontsize=\small,frame=single,label=Example]
  !gapprompt@gap>| !gapinput@S := Semigroup([|
  !gapprompt@>| !gapinput@Bipartition([|
  !gapprompt@>| !gapinput@  [1, 2], [3, 6, -2], [4, 5, -3, -4], [-1, -6], [-5]]),|
  !gapprompt@>| !gapinput@Bipartition([|
  !gapprompt@>| !gapinput@  [1, -4], [2, 3, 4, 5], [6], [-1, -6], [-2, -3], [-5]])]);|
  <bipartition semigroup of degree 6 with 2 generators>
  !gapprompt@gap>| !gapinput@AsSemigroup(IsTransformationSemigroup, S);|
  <transformation semigroup of size 11, degree 12 with 2 generators>
  !gapprompt@gap>| !gapinput@S := Semigroup([|
  !gapprompt@>| !gapinput@Bipartition([|
  !gapprompt@>| !gapinput@  [1, 2], [3, 6, -2], [4, 5, -3, -4], [-1, -6], [-5]]),|
  !gapprompt@>| !gapinput@Bipartition([|
  !gapprompt@>| !gapinput@  [1, -4], [2, 3, 4, 5], [6], [-1, -6], [-2, -3], [-5]])]);|
  <bipartition semigroup of degree 6 with 2 generators>
  !gapprompt@gap>| !gapinput@AsSemigroup(IsTransformationSemigroup, S);|
  <transformation semigroup of size 11, degree 12 with 2 generators>
  !gapprompt@gap>| !gapinput@T := Semigroup(Transformation([2, 2, 3]),|
  !gapprompt@>| !gapinput@                  Transformation([3, 1, 3]));|
  <transformation semigroup of degree 3 with 2 generators>
  !gapprompt@gap>| !gapinput@S := AsSemigroup(IsMatrixOverFiniteFieldSemigroup, GF(5), T);|
  <semigroup of 3x3 matrices over GF(5) with 2 generators>
  !gapprompt@gap>| !gapinput@Size(S);|
  5
\end{Verbatim}
 }

 

\subsection{\textcolor{Chapter }{AsMonoid}}
\logpage{[ 6, 5, 4 ]}\nobreak
\hyperdef{L}{X7B22038F832B9C0F}{}
{\noindent\textcolor{FuncColor}{$\triangleright$\enspace\texttt{AsMonoid({\mdseries\slshape [filt, ]S})\index{AsMonoid@\texttt{AsMonoid}}
\label{AsMonoid}
}\hfill{\scriptsize (operation)}}\\
\textbf{\indent Returns:\ }
A monoid or \texttt{fail}.



 \texttt{AsMonoid(\mbox{\texttt{\mdseries\slshape filt}}, \mbox{\texttt{\mdseries\slshape S}})} is just shorthand for \texttt{Range(IsomorphismMonoid(\mbox{\texttt{\mdseries\slshape filt}}, \mbox{\texttt{\mdseries\slshape S}}))}, when \mbox{\texttt{\mdseries\slshape S}} is a semigroup or monoid; see \texttt{IsomorphismMonoid} (\ref{IsomorphismMonoid}) for more details. 

 If the first argument \mbox{\texttt{\mdseries\slshape filt}} is omitted and the semigroup \mbox{\texttt{\mdseries\slshape S}} is mathematically a monoid which does not belong to the category of monoids in \textsf{GAP}, then \texttt{AsMonoid} returns a monoid (in the category of monoids) isomorphic to \mbox{\texttt{\mdseries\slshape S}} and of the same type as \mbox{\texttt{\mdseries\slshape S}}. If \mbox{\texttt{\mdseries\slshape S}} is already in the category of monoids and the first argument \mbox{\texttt{\mdseries\slshape filt}} is omitted, then \mbox{\texttt{\mdseries\slshape S}} is returned. 

 If the first argument \mbox{\texttt{\mdseries\slshape filt}} is omitted and the semigroup \mbox{\texttt{\mdseries\slshape S}} is not a monoid, i.e. it does not satisfy \texttt{IsMonoidAsSemigroup} (\ref{IsMonoidAsSemigroup}), then \texttt{fail} is returned. 
\begin{Verbatim}[commandchars=!@|,fontsize=\small,frame=single,label=Example]
  !gapprompt@gap>| !gapinput@S := Semigroup(Transformation([1, 4, 6, 2, 5, 3, 7, 8, 9, 9]),|
  !gapprompt@>| !gapinput@                  Transformation([6, 3, 2, 7, 5, 1, 8, 8, 9, 9]));;|
  !gapprompt@gap>| !gapinput@AsMonoid(S);|
  <transformation monoid of degree 8 with 2 generators>
  !gapprompt@gap>| !gapinput@AsSemigroup(IsBooleanMatSemigroup, S);|
  <semigroup of 10x10 boolean matrices with 2 generators>
  !gapprompt@gap>| !gapinput@AsMonoid(IsBooleanMatMonoid, S);|
  <monoid of 8x8 boolean matrices with 2 generators>
  !gapprompt@gap>| !gapinput@S := Monoid(Bipartition([[1, -1, -3], [2, 3], [-2]]),|
  !gapprompt@>| !gapinput@               Bipartition([[1, -1], [2, 3, -3], [-2]]));|
  <bipartition monoid of degree 3 with 2 generators>
  !gapprompt@gap>| !gapinput@AsMonoid(IsTransformationMonoid, S);|
  <transformation monoid of size 3, degree 3 with 2 generators>
  !gapprompt@gap>| !gapinput@AsMonoid(S);|
  <bipartition monoid of size 3, degree 3 with 2 generators>
\end{Verbatim}
 }

 

\subsection{\textcolor{Chapter }{IsomorphismPermGroup}}
\logpage{[ 6, 5, 5 ]}\nobreak
\hyperdef{L}{X80B7B1C783AA1567}{}
{\noindent\textcolor{FuncColor}{$\triangleright$\enspace\texttt{IsomorphismPermGroup({\mdseries\slshape S})\index{IsomorphismPermGroup@\texttt{IsomorphismPermGroup}}
\label{IsomorphismPermGroup}
}\hfill{\scriptsize (attribute)}}\\
\textbf{\indent Returns:\ }
 An isomorphism. 



 If the semigroup \mbox{\texttt{\mdseries\slshape S}} is mathematically a group, so that it satisfies \texttt{IsGroupAsSemigroup} (\ref{IsGroupAsSemigroup}), then \texttt{IsomorphismPermGroup} returns an isomorphism to a permutation group.

 If \mbox{\texttt{\mdseries\slshape S}} is not a group then an error is given.

 See also \texttt{IsomorphismPermGroup} (\textbf{Reference: IsomorphismPermGroup}). 
\begin{Verbatim}[commandchars=!@|,fontsize=\small,frame=single,label=Example]
  !gapprompt@gap>| !gapinput@S := Semigroup(Transformation([2, 2, 3, 4, 6, 8, 5, 5]),|
  !gapprompt@>| !gapinput@                  Transformation([3, 3, 8, 2, 5, 6, 4, 4]));;|
  !gapprompt@gap>| !gapinput@IsGroupAsSemigroup(S);|
  true
  !gapprompt@gap>| !gapinput@iso := IsomorphismPermGroup(S);;|
  !gapprompt@gap>| !gapinput@Source(iso) = S and Range(iso) = Group([(5, 6, 8), (2, 3, 8, 4)]);|
  true
  !gapprompt@gap>| !gapinput@StructureDescription(Range(IsomorphismPermGroup(S)));|
  "S6"
  !gapprompt@gap>| !gapinput@S := Range(IsomorphismPartialPermSemigroup(SymmetricGroup(4)));|
  <partial perm group of size 24, rank 4 with 2 generators>
  !gapprompt@gap>| !gapinput@Range(IsomorphismPermGroup(S));|
  Group([ (1,2,3,4), (1,2) ])
  !gapprompt@gap>| !gapinput@G := GroupOfUnits(PartitionMonoid(4));|
  <block bijection group of degree 4 with 2 generators>
  !gapprompt@gap>| !gapinput@StructureDescription(G);|
  "S4"
  !gapprompt@gap>| !gapinput@iso := IsomorphismPermGroup(G);;|
  !gapprompt@gap>| !gapinput@RespectsMultiplication(iso);|
  true
  !gapprompt@gap>| !gapinput@inv := InverseGeneralMapping(iso);;|
  !gapprompt@gap>| !gapinput@ForAll(G, x -> (x ^ iso) ^ inv = x);|
  true
  !gapprompt@gap>| !gapinput@ForAll(G, x -> ForAll(G, y -> (x * y) ^ iso = x ^ iso * y ^ iso));|
  true
\end{Verbatim}
 }

 

\subsection{\textcolor{Chapter }{RZMSNormalization}}
\logpage{[ 6, 5, 6 ]}\nobreak
\hyperdef{L}{X870210EA7912B52A}{}
{\noindent\textcolor{FuncColor}{$\triangleright$\enspace\texttt{RZMSNormalization({\mdseries\slshape R})\index{RZMSNormalization@\texttt{RZMSNormalization}}
\label{RZMSNormalization}
}\hfill{\scriptsize (attribute)}}\\
\textbf{\indent Returns:\ }
An isomorphism.



 If \mbox{\texttt{\mdseries\slshape R}} is a Rees 0\texttt{\symbol{45}}matrix semigroup $M^{0}[I, T, \Lambda; P]$ then \texttt{RZMSNormalization} returns an isomorphism from \mbox{\texttt{\mdseries\slshape R}} to a \emph{normalized} Rees 0\texttt{\symbol{45}}matrix semigroup $S = M^{0}[I, T, \Lambda; Q]$. The structure matrix $Q$ is obtained by \emph{normalizing} the matrix $P$ (see \texttt{Matrix} (\textbf{Reference: Matrix})) and has the following properties: 
\begin{itemize}
\item  The matrix $Q$ is in block diagonal form, and the blocks are ordered by decreasing size along
the leading diagonal (the size of a block is defined to be the number of rows
it contains multiplied by the number of columns it contains). 

 If the index sets $I$ and $\Lambda$ are partitioned into $k$ parts according to the \texttt{RZMSConnectedComponents} (\ref{RZMSConnectedComponents}) of $S$, giving a disjoint union $I=I_1\cup\ldots\cup I_k$ and $\Lambda=\Lambda_1\cup\ldots\cup\Lambda_k$, then the $r$th block corresponds to the sub\texttt{\symbol{45}}matrix $Q_{r}$ of $Q$ defined by $I_{r}$ and $\Lambda_{r}$. 
\item  The first non\texttt{\symbol{45}}zero entry in a row occurs no sooner than the
first non\texttt{\symbol{45}}zero entry in any previous row. 
\item  The first non\texttt{\symbol{45}}zero entry in a column occurs no sooner than
the first non\texttt{\symbol{45}}zero entry in any previous column. 
\item  The previous two items imply that if the matrix $P$ has any rows/columns consisting entirely of zeroes, then these will become the
final rows/columns of $Q$. 
\end{itemize}
 Furthermore, if $T$ is a group (i.e. a semigroup for which \texttt{IsGroupAsSemigroup} (\ref{IsGroupAsSemigroup}) returns \texttt{true}), then the non\texttt{\symbol{45}}zero entries of the structure matrix $Q$ are chosen such that the following hold: 
\begin{itemize}
\item  The first non\texttt{\symbol{45}}zero entry of every row and every column is
equal to the identity of $T$. 
\item  For each $r$, let $Q_{r}$ be the sub\texttt{\symbol{45}}matrix of $Q$ defined by $I_r$ and $\Lambda_r$ (as above), and let $T_r$ be the subsemigroup of $T$ generated by the non\texttt{\symbol{45}}zero entries of $Q_{r}$. Then the idempotent generated subsemigroup of $S$ is equal to: 
\begin{itemize}
\item  $\bigcup_{r=1}^{k} M^{0}[I_r, T_r, \Lambda_r, Q_r]$, where the zeroes of these Rees 0\texttt{\symbol{45}}matrix semigroups are
all identified with the zero of $S$. 
\end{itemize}
 
\end{itemize}
 The normalization given by \texttt{RZMSNormalization} is based on Theorem 2 of \cite{Graham1968Graph} and is sometimes called \emph{Graham normal form}. Note that isomorphic Rees 0\texttt{\symbol{45}}matrix semigroups can have
normalizations which are not equal.

 
\begin{Verbatim}[commandchars=!@|,fontsize=\small,frame=single,label=Example]
  !gapprompt@gap>| !gapinput@R := ReesZeroMatrixSemigroup(Group(()),|
  !gapprompt@>| !gapinput@[[0, (), 0],|
  !gapprompt@>| !gapinput@ [(), 0, 0],|
  !gapprompt@>| !gapinput@ [0, 0, ()]]);|
  <Rees 0-matrix semigroup 3x3 over Group(())>
  !gapprompt@gap>| !gapinput@iso := RZMSNormalization(R);|
  <Rees 0-matrix semigroup 3x3 over Group(())> ->
  <Rees 0-matrix semigroup 3x3 over Group(())>
  !gapprompt@gap>| !gapinput@S := Range(iso);|
  <Rees 0-matrix semigroup 3x3 over Group(())>
  !gapprompt@gap>| !gapinput@Matrix(S);|
  [ [ (), 0, 0 ], [ 0, (), 0 ], [ 0, 0, () ] ]
  !gapprompt@gap>| !gapinput@R := ReesZeroMatrixSemigroup(SymmetricGroup(4),|
  !gapprompt@>| !gapinput@[[0, 0, 0, (1, 3, 2)],|
  !gapprompt@>| !gapinput@ [(2, 3), 0, 0, 0],|
  !gapprompt@>| !gapinput@ [0, 0, (1, 3), (1, 2)],|
  !gapprompt@>| !gapinput@ [0, (4, 1, 2, 3), 0, 0]]);|
  <Rees 0-matrix semigroup 4x4 over Sym( [ 1 .. 4 ] )>
  !gapprompt@gap>| !gapinput@S := Range(RZMSNormalization(R));|
  <Rees 0-matrix semigroup 4x4 over Sym( [ 1 .. 4 ] )>
  !gapprompt@gap>| !gapinput@Matrix(S);|
  [ [ (), (), 0, 0 ], [ 0, (), 0, 0 ], [ 0, 0, (), 0 ], [ 0, 0, 0, () ]
   ]
\end{Verbatim}
 }

 

\subsection{\textcolor{Chapter }{RMSNormalization}}
\logpage{[ 6, 5, 7 ]}\nobreak
\hyperdef{L}{X80DE617E841E5BA0}{}
{\noindent\textcolor{FuncColor}{$\triangleright$\enspace\texttt{RMSNormalization({\mdseries\slshape R})\index{RMSNormalization@\texttt{RMSNormalization}}
\label{RMSNormalization}
}\hfill{\scriptsize (attribute)}}\\
\textbf{\indent Returns:\ }
An isomorphism.



 If \mbox{\texttt{\mdseries\slshape R}} is a Rees matrix semigroup over a group \texttt{G} (i.e. a semigroup for which \texttt{IsGroupAsSemigroup} (\ref{IsGroupAsSemigroup}) returns \texttt{true}), then \texttt{RMSNormalization} returns an isomorphism from \mbox{\texttt{\mdseries\slshape R}} to a \emph{normalized} Rees matrix semigroup \texttt{S} over \texttt{G}. 

 The semigroup \texttt{S} is normalized in the sense that the first entry of each row and column of the \texttt{Matrix} (\textbf{Reference: Matrix}) of \texttt{S} is the identity element of \texttt{G}. 
\begin{Verbatim}[commandchars=!@|,fontsize=\small,frame=single,label=Example]
  !gapprompt@gap>| !gapinput@R := ReesMatrixSemigroup(SymmetricGroup(4),|
  !gapprompt@>| !gapinput@[[(1, 2), (2, 4, 3), (2, 1, 4)],|
  !gapprompt@>| !gapinput@ [(1, 3, 2), (1, 2)(3, 4), ()],|
  !gapprompt@>| !gapinput@ [(2, 3), (1, 3, 2, 4), (2, 3)]]);|
  <Rees matrix semigroup 3x3 over Sym( [ 1 .. 4 ] )>
  !gapprompt@gap>| !gapinput@iso := RMSNormalization(R);|
  <Rees matrix semigroup 3x3 over Sym( [ 1 .. 4 ] )> ->
  <Rees matrix semigroup 3x3 over Sym( [ 1 .. 4 ] )>
  !gapprompt@gap>| !gapinput@S := Range(iso);|
  <Rees matrix semigroup 3x3 over Sym( [ 1 .. 4 ] )>
  !gapprompt@gap>| !gapinput@Matrix(S);|
  [ [ (), (), () ], [ (), (1,2), (1,4,2,3) ], [ (), (1,4,2,3), (2,4) ] ]
\end{Verbatim}
 }

 

\subsection{\textcolor{Chapter }{IsomorphismReesMatrixSemigroup (for a semigroup)}}
\logpage{[ 6, 5, 8 ]}\nobreak
\hyperdef{L}{X7E2ECC577A1CF7CA}{}
{\noindent\textcolor{FuncColor}{$\triangleright$\enspace\texttt{IsomorphismReesMatrixSemigroup({\mdseries\slshape S})\index{IsomorphismReesMatrixSemigroup@\texttt{IsomorphismReesMatrixSemigroup}!for a semigroup}
\label{IsomorphismReesMatrixSemigroup:for a semigroup}
}\hfill{\scriptsize (attribute)}}\\
\noindent\textcolor{FuncColor}{$\triangleright$\enspace\texttt{IsomorphismReesZeroMatrixSemigroup({\mdseries\slshape S})\index{IsomorphismReesZeroMatrixSemigroup@\texttt{IsomorphismReesZeroMatrixSemigroup}}
\label{IsomorphismReesZeroMatrixSemigroup}
}\hfill{\scriptsize (attribute)}}\\
\noindent\textcolor{FuncColor}{$\triangleright$\enspace\texttt{IsomorphismReesMatrixSemigroupOverPermGroup({\mdseries\slshape S})\index{IsomorphismReesMatrixSemigroupOverPermGroup@\texttt{Isomorphism}\-\texttt{Rees}\-\texttt{Matrix}\-\texttt{Semigroup}\-\texttt{Over}\-\texttt{Perm}\-\texttt{Group}}
\label{IsomorphismReesMatrixSemigroupOverPermGroup}
}\hfill{\scriptsize (attribute)}}\\
\noindent\textcolor{FuncColor}{$\triangleright$\enspace\texttt{IsomorphismReesZeroMatrixSemigroupOverPermGroup({\mdseries\slshape S})\index{IsomorphismReesZeroMatrixSemigroupOverPermGroup@\texttt{Isomorphism}\-\texttt{Rees}\-\texttt{Zero}\-\texttt{Matrix}\-\texttt{Semigroup}\-\texttt{Over}\-\texttt{Perm}\-\texttt{Group}}
\label{IsomorphismReesZeroMatrixSemigroupOverPermGroup}
}\hfill{\scriptsize (attribute)}}\\
\textbf{\indent Returns:\ }
An isomorphism.



 If the semigroup \mbox{\texttt{\mdseries\slshape S}} is finite and simple, then \texttt{IsomorphismReesMatrixSemigroup} returns an isomorphism to a Rees matrix semigroup over some group (usually a
permutation group), and \texttt{IsomorphismReesMatrixSemigroupOverPermGroup} returns an isomorphism to a Rees matrix semigroup over a permutation group. 

 If \mbox{\texttt{\mdseries\slshape S}} is finite and 0\texttt{\symbol{45}}simple, then \texttt{IsomorphismReesZeroMatrixSemigroup} returns an isomorphism to a Rees 0\texttt{\symbol{45}}matrix semigroup over
some group (usually a permutation group), and \texttt{IsomorphismReesZeroMatrixSemigroupOverPermGroup} returns an isomorphism to a Rees 0\texttt{\symbol{45}}matrix semigroup over a
permutation group. 

 See also \texttt{InjectionPrincipalFactor} (\ref{InjectionPrincipalFactor}). 
\begin{Verbatim}[commandchars=!@|,fontsize=\small,frame=single,label=Example]
  !gapprompt@gap>| !gapinput@S := Semigroup(PartialPerm([1]));|
  <trivial partial perm group of rank 1 with 1 generator>
  !gapprompt@gap>| !gapinput@iso := IsomorphismReesMatrixSemigroup(S);;|
  !gapprompt@gap>| !gapinput@Source(iso) = S;|
  true
  !gapprompt@gap>| !gapinput@Range(iso);|
  <Rees matrix semigroup 1x1 over Group(())>
  !gapprompt@gap>| !gapinput@S := Semigroup(PartialPerm([1]), PartialPerm([]));|
  <partial perm monoid of rank 1 with 2 generators>
  !gapprompt@gap>| !gapinput@Range(IsomorphismReesZeroMatrixSemigroup(S));|
  <Rees 0-matrix semigroup 1x1 over Group(())>
\end{Verbatim}
 }

 

\subsection{\textcolor{Chapter }{AntiIsomorphismDualFpSemigroup}}
\logpage{[ 6, 5, 9 ]}\nobreak
\hyperdef{L}{X820BB66381737F2D}{}
{\noindent\textcolor{FuncColor}{$\triangleright$\enspace\texttt{AntiIsomorphismDualFpSemigroup({\mdseries\slshape S})\index{AntiIsomorphismDualFpSemigroup@\texttt{AntiIsomorphismDualFpSemigroup}}
\label{AntiIsomorphismDualFpSemigroup}
}\hfill{\scriptsize (attribute)}}\\
\noindent\textcolor{FuncColor}{$\triangleright$\enspace\texttt{AntiIsomorphismDualFpMonoid({\mdseries\slshape S})\index{AntiIsomorphismDualFpMonoid@\texttt{AntiIsomorphismDualFpMonoid}}
\label{AntiIsomorphismDualFpMonoid}
}\hfill{\scriptsize (attribute)}}\\
\textbf{\indent Returns:\ }
A finitely presented semigroup or monoid.



 \texttt{AntiIsomorphismDualFpSemigroup} returns an anti\texttt{\symbol{45}}isomorphism (\texttt{MappingByFunction} (\textbf{Reference: MappingByFunction})) from the finitely presented semigroup \mbox{\texttt{\mdseries\slshape S}} to another finitely presented semigroup. The range finitely presented
semigroup is obtained from \mbox{\texttt{\mdseries\slshape S}} by reversing the relations of \mbox{\texttt{\mdseries\slshape S}}.

 \texttt{AntiIsomorphismDualFpMonoid} works analogously when \mbox{\texttt{\mdseries\slshape S}} is a finitely presented monoid, and the range of the returned
anti\texttt{\symbol{45}}isomorphism is a finitely presented monoid. 
\begin{Verbatim}[commandchars=!@|,fontsize=\small,frame=single,label=Example]
  !gapprompt@gap>| !gapinput@F := FreeSemigroup("a", "b");|
  <free semigroup on the generators [ a, b ]>
  !gapprompt@gap>| !gapinput@AssignGeneratorVariables(F);|
  !gapprompt@gap>| !gapinput@R := [[a ^ 3, a], [b ^ 2, b], [(a * b) ^ 2, a]];|
  [ [ a^3, a ], [ b^2, b ], [ (a*b)^2, a ] ]
  !gapprompt@gap>| !gapinput@S := F / R;|
  <fp semigroup with 2 generators and 3 relations of length 14>
  !gapprompt@gap>| !gapinput@map := AntiIsomorphismDualFpSemigroup(S);|
  MappingByFunction( <fp semigroup with 2 generators and
    3 relations of length 14>, <fp semigroup with 2 generators and
    3 relations of length 14>
   , function( x ) ... end, function( x ) ... end )
  !gapprompt@gap>| !gapinput@RelationsOfFpSemigroup(Range(map));|
  [ [ a^3, a ], [ b^2, b ], [ (b*a)^2, a ] ]
\end{Verbatim}
 }

 

\subsection{\textcolor{Chapter }{EmbeddingFpMonoid}}
\logpage{[ 6, 5, 10 ]}\nobreak
\hyperdef{L}{X7873016586653A44}{}
{\noindent\textcolor{FuncColor}{$\triangleright$\enspace\texttt{EmbeddingFpMonoid({\mdseries\slshape S})\index{EmbeddingFpMonoid@\texttt{EmbeddingFpMonoid}}
\label{EmbeddingFpMonoid}
}\hfill{\scriptsize (attribute)}}\\
\textbf{\indent Returns:\ }
A finitely presented monoid.



 \texttt{EmbeddingFpMonoid} returns an embedding (\texttt{SemigroupHomomorphismByImages} (\ref{SemigroupHomomorphismByImages:for two semigroups and two
    lists})) from the finitely presented semigroup \mbox{\texttt{\mdseries\slshape S}} into a finitely presented monoid. If \mbox{\texttt{\mdseries\slshape S}} satisfies \texttt{IsMonoidAsSemigroup} (\ref{IsMonoidAsSemigroup}), then the mapping returned by this function is an isomorphism (the same
isomorphism as \texttt{IsomorphismFpMonoid} (\textbf{Reference: IsomorphismFpMonoid})). If \mbox{\texttt{\mdseries\slshape S}} is not a monoid, then the range is isomorphic to \mbox{\texttt{\mdseries\slshape S}} with an identity adjoined (a new element not previously in \mbox{\texttt{\mdseries\slshape S}}). The embedded copy of \mbox{\texttt{\mdseries\slshape S}} in the range can be recovered using \texttt{Image} (\textbf{Reference: Image}). 
\begin{Verbatim}[commandchars=!@|,fontsize=\small,frame=single,label=Example]
  !gapprompt@gap>| !gapinput@F := FreeSemigroup("a", "b");|
  <free semigroup on the generators [ a, b ]>
  !gapprompt@gap>| !gapinput@AssignGeneratorVariables(F);|
  !gapprompt@gap>| !gapinput@R := [[a ^ 3, a], [b ^ 2, b], [(a * b) ^ 2, a]];|
  [ [ a^3, a ], [ b^2, b ], [ (a*b)^2, a ] ]
  !gapprompt@gap>| !gapinput@S := F / R;|
  <fp semigroup with 2 generators and 3 relations of length 14>
  !gapprompt@gap>| !gapinput@Size(S);|
  3
  !gapprompt@gap>| !gapinput@IsMonoidAsSemigroup(S);|
  false
  !gapprompt@gap>| !gapinput@map := EmbeddingFpMonoid(S);|
  <fp semigroup with 2 generators and 3 relations of length 14> ->
  <fp monoid with 2 generators and 3 relations of length 14>
  !gapprompt@gap>| !gapinput@Size(Range(map));|
  4
\end{Verbatim}
 }

 }

   
\section{\textcolor{Chapter }{Random semigroups}}\logpage{[ 6, 6, 0 ]}
\hyperdef{L}{X7C3F130B8362D55A}{}
{
  

\subsection{\textcolor{Chapter }{RandomSemigroup}}
\logpage{[ 6, 6, 1 ]}\nobreak
\hyperdef{L}{X789DE9AB79FCFEB5}{}
{\noindent\textcolor{FuncColor}{$\triangleright$\enspace\texttt{RandomSemigroup({\mdseries\slshape arg...})\index{RandomSemigroup@\texttt{RandomSemigroup}}
\label{RandomSemigroup}
}\hfill{\scriptsize (function)}}\\
\noindent\textcolor{FuncColor}{$\triangleright$\enspace\texttt{RandomMonoid({\mdseries\slshape arg...})\index{RandomMonoid@\texttt{RandomMonoid}}
\label{RandomMonoid}
}\hfill{\scriptsize (function)}}\\
\noindent\textcolor{FuncColor}{$\triangleright$\enspace\texttt{RandomInverseSemigroup({\mdseries\slshape arg...})\index{RandomInverseSemigroup@\texttt{RandomInverseSemigroup}}
\label{RandomInverseSemigroup}
}\hfill{\scriptsize (function)}}\\
\noindent\textcolor{FuncColor}{$\triangleright$\enspace\texttt{RandomInverseMonoid({\mdseries\slshape arg...})\index{RandomInverseMonoid@\texttt{RandomInverseMonoid}}
\label{RandomInverseMonoid}
}\hfill{\scriptsize (function)}}\\
\textbf{\indent Returns:\ }
A semigroup.



 The operations described in this section can be used to generate semigroups,
in some sense, at random. There is no guarantee given about the distribution
of these semigroups, and this is only intended as a means of generating
semigroups for testing and other similar purposes. 

 Roughly speaking, the arguments of \texttt{RandomSemigroup} are a filter specifying the type of the semigroup to be returned, together
with some further parameters that describe some attributes of the semigroup to
be returned. For instance, we may want to specify the number of generators,
and, say, the degree, or dimension, of the elements, where appropriate. The
arguments of \texttt{RandomMonoid}, \texttt{RandomInverseSemigroup}, and \texttt{RandomInverseMonoid} are analogous. 

 If no arguments are specified, then they are all chosen at random, for a truly
random experience. 

 The first argument, if present, should be a filter \mbox{\texttt{\mdseries\slshape filter}}. For \texttt{RandomSemigroup} and \texttt{RandomInverseSemigroup} the filter \mbox{\texttt{\mdseries\slshape filter}} must be of the form \texttt{IsXSemigroup}. For example, \texttt{IsTransformationSemigroup} (\textbf{Reference: IsTransformationSemigroup}), \texttt{IsFpSemigroup} (\textbf{Reference: IsFpSemigroup}), and \texttt{IsPBRSemigroup} (\ref{IsPBRSemigroup}) are some possible values for \mbox{\texttt{\mdseries\slshape filter}}. For \texttt{RandomMonoid} and \texttt{RandomInverseMonoid} the argument \mbox{\texttt{\mdseries\slshape filter}} must be of the form \texttt{IsXMonoid}; such as \texttt{IsBipartitionMonoid} (\ref{IsBipartitionMonoid}) or \texttt{IsBooleanMatMonoid} (\ref{IsBooleanMatMonoid}). 

 Suppose that the first argument \mbox{\texttt{\mdseries\slshape filter}} is \texttt{IsFpSemigroup} (\textbf{Reference: IsFpSemigroup}). Then the only other arguments that can be specified is (and this argument is
also optional): 
\begin{description}
\item[{number of generators}]  The second argument, if present, should be a positive integer \mbox{\texttt{\mdseries\slshape m}} indicating the number of generators that the semigroup should have. If the
second argument \mbox{\texttt{\mdseries\slshape m}} is not specified, then a number is selected at random. 
\end{description}
 If \mbox{\texttt{\mdseries\slshape filter}} is a filter such as \texttt{IsTransformationSemigroup} (\textbf{Reference: IsTransformationSemigroup}) or \texttt{IsIntegerMatrixSemigroup} (\ref{IsIntegerMatrixSemigroup}),  then a further argument can be specified: 
\begin{description}
\item[{degree / dimension}]  The third argument, if present, should be a positive integer \mbox{\texttt{\mdseries\slshape n}}, which specifies the degree or dimension of the generators. For example, if
the first argument \mbox{\texttt{\mdseries\slshape filter}} is \texttt{IsTransformationSemigroup}, then the value of this argument is the degree of the transformations in the
returned semigroup; or if \mbox{\texttt{\mdseries\slshape filter}} is \texttt{IsMatrixOverFiniteFieldSemigroup}, then this argument is the dimension of the matrices in the returned
semigroup. 
\end{description}
 If \mbox{\texttt{\mdseries\slshape filter}} is \texttt{IsTropicalMaxPlusMatrixSemigroup} (\ref{IsTropicalMaxPlusMatrixSemigroup}), for example, then a fourth argument can be given (or not!): 
\begin{description}
\item[{threshold}]  The fourth argument, if present, should be a positive integer \mbox{\texttt{\mdseries\slshape t}}, which specifies the threshold of the semiring over which the matrices in the
returned semigroup are defined. 
\end{description}
 You get the idea, the error messages are self\texttt{\symbol{45}}explanatory,
and \texttt{RandomSemigroup} works for most of the type of semigroups defined in \textsf{GAP}.

 \texttt{RandomMonoid} is similar to \texttt{RandomSemigroup} except it returns a monoid. Likewise, \texttt{RandomInverseSemigroup} and \texttt{RandomInverseMonoid} return inverse semigroups and monoids, respectively. 
\begin{Verbatim}[commandchars=!@|,fontsize=\small,frame=single,label=Example]
  !gapprompt@gap>| !gapinput@RandomSemigroup();|
  <semigroup of 10x10 max-plus matrices with 12 generators>
  !gapprompt@gap>| !gapinput@RandomMonoid(IsTransformationMonoid);|
  <transformation monoid of degree 9 with 7 generators>
  !gapprompt@gap>| !gapinput@RandomMonoid(IsPartialPermMonoid, 2);|
  <partial perm monoid of rank 17 with 2 generators>
  !gapprompt@gap>| !gapinput@RandomMonoid(IsPartialPermMonoid, 2, 3);|
  <partial perm monoid of rank 3 with 2 generators>
  !gapprompt@gap>| !gapinput@RandomInverseSemigroup(IsTropicalMinPlusMatrixSemigroup);|
  <semigroup of 6x6 tropical min-plus matrices with 14 generators>
  !gapprompt@gap>| !gapinput@RandomInverseSemigroup(IsTropicalMinPlusMatrixSemigroup, 1);|
  <semigroup of 6x6 tropical min-plus matrices with 14 generators>
  !gapprompt@gap>| !gapinput@RandomSemigroup(IsTropicalMinPlusMatrixSemigroup, 2);|
  <semigroup of 11x11 tropical min-plus matrices with 2 generators>
  !gapprompt@gap>| !gapinput@RandomSemigroup(IsTropicalMinPlusMatrixSemigroup, 2, 1);|
  <semigroup of 1x1 tropical min-plus matrices with 2 generators>
  !gapprompt@gap>| !gapinput@RandomSemigroup(IsTropicalMinPlusMatrixSemigroup, 2, 1, 3);|
  !gapprompt@gap>| !gapinput@last.1;|
  Matrix(IsTropicalMinPlusMatrix, [[infinity]], 3)
  !gapprompt@gap>| !gapinput@RandomSemigroup(IsNTPMatrixSemigroup, 2, 1, 3, 4);|
  <semigroup of 1x1 ntp matrices with 2 generators>
  !gapprompt@gap>| !gapinput@last.1;|
  Matrix(IsNTPMatrix, [[2]], 3, 4)
  !gapprompt@gap>| !gapinput@RandomSemigroup(IsReesMatrixSemigroup, 2, 2);|
  <Rees matrix semigroup 2x2 over
    <permutation group of size 659 with 1 generator>>
  !gapprompt@gap>| !gapinput@RandomSemigroup(IsReesZeroMatrixSemigroup, 2, 2, Group((1, 2), (3, 4)));|
  <Rees 0-matrix semigroup 2x2 over Group([ (1,2), (3,4) ])>
  !gapprompt@gap>| !gapinput@RandomInverseMonoid(IsMatrixOverFiniteFieldMonoid, 2, 2);|
  <monoid of 3x3 matrices over GF(421^4) with 3 generators>
  !gapprompt@gap>| !gapinput@RandomInverseMonoid(IsMatrixOverFiniteFieldMonoid, 2, 2, GF(7));|
  <monoid of 3x3 matrices over GF(7) with 2 generators>
  !gapprompt@gap>| !gapinput@RandomSemigroup(IsBipartitionSemigroup, 5, 5);|
  <bipartition semigroup of degree 5 with 5 generators>
  !gapprompt@gap>| !gapinput@RandomMonoid(IsBipartitionMonoid, 5, 5);|
  <bipartition monoid of degree 5 with 5 generators>
  !gapprompt@gap>| !gapinput@RandomSemigroup(IsBooleanMatSemigroup);|
  <semigroup of 3x3 boolean matrices with 18 generators>
  !gapprompt@gap>| !gapinput@RandomMonoid(IsBooleanMatMonoid);|
  <monoid of 11x11 boolean matrices with 19 generators>
\end{Verbatim}
 }

 }

   }

 
\chapter{\textcolor{Chapter }{ Standard examples }}\label{Standard examples}
\logpage{[ 7, 0, 0 ]}
\hyperdef{L}{X7C76D1DC7DAF03D3}{}
{
  In this chapter we describe some standard families of examples of semigroups
and monoids which are available in the \textsf{Semigroups} package. 

   
\section{\textcolor{Chapter }{ Transformation semigroups }}\label{Transformation semigroups}
\logpage{[ 7, 1, 0 ]}
\hyperdef{L}{X7E42E8337A78B076}{}
{
  In this section, we describe the operations in \textsf{Semigroups} that can be used to create transformation semigroups belonging to several
standard classes of example. See  (\textbf{Reference: Transformations}) for more information about transformations. 

\subsection{\textcolor{Chapter }{CatalanMonoid}}
\logpage{[ 7, 1, 1 ]}\nobreak
\hyperdef{L}{X84C4C81380B0239D}{}
{\noindent\textcolor{FuncColor}{$\triangleright$\enspace\texttt{CatalanMonoid({\mdseries\slshape n})\index{CatalanMonoid@\texttt{CatalanMonoid}}
\label{CatalanMonoid}
}\hfill{\scriptsize (operation)}}\\
\textbf{\indent Returns:\ }
A transformation monoid.



 If \mbox{\texttt{\mdseries\slshape n}} is a positive integer, then this operation returns the Catalan monoid of
degree \mbox{\texttt{\mdseries\slshape n}}. The \emph{Catalan monoid} is the semigroup of the order\texttt{\symbol{45}}preserving and
order\texttt{\symbol{45}}decreasing transformations of \texttt{[1 .. n]} with the usual ordering. 

  The Catalan monoid is generated by the $\mbox{\texttt{\mdseries\slshape n - 1}}$ transformations $f_i$: 
\[ \left( \begin{array}{cccccccccc} 1&2&3&\cdots& i &i + 1& i + 2 & \cdots & n\\
1&2&3&\cdots& i &i & i + 2 & \cdots & n \end{array}\right), \qquad \]
 where $i = 1, \ldots, n - 1$ and has size equal to the $n$th Catalan number.  
\begin{Verbatim}[commandchars=!@|,fontsize=\small,frame=single,label=Example]
  !gapprompt@gap>| !gapinput@S := CatalanMonoid(6);|
  <transformation monoid of degree 6 with 5 generators>
  !gapprompt@gap>| !gapinput@Size(S);|
  132
\end{Verbatim}
 }

 

\subsection{\textcolor{Chapter }{EndomorphismsPartition}}
\logpage{[ 7, 1, 2 ]}\nobreak
\hyperdef{L}{X85C1D4307D0F5FF7}{}
{\noindent\textcolor{FuncColor}{$\triangleright$\enspace\texttt{EndomorphismsPartition({\mdseries\slshape list})\index{EndomorphismsPartition@\texttt{EndomorphismsPartition}}
\label{EndomorphismsPartition}
}\hfill{\scriptsize (operation)}}\\
\textbf{\indent Returns:\ }
A transformation monoid.



 If \mbox{\texttt{\mdseries\slshape list}} is a list of positive integers, then \texttt{EndomorphismsPartition} returns a monoid of endomorphisms preserving a partition of \texttt{[1 .. Sum(\mbox{\texttt{\mdseries\slshape list}})]} with a part of length \texttt{\mbox{\texttt{\mdseries\slshape list}}[i]} for every \texttt{i}. For example, if \texttt{\mbox{\texttt{\mdseries\slshape list}} = [1, 2, 3]}, then \texttt{EndomorphismsPartition} returns the monoid of endomorphisms of the partition \texttt{[[1], [2, 3], [4, 5, 6]]}. 

 If \texttt{f} is a transformation of \texttt{[1 .. n]}, then it is an \textsc{endomorphism} of a partition \texttt{P} on \texttt{[1 .. n]} if \texttt{(i, j)} in \texttt{P} implies that \texttt{(i \texttt{\symbol{94}} f, j \texttt{\symbol{94}} f)} is in \texttt{P}. 

 \texttt{EndomorphismsPartition} returns a monoid with a minimal size generating set, as described in \cite{Araujo2015aa}. 
\begin{Verbatim}[commandchars=!@|,fontsize=\small,frame=single,label=Example]
  !gapprompt@gap>| !gapinput@S := EndomorphismsPartition([3, 3, 3]);|
  <transformation semigroup of degree 9 with 4 generators>
  !gapprompt@gap>| !gapinput@Size(S);|
  531441
\end{Verbatim}
 }

 

\subsection{\textcolor{Chapter }{PartialTransformationMonoid}}
\logpage{[ 7, 1, 3 ]}\nobreak
\hyperdef{L}{X808A27F87E5AC598}{}
{\noindent\textcolor{FuncColor}{$\triangleright$\enspace\texttt{PartialTransformationMonoid({\mdseries\slshape n})\index{PartialTransformationMonoid@\texttt{PartialTransformationMonoid}}
\label{PartialTransformationMonoid}
}\hfill{\scriptsize (operation)}}\\
\textbf{\indent Returns:\ }
A transformation monoid.



 If \mbox{\texttt{\mdseries\slshape n}} is a positive integer, then this function returns a semigroup of
transformations on \texttt{\mbox{\texttt{\mdseries\slshape n}} + 1} points which is isomorphic to the semigroup consisting of all partial
transformation on \mbox{\texttt{\mdseries\slshape n}} points. This monoid has \texttt{(\mbox{\texttt{\mdseries\slshape n}} + 1) \texttt{\symbol{94}} \mbox{\texttt{\mdseries\slshape n}}} elements.  
\begin{Verbatim}[commandchars=!@|,fontsize=\small,frame=single,label=Example]
  !gapprompt@gap>| !gapinput@S := PartialTransformationMonoid(5);|
  <regular transformation monoid of degree 6 with 4 generators>
  !gapprompt@gap>| !gapinput@Size(S);|
  7776
\end{Verbatim}
 }

 

\subsection{\textcolor{Chapter }{SingularTransformationSemigroup}}
\logpage{[ 7, 1, 4 ]}\nobreak
\hyperdef{L}{X7894EE357D103806}{}
{\noindent\textcolor{FuncColor}{$\triangleright$\enspace\texttt{SingularTransformationSemigroup({\mdseries\slshape n})\index{SingularTransformationSemigroup@\texttt{SingularTransformationSemigroup}}
\label{SingularTransformationSemigroup}
}\hfill{\scriptsize (operation)}}\\
\noindent\textcolor{FuncColor}{$\triangleright$\enspace\texttt{SingularTransformationMonoid({\mdseries\slshape n})\index{SingularTransformationMonoid@\texttt{SingularTransformationMonoid}}
\label{SingularTransformationMonoid}
}\hfill{\scriptsize (operation)}}\\
\textbf{\indent Returns:\ }
The semigroup of non\texttt{\symbol{45}}invertible transformations.



 If \mbox{\texttt{\mdseries\slshape n}} is a integer greater than 1, then this function returns the semigroup of
non\texttt{\symbol{45}}invertible transformations, which is generated by the \texttt{\mbox{\texttt{\mdseries\slshape n}}(\mbox{\texttt{\mdseries\slshape n}} \texttt{\symbol{45}} 1)} idempotents of degree \mbox{\texttt{\mdseries\slshape n}} and rank \texttt{\mbox{\texttt{\mdseries\slshape n}} \texttt{\symbol{45}} 1} and has $n ^ n - n!$ elements. 
\begin{Verbatim}[commandchars=!@|,fontsize=\small,frame=single,label=Example]
  !gapprompt@gap>| !gapinput@S := SingularTransformationSemigroup(4);|
  <regular transformation semigroup ideal of degree 4 with 1 generator>
  !gapprompt@gap>| !gapinput@Size(S);|
  232
\end{Verbatim}
 }

 
\subsection{\textcolor{Chapter }{Semigroups of order\texttt{\symbol{45}}preserving transformations}}\logpage{[ 7, 1, 5 ]}
\hyperdef{L}{X80E80A0A83B57483}{}
{
\noindent\textcolor{FuncColor}{$\triangleright$\enspace\texttt{OrderEndomorphisms({\mdseries\slshape n})\index{OrderEndomorphisms@\texttt{OrderEndomorphisms}!monoid of order preserving transformations}
\label{OrderEndomorphisms:monoid of order preserving transformations}
}\hfill{\scriptsize (operation)}}\\
\noindent\textcolor{FuncColor}{$\triangleright$\enspace\texttt{SingularOrderEndomorphisms({\mdseries\slshape n})\index{SingularOrderEndomorphisms@\texttt{SingularOrderEndomorphisms}}
\label{SingularOrderEndomorphisms}
}\hfill{\scriptsize (operation)}}\\
\noindent\textcolor{FuncColor}{$\triangleright$\enspace\texttt{OrderAntiEndomorphisms({\mdseries\slshape n})\index{OrderAntiEndomorphisms@\texttt{OrderAntiEndomorphisms}}
\label{OrderAntiEndomorphisms}
}\hfill{\scriptsize (operation)}}\\
\noindent\textcolor{FuncColor}{$\triangleright$\enspace\texttt{PartialOrderEndomorphisms({\mdseries\slshape n})\index{PartialOrderEndomorphisms@\texttt{PartialOrderEndomorphisms}}
\label{PartialOrderEndomorphisms}
}\hfill{\scriptsize (operation)}}\\
\noindent\textcolor{FuncColor}{$\triangleright$\enspace\texttt{PartialOrderAntiEndomorphisms({\mdseries\slshape n})\index{PartialOrderAntiEndomorphisms@\texttt{PartialOrderAntiEndomorphisms}}
\label{PartialOrderAntiEndomorphisms}
}\hfill{\scriptsize (operation)}}\\
\textbf{\indent Returns:\ }
 A semigroup of transformations related to a linear order. 



 
\begin{description}
\item[{\texttt{OrderEndomorphisms(\mbox{\texttt{\mdseries\slshape n}})}}]  \texttt{OrderEndomorphisms(\mbox{\texttt{\mdseries\slshape n}})} returns the monoid of transformations that preserve the usual order on $\{1, 2, \ldots, n\}$, where \mbox{\texttt{\mdseries\slshape n}} is a positive integer.  \texttt{OrderEndomorphisms(\mbox{\texttt{\mdseries\slshape n}})} is generated by the $\mbox{\texttt{\mdseries\slshape n + 1}}$ transformations: 
\[ \left( \begin{array}{ccccccccc} 1&2&3&\cdots&n-1& n\\ 1&1&2&\cdots&n-2&n-1
\end{array}\right), \qquad \left( \begin{array}{ccccccccc} 1&2&\cdots&i-1& i&
i+1&i+2&\cdots &n\\ 1&2&\cdots&i-1& i+1&i+1&i+2&\cdots &n\\ \end{array}\right) \]
 where $i=0,\ldots,n-1$, and has ${2n-1\choose n-1}$ elements.  
\item[{\texttt{SingularOrderEndomorphisms(\mbox{\texttt{\mdseries\slshape n}})}}]  \texttt{SingularOrderEndomorphisms(\mbox{\texttt{\mdseries\slshape n}})} returns the ideal of \texttt{OrderEndomorphisms(\mbox{\texttt{\mdseries\slshape n}})} consisting of the non\texttt{\symbol{45}}invertible elements, when \mbox{\texttt{\mdseries\slshape n}} is at least \texttt{2}. The only invertible element in \texttt{OrderEndomorphisms(\mbox{\texttt{\mdseries\slshape n}})} is the identity transformation.  Therefore \texttt{SingularOrderEndomorphisms(\mbox{\texttt{\mdseries\slshape n}})} has ${2n-1\choose n-1} - 1$ elements.  
\item[{\texttt{OrderAntiEndomorphisms(\mbox{\texttt{\mdseries\slshape n}})}}]  \texttt{OrderAntiEndomorphisms(\mbox{\texttt{\mdseries\slshape n}})} returns the monoid of transformations that preserve or reverse the usual order
on $\{1, 2, \ldots, n\}$, where \mbox{\texttt{\mdseries\slshape n}} is a positive integer. \texttt{OrderAntiEndomorphisms(\mbox{\texttt{\mdseries\slshape n}})} is generated by the generators of \texttt{OrderEndomorphisms(\mbox{\texttt{\mdseries\slshape n}})} along with the bijective transformation that reverses the order on $\{1, 2, \ldots, n\}$.  The monoid \texttt{OrderAntiEndomorphisms(\mbox{\texttt{\mdseries\slshape n}})} has ${2n-1\choose n-1} - n$ elements.  
\item[{\texttt{PartialOrderEndomorphisms(\mbox{\texttt{\mdseries\slshape n}})}}]  \texttt{PartialOrderEndomorphisms(\mbox{\texttt{\mdseries\slshape n}})} returns a monoid of transformations on \texttt{\mbox{\texttt{\mdseries\slshape n}} + 1} points that is isomorphic to the monoid consisting of all partial
transformations that preserve the usual order on $\{1, 2, \ldots, n\}$. 
\item[{\texttt{PartialOrderAntiEndomorphisms(\mbox{\texttt{\mdseries\slshape n}})}}]  \texttt{PartialAntiOrderEndomorphisms(\mbox{\texttt{\mdseries\slshape n}})} returns a monoid of transformations on \texttt{\mbox{\texttt{\mdseries\slshape n}} + 1} points that is isomorphic to the monoid consisting of all partial
transformations that preserve or reverse the usual order on $\{1, 2, \ldots, n\}$. 
\end{description}
 
\begin{Verbatim}[commandchars=!@|,fontsize=\small,frame=single,label=Example]
  !gapprompt@gap>| !gapinput@S := OrderEndomorphisms(5);|
  <regular transformation monoid of degree 5 with 5 generators>
  !gapprompt@gap>| !gapinput@IsIdempotentGenerated(S);|
  true
  !gapprompt@gap>| !gapinput@Size(S) = Binomial(2 * 5 - 1, 5 - 1);|
  true
  !gapprompt@gap>| !gapinput@Difference(S, SingularOrderEndomorphisms(5));|
  [ IdentityTransformation ]
  !gapprompt@gap>| !gapinput@SingularOrderEndomorphisms(10);|
  <regular transformation semigroup ideal of degree 10 with 1 generator>
  !gapprompt@gap>| !gapinput@T := OrderAntiEndomorphisms(4);|
  <regular transformation monoid of degree 4 with 5 generators>
  !gapprompt@gap>| !gapinput@Transformation([4, 2, 2, 1]) in T;|
  true
  !gapprompt@gap>| !gapinput@U := PartialOrderEndomorphisms(6);|
  <regular transformation monoid of degree 7 with 12 generators>
  !gapprompt@gap>| !gapinput@V := PartialOrderAntiEndomorphisms(6);|
  <regular transformation monoid of degree 7 with 13 generators>
  !gapprompt@gap>| !gapinput@IsSubsemigroup(V, U);|
  true
\end{Verbatim}
 }

 

\subsection{\textcolor{Chapter }{EndomorphismMonoid (for a digraph)}}
\logpage{[ 7, 1, 6 ]}\nobreak
\hyperdef{L}{X868955247F2AFAA5}{}
{\noindent\textcolor{FuncColor}{$\triangleright$\enspace\texttt{EndomorphismMonoid({\mdseries\slshape digraph})\index{EndomorphismMonoid@\texttt{EndomorphismMonoid}!for a digraph}
\label{EndomorphismMonoid:for a digraph}
}\hfill{\scriptsize (attribute)}}\\
\noindent\textcolor{FuncColor}{$\triangleright$\enspace\texttt{EndomorphismMonoid({\mdseries\slshape digraph, colors})\index{EndomorphismMonoid@\texttt{EndomorphismMonoid}!for a digraph and vertex coloring}
\label{EndomorphismMonoid:for a digraph and vertex coloring}
}\hfill{\scriptsize (operation)}}\\
\textbf{\indent Returns:\ }
 A monoid.



 An endomorphism of \mbox{\texttt{\mdseries\slshape digraph}} is a homomorphism \texttt{DigraphHomomorphism} (\textbf{Digraphs: DigraphHomomorphism}) from \mbox{\texttt{\mdseries\slshape digraph}} back to itself.

 \texttt{EndomorphismMonoid}, called with a single argument, returns the monoid of all endomorphisms of \mbox{\texttt{\mdseries\slshape digraph}}. 

 If the \mbox{\texttt{\mdseries\slshape colors}} argument is specified, then it will return the monoid of endomorphisms which
respect the given colouring. The colouring \mbox{\texttt{\mdseries\slshape colors}} can be in one of two forms: 

 
\begin{itemize}
\item  A list of positive integers of size the number of vertices of \mbox{\texttt{\mdseries\slshape digraph}}, where \mbox{\texttt{\mdseries\slshape colors}}\texttt{[i]} is the colour of vertex \texttt{i}. 
\item  A list of lists, such that \mbox{\texttt{\mdseries\slshape colors}}\texttt{[i]} is a list of all vertices with colour \texttt{i}. 
\end{itemize}
 See also \texttt{GeneratorsOfEndomorphismMonoid} (\textbf{Digraphs: GeneratorsOfEndomorphismMonoid}). Note that the performance of \texttt{EndomorphismMonoid} may differ from that of \texttt{GeneratorsOfEndomorphismMonoid} (\textbf{Digraphs: GeneratorsOfEndomorphismMonoid}) since the former incrementally adds newly discovered endomorphisms to the
monoid using \texttt{ClosureMonoid} (\ref{ClosureMonoid}). 
\begin{Verbatim}[commandchars=!@|,fontsize=\small,frame=single,label=Example]
  !gapprompt@gap>| !gapinput@gr := Digraph(List([1 .. 3], x -> [1 .. 3]));;|
  !gapprompt@gap>| !gapinput@EndomorphismMonoid(gr);|
  <transformation monoid of degree 3 with 3 generators>
  !gapprompt@gap>| !gapinput@gr := CompleteDigraph(3);;|
  !gapprompt@gap>| !gapinput@EndomorphismMonoid(gr);|
  <transformation group of size 6, degree 3 with 2 generators>
  !gapprompt@gap>| !gapinput@S := EndomorphismMonoid(gr, [1, 2, 2]);;|
  !gapprompt@gap>| !gapinput@IsGroupAsSemigroup(S);|
  true
  !gapprompt@gap>| !gapinput@Size(S);|
  2
  !gapprompt@gap>| !gapinput@S := EndomorphismMonoid(gr, [[1], [2, 3]]);;|
  !gapprompt@gap>| !gapinput@S := EndomorphismMonoid(gr, [1, 2, 2]);;|
  !gapprompt@gap>| !gapinput@IsGroupAsSemigroup(S);|
  true
\end{Verbatim}
 }

 }

   
\section{\textcolor{Chapter }{ Semigroups of partial permutations }}\label{Semigroups of partial permutations}
\logpage{[ 7, 2, 0 ]}
\hyperdef{L}{X862BA1C67AA1C77C}{}
{
  In this section, we describe the operations in \textsf{Semigroups} that can be used to create semigroups of partial permutations belonging to
several standard classes of example. See  (\textbf{Reference: Partial permutations}) for more information about partial permutations. 

\subsection{\textcolor{Chapter }{MunnSemigroup}}
\logpage{[ 7, 2, 1 ]}\nobreak
\hyperdef{L}{X78FBE6DD7BCA30C1}{}
{\noindent\textcolor{FuncColor}{$\triangleright$\enspace\texttt{MunnSemigroup({\mdseries\slshape S})\index{MunnSemigroup@\texttt{MunnSemigroup}}
\label{MunnSemigroup}
}\hfill{\scriptsize (attribute)}}\\
\textbf{\indent Returns:\ }
The Munn semigroup of a semilattice.



 If \mbox{\texttt{\mdseries\slshape S}} is a semilattice, then \texttt{MunnSemigroup} returns the inverse semigroup of partial permutations of isomorphisms of
principal ideals of \mbox{\texttt{\mdseries\slshape S}}; called the \emph{Munn semigroup} of \mbox{\texttt{\mdseries\slshape S}}.

 This function was written jointly by J. D. Mitchell, Yann P{\a'e}resse (St
Andrews), Yanhui Wang (York). 
\begin{Verbatim}[commandchars=!@|,fontsize=\small,frame=single,label=Example]
  !gapprompt@gap>| !gapinput@S := InverseSemigroup([|
  !gapprompt@>| !gapinput@PartialPerm([1, 2, 3, 4, 5, 6, 7, 10], [4, 6, 7, 3, 8, 2, 9, 5]),|
  !gapprompt@>| !gapinput@PartialPerm([1, 2, 7, 9], [5, 6, 4, 3])]);|
  <inverse partial perm semigroup of rank 10 with 2 generators>
  !gapprompt@gap>| !gapinput@T := IdempotentGeneratedSubsemigroup(S);;|
  !gapprompt@gap>| !gapinput@M := MunnSemigroup(T);|
  <inverse partial perm semigroup of rank 60 with 7 generators>
  !gapprompt@gap>| !gapinput@NrIdempotents(M);|
  60
  !gapprompt@gap>| !gapinput@NrIdempotents(S);|
  60
\end{Verbatim}
 }

 

\subsection{\textcolor{Chapter }{RookMonoid}}
\logpage{[ 7, 2, 2 ]}\nobreak
\hyperdef{L}{X82D9619B7845CAEB}{}
{\noindent\textcolor{FuncColor}{$\triangleright$\enspace\texttt{RookMonoid({\mdseries\slshape n})\index{RookMonoid@\texttt{RookMonoid}}
\label{RookMonoid}
}\hfill{\scriptsize (operation)}}\\
\textbf{\indent Returns:\ }
An inverse monoid of partial permutations.



 \texttt{RookMonoid} is a synonym for \texttt{SymmetricInverseMonoid} (\textbf{Reference: SymmetricInverseMonoid}). 
\begin{Verbatim}[commandchars=!@|,fontsize=\small,frame=single,label=Example]
  !gapprompt@gap>| !gapinput@S := RookMonoid(4);|
  <symmetric inverse monoid of degree 4>
  !gapprompt@gap>| !gapinput@S = SymmetricInverseMonoid(4);|
  true
\end{Verbatim}
 }

 
\subsection{\textcolor{Chapter }{Inverse monoids of order\texttt{\symbol{45}}preserving partial permutations}}\logpage{[ 7, 2, 3 ]}
\hyperdef{L}{X85D841AE83DF101C}{}
{
\noindent\textcolor{FuncColor}{$\triangleright$\enspace\texttt{POI({\mdseries\slshape n})\index{POI@\texttt{POI}!monoid of order preserving partial perms}
\label{POI:monoid of order preserving partial perms}
}\hfill{\scriptsize (operation)}}\\
\noindent\textcolor{FuncColor}{$\triangleright$\enspace\texttt{PODI({\mdseries\slshape n})\index{PODI@\texttt{PODI}!monoid of order preserving or reversing partial perms}
\label{PODI:monoid of order preserving or reversing partial perms}
}\hfill{\scriptsize (operation)}}\\
\noindent\textcolor{FuncColor}{$\triangleright$\enspace\texttt{POPI({\mdseries\slshape n})\index{POPI@\texttt{POPI}!monoid of orientation preserving partial perms}
\label{POPI:monoid of orientation preserving partial perms}
}\hfill{\scriptsize (operation)}}\\
\noindent\textcolor{FuncColor}{$\triangleright$\enspace\texttt{PORI({\mdseries\slshape n})\index{PORI@\texttt{PORI}!monoid of orientation preserving or reversing partial perms}
\label{PORI:monoid of orientation preserving or reversing partial perms}
}\hfill{\scriptsize (operation)}}\\
\textbf{\indent Returns:\ }
 An inverse monoid of partial permutations related to a linear order. 



 
\begin{description}
\item[{\texttt{POI(\mbox{\texttt{\mdseries\slshape n}})}}]  \texttt{POI(\mbox{\texttt{\mdseries\slshape n}})} returns the inverse monoid of partial permutations that preserve the usual
order on $\{1, 2, \ldots, n\}$, where \mbox{\texttt{\mdseries\slshape n}} is a positive integer.  \texttt{POI(\mbox{\texttt{\mdseries\slshape n}})} is generated by the $\mbox{\texttt{\mdseries\slshape n}}$ partial permutations: 
\[ \left( \begin{array}{ccccc} 1&2&3&\cdots&n\\ -&1&2&\cdots&n-1
\end{array}\right), \qquad \left( \begin{array}{ccccccccc} 1&2&\cdots&i-1& i&
i+1&i+2&\cdots &n\\ 1&2&\cdots&i-1& i+1&-&i+2&\cdots&n\\ \end{array}\right) \]
 where $i = 1, \ldots, n - 1$, and has ${2n\choose n}$ elements.  
\item[{\texttt{PODI(\mbox{\texttt{\mdseries\slshape n}})}}]  \texttt{PODI(\mbox{\texttt{\mdseries\slshape n}})} returns the inverse monoid of partial permutations that preserve or reverse
the usual order on $\{1, 2, \ldots, n\}$, where \mbox{\texttt{\mdseries\slshape n}} is a positive integer. \texttt{PODI(\mbox{\texttt{\mdseries\slshape n}})} is generated by the generators of \texttt{POI(\mbox{\texttt{\mdseries\slshape n}})}, along with the permutation that reverses the usual order on $\{1, 2, \ldots, n\}$.  \texttt{PODI(\mbox{\texttt{\mdseries\slshape n}})} has ${2n\choose n} - n^{2} - 1$ elements.  
\item[{\texttt{POPI(\mbox{\texttt{\mdseries\slshape n}})}}]  \texttt{POPI(\mbox{\texttt{\mdseries\slshape n}})} returns the inverse monoid of partial permutations that preserve the
orientation of $\{1,2,\ldots, n\}$, where $n$ is a positive integer.  \texttt{POPI(\mbox{\texttt{\mdseries\slshape n}})} is generated by the partial permutations: 
\[ \left( \begin{array}{ccccc} 1&2&\cdots&n-1&n\\ 2&3&\cdots&n&1
\end{array}\right),\qquad \left( \begin{array}{cccccc} 1&2&\cdots&n-2&n-1&n\\
1&2&\cdots&n-2&n&- \end{array}\right), \]
 and has $1+\frac{n}{2}{2n\choose n}$ elements.  
\item[{\texttt{PORI(\mbox{\texttt{\mdseries\slshape n}})}}]  \texttt{PORI(\mbox{\texttt{\mdseries\slshape n}})} returns the inverse monoid of partial permutations that preserve or reverse
the orientation of $\{1, 2, \ldots, n\}$, where $n$ is a positive integer. \texttt{PORI(\mbox{\texttt{\mdseries\slshape n}})} is generated by the generators of \texttt{POPI(\mbox{\texttt{\mdseries\slshape n}})}, along with the permutation that reverses the usual order on $\{1, 2, \ldots, n\}$.  \texttt{PORI(\mbox{\texttt{\mdseries\slshape n}})} has $\frac{n}{2}{2n\choose n} - n(n + 1)$ elements.  
\end{description}
 
\begin{Verbatim}[commandchars=!@|,fontsize=\small,frame=single,label=Example]
  !gapprompt@gap>| !gapinput@S := PORI(10);|
  <inverse partial perm monoid of rank 10 with 3 generators>
  !gapprompt@gap>| !gapinput@S := POPI(10);|
  <inverse partial perm monoid of rank 10 with 2 generators>
  !gapprompt@gap>| !gapinput@Size(S) = 1 + 5 * Binomial(20, 10);|
  true
  !gapprompt@gap>| !gapinput@S := PODI(10);|
  <inverse partial perm monoid of rank 10 with 11 generators>
  !gapprompt@gap>| !gapinput@S := POI(10);|
  <inverse partial perm monoid of rank 10 with 10 generators>
  !gapprompt@gap>| !gapinput@Size(S) = Binomial(20, 10);|
  true
  !gapprompt@gap>| !gapinput@IsSubsemigroup(PORI(10), PODI(10))|
  !gapprompt@>| !gapinput@and IsSubsemigroup(PORI(10), POPI(10))|
  !gapprompt@>| !gapinput@and IsSubsemigroup(PODI(10), POI(10))|
  !gapprompt@>| !gapinput@and IsSubsemigroup(POPI(10), POI(10));|
  true
\end{Verbatim}
 }

 }

   
\section{\textcolor{Chapter }{ Semigroups of bipartitions }}\label{Semigroups of bipartitions}
\logpage{[ 7, 3, 0 ]}
\hyperdef{L}{X876C963F830719E2}{}
{
  In this section, we describe the operations in \textsf{Semigroups} that can be used to create bipartition semigroups belonging to several
standard classes of example. See Chapter \ref{Bipartitions and blocks} for more information about bipartitions. 

\subsection{\textcolor{Chapter }{PartitionMonoid}}
\logpage{[ 7, 3, 1 ]}\nobreak
\hyperdef{L}{X7E4B61FF7CCFD74A}{}
{\noindent\textcolor{FuncColor}{$\triangleright$\enspace\texttt{PartitionMonoid({\mdseries\slshape n})\index{PartitionMonoid@\texttt{PartitionMonoid}}
\label{PartitionMonoid}
}\hfill{\scriptsize (operation)}}\\
\noindent\textcolor{FuncColor}{$\triangleright$\enspace\texttt{RookPartitionMonoid({\mdseries\slshape n})\index{RookPartitionMonoid@\texttt{RookPartitionMonoid}}
\label{RookPartitionMonoid}
}\hfill{\scriptsize (operation)}}\\
\noindent\textcolor{FuncColor}{$\triangleright$\enspace\texttt{SingularPartitionMonoid({\mdseries\slshape n})\index{SingularPartitionMonoid@\texttt{SingularPartitionMonoid}}
\label{SingularPartitionMonoid}
}\hfill{\scriptsize (operation)}}\\
\textbf{\indent Returns:\ }
A bipartition monoid.



 If \mbox{\texttt{\mdseries\slshape n}} is a non\texttt{\symbol{45}}negative integer, then this operation returns the
partition monoid of degree \mbox{\texttt{\mdseries\slshape n}}. The \emph{partition monoid of degree \mbox{\texttt{\mdseries\slshape n}}} is the monoid consisting of all the bipartitions of degree \mbox{\texttt{\mdseries\slshape n}}. 

 \texttt{SingularPartitionMonoid} returns the ideal of the partition monoid consisting of the
non\texttt{\symbol{45}}invertible elements (i.e. those not in the group of
units), when \mbox{\texttt{\mdseries\slshape n}} is positive. 

 If \mbox{\texttt{\mdseries\slshape n}} is positive, then \texttt{RookPartitionMonoid} returns submonoid of the partition monoid of degree \texttt{\mbox{\texttt{\mdseries\slshape n}} + 1} consisting of those bipartitions with \texttt{\mbox{\texttt{\mdseries\slshape n}} + 1} and \texttt{\texttt{\symbol{45}}\mbox{\texttt{\mdseries\slshape n}} \texttt{\symbol{45}} 1} in the same block; see \cite{Halverson2005PartitionAlgebras}, \cite{Grood2006aa}, and \cite{East2019aa}. 
\begin{Verbatim}[commandchars=!@|,fontsize=\small,frame=single,label=Example]
  !gapprompt@gap>| !gapinput@S := PartitionMonoid(4);|
  <regular bipartition *-monoid of size 4140, degree 4 with 4
   generators>
  !gapprompt@gap>| !gapinput@Size(S);|
  4140
  !gapprompt@gap>| !gapinput@T := SingularPartitionMonoid(4);|
  <regular bipartition *-semigroup ideal of degree 4 with 1 generator>
  !gapprompt@gap>| !gapinput@Size(S) - Size(T) = Factorial(4);|
  true
  !gapprompt@gap>| !gapinput@S := RookPartitionMonoid(4);|
  <regular bipartition *-monoid of degree 5 with 5 generators>
  !gapprompt@gap>| !gapinput@Size(S);|
  21147
\end{Verbatim}
 }

 

\subsection{\textcolor{Chapter }{BrauerMonoid}}
\logpage{[ 7, 3, 2 ]}\nobreak
\hyperdef{L}{X79D33B2E7BA3073A}{}
{\noindent\textcolor{FuncColor}{$\triangleright$\enspace\texttt{BrauerMonoid({\mdseries\slshape n})\index{BrauerMonoid@\texttt{BrauerMonoid}}
\label{BrauerMonoid}
}\hfill{\scriptsize (operation)}}\\
\noindent\textcolor{FuncColor}{$\triangleright$\enspace\texttt{PartialBrauerMonoid({\mdseries\slshape n})\index{PartialBrauerMonoid@\texttt{PartialBrauerMonoid}}
\label{PartialBrauerMonoid}
}\hfill{\scriptsize (operation)}}\\
\noindent\textcolor{FuncColor}{$\triangleright$\enspace\texttt{SingularBrauerMonoid({\mdseries\slshape n})\index{SingularBrauerMonoid@\texttt{SingularBrauerMonoid}}
\label{SingularBrauerMonoid}
}\hfill{\scriptsize (operation)}}\\
\textbf{\indent Returns:\ }
A bipartition monoid.



 If \mbox{\texttt{\mdseries\slshape n}} is a non\texttt{\symbol{45}}negative integer, then this operation returns the
Brauer monoid of degree \mbox{\texttt{\mdseries\slshape n}}. The \emph{Brauer monoid} is the submonoid of the partition monoid consisting of those bipartitions
where the size of every block is 2. 

 \texttt{PartialBrauerMonoid} returns the partial Brauer monoid, which is the submonoid of the partition
monoid consisting of those bipartitions where the size of every block is \emph{at most} 2. The partial Brauer monoid contains the Brauer monoid as a submonoid. 

 \texttt{SingularBrauerMonoid} returns the ideal of the Brauer monoid consisting of the
non\texttt{\symbol{45}}invertible elements (i.e. those not in the group of
units), when \mbox{\texttt{\mdseries\slshape n}} is at least 2. 
\begin{Verbatim}[commandchars=!@|,fontsize=\small,frame=single,label=Example]
  !gapprompt@gap>| !gapinput@S := BrauerMonoid(4);|
  <regular bipartition *-monoid of degree 4 with 3 generators>
  !gapprompt@gap>| !gapinput@IsSubsemigroup(S, JonesMonoid(4));|
  true
  !gapprompt@gap>| !gapinput@Size(S);|
  105
  !gapprompt@gap>| !gapinput@SingularBrauerMonoid(8);|
  <regular bipartition *-semigroup ideal of degree 8 with 1 generator>
  !gapprompt@gap>| !gapinput@S := PartialBrauerMonoid(3);|
  <regular bipartition *-monoid of degree 3 with 8 generators>
  !gapprompt@gap>| !gapinput@IsSubsemigroup(S, BrauerMonoid(3));|
  true
  !gapprompt@gap>| !gapinput@Size(S);|
  76
\end{Verbatim}
 }

 

\subsection{\textcolor{Chapter }{JonesMonoid}}
\logpage{[ 7, 3, 3 ]}\nobreak
\hyperdef{L}{X8378FC8B840B9706}{}
{\noindent\textcolor{FuncColor}{$\triangleright$\enspace\texttt{JonesMonoid({\mdseries\slshape n})\index{JonesMonoid@\texttt{JonesMonoid}}
\label{JonesMonoid}
}\hfill{\scriptsize (operation)}}\\
\noindent\textcolor{FuncColor}{$\triangleright$\enspace\texttt{TemperleyLiebMonoid({\mdseries\slshape n})\index{TemperleyLiebMonoid@\texttt{TemperleyLiebMonoid}}
\label{TemperleyLiebMonoid}
}\hfill{\scriptsize (operation)}}\\
\noindent\textcolor{FuncColor}{$\triangleright$\enspace\texttt{SingularJonesMonoid({\mdseries\slshape n})\index{SingularJonesMonoid@\texttt{SingularJonesMonoid}}
\label{SingularJonesMonoid}
}\hfill{\scriptsize (operation)}}\\
\textbf{\indent Returns:\ }
A bipartition monoid.



 If \mbox{\texttt{\mdseries\slshape n}} is a non\texttt{\symbol{45}}negative integer, then this operation returns the
Jones monoid of degree \mbox{\texttt{\mdseries\slshape n}}. The \emph{Jones monoid} is the subsemigroup of the Brauer monoid consisting of those bipartitions that
are planar; see \texttt{PlanarPartitionMonoid} (\ref{PlanarPartitionMonoid}). The Jones monoid is sometimes referred to as the \textsc{Temperley\texttt{\symbol{45}}Lieb monoid}. 

 \texttt{SingularJonesMonoid} returns the ideal of the Jones monoid consisting of the
non\texttt{\symbol{45}}invertible elements (i.e. those not in the group of
units), when \mbox{\texttt{\mdseries\slshape n}} is at least 2. 
\begin{Verbatim}[commandchars=!@|,fontsize=\small,frame=single,label=Example]
  !gapprompt@gap>| !gapinput@S := JonesMonoid(4);|
  <regular bipartition *-monoid of degree 4 with 3 generators>
  !gapprompt@gap>| !gapinput@S = TemperleyLiebMonoid(4);|
  true
  !gapprompt@gap>| !gapinput@SingularJonesMonoid(8);|
  <regular bipartition *-semigroup ideal of degree 8 with 1 generator>
\end{Verbatim}
 }

 

\subsection{\textcolor{Chapter }{PartialJonesMonoid}}
\logpage{[ 7, 3, 4 ]}\nobreak
\hyperdef{L}{X8458B0F7874484CE}{}
{\noindent\textcolor{FuncColor}{$\triangleright$\enspace\texttt{PartialJonesMonoid({\mdseries\slshape n})\index{PartialJonesMonoid@\texttt{PartialJonesMonoid}}
\label{PartialJonesMonoid}
}\hfill{\scriptsize (operation)}}\\
\textbf{\indent Returns:\ }
A bipartition monoid.



 If \mbox{\texttt{\mdseries\slshape n}} is a non\texttt{\symbol{45}}negative integer, then \texttt{PartialJonesMonoid} returns the partial Jones monoid of degree \mbox{\texttt{\mdseries\slshape n}}. The \emph{partial Jones monoid} is a subsemigroup of the partial Brauer monoid. An element of the partial
Brauer monoid is contained in the partial Jones monoid if the partition that
it defines is finer than the partition defined by some element of the Jones
monoid. In other words, an element of the partial Jones monoid can be formed
from some element \texttt{x} of the Jones monoid by replacing some blocks \texttt{[a, b]} of \texttt{x} by singleton blocks \texttt{[a], [b]}. 

 Note that, in general, the partial Jones monoid of degree \mbox{\texttt{\mdseries\slshape n}} is strictly contained in the Motzkin monoid of the same degree. 

 See \texttt{PartialBrauerMonoid} (\ref{PartialBrauerMonoid}), \texttt{JonesMonoid} (\ref{JonesMonoid}), and \texttt{MotzkinMonoid} (\ref{MotzkinMonoid}). 
\begin{Verbatim}[commandchars=!@|,fontsize=\small,frame=single,label=Example]
  !gapprompt@gap>| !gapinput@S := PartialJonesMonoid(4);|
  <regular bipartition *-monoid of degree 4 with 7 generators>
  !gapprompt@gap>| !gapinput@T := JonesMonoid(4);|
  <regular bipartition *-monoid of degree 4 with 3 generators>
  !gapprompt@gap>| !gapinput@U := MotzkinMonoid(4);|
  <regular bipartition *-monoid of degree 4 with 8 generators>
  !gapprompt@gap>| !gapinput@IsSubsemigroup(U, S);|
  true
  !gapprompt@gap>| !gapinput@IsSubsemigroup(S, T);|
  true
  !gapprompt@gap>| !gapinput@Size(U);|
  323
  !gapprompt@gap>| !gapinput@Size(S);|
  143
  !gapprompt@gap>| !gapinput@Size(T);|
  14
\end{Verbatim}
 }

 

\subsection{\textcolor{Chapter }{AnnularJonesMonoid}}
\logpage{[ 7, 3, 5 ]}\nobreak
\hyperdef{L}{X7DB8CB067CBE1254}{}
{\noindent\textcolor{FuncColor}{$\triangleright$\enspace\texttt{AnnularJonesMonoid({\mdseries\slshape n})\index{AnnularJonesMonoid@\texttt{AnnularJonesMonoid}}
\label{AnnularJonesMonoid}
}\hfill{\scriptsize (operation)}}\\
\textbf{\indent Returns:\ }
A bipartition monoid.



 If \mbox{\texttt{\mdseries\slshape n}} is a non\texttt{\symbol{45}}negative integer, then \texttt{AnnularJonesMonoid} returns the annular Jones monoid of degree \mbox{\texttt{\mdseries\slshape n}}. The \emph{annular Jones monoid} is the subsemigroup of the partition monoid consisting of all annular
bipartitions whose blocks have size 2 (annular bipartitions are defined in
Chapter \ref{Bipartitions and blocks}). See \texttt{BrauerMonoid} (\ref{BrauerMonoid}). 
\begin{Verbatim}[commandchars=!@|,fontsize=\small,frame=single,label=Example]
  !gapprompt@gap>| !gapinput@S := AnnularJonesMonoid(4);|
  <regular bipartition *-monoid of degree 4 with 2 generators>
\end{Verbatim}
 }

 

\subsection{\textcolor{Chapter }{MotzkinMonoid}}
\logpage{[ 7, 3, 6 ]}\nobreak
\hyperdef{L}{X8375152F7AB52B7B}{}
{\noindent\textcolor{FuncColor}{$\triangleright$\enspace\texttt{MotzkinMonoid({\mdseries\slshape n})\index{MotzkinMonoid@\texttt{MotzkinMonoid}}
\label{MotzkinMonoid}
}\hfill{\scriptsize (operation)}}\\
\textbf{\indent Returns:\ }
A bipartition monoid.



 If \mbox{\texttt{\mdseries\slshape n}} is a non\texttt{\symbol{45}}negative integer, then this operation returns the
Motzkin monoid of degree \mbox{\texttt{\mdseries\slshape n}}. The \emph{Motzkin monoid} is the subsemigroup of the partial Brauer monoid consisting of those
bipartitions that are planar (planar bipartitions are defined in Chapter \ref{Bipartitions and blocks}). 

 Note that the Motzkin monoid of degree \mbox{\texttt{\mdseries\slshape n}} contains the partial Jones monoid of degree \mbox{\texttt{\mdseries\slshape n}}, but in general, these monoids are not equal; see \texttt{PartialJonesMonoid} (\ref{PartialJonesMonoid}). 
\begin{Verbatim}[commandchars=!@|,fontsize=\small,frame=single,label=Example]
  !gapprompt@gap>| !gapinput@S := MotzkinMonoid(4);|
  <regular bipartition *-monoid of degree 4 with 8 generators>
  !gapprompt@gap>| !gapinput@T := PartialJonesMonoid(4);|
  <regular bipartition *-monoid of degree 4 with 7 generators>
  !gapprompt@gap>| !gapinput@IsSubsemigroup(S, T);|
  true
  !gapprompt@gap>| !gapinput@Size(S);|
  323
  !gapprompt@gap>| !gapinput@Size(T);|
  143
\end{Verbatim}
 }

 

\subsection{\textcolor{Chapter }{DualSymmetricInverseSemigroup}}
\logpage{[ 7, 3, 7 ]}\nobreak
\hyperdef{L}{X83C7587C81B985BA}{}
{\noindent\textcolor{FuncColor}{$\triangleright$\enspace\texttt{DualSymmetricInverseSemigroup({\mdseries\slshape n})\index{DualSymmetricInverseSemigroup@\texttt{DualSymmetricInverseSemigroup}}
\label{DualSymmetricInverseSemigroup}
}\hfill{\scriptsize (operation)}}\\
\noindent\textcolor{FuncColor}{$\triangleright$\enspace\texttt{DualSymmetricInverseMonoid({\mdseries\slshape n})\index{DualSymmetricInverseMonoid@\texttt{DualSymmetricInverseMonoid}}
\label{DualSymmetricInverseMonoid}
}\hfill{\scriptsize (operation)}}\\
\noindent\textcolor{FuncColor}{$\triangleright$\enspace\texttt{SingularDualSymmetricInverseMonoid({\mdseries\slshape n})\index{SingularDualSymmetricInverseMonoid@\texttt{SingularDualSymmetricInverseMonoid}}
\label{SingularDualSymmetricInverseMonoid}
}\hfill{\scriptsize (operation)}}\\
\noindent\textcolor{FuncColor}{$\triangleright$\enspace\texttt{PartialDualSymmetricInverseMonoid({\mdseries\slshape n})\index{PartialDualSymmetricInverseMonoid@\texttt{PartialDualSymmetricInverseMonoid}}
\label{PartialDualSymmetricInverseMonoid}
}\hfill{\scriptsize (operation)}}\\
\textbf{\indent Returns:\ }
An inverse bipartition monoid.



 If \mbox{\texttt{\mdseries\slshape n}} is a positive integer, then the operations \texttt{DualSymmetricInverseSemigroup} and \texttt{DualSymmetricInverseMonoid} return the dual symmetric inverse monoid of degree \mbox{\texttt{\mdseries\slshape n}}, which is the subsemigroup of the partition monoid consisting of the block
bijections of degree \mbox{\texttt{\mdseries\slshape n}}.

 \texttt{SingularDualSymmetricInverseMonoid} returns the ideal of the dual symmetric inverse monoid consisting of the
non\texttt{\symbol{45}}invertible elements (i.e. those not in the group of
units), when \mbox{\texttt{\mdseries\slshape n}} is at least 2.

 \texttt{PartialDualSymmetricInverseMonoid} returns the submonoid of the dual symmetric inverse monoid of degree \texttt{\mbox{\texttt{\mdseries\slshape n}} + 1} consisting of those block bijections with \texttt{\mbox{\texttt{\mdseries\slshape n}} + 1} and \texttt{\texttt{\symbol{45}}\mbox{\texttt{\mdseries\slshape n}} \texttt{\symbol{45}} 1} in the same block; see \cite{Kudryavtseva2011aa} and \cite{Kudryavtseva2015aa}. 

 See \texttt{IsBlockBijection} (\ref{IsBlockBijection}). 
\begin{Verbatim}[commandchars=!@|,fontsize=\small,frame=single,label=Example]
  !gapprompt@gap>| !gapinput@Number(PartitionMonoid(3), IsBlockBijection);|
  25
  !gapprompt@gap>| !gapinput@S := DualSymmetricInverseSemigroup(3);|
  <inverse block bijection monoid of degree 3 with 3 generators>
  !gapprompt@gap>| !gapinput@Size(S);|
  25
  !gapprompt@gap>| !gapinput@S := PartialDualSymmetricInverseMonoid(5);|
  <inverse block bijection monoid of degree 6 with 4 generators>
  !gapprompt@gap>| !gapinput@Size(S);|
  29072
\end{Verbatim}
 }

 

\subsection{\textcolor{Chapter }{UniformBlockBijectionMonoid}}
\logpage{[ 7, 3, 8 ]}\nobreak
\hyperdef{L}{X8301C61384168D6F}{}
{\noindent\textcolor{FuncColor}{$\triangleright$\enspace\texttt{UniformBlockBijectionMonoid({\mdseries\slshape n})\index{UniformBlockBijectionMonoid@\texttt{UniformBlockBijectionMonoid}}
\label{UniformBlockBijectionMonoid}
}\hfill{\scriptsize (operation)}}\\
\noindent\textcolor{FuncColor}{$\triangleright$\enspace\texttt{FactorisableDualSymmetricInverseMonoid({\mdseries\slshape n})\index{FactorisableDualSymmetricInverseMonoid@\texttt{Factorisable}\-\texttt{Dual}\-\texttt{Symmetric}\-\texttt{Inverse}\-\texttt{Monoid}}
\label{FactorisableDualSymmetricInverseMonoid}
}\hfill{\scriptsize (operation)}}\\
\noindent\textcolor{FuncColor}{$\triangleright$\enspace\texttt{SingularUniformBlockBijectionMonoid({\mdseries\slshape n})\index{SingularUniformBlockBijectionMonoid@\texttt{SingularUniformBlockBijectionMonoid}}
\label{SingularUniformBlockBijectionMonoid}
}\hfill{\scriptsize (operation)}}\\
\noindent\textcolor{FuncColor}{$\triangleright$\enspace\texttt{PartialUniformBlockBijectionMonoid({\mdseries\slshape n})\index{PartialUniformBlockBijectionMonoid@\texttt{PartialUniformBlockBijectionMonoid}}
\label{PartialUniformBlockBijectionMonoid}
}\hfill{\scriptsize (operation)}}\\
\noindent\textcolor{FuncColor}{$\triangleright$\enspace\texttt{SingularFactorisableDualSymmetricInverseMonoid({\mdseries\slshape n})\index{SingularFactorisableDualSymmetricInverseMonoid@\texttt{Singular}\-\texttt{Factorisable}\-\texttt{Dual}\-\texttt{Symmetric}\-\texttt{Inverse}\-\texttt{Monoid}}
\label{SingularFactorisableDualSymmetricInverseMonoid}
}\hfill{\scriptsize (operation)}}\\
\noindent\textcolor{FuncColor}{$\triangleright$\enspace\texttt{PlanarUniformBlockBijectionMonoid({\mdseries\slshape n})\index{PlanarUniformBlockBijectionMonoid@\texttt{PlanarUniformBlockBijectionMonoid}}
\label{PlanarUniformBlockBijectionMonoid}
}\hfill{\scriptsize (operation)}}\\
\noindent\textcolor{FuncColor}{$\triangleright$\enspace\texttt{SingularPlanarUniformBlockBijectionMonoid({\mdseries\slshape n})\index{SingularPlanarUniformBlockBijectionMonoid@\texttt{Singular}\-\texttt{Planar}\-\texttt{Uniform}\-\texttt{Block}\-\texttt{Bijection}\-\texttt{Monoid}}
\label{SingularPlanarUniformBlockBijectionMonoid}
}\hfill{\scriptsize (operation)}}\\
\textbf{\indent Returns:\ }
An inverse bipartition monoid.



 If \mbox{\texttt{\mdseries\slshape n}} is a positive integer, then this operation returns the uniform block bijection
monoid of degree \mbox{\texttt{\mdseries\slshape n}}. The \emph{uniform block bijection monoid} is the submonoid of the partition monoid consisting of the block bijections of
degree $n$  where the number of positive integers in a block equals the number of negative
integers in that block. The uniform block bijection monoid is also referred to
as the \emph{factorisable dual symmetric inverse monoid}. 

 \texttt{SingularUniformBlockBijectionMonoid} returns the ideal of the uniform block bijection monoid consisting of the
non\texttt{\symbol{45}}invertible elements (i.e. those not in the group of
units), when \mbox{\texttt{\mdseries\slshape n}} is at least 2. 

 \texttt{PlanarUniformBlockBijectionMonoid} returns the submonoid of the uniform block bijection monoid consisting of the
planar elements (i.e. those in the planar partition monoid, see \texttt{PlanarPartitionMonoid} (\ref{PlanarPartitionMonoid})). 

 \texttt{SingularPlanarUniformBlockBijectionMonoid} returns the ideal of the planar uniform block bijection monoid consisting of
the non\texttt{\symbol{45}}invertible elements (i.e. those not in the group of
units), when \mbox{\texttt{\mdseries\slshape n}} is at least 2. 

  \texttt{PartialUniformBlockBijectionMonoid} returns the submonoid of the uniform block bijection monoid of degree \texttt{\mbox{\texttt{\mdseries\slshape n}} + 1} consisting of those uniform block bijection with \texttt{\mbox{\texttt{\mdseries\slshape n}} + 1} and \texttt{\texttt{\symbol{45}}\mbox{\texttt{\mdseries\slshape n}} \texttt{\symbol{45}} 1} in the same block. 

 See \texttt{IsUniformBlockBijection} (\ref{IsUniformBlockBijection}). 
\begin{Verbatim}[commandchars=!@|,fontsize=\small,frame=single,label=Example]
  !gapprompt@gap>| !gapinput@S := UniformBlockBijectionMonoid(4);|
  <inverse block bijection monoid of degree 4 with 3 generators>
  !gapprompt@gap>| !gapinput@Size(PlanarUniformBlockBijectionMonoid(8));|
  128
  !gapprompt@gap>| !gapinput@S := DualSymmetricInverseMonoid(4);|
  <inverse block bijection monoid of degree 4 with 3 generators>
  !gapprompt@gap>| !gapinput@IsFactorisableInverseMonoid(S);|
  false
  !gapprompt@gap>| !gapinput@S := UniformBlockBijectionMonoid(4);|
  <inverse block bijection monoid of degree 4 with 3 generators>
  !gapprompt@gap>| !gapinput@IsFactorisableInverseMonoid(S);|
  true
  !gapprompt@gap>| !gapinput@S := AsSemigroup(IsBipartitionSemigroup,|
  !gapprompt@>| !gapinput@                    SymmetricInverseMonoid(5));|
  <inverse bipartition monoid of degree 5 with 3 generators>
  !gapprompt@gap>| !gapinput@IsFactorisableInverseMonoid(S);|
  true
  !gapprompt@gap>| !gapinput@S := PartialUniformBlockBijectionMonoid(5);|
  <inverse block bijection monoid of degree 6 with 4 generators>
  !gapprompt@gap>| !gapinput@NrIdempotents(S);|
  203
  !gapprompt@gap>| !gapinput@IsFactorisableInverseMonoid(S);|
  true
\end{Verbatim}
 }

 

\subsection{\textcolor{Chapter }{PlanarPartitionMonoid}}
\logpage{[ 7, 3, 9 ]}\nobreak
\hyperdef{L}{X8444092A7967A029}{}
{\noindent\textcolor{FuncColor}{$\triangleright$\enspace\texttt{PlanarPartitionMonoid({\mdseries\slshape n})\index{PlanarPartitionMonoid@\texttt{PlanarPartitionMonoid}}
\label{PlanarPartitionMonoid}
}\hfill{\scriptsize (operation)}}\\
\noindent\textcolor{FuncColor}{$\triangleright$\enspace\texttt{SingularPlanarPartitionMonoid({\mdseries\slshape n})\index{SingularPlanarPartitionMonoid@\texttt{SingularPlanarPartitionMonoid}}
\label{SingularPlanarPartitionMonoid}
}\hfill{\scriptsize (operation)}}\\
\textbf{\indent Returns:\ }
A bipartition monoid.



 If \mbox{\texttt{\mdseries\slshape n}} is a positive integer, then this operation returns the planar partition monoid
of degree \mbox{\texttt{\mdseries\slshape n}} which is the monoid consisting of all the planar bipartitions of degree \mbox{\texttt{\mdseries\slshape n}} (planar bipartitions are defined in Chapter \ref{Bipartitions and blocks}). 

 \texttt{SingularPlanarPartitionMonoid} returns the ideal of the planar partition monoid consisting of the
non\texttt{\symbol{45}}invertible elements (i.e. those not in the group of
units). 

 
\begin{Verbatim}[commandchars=!@|,fontsize=\small,frame=single,label=Example]
  !gapprompt@gap>| !gapinput@S := PlanarPartitionMonoid(3);|
  <regular bipartition *-monoid of degree 3 with 5 generators>
  !gapprompt@gap>| !gapinput@Size(S);|
  132
  !gapprompt@gap>| !gapinput@T := SingularPlanarPartitionMonoid(3);|
  <regular bipartition *-semigroup ideal of degree 3 with 1 generator>
  !gapprompt@gap>| !gapinput@Size(T);|
  131
  !gapprompt@gap>| !gapinput@Difference(S, T);|
  [ <block bijection: [ 1, -1 ], [ 2, -2 ], [ 3, -3 ]> ]
\end{Verbatim}
 }

 

\subsection{\textcolor{Chapter }{ModularPartitionMonoid}}
\logpage{[ 7, 3, 10 ]}\nobreak
\hyperdef{L}{X7F208DC584C0B9D1}{}
{\noindent\textcolor{FuncColor}{$\triangleright$\enspace\texttt{ModularPartitionMonoid({\mdseries\slshape m, n})\index{ModularPartitionMonoid@\texttt{ModularPartitionMonoid}}
\label{ModularPartitionMonoid}
}\hfill{\scriptsize (operation)}}\\
\noindent\textcolor{FuncColor}{$\triangleright$\enspace\texttt{SingularModularPartitionMonoid({\mdseries\slshape m, n})\index{SingularModularPartitionMonoid@\texttt{SingularModularPartitionMonoid}}
\label{SingularModularPartitionMonoid}
}\hfill{\scriptsize (operation)}}\\
\noindent\textcolor{FuncColor}{$\triangleright$\enspace\texttt{PlanarModularPartitionMonoid({\mdseries\slshape m, n})\index{PlanarModularPartitionMonoid@\texttt{PlanarModularPartitionMonoid}}
\label{PlanarModularPartitionMonoid}
}\hfill{\scriptsize (operation)}}\\
\noindent\textcolor{FuncColor}{$\triangleright$\enspace\texttt{SingularPlanarModularPartitionMonoid({\mdseries\slshape m, n})\index{SingularPlanarModularPartitionMonoid@\texttt{Singular}\-\texttt{Planar}\-\texttt{Modular}\-\texttt{Partition}\-\texttt{Monoid}}
\label{SingularPlanarModularPartitionMonoid}
}\hfill{\scriptsize (operation)}}\\
\textbf{\indent Returns:\ }
A bipartition monoid.



 If \mbox{\texttt{\mdseries\slshape m}} and \mbox{\texttt{\mdseries\slshape n}} are positive integers, then this operation returns the
modular\texttt{\symbol{45}}\mbox{\texttt{\mdseries\slshape m}} partition monoid of degree \mbox{\texttt{\mdseries\slshape n}}. The \emph{modular\texttt{\symbol{45}}}\mbox{\texttt{\mdseries\slshape m}} \emph{partition monoid} is the submonoid of the partition monoid such that the numbers of positive and
negative integers contained in each block are congruent mod \mbox{\texttt{\mdseries\slshape m}}. 

 \texttt{SingularModularPartitionMonoid} returns the ideal of the modular partition monoid consisting of the
non\texttt{\symbol{45}}invertible elements (i.e. those not in the group of
units), when either \mbox{\texttt{\mdseries\slshape m = n = 1}} or \mbox{\texttt{\mdseries\slshape m, n {\textgreater} 1}}. 

 \texttt{PlanarModularPartitionMonoid} returns the submonoid of the modular\texttt{\symbol{45}}\mbox{\texttt{\mdseries\slshape m}} partition monoid consisting of the planar elements (i.e. those in the planar
partition monoid, see \texttt{PlanarPartitionMonoid} (\ref{PlanarPartitionMonoid})). 

 \texttt{SingularPlanarModularPartitionMonoid} returns the ideal of the planar modular partition monoid consisting of the
non\texttt{\symbol{45}}invertible elements (i.e. those not in the group of
units), when either \mbox{\texttt{\mdseries\slshape m = n = 1}} or \mbox{\texttt{\mdseries\slshape m, n {\textgreater} 1}}.  
\begin{Verbatim}[commandchars=!@|,fontsize=\small,frame=single,label=Example]
  !gapprompt@gap>| !gapinput@S := ModularPartitionMonoid(3, 6);|
  <regular bipartition *-monoid of degree 6 with 4 generators>
  !gapprompt@gap>| !gapinput@Size(S);|
  36243
  !gapprompt@gap>| !gapinput@S := SingularModularPartitionMonoid(1, 1);|
  <commutative inverse bipartition semigroup ideal of degree 1 with
    1 generator>
  !gapprompt@gap>| !gapinput@Size(SingularModularPartitionMonoid(2, 4));|
  355
  !gapprompt@gap>| !gapinput@S := PlanarModularPartitionMonoid(4, 9);|
  <regular bipartition *-monoid of degree 9 with 14 generators>
  !gapprompt@gap>| !gapinput@Size(S);|
  1795
  !gapprompt@gap>| !gapinput@S := SingularPlanarModularPartitionMonoid(3, 5);|
  <regular bipartition *-semigroup ideal of degree 5 with 1 generator>
  !gapprompt@gap>| !gapinput@Size(SingularPlanarModularPartitionMonoid(1, 2));|
  13
\end{Verbatim}
 }

 

\subsection{\textcolor{Chapter }{ApsisMonoid}}
\logpage{[ 7, 3, 11 ]}\nobreak
\hyperdef{L}{X7C82B25F8441928E}{}
{\noindent\textcolor{FuncColor}{$\triangleright$\enspace\texttt{ApsisMonoid({\mdseries\slshape m, n})\index{ApsisMonoid@\texttt{ApsisMonoid}}
\label{ApsisMonoid}
}\hfill{\scriptsize (operation)}}\\
\noindent\textcolor{FuncColor}{$\triangleright$\enspace\texttt{SingularApsisMonoid({\mdseries\slshape m, n})\index{SingularApsisMonoid@\texttt{SingularApsisMonoid}}
\label{SingularApsisMonoid}
}\hfill{\scriptsize (operation)}}\\
\noindent\textcolor{FuncColor}{$\triangleright$\enspace\texttt{CrossedApsisMonoid({\mdseries\slshape m, n})\index{CrossedApsisMonoid@\texttt{CrossedApsisMonoid}}
\label{CrossedApsisMonoid}
}\hfill{\scriptsize (operation)}}\\
\noindent\textcolor{FuncColor}{$\triangleright$\enspace\texttt{SingularCrossedApsisMonoid({\mdseries\slshape m, n})\index{SingularCrossedApsisMonoid@\texttt{SingularCrossedApsisMonoid}}
\label{SingularCrossedApsisMonoid}
}\hfill{\scriptsize (operation)}}\\
\textbf{\indent Returns:\ }
A bipartition monoid.



 If \mbox{\texttt{\mdseries\slshape m}} and \mbox{\texttt{\mdseries\slshape n}} are positive integers, then this operation returns the \mbox{\texttt{\mdseries\slshape m}}\texttt{\symbol{45}}apsis monoid of degree \mbox{\texttt{\mdseries\slshape n}}. The \mbox{\texttt{\mdseries\slshape m}}\emph{\texttt{\symbol{45}}apsis monoid} is the monoid of bipartitions generated when the diapses in generators of the
Jones monoid are replaced with \mbox{\texttt{\mdseries\slshape m}}\texttt{\symbol{45}}apses. Note that an \mbox{\texttt{\mdseries\slshape m}}\emph{\texttt{\symbol{45}}apsis} is a block that contains precisely \mbox{\texttt{\mdseries\slshape m}} consecutive integers. 

 \texttt{SingularApsisMonoid} returns the ideal of the apsis monoid consisting of the
non\texttt{\symbol{45}}invertible elements (i.e. those not in the group of
units), when \mbox{\texttt{\mdseries\slshape m}} $\leq$ \mbox{\texttt{\mdseries\slshape n}}. 

 \texttt{CrossedApsisGeneratedMonoid} returns the semigroup generated by the symmetric group of degree \mbox{\texttt{\mdseries\slshape n}} and the \mbox{\texttt{\mdseries\slshape m}}\texttt{\symbol{45}}apsis monoid of degree \mbox{\texttt{\mdseries\slshape n}}. 

 \texttt{SingularCrossedApsisMonoid} returns the ideal of the crossed apsis monoid consisting of the
non\texttt{\symbol{45}}invertible elements (i.e. those not in the group of
units), when \mbox{\texttt{\mdseries\slshape m {\textless}= n}}.  
\begin{Verbatim}[commandchars=!@|,fontsize=\small,frame=single,label=Example]
  !gapprompt@gap>| !gapinput@S := ApsisMonoid(3, 7);|
  <regular bipartition *-monoid of degree 7 with 5 generators>
  !gapprompt@gap>| !gapinput@Size(S);|
  320
  !gapprompt@gap>| !gapinput@T := SingularApsisMonoid(3, 7);|
  <regular bipartition *-semigroup ideal of degree 7 with 1 generator>
  !gapprompt@gap>| !gapinput@Difference(S, T) = [One(S)];|
  true
  !gapprompt@gap>| !gapinput@Size(CrossedApsisMonoid(2, 5));|
  945
  !gapprompt@gap>| !gapinput@SingularCrossedApsisMonoid(4, 6);|
  <regular bipartition *-semigroup ideal of degree 6 with 1 generator>
\end{Verbatim}
 }

 }

   
\section{\textcolor{Chapter }{ Standard PBR semigroups }}\label{Standard PBR semigroups}
\logpage{[ 7, 4, 0 ]}
\hyperdef{L}{X874C945E7C61A969}{}
{
  In this section, we describe the operations in \textsf{Semigroups} that can be used to create standard examples of semigroups of partitioned
binary relations (PBRs). See Chapter \ref{Partitioned binary relations (PBRs)} for more information about PBRs. 

\subsection{\textcolor{Chapter }{FullPBRMonoid}}
\logpage{[ 7, 4, 1 ]}\nobreak
\hyperdef{L}{X7DBB30AA83663CE8}{}
{\noindent\textcolor{FuncColor}{$\triangleright$\enspace\texttt{FullPBRMonoid({\mdseries\slshape n})\index{FullPBRMonoid@\texttt{FullPBRMonoid}}
\label{FullPBRMonoid}
}\hfill{\scriptsize (operation)}}\\
\textbf{\indent Returns:\ }
A PBR monoid.



 If \mbox{\texttt{\mdseries\slshape n}} is a positive integer not greater than \texttt{2}, then this operation returns the monoid consisting of all of the partitioned
binary relations (PBRs) of degree \mbox{\texttt{\mdseries\slshape n}}; called the \emph{full PBR monoid}. There are \texttt{2 \texttt{\symbol{94}} ((2 * n) \texttt{\symbol{94}} 2)} PBRs of degree \mbox{\texttt{\mdseries\slshape n}}. The full PBR monoid of degree \mbox{\texttt{\mdseries\slshape n}} is currently too large to compute when $\mbox{\texttt{\mdseries\slshape n}} \geq 3$. 

 The full PBR monoid is not regular in general. 
\begin{Verbatim}[commandchars=!@|,fontsize=\small,frame=single,label=Example]
  !gapprompt@gap>| !gapinput@S := FullPBRMonoid(1);|
  <pbr monoid of degree 1 with 4 generators>
  !gapprompt@gap>| !gapinput@S := FullPBRMonoid(2);|
  <pbr monoid of degree 2 with 10 generators>
\end{Verbatim}
 }

 }

   
\section{\textcolor{Chapter }{ Semigroups of matrices over a finite field }}\label{Semigroups of matrices of a finite field}
\logpage{[ 7, 5, 0 ]}
\hyperdef{L}{X857DBF537A9A9976}{}
{
  In this section, we describe the operations in \textsf{Semigroups} that can be used to create semigroups of matrices over a finite field that
belonging to several standard classes of example. See the section \hyperref[Matrices over finite fields]{`Matrices over finite fields'} for more information about matrices over a finite field. 

\subsection{\textcolor{Chapter }{FullMatrixMonoid}}
\logpage{[ 7, 5, 1 ]}\nobreak
\hyperdef{L}{X7D4B473A7D7735E3}{}
{\noindent\textcolor{FuncColor}{$\triangleright$\enspace\texttt{FullMatrixMonoid({\mdseries\slshape d, q})\index{FullMatrixMonoid@\texttt{FullMatrixMonoid}}
\label{FullMatrixMonoid}
}\hfill{\scriptsize (operation)}}\\
\noindent\textcolor{FuncColor}{$\triangleright$\enspace\texttt{GeneralLinearMonoid({\mdseries\slshape d, q})\index{GeneralLinearMonoid@\texttt{GeneralLinearMonoid}}
\label{GeneralLinearMonoid}
}\hfill{\scriptsize (operation)}}\\
\noindent\textcolor{FuncColor}{$\triangleright$\enspace\texttt{GLM({\mdseries\slshape d, q})\index{GLM@\texttt{GLM}}
\label{GLM}
}\hfill{\scriptsize (operation)}}\\
\textbf{\indent Returns:\ }
A matrix monoid.



 These operations return the full matrix monoid of \mbox{\texttt{\mdseries\slshape d}} by \mbox{\texttt{\mdseries\slshape d}} matrices over the field with \mbox{\texttt{\mdseries\slshape q}} elements. The \emph{full matrix monoid}, also known as the \emph{general linear monoid}, with these parameters, is the monoid consisting of all \mbox{\texttt{\mdseries\slshape d}} by \mbox{\texttt{\mdseries\slshape d}} matrices with entries from the field \texttt{GF(\mbox{\texttt{\mdseries\slshape q}})}. This monoid has \texttt{\mbox{\texttt{\mdseries\slshape q}} \texttt{\symbol{94}} (\mbox{\texttt{\mdseries\slshape d}} \texttt{\symbol{94}} 2)} elements. 
\begin{Verbatim}[commandchars=!@|,fontsize=\small,frame=single,label=Example]
  !gapprompt@gap>| !gapinput@S := FullMatrixMonoid(2, 4);|
  <general linear monoid 2x2 over GF(2^2)>
  !gapprompt@gap>| !gapinput@Size(S);|
  256
  !gapprompt@gap>| !gapinput@S = GeneralLinearMonoid(2, 4);|
  true
  !gapprompt@gap>| !gapinput@GLM(2, 2);|
  <general linear monoid 2x2 over GF(2)>
\end{Verbatim}
 }

 

\subsection{\textcolor{Chapter }{SpecialLinearMonoid}}
\logpage{[ 7, 5, 2 ]}\nobreak
\hyperdef{L}{X785924807B60F187}{}
{\noindent\textcolor{FuncColor}{$\triangleright$\enspace\texttt{SpecialLinearMonoid({\mdseries\slshape d, q})\index{SpecialLinearMonoid@\texttt{SpecialLinearMonoid}}
\label{SpecialLinearMonoid}
}\hfill{\scriptsize (operation)}}\\
\noindent\textcolor{FuncColor}{$\triangleright$\enspace\texttt{SLM({\mdseries\slshape d, q})\index{SLM@\texttt{SLM}}
\label{SLM}
}\hfill{\scriptsize (operation)}}\\
\textbf{\indent Returns:\ }
A matrix monoid.



 These operations return the special linear monoid of \mbox{\texttt{\mdseries\slshape d}} by \mbox{\texttt{\mdseries\slshape d}} matrices over the field with \mbox{\texttt{\mdseries\slshape q}} elements. The \emph{special linear monoid} is the monoid consisting of all \mbox{\texttt{\mdseries\slshape d}} by \mbox{\texttt{\mdseries\slshape d}} matrices with entries from the field \texttt{GF(\mbox{\texttt{\mdseries\slshape q}})} that have determinant \texttt{0} or \texttt{1}. In other words, the special linear monoid is formed from the general linear
monoid of the same parameters by replacing its group of units (the general
linear group) by the special linear group. 
\begin{Verbatim}[commandchars=!@|,fontsize=\small,frame=single,label=Example]
  !gapprompt@gap>| !gapinput@S := SpecialLinearMonoid(2, 4);|
  <regular monoid of 2x2 matrices over GF(2^2) with 3 generators>
  !gapprompt@gap>| !gapinput@S = SLM(2, 4);|
  true
  !gapprompt@gap>| !gapinput@Size(S);|
  136
\end{Verbatim}
 }

  

\subsection{\textcolor{Chapter }{IsFullMatrixMonoid}}
\logpage{[ 7, 5, 3 ]}\nobreak
\hyperdef{L}{X860B2A4382CA8F87}{}
{\noindent\textcolor{FuncColor}{$\triangleright$\enspace\texttt{IsFullMatrixMonoid({\mdseries\slshape S})\index{IsFullMatrixMonoid@\texttt{IsFullMatrixMonoid}}
\label{IsFullMatrixMonoid}
}\hfill{\scriptsize (property)}}\\
\noindent\textcolor{FuncColor}{$\triangleright$\enspace\texttt{IsGeneralLinearMonoid({\mdseries\slshape S})\index{IsGeneralLinearMonoid@\texttt{IsGeneralLinearMonoid}}
\label{IsGeneralLinearMonoid}
}\hfill{\scriptsize (property)}}\\


 \texttt{IsFullMatrixMonoid} and \texttt{IsGeneralLinearMonoid} return \texttt{true} if the semigroup \texttt{S} was created using either of the commands \texttt{FullMatrixMonoid} (\ref{FullMatrixMonoid}) or \texttt{GeneralLinearMonoid} (\ref{GeneralLinearMonoid}) and \texttt{false} otherwise. 
\begin{Verbatim}[commandchars=!@|,fontsize=\small,frame=single,label=Example]
  !gapprompt@gap>| !gapinput@S := RandomSemigroup(IsTransformationSemigroup, 4, 4);;|
  !gapprompt@gap>| !gapinput@IsFullMatrixMonoid(S);|
  false
  !gapprompt@gap>| !gapinput@S := GeneralLinearMonoid(3, 3);|
  <general linear monoid 3x3 over GF(3)>
  !gapprompt@gap>| !gapinput@IsFullMatrixMonoid(S);|
  true
\end{Verbatim}
 }

 }

   
\section{\textcolor{Chapter }{ Semigroups of boolean matrices }}\label{Semigroups of boolean matrices}
\logpage{[ 7, 6, 0 ]}
\hyperdef{L}{X85BACB7F81660ECC}{}
{
  In this section, we describe the operations in \textsf{Semigroups} that can be used to create semigroups of boolean matrices belonging to several
standard classes of example. See the section \hyperref[Boolean matrices]{`Boolean matrices'} for more information about boolean matrices. 

\subsection{\textcolor{Chapter }{FullBooleanMatMonoid}}
\logpage{[ 7, 6, 1 ]}\nobreak
\hyperdef{L}{X7B20103D84E010EF}{}
{\noindent\textcolor{FuncColor}{$\triangleright$\enspace\texttt{FullBooleanMatMonoid({\mdseries\slshape d})\index{FullBooleanMatMonoid@\texttt{FullBooleanMatMonoid}}
\label{FullBooleanMatMonoid}
}\hfill{\scriptsize (operation)}}\\
\textbf{\indent Returns:\ }
The monoid of all boolean matrices of dimension \mbox{\texttt{\mdseries\slshape d}}.



 If \mbox{\texttt{\mdseries\slshape d}} is a positive integer less than or equal to \texttt{5}, then this operation returns the full boolean matrix monoid of dimension \mbox{\texttt{\mdseries\slshape d}}. The \emph{full boolean matrix monoid of dimension \mbox{\texttt{\mdseries\slshape d}}} is the monoid consisting of all \mbox{\texttt{\mdseries\slshape d}} by \mbox{\texttt{\mdseries\slshape d}} boolean matrices, and has \texttt{2 \texttt{\symbol{94}} (\mbox{\texttt{\mdseries\slshape n}} \texttt{\symbol{94}} 2)} matrices. 

 \texttt{FullBooleanMatMonoid} returns a monoid with a generating set that is minimal in size. These
generating sets are pre\texttt{\symbol{45}}computed. 
\begin{Verbatim}[commandchars=!@|,fontsize=\small,frame=single,label=Example]
  !gapprompt@gap>| !gapinput@S := FullBooleanMatMonoid(3);|
  <monoid of 3x3 boolean matrices with 5 generators>
  !gapprompt@gap>| !gapinput@Size(S);|
  512
\end{Verbatim}
 }

 

\subsection{\textcolor{Chapter }{RegularBooleanMatMonoid}}
\logpage{[ 7, 6, 2 ]}\nobreak
\hyperdef{L}{X7A43263981F2F2AF}{}
{\noindent\textcolor{FuncColor}{$\triangleright$\enspace\texttt{RegularBooleanMatMonoid({\mdseries\slshape d})\index{RegularBooleanMatMonoid@\texttt{RegularBooleanMatMonoid}}
\label{RegularBooleanMatMonoid}
}\hfill{\scriptsize (operation)}}\\
\textbf{\indent Returns:\ }
A monoid of boolean matrices.



 If \mbox{\texttt{\mdseries\slshape d}} is a positive integer, then \texttt{RegularBooleanMatMonoid} returns the monoid generated by the regular \mbox{\texttt{\mdseries\slshape d}} by \mbox{\texttt{\mdseries\slshape d}} boolean matrices. Note that this monoid is \emph{not} regular in general. \texttt{RegularBooleanMatMonoid(\mbox{\texttt{\mdseries\slshape d}})} is generated by the four boolean matrices \texttt{A, B, C, D}, whose \texttt{true} entries are: 
\begin{itemize}
\item  \texttt{A[i][i + 1]} and \texttt{A[n][1]}, for $i \in \{1, \ldots, n - 1\}$; 
\item  \texttt{B[1][2]}, \texttt{B[2][1]}, and \texttt{B[i][i]} for $i \in \{3, \ldots, n\}$; 
\item  \texttt{C[1][2]} and \texttt{C[i][i]}, for $i \in \{2, \ldots, n - 1\}$; and 
\item  \texttt{D[1][2]}, \texttt{D[i][i]}, for $i \in \{2, \ldots, n\}$, and \texttt{D[n][1]}. 
\end{itemize}
 This monoid has nearly \texttt{2 \texttt{\symbol{94}} (n \texttt{\symbol{94}} 2)} elements. 
\begin{Verbatim}[commandchars=!@|,fontsize=\small,frame=single,label=Example]
  !gapprompt@gap>| !gapinput@S := RegularBooleanMatMonoid(3);|
  <monoid of 3x3 boolean matrices with 4 generators>
  !gapprompt@gap>| !gapinput@Size(S);|
  506
\end{Verbatim}
 }

 

\subsection{\textcolor{Chapter }{ReflexiveBooleanMatMonoid}}
\logpage{[ 7, 6, 3 ]}\nobreak
\hyperdef{L}{X78DF50747A28098C}{}
{\noindent\textcolor{FuncColor}{$\triangleright$\enspace\texttt{ReflexiveBooleanMatMonoid({\mdseries\slshape d})\index{ReflexiveBooleanMatMonoid@\texttt{ReflexiveBooleanMatMonoid}}
\label{ReflexiveBooleanMatMonoid}
}\hfill{\scriptsize (operation)}}\\
\textbf{\indent Returns:\ }
A monoid of boolean matrices.



 If \mbox{\texttt{\mdseries\slshape d}} is a positive integer less than or equal to \texttt{5}, then this operation returns the monoid consisting of all reflexive \mbox{\texttt{\mdseries\slshape d}} by \mbox{\texttt{\mdseries\slshape d}} boolean matrices. A boolean matrix \texttt{mat} is \emph{reflexive} if each entry of its leading diagonal is \texttt{true}, i.e. if \texttt{mat[i][i]} is \texttt{true} for all $i \in \{1, \ldots, d\}$. 

 The generating sets for the monoids returned by \texttt{ReflexiveBooleanMatMonoid} are pre\texttt{\symbol{45}}computed, and read from a file. Small generating
sets are not known for $\mbox{\texttt{\mdseries\slshape d}} \geq 6$. 
\begin{Verbatim}[commandchars=!@|,fontsize=\small,frame=single,label=Example]
  !gapprompt@gap>| !gapinput@S := ReflexiveBooleanMatMonoid(3);|
  <monoid of 3x3 boolean matrices with 8 generators>
  !gapprompt@gap>| !gapinput@Size(S);|
  64
\end{Verbatim}
 }

 

\subsection{\textcolor{Chapter }{HallMonoid}}
\logpage{[ 7, 6, 4 ]}\nobreak
\hyperdef{L}{X79EF0EA68782CFCA}{}
{\noindent\textcolor{FuncColor}{$\triangleright$\enspace\texttt{HallMonoid({\mdseries\slshape d})\index{HallMonoid@\texttt{HallMonoid}}
\label{HallMonoid}
}\hfill{\scriptsize (operation)}}\\
\textbf{\indent Returns:\ }
A monoid of boolean matrices.



 If \mbox{\texttt{\mdseries\slshape d}} is a positive integer less than or equal to \texttt{5}, then this operation returns the monoid consisting Hall matrices of degree \mbox{\texttt{\mdseries\slshape d}}. A \emph{Hall matrix} is a boolean matrix in which every column and every row contains at least one \texttt{true} entry. Equivalently, a Hall matrix is a boolean matrix than contains a
permutation. 

 A Hall matrix of dimension \mbox{\texttt{\mdseries\slshape d}} corresponds to a solution to Hall's Marriage Problem, when there are two
collection of \mbox{\texttt{\mdseries\slshape d}} people. Thus the number of solutions to Hall's Marriage Problem in this
instance is the number of elements of \texttt{HallMonoid(\mbox{\texttt{\mdseries\slshape d}})}. 

 The operation \texttt{HallMonoid} returns a monoid with a generating set that is minimal in size. These
generating sets are pre\texttt{\symbol{45}}computed, and a minimal generating
set is not known for larger dimensions. 
\begin{Verbatim}[commandchars=!@|,fontsize=\small,frame=single,label=Example]
  !gapprompt@gap>| !gapinput@S := HallMonoid(3);|
  <monoid of 3x3 boolean matrices with 4 generators>
  !gapprompt@gap>| !gapinput@Size(S);|
  247
\end{Verbatim}
 }

 

\subsection{\textcolor{Chapter }{GossipMonoid}}
\logpage{[ 7, 6, 5 ]}\nobreak
\hyperdef{L}{X7F083600787C78FF}{}
{\noindent\textcolor{FuncColor}{$\triangleright$\enspace\texttt{GossipMonoid({\mdseries\slshape d})\index{GossipMonoid@\texttt{GossipMonoid}}
\label{GossipMonoid}
}\hfill{\scriptsize (operation)}}\\
\textbf{\indent Returns:\ }
A monoid of boolean matrices.



 If \mbox{\texttt{\mdseries\slshape d}} is a positive integer, then this operation returns the \mbox{\texttt{\mdseries\slshape d}} by \mbox{\texttt{\mdseries\slshape d}} gossip monoid. The \emph{gossip monoid} is defined to be the monoid generated by the collection of all \mbox{\texttt{\mdseries\slshape d}} by \mbox{\texttt{\mdseries\slshape d}} boolean matrices that define an equivalence relation; see \texttt{IsEquivalenceBooleanMat} (\ref{IsEquivalenceBooleanMat}). 

 For $\mbox{\texttt{\mdseries\slshape d}} \geq 2$, \texttt{GossipMonoid(\mbox{\texttt{\mdseries\slshape d}})} returns a monoid with ${d \choose 2}$ generators. The generating set is the collection of boolean matrices that
define an equivalence relation that has one equivalence class of size \texttt{2}, and no other non\texttt{\symbol{45}}trivial equivalence classes. Note that
this generating set is strictly contained within the collection of all
equivalence relation boolean matrices. 

 The number of elements of \texttt{GossipMonoid(\mbox{\texttt{\mdseries\slshape d}})} is known for some small values of \mbox{\texttt{\mdseries\slshape d}} {\textemdash} see \cite{Brouwer2015aa} for more information about the gossip monoid, and its size for $\mbox{\texttt{\mdseries\slshape d}} \leq 9$. 
\begin{Verbatim}[commandchars=!@|,fontsize=\small,frame=single,label=Example]
  !gapprompt@gap>| !gapinput@S := GossipMonoid(3);|
  <monoid of 3x3 boolean matrices with 3 generators>
  !gapprompt@gap>| !gapinput@Size(S);|
  11
\end{Verbatim}
 }

 

\subsection{\textcolor{Chapter }{TriangularBooleanMatMonoid}}
\logpage{[ 7, 6, 6 ]}\nobreak
\hyperdef{L}{X81BBCF2E84239521}{}
{\noindent\textcolor{FuncColor}{$\triangleright$\enspace\texttt{TriangularBooleanMatMonoid({\mdseries\slshape d})\index{TriangularBooleanMatMonoid@\texttt{TriangularBooleanMatMonoid}}
\label{TriangularBooleanMatMonoid}
}\hfill{\scriptsize (operation)}}\\
\noindent\textcolor{FuncColor}{$\triangleright$\enspace\texttt{UnitriangularBooleanMatMonoid({\mdseries\slshape d})\index{UnitriangularBooleanMatMonoid@\texttt{UnitriangularBooleanMatMonoid}}
\label{UnitriangularBooleanMatMonoid}
}\hfill{\scriptsize (operation)}}\\
\textbf{\indent Returns:\ }
A monoid of boolean matrices.



 If \mbox{\texttt{\mdseries\slshape d}} is a positive integer, then \texttt{TriangularBooleanMatMonoid} returns the monoid consisting of the upper\texttt{\symbol{45}}triangular \mbox{\texttt{\mdseries\slshape d}} by \mbox{\texttt{\mdseries\slshape d}} boolean matrices. A boolean matrix is \emph{upper\texttt{\symbol{45}}triangular} if the entry in row \texttt{i}, column \texttt{j} is \texttt{false} whenever \texttt{i {\textgreater} j}. 

 \texttt{UnitriangularBooleanMatMonoid} returns the subsemigroup of the \texttt{TriangularBooleanMatMonoid} that consists of reflexive upper\texttt{\symbol{45}}triangular boolean
matrices; see \texttt{ReflexiveBooleanMatMonoid} (\ref{ReflexiveBooleanMatMonoid}). 
\begin{Verbatim}[commandchars=!@|,fontsize=\small,frame=single,label=Example]
  !gapprompt@gap>| !gapinput@S := TriangularBooleanMatMonoid(3);|
  <monoid of 3x3 boolean matrices with 6 generators>
  !gapprompt@gap>| !gapinput@Size(S);|
  64
  !gapprompt@gap>| !gapinput@T := UnitriangularBooleanMatMonoid(4);|
  <monoid of 4x4 boolean matrices with 6 generators>
  !gapprompt@gap>| !gapinput@Size(T);|
  64
\end{Verbatim}
 }

 }

   
\section{\textcolor{Chapter }{ Semigroups of matrices over a semiring }}\label{Semigroups of matrices over a semiring}
\logpage{[ 7, 7, 0 ]}
\hyperdef{L}{X7F3D0AEE79AA8C98}{}
{
  In this section, we describe the operations in \textsf{Semigroups} that can be used to create semigroups of matices over a semiring that belong
to several standard classes of example. See Chapter \ref{Matrices over semirings} for more information about matrices over a semiring. 

\subsection{\textcolor{Chapter }{FullTropicalMaxPlusMonoid}}
\logpage{[ 7, 7, 1 ]}\nobreak
\hyperdef{L}{X81E937B6852A9C69}{}
{\noindent\textcolor{FuncColor}{$\triangleright$\enspace\texttt{FullTropicalMaxPlusMonoid({\mdseries\slshape d, t})\index{FullTropicalMaxPlusMonoid@\texttt{FullTropicalMaxPlusMonoid}}
\label{FullTropicalMaxPlusMonoid}
}\hfill{\scriptsize (operation)}}\\
\textbf{\indent Returns:\ }
A monoid of tropical max plus matrices.



 If \texttt{\mbox{\texttt{\mdseries\slshape d}} = 2} and \mbox{\texttt{\mdseries\slshape t}} is a positive integer, then \texttt{FullTropicalMaxPlusMonoid} returns the monoid consisting of all \mbox{\texttt{\mdseries\slshape d}} by \mbox{\texttt{\mdseries\slshape d}} matrices with entries from the tropical max\texttt{\symbol{45}}plus semiring
with threshold \mbox{\texttt{\mdseries\slshape t}}. A small generating set for larger values of \mbox{\texttt{\mdseries\slshape d}} is not currently known. 

 This monoid contains \texttt{(\mbox{\texttt{\mdseries\slshape t}} + 2) \texttt{\symbol{94}} (\mbox{\texttt{\mdseries\slshape d}} \texttt{\symbol{94}} 2)} elements. 
\begin{Verbatim}[commandchars=!@|,fontsize=\small,frame=single,label=Example]
  !gapprompt@gap>| !gapinput@S := FullTropicalMaxPlusMonoid(2, 5);|
  <monoid of 2x2 tropical max-plus matrices with 24 generators>
  !gapprompt@gap>| !gapinput@Size(S);|
  2401
  !gapprompt@gap>| !gapinput@(5 + 2) ^ (2 ^ 2);|
  2401
\end{Verbatim}
 }

 

\subsection{\textcolor{Chapter }{FullTropicalMinPlusMonoid}}
\logpage{[ 7, 7, 2 ]}\nobreak
\hyperdef{L}{X85EDC03180768931}{}
{\noindent\textcolor{FuncColor}{$\triangleright$\enspace\texttt{FullTropicalMinPlusMonoid({\mdseries\slshape d, t})\index{FullTropicalMinPlusMonoid@\texttt{FullTropicalMinPlusMonoid}}
\label{FullTropicalMinPlusMonoid}
}\hfill{\scriptsize (operation)}}\\
\textbf{\indent Returns:\ }
A monoid of tropical min plus matrices.



 If \mbox{\texttt{\mdseries\slshape d}} is equal to \texttt{2} or \texttt{3}, and \mbox{\texttt{\mdseries\slshape t}} is a positive integer, then \texttt{FullTropicalMinPlusMonoid} returns the monoid consisting of all \mbox{\texttt{\mdseries\slshape d}} by \mbox{\texttt{\mdseries\slshape d}} matrices with entries from the tropical min\texttt{\symbol{45}}plus semiring
with threshold \mbox{\texttt{\mdseries\slshape t}}. A small generating set for larger values of \mbox{\texttt{\mdseries\slshape d}} is not currently known. 

 This monoid contains \texttt{(\mbox{\texttt{\mdseries\slshape t}} + 2) \texttt{\symbol{94}} (\mbox{\texttt{\mdseries\slshape d}} \texttt{\symbol{94}} 2)} elements. 
\begin{Verbatim}[commandchars=!@|,fontsize=\small,frame=single,label=Example]
  !gapprompt@gap>| !gapinput@S := FullTropicalMinPlusMonoid(2, 3);|
  <monoid of 2x2 tropical min-plus matrices with 7 generators>
  !gapprompt@gap>| !gapinput@Size(S);|
  625
  !gapprompt@gap>| !gapinput@(3 + 2) ^ (2 ^ 2);|
  625
\end{Verbatim}
 }

 }

   
\section{\textcolor{Chapter }{ Examples in various representations }}\label{Examples in various representations}
\logpage{[ 7, 8, 0 ]}
\hyperdef{L}{X7ED2F2577CD6B578}{}
{
  In this section, we describe the functions in \textsf{Semigroups} that can be used to create standard semigroups in various representations. For
all of these examples, the default representation is as a semigroup of
transformations. In general, these functions do not return a representation of
minimal degree. 

\subsection{\textcolor{Chapter }{TrivialSemigroup}}
\logpage{[ 7, 8, 1 ]}\nobreak
\hyperdef{L}{X82B07E907B3A55F0}{}
{\noindent\textcolor{FuncColor}{$\triangleright$\enspace\texttt{TrivialSemigroup({\mdseries\slshape [filt][,] [deg]})\index{TrivialSemigroup@\texttt{TrivialSemigroup}}
\label{TrivialSemigroup}
}\hfill{\scriptsize (function)}}\\
\textbf{\indent Returns:\ }
 A trivial semigroup. 



 A \textsc{trivial} semigroup is a semigroup with precisely one element. This function returns a
trivial semigroup in the representation given by the filter \mbox{\texttt{\mdseries\slshape filter}}, and (if possible) with the degree of the representation given by the
non\texttt{\symbol{45}}negative integer \mbox{\texttt{\mdseries\slshape deg}}. 

 The optional argument \mbox{\texttt{\mdseries\slshape filt}} may be one of the following: 
\begin{itemize}
\item \texttt{IsTransformationSemigroup} (the default, if \mbox{\texttt{\mdseries\slshape filt}} is not specified),
\item \texttt{IsPartialPermSemigroup},
\item \texttt{IsBipartitionSemigroup},
\item \texttt{IsBlockBijectionSemigroup},
\item \texttt{IsPBRSemigroup},
\item \texttt{IsBooleanMatSemigroup}.
\end{itemize}
 If the optional argument \mbox{\texttt{\mdseries\slshape deg}} is not specified, then the smallest possible degree will be used. 

 
\begin{Verbatim}[commandchars=!@|,fontsize=\small,frame=single,label=Example]
  !gapprompt@gap>| !gapinput@S := TrivialSemigroup();|
  <trivial transformation group of degree 0 with 1 generator>
  !gapprompt@gap>| !gapinput@Size(S);|
  1
  !gapprompt@gap>| !gapinput@S := TrivialSemigroup(3);|
  <trivial transformation group of degree 3 with 1 generator>
  !gapprompt@gap>| !gapinput@S := TrivialSemigroup(IsBipartitionSemigroup, 2);|
  <trivial block bijection group of degree 2 with 1 generator>
  !gapprompt@gap>| !gapinput@Elements(S);|
  [ <block bijection: [ 1, 2, -1, -2 ]> ]
\end{Verbatim}
 }

 

\subsection{\textcolor{Chapter }{MonogenicSemigroup}}
\logpage{[ 7, 8, 2 ]}\nobreak
\hyperdef{L}{X8411EBD97A220921}{}
{\noindent\textcolor{FuncColor}{$\triangleright$\enspace\texttt{MonogenicSemigroup({\mdseries\slshape [filt, ]m, r})\index{MonogenicSemigroup@\texttt{MonogenicSemigroup}}
\label{MonogenicSemigroup}
}\hfill{\scriptsize (function)}}\\
\textbf{\indent Returns:\ }
 A monogenic semigroup with index \mbox{\texttt{\mdseries\slshape m}} and period \mbox{\texttt{\mdseries\slshape r}}. 



 If \mbox{\texttt{\mdseries\slshape m}} and \mbox{\texttt{\mdseries\slshape r}} are positive integers, then this function returns a monogenic semigroup \texttt{S} with index \mbox{\texttt{\mdseries\slshape m}} and period \mbox{\texttt{\mdseries\slshape r}} in the representation given by the filter \mbox{\texttt{\mdseries\slshape filt}}. 

 The optional argument \mbox{\texttt{\mdseries\slshape filt}} may be one of the following: 
\begin{itemize}
\item \texttt{IsTransformationSemigroup} (the default, if \mbox{\texttt{\mdseries\slshape filt}} is not specified),
\item \texttt{IsPartialPermSemigroup},
\item \texttt{IsBipartitionSemigroup},
\item \texttt{IsBlockBijectionSemigroup},
\item \texttt{IsPBRSemigroup},
\item \texttt{IsBooleanMatSemigroup}.
\end{itemize}
 The semigroup \texttt{S} is generated by a single element, $f$. \texttt{S} consists of the elements $f, f ^ {2}, \ldots, f ^ {m}, \ldots, f ^ {m + r - 1}$. The minimal ideal of \texttt{S} consists of the elements $f ^ {m}, \ldots, f ^ {m + r - 1}$ and is isomorphic to the cyclic group of order $r$. 

 See \texttt{IsMonogenicSemigroup} (\ref{IsMonogenicSemigroup}) for more information about monogenic semigroups. 
\begin{Verbatim}[commandchars=!@|,fontsize=\small,frame=single,label=Example]
  !gapprompt@gap>| !gapinput@S := MonogenicSemigroup(5, 3);|
  <commutative non-regular transformation semigroup of size 7, degree 8
   with 1 generator>
  !gapprompt@gap>| !gapinput@IsMonogenicSemigroup(S);|
  true
  !gapprompt@gap>| !gapinput@I := MinimalIdeal(S);;|
  !gapprompt@gap>| !gapinput@IsGroupAsSemigroup(I);|
  true
  !gapprompt@gap>| !gapinput@StructureDescription(I);|
  "C3"
  !gapprompt@gap>| !gapinput@S := MonogenicSemigroup(IsBlockBijectionSemigroup, 9, 1);|
  <commutative non-regular block bijection semigroup of size 9,
   degree 10 with 1 generator>
\end{Verbatim}
 }

 

\subsection{\textcolor{Chapter }{RectangularBand}}
\logpage{[ 7, 8, 3 ]}\nobreak
\hyperdef{L}{X7E4DFDE27BF8B8F7}{}
{\noindent\textcolor{FuncColor}{$\triangleright$\enspace\texttt{RectangularBand({\mdseries\slshape [filt, ]m, n})\index{RectangularBand@\texttt{RectangularBand}}
\label{RectangularBand}
}\hfill{\scriptsize (function)}}\\
\textbf{\indent Returns:\ }
 An \mbox{\texttt{\mdseries\slshape m}} by \mbox{\texttt{\mdseries\slshape n}} rectangular band. 



 If \mbox{\texttt{\mdseries\slshape m}} and \mbox{\texttt{\mdseries\slshape n}} are positive integers, then this function returns a semigroup isomorphic to an \mbox{\texttt{\mdseries\slshape m}} by \mbox{\texttt{\mdseries\slshape n}} rectangular band, in the representation given by the filter \mbox{\texttt{\mdseries\slshape filt}}. 

 The optional argument \mbox{\texttt{\mdseries\slshape filt}} may be one of the following: 
\begin{itemize}
\item \texttt{IsTransformationSemigroup} (the default, if \mbox{\texttt{\mdseries\slshape filt}} is not specified),
\item \texttt{IsBipartitionSemigroup},
\item \texttt{IsPBRSemigroup},
\item \texttt{IsBooleanMatSemigroup},
\item \texttt{IsReesMatrixSemigroup}.
\end{itemize}
 See \texttt{IsRectangularBand} (\ref{IsRectangularBand}) for more information about rectangular bands. 
\begin{Verbatim}[commandchars=!@|,fontsize=\small,frame=single,label=Example]
  !gapprompt@gap>| !gapinput@T := RectangularBand(5, 6);|
  <regular transformation semigroup of size 30, degree 10 with 6
   generators>
  !gapprompt@gap>| !gapinput@IsRectangularBand(T);|
  true
  !gapprompt@gap>| !gapinput@S := RectangularBand(IsReesMatrixSemigroup, 4, 8);|
  <Rees matrix semigroup 4x8 over Group(())>
  !gapprompt@gap>| !gapinput@IsRectangularBand(S);|
  true
  !gapprompt@gap>| !gapinput@IsCompletelySimpleSemigroup(S) and IsHTrivial(S);|
  true
\end{Verbatim}
 }

 

\subsection{\textcolor{Chapter }{FreeSemilattice}}
\logpage{[ 7, 8, 4 ]}\nobreak
\hyperdef{L}{X7982E0667ECEB265}{}
{\noindent\textcolor{FuncColor}{$\triangleright$\enspace\texttt{FreeSemilattice({\mdseries\slshape [filt, ]n})\index{FreeSemilattice@\texttt{FreeSemilattice}}
\label{FreeSemilattice}
}\hfill{\scriptsize (function)}}\\
\textbf{\indent Returns:\ }
 A free semilattice with \mbox{\texttt{\mdseries\slshape n}} generators. 



 If \mbox{\texttt{\mdseries\slshape n}} is a positive integer, then this function returns a free semilattice with \mbox{\texttt{\mdseries\slshape n}} generators in the representation given by the filter \mbox{\texttt{\mdseries\slshape filt}}. The optional argument \mbox{\texttt{\mdseries\slshape filt}} may be one of the following: 
\begin{itemize}
\item \texttt{IsTransformationSemigroup} (the default, if \mbox{\texttt{\mdseries\slshape filt}} is not specified),
\item \texttt{IsTransformationMonoid},
\item \texttt{IsPartialPermSemigroup},
\item \texttt{IsPartialPermMonoid},
\item \texttt{IsFpSemigroup},
\item \texttt{IsFpMonoid},
\item \texttt{IsBipartitionSemigroup},
\item \texttt{IsBipartitionMonoid},
\item \texttt{IsPBRSemigroup},
\item \texttt{IsPBRMonoid},
\item \texttt{IsBooleanMatSemigroup},
\item \texttt{IsBooleanMatMonoid},
\item \texttt{IsNTPMatrixSemigroup},
\item \texttt{IsNTPMatrixMonoid},
\item \texttt{IsMaxPlusMatrixSemigroup},
\item \texttt{IsMaxPlusMatrixMonoid},
\item \texttt{IsMinPlusMatrixSemigroup},
\item \texttt{IsMinPlusMatrixMonoid},
\item \texttt{IsTropicalMaxPlusMatrixSemigroup},
\item \texttt{IsTropicalMaxPlusMatrixMonoid},
\item \texttt{IsTropicalMinPlusMatrixSemigroup},
\item \texttt{IsTropicalMinPlusMatrixMonoid},
\item \texttt{IsProjectiveMaxPlusMatrixSemigroup},
\item \texttt{IsProjectiveMaxPlusMatrixMonoid},
\item \texttt{IsIntegerMatrixSemigroup.}
\item \texttt{IsIntegerMatrixMonoid.}
\end{itemize}
 
\begin{Verbatim}[commandchars=!@|,fontsize=\small,frame=single,label=Example]
  !gapprompt@gap>| !gapinput@S := FreeSemilattice(IsTransformationSemigroup, 5);|
  <inverse transformation semigroup of size 31, degree 6 with 5 
   generators>
  !gapprompt@gap>| !gapinput@T := FreeSemilattice(IsPartialPermSemigroup, 3);|
  <inverse partial perm semigroup of size 7, rank 3 with 3 generators>
  !gapprompt@gap>| !gapinput@U := FreeSemilattice(IsBooleanMatSemigroup, 4);|
  <inverse semigroup of size 15, 5x5 boolean matrices with 4 generators>
\end{Verbatim}
 }

 

\subsection{\textcolor{Chapter }{ZeroSemigroup}}
\logpage{[ 7, 8, 5 ]}\nobreak
\hyperdef{L}{X801FC1D97D832A6F}{}
{\noindent\textcolor{FuncColor}{$\triangleright$\enspace\texttt{ZeroSemigroup({\mdseries\slshape [filt, ]n})\index{ZeroSemigroup@\texttt{ZeroSemigroup}}
\label{ZeroSemigroup}
}\hfill{\scriptsize (function)}}\\
\textbf{\indent Returns:\ }
 A zero semigroup of order \mbox{\texttt{\mdseries\slshape n}}. 



 If \mbox{\texttt{\mdseries\slshape n}} is a positive integer, then this function returns a zero semigroup of order \mbox{\texttt{\mdseries\slshape n}} in the representation given by the filter \mbox{\texttt{\mdseries\slshape filt}}. 

 The optional argument \mbox{\texttt{\mdseries\slshape filt}} may be one of the following: 
\begin{itemize}
\item \texttt{IsTransformationSemigroup} (the default, if \mbox{\texttt{\mdseries\slshape filt}} is not specified),
\item \texttt{IsPartialPermSemigroup},
\item \texttt{IsBipartitionSemigroup},
\item \texttt{IsBlockBijectionSemigroup},
\item \texttt{IsPBRSemigroup},
\item \texttt{IsBooleanMatSemigroup},
\item \texttt{IsReesZeroMatrixSemigroup} (provided that \texttt{\mbox{\texttt{\mdseries\slshape n}} {\textgreater} 1}).
\end{itemize}
 See \texttt{IsZeroSemigroup} (\ref{IsZeroSemigroup}) for more information about zero semigroups. 
\begin{Verbatim}[commandchars=!@|,fontsize=\small,frame=single,label=Example]
  !gapprompt@gap>| !gapinput@S := ZeroSemigroup(5);|
  <commutative non-regular transformation semigroup of size 5, degree 5
   with 4 generators>
  !gapprompt@gap>| !gapinput@IsZeroSemigroup(S);|
  true
  !gapprompt@gap>| !gapinput@S := ZeroSemigroup(IsPartialPermSemigroup, 15);|
  <commutative non-regular partial perm semigroup of size 15, rank 14
   with 14 generators>
  !gapprompt@gap>| !gapinput@Size(S);|
  15
  !gapprompt@gap>| !gapinput@z := MultiplicativeZero(S);|
  <empty partial perm>
  !gapprompt@gap>| !gapinput@IsZeroSemigroup(S);|
  true
  !gapprompt@gap>| !gapinput@ForAll(S, x -> ForAll(S, y -> x * y = z));|
  true
\end{Verbatim}
 }

 

\subsection{\textcolor{Chapter }{LeftZeroSemigroup}}
\logpage{[ 7, 8, 6 ]}\nobreak
\hyperdef{L}{X8672CFA47CA620B2}{}
{\noindent\textcolor{FuncColor}{$\triangleright$\enspace\texttt{LeftZeroSemigroup({\mdseries\slshape [filt, ]n})\index{LeftZeroSemigroup@\texttt{LeftZeroSemigroup}}
\label{LeftZeroSemigroup}
}\hfill{\scriptsize (function)}}\\
\noindent\textcolor{FuncColor}{$\triangleright$\enspace\texttt{RightZeroSemigroup({\mdseries\slshape [filt, ]n})\index{RightZeroSemigroup@\texttt{RightZeroSemigroup}}
\label{RightZeroSemigroup}
}\hfill{\scriptsize (function)}}\\
\textbf{\indent Returns:\ }
 A left zero (or right zero) semigroup of order \mbox{\texttt{\mdseries\slshape n}}. 



 If \mbox{\texttt{\mdseries\slshape n}} is a positive integer, then this function returns a left zero (or right zero,
as appropriate) semigroup of order \mbox{\texttt{\mdseries\slshape n}} in the representation given by the filter \mbox{\texttt{\mdseries\slshape filt}}. If \mbox{\texttt{\mdseries\slshape filt}} is not specified then the default representation is \texttt{IsTransformationSemigroup}. 

 The function \texttt{LeftZeroSemigroup([\mbox{\texttt{\mdseries\slshape filt}},] \mbox{\texttt{\mdseries\slshape n}})} simply calls \texttt{RectangularBand([\mbox{\texttt{\mdseries\slshape filt}},] \mbox{\texttt{\mdseries\slshape n}}, 1)} and the function \texttt{RightZeroSemigroup([\mbox{\texttt{\mdseries\slshape filt}},] \mbox{\texttt{\mdseries\slshape n}})} simply calls \texttt{RectangularBand([\mbox{\texttt{\mdseries\slshape filt}},] 1, \mbox{\texttt{\mdseries\slshape n}})}. 

 For more information about \texttt{RectangularBand}, including its permitted values of \mbox{\texttt{\mdseries\slshape filt}}, see \texttt{RectangularBand} (\ref{RectangularBand}). See \texttt{IsLeftZeroSemigroup} (\ref{IsLeftZeroSemigroup}) and \texttt{IsRightZeroSemigroup} (\ref{IsRightZeroSemigroup}) for more information about left zero and right zero semigroups. 
\begin{Verbatim}[commandchars=!@|,fontsize=\small,frame=single,label=Example]
  !gapprompt@gap>| !gapinput@S := LeftZeroSemigroup(20);|
  <transformation semigroup of degree 6 with 20 generators>
  !gapprompt@gap>| !gapinput@IsLeftZeroSemigroup(S);|
  true
  !gapprompt@gap>| !gapinput@ForAll(Tuples(S, 2), p -> p[1] * p[2] = p[1]);|
  true
  !gapprompt@gap>| !gapinput@S := RightZeroSemigroup(IsBipartitionSemigroup, 5);|
  <regular bipartition semigroup of size 5, degree 3 with 5 generators>
  !gapprompt@gap>| !gapinput@IsRightZeroSemigroup(S);|
  true
\end{Verbatim}
 }

 

\subsection{\textcolor{Chapter }{BrandtSemigroup}}
\logpage{[ 7, 8, 7 ]}\nobreak
\hyperdef{L}{X7E2B20C77D47F7FB}{}
{\noindent\textcolor{FuncColor}{$\triangleright$\enspace\texttt{BrandtSemigroup({\mdseries\slshape [[filt, ]G, ]n})\index{BrandtSemigroup@\texttt{BrandtSemigroup}}
\label{BrandtSemigroup}
}\hfill{\scriptsize (function)}}\\
\textbf{\indent Returns:\ }
 An \mbox{\texttt{\mdseries\slshape n}} by \mbox{\texttt{\mdseries\slshape n}} Brandt semigroup over the group \mbox{\texttt{\mdseries\slshape G}}. 



 If \mbox{\texttt{\mdseries\slshape n}} is a positive integer, then this function returns an \mbox{\texttt{\mdseries\slshape n}} by \mbox{\texttt{\mdseries\slshape n}} Brandt semigroup over the group \mbox{\texttt{\mdseries\slshape G}} in the representation given by the filter \mbox{\texttt{\mdseries\slshape filt}}. 

 The optional argument \mbox{\texttt{\mdseries\slshape filt}} can be any of the following: 
\begin{itemize}
\item \texttt{IsPartialPermSemigroup} (the default, if \mbox{\texttt{\mdseries\slshape filt}} is not specified),
\item \texttt{IsReesZeroMatrixSemigroup},
\item \texttt{IsTransformationSemigroup},
\item \texttt{IsBipartitionSemigroup},
\item \texttt{IsPBRSemigroup},
\item \texttt{IsBooleanMatSemigroup},
\item \texttt{IsNTPMatrixSemigroup},
\item \texttt{IsMaxPlusMatrixSemigroup},
\item \texttt{IsMinPlusMatrixSemigroup},
\item \texttt{IsTropicalMaxPlusMatrixSemigroup},
\item \texttt{IsTropicalMinPlusMatrixSemigroup},
\item \texttt{IsProjectiveMaxPlusMatrixSemigroup},
\item \texttt{IsIntegerMatrixSemigroup.}
\end{itemize}
 The optional argument \mbox{\texttt{\mdseries\slshape G}} defaults to a trivial permutation group. If present \mbox{\texttt{\mdseries\slshape G}} must be a permutation group, unless \mbox{\texttt{\mdseries\slshape filt}} is \texttt{IsReesZeroMatrixSemigroup} when \mbox{\texttt{\mdseries\slshape G}} may be any type of finite group. 

 See \texttt{IsBrandtSemigroup} (\ref{IsBrandtSemigroup}) for more information about Brandt semigroups. 
\begin{Verbatim}[commandchars=!@|,fontsize=\small,frame=single,label=Example]
  !gapprompt@gap>| !gapinput@S := BrandtSemigroup(5);|
  <0-simple inverse partial perm semigroup of rank 5 with 4 generators>
  !gapprompt@gap>| !gapinput@IsBrandtSemigroup(S);|
  true
  !gapprompt@gap>| !gapinput@S := BrandtSemigroup(IsTransformationSemigroup, 15);|
  <0-simple transformation semigroup of degree 16 with 28 generators>
  !gapprompt@gap>| !gapinput@Size(S);|
  226
  !gapprompt@gap>| !gapinput@MultiplicativeZero(S);|
  Transformation( [ 16, 16, 16, 16, 16, 16, 16, 16, 16, 16, 16, 16, 16,
    16, 16, 16 ] )
  !gapprompt@gap>| !gapinput@S := BrandtSemigroup(Group((1, 2)), 3);|
  <0-simple inverse partial perm semigroup of rank 6 with 3 generators>
  !gapprompt@gap>| !gapinput@S := BrandtSemigroup(IsTransformationSemigroup, Group((1, 2)), 3);|
  <0-simple transformation semigroup of degree 7 with 5 generators>
  !gapprompt@gap>| !gapinput@S := BrandtSemigroup(IsReesZeroMatrixSemigroup,|
  !gapprompt@>| !gapinput@                        DihedralGroup(4),|
  !gapprompt@>| !gapinput@                        2);|
  <Rees 0-matrix semigroup 2x2 over <pc group of size 4 with
   2 generators>>
\end{Verbatim}
 }

 }

   
\section{\textcolor{Chapter }{ Free bands }}\label{Free bands}
\logpage{[ 7, 9, 0 ]}
\hyperdef{L}{X7BB29A6779E8066A}{}
{
  This chapter describes the functions in \textsf{Semigroups} for dealing with free bands. This part of the manual and the functions
described herein were originally written by Julius Jonu{\v s}as, with later
additions by Reinis Cirpons, Tom Conti\texttt{\symbol{45}}Leslie, and Murray
Whyte

 A semigroup $B$ is a \emph{free band} on a non\texttt{\symbol{45}}empty set $X$ if $B$ is a band with a map $ f $ from $ B $ to $X$ such that for every band $ S $ and every map $ g $ from $X$ to $ B $ there exists a unique homomorphism $ g'$ from $B$ to $S$ such that $fg' = g$. The free band on a set $X$ is unique up to isomorphism. Moreover, by the universal property, every band
can be expressed as a quotient of a free band.

 For an alternative description of a free band. Suppose that $ X $ is a non\texttt{\symbol{45}}empty set and $ X ^ + $ a free semigroup on $ X $. Also suppose that $ b $ is the smallest congurance on $ X ^ + $ containing the set 
\[ \{(w ^ 2, w) : w \in X ^ + \}. \]
 Then the free band on $ X $ is isomorphic to the quotient of $ X ^ + $ by $ b $. See Section 4.5 of \cite{Howie1995aa} for more information on free bands. 

\subsection{\textcolor{Chapter }{FreeBand (for a given rank)}}
\logpage{[ 7, 9, 1 ]}\nobreak
\hyperdef{L}{X7B2A65F382DB36EC}{}
{\noindent\textcolor{FuncColor}{$\triangleright$\enspace\texttt{FreeBand({\mdseries\slshape rank[, name]})\index{FreeBand@\texttt{FreeBand}!for a given rank}
\label{FreeBand:for a given rank}
}\hfill{\scriptsize (function)}}\\
\noindent\textcolor{FuncColor}{$\triangleright$\enspace\texttt{FreeBand({\mdseries\slshape name1, name2, ...})\index{FreeBand@\texttt{FreeBand}!for a list of distinct names}
\label{FreeBand:for a list of distinct names}
}\hfill{\scriptsize (function)}}\\
\noindent\textcolor{FuncColor}{$\triangleright$\enspace\texttt{FreeBand({\mdseries\slshape names})\index{FreeBand@\texttt{FreeBand}!for various distinct names}
\label{FreeBand:for various distinct names}
}\hfill{\scriptsize (function)}}\\
\textbf{\indent Returns:\ }
 A free band.



 Returns a free band on \mbox{\texttt{\mdseries\slshape rank}} generators, for a positive integer \mbox{\texttt{\mdseries\slshape rank}}. If \mbox{\texttt{\mdseries\slshape rank}} is not specified, the number of \mbox{\texttt{\mdseries\slshape names}} is used. In the second and third forms, the \mbox{\texttt{\mdseries\slshape names}} must be distinct. The resulting semigroup is always finite. 
\begin{Verbatim}[commandchars=!@|,fontsize=\small,frame=single,label=Example]
  !gapprompt@gap>| !gapinput@FreeBand(6);|
  <free band on the generators [ x1, x2, x3, x4, x5, x6 ]>
  !gapprompt@gap>| !gapinput@FreeBand(6, "b");|
  <free band on the generators [ b1, b2, b3, b4, b5, b6 ]>
  !gapprompt@gap>| !gapinput@FreeBand("a", "b", "c");|
  <free band on the generators [ a, b, c ]>
  !gapprompt@gap>| !gapinput@FreeBand("a", "b", "c");|
  <free band on the generators [ a, b, c ]>
  !gapprompt@gap>| !gapinput@S := FreeBand(["a", "b", "c"]);|
  <free band on the generators [ a, b, c ]>
  !gapprompt@gap>| !gapinput@Size(S);|
  159
  !gapprompt@gap>| !gapinput@gens := Generators(S);|
  [ a, b, c ]
  !gapprompt@gap>| !gapinput@S.1 * S.2;|
  ab
\end{Verbatim}
 }

 

\subsection{\textcolor{Chapter }{IsFreeBandCategory}}
\logpage{[ 7, 9, 2 ]}\nobreak
\hyperdef{L}{X7F5658DC7E56C4A6}{}
{\noindent\textcolor{FuncColor}{$\triangleright$\enspace\texttt{IsFreeBandCategory\index{IsFreeBandCategory@\texttt{IsFreeBandCategory}}
\label{IsFreeBandCategory}
}\hfill{\scriptsize (Category)}}\\


 \texttt{IsFreeBandCategory} is the category of semigroups created using \texttt{FreeBand} (\ref{FreeBand:for a given rank}). 
\begin{Verbatim}[commandchars=!@|,fontsize=\small,frame=single,label=Example]
  !gapprompt@gap>| !gapinput@IsFreeBandCategory(FreeBand(3));|
  true
  !gapprompt@gap>| !gapinput@IsFreeBand(SymmetricGroup(6));|
  false
\end{Verbatim}
 }

 

\subsection{\textcolor{Chapter }{IsFreeBand (for a given semigroup)}}
\logpage{[ 7, 9, 3 ]}\nobreak
\hyperdef{L}{X7B1CD5FC7E034B88}{}
{\noindent\textcolor{FuncColor}{$\triangleright$\enspace\texttt{IsFreeBand({\mdseries\slshape S})\index{IsFreeBand@\texttt{IsFreeBand}!for a given semigroup}
\label{IsFreeBand:for a given semigroup}
}\hfill{\scriptsize (property)}}\\
\textbf{\indent Returns:\ }
\texttt{true} or \texttt{false}. 



 \texttt{IsFreeBand} returns \texttt{true} if the given semigroup \mbox{\texttt{\mdseries\slshape S}} is a free band. 
\begin{Verbatim}[commandchars=!@|,fontsize=\small,frame=single,label=Example]
  !gapprompt@gap>| !gapinput@IsFreeBand(FreeBand(3));|
  true
  !gapprompt@gap>| !gapinput@IsFreeBand(SymmetricGroup(6));|
  false
  !gapprompt@gap>| !gapinput@IsFreeBand(FullTransformationMonoid(7));|
  false
\end{Verbatim}
 }

 

\subsection{\textcolor{Chapter }{IsFreeBandElement}}
\logpage{[ 7, 9, 4 ]}\nobreak
\hyperdef{L}{X7DECF69087BB3B16}{}
{\noindent\textcolor{FuncColor}{$\triangleright$\enspace\texttt{IsFreeBandElement\index{IsFreeBandElement@\texttt{IsFreeBandElement}}
\label{IsFreeBandElement}
}\hfill{\scriptsize (Category)}}\\


 \texttt{IsFreeBandElement} is a \texttt{Category} containing the elements of a free band. 
\begin{Verbatim}[commandchars=!@|,fontsize=\small,frame=single,label=Example]
  !gapprompt@gap>| !gapinput@IsFreeBandElement(Generators(FreeBand(4))[1]);|
  true
  !gapprompt@gap>| !gapinput@IsFreeBandElement(Transformation([1, 3, 4, 1]));|
  false
  !gapprompt@gap>| !gapinput@IsFreeBandElement((1, 2, 3, 4));|
  false
\end{Verbatim}
 }

 

\subsection{\textcolor{Chapter }{IsFreeBandElementCollection}}
\logpage{[ 7, 9, 5 ]}\nobreak
\hyperdef{L}{X842839C87DAAA43C}{}
{\noindent\textcolor{FuncColor}{$\triangleright$\enspace\texttt{IsFreeBandElementCollection\index{IsFreeBandElementCollection@\texttt{IsFreeBandElementCollection}}
\label{IsFreeBandElementCollection}
}\hfill{\scriptsize (Category)}}\\


 Every collection of elements of a free band belongs to the category \texttt{IsFreeBandElementCollection}. For example, every free band belongs to \texttt{IsFreeBandElementCollection}. }

 

\subsection{\textcolor{Chapter }{IsFreeBandSubsemigroup}}
\logpage{[ 7, 9, 6 ]}\nobreak
\hyperdef{L}{X7AEF4CD1857E7DCC}{}
{\noindent\textcolor{FuncColor}{$\triangleright$\enspace\texttt{IsFreeBandSubsemigroup\index{IsFreeBandSubsemigroup@\texttt{IsFreeBandSubsemigroup}}
\label{IsFreeBandSubsemigroup}
}\hfill{\scriptsize (filter)}}\\


 \texttt{IsFreeBandSubsemigroup} is a synonym for \texttt{IsSemigroup} and \texttt{IsFreeBandElementCollection}. 
\begin{Verbatim}[commandchars=!@|,fontsize=\small,frame=single,label=Example]
  !gapprompt@gap>| !gapinput@S := FreeBand(2);|
  <free band on the generators [ x1, x2 ]>
  !gapprompt@gap>| !gapinput@x := S.1;|
  x1
  !gapprompt@gap>| !gapinput@y := S.2;|
  x2
  !gapprompt@gap>| !gapinput@new := Semigroup([x * y, x]);|
  <semigroup with 2 generators>
  !gapprompt@gap>| !gapinput@IsFreeBand(new);|
  false
  !gapprompt@gap>| !gapinput@IsFreeBandSubsemigroup(new);|
  true
\end{Verbatim}
 }

 

\subsection{\textcolor{Chapter }{ContentOfFreeBandElement}}
\logpage{[ 7, 9, 7 ]}\nobreak
\hyperdef{L}{X808CAEC17BF271D1}{}
{\noindent\textcolor{FuncColor}{$\triangleright$\enspace\texttt{ContentOfFreeBandElement({\mdseries\slshape x})\index{ContentOfFreeBandElement@\texttt{ContentOfFreeBandElement}}
\label{ContentOfFreeBandElement}
}\hfill{\scriptsize (attribute)}}\\
\noindent\textcolor{FuncColor}{$\triangleright$\enspace\texttt{ContentOfFreeBandElementCollection({\mdseries\slshape coll})\index{ContentOfFreeBandElementCollection@\texttt{ContentOfFreeBandElementCollection}}
\label{ContentOfFreeBandElementCollection}
}\hfill{\scriptsize (attribute)}}\\
\textbf{\indent Returns:\ }
 A list of integers 



 The content of a free band element \mbox{\texttt{\mdseries\slshape x}} is the set of generators appearing in the word representing the element \mbox{\texttt{\mdseries\slshape x}} of the free band. 

 The function \texttt{ContentOfFreeBandElement} returns the content of free band element \mbox{\texttt{\mdseries\slshape x}} represented as a list of integers, where \texttt{1} represents the first generator, \texttt{2} the second generator, and so on. 

 The function \texttt{ContentOfFreeBandElementCollection} returns the the least list \texttt{C} for the collection of free band elements \mbox{\texttt{\mdseries\slshape coll}} such that the content of every element in \mbox{\texttt{\mdseries\slshape coll}} is contained in \texttt{C}. 
\begin{Verbatim}[commandchars=!@|,fontsize=\small,frame=single,label=Example]
  !gapprompt@gap>| !gapinput@S := FreeBand(2);|
  <free band on the generators [ x1, x2 ]>
  !gapprompt@gap>| !gapinput@x := S.1;|
  x1
  !gapprompt@gap>| !gapinput@y := S.2;|
  x2
  !gapprompt@gap>| !gapinput@ContentOfFreeBandElement(x);|
  [ 1 ]
  !gapprompt@gap>| !gapinput@ContentOfFreeBandElement(x * y);|
  [ 1, 2 ]
  !gapprompt@gap>| !gapinput@ContentOfFreeBandElement(x * y * x);|
  [ 1, 2 ]
  !gapprompt@gap>| !gapinput@ContentOfFreeBandElementCollection([x, y]);|
  [ 1, 2 ]
\end{Verbatim}
 }

 

\subsection{\textcolor{Chapter }{EqualInFreeBand}}
\logpage{[ 7, 9, 8 ]}\nobreak
\hyperdef{L}{X7CD9426180587CA4}{}
{\noindent\textcolor{FuncColor}{$\triangleright$\enspace\texttt{EqualInFreeBand({\mdseries\slshape u, v})\index{EqualInFreeBand@\texttt{EqualInFreeBand}}
\label{EqualInFreeBand}
}\hfill{\scriptsize (operation)}}\\


 This operation takes a pair \mbox{\texttt{\mdseries\slshape u}} and \mbox{\texttt{\mdseries\slshape v}} of lists of positive integers or strings, representing words in a free
semigroup.

 Where \texttt{F} is a free band over some alphabet containing the letters occurring in \mbox{\texttt{\mdseries\slshape u}} and \mbox{\texttt{\mdseries\slshape v}}, this operation returns \texttt{true} if \mbox{\texttt{\mdseries\slshape u}} and \mbox{\texttt{\mdseries\slshape v}} are equal in \texttt{F}, and \texttt{false} otherwise.

 Note that this operation is for lists and strings, as opposed to \texttt{FreeBandElement} objects.

 This is an implementation of an algorithm described by Jakub Radoszewski and
Wojciech Rytter in \cite{Radoszewski2010aa}. 
\begin{Verbatim}[commandchars=!@|,fontsize=\small,frame=single,label=Example]
  !gapprompt@gap>| !gapinput@EqualInFreeBand("aa", "a");|
  true
  !gapprompt@gap>| !gapinput@EqualInFreeBand("abcacba", "abcba");|
  true
  !gapprompt@gap>| !gapinput@EqualInFreeBand("aab", "aac");|
  false
  !gapprompt@gap>| !gapinput@EqualInFreeBand([1, 3, 3], [2]);|
  false
\end{Verbatim}
 }

 

\subsection{\textcolor{Chapter }{GreensDClassOfElement (for a free band and element)}}
\logpage{[ 7, 9, 9 ]}\nobreak
\hyperdef{L}{X85DC5D50875E55D6}{}
{\noindent\textcolor{FuncColor}{$\triangleright$\enspace\texttt{GreensDClassOfElement({\mdseries\slshape S, x})\index{GreensDClassOfElement@\texttt{GreensDClassOfElement}!for a free band and element}
\label{GreensDClassOfElement:for a free band and element}
}\hfill{\scriptsize (operation)}}\\
\textbf{\indent Returns:\ }
 A Green's $\mathcal{D}$\texttt{\symbol{45}}class 



 Let \mbox{\texttt{\mdseries\slshape S}} be a free band. Two elements of \mbox{\texttt{\mdseries\slshape  S }} are $\mathcal{D}$\texttt{\symbol{45}}related if and only if they have the same content i.e. the
set of generators appearing in any factorization of the elements. Therefore, a $\mathcal{D}$\texttt{\symbol{45}}class of a free band element \mbox{\texttt{\mdseries\slshape  x }} is the set of elements of \mbox{\texttt{\mdseries\slshape  S }} which have the same content as \mbox{\texttt{\mdseries\slshape  x }}. 
\begin{Verbatim}[commandchars=!@|,fontsize=\small,frame=single,label=Example]
  !gapprompt@gap>| !gapinput@S := FreeBand(3, "b");|
  <free band on the generators [ b1, b2, b3 ]>
  !gapprompt@gap>| !gapinput@x := S.1 * S.2;|
  b1b2
  !gapprompt@gap>| !gapinput@D := GreensDClassOfElement(S, x);|
  <Green's D-class: b1b2>
  !gapprompt@gap>| !gapinput@IsGreensDClass(D);|
  true
\end{Verbatim}
 }

 
\subsection{\textcolor{Chapter }{Operators}}\logpage{[ 7, 9, 10 ]}
\hyperdef{L}{X7AD6F77E7D95C996}{}
{
  The following operators are also included for free band elements: 
\begin{description}
\item[{\texttt{\mbox{\texttt{\mdseries\slshape u}} * \mbox{\texttt{\mdseries\slshape v}}}}]  returns the product of two free band elements \mbox{\texttt{\mdseries\slshape u}} and \mbox{\texttt{\mdseries\slshape v}}. 
\item[{\texttt{\mbox{\texttt{\mdseries\slshape u}} = \mbox{\texttt{\mdseries\slshape v}} }}]  checks if two free band elements are equal. 
\item[{\texttt{\mbox{\texttt{\mdseries\slshape u}} {\textless} \mbox{\texttt{\mdseries\slshape v}} }}]  compares the sizes of the internal representations of two free band elements. 
\end{description}
 }

 }

   
\section{\textcolor{Chapter }{ Graph inverse semigroups }}\label{Graph inverse semigroups}
\logpage{[ 7, 10, 0 ]}
\hyperdef{L}{X850B10D783053100}{}
{
  In this chapter we describe a class of semigroups arising from directed
graphs. 

 The functionality in \textsf{Semigroups} for graph inverse semigroups was written jointly by Zak Mesyan (UCCS) and J.
D. Mitchell (St Andrews). The functionality for graph inverse semigroup
congruences was written by Marina Anagnostopoulou\texttt{\symbol{45}}Merkouri
(St Andrews). 

 

\subsection{\textcolor{Chapter }{GraphInverseSemigroup}}
\logpage{[ 7, 10, 1 ]}\nobreak
\hyperdef{L}{X7A9EEFD386D6F630}{}
{\noindent\textcolor{FuncColor}{$\triangleright$\enspace\texttt{GraphInverseSemigroup({\mdseries\slshape E})\index{GraphInverseSemigroup@\texttt{GraphInverseSemigroup}}
\label{GraphInverseSemigroup}
}\hfill{\scriptsize (operation)}}\\
\textbf{\indent Returns:\ }
A graph inverse semigroup.



 If \mbox{\texttt{\mdseries\slshape E}} is a digraph (i.e. it satisfies \texttt{IsDigraph} (\textbf{Digraphs: IsDigraph})), then \texttt{GraphInverseSemigroup} returns the graph inverse semigroup $G(\mbox{\texttt{\mdseries\slshape E}})$ where, roughly speaking, elements correspond to paths in the graph \mbox{\texttt{\mdseries\slshape E}}. 

 Let us describe \mbox{\texttt{\mdseries\slshape E}} as a digraph $\mbox{\texttt{\mdseries\slshape E}} = (E ^ 0, E ^ 1, r, s)$, where $E^0$ is the set of vertices, $E^1$ is the set of edges, and $r$ and $s$ are functions $E^1 \to E^0$ giving the \emph{range} and \emph{source} of an edge, respectively. The \emph{graph inverse semigroup $G(\mbox{\texttt{\mdseries\slshape E}})$ of $E$} is the semigroup\texttt{\symbol{45}}with\texttt{\symbol{45}}zero generated by
the sets $\mbox{\texttt{\mdseries\slshape E}} ^ 0$ and $\mbox{\texttt{\mdseries\slshape E}} ^ 1$, together with a set of variables $\{e ^ {-1} \mid e \in \mbox{\texttt{\mdseries\slshape E}} ^ 1\}$, satisfying the following relations for all $v, w \in \mbox{\texttt{\mdseries\slshape E}} ^ 0$ and $e, f \in \mbox{\texttt{\mdseries\slshape E}} ^ 1$: 
\begin{description}
\item[{(V)}] $vw = \delta_{v,w} \cdot v$,
\item[{(E1)}] $s(e) \cdot e=e \cdot r(e)=e$,
\item[{(E2)}] $r(e) \cdot e^{-1} = e^{-1} \cdot s(e) =e^{-1}$,
\item[{(CK1)}] $e^{-1} \cdot f = \delta_{e,f} \cdot r(e)$.
\end{description}
 (Here $\delta$ is the Kronecker delta.) We define $v^{-1}=v$ for each $v \in E^0$, and for any path $y=e_1\dots e_n$ ($e_1\dots e_n \in E^1$) we let $y^{-1} = e_n^{-1} \dots e_1^{-1}$. With this notation, every nonzero element of $G(E)$ can be written uniquely as $xy^{-1}$ for some paths $x, y$ in $E$, by the CK1 relation. 

 For a more complete description, see \cite{Mesyan2016}. 

 
\begin{Verbatim}[commandchars=!@|,fontsize=\small,frame=single,label=Example]
  !gapprompt@gap>| !gapinput@gr := Digraph([[2, 5, 8, 10], [2, 3, 4, 5, 6, 8, 9, 10], [1],|
  !gapprompt@>| !gapinput@                  [3, 5, 7, 8, 10], [2, 5, 7], [3, 6, 7, 9, 10],|
  !gapprompt@>| !gapinput@                  [1, 4], [1, 5, 9], [1, 2, 7, 8], [3, 5]]);|
  <immutable digraph with 10 vertices, 37 edges>
  !gapprompt@gap>| !gapinput@S := GraphInverseSemigroup(gr);|
  <infinite graph inverse semigroup with 10 vertices, 37 edges>
  !gapprompt@gap>| !gapinput@GeneratorsOfInverseSemigroup(S);|
  [ e_1, e_2, e_3, e_4, e_5, e_6, e_7, e_8, e_9, e_10, e_11, e_12,
    e_13, e_14, e_15, e_16, e_17, e_18, e_19, e_20, e_21, e_22, e_23,
    e_24, e_25, e_26, e_27, e_28, e_29, e_30, e_31, e_32, e_33, e_34,
    e_35, e_36, e_37, v_1, v_2, v_3, v_4, v_5, v_6, v_7, v_8, v_9, v_10
   ]
  !gapprompt@gap>| !gapinput@AssignGeneratorVariables(S);|
  !gapprompt@gap>| !gapinput@e_1 * e_1 ^ -1;|
  e_1e_1^-1
  !gapprompt@gap>| !gapinput@e_1 ^ -1 * e_1 ^ -1;|
  0
  !gapprompt@gap>| !gapinput@e_1 ^ -1 * e_1;|
  v_2
\end{Verbatim}
 }

 

\subsection{\textcolor{Chapter }{Range (for a graph inverse semigroup element)}}
\logpage{[ 7, 10, 2 ]}\nobreak
\hyperdef{L}{X8187F0FF784A82CD}{}
{\noindent\textcolor{FuncColor}{$\triangleright$\enspace\texttt{Range({\mdseries\slshape x})\index{Range@\texttt{Range}!for a graph inverse semigroup element}
\label{Range:for a graph inverse semigroup element}
}\hfill{\scriptsize (attribute)}}\\
\noindent\textcolor{FuncColor}{$\triangleright$\enspace\texttt{Source({\mdseries\slshape x})\index{Source@\texttt{Source}!for a graph inverse semigroup element}
\label{Source:for a graph inverse semigroup element}
}\hfill{\scriptsize (attribute)}}\\
\textbf{\indent Returns:\ }
A graph inverse semigroup element.



 If \mbox{\texttt{\mdseries\slshape x}} is an element of a graph inverse semigroup (i.e. it satisfies \texttt{IsGraphInverseSemigroupElement} (\ref{IsGraphInverseSemigroupElement})), then \texttt{Range} and \texttt{Source} give, respectively, the start and end vertices of \mbox{\texttt{\mdseries\slshape x}} when viewed as a path in the digraph over which the semigroup is defined. 

 For a fuller description, see \texttt{GraphInverseSemigroup} (\ref{GraphInverseSemigroup}). 

 
\begin{Verbatim}[commandchars=!@|,fontsize=\small,frame=single,label=Example]
  !gapprompt@gap>| !gapinput@gr := Digraph([[], [1], [3]]);;|
  !gapprompt@gap>| !gapinput@S := GraphInverseSemigroup(gr);;|
  !gapprompt@gap>| !gapinput@e := S.1;|
  e_1
  !gapprompt@gap>| !gapinput@Source(e);|
  v_2
  !gapprompt@gap>| !gapinput@Range(e);|
  v_1
\end{Verbatim}
 }

 

\subsection{\textcolor{Chapter }{IsVertex (for a graph inverse semigroup element)}}
\logpage{[ 7, 10, 3 ]}\nobreak
\hyperdef{L}{X7DEE927C83D4DFDD}{}
{\noindent\textcolor{FuncColor}{$\triangleright$\enspace\texttt{IsVertex({\mdseries\slshape x})\index{IsVertex@\texttt{IsVertex}!for a graph inverse semigroup element}
\label{IsVertex:for a graph inverse semigroup element}
}\hfill{\scriptsize (operation)}}\\
\textbf{\indent Returns:\ }
\texttt{true} or \texttt{false}.



 If \mbox{\texttt{\mdseries\slshape x}} is an element of a graph inverse semigroup (i.e. it satisfies \texttt{IsGraphInverseSemigroupElement} (\ref{IsGraphInverseSemigroupElement})), then this attribute returns \texttt{true} if \mbox{\texttt{\mdseries\slshape x}} corresponds to a vertex in the digraph over which the semigroup is defined,
and \texttt{false} otherwise. 

 For a fuller description, see \texttt{GraphInverseSemigroup} (\ref{GraphInverseSemigroup}). 

 
\begin{Verbatim}[commandchars=!@|,fontsize=\small,frame=single,label=Example]
  !gapprompt@gap>| !gapinput@gr := Digraph([[], [1], [3]]);;|
  !gapprompt@gap>| !gapinput@S := GraphInverseSemigroup(gr);;|
  !gapprompt@gap>| !gapinput@e := S.1;|
  e_1
  !gapprompt@gap>| !gapinput@IsVertex(e);|
  false
  !gapprompt@gap>| !gapinput@v := S.3;|
  v_1
  !gapprompt@gap>| !gapinput@IsVertex(v);|
  true
  !gapprompt@gap>| !gapinput@z := v * e;|
  0
  !gapprompt@gap>| !gapinput@IsVertex(z);|
  false
\end{Verbatim}
 }

 

\subsection{\textcolor{Chapter }{IsGraphInverseSemigroup}}
\logpage{[ 7, 10, 4 ]}\nobreak
\hyperdef{L}{X7BFDF88B799B05A0}{}
{\noindent\textcolor{FuncColor}{$\triangleright$\enspace\texttt{IsGraphInverseSemigroup({\mdseries\slshape x})\index{IsGraphInverseSemigroup@\texttt{IsGraphInverseSemigroup}}
\label{IsGraphInverseSemigroup}
}\hfill{\scriptsize (filter)}}\\
\noindent\textcolor{FuncColor}{$\triangleright$\enspace\texttt{IsGraphInverseSemigroupElement({\mdseries\slshape x})\index{IsGraphInverseSemigroupElement@\texttt{IsGraphInverseSemigroupElement}}
\label{IsGraphInverseSemigroupElement}
}\hfill{\scriptsize (filter)}}\\
\textbf{\indent Returns:\ }
\texttt{true} or \texttt{false}.



 The category \texttt{IsGraphInverseSemigroup} contains any semigroup defined over a digraph using the \texttt{GraphInverseSemigroup} (\ref{GraphInverseSemigroup}) operation. The category \texttt{IsGraphInverseSemigroupElement} contains any element contained in such a semigroup. 

 
\begin{Verbatim}[commandchars=!@|,fontsize=\small,frame=single,label=Example]
  !gapprompt@gap>| !gapinput@gr := Digraph([[], [1], [3]]);;|
  !gapprompt@gap>| !gapinput@S := GraphInverseSemigroup(gr);|
  <infinite graph inverse semigroup with 3 vertices, 2 edges>
  !gapprompt@gap>| !gapinput@IsGraphInverseSemigroup(S);|
  true
  !gapprompt@gap>| !gapinput@x := GeneratorsOfSemigroup(S)[1];|
  e_1
  !gapprompt@gap>| !gapinput@IsGraphInverseSemigroupElement(x);|
  true
\end{Verbatim}
 }

 

\subsection{\textcolor{Chapter }{GraphOfGraphInverseSemigroup}}
\logpage{[ 7, 10, 5 ]}\nobreak
\hyperdef{L}{X7BE287A385A058BC}{}
{\noindent\textcolor{FuncColor}{$\triangleright$\enspace\texttt{GraphOfGraphInverseSemigroup({\mdseries\slshape S})\index{GraphOfGraphInverseSemigroup@\texttt{GraphOfGraphInverseSemigroup}}
\label{GraphOfGraphInverseSemigroup}
}\hfill{\scriptsize (attribute)}}\\
\textbf{\indent Returns:\ }
A digraph.



 If \mbox{\texttt{\mdseries\slshape S}} is a graph inverse semigroup (i.e. it satisfies \texttt{IsGraphInverseSemigroup} (\ref{IsGraphInverseSemigroup})), then this attribute returns the original digraph over which \mbox{\texttt{\mdseries\slshape S}} was defined (most likely the argument given to \texttt{GraphInverseSemigroup} (\ref{GraphInverseSemigroup}) to create \mbox{\texttt{\mdseries\slshape S}}). 

 
\begin{Verbatim}[commandchars=!@|,fontsize=\small,frame=single,label=Example]
  !gapprompt@gap>| !gapinput@gr := Digraph([[], [1], [3]]);|
  <immutable digraph with 3 vertices, 2 edges>
  !gapprompt@gap>| !gapinput@S := GraphInverseSemigroup(gr);;|
  !gapprompt@gap>| !gapinput@GraphOfGraphInverseSemigroup(S);|
  <immutable digraph with 3 vertices, 2 edges>
\end{Verbatim}
 }

 

\subsection{\textcolor{Chapter }{IsGraphInverseSemigroupElementCollection}}
\logpage{[ 7, 10, 6 ]}\nobreak
\hyperdef{L}{X870128E4845D6ABD}{}
{\noindent\textcolor{FuncColor}{$\triangleright$\enspace\texttt{IsGraphInverseSemigroupElementCollection\index{IsGraphInverseSemigroupElementCollection@\texttt{IsGraph}\-\texttt{Inverse}\-\texttt{Semigroup}\-\texttt{Element}\-\texttt{Collection}}
\label{IsGraphInverseSemigroupElementCollection}
}\hfill{\scriptsize (Category)}}\\


 Every collection of elements of a graph inverse semigroup belongs to the
category \texttt{IsGraphInverseSemigroupElementCollection}. For example, every graph inverse semigroup belongs to \texttt{IsGraphInverseSemigroupElementCollection}. }

 

\subsection{\textcolor{Chapter }{IsGraphInverseSubsemigroup}}
\logpage{[ 7, 10, 7 ]}\nobreak
\hyperdef{L}{X7BC6D5107ED09DBA}{}
{\noindent\textcolor{FuncColor}{$\triangleright$\enspace\texttt{IsGraphInverseSubsemigroup\index{IsGraphInverseSubsemigroup@\texttt{IsGraphInverseSubsemigroup}}
\label{IsGraphInverseSubsemigroup}
}\hfill{\scriptsize (filter)}}\\


 \texttt{IsGraphInverseSubsemigroup} is a synonym for \texttt{IsSemigroup} and \texttt{IsInverseSemigroup} and \texttt{IsGraphInverseSemigroupElementCollection}. 

 See \texttt{IsGraphInverseSemigroupElementCollection} (\ref{IsGraphInverseSemigroupElementCollection}) and \texttt{IsInverseSemigroup} (\textbf{Reference: IsInverseSemigroup}). 
\begin{Verbatim}[commandchars=!@|,fontsize=\small,frame=single,label=Example]
  !gapprompt@gap>| !gapinput@gr := Digraph([[], [1], [2]]);|
  <immutable digraph with 3 vertices, 2 edges>
  !gapprompt@gap>| !gapinput@S := GraphInverseSemigroup(gr);|
  <finite graph inverse semigroup with 3 vertices, 2 edges>
  !gapprompt@gap>| !gapinput@Elements(S);|
  [ e_2^-1, e_1^-1, e_1^-1e_2^-1, 0, e_1, e_1e_1^-1, e_1e_1^-1e_2^-1,
    e_2, e_2e_2^-1, e_2e_1, e_2e_1e_1^-1, e_2e_1e_1^-1e_2^-1, v_1, v_2,
    v_3 ]
  !gapprompt@gap>| !gapinput@T := InverseSemigroup(Elements(S){[3, 5]});;|
  !gapprompt@gap>| !gapinput@IsGraphInverseSubsemigroup(T);|
  true
\end{Verbatim}
 }

 

\subsection{\textcolor{Chapter }{VerticesOfGraphInverseSemigroup}}
\logpage{[ 7, 10, 8 ]}\nobreak
\hyperdef{L}{X7DF1ACC27CC998EB}{}
{\noindent\textcolor{FuncColor}{$\triangleright$\enspace\texttt{VerticesOfGraphInverseSemigroup({\mdseries\slshape S})\index{VerticesOfGraphInverseSemigroup@\texttt{VerticesOfGraphInverseSemigroup}}
\label{VerticesOfGraphInverseSemigroup}
}\hfill{\scriptsize (attribute)}}\\
\textbf{\indent Returns:\ }
A list.



 If \mbox{\texttt{\mdseries\slshape S}} is a graph inverse semigroup (i.e. it satisfies \texttt{IsGraphInverseSemigroup} (\ref{IsGraphInverseSemigroup})), then this attribute returns the list of vertices of \mbox{\texttt{\mdseries\slshape S}}. 
\begin{Verbatim}[commandchars=!@|,fontsize=\small,frame=single,label=Example]
  !gapprompt@gap>| !gapinput@D := Digraph([[3, 4], [3, 4], [4], []]);|
  <immutable digraph with 4 vertices, 5 edges>
  !gapprompt@gap>| !gapinput@S := GraphInverseSemigroup(D);|
  <finite graph inverse semigroup with 4 vertices, 5 edges>
  !gapprompt@gap>| !gapinput@VerticesOfGraphInverseSemigroup(S);|
  [ v_1, v_2, v_3, v_4 ]
  !gapprompt@gap>| !gapinput@D := ChainDigraph(12);|
  <immutable chain digraph with 12 vertices>
  !gapprompt@gap>| !gapinput@S := GraphInverseSemigroup(D);|
  <finite graph inverse semigroup with 12 vertices, 11 edges>
  !gapprompt@gap>| !gapinput@VerticesOfGraphInverseSemigroup(S);|
  [ v_1, v_2, v_3, v_4, v_5, v_6, v_7, v_8, v_9, v_10, v_11, v_12 ]
\end{Verbatim}
 }

 

\subsection{\textcolor{Chapter }{IndexOfVertexOfGraphInverseSemigroup}}
\logpage{[ 7, 10, 9 ]}\nobreak
\hyperdef{L}{X87500BC782212D4A}{}
{\noindent\textcolor{FuncColor}{$\triangleright$\enspace\texttt{IndexOfVertexOfGraphInverseSemigroup({\mdseries\slshape v})\index{IndexOfVertexOfGraphInverseSemigroup@\texttt{Index}\-\texttt{Of}\-\texttt{Vertex}\-\texttt{Of}\-\texttt{Graph}\-\texttt{Inverse}\-\texttt{Semigroup}}
\label{IndexOfVertexOfGraphInverseSemigroup}
}\hfill{\scriptsize (attribute)}}\\
\textbf{\indent Returns:\ }
A positive integer.



 If \mbox{\texttt{\mdseries\slshape v}} is a vertex of a graph inverse semigroup (i.e. it satisfies \texttt{IsGraphInverseSemigroup} (\ref{IsGraphInverseSemigroup})), then this attribute returns the index of this vertex in \mbox{\texttt{\mdseries\slshape S}}.

 
\begin{Verbatim}[commandchars=!@|,fontsize=\small,frame=single,label=Example]
  !gapprompt@gap>| !gapinput@D := Digraph([[3, 4], [3, 4], [4], []]);|
  <immutable digraph with 4 vertices, 5 edges>
  !gapprompt@gap>| !gapinput@S := GraphInverseSemigroup(D);|
  <finite graph inverse semigroup with 4 vertices, 5 edges>
  !gapprompt@gap>| !gapinput@IndexOfVertexOfGraphInverseSemigroup(v_1);|
  1
  !gapprompt@gap>| !gapinput@IndexOfVertexOfGraphInverseSemigroup(v_3);|
  3
\end{Verbatim}
 }

 }

   
\section{\textcolor{Chapter }{ Free inverse semigroups }}\label{Free inverse semigroups}
\logpage{[ 7, 11, 0 ]}
\hyperdef{L}{X7E51292C8755DCF2}{}
{
  This chapter describes the functions in \textsf{Semigroups} for dealing with free inverse semigroups. This part of the manual and the
functions described herein were written by Julius Jonu{\v s}as.

 An inverse semigroup $F$ is said to be \emph{free} on a non\texttt{\symbol{45}}empty set $X$ if there is a map $f$ from $F$ to $X$ such that for every inverse semigroup $S$ and a map $g$ from $X$ to $S$ there exists a unique homomorphism $g'$ from $F$ to $S$ such that $fg' = g$. Moreover, by this universal property, every inverse semigroup can be
expressed as a quotient of a free inverse semigroup. 

 The internal representation of an element of a free inverse semigroup uses a
Munn tree. A \emph{Munn tree} is a directed tree with distinguished start and terminal vertices and where
the edges are labeled by generators so that two edges labeled by the same
generator are only incident to the same vertex if one of the edges is coming
in and the other is leaving the vertex. For more information regarding free
inverse semigroups and the Munn representations see Section 5.10 of \cite{Howie1995aa}. 

 See also  (\textbf{Reference: Inverse semigroups and monoids}),  (\textbf{Reference: Partial permutations}) and  (\textbf{Reference: Free Groups, Monoids and Semigroups}). 

 An element of a free inverse semigroup in \textsf{Semigroups} is displayed, by default, as a shortest word corresponding to the element.
However, there might be more than one word of the minimum length. For example,
if $x$ and $y$ are generators of a free inverse semigroups, then 
\[xyy ^ {-1}xx ^ {-1}x ^ {-1} = xxx ^ {-1}yy ^ {-1}x ^ {-1}.\]
 See \texttt{MinimalWord} (\ref{MinimalWord:for free inverse semigroup element}). Therefore we provide a another method for printing elements of a free
inverse semigroup: a unique canonical form. Suppose an element of a free
inverse semigroup is given as a Munn tree. Let $L$ be the set of words corresponding to the shortest paths from the start vertex
to the leaves of the tree. Also let $w$ be the word corresponding to the shortest path from the start vertex to the
terminal vertex. The word $vv ^ {-1}$ is an idempotent for every $v$ in $L$. The canonical form is given by multiplying these idempotents, in shortlex
order, and then postmultiplying by $w$. For example, consider the word $xyy ^ {-1}xx ^ {-1}x ^ {-1}$ again. The words corresponding to the paths to the leaves are in this case $xx$ and $xy$. And $w$ is an empty word since start and terminal vertices are the same. Therefore,
the canonical form is 
\[xxx ^ {-1}x ^ {-1}xyy ^ {-1}x ^ {-1}.\]
 See \texttt{CanonicalForm} (\ref{CanonicalForm:for a free inverse semigroup element}). 

\subsection{\textcolor{Chapter }{FreeInverseSemigroup (for a given rank)}}
\logpage{[ 7, 11, 1 ]}\nobreak
\hyperdef{L}{X7F3F9DED8003CBD0}{}
{\noindent\textcolor{FuncColor}{$\triangleright$\enspace\texttt{FreeInverseSemigroup({\mdseries\slshape rank[, name]})\index{FreeInverseSemigroup@\texttt{FreeInverseSemigroup}!for a given rank}
\label{FreeInverseSemigroup:for a given rank}
}\hfill{\scriptsize (function)}}\\
\noindent\textcolor{FuncColor}{$\triangleright$\enspace\texttt{FreeInverseSemigroup({\mdseries\slshape name1, name2, ...})\index{FreeInverseSemigroup@\texttt{FreeInverseSemigroup}!for a list of names}
\label{FreeInverseSemigroup:for a list of names}
}\hfill{\scriptsize (function)}}\\
\noindent\textcolor{FuncColor}{$\triangleright$\enspace\texttt{FreeInverseSemigroup({\mdseries\slshape names})\index{FreeInverseSemigroup@\texttt{FreeInverseSemigroup}!for various names}
\label{FreeInverseSemigroup:for various names}
}\hfill{\scriptsize (function)}}\\
\textbf{\indent Returns:\ }
 A free inverse semigroup. 



 Returns a free inverse semigroup on \mbox{\texttt{\mdseries\slshape rank}} generators, where \mbox{\texttt{\mdseries\slshape rank}} is a positive integer. If \mbox{\texttt{\mdseries\slshape rank}} is not specified, the number of \mbox{\texttt{\mdseries\slshape names}} is used. In the second and third forms, the \mbox{\texttt{\mdseries\slshape names}} must be distinct. If \texttt{S} is a free inverse semigroup, then the generators can be accessed by \texttt{S.1}, \texttt{S.2} and so on. 
\begin{Verbatim}[commandchars=!@|,fontsize=\small,frame=single,label=Example]
  !gapprompt@gap>| !gapinput@S := FreeInverseSemigroup(7);|
  <free inverse semigroup on the generators
  [ x1, x2, x3, x4, x5, x6, x7 ]>
  !gapprompt@gap>| !gapinput@S := FreeInverseSemigroup(7, "s");|
  <free inverse semigroup on the generators
  [ s1, s2, s3, s4, s5, s6, s7 ]>
  !gapprompt@gap>| !gapinput@S := FreeInverseSemigroup("a", "b", "c");|
  <free inverse semigroup on the generators [ a, b, c ]>
  !gapprompt@gap>| !gapinput@S := FreeInverseSemigroup(["a", "b", "c"]);|
  <free inverse semigroup on the generators [ a, b, c ]>
  !gapprompt@gap>| !gapinput@S.1;|
  a
  !gapprompt@gap>| !gapinput@S.2;|
  b
\end{Verbatim}
 }

 

\subsection{\textcolor{Chapter }{IsFreeInverseSemigroupCategory}}
\logpage{[ 7, 11, 2 ]}\nobreak
\hyperdef{L}{X7CE4CFD886220179}{}
{\noindent\textcolor{FuncColor}{$\triangleright$\enspace\texttt{IsFreeInverseSemigroupCategory({\mdseries\slshape obj})\index{IsFreeInverseSemigroupCategory@\texttt{IsFreeInverseSemigroupCategory}}
\label{IsFreeInverseSemigroupCategory}
}\hfill{\scriptsize (Category)}}\\


 Every free inverse semigroup in \textsf{GAP} created by \texttt{FreeInverseSemigroup} (\ref{FreeInverseSemigroup:for a given rank}) belongs to the category \texttt{IsFreeInverseSemigroup}. Basic operations for a free inverse semigroup are: \texttt{GeneratorsOfInverseSemigroup} (\textbf{Reference: GeneratorsOfInverseSemigroup}) and \texttt{GeneratorsOfSemigroup} (\textbf{Reference: GeneratorsOfSemigroup}). Elements of a free inverse semigroup belong to the category \texttt{IsFreeInverseSemigroupElement} (\ref{IsFreeInverseSemigroupElement}). }

 

\subsection{\textcolor{Chapter }{IsFreeInverseSemigroup}}
\logpage{[ 7, 11, 3 ]}\nobreak
\hyperdef{L}{X7B91643B827DA6DB}{}
{\noindent\textcolor{FuncColor}{$\triangleright$\enspace\texttt{IsFreeInverseSemigroup({\mdseries\slshape S})\index{IsFreeInverseSemigroup@\texttt{IsFreeInverseSemigroup}}
\label{IsFreeInverseSemigroup}
}\hfill{\scriptsize (property)}}\\
\textbf{\indent Returns:\ }
 \texttt{true} or \texttt{false} 



 Attempts to determine whether the given semigroup \mbox{\texttt{\mdseries\slshape S}} is a free inverse semigroup. }

 

\subsection{\textcolor{Chapter }{IsFreeInverseSemigroupElement}}
\logpage{[ 7, 11, 4 ]}\nobreak
\hyperdef{L}{X7999FE0286283CC2}{}
{\noindent\textcolor{FuncColor}{$\triangleright$\enspace\texttt{IsFreeInverseSemigroupElement\index{IsFreeInverseSemigroupElement@\texttt{IsFreeInverseSemigroupElement}}
\label{IsFreeInverseSemigroupElement}
}\hfill{\scriptsize (Category)}}\\


 Every element of a free inverse semigroup belongs to the category \texttt{IsFreeInverseSemigroupElement}. }

 

\subsection{\textcolor{Chapter }{IsFreeInverseSemigroupElementCollection}}
\logpage{[ 7, 11, 5 ]}\nobreak
\hyperdef{L}{X813A291779726739}{}
{\noindent\textcolor{FuncColor}{$\triangleright$\enspace\texttt{IsFreeInverseSemigroupElementCollection\index{IsFreeInverseSemigroupElementCollection@\texttt{IsFree}\-\texttt{Inverse}\-\texttt{Semigroup}\-\texttt{Element}\-\texttt{Collection}}
\label{IsFreeInverseSemigroupElementCollection}
}\hfill{\scriptsize (Category)}}\\


 Every collection of elements of a free inverse semigroup belongs to the
category \texttt{IsFreeInverseSemigroupElementCollection}. For example, every free inverse semigroup belongs to \texttt{IsFreeInverseSemigroupElementCollection}. }

 

\subsection{\textcolor{Chapter }{CanonicalForm (for a free inverse semigroup element)}}
\logpage{[ 7, 11, 6 ]}\nobreak
\hyperdef{L}{X7DB7DCEC7E0FE9A3}{}
{\noindent\textcolor{FuncColor}{$\triangleright$\enspace\texttt{CanonicalForm({\mdseries\slshape w})\index{CanonicalForm@\texttt{CanonicalForm}!for a free inverse semigroup element}
\label{CanonicalForm:for a free inverse semigroup element}
}\hfill{\scriptsize (attribute)}}\\
\textbf{\indent Returns:\ }
 A string. 



 Every element of a free inverse semigroup has a unique canonical form. If \mbox{\texttt{\mdseries\slshape w}} is such an element, then \texttt{CanonicalForm} returns the canonical form of \mbox{\texttt{\mdseries\slshape w}} as a string. 
\begin{Verbatim}[commandchars=!@|,fontsize=\small,frame=single,label=Example]
  !gapprompt@gap>| !gapinput@S := FreeInverseSemigroup(3);|
  <free inverse semigroup on the generators [ x1, x2, x3 ]>
  !gapprompt@gap>| !gapinput@x := S.1; y := S.2;|
  x1
  x2
  !gapprompt@gap>| !gapinput@CanonicalForm(x ^ 3 * y ^ 3);|
  "x1x1x1x2x2x2x2^-1x2^-1x2^-1x1^-1x1^-1x1^-1x1x1x1x2x2x2"
\end{Verbatim}
 }

 

\subsection{\textcolor{Chapter }{MinimalWord (for free inverse semigroup element)}}
\logpage{[ 7, 11, 7 ]}\nobreak
\hyperdef{L}{X87BB5D047EB7C2BF}{}
{\noindent\textcolor{FuncColor}{$\triangleright$\enspace\texttt{MinimalWord({\mdseries\slshape w})\index{MinimalWord@\texttt{MinimalWord}!for free inverse semigroup element}
\label{MinimalWord:for free inverse semigroup element}
}\hfill{\scriptsize (attribute)}}\\
\textbf{\indent Returns:\ }
 A string. 



 For an element \mbox{\texttt{\mdseries\slshape w}} of a free inverse semigroup \texttt{S}, \texttt{MinimalWord} returns a word of minimal length equal to \mbox{\texttt{\mdseries\slshape w}} in \texttt{S} as a string.

 Note that there maybe more than one word of minimal length which is equal to \mbox{\texttt{\mdseries\slshape w}} in \texttt{S}. 
\begin{Verbatim}[commandchars=!@|,fontsize=\small,frame=single,label=Example]
  !gapprompt@gap>| !gapinput@S := FreeInverseSemigroup(3);|
  <free inverse semigroup on the generators [ x1, x2, x3 ]>
  !gapprompt@gap>| !gapinput@x := S.1;|
  x1
  !gapprompt@gap>| !gapinput@y := S.2;|
  x2
  !gapprompt@gap>| !gapinput@MinimalWord(x ^ 3 * y ^ 3);|
  "x1*x1*x1*x2*x2*x2"
\end{Verbatim}
 }

 
\subsection{\textcolor{Chapter }{Displaying free inverse semigroup elements }}\logpage{[ 7, 11, 8 ]}
\hyperdef{L}{X8073A2387A42B52D}{}
{
   There is a way to change how \textsf{GAP} displays free inverse semigroup elements using the user preference \texttt{FreeInverseSemigroupElementDisplay}. See \texttt{UserPreference} (\textbf{Reference: UserPreference}) for more information about user preferences.

 There are two possible values for \texttt{FreeInverseSemigroupElementDisplay}: 
\begin{description}
\item[{minimal }]  With this option selected, \textsf{GAP} will display a shortest word corresponding to the free inverse semigroup
element. However, this shortest word is not unique. This is a default setting. 
\item[{canonical}]  With this option selected, \textsf{GAP} will display a free inverse semigroup element in the canonical form. 
\end{description}
 
\begin{Verbatim}[commandchars=!@|,fontsize=\small,frame=single,label=Example]
  !gapprompt@gap>| !gapinput@SetUserPreference("semigroups",|
  !gapprompt@>| !gapinput@                     "FreeInverseSemigroupElementDisplay",|
  !gapprompt@>| !gapinput@                     "minimal");|
  !gapprompt@gap>| !gapinput@S := FreeInverseSemigroup(2);|
  <free inverse semigroup on the generators [ x1, x2 ]>
  !gapprompt@gap>| !gapinput@S.1 * S.2;|
  x1*x2
  !gapprompt@gap>| !gapinput@SetUserPreference("semigroups",|
  !gapprompt@>| !gapinput@                     "FreeInverseSemigroupElementDisplay",|
  !gapprompt@>| !gapinput@                     "canonical");|
  !gapprompt@gap>| !gapinput@S.1 * S.2;|
  x1x2x2^-1x1^-1x1x2
\end{Verbatim}
 }

 
\subsection{\textcolor{Chapter }{Operators for free inverse semigroup elements }}\label{Operators and operations for free inverse semigroup elements}
\logpage{[ 7, 11, 9 ]}
\hyperdef{L}{X7A55FD9A7DF21C60}{}
{
  
\begin{description}
\item[{\texttt{\mbox{\texttt{\mdseries\slshape w}} \texttt{\symbol{94}} \texttt{\symbol{45}}1}}]  returns the semigroup inverse of the free inverse semigroup element \mbox{\texttt{\mdseries\slshape w}}. 
\item[{\texttt{\mbox{\texttt{\mdseries\slshape u}} * \mbox{\texttt{\mdseries\slshape v}}}}]  returns the product of two free inverse semigroup elements \mbox{\texttt{\mdseries\slshape u}} and \mbox{\texttt{\mdseries\slshape v}}. 
\item[{\texttt{\mbox{\texttt{\mdseries\slshape u}} = \mbox{\texttt{\mdseries\slshape v}} }}]  checks if two free inverse semigroup elements are equal, by comparing their
canonical forms. 
\end{description}
 }

 }

 }

 
\chapter{\textcolor{Chapter }{ Standard constructions }}\label{Standard constructions}
\logpage{[ 8, 0, 0 ]}
\hyperdef{L}{X86EE8DC987BA646E}{}
{
  In this chapter we describe some standard ways of constructing semigroups and
monoids from other semigroups that are available in the \textsf{Semigroups} package. 

 
\section{\textcolor{Chapter }{ Products of semigroups }}\label{Products of semigroups}
\logpage{[ 8, 1, 0 ]}
\hyperdef{L}{X79546641809113CE}{}
{
  In this section, we describe the functions in \textsf{Semigroups} that can be used to create various products of semigroups. 

\subsection{\textcolor{Chapter }{DirectProduct}}
\logpage{[ 8, 1, 1 ]}\nobreak
\hyperdef{L}{X861BA02C7902A4F4}{}
{\noindent\textcolor{FuncColor}{$\triangleright$\enspace\texttt{DirectProduct({\mdseries\slshape S[, T, ...]})\index{DirectProduct@\texttt{DirectProduct}}
\label{DirectProduct}
}\hfill{\scriptsize (function)}}\\
\noindent\textcolor{FuncColor}{$\triangleright$\enspace\texttt{DirectProductOp({\mdseries\slshape list, S})\index{DirectProductOp@\texttt{DirectProductOp}}
\label{DirectProductOp}
}\hfill{\scriptsize (operation)}}\\
\textbf{\indent Returns:\ }
A transformation semigroup.



 The function \texttt{DirectProduct} takes an arbitrary positive number of finite semigroups, and returns a
semigroup that is isomorphic to their direct product. 

 If these finite semigroups are all partial perm semigroups, all bipartition
semigroups, or all PBR semigroups, then \texttt{DirectProduct} returns a semigroup of the same type. Otherwise, \texttt{DirectProduct} returns a transformation semigroup. 

 The operation \texttt{DirectProductOp} is included for consistency with the \textsf{GAP} library (see \texttt{DirectProductOp} (\textbf{Reference: DirectProductOp})). It takes exactly two arguments, namely a non\texttt{\symbol{45}}empty list \mbox{\texttt{\mdseries\slshape list}} of semigroups and one of these semigroups, \mbox{\texttt{\mdseries\slshape S}}, and returns the same result as \texttt{CallFuncList(DirectProduct, \mbox{\texttt{\mdseries\slshape list}})}. 

 If \texttt{D} is the direct product of a collection of semigroups, then an embedding of the \texttt{i}th factor into \texttt{D} can be accessed with the command \texttt{Embedding(D, i)}, and a projection of \texttt{D} onto its \texttt{i}th factor can be accessed with the command \texttt{Projection(D, i)}; see \texttt{Embedding} (\textbf{Reference: Embedding}) and \texttt{Projection} (\textbf{Reference: Projection}) for more information. 
\begin{Verbatim}[commandchars=!@|,fontsize=\small,frame=single,label=Example]
  !gapprompt@gap>| !gapinput@S := InverseMonoid([PartialPerm([2, 1])]);;|
  !gapprompt@gap>| !gapinput@T := InverseMonoid([PartialPerm([1, 2, 3])]);;|
  !gapprompt@gap>| !gapinput@D := DirectProduct(S, T);|
  <commutative inverse partial perm monoid of rank 5 with 1 generator>
  !gapprompt@gap>| !gapinput@Elements(D);|
  [ <identity partial perm on [ 1, 2, 3, 4, 5 ]>, (1,2)(3)(4)(5) ]
  !gapprompt@gap>| !gapinput@S := PartitionMonoid(2);;|
  !gapprompt@gap>| !gapinput@D := DirectProduct(S, S, S);; IsRegularSemigroup(D);; D;|
  <regular bipartition monoid of size 3375, degree 6 with 9 generators>
  !gapprompt@gap>| !gapinput@S := Semigroup([PartialPerm([2, 5, 0, 1, 3]),|
  !gapprompt@>| !gapinput@                   PartialPerm([5, 2, 4, 3])]);;|
  !gapprompt@gap>| !gapinput@T := Semigroup([Bipartition([[1, -2], [2], [3, -1, -3]])]);;|
  !gapprompt@gap>| !gapinput@D := DirectProduct(S, T);|
  <transformation semigroup of size 122, degree 9 with 63 generators>
  !gapprompt@gap>| !gapinput@Size(D) = Size(S) * Size(T);|
  true
\end{Verbatim}
 }

 

\subsection{\textcolor{Chapter }{WreathProduct}}
\logpage{[ 8, 1, 2 ]}\nobreak
\hyperdef{L}{X8786EFBC78D7D6ED}{}
{\noindent\textcolor{FuncColor}{$\triangleright$\enspace\texttt{WreathProduct({\mdseries\slshape M, S})\index{WreathProduct@\texttt{WreathProduct}}
\label{WreathProduct}
}\hfill{\scriptsize (operation)}}\\
\textbf{\indent Returns:\ }
A transformation semigroup.



 If \mbox{\texttt{\mdseries\slshape M}} is a transformation monoid (or a permutation group) of degree \texttt{m}, and \mbox{\texttt{\mdseries\slshape S}} is a transformation semigroup (or permutation group) of degree \texttt{s}, the operation \texttt{WreathProduct(\mbox{\texttt{\mdseries\slshape M}}, \mbox{\texttt{\mdseries\slshape S}})} returns the wreath product of \mbox{\texttt{\mdseries\slshape M}} and \mbox{\texttt{\mdseries\slshape S}}, in terms of an embedding in the full transformation monoid of degree \texttt{m * s}. 
\begin{Verbatim}[commandchars=!@|,fontsize=\small,frame=single,label=Example]
  !gapprompt@gap>| !gapinput@T := FullTransformationMonoid(3);;|
  !gapprompt@gap>| !gapinput@C := Group((1, 3));;|
  !gapprompt@gap>| !gapinput@W := WreathProduct(T, C);;|
  !gapprompt@gap>| !gapinput@Size(W);|
  39366
  !gapprompt@gap>| !gapinput@WW := WreathProduct(C, T);;|
  !gapprompt@gap>| !gapinput@Size(WW);|
  216
\end{Verbatim}
 }

 }

   
\section{\textcolor{Chapter }{ Dual semigroups }}\logpage{[ 8, 2, 0 ]}
\hyperdef{L}{X7F035EC07AA7CD97}{}
{
  The \emph{dual semigroup} of a semigroup \texttt{S} is the semigroup with the same underlying set of elements but with reversed
multiplication; this is anti\texttt{\symbol{45}}isomorphic to \texttt{S}. In \textsf{Semigroups} a semigroup and its dual are represented with disjoint sets of elements. 

\subsection{\textcolor{Chapter }{DualSemigroup}}
\logpage{[ 8, 2, 1 ]}\nobreak
\hyperdef{L}{X79F2643C8642A3B0}{}
{\noindent\textcolor{FuncColor}{$\triangleright$\enspace\texttt{DualSemigroup({\mdseries\slshape S})\index{DualSemigroup@\texttt{DualSemigroup}}
\label{DualSemigroup}
}\hfill{\scriptsize (attribute)}}\\
\textbf{\indent Returns:\ }
The dual semigroup of the given semigroup.



 The dual semigroup of a semigroup \mbox{\texttt{\mdseries\slshape S}} is the semigroup with the same underlying set as \mbox{\texttt{\mdseries\slshape S}}, but with multiplication reversed; this is anti\texttt{\symbol{45}}isomorphic
to \mbox{\texttt{\mdseries\slshape S}}. This attribute returns a semigroup isomorphic to the dual semigroup of \mbox{\texttt{\mdseries\slshape S}}. 
\begin{Verbatim}[commandchars=!@|,fontsize=\small,frame=single,label=Example]
  
  !gapprompt@gap>| !gapinput@S := Semigroup([Transformation([1, 4, 3, 2, 2]),|
  !gapprompt@>| !gapinput@Transformation([5, 4, 4, 1, 2])]);;|
  !gapprompt@gap>| !gapinput@D := DualSemigroup(S);|
  <dual semigroup of <transformation semigroup of degree 5 with 2
   generators>>
  !gapprompt@gap>| !gapinput@Size(S) = Size(D);|
  true
  !gapprompt@gap>| !gapinput@NrDClasses(S) = NrDClasses(D);|
  true
\end{Verbatim}
 }

 

\subsection{\textcolor{Chapter }{IsDualSemigroupRep}}
\logpage{[ 8, 2, 2 ]}\nobreak
\hyperdef{L}{X83403224821CD079}{}
{\noindent\textcolor{FuncColor}{$\triangleright$\enspace\texttt{IsDualSemigroupRep({\mdseries\slshape sgrp})\index{IsDualSemigroupRep@\texttt{IsDualSemigroupRep}}
\label{IsDualSemigroupRep}
}\hfill{\scriptsize (Category)}}\\
\textbf{\indent Returns:\ }
Returns \texttt{true} if \mbox{\texttt{\mdseries\slshape sgrp}} lies in the category of dual semigroups.



 Semigroups created using \texttt{DualSemigroup} (\ref{DualSemigroup}) normally lie in this category. The exception is semigroups which are the dual
of semigroups already lying in this category. That is, a semigroup lies in the
category \texttt{IsDualSemigroupRep} if and only if the corresponding dual semigroup does not. Note that this is
not a Representation in the GAP sense, and will likely be renamed in a future
major release of the package. 
\begin{Verbatim}[commandchars=!@|,fontsize=\small,frame=single,label=Example]
  
  !gapprompt@gap>| !gapinput@S := Semigroup([Transformation([3, 5, 1, 1, 2]),|
  !gapprompt@>| !gapinput@Transformation([1, 2, 4, 4, 3])]);|
  <transformation semigroup of degree 5 with 2 generators>
  !gapprompt@gap>| !gapinput@D := DualSemigroup(S);|
  <dual semigroup of <transformation semigroup of degree 5 with 2
   generators>>
  !gapprompt@gap>| !gapinput@IsDualSemigroupRep(D);|
  true
  !gapprompt@gap>| !gapinput@R := DualSemigroup(D);|
  <transformation semigroup of degree 5 with 2 generators>
  !gapprompt@gap>| !gapinput@IsDualSemigroupRep(R);|
  false
  !gapprompt@gap>| !gapinput@R = S;|
  true
  !gapprompt@gap>| !gapinput@T := Range(IsomorphismTransformationSemigroup(D));|
  <transformation semigroup of size 16, degree 17 with 2 generators>
  !gapprompt@gap>| !gapinput@IsDualSemigroupRep(T);|
  false
  !gapprompt@gap>| !gapinput@x := Representative(D);|
  <Transformation( [ 3, 5, 1, 1, 2 ] ) in the dual semigroup>
  !gapprompt@gap>| !gapinput@V := Semigroup(x);|
  <dual semigroup of <commutative transformation semigroup of degree 5
   with 1 generator>>
  !gapprompt@gap>| !gapinput@IsDualSemigroupRep(V);|
  true
\end{Verbatim}
 }

 

\subsection{\textcolor{Chapter }{IsDualSemigroupElement}}
\logpage{[ 8, 2, 3 ]}\nobreak
\hyperdef{L}{X79BAAA397FC1FA2E}{}
{\noindent\textcolor{FuncColor}{$\triangleright$\enspace\texttt{IsDualSemigroupElement({\mdseries\slshape elt})\index{IsDualSemigroupElement@\texttt{IsDualSemigroupElement}}
\label{IsDualSemigroupElement}
}\hfill{\scriptsize (Category)}}\\
\textbf{\indent Returns:\ }
Returns \texttt{true} if \mbox{\texttt{\mdseries\slshape elt}} has the representation of a dual semigroup element.



 Elements of a dual semigroup obtained using \texttt{AntiIsomorphismDualSemigroup} (\ref{AntiIsomorphismDualSemigroup}) normally lie in this category. The exception is elements obtained by applying
the map \texttt{AntiIsomorphismDualSemigroup} (\ref{AntiIsomorphismDualSemigroup}) to elements already in this category. That is, the elements of a semigroup lie
in the category \texttt{IsDualSemigroupElement} if and only if the elements of the corresponding dual semigroup do not. 
\begin{Verbatim}[commandchars=!@|,fontsize=\small,frame=single,label=Example]
  
  !gapprompt@gap>| !gapinput@S := SingularPartitionMonoid(4);;|
  !gapprompt@gap>| !gapinput@D := DualSemigroup(S);;|
  !gapprompt@gap>| !gapinput@s := GeneratorsOfSemigroup(S)[1];;|
  !gapprompt@gap>| !gapinput@map := AntiIsomorphismDualSemigroup(S);;|
  !gapprompt@gap>| !gapinput@t := s ^ map;|
  <<block bijection: [ 1, 2, -1, -2 ], [ 3, -3 ], [ 4, -4 ]>
    in the dual semigroup>
  !gapprompt@gap>| !gapinput@IsDualSemigroupElement(t);|
  true
  !gapprompt@gap>| !gapinput@inv := InverseGeneralMapping(map);;|
  !gapprompt@gap>| !gapinput@x := t ^ inv;|
  <block bijection: [ 1, 2, -1, -2 ], [ 3, -3 ], [ 4, -4 ]>
  !gapprompt@gap>| !gapinput@IsDualSemigroupElement(x);|
  false
\end{Verbatim}
 }

 

\subsection{\textcolor{Chapter }{AntiIsomorphismDualSemigroup}}
\logpage{[ 8, 2, 4 ]}\nobreak
\hyperdef{L}{X7CB64FA378EC715B}{}
{\noindent\textcolor{FuncColor}{$\triangleright$\enspace\texttt{AntiIsomorphismDualSemigroup({\mdseries\slshape S})\index{AntiIsomorphismDualSemigroup@\texttt{AntiIsomorphismDualSemigroup}}
\label{AntiIsomorphismDualSemigroup}
}\hfill{\scriptsize (attribute)}}\\
\textbf{\indent Returns:\ }
 An anti\texttt{\symbol{45}}isomorphism from \mbox{\texttt{\mdseries\slshape S}} to the corresponding dual semigroup. 



 The dual semigroup of \mbox{\texttt{\mdseries\slshape S}} mathematically has the same underlying set as \mbox{\texttt{\mdseries\slshape S}}, but is represented with a different set of elements in \textsf{Semigroups}. This function returns a mapping which is an
anti\texttt{\symbol{45}}isomorphism from \mbox{\texttt{\mdseries\slshape S}} to its dual. 
\begin{Verbatim}[commandchars=!@|,fontsize=\small,frame=single,label=Example]
  
  !gapprompt@gap>| !gapinput@S := PartitionMonoid(3);|
  <regular bipartition *-monoid of size 203, degree 3 with 4 generators>
  !gapprompt@gap>| !gapinput@map := AntiIsomorphismDualSemigroup(S);|
  MappingByFunction( <regular bipartition *-monoid of size 203,
   degree 3 with 4 generators>, <dual semigroup of
  <regular bipartition *-monoid of size 203, degree 3 with 4 generators>
   >, function( x ) ... end, function( x ) ... end )
  !gapprompt@gap>| !gapinput@inv := InverseGeneralMapping(map);;|
  !gapprompt@gap>| !gapinput@x := Bipartition([[1, -2], [2, -3], [3, -1]]);|
  <block bijection: [ 1, -2 ], [ 2, -3 ], [ 3, -1 ]>
  !gapprompt@gap>| !gapinput@y := Bipartition([[1], [2, -2], [3, -3], [-1]]);|
  <bipartition: [ 1 ], [ 2, -2 ], [ 3, -3 ], [ -1 ]>
  !gapprompt@gap>| !gapinput@(x ^ map) * (y ^ map) = (y * x) ^ map;|
  true
  !gapprompt@gap>| !gapinput@x ^ map;|
  <<block bijection: [ 1, -2 ], [ 2, -3 ], [ 3, -1 ]>
    in the dual semigroup>
\end{Verbatim}
 }

 }

   
\section{\textcolor{Chapter }{ Strong semilattices of semigroups }}\label{Strong semilattices of semigroups}
\logpage{[ 8, 3, 0 ]}
\hyperdef{L}{X7BEA92E67A6D349A}{}
{
  In this section, we describe how \textsf{Semigroups} can be used to create and manipulate strong semilattices of semigroups (SSSs).
Strong semilattices of semigroups are described, for example, in Section 4.1
of \cite{Howie1995aa}. They consist of a meet\texttt{\symbol{45}}semilattice $Y$ along with a collection of semigroups $S_a$ for each $a$ in $Y$, and a collection of homomorphisms $f_{ab} : S_a \rightarrow S_b$ for each $a$ and $b$ in $Y$ such that $a \geq b$. 

 The product of two elements $x \in S_a, y \in S_b$ is defined to lie in the semigroup $S_c$, corresponding to the meet $c$ of $a, b \in Y$. The exact element of $S_c$ equal to the product is obtained using the homomorphisms of the SSS: $xy = (x f_{ac}) (y f_{bc})$. 

\subsection{\textcolor{Chapter }{StrongSemilatticeOfSemigroups}}
\logpage{[ 8, 3, 1 ]}\nobreak
\hyperdef{L}{X82C3F9C9861EEDFE}{}
{\noindent\textcolor{FuncColor}{$\triangleright$\enspace\texttt{StrongSemilatticeOfSemigroups({\mdseries\slshape D, L, H})\index{StrongSemilatticeOfSemigroups@\texttt{StrongSemilatticeOfSemigroups}}
\label{StrongSemilatticeOfSemigroups}
}\hfill{\scriptsize (operation)}}\\
\textbf{\indent Returns:\ }
 A strong semilattice of semigroups. 



 If \mbox{\texttt{\mdseries\slshape D}} is a digraph, \mbox{\texttt{\mdseries\slshape L}} is a list of semigroups, and \mbox{\texttt{\mdseries\slshape H}} is a list of lists of maps, then this function returns a corresponding \texttt{IsStrongSemilatticeOfSemigroups} object. The format of the arguments is not required to be exactly analogous to
Howie's description above, but consistency amongst the arguments is required: 
\begin{itemize}
\item  \mbox{\texttt{\mdseries\slshape D}} must be a digraph whose \texttt{DigraphReflexiveTransitiveClosure} (\textbf{Digraphs: DigraphReflexiveTransitiveClosure}) is a meet\texttt{\symbol{45}}semilattice. For example, \texttt{Digraph([2, 3], [4], [4], []])} is valid and produces a semilattice where the meet of \texttt{2} and \texttt{3} is \texttt{1}. See \texttt{IsMeetSemilatticeDigraph} (\textbf{Digraphs: IsMeetSemilatticeDigraph}). 
\item  \mbox{\texttt{\mdseries\slshape L}} must contain as many semigroups as there are vertices in \mbox{\texttt{\mdseries\slshape D}}. 
\item  \mbox{\texttt{\mdseries\slshape H}} must be a list with as many elements as there are vertices in \mbox{\texttt{\mdseries\slshape D}}. Each element of \mbox{\texttt{\mdseries\slshape H}} must itself be a (possibly empty) list with as many entries as the
corresponding vertex of \mbox{\texttt{\mdseries\slshape D}} has out\texttt{\symbol{45}}edges. The entries of each sublist must be the
corresponding homomorphisms: for example, if \mbox{\texttt{\mdseries\slshape D}} is entered as above, then \texttt{H[1][2]} must be the homomorphism $f_31$, i.e. \texttt{H[1][2]} is an \texttt{IsMapping} object whose domain is a superset of \texttt{L[3]} and whose range is a subset of \texttt{L[1]}. 
\end{itemize}
 Note that in the example above, the edge $1 \rightarrow 4$ is not entered as part of the argument \mbox{\texttt{\mdseries\slshape D}}, but it is still an edge in the reflexive transitive closure of \mbox{\texttt{\mdseries\slshape D}}. When creating the object, \textsf{GAP} creates the homomorphism $f_{41}$ by composing the mappings along paths that lead from 4 to 1, and checks that
composing along all possible paths produces the same result. }

 

\subsection{\textcolor{Chapter }{SSSE}}
\logpage{[ 8, 3, 2 ]}\nobreak
\hyperdef{L}{X798DE3E581978834}{}
{\noindent\textcolor{FuncColor}{$\triangleright$\enspace\texttt{SSSE({\mdseries\slshape SSS, n, x})\index{SSSE@\texttt{SSSE}}
\label{SSSE}
}\hfill{\scriptsize (operation)}}\\
\textbf{\indent Returns:\ }
 An element of a strong semilattice of semigroups. 



 If \mbox{\texttt{\mdseries\slshape n}} is a vertex of the underlying semilattice of the strong semilattice of
semigroups \mbox{\texttt{\mdseries\slshape SSS}}, and if \mbox{\texttt{\mdseries\slshape x}} is an element of the \mbox{\texttt{\mdseries\slshape n}}th semigroup of \mbox{\texttt{\mdseries\slshape SSS}}, then this function returns the element of \mbox{\texttt{\mdseries\slshape SSS}} which lies in semigroup number \mbox{\texttt{\mdseries\slshape n}} and which corresponds to the element \mbox{\texttt{\mdseries\slshape x}} in that semigroup. 

 This function returns an \texttt{IsSSSE} (\ref{IsSSSE}) object. SSSEs from the same strong semilattice of semigroups can be compared
and multiplied. 
\begin{Verbatim}[commandchars=!@|,fontsize=\small,frame=single,label=Example]
  !gapprompt@gap>| !gapinput@D := Digraph([[2, 3], [4], [4], []]);;|
  !gapprompt@gap>| !gapinput@S4 := FullTransformationMonoid(2);;|
  !gapprompt@gap>| !gapinput@S3 := FullTransformationMonoid(3);;|
  !gapprompt@gap>| !gapinput@pairs := [[Transformation([1, 2]), Transformation([2, 1])]];;|
  !gapprompt@gap>| !gapinput@cong := SemigroupCongruence(S4, pairs);;|
  !gapprompt@gap>| !gapinput@S2 := S4 / cong;;|
  !gapprompt@gap>| !gapinput@S1 := TrivialSemigroup();;|
  !gapprompt@gap>| !gapinput@L := [S1, S2, S3, S4];;|
  !gapprompt@gap>| !gapinput@idfn := t -> IdentityTransformation;;|
  !gapprompt@gap>| !gapinput@f21 := SemigroupHomomorphismByFunction(S2, S1, idfn);;|
  !gapprompt@gap>| !gapinput@f31 := SemigroupHomomorphismByFunction(S3, S1, idfn);;|
  !gapprompt@gap>| !gapinput@f42 := HomomorphismQuotientSemigroup(cong);;|
  !gapprompt@gap>| !gapinput@f43 := SemigroupHomomorphismByFunction(S4, S3, IdFunc);;|
  !gapprompt@gap>| !gapinput@H := [[f21, f31], [f42], [f43], []];;|
  !gapprompt@gap>| !gapinput@SSS := StrongSemilatticeOfSemigroups(D, L, H);|
  <strong semilattice of 4 semigroups>
  !gapprompt@gap>| !gapinput@Size(SSS);|
  34
  !gapprompt@gap>| !gapinput@x := SSSE(SSS, 3, Elements(S3)[10]);|
  SSSE(3, Transformation( [ 2, 1, 1 ] ))
  !gapprompt@gap>| !gapinput@y := SSSE(SSS, 4, Elements(S4)[1]);|
  SSSE(4, Transformation( [ 1, 1 ] ))
  !gapprompt@gap>| !gapinput@x * y;|
  SSSE(3, Transformation( [ 1, 1, 1 ] ))
\end{Verbatim}
 }

 

\subsection{\textcolor{Chapter }{IsSSSE}}
\logpage{[ 8, 3, 3 ]}\nobreak
\hyperdef{L}{X7B7B70F37C9C3836}{}
{\noindent\textcolor{FuncColor}{$\triangleright$\enspace\texttt{IsSSSE({\mdseries\slshape obj})\index{IsSSSE@\texttt{IsSSSE}}
\label{IsSSSE}
}\hfill{\scriptsize (filter)}}\\
\textbf{\indent Returns:\ }
 \texttt{true} or \texttt{false}. 



 All elements of an SSS belong in the category \texttt{IsSSSE} (for "Strong Semilattice of Semigroups Element"). }

 

\subsection{\textcolor{Chapter }{IsStrongSemilatticeOfSemigroups}}
\logpage{[ 8, 3, 4 ]}\nobreak
\hyperdef{L}{X838F24247D4DBE18}{}
{\noindent\textcolor{FuncColor}{$\triangleright$\enspace\texttt{IsStrongSemilatticeOfSemigroups({\mdseries\slshape obj})\index{IsStrongSemilatticeOfSemigroups@\texttt{IsStrongSemilatticeOfSemigroups}}
\label{IsStrongSemilatticeOfSemigroups}
}\hfill{\scriptsize (filter)}}\\
\textbf{\indent Returns:\ }
 \texttt{true} or \texttt{false}. 



 Every Strong Semilattice of Semigroups in \textsf{GAP} belongs to the category \texttt{IsStrongSemilatticeOfSemigroups}. Basic operations in this category allow the user to recover the three
essential elements of an SSS object: \texttt{SemilatticeOfStrongSemilatticeOfSemigroups} (\ref{SemilatticeOfStrongSemilatticeOfSemigroups}), \texttt{SemigroupsOfStrongSemilatticeOfSemigroups} (\ref{SemigroupsOfStrongSemilatticeOfSemigroups}), and \texttt{HomomorphismsOfStrongSemilatticeOfSemigroups} (\ref{HomomorphismsOfStrongSemilatticeOfSemigroups}). }

 

\subsection{\textcolor{Chapter }{SemilatticeOfStrongSemilatticeOfSemigroups}}
\logpage{[ 8, 3, 5 ]}\nobreak
\hyperdef{L}{X87100CE6836DE3DB}{}
{\noindent\textcolor{FuncColor}{$\triangleright$\enspace\texttt{SemilatticeOfStrongSemilatticeOfSemigroups({\mdseries\slshape SSS})\index{SemilatticeOfStrongSemilatticeOfSemigroups@\texttt{Semilattice}\-\texttt{Of}\-\texttt{Strong}\-\texttt{Semilattice}\-\texttt{Of}\-\texttt{Semigroups}}
\label{SemilatticeOfStrongSemilatticeOfSemigroups}
}\hfill{\scriptsize (attribute)}}\\
\textbf{\indent Returns:\ }
 A meet\texttt{\symbol{45}}semilattice digraph. 



 If \mbox{\texttt{\mdseries\slshape SSS}} is a strong semilattice of semigroups, this function returns the underlying
semilattice structure as a digraph. Note that this may not be equal to the
digraph passed as input when \mbox{\texttt{\mdseries\slshape SSS}} was created: rather, it is the reflexive transitive closure of the input
digraph. }

 

\subsection{\textcolor{Chapter }{SemigroupsOfStrongSemilatticeOfSemigroups}}
\logpage{[ 8, 3, 6 ]}\nobreak
\hyperdef{L}{X79E6C08D87984579}{}
{\noindent\textcolor{FuncColor}{$\triangleright$\enspace\texttt{SemigroupsOfStrongSemilatticeOfSemigroups({\mdseries\slshape SSS})\index{SemigroupsOfStrongSemilatticeOfSemigroups@\texttt{Semigroups}\-\texttt{Of}\-\texttt{Strong}\-\texttt{Semilattice}\-\texttt{Of}\-\texttt{Semigroups}}
\label{SemigroupsOfStrongSemilatticeOfSemigroups}
}\hfill{\scriptsize (attribute)}}\\
\textbf{\indent Returns:\ }
 A list of semigroups. 



 If \mbox{\texttt{\mdseries\slshape SSS}} is a strong semilattice of semigroups, this function returns the list of
semigroups that make up \mbox{\texttt{\mdseries\slshape SSS}}. The position of a semigroup in the list corresponds to the node of the
semilattice where that semigroup lies. }

 

\subsection{\textcolor{Chapter }{HomomorphismsOfStrongSemilatticeOfSemigroups}}
\logpage{[ 8, 3, 7 ]}\nobreak
\hyperdef{L}{X806655138370ECFF}{}
{\noindent\textcolor{FuncColor}{$\triangleright$\enspace\texttt{HomomorphismsOfStrongSemilatticeOfSemigroups({\mdseries\slshape SSS})\index{HomomorphismsOfStrongSemilatticeOfSemigroups@\texttt{Homomorphisms}\-\texttt{Of}\-\texttt{Strong}\-\texttt{Semilattice}\-\texttt{Of}\-\texttt{Semigroups}}
\label{HomomorphismsOfStrongSemilatticeOfSemigroups}
}\hfill{\scriptsize (attribute)}}\\
\textbf{\indent Returns:\ }
 A list of lists of mappings. 



 If \mbox{\texttt{\mdseries\slshape SSS}} is a strong semilattice of $n$ semigroups, this function returns an $n \times n$ list where the $(i, j)$th entry of the list is the homomorphism $f_{ji}$, provided $i \leq j$ in the semilattice. If this last condition is not true, then the entry is \texttt{fail}. }

 }

   
\section{\textcolor{Chapter }{ McAlister triple semigroups }}\label{McAlister triple semigroups}
\logpage{[ 8, 4, 0 ]}
\hyperdef{L}{X7CC4F6FE87AFE638}{}
{
  In this section, we describe the functions in \textsf{GAP} for creating and computing with McAlister triple semigroups and their
subsemigroups. This implementation is based on the section in Chapter 5 of \cite{Howie1995aa} but differs from the treatment in Howie by using right actions instead of
left. Some definitions found in the documentation are changed for this reason. 

 The importance of the McAlister triple semigroups lies in the fact that they
are exactly the E\texttt{\symbol{45}}unitary inverse semigroups, which are an
important class in the study of inverse semigroups. 

 First we define E\texttt{\symbol{45}}unitary inverse semigroups. It is
standard to denote the subsemigroup of a semigroup consisting of its
idempotents by \texttt{E}. A semigroup \texttt{S} is said to be \emph{E\texttt{\symbol{45}}unitary} if for all \texttt{e} in \texttt{E} and for all \texttt{s} in \texttt{S}: 
\begin{itemize}
\item  \texttt{es} $\in$ \texttt{E} implies \texttt{s} $\in$ \texttt{E}, 
\item  \texttt{se} $\in$ \texttt{E} implies \texttt{s} $\in$ \texttt{E}. 
\end{itemize}
 For inverse semigroups these two conditions are equivalent. We are only
interested in \emph{E\texttt{\symbol{45}}unitary inverse semigroups}. Before defining McAlister triple semigroups we define a McAlister triple. A \emph{McAlister triple} is a triple \texttt{(G,X,Y)} which consists of: 
\begin{itemize}
\item  a partial order \texttt{X}, 
\item  a subset \texttt{Y} of \texttt{X}, 
\item  a group \texttt{G} which acts on \texttt{X}, on the right, by order automorphisms. That means for all \texttt{A,B} $\in$ \texttt{X} and for all \texttt{g} $\in$ \texttt{G}: \texttt{A} $\leq$ \texttt{B} if and only if \texttt{Ag} $\leq$ \texttt{Bg}. 
\end{itemize}
 Furthermore, \texttt{(G,X,Y)} must satisfy the following four properties to be a McAlister triple: 
\begin{description}
\item[{ M1 }]  \texttt{Y} is a subset of \texttt{X} which is a join\texttt{\symbol{45}}semilattice together with the restriction
of the order relation of \texttt{X} to \texttt{Y}. 
\item[{ M2 }]  \texttt{Y} is an order ideal of \texttt{X}. That is to say, for all \texttt{A} $\in$ \texttt{X} and for all \texttt{B} $\in$ \texttt{Y}: if \texttt{A} $\leq$ \texttt{B}, then \texttt{A} $\in$ \texttt{Y}. 
\item[{ M3 }]  Every element of \texttt{X} is the image of some element in \texttt{Y} moved by an element of \texttt{G}. That is to say, for every \texttt{A} $\in$ \texttt{X}, there exists some \texttt{B} $\in$ \texttt{Y} and there exists \texttt{g} $\in$ \texttt{G} such that \texttt{A} = \texttt{Bg}. 
\item[{ M4 }]  Finally, for all \texttt{g} $\in$ \texttt{G}, the intersection \texttt{\texttt{\symbol{123}}yg : y }$\in$\texttt{ Y\texttt{\symbol{125}}} $\cap$ \texttt{Y} is non\texttt{\symbol{45}}empty. 
\end{description}
 We may define an E\texttt{\symbol{45}}unitary inverse semigroup using a
McAlister triple. Given \texttt{(G,X,Y)} let \texttt{M(G,X,Y)} be the set of all pairs \texttt{(A,g)} in \texttt{Y x G} such that \texttt{A} acted on by the inverse of \texttt{g} is in \texttt{Y} together with multiplication defined by 

 \texttt{(A,g)*(B,h) = (Join(A,Bg\texttt{\symbol{94}}\texttt{\symbol{45}}1),hg)} 

 where \texttt{Join} is the natural join operation of the semilattice and \texttt{Bg\texttt{\symbol{94}}\texttt{\symbol{45}}1} is \texttt{B} acted on by the inverse of \texttt{g}. With this operation, \texttt{M(G,X,Y)} is a semigroup which we call a \emph{McAlister triple semigroup} over \texttt{(G,X,Y)}. In fact every McAlister triple semigroup is an E\texttt{\symbol{45}}unitary
inverse semigroup and every E\texttt{\symbol{45}}unitary inverse semigroup is
isomorphic to some McAlister triple semigroup. Note that there need not be a
unique McAlister triple semigroup for a particular McAlister triple because in
general there is more than one way for a group to act on a partial order. 

\subsection{\textcolor{Chapter }{IsMcAlisterTripleSemigroup}}
\logpage{[ 8, 4, 1 ]}\nobreak
\hyperdef{L}{X85C00EB085774624}{}
{\noindent\textcolor{FuncColor}{$\triangleright$\enspace\texttt{IsMcAlisterTripleSemigroup({\mdseries\slshape S})\index{IsMcAlisterTripleSemigroup@\texttt{IsMcAlisterTripleSemigroup}}
\label{IsMcAlisterTripleSemigroup}
}\hfill{\scriptsize (filter)}}\\
\textbf{\indent Returns:\ }
\texttt{true} or \texttt{false}.



 This function returns \texttt{true} if \mbox{\texttt{\mdseries\slshape S}} is a McAlister triple semigroup. A \emph{McAlister triple semigroup} is a special representation of an E\texttt{\symbol{45}}unitary inverse
semigroup \texttt{IsEUnitaryInverseSemigroup} (\ref{IsEUnitaryInverseSemigroup}) created by \texttt{McAlisterTripleSemigroup} (\ref{McAlisterTripleSemigroup}). }

 

\subsection{\textcolor{Chapter }{McAlisterTripleSemigroup}}
\logpage{[ 8, 4, 2 ]}\nobreak
\hyperdef{L}{X7B5FF3A27BB057F2}{}
{\noindent\textcolor{FuncColor}{$\triangleright$\enspace\texttt{McAlisterTripleSemigroup({\mdseries\slshape G, X, Y[, act]})\index{McAlisterTripleSemigroup@\texttt{McAlisterTripleSemigroup}}
\label{McAlisterTripleSemigroup}
}\hfill{\scriptsize (operation)}}\\
\textbf{\indent Returns:\ }
A McAlister triple semigroup.



 The following documentation covers the technical information needed to create
McAlister triple semigroups in GAP, the underlying theory can be read in the
introduction to Chapter \ref{McAlister triple semigroups}. 

 In this implementation the partial order \texttt{X} of a McAlister triple is represented by a \texttt{Digraph} (\textbf{Digraphs: Digraph}) object \mbox{\texttt{\mdseries\slshape X}}. The digraph represents a partial order in the sense that vertices are the
elements of the partial order and the order relation is defined by \texttt{A} $\leq$ \texttt{B} if and only if there is an edge from \texttt{B} to \texttt{A}. The semilattice \texttt{Y} of the McAlister triple should be an induced subdigraph \mbox{\texttt{\mdseries\slshape Y}} of \mbox{\texttt{\mdseries\slshape X}} and the \texttt{DigraphVertexLabels} (\textbf{Digraphs: DigraphVertexLabels}) must correspond to the vertices of \mbox{\texttt{\mdseries\slshape X}} on which \mbox{\texttt{\mdseries\slshape Y}} is induced. That means that the following: 

 \texttt{\mbox{\texttt{\mdseries\slshape Y}} = InducedSubdigraph(\mbox{\texttt{\mdseries\slshape X}}, DigraphVertexLabels(\mbox{\texttt{\mdseries\slshape Y}})) }

 must return \texttt{true}. Herein if we say that a vertex \texttt{A} of \mbox{\texttt{\mdseries\slshape X}} is 'in' \mbox{\texttt{\mdseries\slshape Y}} then we mean there is a vertex of \mbox{\texttt{\mdseries\slshape Y}} whose label is \texttt{A}. Alternatively the user may choose to give the argument \mbox{\texttt{\mdseries\slshape Y}} as the vertices of \mbox{\texttt{\mdseries\slshape X}} on which \mbox{\texttt{\mdseries\slshape Y}} is the induced subdigraph. 

 A McAlister triple semigroup is created from a quadruple \texttt{(\mbox{\texttt{\mdseries\slshape G}}, \mbox{\texttt{\mdseries\slshape X}}, \mbox{\texttt{\mdseries\slshape Y}}, \mbox{\texttt{\mdseries\slshape act}})} where: 
\begin{itemize}
\item  \mbox{\texttt{\mdseries\slshape G}} is a finite group. 
\item  \mbox{\texttt{\mdseries\slshape X}} is a digraph satisfying \texttt{IsPartialOrderDigraph} (\textbf{Digraphs: IsPartialOrderDigraph}). 
\item  \mbox{\texttt{\mdseries\slshape Y}} is a digraph satisfying \texttt{IsJoinSemilatticeDigraph} (\textbf{Digraphs: IsJoinSemilatticeDigraph}) which is an induced subdigraph of \mbox{\texttt{\mdseries\slshape X}} satisfying the aforementioned labeling criteria. Furthermore the \texttt{OutNeighbours} (\textbf{Digraphs: OutNeighbours}) of each vertex of \mbox{\texttt{\mdseries\slshape X}} which is in \mbox{\texttt{\mdseries\slshape Y}} must contain only vertices which are in \mbox{\texttt{\mdseries\slshape Y}}. 
\item  \mbox{\texttt{\mdseries\slshape act}} is a function which takes as its first argument a vertex of the digraph \mbox{\texttt{\mdseries\slshape X}}, its second argument should be an element of \mbox{\texttt{\mdseries\slshape G}}, and it must return a vertex of \mbox{\texttt{\mdseries\slshape X}}. \mbox{\texttt{\mdseries\slshape act}} must be a right action, meaning that \mbox{\texttt{\mdseries\slshape act}}\texttt{(A,gh)=\mbox{\texttt{\mdseries\slshape act}}(\mbox{\texttt{\mdseries\slshape act}}(A,g),h)} holds for all \texttt{A} in \mbox{\texttt{\mdseries\slshape X}} and \texttt{g,h} $\in$ \mbox{\texttt{\mdseries\slshape G}}. Furthermore the permutation representation of this action must be a subgroup
of the automorphism group of \mbox{\texttt{\mdseries\slshape X}}. That means we require the following to return \texttt{true}: 

 \texttt{IsSubgroup(AutomorphismGroup(}\mbox{\texttt{\mdseries\slshape X}}\texttt{), Image(ActionHomomorphism(}\mbox{\texttt{\mdseries\slshape G}}\texttt{, DigraphVertices(}\mbox{\texttt{\mdseries\slshape X}}\texttt{), }\mbox{\texttt{\mdseries\slshape act}}\texttt{));} 

 Furthermore every vertex of \mbox{\texttt{\mdseries\slshape X}} must be in the orbit of some vertex of \mbox{\texttt{\mdseries\slshape X}} which is in \mbox{\texttt{\mdseries\slshape Y}}. Finally, \mbox{\texttt{\mdseries\slshape act}} must fix the vertex of \mbox{\texttt{\mdseries\slshape X}} which is the minimal vertex of \mbox{\texttt{\mdseries\slshape Y}}, i.e. the unique vertex of \mbox{\texttt{\mdseries\slshape Y}} whose only out\texttt{\symbol{45}}neighbour is itself. 
\end{itemize}
 For user convenience, there are multiple versions of \texttt{McAlisterTripleSemigroup}. When the argument \mbox{\texttt{\mdseries\slshape act}} is omitted it is assumed to be \texttt{OnPoints} (\textbf{Reference: OnPoints}). Additionally, the semilattice argument \mbox{\texttt{\mdseries\slshape Y}} may be replaced by a homogeneous list \mbox{\texttt{\mdseries\slshape sub{\textunderscore}ver}} of vertices of \mbox{\texttt{\mdseries\slshape X}}. When \mbox{\texttt{\mdseries\slshape sub{\textunderscore}ver}} is provided, \texttt{McAlisterTripleSemigroup} is called with \mbox{\texttt{\mdseries\slshape Y}} equalling \texttt{InducedSubdigraph(\mbox{\texttt{\mdseries\slshape X}}, \mbox{\texttt{\mdseries\slshape sub{\textunderscore}ver}})} with the appropriate labels. 
\begin{Verbatim}[commandchars=!@|,fontsize=\small,frame=single,label=Example]
  !gapprompt@gap>| !gapinput@x := Digraph([[1], [1, 2], [1, 2, 3], [1, 4], [1, 4, 5]]);|
  <immutable digraph with 5 vertices, 11 edges>
  !gapprompt@gap>| !gapinput@y := InducedSubdigraph(x, [1, 4, 5]);|
  <immutable digraph with 3 vertices, 6 edges>
  !gapprompt@gap>| !gapinput@DigraphVertexLabels(y);|
  [ 1, 4, 5 ]
  !gapprompt@gap>| !gapinput@A := AutomorphismGroup(x);|
  Group([ (2,4)(3,5) ])
  !gapprompt@gap>| !gapinput@S := McAlisterTripleSemigroup(A, x, y, OnPoints);|
  <McAlister triple semigroup over Group([ (2,4)(3,5) ])>
  !gapprompt@gap>| !gapinput@T := McAlisterTripleSemigroup(A, x, y);|
  <McAlister triple semigroup over Group([ (2,4)(3,5) ])>
  !gapprompt@gap>| !gapinput@S = T;|
  false
  !gapprompt@gap>| !gapinput@IsIsomorphicSemigroup(S, T);|
  true
\end{Verbatim}
 }

 

\subsection{\textcolor{Chapter }{McAlisterTripleSemigroupGroup}}
\logpage{[ 8, 4, 3 ]}\nobreak
\hyperdef{L}{X7A54FDB186CD2E94}{}
{\noindent\textcolor{FuncColor}{$\triangleright$\enspace\texttt{McAlisterTripleSemigroupGroup({\mdseries\slshape S})\index{McAlisterTripleSemigroupGroup@\texttt{McAlisterTripleSemigroupGroup}}
\label{McAlisterTripleSemigroupGroup}
}\hfill{\scriptsize (attribute)}}\\
\textbf{\indent Returns:\ }
A group.



 Returns the group used to create the McAlister triple semigroup \mbox{\texttt{\mdseries\slshape S}} via \texttt{McAlisterTripleSemigroup} (\ref{McAlisterTripleSemigroup}). }

 

\subsection{\textcolor{Chapter }{McAlisterTripleSemigroupPartialOrder}}
\logpage{[ 8, 4, 4 ]}\nobreak
\hyperdef{L}{X8046966B7F9A1ED5}{}
{\noindent\textcolor{FuncColor}{$\triangleright$\enspace\texttt{McAlisterTripleSemigroupPartialOrder({\mdseries\slshape S})\index{McAlisterTripleSemigroupPartialOrder@\texttt{McAlister}\-\texttt{Triple}\-\texttt{Semigroup}\-\texttt{Partial}\-\texttt{Order}}
\label{McAlisterTripleSemigroupPartialOrder}
}\hfill{\scriptsize (attribute)}}\\
\textbf{\indent Returns:\ }
A partial order digraph.



 Returns the \texttt{IsPartialOrderDigraph} (\textbf{Digraphs: IsPartialOrderDigraph}) used to create the McAlister triple semigroup \mbox{\texttt{\mdseries\slshape S}} via \texttt{McAlisterTripleSemigroup} (\ref{McAlisterTripleSemigroup}). }

 

\subsection{\textcolor{Chapter }{McAlisterTripleSemigroupSemilattice}}
\logpage{[ 8, 4, 5 ]}\nobreak
\hyperdef{L}{X86C0C3EF84517DAB}{}
{\noindent\textcolor{FuncColor}{$\triangleright$\enspace\texttt{McAlisterTripleSemigroupSemilattice({\mdseries\slshape S})\index{McAlisterTripleSemigroupSemilattice@\texttt{McAlisterTripleSemigroupSemilattice}}
\label{McAlisterTripleSemigroupSemilattice}
}\hfill{\scriptsize (attribute)}}\\
\textbf{\indent Returns:\ }
A join\texttt{\symbol{45}}semilattice digraph.



 Returns the \texttt{IsJoinSemilatticeDigraph} (\textbf{Digraphs: IsJoinSemilatticeDigraph}) used to create the McAlister triple semigroup \mbox{\texttt{\mdseries\slshape S}} via \texttt{McAlisterTripleSemigroup} (\ref{McAlisterTripleSemigroup}). }

 

\subsection{\textcolor{Chapter }{McAlisterTripleSemigroupAction}}
\logpage{[ 8, 4, 6 ]}\nobreak
\hyperdef{L}{X86D6442E85881DEA}{}
{\noindent\textcolor{FuncColor}{$\triangleright$\enspace\texttt{McAlisterTripleSemigroupAction({\mdseries\slshape S})\index{McAlisterTripleSemigroupAction@\texttt{McAlisterTripleSemigroupAction}}
\label{McAlisterTripleSemigroupAction}
}\hfill{\scriptsize (attribute)}}\\
\textbf{\indent Returns:\ }
A function.



 Returns the action used to create the McAlister triple semigroup \mbox{\texttt{\mdseries\slshape S}} via \texttt{McAlisterTripleSemigroup} (\ref{McAlisterTripleSemigroup}). }

 

\subsection{\textcolor{Chapter }{IsMcAlisterTripleSemigroupElement}}
\logpage{[ 8, 4, 7 ]}\nobreak
\hyperdef{L}{X7B4EC9FC82249A83}{}
{\noindent\textcolor{FuncColor}{$\triangleright$\enspace\texttt{IsMcAlisterTripleSemigroupElement({\mdseries\slshape x})\index{IsMcAlisterTripleSemigroupElement@\texttt{IsMcAlisterTripleSemigroupElement}}
\label{IsMcAlisterTripleSemigroupElement}
}\hfill{\scriptsize (filter)}}\\
\noindent\textcolor{FuncColor}{$\triangleright$\enspace\texttt{IsMTSE({\mdseries\slshape x})\index{IsMTSE@\texttt{IsMTSE}}
\label{IsMTSE}
}\hfill{\scriptsize (filter)}}\\
\textbf{\indent Returns:\ }
\texttt{true} or \texttt{false}.



 Returns \texttt{true} if \mbox{\texttt{\mdseries\slshape x}} is an element of a McAlister triple semigroup; in particular, this returns \texttt{true} if \mbox{\texttt{\mdseries\slshape x}} has been created by \texttt{McAlisterTripleSemigroupElement} (\ref{McAlisterTripleSemigroupElement}). The functions \texttt{IsMTSE} and \texttt{IsMcAlisterTripleSemigroupElement} are synonyms. The mathematical description of these objects can be found in
the introduction to Chapter \ref{McAlister triple semigroups}. }

 

\subsection{\textcolor{Chapter }{McAlisterTripleSemigroupElement}}
\logpage{[ 8, 4, 8 ]}\nobreak
\hyperdef{L}{X854BFB1C7BA57985}{}
{\noindent\textcolor{FuncColor}{$\triangleright$\enspace\texttt{McAlisterTripleSemigroupElement({\mdseries\slshape S, A, g})\index{McAlisterTripleSemigroupElement@\texttt{McAlisterTripleSemigroupElement}}
\label{McAlisterTripleSemigroupElement}
}\hfill{\scriptsize (operation)}}\\
\noindent\textcolor{FuncColor}{$\triangleright$\enspace\texttt{MTSE({\mdseries\slshape S, A, g})\index{MTSE@\texttt{MTSE}}
\label{MTSE}
}\hfill{\scriptsize (operation)}}\\
\textbf{\indent Returns:\ }
A McAlister triple semigroup element.



 Returns the \emph{McAlister triple semigroup element} of the McAlister triple semigroup \mbox{\texttt{\mdseries\slshape S}} which corresponds to a label \mbox{\texttt{\mdseries\slshape A}} of a vertex from the \texttt{McAlisterTripleSemigroupSemilattice} (\ref{McAlisterTripleSemigroupSemilattice}) of \mbox{\texttt{\mdseries\slshape S}} and a group element \mbox{\texttt{\mdseries\slshape g}} of the \texttt{McAlisterTripleSemigroupGroup} (\ref{McAlisterTripleSemigroupGroup}) of \mbox{\texttt{\mdseries\slshape S}}. The pair \mbox{\texttt{\mdseries\slshape (A,g)}} only represents an element of \mbox{\texttt{\mdseries\slshape S}} if the following holds: \mbox{\texttt{\mdseries\slshape A}} acted on by the inverse of \mbox{\texttt{\mdseries\slshape g}} (via \texttt{McAlisterTripleSemigroupAction} (\ref{McAlisterTripleSemigroupAction})) is a vertex of the join\texttt{\symbol{45}}semilattice of \mbox{\texttt{\mdseries\slshape S}}. 

 The functions \texttt{MTSE} and \texttt{McAlisterTripleSemigroupElement} are synonyms. 
\begin{Verbatim}[commandchars=!@|,fontsize=\small,frame=single,label=Example]
  !gapprompt@gap>| !gapinput@x := Digraph([[1], [1, 2], [1, 2, 3], [1, 4], [1, 4, 5]]);|
  <immutable digraph with 5 vertices, 11 edges>
  !gapprompt@gap>| !gapinput@y := InducedSubdigraph(x, [1, 2, 3]);|
  <immutable digraph with 3 vertices, 6 edges>
  !gapprompt@gap>| !gapinput@A := AutomorphismGroup(x);|
  Group([ (2,4)(3,5) ])
  !gapprompt@gap>| !gapinput@S := McAlisterTripleSemigroup(A, x, y, OnPoints);|
  <McAlister triple semigroup over Group([ (2,4)(3,5) ])>
  !gapprompt@gap>| !gapinput@T := McAlisterTripleSemigroup(A, x, y);|
  <McAlister triple semigroup over Group([ (2,4)(3,5) ])>
  !gapprompt@gap>| !gapinput@S = T;|
  false
  !gapprompt@gap>| !gapinput@IsIsomorphicSemigroup(S, T);|
  true
  !gapprompt@gap>| !gapinput@a := MTSE(S, 1, (2, 4)(3, 5));|
  (1, (2,4)(3,5))
  !gapprompt@gap>| !gapinput@b := MTSE(S, 2, ());|
  (2, ())
  !gapprompt@gap>| !gapinput@a * a;|
  (1, ())
  !gapprompt@gap>| !gapinput@IsMTSE(a * a);|
  true
  !gapprompt@gap>| !gapinput@a = MTSE(T, 1, (2, 4)(3, 5));|
  false
  !gapprompt@gap>| !gapinput@a * b;|
  (1, (2,4)(3,5))
\end{Verbatim}
 }

 }

   }

 
\chapter{\textcolor{Chapter }{ Ideals }}\label{Ideals}
\logpage{[ 9, 0, 0 ]}
\hyperdef{L}{X83629803819C4A6F}{}
{
  In this chapter we describe the various ways that an ideal of a semigroup can
be created and manipulated in \textsf{Semigroups}. 

 We write \emph{ideal} to mean two\texttt{\symbol{45}}sided ideal everywhere in this chapter. 

 The methods in the \textsf{Semigroups} package apply to any ideal of a semigroup that is created using the function \texttt{SemigroupIdeal} (\ref{SemigroupIdeal}) or \texttt{SemigroupIdealByGenerators}. Anything that can be calculated for a semigroup defined by a generating set
can also be found for an ideal. This works particularly well for regular
ideals, since such an ideal can be represented using a similar data structure
to that used to represent a semigroup defined by a generating set but without
the necessity to find a generating set for the ideal. Many methods for
non\texttt{\symbol{45}}regular ideals rely on first finding a generating set
for the ideal, which can be costly (but not nearly as costly as an exhaustive
enumeration of the elements of the ideal). We plan to improve the
functionality of \textsf{Semigroups} for non\texttt{\symbol{45}}regular ideals in the future. 
\section{\textcolor{Chapter }{ Creating ideals }}\label{Creating ideals}
\logpage{[ 9, 1, 0 ]}
\hyperdef{L}{X82D4D9A578A56A8D}{}
{
  

\subsection{\textcolor{Chapter }{SemigroupIdeal}}
\logpage{[ 9, 1, 1 ]}\nobreak
\hyperdef{L}{X78E15B0184A1DC14}{}
{\noindent\textcolor{FuncColor}{$\triangleright$\enspace\texttt{SemigroupIdeal({\mdseries\slshape S, obj1, obj2, .., .})\index{SemigroupIdeal@\texttt{SemigroupIdeal}}
\label{SemigroupIdeal}
}\hfill{\scriptsize (function)}}\\
\textbf{\indent Returns:\ }
 An ideal of a semigroup. 



 If \mbox{\texttt{\mdseries\slshape obj1}}, \mbox{\texttt{\mdseries\slshape obj2}}, .. . are (any combination) of elements of the semigroup \mbox{\texttt{\mdseries\slshape S}} or collections of elements of \mbox{\texttt{\mdseries\slshape S}} (including subsemigroups and ideals of \mbox{\texttt{\mdseries\slshape S}}), then \texttt{SemigroupIdeal} returns the 2\texttt{\symbol{45}}sided ideal of the semigroup \mbox{\texttt{\mdseries\slshape S}} generated by the union of \mbox{\texttt{\mdseries\slshape obj1}}, \mbox{\texttt{\mdseries\slshape obj2}}, .. .. 

 The \texttt{Parent} (\textbf{Reference: Parent}) of the ideal returned by this function is \mbox{\texttt{\mdseries\slshape S}}. 
\begin{Verbatim}[commandchars=!@|,fontsize=\small,frame=single,label=Example]
  !gapprompt@gap>| !gapinput@S := SymmetricInverseMonoid(10);|
  <symmetric inverse monoid of degree 10>
  !gapprompt@gap>| !gapinput@I := SemigroupIdeal(S, PartialPerm([1, 2]));|
  <inverse partial perm semigroup ideal of rank 10 with 1 generator>
  !gapprompt@gap>| !gapinput@Size(I);|
  4151
  !gapprompt@gap>| !gapinput@I := SemigroupIdeal(S, I, Idempotents(S));|
  <inverse partial perm semigroup ideal of rank 10 with 1025 generators>
\end{Verbatim}
 }

 

\subsection{\textcolor{Chapter }{Ideals (for a semigroup)}}
\logpage{[ 9, 1, 2 ]}\nobreak
\hyperdef{L}{X7AF9B33881D185C6}{}
{\noindent\textcolor{FuncColor}{$\triangleright$\enspace\texttt{Ideals({\mdseries\slshape S})\index{Ideals@\texttt{Ideals}!for a semigroup}
\label{Ideals:for a semigroup}
}\hfill{\scriptsize (attribute)}}\\
\textbf{\indent Returns:\ }
 An list of ideals. 



 If \mbox{\texttt{\mdseries\slshape S}} is a finite non\texttt{\symbol{45}}empty semigroup, then this attribute
returns a list of the non\texttt{\symbol{45}}empty
two\texttt{\symbol{45}}sided ideals of \mbox{\texttt{\mdseries\slshape S}}. 

 The ideals are returned in no particular order, and each ideal uses the
minimum possible number of generators (see \texttt{GeneratorsOfSemigroupIdeal} (\ref{GeneratorsOfSemigroupIdeal})). 
\begin{Verbatim}[commandchars=!@|,fontsize=\small,frame=single,label=Example]
  !gapprompt@gap>| !gapinput@S := Semigroup([Transformation([4, 3, 4, 1]),|
  !gapprompt@>| !gapinput@                   Transformation([4, 3, 2, 2])]);|
  <transformation semigroup of degree 4 with 2 generators>
  !gapprompt@gap>| !gapinput@Ideals(S);|
  [ <non-regular transformation semigroup ideal of degree 4 with
        1 generator>,
   <non-regular transformation semigroup ideal of degree 4 with
       1 generator>,
  <non-regular transformation semigroup ideal of degree 4 with
      2 generators>,
  <regular transformation semigroup ideal of degree 4 with 1 generator>,
  <non-regular transformation semigroup ideal of degree 4 with
    1 generator>,
  <regular transformation semigroup ideal of degree 4 with 1 generator>
  ]
\end{Verbatim}
 }

 }

 
\section{\textcolor{Chapter }{ Attributes of ideals }}\label{Attributes of ideals}
\logpage{[ 9, 2, 0 ]}
\hyperdef{L}{X85D4E72B787B1C49}{}
{
  

\subsection{\textcolor{Chapter }{GeneratorsOfSemigroupIdeal}}
\logpage{[ 9, 2, 1 ]}\nobreak
\hyperdef{L}{X87BB45DB844D41BC}{}
{\noindent\textcolor{FuncColor}{$\triangleright$\enspace\texttt{GeneratorsOfSemigroupIdeal({\mdseries\slshape I})\index{GeneratorsOfSemigroupIdeal@\texttt{GeneratorsOfSemigroupIdeal}}
\label{GeneratorsOfSemigroupIdeal}
}\hfill{\scriptsize (attribute)}}\\
\textbf{\indent Returns:\ }
The generators of an ideal of a semigroup.



 This function returns the generators of the two\texttt{\symbol{45}}sided ideal \mbox{\texttt{\mdseries\slshape I}}, which were used to defined \mbox{\texttt{\mdseries\slshape I}} when it was created. 

 If \mbox{\texttt{\mdseries\slshape I}} is an ideal of a semigroup, then \mbox{\texttt{\mdseries\slshape I}} is defined to be the least 2\texttt{\symbol{45}}sided ideal of a semigroup \texttt{S} containing a set \texttt{J} of elements of \texttt{S}. The set \texttt{J} is said to \emph{generate} \mbox{\texttt{\mdseries\slshape I}}. 

 The notion of the generators of an ideal is distinct from the notion of the
generators of a semigroup or monoid. In particular, the semigroup generated by
the generators of an ideal is not, in general, equal to that ideal. Use \texttt{GeneratorsOfSemigroup} (\textbf{Reference: GeneratorsOfSemigroup}) to obtain a semigroup generating set for an ideal, but beware that this can be
very costly. 
\begin{Verbatim}[commandchars=!@|,fontsize=\small,frame=single,label=Example]
  !gapprompt@gap>| !gapinput@S := Semigroup(|
  !gapprompt@>| !gapinput@Bipartition([[1, 2, 3, 4, -1], [-2, -4], [-3]]),|
  !gapprompt@>| !gapinput@Bipartition([[1, 2, 3, -3], [4], [-1], [-2, -4]]),|
  !gapprompt@>| !gapinput@Bipartition([[1, 3, -2], [2, 4], [-1, -3, -4]]),|
  !gapprompt@>| !gapinput@Bipartition([[1], [2, 3, 4], [-1, -3, -4], [-2]]),|
  !gapprompt@>| !gapinput@Bipartition([[1], [2, 4, -2], [3, -4], [-1], [-3]]));;|
  !gapprompt@gap>| !gapinput@I := SemigroupIdeal(S, S.1 * S.2 * S.5);;|
  !gapprompt@gap>| !gapinput@GeneratorsOfSemigroupIdeal(I);|
  [ <bipartition: [ 1, 2, 3, 4, -4 ], [ -1 ], [ -2 ], [ -3 ]> ]
  !gapprompt@gap>| !gapinput@I = Semigroup(GeneratorsOfSemigroupIdeal(I));|
  false
\end{Verbatim}
 }

 

\subsection{\textcolor{Chapter }{MinimalIdealGeneratingSet}}
\logpage{[ 9, 2, 2 ]}\nobreak
\hyperdef{L}{X8777E71A82C2BAF9}{}
{\noindent\textcolor{FuncColor}{$\triangleright$\enspace\texttt{MinimalIdealGeneratingSet({\mdseries\slshape I})\index{MinimalIdealGeneratingSet@\texttt{MinimalIdealGeneratingSet}}
\label{MinimalIdealGeneratingSet}
}\hfill{\scriptsize (attribute)}}\\
\textbf{\indent Returns:\ }
A minimal set ideal generators of an ideal.



 This function returns a minimal set of elements of the parent of the semigroup
ideal \mbox{\texttt{\mdseries\slshape I}} required to generate \mbox{\texttt{\mdseries\slshape I}} as an ideal. 

 The notion of the generators of an ideal is distinct from the notion of the
generators of a semigroup or monoid. In particular, the semigroup generated by
the generators of an ideal is not, in general, equal to that ideal. Use \texttt{GeneratorsOfSemigroup} (\textbf{Reference: GeneratorsOfSemigroup}) to obtain a semigroup generating set for an ideal, but beware that this can be
very costly. 
\begin{Verbatim}[commandchars=!@|,fontsize=\small,frame=single,label=Example]
  !gapprompt@gap>| !gapinput@S := Monoid([|
  !gapprompt@>| !gapinput@Bipartition([[1, 2, 3, -2], [4], [-1, -4], [-3]]),|
  !gapprompt@>| !gapinput@Bipartition([[1, 4, -2, -4], [2, -1, -3], [3]])]);;|
  !gapprompt@gap>| !gapinput@I := SemigroupIdeal(S, S);;|
  !gapprompt@gap>| !gapinput@MinimalIdealGeneratingSet(I);|
  [ <block bijection: [ 1, -1 ], [ 2, -2 ], [ 3, -3 ], [ 4, -4 ]> ]
\end{Verbatim}
 }

 

\subsection{\textcolor{Chapter }{SupersemigroupOfIdeal}}
\logpage{[ 9, 2, 3 ]}\nobreak
\hyperdef{L}{X7DB8699784FA4114}{}
{\noindent\textcolor{FuncColor}{$\triangleright$\enspace\texttt{SupersemigroupOfIdeal({\mdseries\slshape I})\index{SupersemigroupOfIdeal@\texttt{SupersemigroupOfIdeal}}
\label{SupersemigroupOfIdeal}
}\hfill{\scriptsize (attribute)}}\\
\textbf{\indent Returns:\ }
 An ideal of a semigroup. 



 The \texttt{Parent} (\textbf{Reference: Parent}) of an ideal is the semigroup in which the ideal was created, i.e. the first
argument of \texttt{SemigroupIdeal} (\ref{SemigroupIdeal}) or \texttt{SemigroupIdealByGenerators}. This function returns the semigroup containing the generators of the
semigroup (i.e. \texttt{GeneratorsOfSemigroup} (\textbf{Reference: GeneratorsOfSemigroup})) which are used to compute the ideal. 

 For a regular semigroup ideal, \texttt{SupersemigroupOfIdeal} will always be the top most semigroup used to create any of the predecessors
of the current ideal. For example, if \texttt{S} is a semigroup, \texttt{I} is a regular ideal of \texttt{S}, and \texttt{J} is an ideal of \texttt{I}, then \texttt{Parent(J)} is \texttt{I} and \texttt{SupersemigroupOfIdeal(J)} is \texttt{S}. This is to avoid computing a generating set for \texttt{I}, in this example, which is expensive and unnecessary since \texttt{I} is regular (in which case the Green's relations of \texttt{I} are just restrictions of the Green's relations on \texttt{S}). 

 If \texttt{S} is a semigroup, \texttt{I} is a non\texttt{\symbol{45}}regular ideal of \texttt{S}, \texttt{J} is an ideal of \texttt{I}, then \texttt{SupersemigroupOfIdeal(J)} is \texttt{I}, since we currently have to use \texttt{GeneratorsOfSemigroup(I)} to compute anything about \texttt{I} other than its size and membership. 
\begin{Verbatim}[commandchars=!@|,fontsize=\small,frame=single,label=Example]
  !gapprompt@gap>| !gapinput@S := FullTransformationSemigroup(8);|
  <full transformation monoid of degree 8>
  !gapprompt@gap>| !gapinput@x := Transformation([2, 6, 7, 2, 6, 1, 1, 5]);;|
  !gapprompt@gap>| !gapinput@D := DClass(S, x);|
  <Green's D-class: Transformation( [ 2, 6, 7, 2, 6, 1, 1, 5 ] )>
  !gapprompt@gap>| !gapinput@R := PrincipalFactor(D);|
  <Rees 0-matrix semigroup 1050x56 over Group([ (2,8,7,4,3), (3,4) ])>
  !gapprompt@gap>| !gapinput@S := Semigroup(List([1 .. 10], x -> Random(R)));|
  <subsemigroup of 1050x56 Rees 0-matrix semigroup with 10 generators>
  !gapprompt@gap>| !gapinput@I := SemigroupIdeal(S, MultiplicativeZero(S));|
  <regular Rees 0-matrix semigroup ideal with 1 generator>
  !gapprompt@gap>| !gapinput@SupersemigroupOfIdeal(I);|
  <subsemigroup of 1050x56 Rees 0-matrix semigroup with 10 generators>
  !gapprompt@gap>| !gapinput@J := SemigroupIdeal(I, Representative(MinimalDClass(S)));|
  <regular Rees 0-matrix semigroup ideal with 1 generator>
  !gapprompt@gap>| !gapinput@Parent(J) = I;|
  true
  !gapprompt@gap>| !gapinput@SupersemigroupOfIdeal(J) = I;|
  false
\end{Verbatim}
 }

 }

 }

 
\chapter{\textcolor{Chapter }{ Green's relations }}\label{Green's relations}
\logpage{[ 10, 0, 0 ]}
\hyperdef{L}{X80C6C718801855E9}{}
{
  In this chapter we describe the functions in \textsf{Semigroups} for computing Green's classes and related properties of semigroups.   
\section{\textcolor{Chapter }{ Creating Green's classes and representatives }}\logpage{[ 10, 1, 0 ]}
\hyperdef{L}{X788D6753849BAD7C}{}
{
  In this section, we describe the methods in the \textsf{Semigroups} package for creating Green's classes. 

 
\subsection{\textcolor{Chapter }{XClassOfYClass}}\logpage{[ 10, 1, 1 ]}
\hyperdef{L}{X87558FEF805D24E1}{}
{
\noindent\textcolor{FuncColor}{$\triangleright$\enspace\texttt{DClassOfHClass({\mdseries\slshape class})\index{DClassOfHClass@\texttt{DClassOfHClass}}
\label{DClassOfHClass}
}\hfill{\scriptsize (method)}}\\
\noindent\textcolor{FuncColor}{$\triangleright$\enspace\texttt{DClassOfLClass({\mdseries\slshape class})\index{DClassOfLClass@\texttt{DClassOfLClass}}
\label{DClassOfLClass}
}\hfill{\scriptsize (method)}}\\
\noindent\textcolor{FuncColor}{$\triangleright$\enspace\texttt{DClassOfRClass({\mdseries\slshape class})\index{DClassOfRClass@\texttt{DClassOfRClass}}
\label{DClassOfRClass}
}\hfill{\scriptsize (method)}}\\
\noindent\textcolor{FuncColor}{$\triangleright$\enspace\texttt{LClassOfHClass({\mdseries\slshape class})\index{LClassOfHClass@\texttt{LClassOfHClass}}
\label{LClassOfHClass}
}\hfill{\scriptsize (method)}}\\
\noindent\textcolor{FuncColor}{$\triangleright$\enspace\texttt{RClassOfHClass({\mdseries\slshape class})\index{RClassOfHClass@\texttt{RClassOfHClass}}
\label{RClassOfHClass}
}\hfill{\scriptsize (method)}}\\
\textbf{\indent Returns:\ }
A Green's class.



 \texttt{XClassOfYClass} returns the \texttt{X}\texttt{\symbol{45}}class containing the \texttt{Y}\texttt{\symbol{45}}class \mbox{\texttt{\mdseries\slshape class}} where \texttt{X} and \texttt{Y} should be replaced by an appropriate choice of \texttt{D, H, L,} and \texttt{R}.

 Note that if it is not known to \textsf{GAP} whether or not the representative of \mbox{\texttt{\mdseries\slshape class}} is an element of the semigroup containing \mbox{\texttt{\mdseries\slshape class}}, then no attempt is made to check this.

 The same result can be produced using: 
\begin{Verbatim}[commandchars=!@|,fontsize=\small,frame=single,label=Example]
  First(GreensXClasses(S), x -> Representative(x) in class);
\end{Verbatim}
 but this might be substantially slower. Note that \texttt{XClassOfYClass} is also likely to be faster than 
\begin{Verbatim}[commandchars=!@|,fontsize=\small,frame=single,label=Example]
  GreensXClassOfElement(S, Representative(class));
\end{Verbatim}
 

 \texttt{DClass} can also be used as a synonym for \texttt{DClassOfHClass}, \texttt{DClassOfLClass}, and \texttt{DClassOfRClass}; \texttt{LClass} as a synonym for \texttt{LClassOfHClass}; and \texttt{RClass} as a synonym for \texttt{RClassOfHClass}. See also \texttt{GreensDClassOfElement} (\textbf{Reference: GreensDClassOfElement}) and \texttt{GreensDClassOfElementNC} (\ref{GreensDClassOfElementNC}). 
\begin{Verbatim}[commandchars=!@|,fontsize=\small,frame=single,label=Example]
  !gapprompt@gap>| !gapinput@S := Semigroup(Transformation([1, 3, 2]),|
  !gapprompt@>| !gapinput@                  Transformation([2, 1, 3]),|
  !gapprompt@>| !gapinput@                  Transformation([3, 2, 1]),|
  !gapprompt@>| !gapinput@                  Transformation([1, 3, 1]));;|
  !gapprompt@gap>| !gapinput@R := GreensRClassOfElement(S, Transformation([3, 2, 1]));|
  <Green's R-class: Transformation( [ 3, 2, 1 ] )>
  !gapprompt@gap>| !gapinput@DClassOfRClass(R);|
  <Green's D-class: Transformation( [ 3, 2, 1 ] )>
  !gapprompt@gap>| !gapinput@IsGreensDClass(DClassOfRClass(R));|
  true
  !gapprompt@gap>| !gapinput@S := InverseSemigroup(|
  !gapprompt@>| !gapinput@PartialPerm([2, 6, 7, 0, 0, 9, 0, 1, 0, 5]),|
  !gapprompt@>| !gapinput@PartialPerm([3, 8, 1, 9, 0, 4, 10, 5, 0, 6]));|
  <inverse partial perm semigroup of rank 10 with 2 generators>
  !gapprompt@gap>| !gapinput@x := S.1;|
  [3,7][8,1,2,6,9][10,5]
  !gapprompt@gap>| !gapinput@H := HClass(S, x);|
  <Green's H-class: [3,7][8,1,2,6,9][10,5]>
  !gapprompt@gap>| !gapinput@R := RClassOfHClass(H);|
  <Green's R-class: [3,7][8,1,2,6,9][10,5]>
  !gapprompt@gap>| !gapinput@L := LClass(H);;|
  !gapprompt@gap>| !gapinput@L = LClass(S, PartialPerm([1, 2, 0, 0, 5, 6, 7, 0, 9]));|
  true
  !gapprompt@gap>| !gapinput@DClass(R) = DClass(L);|
  true
  !gapprompt@gap>| !gapinput@DClass(H) = DClass(L);|
  true
\end{Verbatim}
 }

 
\subsection{\textcolor{Chapter }{GreensXClassOfElement}}\logpage{[ 10, 1, 2 ]}
\hyperdef{L}{X81B7AD4C7C552867}{}
{
\noindent\textcolor{FuncColor}{$\triangleright$\enspace\texttt{GreensDClassOfElement({\mdseries\slshape X, f})\index{GreensDClassOfElement@\texttt{GreensDClassOfElement}}
\label{GreensDClassOfElement}
}\hfill{\scriptsize (operation)}}\\
\noindent\textcolor{FuncColor}{$\triangleright$\enspace\texttt{DClass({\mdseries\slshape X, f})\index{DClass@\texttt{DClass}}
\label{DClass}
}\hfill{\scriptsize (operation)}}\\
\noindent\textcolor{FuncColor}{$\triangleright$\enspace\texttt{GreensHClassOfElement({\mdseries\slshape X, f})\index{GreensHClassOfElement@\texttt{GreensHClassOfElement}}
\label{GreensHClassOfElement}
}\hfill{\scriptsize (operation)}}\\
\noindent\textcolor{FuncColor}{$\triangleright$\enspace\texttt{GreensHClassOfElement({\mdseries\slshape R, i, j})\index{GreensHClassOfElement@\texttt{GreensHClassOfElement}!for a Rees matrix semigroup}
\label{GreensHClassOfElement:for a Rees matrix semigroup}
}\hfill{\scriptsize (operation)}}\\
\noindent\textcolor{FuncColor}{$\triangleright$\enspace\texttt{HClass({\mdseries\slshape X, f})\index{HClass@\texttt{HClass}}
\label{HClass}
}\hfill{\scriptsize (operation)}}\\
\noindent\textcolor{FuncColor}{$\triangleright$\enspace\texttt{HClass({\mdseries\slshape R, i, j})\index{HClass@\texttt{HClass}!for a Rees matrix semigroup}
\label{HClass:for a Rees matrix semigroup}
}\hfill{\scriptsize (operation)}}\\
\noindent\textcolor{FuncColor}{$\triangleright$\enspace\texttt{GreensLClassOfElement({\mdseries\slshape X, f})\index{GreensLClassOfElement@\texttt{GreensLClassOfElement}}
\label{GreensLClassOfElement}
}\hfill{\scriptsize (operation)}}\\
\noindent\textcolor{FuncColor}{$\triangleright$\enspace\texttt{LClass({\mdseries\slshape X, f})\index{LClass@\texttt{LClass}}
\label{LClass}
}\hfill{\scriptsize (operation)}}\\
\noindent\textcolor{FuncColor}{$\triangleright$\enspace\texttt{GreensRClassOfElement({\mdseries\slshape X, f})\index{GreensRClassOfElement@\texttt{GreensRClassOfElement}}
\label{GreensRClassOfElement}
}\hfill{\scriptsize (operation)}}\\
\noindent\textcolor{FuncColor}{$\triangleright$\enspace\texttt{RClass({\mdseries\slshape X, f})\index{RClass@\texttt{RClass}}
\label{RClass}
}\hfill{\scriptsize (operation)}}\\
\textbf{\indent Returns:\ }
A Green's class.



 These functions produce essentially the same output as the \textsf{GAP} library functions with the same names; see \texttt{GreensDClassOfElement} (\textbf{Reference: GreensDClassOfElement}). The main difference is that these functions can be applied to a wider class
of objects: 
\begin{description}
\item[{\texttt{GreensDClassOfElement} and \texttt{DClass}}]  \mbox{\texttt{\mdseries\slshape X}} must be a semigroup. 
\item[{\texttt{GreensHClassOfElement} and \texttt{HClass}}]  \mbox{\texttt{\mdseries\slshape X}} can be a semigroup, $\mathcal{R}$\texttt{\symbol{45}}class, $\mathcal{L}$\texttt{\symbol{45}}class, or $\mathcal{D}$\texttt{\symbol{45}}class.  If \mbox{\texttt{\mdseries\slshape R}} is a \mbox{\texttt{\mdseries\slshape IxJ}} Rees matrix semigroup or a Rees 0\texttt{\symbol{45}}matrix semigroup, and \mbox{\texttt{\mdseries\slshape i}} and \mbox{\texttt{\mdseries\slshape j}} are integers of the corresponding index sets, then \texttt{GreensHClassOfElement} returns the $\mathcal{H}$\texttt{\symbol{45}}class in row \mbox{\texttt{\mdseries\slshape i}} and column \mbox{\texttt{\mdseries\slshape j}}. 
\item[{\texttt{GreensLClassOfElement} and \texttt{LClass}}]  \mbox{\texttt{\mdseries\slshape X}} can be a semigroup or $\mathcal{D}$\texttt{\symbol{45}}class. 
\item[{\texttt{GreensRClassOfElement} and \texttt{RClass}}]  \mbox{\texttt{\mdseries\slshape X}} can be a semigroup or $\mathcal{D}$\texttt{\symbol{45}}class. 
\end{description}
 Note that \texttt{GreensXClassOfElement} and \texttt{XClass} are synonyms and have identical output. The shorter command is provided for
the sake of convenience.

 }

 
\subsection{\textcolor{Chapter }{GreensXClassOfElementNC}}\logpage{[ 10, 1, 3 ]}
\hyperdef{L}{X7B44317786571F8B}{}
{
\noindent\textcolor{FuncColor}{$\triangleright$\enspace\texttt{GreensDClassOfElementNC({\mdseries\slshape X, f})\index{GreensDClassOfElementNC@\texttt{GreensDClassOfElementNC}}
\label{GreensDClassOfElementNC}
}\hfill{\scriptsize (operation)}}\\
\noindent\textcolor{FuncColor}{$\triangleright$\enspace\texttt{DClassNC({\mdseries\slshape X, f})\index{DClassNC@\texttt{DClassNC}}
\label{DClassNC}
}\hfill{\scriptsize (operation)}}\\
\noindent\textcolor{FuncColor}{$\triangleright$\enspace\texttt{GreensHClassOfElementNC({\mdseries\slshape X, f})\index{GreensHClassOfElementNC@\texttt{GreensHClassOfElementNC}}
\label{GreensHClassOfElementNC}
}\hfill{\scriptsize (operation)}}\\
\noindent\textcolor{FuncColor}{$\triangleright$\enspace\texttt{HClassNC({\mdseries\slshape X, f})\index{HClassNC@\texttt{HClassNC}}
\label{HClassNC}
}\hfill{\scriptsize (operation)}}\\
\noindent\textcolor{FuncColor}{$\triangleright$\enspace\texttt{GreensLClassOfElementNC({\mdseries\slshape X, f})\index{GreensLClassOfElementNC@\texttt{GreensLClassOfElementNC}}
\label{GreensLClassOfElementNC}
}\hfill{\scriptsize (operation)}}\\
\noindent\textcolor{FuncColor}{$\triangleright$\enspace\texttt{LClassNC({\mdseries\slshape X, f})\index{LClassNC@\texttt{LClassNC}}
\label{LClassNC}
}\hfill{\scriptsize (operation)}}\\
\noindent\textcolor{FuncColor}{$\triangleright$\enspace\texttt{GreensRClassOfElementNC({\mdseries\slshape X, f})\index{GreensRClassOfElementNC@\texttt{GreensRClassOfElementNC}}
\label{GreensRClassOfElementNC}
}\hfill{\scriptsize (operation)}}\\
\noindent\textcolor{FuncColor}{$\triangleright$\enspace\texttt{RClassNC({\mdseries\slshape X, f})\index{RClassNC@\texttt{RClassNC}}
\label{RClassNC}
}\hfill{\scriptsize (operation)}}\\
\textbf{\indent Returns:\ }
A Green's class.



 These functions are essentially the same as \texttt{GreensDClassOfElement} (\ref{GreensDClassOfElement}) except that no effort is made to verify if \mbox{\texttt{\mdseries\slshape f}} is an element of \mbox{\texttt{\mdseries\slshape X}}. More precisely, \texttt{GreensXClassOfElementNC} and \texttt{XClassNC} first check if \mbox{\texttt{\mdseries\slshape f}} has already been shown to be an element of \mbox{\texttt{\mdseries\slshape X}}. If it is not known to \textsf{GAP} if \mbox{\texttt{\mdseries\slshape f}} is an element of \mbox{\texttt{\mdseries\slshape X}}, then no further attempt to verify this is made. 

 Note that \texttt{GreensXClassOfElementNC} and \texttt{XClassNC} are synonyms and have identical output. The shorter command is provided for
the sake of convenience. 

 It can be quicker to compute the class of an element using \texttt{GreensRClassOfElementNC}, say, than using \texttt{GreensRClassOfElement} if it is known \emph{a priori} that \mbox{\texttt{\mdseries\slshape f}} is an element of \mbox{\texttt{\mdseries\slshape X}}. On the other hand, if \mbox{\texttt{\mdseries\slshape f}} is not an element of \mbox{\texttt{\mdseries\slshape X}}, then the results of this computation are unpredictable.

 For example, if 
\begin{Verbatim}[commandchars=!@|,fontsize=\small,frame=single,label=Example]
  x := Transformation([15, 18, 20, 20, 20, 20, 20, 20, 20, 20, 20, 20, 20, 20, 20, 20, 20, 20, 20, 20]);
\end{Verbatim}
 in the semigroup \mbox{\texttt{\mdseries\slshape X}} of order\texttt{\symbol{45}}preserving mappings on 20 points, then 
\begin{Verbatim}[commandchars=!@|,fontsize=\small,frame=single,label=Example]
  GreensRClassOfElementNC(X, x);
\end{Verbatim}
 returns an answer relatively quickly, whereas 
\begin{Verbatim}[commandchars=!@|,fontsize=\small,frame=single,label=Example]
  GreensRClassOfElement(X, x)
\end{Verbatim}
 can take a significant amount of time to return a value.

 See also \texttt{GreensRClassOfElement} (\textbf{Reference: GreensRClassOfElement}) and \texttt{RClassOfHClass} (\ref{RClassOfHClass}). 
\begin{Verbatim}[commandchars=!@|,fontsize=\small,frame=single,label=Example]
  !gapprompt@gap>| !gapinput@S := RandomSemigroup(IsTransformationSemigroup, 2, 1000);;|
  !gapprompt@gap>| !gapinput@x := [1, 1, 2, 2, 2, 1, 1, 1, 1, 1, 2, 2, 2, 2, 1, 1, 2, 2, 1];;|
  !gapprompt@gap>| !gapinput@x := EvaluateWord(Generators(S), x);;|
  !gapprompt@gap>| !gapinput@R := GreensRClassOfElementNC(S, x);;|
  !gapprompt@gap>| !gapinput@Size(R);|
  1
  !gapprompt@gap>| !gapinput@L := GreensLClassOfElementNC(S, x);;|
  !gapprompt@gap>| !gapinput@Size(L);|
  1
  !gapprompt@gap>| !gapinput@x := PartialPerm([2, 3, 4, 5, 0, 0, 6, 8, 10, 11]);;|
  !gapprompt@gap>| !gapinput@L := LClass(POI(11), x);|
  <Green's L-class: [1,2,3,4,5][7,6][9,10,11](8)>
  !gapprompt@gap>| !gapinput@Size(L);|
  165
\end{Verbatim}
 }

 
\subsection{\textcolor{Chapter }{GreensXClasses}}\label{GreensXClasses}
\logpage{[ 10, 1, 4 ]}
\hyperdef{L}{X7D51218A80234DE5}{}
{
\noindent\textcolor{FuncColor}{$\triangleright$\enspace\texttt{GreensDClasses({\mdseries\slshape obj})\index{GreensDClasses@\texttt{GreensDClasses}}
\label{GreensDClasses}
}\hfill{\scriptsize (method)}}\\
\noindent\textcolor{FuncColor}{$\triangleright$\enspace\texttt{DClasses({\mdseries\slshape obj})\index{DClasses@\texttt{DClasses}}
\label{DClasses}
}\hfill{\scriptsize (method)}}\\
\noindent\textcolor{FuncColor}{$\triangleright$\enspace\texttt{GreensHClasses({\mdseries\slshape obj})\index{GreensHClasses@\texttt{GreensHClasses}}
\label{GreensHClasses}
}\hfill{\scriptsize (method)}}\\
\noindent\textcolor{FuncColor}{$\triangleright$\enspace\texttt{HClasses({\mdseries\slshape obj})\index{HClasses@\texttt{HClasses}}
\label{HClasses}
}\hfill{\scriptsize (method)}}\\
\noindent\textcolor{FuncColor}{$\triangleright$\enspace\texttt{GreensJClasses({\mdseries\slshape obj})\index{GreensJClasses@\texttt{GreensJClasses}}
\label{GreensJClasses}
}\hfill{\scriptsize (method)}}\\
\noindent\textcolor{FuncColor}{$\triangleright$\enspace\texttt{JClasses({\mdseries\slshape obj})\index{JClasses@\texttt{JClasses}}
\label{JClasses}
}\hfill{\scriptsize (method)}}\\
\noindent\textcolor{FuncColor}{$\triangleright$\enspace\texttt{GreensLClasses({\mdseries\slshape obj})\index{GreensLClasses@\texttt{GreensLClasses}}
\label{GreensLClasses}
}\hfill{\scriptsize (method)}}\\
\noindent\textcolor{FuncColor}{$\triangleright$\enspace\texttt{LClasses({\mdseries\slshape obj})\index{LClasses@\texttt{LClasses}}
\label{LClasses}
}\hfill{\scriptsize (method)}}\\
\noindent\textcolor{FuncColor}{$\triangleright$\enspace\texttt{GreensRClasses({\mdseries\slshape obj})\index{GreensRClasses@\texttt{GreensRClasses}}
\label{GreensRClasses}
}\hfill{\scriptsize (method)}}\\
\noindent\textcolor{FuncColor}{$\triangleright$\enspace\texttt{RClasses({\mdseries\slshape obj})\index{RClasses@\texttt{RClasses}}
\label{RClasses}
}\hfill{\scriptsize (method)}}\\
\textbf{\indent Returns:\ }
A list of Green's classes. 



 These functions produce essentially the same output as the \textsf{GAP} library functions with the same names; see \texttt{GreensDClasses} (\textbf{Reference: GreensDClasses}). The main difference is that these functions can be applied to a wider class
of objects: 
\begin{description}
\item[{\texttt{GreensDClasses} and \texttt{DClasses}}]  \mbox{\texttt{\mdseries\slshape X}} should be a semigroup. 
\item[{\texttt{GreensHClasses} and \texttt{HClasses}}]  \mbox{\texttt{\mdseries\slshape X}} can be a semigroup, $\mathcal{R}$\texttt{\symbol{45}}class, $\mathcal{L}$\texttt{\symbol{45}}class, or $\mathcal{D}$\texttt{\symbol{45}}class. 
\item[{\texttt{GreensLClasses} and \texttt{LClasses}}]  \mbox{\texttt{\mdseries\slshape X}} can be a semigroup or $\mathcal{D}$\texttt{\symbol{45}}class. 
\item[{\texttt{GreensRClasses} and \texttt{RClasses}}]  \mbox{\texttt{\mdseries\slshape X}} can be a semigroup or $\mathcal{D}$\texttt{\symbol{45}}class. 
\end{description}
 Note that \texttt{GreensXClasses} and \texttt{XClasses} are synonyms and have identical output. The shorter command is provided for
the sake of convenience.

 See also \texttt{DClassReps} (\ref{DClassReps}), \texttt{IteratorOfDClassReps} (\ref{IteratorOfDClassReps}), \texttt{IteratorOfDClasses} (\ref{IteratorOfDClasses}), and \texttt{NrDClasses} (\ref{NrDClasses}). 
\begin{Verbatim}[commandchars=!@|,fontsize=\small,frame=single,label=Example]
  !gapprompt@gap>| !gapinput@S := Semigroup(Transformation([3, 4, 4, 4]),|
  !gapprompt@>| !gapinput@                  Transformation([4, 3, 1, 2]));;|
  !gapprompt@gap>| !gapinput@GreensDClasses(S);|
  [ <Green's D-class: Transformation( [ 3, 4, 4, 4 ] )>,
    <Green's D-class: Transformation( [ 4, 3, 1, 2 ] )>,
    <Green's D-class: Transformation( [ 4, 4, 4, 4 ] )> ]
  !gapprompt@gap>| !gapinput@GreensRClasses(S);|
  [ <Green's R-class: Transformation( [ 3, 4, 4, 4 ] )>,
    <Green's R-class: Transformation( [ 4, 3, 1, 2 ] )>,
    <Green's R-class: Transformation( [ 4, 4, 4, 4 ] )>,
    <Green's R-class: Transformation( [ 4, 4, 3, 4 ] )>,
    <Green's R-class: Transformation( [ 4, 3, 4, 4 ] )>,
    <Green's R-class: Transformation( [ 4, 4, 4, 3 ] )> ]
  !gapprompt@gap>| !gapinput@D := GreensDClasses(S)[1];|
  <Green's D-class: Transformation( [ 3, 4, 4, 4 ] )>
  !gapprompt@gap>| !gapinput@GreensLClasses(D);|
  [ <Green's L-class: Transformation( [ 3, 4, 4, 4 ] )>,
    <Green's L-class: Transformation( [ 1, 2, 2, 2 ] )> ]
  !gapprompt@gap>| !gapinput@GreensRClasses(D);|
  [ <Green's R-class: Transformation( [ 3, 4, 4, 4 ] )>,
    <Green's R-class: Transformation( [ 4, 4, 3, 4 ] )>,
    <Green's R-class: Transformation( [ 4, 3, 4, 4 ] )>,
    <Green's R-class: Transformation( [ 4, 4, 4, 3 ] )> ]
  !gapprompt@gap>| !gapinput@R := GreensRClasses(D)[1];|
  <Green's R-class: Transformation( [ 3, 4, 4, 4 ] )>
  !gapprompt@gap>| !gapinput@GreensHClasses(R);|
  [ <Green's H-class: Transformation( [ 3, 4, 4, 4 ] )>,
    <Green's H-class: Transformation( [ 1, 2, 2, 2 ] )> ]
  !gapprompt@gap>| !gapinput@S := InverseSemigroup([|
  !gapprompt@>| !gapinput@PartialPerm([2, 4, 1]), PartialPerm([3, 0, 4, 1])]);;|
  !gapprompt@gap>| !gapinput@GreensDClasses(S);|
  [ <Green's D-class: <identity partial perm on [ 1, 2, 4 ]>>,
    <Green's D-class: <identity partial perm on [ 1, 3, 4 ]>>,
    <Green's D-class: <identity partial perm on [ 1, 3 ]>>,
    <Green's D-class: <identity partial perm on [ 4 ]>>,
    <Green's D-class: <empty partial perm>> ]
  !gapprompt@gap>| !gapinput@GreensLClasses(S);|
  [ <Green's L-class: <identity partial perm on [ 1, 2, 4 ]>>,
    <Green's L-class: [4,2,1,3]>,
    <Green's L-class: <identity partial perm on [ 1, 3, 4 ]>>,
    <Green's L-class: <identity partial perm on [ 1, 3 ]>>,
    <Green's L-class: [3,1,2]>, <Green's L-class: [1,4][3,2]>,
    <Green's L-class: [1,3,4]>, <Green's L-class: [3,1,4]>,
    <Green's L-class: [1,2](3)>,
    <Green's L-class: <identity partial perm on [ 4 ]>>,
    <Green's L-class: [4,1]>, <Green's L-class: [4,3]>,
    <Green's L-class: [4,2]>, <Green's L-class: <empty partial perm>> ]
  !gapprompt@gap>| !gapinput@D := GreensDClasses(S)[3];|
  <Green's D-class: <identity partial perm on [ 1, 3 ]>>
  !gapprompt@gap>| !gapinput@GreensLClasses(D);|
  [ <Green's L-class: <identity partial perm on [ 1, 3 ]>>,
    <Green's L-class: [3,1,2]>, <Green's L-class: [1,4][3,2]>,
    <Green's L-class: [1,3,4]>, <Green's L-class: [3,1,4]>,
    <Green's L-class: [1,2](3)> ]
  !gapprompt@gap>| !gapinput@GreensRClasses(D);|
  [ <Green's R-class: <identity partial perm on [ 1, 3 ]>>,
    <Green's R-class: [2,1,3]>, <Green's R-class: [2,3][4,1]>,
    <Green's R-class: [4,3,1]>, <Green's R-class: [4,1,3]>,
    <Green's R-class: [2,1](3)> ]
\end{Verbatim}
 }

 
\subsection{\textcolor{Chapter }{XClassReps}}\logpage{[ 10, 1, 5 ]}
\hyperdef{L}{X865387A87FAAC395}{}
{
\noindent\textcolor{FuncColor}{$\triangleright$\enspace\texttt{DClassReps({\mdseries\slshape obj})\index{DClassReps@\texttt{DClassReps}}
\label{DClassReps}
}\hfill{\scriptsize (attribute)}}\\
\noindent\textcolor{FuncColor}{$\triangleright$\enspace\texttt{HClassReps({\mdseries\slshape obj})\index{HClassReps@\texttt{HClassReps}}
\label{HClassReps}
}\hfill{\scriptsize (attribute)}}\\
\noindent\textcolor{FuncColor}{$\triangleright$\enspace\texttt{LClassReps({\mdseries\slshape obj})\index{LClassReps@\texttt{LClassReps}}
\label{LClassReps}
}\hfill{\scriptsize (attribute)}}\\
\noindent\textcolor{FuncColor}{$\triangleright$\enspace\texttt{RClassReps({\mdseries\slshape obj})\index{RClassReps@\texttt{RClassReps}}
\label{RClassReps}
}\hfill{\scriptsize (attribute)}}\\
\textbf{\indent Returns:\ }
A list of representatives.



 \texttt{XClassReps} returns a list of the representatives of the Green's classes of \mbox{\texttt{\mdseries\slshape obj}}, which can be a semigroup, $\mathcal{D}$\texttt{\symbol{45}}, $\mathcal{L}$\texttt{\symbol{45}}, or $\mathcal{R}$\texttt{\symbol{45}}class where appropriate.

 The same output can be obtained by calling, for example: 
\begin{Verbatim}[commandchars=!@|,fontsize=\small,frame=single,label=Example]
  List(GreensXClasses(obj), Representative);
\end{Verbatim}
 Note that if the Green's classes themselves are not required, then \texttt{XClassReps} will return an answer more quickly than the above, since the Green's class
objects are not created.

 See also \texttt{GreensDClasses} (\ref{GreensDClasses}), \texttt{IteratorOfDClassReps} (\ref{IteratorOfDClassReps}), \texttt{IteratorOfDClasses} (\ref{IteratorOfDClasses}), and \texttt{NrDClasses} (\ref{NrDClasses}). 
\begin{Verbatim}[commandchars=!@|,fontsize=\small,frame=single,label=Example]
  !gapprompt@gap>| !gapinput@S := Semigroup(Transformation([3, 4, 4, 4]),|
  !gapprompt@>| !gapinput@                  Transformation([4, 3, 1, 2]));;|
  !gapprompt@gap>| !gapinput@DClassReps(S);|
  [ Transformation( [ 3, 4, 4, 4 ] ), Transformation( [ 4, 3, 1, 2 ] ),
    Transformation( [ 4, 4, 4, 4 ] ) ]
  !gapprompt@gap>| !gapinput@LClassReps(S);|
  [ Transformation( [ 3, 4, 4, 4 ] ), Transformation( [ 1, 2, 2, 2 ] ),
    Transformation( [ 4, 3, 1, 2 ] ), Transformation( [ 4, 4, 4, 4 ] ),
    Transformation( [ 2, 2, 2, 2 ] ), Transformation( [ 3, 3, 3, 3 ] ),
    Transformation( [ 1, 1, 1, 1 ] ) ]
  !gapprompt@gap>| !gapinput@D := GreensDClasses(S)[1];|
  <Green's D-class: Transformation( [ 3, 4, 4, 4 ] )>
  !gapprompt@gap>| !gapinput@LClassReps(D);|
  [ Transformation( [ 3, 4, 4, 4 ] ), Transformation( [ 1, 2, 2, 2 ] ) ]
  !gapprompt@gap>| !gapinput@RClassReps(D);|
  [ Transformation( [ 3, 4, 4, 4 ] ), Transformation( [ 4, 4, 3, 4 ] ),
    Transformation( [ 4, 3, 4, 4 ] ), Transformation( [ 4, 4, 4, 3 ] ) ]
  !gapprompt@gap>| !gapinput@R := GreensRClasses(D)[1];;|
  !gapprompt@gap>| !gapinput@HClassReps(R);|
  [ Transformation( [ 3, 4, 4, 4 ] ), Transformation( [ 1, 2, 2, 2 ] ) ]
  !gapprompt@gap>| !gapinput@S := SymmetricInverseSemigroup(6);;|
  !gapprompt@gap>| !gapinput@e := InverseSemigroup(Idempotents(S));;|
  !gapprompt@gap>| !gapinput@M := MunnSemigroup(e);;|
  !gapprompt@gap>| !gapinput@L := LClassNC(M, PartialPerm([51, 63], [51, 47]));;|
  !gapprompt@gap>| !gapinput@HClassReps(L);|
  [ <identity partial perm on [ 47, 51 ]>, [27,47](51), [50,47](51),
    [64,47](51), [63,47](51), [59,47](51) ]
\end{Verbatim}
 }

 

\subsection{\textcolor{Chapter }{MinimalDClass}}
\logpage{[ 10, 1, 6 ]}\nobreak
\hyperdef{L}{X81E5A04F7DA3A1E1}{}
{\noindent\textcolor{FuncColor}{$\triangleright$\enspace\texttt{MinimalDClass({\mdseries\slshape S})\index{MinimalDClass@\texttt{MinimalDClass}}
\label{MinimalDClass}
}\hfill{\scriptsize (attribute)}}\\
\textbf{\indent Returns:\ }
The minimal $\mathcal{D}$\texttt{\symbol{45}}class of a semigroup.



 The minimal ideal of a semigroup is the least ideal with respect to
containment. \texttt{MinimalDClass} returns the $\mathcal{D}$\texttt{\symbol{45}}class corresponding to the minimal ideal of the semigroup \mbox{\texttt{\mdseries\slshape S}}. Equivalently, \texttt{MinimalDClass} returns the minimal $\mathcal{D}$\texttt{\symbol{45}}class with respect to the partial order of $\mathcal{D}$\texttt{\symbol{45}}classes.

 It is significantly easier to find the minimal $\mathcal{D}$\texttt{\symbol{45}}class of a semigroup, than to find its $\mathcal{D}$\texttt{\symbol{45}}classes. 

 See also \texttt{PartialOrderOfDClasses} (\ref{PartialOrderOfDClasses}), \texttt{IsGreensLessThanOrEqual} (\textbf{Reference: IsGreensLessThanOrEqual}), \texttt{MinimalIdeal} (\ref{MinimalIdeal}) and \texttt{RepresentativeOfMinimalIdeal} (\ref{RepresentativeOfMinimalIdeal}). 
\begin{Verbatim}[commandchars=!@|,fontsize=\small,frame=single,label=Example]
  !gapprompt@gap>| !gapinput@D := MinimalDClass(JonesMonoid(8));|
  <Green's D-class: <bipartition: [ 1, 2 ], [ 3, 4 ], [ 5, 6 ],
    [ 7, 8 ], [ -1, -2 ], [ -3, -4 ], [ -5, -6 ], [ -7, -8 ]>>
  !gapprompt@gap>| !gapinput@S := InverseSemigroup(|
  !gapprompt@>| !gapinput@PartialPerm([1, 2, 3, 5, 7, 8, 9], [2, 6, 9, 1, 5, 3, 8]),|
  !gapprompt@>| !gapinput@PartialPerm([1, 3, 4, 5, 7, 8, 9], [9, 4, 10, 5, 6, 7, 1]));;|
  !gapprompt@gap>| !gapinput@MinimalDClass(S);|
  <Green's D-class: <empty partial perm>>
\end{Verbatim}
 }

 
\subsection{\textcolor{Chapter }{MaximalXClasses}}\logpage{[ 10, 1, 7 ]}
\hyperdef{L}{X834172F4787A565B}{}
{
\noindent\textcolor{FuncColor}{$\triangleright$\enspace\texttt{MaximalDClasses({\mdseries\slshape S})\index{MaximalDClasses@\texttt{MaximalDClasses}}
\label{MaximalDClasses}
}\hfill{\scriptsize (attribute)}}\\
\noindent\textcolor{FuncColor}{$\triangleright$\enspace\texttt{MaximalLClasses({\mdseries\slshape S})\index{MaximalLClasses@\texttt{MaximalLClasses}}
\label{MaximalLClasses}
}\hfill{\scriptsize (attribute)}}\\
\noindent\textcolor{FuncColor}{$\triangleright$\enspace\texttt{MaximalRClasses({\mdseries\slshape S})\index{MaximalRClasses@\texttt{MaximalRClasses}}
\label{MaximalRClasses}
}\hfill{\scriptsize (attribute)}}\\
\textbf{\indent Returns:\ }
The maximal $\mathcal{D}$, $\mathcal{L}$, or $\mathcal{R}$\texttt{\symbol{45}}classes of a semigroup.



 Let \texttt{X} be one of Green's $\mathcal{D}$\texttt{\symbol{45}}, $\mathcal{L}$\texttt{\symbol{45}}, or $\mathcal{R}$\texttt{\symbol{45}}relations. Then \texttt{MaximalXClasses} returns the maximal Green's \texttt{X}\texttt{\symbol{45}}classes with respect to the partial order of \texttt{X}\texttt{\symbol{45}}classes. 

 See also \texttt{PartialOrderOfDClasses} (\ref{PartialOrderOfDClasses}), \texttt{IsGreensLessThanOrEqual} (\textbf{Reference: IsGreensLessThanOrEqual}), and \texttt{MinimalDClass} (\ref{MinimalDClass}). 
\begin{Verbatim}[commandchars=!@|,fontsize=\small,frame=single,label=Example]
  !gapprompt@gap>| !gapinput@MaximalDClasses(BrauerMonoid(8));|
  [ <Green's D-class: <block bijection: [ 1, -1 ], [ 2, -2 ],
        [ 3, -3 ], [ 4, -4 ], [ 5, -5 ], [ 6, -6 ], [ 7, -7 ],
        [ 8, -8 ]>> ]
  !gapprompt@gap>| !gapinput@MaximalDClasses(FullTransformationMonoid(5));|
  [ <Green's D-class: IdentityTransformation> ]
  !gapprompt@gap>| !gapinput@S := Semigroup(|
  !gapprompt@>| !gapinput@PartialPerm([1, 2, 3, 4, 5, 6, 7], [3, 8, 1, 4, 5, 6, 7]),|
  !gapprompt@>| !gapinput@PartialPerm([1, 2, 3, 6, 8], [2, 6, 7, 1, 5]),|
  !gapprompt@>| !gapinput@PartialPerm([1, 2, 3, 4, 6, 8], [4, 3, 2, 7, 6, 5]),|
  !gapprompt@>| !gapinput@PartialPerm([1, 2, 4, 5, 6, 7, 8], [7, 1, 4, 2, 5, 6, 3]));;|
  !gapprompt@gap>| !gapinput@MaximalDClasses(S);|
  [ <Green's D-class: [2,8](1,3)(4)(5)(6)(7)>,
    <Green's D-class: [8,3](1,7,6,5,2)(4)> ]
\end{Verbatim}
 }

 

\subsection{\textcolor{Chapter }{NrRegularDClasses}}
\logpage{[ 10, 1, 8 ]}\nobreak
\hyperdef{L}{X7AA3F0A77D0043FB}{}
{\noindent\textcolor{FuncColor}{$\triangleright$\enspace\texttt{NrRegularDClasses({\mdseries\slshape S})\index{NrRegularDClasses@\texttt{NrRegularDClasses}}
\label{NrRegularDClasses}
}\hfill{\scriptsize (attribute)}}\\
\noindent\textcolor{FuncColor}{$\triangleright$\enspace\texttt{RegularDClasses({\mdseries\slshape S})\index{RegularDClasses@\texttt{RegularDClasses}}
\label{RegularDClasses}
}\hfill{\scriptsize (attribute)}}\\
\textbf{\indent Returns:\ }
 A positive integer, or a list. 



 \texttt{NrRegularDClasses} returns the number of regular $\mathcal{D}$\texttt{\symbol{45}}classes of the semigroup \mbox{\texttt{\mdseries\slshape S}}.

 \texttt{RegularDClasses} returns a list of the regular $\mathcal{D}$\texttt{\symbol{45}}classes of the semigroup \mbox{\texttt{\mdseries\slshape S}}. 

 See also \texttt{IsRegularGreensClass} (\ref{IsRegularGreensClass}) and \texttt{IsRegularDClass} (\textbf{Reference: IsRegularDClass}). 
\begin{Verbatim}[commandchars=!@|,fontsize=\small,frame=single,label=Example]
  !gapprompt@gap>| !gapinput@S := Semigroup(Transformation([1, 3, 4, 1, 3, 5]),|
  !gapprompt@>| !gapinput@                  Transformation([5, 1, 6, 1, 6, 3]));;|
  !gapprompt@gap>| !gapinput@NrRegularDClasses(S);|
  3
  !gapprompt@gap>| !gapinput@NrDClasses(S);|
  7
  !gapprompt@gap>| !gapinput@AsSet(RegularDClasses(S));|
  [ <Green's D-class: Transformation( [ 1, 4, 1, 1, 4, 3 ] )>,
    <Green's D-class: Transformation( [ 1, 1, 1, 1, 1 ] )>,
    <Green's D-class: Transformation( [ 1, 1, 1, 1, 1, 1 ] )> ]
\end{Verbatim}
 }

 
\subsection{\textcolor{Chapter }{NrXClasses}}\logpage{[ 10, 1, 9 ]}
\hyperdef{L}{X7E45FD9F7BADDFBD}{}
{
\noindent\textcolor{FuncColor}{$\triangleright$\enspace\texttt{NrDClasses({\mdseries\slshape obj})\index{NrDClasses@\texttt{NrDClasses}}
\label{NrDClasses}
}\hfill{\scriptsize (attribute)}}\\
\noindent\textcolor{FuncColor}{$\triangleright$\enspace\texttt{NrHClasses({\mdseries\slshape obj})\index{NrHClasses@\texttt{NrHClasses}}
\label{NrHClasses}
}\hfill{\scriptsize (attribute)}}\\
\noindent\textcolor{FuncColor}{$\triangleright$\enspace\texttt{NrLClasses({\mdseries\slshape obj})\index{NrLClasses@\texttt{NrLClasses}}
\label{NrLClasses}
}\hfill{\scriptsize (attribute)}}\\
\noindent\textcolor{FuncColor}{$\triangleright$\enspace\texttt{NrRClasses({\mdseries\slshape obj})\index{NrRClasses@\texttt{NrRClasses}}
\label{NrRClasses}
}\hfill{\scriptsize (attribute)}}\\
\textbf{\indent Returns:\ }
 A positive integer. 



 \texttt{NrXClasses} returns the number of Green's classes in \mbox{\texttt{\mdseries\slshape obj}} where \mbox{\texttt{\mdseries\slshape obj}} can be a semigroup, $\mathcal{D}$\texttt{\symbol{45}}, $\mathcal{L}$\texttt{\symbol{45}}, or $\mathcal{R}$\texttt{\symbol{45}}class where appropriate. If the actual Green's classes are
not required, then it is more efficient to use 
\begin{Verbatim}[commandchars=!@|,fontsize=\small,frame=single,label=Example]
  NrHClasses(obj)
\end{Verbatim}
 than 
\begin{Verbatim}[commandchars=!@|,fontsize=\small,frame=single,label=Example]
  Length(HClasses(obj))
\end{Verbatim}
 since the Green's classes themselves are not created when \texttt{NrXClasses} is called. 

 See also \texttt{GreensRClasses} (\ref{GreensRClasses}), \texttt{GreensRClasses} (\textbf{Reference: GreensRClasses}), \texttt{IteratorOfRClasses} (\ref{IteratorOfRClasses}), and 
\begin{Verbatim}[commandchars=!@|,fontsize=\small,frame=single,label=Example]
  !gapprompt@gap>| !gapinput@S := Semigroup(|
  !gapprompt@>| !gapinput@ Transformation([1, 2, 5, 4, 3, 8, 7, 6]),|
  !gapprompt@>| !gapinput@ Transformation([1, 6, 3, 4, 7, 2, 5, 8]),|
  !gapprompt@>| !gapinput@ Transformation([2, 1, 6, 7, 8, 3, 4, 5]),|
  !gapprompt@>| !gapinput@ Transformation([3, 2, 3, 6, 1, 6, 1, 2]),|
  !gapprompt@>| !gapinput@ Transformation([5, 2, 3, 6, 3, 4, 7, 4]));;|
  !gapprompt@gap>| !gapinput@x := Transformation([2, 5, 4, 7, 4, 3, 6, 3]);;|
  !gapprompt@gap>| !gapinput@R := RClass(S, x);|
  <Green's R-class: Transformation( [ 2, 5, 4, 7, 4, 3, 6, 3 ] )>
  !gapprompt@gap>| !gapinput@NrHClasses(R);|
  12
  !gapprompt@gap>| !gapinput@D := DClass(R);|
  <Green's D-class: Transformation( [ 2, 5, 4, 7, 4, 3, 6, 3 ] )>
  !gapprompt@gap>| !gapinput@NrHClasses(D);|
  72
  !gapprompt@gap>| !gapinput@L := LClass(S, x);|
  <Green's L-class: Transformation( [ 2, 5, 4, 7, 4, 3, 6, 3 ] )>
  !gapprompt@gap>| !gapinput@NrHClasses(L);|
  6
  !gapprompt@gap>| !gapinput@NrHClasses(S);|
  1555
  !gapprompt@gap>| !gapinput@S := Semigroup(Transformation([4, 6, 5, 2, 1, 3]),|
  !gapprompt@>| !gapinput@                  Transformation([6, 3, 2, 5, 4, 1]),|
  !gapprompt@>| !gapinput@                  Transformation([1, 2, 4, 3, 5, 6]),|
  !gapprompt@>| !gapinput@                  Transformation([3, 5, 6, 1, 2, 3]),|
  !gapprompt@>| !gapinput@                  Transformation([5, 3, 6, 6, 6, 2]),|
  !gapprompt@>| !gapinput@                  Transformation([2, 3, 2, 6, 4, 6]),|
  !gapprompt@>| !gapinput@                  Transformation([2, 1, 2, 2, 2, 4]),|
  !gapprompt@>| !gapinput@                  Transformation([4, 4, 1, 2, 1, 2]));;|
  !gapprompt@gap>| !gapinput@NrRClasses(S);|
  150
  !gapprompt@gap>| !gapinput@Size(S);|
  6342
  !gapprompt@gap>| !gapinput@x := Transformation([1, 3, 3, 1, 3, 5]);;|
  !gapprompt@gap>| !gapinput@D := DClass(S, x);|
  <Green's D-class: Transformation( [ 1, 3, 3, 1, 3, 5 ] )>
  !gapprompt@gap>| !gapinput@NrRClasses(D);|
  87
  !gapprompt@gap>| !gapinput@S := SymmetricInverseSemigroup(10);;|
  !gapprompt@gap>| !gapinput@NrDClasses(S); NrRClasses(S); NrHClasses(S); NrLClasses(S);|
  11
  1024
  184756
  1024
  !gapprompt@gap>| !gapinput@S := POPI(10);;|
  !gapprompt@gap>| !gapinput@NrDClasses(S);|
  11
  !gapprompt@gap>| !gapinput@NrRClasses(S);|
  1024
\end{Verbatim}
 }

 
\subsection{\textcolor{Chapter }{PartialOrderOfXClasses}}\logpage{[ 10, 1, 10 ]}
\hyperdef{L}{X8140814084748101}{}
{
\noindent\textcolor{FuncColor}{$\triangleright$\enspace\texttt{PartialOrderOfDClasses({\mdseries\slshape S})\index{PartialOrderOfDClasses@\texttt{PartialOrderOfDClasses}}
\label{PartialOrderOfDClasses}
}\hfill{\scriptsize (attribute)}}\\
\noindent\textcolor{FuncColor}{$\triangleright$\enspace\texttt{PartialOrderOfLClasses({\mdseries\slshape S})\index{PartialOrderOfLClasses@\texttt{PartialOrderOfLClasses}}
\label{PartialOrderOfLClasses}
}\hfill{\scriptsize (attribute)}}\\
\noindent\textcolor{FuncColor}{$\triangleright$\enspace\texttt{PartialOrderOfRClasses({\mdseries\slshape S})\index{PartialOrderOfRClasses@\texttt{PartialOrderOfRClasses}}
\label{PartialOrderOfRClasses}
}\hfill{\scriptsize (attribute)}}\\
\textbf{\indent Returns:\ }
 A digraph. 



 Let \texttt{X} be one of Green's $\mathcal{D}$\texttt{\symbol{45}}, $\mathcal{L}$\texttt{\symbol{45}}, or $\mathcal{R}$\texttt{\symbol{45}}relations. Then \texttt{PartialOrderOfXClasses} returns a digraph \texttt{D} where \texttt{OutNeighbours(D)[i]} contains every \texttt{j} such that \texttt{GreensXClasses(S)[j]} is immediately less than \texttt{GreensXClasses(S)[i]} in the partial order of \texttt{X}\texttt{\symbol{45}}classes of \mbox{\texttt{\mdseries\slshape S}}. The reflexive transitive closure of the digraph \texttt{D} is the partial order of \texttt{X}\texttt{\symbol{45}}classes of \mbox{\texttt{\mdseries\slshape S}} (in the sense of the \textsf{digraphs} package). 

 The partial order on the \texttt{X}\texttt{\symbol{45}}classes is defined as follows. 
\begin{description}
\item[{Green's $\mathcal{D}$\texttt{\symbol{45}}relation:}]  $x\leq y$ if and only if $S ^ 1xS ^ 1$ is a subset of $S ^ 1yS ^ 1$. 
\item[{Green's $\mathcal{L}$\texttt{\symbol{45}}relation:}]  $x\leq y$ if and only if $S ^ 1x$ is a subset of $S ^ 1y$. 
\item[{Green's $\mathcal{R}$\texttt{\symbol{45}}relation:}]  $x\leq y$ if and only if $xS ^ 1$ is a subset of $yS ^ 1$. 
\end{description}
 See also \texttt{GreensDClasses} (\ref{GreensDClasses}), \texttt{GreensDClasses} (\textbf{Reference: GreensDClasses}), \texttt{IsGreensLessThanOrEqual} (\textbf{Reference: IsGreensLessThanOrEqual}), and \texttt{\texttt{\symbol{92}}{\textless}} (\ref{bSlash<:for Green's classes}). 
\begin{Verbatim}[commandchars=!@|,fontsize=\small,frame=single,label=Example]
  !gapprompt@gap>| !gapinput@S := Semigroup(Transformation([2, 4, 1, 2]),|
  !gapprompt@>| !gapinput@                  Transformation([3, 3, 4, 1]));;|
  !gapprompt@gap>| !gapinput@PartialOrderOfDClasses(S);|
  <immutable digraph with 4 vertices, 3 edges>
  !gapprompt@gap>| !gapinput@IsGreensLessThanOrEqual(GreensDClasses(S)[1],|
  !gapprompt@>| !gapinput@                           GreensDClasses(S)[2]);|
  false
  !gapprompt@gap>| !gapinput@IsGreensLessThanOrEqual(GreensDClasses(S)[2],|
  !gapprompt@>| !gapinput@                           GreensDClasses(S)[1]);|
  false
  !gapprompt@gap>| !gapinput@IsGreensLessThanOrEqual(GreensDClasses(S)[3],|
  !gapprompt@>| !gapinput@                           GreensDClasses(S)[1]);|
  true
  !gapprompt@gap>| !gapinput@S := InverseSemigroup(|
  !gapprompt@>| !gapinput@PartialPerm([1, 2, 3], [1, 3, 4]),|
  !gapprompt@>| !gapinput@PartialPerm([1, 3, 5], [5, 1, 3]));;|
  !gapprompt@gap>| !gapinput@Size(S);|
  58
  !gapprompt@gap>| !gapinput@PartialOrderOfDClasses(S);|
  <immutable digraph with 5 vertices, 4 edges>
  !gapprompt@gap>| !gapinput@IsGreensLessThanOrEqual(GreensDClasses(S)[1],|
  !gapprompt@>| !gapinput@                           GreensDClasses(S)[2]);|
  false
  !gapprompt@gap>| !gapinput@IsGreensLessThanOrEqual(GreensDClasses(S)[5],|
  !gapprompt@>| !gapinput@                           GreensDClasses(S)[2]);|
  true
  !gapprompt@gap>| !gapinput@IsGreensLessThanOrEqual(GreensDClasses(S)[3],|
  !gapprompt@>| !gapinput@                           GreensDClasses(S)[4]);|
  false
  !gapprompt@gap>| !gapinput@IsGreensLessThanOrEqual(GreensDClasses(S)[4],|
  !gapprompt@>| !gapinput@                           GreensDClasses(S)[3]);|
  true
\end{Verbatim}
 }

 

\subsection{\textcolor{Chapter }{LengthOfLongestDClassChain}}
\logpage{[ 10, 1, 11 ]}\nobreak
\hyperdef{L}{X83B0EDA57F1D2F97}{}
{\noindent\textcolor{FuncColor}{$\triangleright$\enspace\texttt{LengthOfLongestDClassChain({\mdseries\slshape S})\index{LengthOfLongestDClassChain@\texttt{LengthOfLongestDClassChain}}
\label{LengthOfLongestDClassChain}
}\hfill{\scriptsize (attribute)}}\\
\textbf{\indent Returns:\ }
 A non\texttt{\symbol{45}}negative integer. 



 If \mbox{\texttt{\mdseries\slshape S}} is a semigroup, then \texttt{LengthOfLongestDClassChain} returns the length of the longest chain in the partial order defined by \texttt{PartialOrderOfDClasses(\mbox{\texttt{\mdseries\slshape S}})}. See \texttt{PartialOrderOfDClasses} (\ref{PartialOrderOfDClasses}). 

 The partial order on the $\mathcal{D}$\texttt{\symbol{45}}classes is defined by $x\leq y$ if and only if $S ^ 1xS ^ 1$ is a subset of $S ^ 1yS ^ 1$. A \emph{chain} of $\mathcal{D}$\texttt{\symbol{45}}classes is a collection of \texttt{n} $\mathcal{D}$\texttt{\symbol{45}}classes $D_{1}, D_{2}, \ldots D_{n}$ such that $D_{1} < D_{2} < \cdots < D_{n}$. The \emph{length} of such a chain is \texttt{n \texttt{\symbol{45}} 1}. 
\begin{Verbatim}[commandchars=!@|,fontsize=\small,frame=single,label=Example]
  !gapprompt@gap>| !gapinput@S := TrivialSemigroup();;|
  !gapprompt@gap>| !gapinput@LengthOfLongestDClassChain(S);|
  0
  !gapprompt@gap>| !gapinput@T := ZeroSemigroup(5);;|
  !gapprompt@gap>| !gapinput@LengthOfLongestDClassChain(T);|
  1
  !gapprompt@gap>| !gapinput@U := MonogenicSemigroup(14, 7);;|
  !gapprompt@gap>| !gapinput@LengthOfLongestDClassChain(U);|
  13
  !gapprompt@gap>| !gapinput@V := FullTransformationMonoid(6);|
  <full transformation monoid of degree 6>
  !gapprompt@gap>| !gapinput@LengthOfLongestDClassChain(V);|
  5
\end{Verbatim}
 }

 

\subsection{\textcolor{Chapter }{IsGreensDGreaterThanFunc}}
\logpage{[ 10, 1, 12 ]}\nobreak
\hyperdef{L}{X7E872C5381D0DD8A}{}
{\noindent\textcolor{FuncColor}{$\triangleright$\enspace\texttt{IsGreensDGreaterThanFunc({\mdseries\slshape S})\index{IsGreensDGreaterThanFunc@\texttt{IsGreensDGreaterThanFunc}}
\label{IsGreensDGreaterThanFunc}
}\hfill{\scriptsize (attribute)}}\\
\textbf{\indent Returns:\ }
A function.



 \texttt{IsGreensDGreaterThanFunc(\mbox{\texttt{\mdseries\slshape S}})} returns a function \texttt{func} such that for any two elements \texttt{x} and \texttt{y} of \mbox{\texttt{\mdseries\slshape S}}, \texttt{func(x, y)} return \texttt{true} if the $\mathcal{D}$\texttt{\symbol{45}}class of \texttt{x} in \mbox{\texttt{\mdseries\slshape S}} is greater than or equal to the $\mathcal{D}$\texttt{\symbol{45}}class of \texttt{y} in \mbox{\texttt{\mdseries\slshape S}} under the usual ordering of Green's $\mathcal{D}$\texttt{\symbol{45}}classes of a semigroup. 
\begin{Verbatim}[commandchars=!@|,fontsize=\small,frame=single,label=Example]
  !gapprompt@gap>| !gapinput@S := Semigroup(Transformation([1, 3, 4, 1, 3]),|
  !gapprompt@>| !gapinput@                  Transformation([2, 4, 1, 5, 5]),|
  !gapprompt@>| !gapinput@                  Transformation([2, 5, 3, 5, 3]),|
  !gapprompt@>| !gapinput@                  Transformation([5, 5, 1, 1, 3]));;|
  !gapprompt@gap>| !gapinput@reps := ShallowCopy(AsSet(DClassReps(S)));|
  [ Transformation( [ 1, 1, 1, 1, 1 ] ),
    Transformation( [ 1, 3, 1, 3, 3 ] ),
    Transformation( [ 1, 3, 4, 1, 3 ] ),
    Transformation( [ 2, 4, 1, 5, 5 ] ) ]
  !gapprompt@gap>| !gapinput@Sort(reps, IsGreensDGreaterThanFunc(S));|
  !gapprompt@gap>| !gapinput@reps;|
  [ Transformation( [ 2, 4, 1, 5, 5 ] ),
    Transformation( [ 1, 3, 4, 1, 3 ] ),
    Transformation( [ 1, 3, 1, 3, 3 ] ),
    Transformation( [ 1, 1, 1, 1, 1 ] ) ]
  !gapprompt@gap>| !gapinput@IsGreensLessThanOrEqual(DClass(S, reps[2]),|
  !gapprompt@>| !gapinput@                           DClass(S, reps[1]));|
  true
  !gapprompt@gap>| !gapinput@S := DualSymmetricInverseMonoid(4);;|
  !gapprompt@gap>| !gapinput@IsGreensDGreaterThanFunc(S)(S.1, S.3);|
  true
  !gapprompt@gap>| !gapinput@IsGreensDGreaterThanFunc(S)(S.3, S.1);|
  false
  !gapprompt@gap>| !gapinput@IsGreensLessThanOrEqual(DClass(S, S.3),|
  !gapprompt@>| !gapinput@                           DClass(S, S.1));|
  true
  !gapprompt@gap>| !gapinput@IsGreensLessThanOrEqual(DClass(S, S.1),|
  !gapprompt@>| !gapinput@                           DClass(S, S.3));|
  false
\end{Verbatim}
 }

 }

   
\section{\textcolor{Chapter }{ Iterators and enumerators of classes and representatives }}\logpage{[ 10, 2, 0 ]}
\hyperdef{L}{X819CCBD67FD27115}{}
{
  In this section, we describe the methods in the \textsf{Semigroups} package for incrementally determining Green's classes or their
representatives. 

 
\subsection{\textcolor{Chapter }{IteratorOfXClassReps}}\logpage{[ 10, 2, 1 ]}
\hyperdef{L}{X8566F84A7F6D4193}{}
{
\noindent\textcolor{FuncColor}{$\triangleright$\enspace\texttt{IteratorOfDClassReps({\mdseries\slshape S})\index{IteratorOfDClassReps@\texttt{IteratorOfDClassReps}}
\label{IteratorOfDClassReps}
}\hfill{\scriptsize (operation)}}\\
\noindent\textcolor{FuncColor}{$\triangleright$\enspace\texttt{IteratorOfHClassReps({\mdseries\slshape S})\index{IteratorOfHClassReps@\texttt{IteratorOfHClassReps}}
\label{IteratorOfHClassReps}
}\hfill{\scriptsize (operation)}}\\
\textbf{\indent Returns:\ }
 An iterator. 



 Returns an iterator of the representatives of the Green's classes contained in
the semigroup \mbox{\texttt{\mdseries\slshape S}}. See  (\textbf{Reference: Iterators}) for more information on iterators.

 See also \texttt{GreensRClasses} (\textbf{Reference: GreensRClasses}), \texttt{GreensRClasses} (\ref{GreensRClasses}), and \texttt{IteratorOfRClasses} (\ref{IteratorOfRClasses}).

 }

 
\subsection{\textcolor{Chapter }{IteratorOfXClasses}}\logpage{[ 10, 2, 2 ]}
\hyperdef{L}{X867D7B8982915960}{}
{
\noindent\textcolor{FuncColor}{$\triangleright$\enspace\texttt{IteratorOfDClasses({\mdseries\slshape S})\index{IteratorOfDClasses@\texttt{IteratorOfDClasses}}
\label{IteratorOfDClasses}
}\hfill{\scriptsize (operation)}}\\
\noindent\textcolor{FuncColor}{$\triangleright$\enspace\texttt{IteratorOfRClasses({\mdseries\slshape S})\index{IteratorOfRClasses@\texttt{IteratorOfRClasses}}
\label{IteratorOfRClasses}
}\hfill{\scriptsize (operation)}}\\
\textbf{\indent Returns:\ }
 An iterator. 



 Returns an iterator of the Green's classes in the semigroup \mbox{\texttt{\mdseries\slshape S}}. See  (\textbf{Reference: Iterators}) for more information on iterators.

 This function is useful if you are, for example, looking for an $\mathcal{R}$\texttt{\symbol{45}}class of a semigroup with a particular property but do not
necessarily want to compute all of the $\mathcal{R}$\texttt{\symbol{45}}classes.

 See also \texttt{GreensRClasses} (\ref{GreensRClasses}), \texttt{GreensRClasses} (\textbf{Reference: GreensRClasses}), and \texttt{NrRClasses} (\ref{NrRClasses}). 

 The transformation semigroup in the example below has 25147892 elements but it
only takes a fraction of a second to find a non\texttt{\symbol{45}}trivial $\mathcal{R}$\texttt{\symbol{45}}class. The inverse semigroup of partial permutations in
the example below has size 158122047816 but it only takes a fraction of a
second to find an $\mathcal{R}$\texttt{\symbol{45}}class with more than 1000 elements. 
\begin{Verbatim}[commandchars=!@|,fontsize=\small,frame=single,label=Example]
  !gapprompt@gap>| !gapinput@gens := [Transformation([2, 4, 1, 5, 4, 4, 7, 3, 8, 1]),|
  !gapprompt@>| !gapinput@            Transformation([3, 2, 8, 8, 4, 4, 8, 6, 5, 7]),|
  !gapprompt@>| !gapinput@            Transformation([4, 10, 6, 6, 1, 2, 4, 10, 9, 7]),|
  !gapprompt@>| !gapinput@            Transformation([6, 2, 2, 4, 9, 9, 5, 10, 1, 8]),|
  !gapprompt@>| !gapinput@            Transformation([6, 4, 1, 6, 6, 8, 9, 6, 2, 2]),|
  !gapprompt@>| !gapinput@            Transformation([6, 8, 1, 10, 6, 4, 9, 1, 9, 4]),|
  !gapprompt@>| !gapinput@            Transformation([8, 6, 2, 3, 3, 4, 8, 6, 2, 9]),|
  !gapprompt@>| !gapinput@            Transformation([9, 1, 2, 8, 1, 5, 9, 9, 9, 5]),|
  !gapprompt@>| !gapinput@            Transformation([9, 3, 1, 5, 10, 3, 4, 6, 10, 2]),|
  !gapprompt@>| !gapinput@            Transformation([10, 7, 3, 7, 1, 9, 8, 8, 4, 10])];;|
  !gapprompt@gap>| !gapinput@S := Semigroup(gens);;|
  !gapprompt@gap>| !gapinput@iter := IteratorOfRClasses(S);|
  <iterator>
  !gapprompt@gap>| !gapinput@for R in iter do|
  !gapprompt@>| !gapinput@  if Size(R) > 1 then|
  !gapprompt@>| !gapinput@    break;|
  !gapprompt@>| !gapinput@  fi;|
  !gapprompt@>| !gapinput@od;|
  !gapprompt@gap>| !gapinput@R;|
  <Green's R-class: Transformation( [ 6, 4, 1, 6, 6, 8, 9, 6, 2, 2 ] )>
  !gapprompt@gap>| !gapinput@Size(R);|
  21600
  !gapprompt@gap>| !gapinput@S := InverseSemigroup(|
  !gapprompt@>| !gapinput@ PartialPerm([1, 2, 3, 4, 5, 6, 7, 10, 11, 19, 20],|
  !gapprompt@>| !gapinput@             [19, 4, 11, 15, 3, 20, 1, 14, 8, 13, 17]),|
  !gapprompt@>| !gapinput@ PartialPerm([1, 2, 3, 4, 6, 7, 8, 14, 15, 16, 17],|
  !gapprompt@>| !gapinput@             [15, 14, 20, 19, 4, 5, 1, 13, 11, 10, 3]),|
  !gapprompt@>| !gapinput@ PartialPerm([1, 2, 4, 6, 7, 8, 9, 10, 14, 15, 18],|
  !gapprompt@>| !gapinput@             [7, 2, 17, 10, 1, 19, 9, 3, 11, 16, 18]),|
  !gapprompt@>| !gapinput@ PartialPerm([1, 2, 3, 4, 5, 7, 8, 9, 11, 12, 13, 16],|
  !gapprompt@>| !gapinput@             [8, 3, 18, 1, 4, 13, 12, 7, 19, 20, 2, 11]),|
  !gapprompt@>| !gapinput@ PartialPerm([1, 2, 3, 4, 5, 6, 7, 9, 11, 15, 16, 17, 20],|
  !gapprompt@>| !gapinput@             [7, 17, 13, 4, 6, 9, 18, 10, 11, 19, 5, 2, 8]),|
  !gapprompt@>| !gapinput@ PartialPerm([1, 3, 4, 5, 6, 7, 8, 9, 10, 11, 12, 15, 18],|
  !gapprompt@>| !gapinput@             [10, 20, 11, 7, 13, 8, 4, 9, 2, 18, 17, 6, 15]),|
  !gapprompt@>| !gapinput@ PartialPerm([1, 2, 3, 4, 5, 6, 7, 8, 9, 11, 13, 14, 17, 18],|
  !gapprompt@>| !gapinput@             [10, 20, 18, 1, 14, 16, 9, 5, 15, 4, 8, 12, 19, 11]),|
  !gapprompt@>| !gapinput@ PartialPerm([1, 2, 3, 4, 5, 6, 7, 9, 10, 11, 12, 15, 16, 19, 20],|
  !gapprompt@>| !gapinput@             [13, 6, 1, 2, 11, 7, 16, 18, 9, 10, 4, 14, 15, 5, 17]),|
  !gapprompt@>| !gapinput@ PartialPerm([1, 2, 3, 4, 6, 7, 8, 9, 10, 11, 12, 14, 15, 16, 20],|
  !gapprompt@>| !gapinput@             [5, 3, 12, 9, 20, 15, 8, 16, 13, 1, 17, 11, 14, 10, 2]),|
  !gapprompt@>| !gapinput@ PartialPerm([1, 2, 3, 4, 6, 7, 8, 9, 10, 11, 13, 17, 18, 19, 20],|
  !gapprompt@>| !gapinput@             [8, 3, 9, 20, 2, 12, 14, 15, 4, 18, 13, 1, 17, 19, 5]));;|
  !gapprompt@gap>| !gapinput@iter := IteratorOfRClasses(S);|
  <iterator>
  !gapprompt@gap>| !gapinput@repeat|
  !gapprompt@>| !gapinput@  R := NextIterator(iter);|
  !gapprompt@>| !gapinput@until Size(R) > 1000;|
  !gapprompt@gap>| !gapinput@R;|
  <Green's R-class: [8,19,14][11,4][13,15,5][17,20]>
  !gapprompt@gap>| !gapinput@Size(R);|
  10020240
\end{Verbatim}
 }

 }

   
\section{\textcolor{Chapter }{ Properties of Green's classes }}\logpage{[ 10, 3, 0 ]}
\hyperdef{L}{X820EF2BA7D5D53B4}{}
{
  In this section, we describe the properties and operators of Green's classes
that are available in the \textsf{Semigroups} package. 

 
\subsection{\textcolor{Chapter }{Less than for Green's classes}}\logpage{[ 10, 3, 1 ]}
\hyperdef{L}{X85F30ACF86C3A733}{}
{
\noindent\textcolor{FuncColor}{$\triangleright$\enspace\texttt{\texttt{\symbol{92}}{\textless}({\mdseries\slshape left\texttt{\symbol{45}}expr, right\texttt{\symbol{45}}expr})\index{\texttt{\symbol{92}}{\textless}@\texttt{\texttt{\symbol{92}}{\textless}}!for Green's classes}
\label{bSlash<:for Green's classes}
}\hfill{\scriptsize (method)}}\\
\textbf{\indent Returns:\ }
\texttt{true} or \texttt{false}.



 The Green's class \mbox{\texttt{\mdseries\slshape left\texttt{\symbol{45}}expr}} is less than or equal to \mbox{\texttt{\mdseries\slshape right\texttt{\symbol{45}}expr}} if they belong to the same semigroup and the representative of \mbox{\texttt{\mdseries\slshape left\texttt{\symbol{45}}expr}} is less than the representative of \mbox{\texttt{\mdseries\slshape right\texttt{\symbol{45}}expr}} under \texttt{{\textless}}; see also \texttt{Representative} (\textbf{Reference: Representative}). 

 Please note that this is not the usual order on the Green's classes of a
semigroup as defined in  (\textbf{Reference: Green's Relations}). See also \texttt{IsGreensLessThanOrEqual} (\textbf{Reference: IsGreensLessThanOrEqual}). 
\begin{Verbatim}[commandchars=!@|,fontsize=\small,frame=single,label=Example]
  !gapprompt@gap>| !gapinput@S := FullTransformationSemigroup(4);;|
  !gapprompt@gap>| !gapinput@A := GreensRClassOfElement(S, Transformation([2, 1, 3, 1]));|
  <Green's R-class: Transformation( [ 2, 1, 3, 1 ] )>
  !gapprompt@gap>| !gapinput@B := GreensRClassOfElement(S, Transformation([1, 2, 3, 4]));|
  <Green's R-class: IdentityTransformation>
  !gapprompt@gap>| !gapinput@A < B;|
  false
  !gapprompt@gap>| !gapinput@B < A;|
  true
  !gapprompt@gap>| !gapinput@IsGreensLessThanOrEqual(A, B);|
  true
  !gapprompt@gap>| !gapinput@IsGreensLessThanOrEqual(B, A);|
  false
  !gapprompt@gap>| !gapinput@S := SymmetricInverseSemigroup(4);;|
  !gapprompt@gap>| !gapinput@A := GreensJClassOfElement(S, PartialPerm([1, 3, 4]));;|
  !gapprompt@gap>| !gapinput@B := GreensJClassOfElement(S, PartialPerm([3, 1]));;|
  !gapprompt@gap>| !gapinput@A < B;|
  true
  !gapprompt@gap>| !gapinput@B < A;|
  false
  !gapprompt@gap>| !gapinput@IsGreensLessThanOrEqual(A, B);|
  false
  !gapprompt@gap>| !gapinput@IsGreensLessThanOrEqual(B, A);|
  true
\end{Verbatim}
 }

 

\subsection{\textcolor{Chapter }{IsRegularGreensClass}}
\logpage{[ 10, 3, 2 ]}\nobreak
\hyperdef{L}{X859DD1C079C80DCC}{}
{\noindent\textcolor{FuncColor}{$\triangleright$\enspace\texttt{IsRegularGreensClass({\mdseries\slshape class})\index{IsRegularGreensClass@\texttt{IsRegularGreensClass}}
\label{IsRegularGreensClass}
}\hfill{\scriptsize (property)}}\\
\textbf{\indent Returns:\ }
 \texttt{true} or \texttt{false}. 



 This function returns \texttt{true} if \mbox{\texttt{\mdseries\slshape class}} is a regular Green's class and \texttt{false} if it is not. See also \texttt{IsRegularDClass} (\textbf{Reference: IsRegularDClass}), \texttt{IsGroupHClass} (\textbf{Reference: IsGroupHClass}), \texttt{GroupHClassOfGreensDClass} (\textbf{Reference: GroupHClassOfGreensDClass}), \texttt{GroupHClass} (\ref{GroupHClass}), \texttt{NrIdempotents} (\ref{NrIdempotents}), \texttt{Idempotents} (\ref{Idempotents}), and \texttt{IsRegularSemigroupElement} (\textbf{Reference: IsRegularSemigroupElement}). 

 The function \texttt{IsRegularDClass} produces the same output as the \textsf{GAP} library functions with the same name; see \texttt{IsRegularDClass} (\textbf{Reference: IsRegularDClass}). 
\begin{Verbatim}[commandchars=!@|,fontsize=\small,frame=single,label=Example]
  !gapprompt@gap>| !gapinput@S := Monoid(Transformation([10, 8, 7, 4, 1, 4, 10, 10, 7, 2]),|
  !gapprompt@>| !gapinput@               Transformation([5, 2, 5, 5, 9, 10, 8, 3, 8, 10]));;|
  !gapprompt@gap>| !gapinput@f := Transformation([1, 1, 10, 8, 8, 8, 1, 1, 10, 8]);;|
  !gapprompt@gap>| !gapinput@R := RClass(S, f);;|
  !gapprompt@gap>| !gapinput@IsRegularGreensClass(R);|
  true
  !gapprompt@gap>| !gapinput@S := Monoid(Transformation([2, 3, 4, 5, 1, 8, 7, 6, 2, 7]),|
  !gapprompt@>| !gapinput@               Transformation([3, 8, 7, 4, 1, 4, 3, 3, 7, 2]));;|
  !gapprompt@gap>| !gapinput@f := Transformation([3, 8, 7, 4, 1, 4, 3, 3, 7, 2]);;|
  !gapprompt@gap>| !gapinput@R := RClass(S, f);;|
  !gapprompt@gap>| !gapinput@IsRegularGreensClass(R);|
  false
  !gapprompt@gap>| !gapinput@NrIdempotents(R);|
  0
  !gapprompt@gap>| !gapinput@S := Semigroup(Transformation([2, 1, 3, 1]),|
  !gapprompt@>| !gapinput@                  Transformation([3, 1, 2, 1]),|
  !gapprompt@>| !gapinput@                  Transformation([4, 2, 3, 3]));;|
  !gapprompt@gap>| !gapinput@f := Transformation([4, 2, 3, 3]);;|
  !gapprompt@gap>| !gapinput@L := GreensLClassOfElement(S, f);;|
  !gapprompt@gap>| !gapinput@IsRegularGreensClass(L);|
  false
  !gapprompt@gap>| !gapinput@R := GreensRClassOfElement(S, f);;|
  !gapprompt@gap>| !gapinput@IsRegularGreensClass(R);|
  false
  !gapprompt@gap>| !gapinput@g := Transformation([4, 4, 4, 4]);;|
  !gapprompt@gap>| !gapinput@IsRegularSemigroupElement(S, g);|
  true
  !gapprompt@gap>| !gapinput@IsRegularGreensClass(LClass(S, g));|
  true
  !gapprompt@gap>| !gapinput@IsRegularGreensClass(RClass(S, g));|
  true
  !gapprompt@gap>| !gapinput@IsRegularDClass(DClass(S, g));|
  true
  !gapprompt@gap>| !gapinput@DClass(S, g) = RClass(S, g);|
  false
\end{Verbatim}
 }

 

\subsection{\textcolor{Chapter }{IsGreensClassNC}}
\logpage{[ 10, 3, 3 ]}\nobreak
\hyperdef{L}{X7E9BD34B8021045A}{}
{\noindent\textcolor{FuncColor}{$\triangleright$\enspace\texttt{IsGreensClassNC({\mdseries\slshape class})\index{IsGreensClassNC@\texttt{IsGreensClassNC}}
\label{IsGreensClassNC}
}\hfill{\scriptsize (property)}}\\
\textbf{\indent Returns:\ }
\texttt{true} or \texttt{false}.



 A Green's class \mbox{\texttt{\mdseries\slshape class}} of a semigroup \texttt{S} satisfies \texttt{IsGreensClassNC} if it was not known to \textsf{GAP} that the representative of \mbox{\texttt{\mdseries\slshape class}} was an element of \texttt{S} at the point that \mbox{\texttt{\mdseries\slshape class}} was created. }

 }

   
\section{\textcolor{Chapter }{ Attributes of Green's classes }}\logpage{[ 10, 4, 0 ]}
\hyperdef{L}{X855723B17D4AAF8F}{}
{
  In this section, we describe the attributes of Green's classes that are
available in the \textsf{Semigroups} package. 

 

\subsection{\textcolor{Chapter }{GroupHClass}}
\logpage{[ 10, 4, 1 ]}\nobreak
\hyperdef{L}{X8723756387DD4C0F}{}
{\noindent\textcolor{FuncColor}{$\triangleright$\enspace\texttt{GroupHClass({\mdseries\slshape class})\index{GroupHClass@\texttt{GroupHClass}}
\label{GroupHClass}
}\hfill{\scriptsize (attribute)}}\\
\textbf{\indent Returns:\ }
 A group $\mathcal{H}$\texttt{\symbol{45}}class of the $\mathcal{D}$\texttt{\symbol{45}}class \mbox{\texttt{\mdseries\slshape class}} if it is regular and \texttt{fail} if it is not. 



 \texttt{GroupHClass} is a synonym for \texttt{GroupHClassOfGreensDClass} (\textbf{Reference: GroupHClassOfGreensDClass}). 

 See also \texttt{IsGroupHClass} (\textbf{Reference: IsGroupHClass}), \texttt{IsRegularDClass} (\textbf{Reference: IsRegularDClass}), \texttt{IsRegularGreensClass} (\ref{IsRegularGreensClass}), and \texttt{IsRegularSemigroup} (\ref{IsRegularSemigroup}). 
\begin{Verbatim}[commandchars=!@|,fontsize=\small,frame=single,label=Example]
  !gapprompt@gap>| !gapinput@S := Semigroup(Transformation([2, 6, 7, 2, 6, 1, 1, 5]),|
  !gapprompt@>| !gapinput@                  Transformation([3, 8, 1, 4, 5, 6, 7, 1]));;|
  !gapprompt@gap>| !gapinput@IsRegularSemigroup(S);|
  false
  !gapprompt@gap>| !gapinput@iter := IteratorOfDClasses(S);;|
  !gapprompt@gap>| !gapinput@repeat D := NextIterator(iter); until IsRegularDClass(D);|
  !gapprompt@gap>| !gapinput@D;|
  <Green's D-class: Transformation( [ 6, 1, 1, 6, 1, 2, 2, 6 ] )>
  !gapprompt@gap>| !gapinput@NrIdempotents(D);|
  12
  !gapprompt@gap>| !gapinput@NrRClasses(D);|
  8
  !gapprompt@gap>| !gapinput@NrLClasses(D);|
  4
  !gapprompt@gap>| !gapinput@GroupHClass(D);|
  <Green's H-class: Transformation( [ 1, 2, 2, 1, 2, 6, 6, 1 ] )>
  !gapprompt@gap>| !gapinput@GroupHClassOfGreensDClass(D);|
  <Green's H-class: Transformation( [ 1, 2, 2, 1, 2, 6, 6, 1 ] )>
  !gapprompt@gap>| !gapinput@StructureDescription(GroupHClass(D));|
  "S3"
  !gapprompt@gap>| !gapinput@repeat D := NextIterator(iter); until not IsRegularDClass(D);|
  !gapprompt@gap>| !gapinput@D;|
  <Green's D-class: Transformation( [ 7, 5, 2, 2, 6, 1, 1, 2 ] )>
  !gapprompt@gap>| !gapinput@IsRegularDClass(D);|
  false
  !gapprompt@gap>| !gapinput@GroupHClass(D);|
  fail
  !gapprompt@gap>| !gapinput@S := InverseSemigroup(|
  !gapprompt@>| !gapinput@PartialPerm([2, 1, 6, 0, 3]), PartialPerm([3, 5, 2, 0, 0, 6]));;|
  !gapprompt@gap>| !gapinput@x := PartialPerm([1 .. 3], [6, 3, 1]);;|
  !gapprompt@gap>| !gapinput@First(DClasses(S), x -> not IsTrivial(GroupHClass(x)));|
  <Green's D-class: <identity partial perm on [ 1, 2 ]>>
  !gapprompt@gap>| !gapinput@StructureDescription(GroupHClass(last));|
  "C2"
\end{Verbatim}
 }

 

\subsection{\textcolor{Chapter }{SchutzenbergerGroup}}
\logpage{[ 10, 4, 2 ]}\nobreak
\hyperdef{L}{X84F1321E8217D2A8}{}
{\noindent\textcolor{FuncColor}{$\triangleright$\enspace\texttt{SchutzenbergerGroup({\mdseries\slshape class})\index{SchutzenbergerGroup@\texttt{SchutzenbergerGroup}}
\label{SchutzenbergerGroup}
}\hfill{\scriptsize (attribute)}}\\
\textbf{\indent Returns:\ }
 A group. 



 \texttt{SchutzenbergerGroup} returns the generalized Schutzenberger group (defined below) of the $\mathcal{R}$\texttt{\symbol{45}}, $\mathcal{D}$\texttt{\symbol{45}}, $\mathcal{L}$\texttt{\symbol{45}}, or $\mathcal{H}$\texttt{\symbol{45}}class \mbox{\texttt{\mdseries\slshape class}}. 

 If \texttt{f} is an element of a semigroup of transformations or partial permutations and \texttt{im(f)} denotes the image of \texttt{f}, then the \emph{generalized Schutzenberger group} of \texttt{im(f)} is the permutation group  
\[ \{\:g|_{\textrm{im}(f)}\::\:\textrm{im}(f*g)=\textrm{im}(f)\:\}. \]
  

 The generalized Schutzenberger group of the kernel \texttt{ker(f)} of a transformation \texttt{f} or the domain \texttt{dom(f)} of a partial permutation \texttt{f} is defined analogously. 

 The generalized Schutzenberger group of a Green's class is then defined as
follows. 
\begin{description}
\item[{$\mathcal{R}$\texttt{\symbol{45}}class}]  The generalized Schutzenberger group of the image or range of the
representative of the $\mathcal{R}$\texttt{\symbol{45}}class. 
\item[{$\mathcal{L}$\texttt{\symbol{45}}class}]  The generalized Schutzenberger group of the kernel or domain of the
representative of the $\mathcal{L}$\texttt{\symbol{45}}class. 
\item[{$\mathcal{H}$\texttt{\symbol{45}}class}]  The intersection of the generalized Schutzenberger groups of the $\mathcal{R}$\texttt{\symbol{45}} and $\mathcal{L}$\texttt{\symbol{45}}class containing the $\mathcal{H}$\texttt{\symbol{45}}class. 
\item[{$\mathcal{D}$\texttt{\symbol{45}}class}]  The intersection of the generalized Schutzenberger groups of the $\mathcal{R}$\texttt{\symbol{45}} and $\mathcal{L}$\texttt{\symbol{45}}class containing the representative of the $\mathcal{D}$\texttt{\symbol{45}}class. 
\end{description}
 The output of this attribute is difficult to describe for other types of
semigroup. However, a general description is given in \cite{Mitchell2019aa}. 
\begin{Verbatim}[commandchars=!@|,fontsize=\small,frame=single,label=Example]
  !gapprompt@gap>| !gapinput@S := Semigroup(Transformation([4, 4, 3, 5, 3]),|
  !gapprompt@>| !gapinput@                  Transformation([5, 1, 1, 4, 1]),|
  !gapprompt@>| !gapinput@                  Transformation([5, 5, 4, 4, 5]));;|
  !gapprompt@gap>| !gapinput@f := Transformation([5, 5, 4, 4, 5]);;|
  !gapprompt@gap>| !gapinput@SchutzenbergerGroup(RClass(S, f));|
  Group([ (4,5) ])
  !gapprompt@gap>| !gapinput@S := InverseSemigroup(|
  !gapprompt@>| !gapinput@PartialPerm([1, 2, 3, 7],|
  !gapprompt@>| !gapinput@            [9, 2, 4, 8]),|
  !gapprompt@>| !gapinput@PartialPerm([1, 2, 6, 7, 8, 9, 10],|
  !gapprompt@>| !gapinput@            [6, 8, 4, 5, 9, 1, 3]),|
  !gapprompt@>| !gapinput@PartialPerm([1, 2, 3, 5, 6, 7, 8, 9],|
  !gapprompt@>| !gapinput@            [7, 4, 1, 6, 9, 5, 2, 3]));;|
  !gapprompt@gap>| !gapinput@List(DClasses(S), SchutzenbergerGroup);|
  [ Group(()), Group(()), Group(()), Group(()), Group([ (4,9) ]),
    Group(()), Group(()), Group([ (5,8,6), (5,8) ]), Group(()),
    Group(()), Group(()), Group(()), Group(()), Group(()),
    Group([ (1,7,5,6,9,3) ]), Group([ (1,6)(3,5) ]), Group(()),
    Group(()), Group(()), Group(()), Group(()), Group(()), Group(()) ]
\end{Verbatim}
 }

 

\subsection{\textcolor{Chapter }{StructureDescriptionSchutzenbergerGroups}}
\logpage{[ 10, 4, 3 ]}\nobreak
\hyperdef{L}{X81202126806443F9}{}
{\noindent\textcolor{FuncColor}{$\triangleright$\enspace\texttt{StructureDescriptionSchutzenbergerGroups({\mdseries\slshape S})\index{StructureDescriptionSchutzenbergerGroups@\texttt{Structure}\-\texttt{Description}\-\texttt{Schutzenberger}\-\texttt{Groups}}
\label{StructureDescriptionSchutzenbergerGroups}
}\hfill{\scriptsize (attribute)}}\\
\textbf{\indent Returns:\ }
Distinct structure descriptions of the Schutzenberger groups of a semigroup.



 \texttt{StructureDescriptionSchutzenbergerGroups} returns the distinct values of \texttt{StructureDescription} (\textbf{Reference: StructureDescription}) when it is applied to the Schutzenberger groups of the $\mathcal{R}$\texttt{\symbol{45}}classes of the semigroup \mbox{\texttt{\mdseries\slshape S}}. 
\begin{Verbatim}[commandchars=!@|,fontsize=\small,frame=single,label=Example]
  !gapprompt@gap>| !gapinput@S := Semigroup([|
  !gapprompt@>| !gapinput@ PartialPerm([1, 2, 3], [2, 5, 4]),|
  !gapprompt@>| !gapinput@ PartialPerm([1, 2, 3], [4, 1, 2]),|
  !gapprompt@>| !gapinput@ PartialPerm([1, 2, 3], [5, 2, 3]),|
  !gapprompt@>| !gapinput@ PartialPerm([1, 2, 4, 5], [2, 1, 4, 3]),|
  !gapprompt@>| !gapinput@ PartialPerm([1, 2, 5], [2, 3, 5]),|
  !gapprompt@>| !gapinput@ PartialPerm([1, 2, 3, 5], [2, 3, 5, 4]),|
  !gapprompt@>| !gapinput@ PartialPerm([1, 2, 3, 5], [4, 2, 5, 1]),|
  !gapprompt@>| !gapinput@ PartialPerm([1, 2, 3, 5], [5, 2, 4, 3]),|
  !gapprompt@>| !gapinput@ PartialPerm([1, 2, 5], [5, 4, 3])]);;|
  !gapprompt@gap>| !gapinput@StructureDescriptionSchutzenbergerGroups(S);|
  [ "1", "C2", "S3" ]
  !gapprompt@gap>| !gapinput@S := Monoid(|
  !gapprompt@>| !gapinput@Bipartition([[1, 2, 5, -1, -2], [3, 4, -3, -5], [-4]]),|
  !gapprompt@>| !gapinput@Bipartition([[1, 2, -2], [3, -1], [4], [5], [-3, -4], [-5]]),|
  !gapprompt@>| !gapinput@Bipartition([[1], [2, 3, -5], [4, -3], [5, -2], [-1, -4]]));|
  <bipartition monoid of degree 5 with 3 generators>
  !gapprompt@gap>| !gapinput@StructureDescriptionSchutzenbergerGroups(S);|
  [ "1", "C2" ]
\end{Verbatim}
 }

 

\subsection{\textcolor{Chapter }{StructureDescriptionMaximalSubgroups}}
\logpage{[ 10, 4, 4 ]}\nobreak
\hyperdef{L}{X838F43FE79A8C678}{}
{\noindent\textcolor{FuncColor}{$\triangleright$\enspace\texttt{StructureDescriptionMaximalSubgroups({\mdseries\slshape S})\index{StructureDescriptionMaximalSubgroups@\texttt{Structure}\-\texttt{Description}\-\texttt{Maximal}\-\texttt{Subgroups}}
\label{StructureDescriptionMaximalSubgroups}
}\hfill{\scriptsize (attribute)}}\\
\textbf{\indent Returns:\ }
Distinct structure descriptions of the maximal subgroups of a semigroup.



 \texttt{StructureDescriptionMaximalSubgroups} returns the distinct values of \texttt{StructureDescription} (\textbf{Reference: StructureDescription}) when it is applied to the maximal subgroups of the semigroup \mbox{\texttt{\mdseries\slshape S}}. 
\begin{Verbatim}[commandchars=!@|,fontsize=\small,frame=single,label=Example]
  !gapprompt@gap>| !gapinput@S := DualSymmetricInverseSemigroup(6);|
  <inverse block bijection monoid of degree 6 with 3 generators>
  !gapprompt@gap>| !gapinput@StructureDescriptionMaximalSubgroups(S);|
  [ "1", "C2", "S3", "S4", "S5", "S6" ]
  !gapprompt@gap>| !gapinput@S := Semigroup(|
  !gapprompt@>| !gapinput@ PartialPerm([1, 3, 4, 5, 8],|
  !gapprompt@>| !gapinput@             [8, 3, 9, 4, 5]),|
  !gapprompt@>| !gapinput@ PartialPerm([1, 2, 3, 4, 8],|
  !gapprompt@>| !gapinput@             [10, 4, 1, 9, 6]),|
  !gapprompt@>| !gapinput@ PartialPerm([1, 2, 3, 4, 5, 6, 7, 10],|
  !gapprompt@>| !gapinput@             [4, 1, 6, 7, 5, 3, 2, 10]),|
  !gapprompt@>| !gapinput@ PartialPerm([1, 2, 3, 4, 6, 8, 10],|
  !gapprompt@>| !gapinput@             [4, 9, 10, 3, 1, 5, 2]));;|
  !gapprompt@gap>| !gapinput@StructureDescriptionMaximalSubgroups(S);|
  [ "1", "C2", "C3", "C4" ]
\end{Verbatim}
 }

 

\subsection{\textcolor{Chapter }{MultiplicativeNeutralElement (for an H-class)}}
\logpage{[ 10, 4, 5 ]}\nobreak
\hyperdef{L}{X8459E4067C5773AD}{}
{\noindent\textcolor{FuncColor}{$\triangleright$\enspace\texttt{MultiplicativeNeutralElement({\mdseries\slshape H})\index{MultiplicativeNeutralElement@\texttt{MultiplicativeNeutralElement}!for an H-class}
\label{MultiplicativeNeutralElement:for an H-class}
}\hfill{\scriptsize (method)}}\\
\textbf{\indent Returns:\ }
A semigroup element or \texttt{fail}.



 If the $\mathcal{H}$\texttt{\symbol{45}}class \mbox{\texttt{\mdseries\slshape H}} of a semigroup \texttt{S} is a subgroup of \texttt{S}, then \texttt{MultiplicativeNeutralElement} returns the identity of \mbox{\texttt{\mdseries\slshape H}}. If \mbox{\texttt{\mdseries\slshape H}} is not a subgroup of \texttt{S}, then \texttt{fail} is returned. 
\begin{Verbatim}[commandchars=!@|,fontsize=\small,frame=single,label=Example]
  !gapprompt@gap>| !gapinput@S := Semigroup([PartialPerm([1, 5, 2]),|
  !gapprompt@>| !gapinput@ PartialPerm([2, 0, 4]), PartialPerm([4, 1, 5]),|
  !gapprompt@>| !gapinput@ PartialPerm([1, 0, 3, 0, 4]), PartialPerm([1, 2, 0, 3, 5]),|
  !gapprompt@>| !gapinput@ PartialPerm([1, 3, 2, 0, 5]), PartialPerm([5, 0, 0, 4, 3])]);;|
  !gapprompt@gap>| !gapinput@H := HClass(S, PartialPerm([1, 2]));;|
  !gapprompt@gap>| !gapinput@MultiplicativeNeutralElement(H);|
  <identity partial perm on [ 1, 2 ]>
  !gapprompt@gap>| !gapinput@H := HClass(S, PartialPerm([1, 4]));;|
  !gapprompt@gap>| !gapinput@MultiplicativeNeutralElement(H);|
  fail
\end{Verbatim}
 }

 

\subsection{\textcolor{Chapter }{StructureDescription (for an H-class)}}
\logpage{[ 10, 4, 6 ]}\nobreak
\hyperdef{L}{X85B34FFB82C83127}{}
{\noindent\textcolor{FuncColor}{$\triangleright$\enspace\texttt{StructureDescription({\mdseries\slshape class})\index{StructureDescription@\texttt{StructureDescription}!for an H-class}
\label{StructureDescription:for an H-class}
}\hfill{\scriptsize (attribute)}}\\
\textbf{\indent Returns:\ }
A string or \texttt{fail}.



 \texttt{StructureDescription} returns the value of \texttt{StructureDescription} (\textbf{Reference: StructureDescription}) when it is applied to a group isomorphic to the group $\mathcal{H}$\texttt{\symbol{45}}class \mbox{\texttt{\mdseries\slshape class}}. If \mbox{\texttt{\mdseries\slshape class}} is not a group $\mathcal{H}$\texttt{\symbol{45}}class, then \texttt{fail} is returned. 
\begin{Verbatim}[commandchars=!@|,fontsize=\small,frame=single,label=Example]
  !gapprompt@gap>| !gapinput@S := Semigroup(|
  !gapprompt@>| !gapinput@PartialPerm([1, 2, 3, 4, 6, 7, 8, 9],|
  !gapprompt@>| !gapinput@            [1, 9, 4, 3, 5, 2, 10, 7]),|
  !gapprompt@>| !gapinput@PartialPerm([1, 2, 4, 7, 8, 9],|
  !gapprompt@>| !gapinput@            [6, 2, 4, 9, 1, 3]));;|
  !gapprompt@gap>| !gapinput@H := HClass(S, PartialPerm([1, 2, 3, 4, 7, 9],|
  !gapprompt@>| !gapinput@                              [1, 7, 3, 4, 9, 2]));;|
  !gapprompt@gap>| !gapinput@StructureDescription(H);|
  "C6"
\end{Verbatim}
 }

 

\subsection{\textcolor{Chapter }{InjectionPrincipalFactor}}
\logpage{[ 10, 4, 7 ]}\nobreak
\hyperdef{L}{X7EBB4F1981CC2AE9}{}
{\noindent\textcolor{FuncColor}{$\triangleright$\enspace\texttt{InjectionPrincipalFactor({\mdseries\slshape D})\index{InjectionPrincipalFactor@\texttt{InjectionPrincipalFactor}}
\label{InjectionPrincipalFactor}
}\hfill{\scriptsize (attribute)}}\\
\noindent\textcolor{FuncColor}{$\triangleright$\enspace\texttt{InjectionNormalizedPrincipalFactor({\mdseries\slshape D})\index{InjectionNormalizedPrincipalFactor@\texttt{InjectionNormalizedPrincipalFactor}}
\label{InjectionNormalizedPrincipalFactor}
}\hfill{\scriptsize (attribute)}}\\
\noindent\textcolor{FuncColor}{$\triangleright$\enspace\texttt{IsomorphismReesMatrixSemigroup({\mdseries\slshape D})\index{IsomorphismReesMatrixSemigroup@\texttt{IsomorphismReesMatrixSemigroup}!for a D-class}
\label{IsomorphismReesMatrixSemigroup:for a D-class}
}\hfill{\scriptsize (attribute)}}\\
\textbf{\indent Returns:\ }
A injective mapping.



 If the $\mathcal{D}$\texttt{\symbol{45}}class \mbox{\texttt{\mdseries\slshape D}} is a subsemigroup of a semigroup \texttt{S}, then the \emph{principal factor} of \mbox{\texttt{\mdseries\slshape D}} is just \mbox{\texttt{\mdseries\slshape D}} itself. If \mbox{\texttt{\mdseries\slshape D}} is not a subsemigroup of \texttt{S}, then the principal factor of \mbox{\texttt{\mdseries\slshape D}} is the semigroup with elements \mbox{\texttt{\mdseries\slshape D}} and a new element \texttt{0} with multiplication of $x,y\in D$ defined by:  
\[ xy=\left\{\begin{array}{ll} x*y\ (\textrm{in }S)&\textrm{if }x*y\in D\\
0&\textrm{if }xy\not\in D. \end{array}\right. \]
   \texttt{InjectionPrincipalFactor} returns an injective function from the $\mathcal{D}$\texttt{\symbol{45}}class \mbox{\texttt{\mdseries\slshape D}} to a Rees (0\texttt{\symbol{45}})matrix semigroup, which contains the
principal factor of \mbox{\texttt{\mdseries\slshape D}} as a subsemigroup. 

 If \mbox{\texttt{\mdseries\slshape D}} is a subsemigroup of its parent semigroup, then the function returned by \texttt{InjectionPrincipalFactor} or \texttt{IsomorphismReesMatrixSemigroup} is an isomorphism from \mbox{\texttt{\mdseries\slshape D}} to a Rees matrix semigroup; see \texttt{ReesMatrixSemigroup} (\textbf{Reference: ReesMatrixSemigroup}).

 If \mbox{\texttt{\mdseries\slshape D}} is not a semigroup, then the function returned by \texttt{InjectionPrincipalFactor} is an injective function from \mbox{\texttt{\mdseries\slshape D}} to a Rees 0\texttt{\symbol{45}}matrix semigroup isomorphic to the principal
factor of \mbox{\texttt{\mdseries\slshape D}}; see \texttt{ReesZeroMatrixSemigroup} (\textbf{Reference: ReesZeroMatrixSemigroup}). In this case, \texttt{IsomorphismReesMatrixSemigroup} and \texttt{IsomorphismReesZeroMatrixSemigroup} returns an error.

 \texttt{InjectionNormalizedPrincipalFactor} returns the composition of \texttt{InjectionPrincipalFactor} with \texttt{RZMSNormalization} (\ref{RZMSNormalization}) or \texttt{RMSNormalization} (\ref{RMSNormalization}) as appropriate.

 See also \texttt{PrincipalFactor} (\ref{PrincipalFactor}). 
\begin{Verbatim}[commandchars=!@|,fontsize=\small,frame=single,label=Example]
  !gapprompt@gap>| !gapinput@S := InverseSemigroup(|
  !gapprompt@>| !gapinput@PartialPerm([1, 2, 3, 6, 8, 10],|
  !gapprompt@>| !gapinput@            [2, 6, 7, 9, 1, 5]),|
  !gapprompt@>| !gapinput@PartialPerm([1, 2, 3, 4, 6, 7, 8, 10],|
  !gapprompt@>| !gapinput@            [3, 8, 1, 9, 4, 10, 5, 6]));;|
  !gapprompt@gap>| !gapinput@x := PartialPerm([1, 2, 5, 6, 7, 9],|
  !gapprompt@>| !gapinput@                    [1, 2, 5, 6, 7, 9]);;|
  !gapprompt@gap>| !gapinput@D := GreensDClassOfElement(S, x);|
  <Green's D-class: <identity partial perm on [ 1, 2, 5, 6, 7, 9 ]>>
  !gapprompt@gap>| !gapinput@R := Range(InjectionPrincipalFactor(D));|
  <Rees 0-matrix semigroup 3x3 over Group(())>
  !gapprompt@gap>| !gapinput@MatrixOfReesZeroMatrixSemigroup(R);|
  [ [ (), 0, 0 ], [ 0, (), 0 ], [ 0, 0, () ] ]
  !gapprompt@gap>| !gapinput@Size(R);|
  10
  !gapprompt@gap>| !gapinput@Size(D);|
  9
  !gapprompt@gap>| !gapinput@S := Semigroup(|
  !gapprompt@>| !gapinput@Bipartition([[1, 2, 3, -3, -5], [4], [5, -2], [-1, -4]]),|
  !gapprompt@>| !gapinput@Bipartition([[1, 3, 5], [2, 4, -3], [-1, -2, -4, -5]]),|
  !gapprompt@>| !gapinput@Bipartition([[1, 5, -2, -4], [2, 3, 4, -1, -5], [-3]]),|
  !gapprompt@>| !gapinput@Bipartition([[1, 5, -1, -2, -3], [2, 4, -4], [3, -5]]));;|
  !gapprompt@gap>| !gapinput@D := GreensDClassOfElement(S,|
  !gapprompt@>| !gapinput@Bipartition([[1, 5, -2, -4], [2, 3, 4, -1, -5], [-3]]));|
  <Green's D-class: <bipartition: [ 1, 5, -2, -4 ], [ 2, 3, 4, -1, -5 ]
     , [ -3 ]>>
  !gapprompt@gap>| !gapinput@InjectionNormalizedPrincipalFactor(D);|
  MappingByFunction( <Green's D-class: <bipartition: [ 1, 5, -2, -4 ],
    [ 2, 3, 4, -1, -5 ], [ -3 ]>>, <Rees matrix semigroup 1x1 over
    Group([ (1,2) ])>, function( x ) ... end, function( x ) ... end )
\end{Verbatim}
 }

 

\subsection{\textcolor{Chapter }{PrincipalFactor}}
\logpage{[ 10, 4, 8 ]}\nobreak
\hyperdef{L}{X86C6D777847AAEC7}{}
{\noindent\textcolor{FuncColor}{$\triangleright$\enspace\texttt{PrincipalFactor({\mdseries\slshape D})\index{PrincipalFactor@\texttt{PrincipalFactor}}
\label{PrincipalFactor}
}\hfill{\scriptsize (attribute)}}\\
\noindent\textcolor{FuncColor}{$\triangleright$\enspace\texttt{NormalizedPrincipalFactor({\mdseries\slshape D})\index{NormalizedPrincipalFactor@\texttt{NormalizedPrincipalFactor}}
\label{NormalizedPrincipalFactor}
}\hfill{\scriptsize (attribute)}}\\
\textbf{\indent Returns:\ }
A Rees (0\texttt{\symbol{45}})matrix semigroup.



 If \mbox{\texttt{\mdseries\slshape D}} is a $\mathcal{D}$\texttt{\symbol{45}}class of semigroup, then \texttt{PrincipalFactor(\mbox{\texttt{\mdseries\slshape D}})} is just shorthand for \texttt{Range(InjectionPrincipalFactor(\mbox{\texttt{\mdseries\slshape D}}))}, and \texttt{NormalizedPrincipalFactor(\mbox{\texttt{\mdseries\slshape D}})} is shorthand for \texttt{Range(InjectionNormalizedPrincipalFactor(\mbox{\texttt{\mdseries\slshape D}}))}. 

 See \texttt{InjectionPrincipalFactor} (\ref{InjectionPrincipalFactor}) and \texttt{InjectionNormalizedPrincipalFactor} (\ref{InjectionNormalizedPrincipalFactor}) for more details. 
\begin{Verbatim}[commandchars=!@|,fontsize=\small,frame=single,label=Example]
  !gapprompt@gap>| !gapinput@S := Semigroup([PartialPerm([1, 2, 3], [1, 3, 4]),|
  !gapprompt@>| !gapinput@ PartialPerm([1, 2, 3], [2, 5, 3]),|
  !gapprompt@>| !gapinput@ PartialPerm([1, 2, 3, 4], [2, 4, 1, 5]),|
  !gapprompt@>| !gapinput@ PartialPerm([1, 3, 5], [5, 1, 3])]);;|
  !gapprompt@gap>| !gapinput@PrincipalFactor(MinimalDClass(S));|
  <Rees matrix semigroup 1x1 over Group(())>
  !gapprompt@gap>| !gapinput@MultiplicativeZero(S);|
  <empty partial perm>
  !gapprompt@gap>| !gapinput@S := Semigroup(|
  !gapprompt@>| !gapinput@Bipartition([[1, 2, 3, 4, 5, -1, -3], [-2, -5], [-4]]),|
  !gapprompt@>| !gapinput@Bipartition([[1, -5], [2, 3, 4, 5, -1, -3], [-2, -4]]),|
  !gapprompt@>| !gapinput@Bipartition([[1, 5, -4], [2, 4, -1, -5], [3, -2, -3]]));;|
  !gapprompt@gap>| !gapinput@D := MinimalDClass(S);|
  <Green's D-class: <bipartition: [ 1, 2, 3, 4, 5, -1, -3 ],
    [ -2, -5 ], [ -4 ]>>
  !gapprompt@gap>| !gapinput@NormalizedPrincipalFactor(D);|
  <Rees matrix semigroup 1x5 over Group(())>
\end{Verbatim}
 }

 }

   
\section{\textcolor{Chapter }{ Operations for Green's relations and classes }}\logpage{[ 10, 5, 0 ]}
\hyperdef{L}{X802E2BC9828341A2}{}
{
  In this section, we describe some operations related to Green's classes that
are available in the \textsf{Semigroups} package. 

 

\subsection{\textcolor{Chapter }{LeftGreensMultiplier}}
\logpage{[ 10, 5, 1 ]}\nobreak
\hyperdef{L}{X7EDE3F03879B2B12}{}
{\noindent\textcolor{FuncColor}{$\triangleright$\enspace\texttt{LeftGreensMultiplier({\mdseries\slshape S, a, b})\index{LeftGreensMultiplier@\texttt{LeftGreensMultiplier}}
\label{LeftGreensMultiplier}
}\hfill{\scriptsize (operation)}}\\
\noindent\textcolor{FuncColor}{$\triangleright$\enspace\texttt{RightGreensMultiplier({\mdseries\slshape S, a, b})\index{RightGreensMultiplier@\texttt{RightGreensMultiplier}}
\label{RightGreensMultiplier}
}\hfill{\scriptsize (operation)}}\\
\textbf{\indent Returns:\ }
 An element. 



 If \mbox{\texttt{\mdseries\slshape S}} is a semigroup, and \mbox{\texttt{\mdseries\slshape a}} and \mbox{\texttt{\mdseries\slshape b}} are $\mathcal{L}$\texttt{\symbol{45}}related elements of \mbox{\texttt{\mdseries\slshape S}}, then \texttt{LeftGreensMultiplier} returns an element \texttt{s} such that \texttt{s * \mbox{\texttt{\mdseries\slshape a}} = \mbox{\texttt{\mdseries\slshape b}}}. The element \texttt{s} is of the same type as the elements of \mbox{\texttt{\mdseries\slshape S}} but may or may not be an element of \mbox{\texttt{\mdseries\slshape S}}. In particular, if \mbox{\texttt{\mdseries\slshape S}} is not a monoid and \texttt{\mbox{\texttt{\mdseries\slshape a}} = \mbox{\texttt{\mdseries\slshape b}}}, then \texttt{One(GeneratorsOfSemigroup(S))} or an adjoined identity may be returned. Even if \texttt{\mbox{\texttt{\mdseries\slshape a}} {\textless}{\textgreater} \mbox{\texttt{\mdseries\slshape b}}}, then it is not guaranteed that the returned element \texttt{s} will belong to \mbox{\texttt{\mdseries\slshape S}}. It is guaranteed that the left action of \texttt{s} on the elements of the $\mathcal{L}$\texttt{\symbol{45}}class of \mbox{\texttt{\mdseries\slshape a}} is the same as the left action of an element of \mbox{\texttt{\mdseries\slshape S}} with the identity adjoined.

 \texttt{LeftGreensMultiplier} gives an error if \mbox{\texttt{\mdseries\slshape a}} and \mbox{\texttt{\mdseries\slshape b}} are not $\mathcal{L}$\texttt{\symbol{45}}related elements of \mbox{\texttt{\mdseries\slshape S}}.

 The operation \texttt{RightGreensMultiplier} is defined analogously. 
\begin{Verbatim}[commandchars=!@|,fontsize=\small,frame=single,label=Example]
  !gapprompt@gap>| !gapinput@S := Semigroup(Transformation([4, 4, 3, 5, 3]),|
  !gapprompt@>| !gapinput@                  Transformation([5, 1, 1, 4, 1]),|
  !gapprompt@>| !gapinput@                  Transformation([5, 5, 4, 4, 5]));;|
  !gapprompt@gap>| !gapinput@a := Transformation([5, 5, 4, 4, 5]);;|
  !gapprompt@gap>| !gapinput@LeftGreensMultiplier(S, a, a);|
  Transformation( [ 1, 1, 3, 3, 1 ] )
  !gapprompt@gap>| !gapinput@RightGreensMultiplier(S, a, a);|
  Transformation( [ 5, 5, 5, 4, 5 ] )
  !gapprompt@gap>| !gapinput@b := Transformation([5, 4, 4, 5, 4]);|
  Transformation( [ 5, 4, 4, 5, 4 ] )
  !gapprompt@gap>| !gapinput@s := LeftGreensMultiplier(S, a, b);|
  Transformation( [ 1, 3, 3, 1, 3 ] )
  !gapprompt@gap>| !gapinput@s * a;|
  Transformation( [ 5, 4, 4, 5, 4 ] )
  !gapprompt@gap>| !gapinput@b := Transformation([4, 4, 5, 5, 4]);|
  Transformation( [ 4, 4, 5, 5, 4 ] )
  !gapprompt@gap>| !gapinput@s := RightGreensMultiplier(S, a, b);|
  Transformation( [ 4, 4, 4, 5, 4 ] )
  !gapprompt@gap>| !gapinput@a * s = b;|
  true
\end{Verbatim}
 }

 }

   }

 
\chapter{\textcolor{Chapter }{ Attributes and operations for semigroups }}\label{Attributes and operations for semigroups}
\logpage{[ 11, 0, 0 ]}
\hyperdef{L}{X7C75B1DB81C7779B}{}
{
  In this chapter we describe the methods that are available in \textsf{Semigroups} for determining the attributes of a semigroup, and the operations which can be
applied to a semigroup.   
\section{\textcolor{Chapter }{ Accessing the elements of a semigroup }}\logpage{[ 11, 1, 0 ]}
\hyperdef{L}{X7AE0630287B8A757}{}
{
  

\subsection{\textcolor{Chapter }{AsListCanonical}}
\logpage{[ 11, 1, 1 ]}\nobreak
\hyperdef{L}{X7AC3FAA5826516AD}{}
{\noindent\textcolor{FuncColor}{$\triangleright$\enspace\texttt{AsListCanonical({\mdseries\slshape S})\index{AsListCanonical@\texttt{AsListCanonical}}
\label{AsListCanonical}
}\hfill{\scriptsize (attribute)}}\\
\noindent\textcolor{FuncColor}{$\triangleright$\enspace\texttt{EnumeratorCanonical({\mdseries\slshape S})\index{EnumeratorCanonical@\texttt{EnumeratorCanonical}}
\label{EnumeratorCanonical}
}\hfill{\scriptsize (attribute)}}\\
\noindent\textcolor{FuncColor}{$\triangleright$\enspace\texttt{IteratorCanonical({\mdseries\slshape S})\index{IteratorCanonical@\texttt{IteratorCanonical}}
\label{IteratorCanonical}
}\hfill{\scriptsize (operation)}}\\
\textbf{\indent Returns:\ }
A list, enumerator, or iterator.



 When the argument \mbox{\texttt{\mdseries\slshape S}} is a semigroup satisfying \texttt{CanUseFroidurePin} (\ref{CanUseFroidurePin}), \texttt{AsListCanonical} returns a list of the elements of \mbox{\texttt{\mdseries\slshape S}} in the order they are enumerated by the Froidure\texttt{\symbol{45}}Pin
Algorithm. This is the same as the order used to index the elements of \mbox{\texttt{\mdseries\slshape S}} in \texttt{RightCayleyDigraph} (\ref{RightCayleyDigraph}) and \texttt{LeftCayleyDigraph} (\ref{LeftCayleyDigraph}). 

 \texttt{EnumeratorCanonical} and \texttt{IteratorCanonical} return an enumerator and an iterator where the elements are ordered in the
same way as \texttt{AsListCanonical}. Using \texttt{EnumeratorCanonical} or \texttt{IteratorCanonical} will often use less memory than \texttt{AsListCanonical}, but may have slightly worse performance if the elements of the semigroup are
looped over repeatedly. \texttt{EnumeratorCanonical} returns the same list as \texttt{AsListCanonical} if \texttt{AsListCanonical} has ever been called for \mbox{\texttt{\mdseries\slshape S}}.

 If \mbox{\texttt{\mdseries\slshape S}} is an acting semigroup, then the value returned by \texttt{AsList} may not equal the value returned by \texttt{AsListCanonical}. \texttt{AsListCanonical} exists so that there is a method for obtaining the elements of \mbox{\texttt{\mdseries\slshape S}} in the particular order used by \texttt{RightCayleyDigraph} (\ref{RightCayleyDigraph}) and \texttt{LeftCayleyDigraph} (\ref{LeftCayleyDigraph}).

 See also \texttt{PositionCanonical} (\ref{PositionCanonical}). 
\begin{Verbatim}[commandchars=!@|,fontsize=\small,frame=single,label=Example]
  !gapprompt@gap>| !gapinput@S := Semigroup(Transformation([1, 3, 2]));;|
  !gapprompt@gap>| !gapinput@AsListCanonical(S);|
  [ Transformation( [ 1, 3, 2 ] ), IdentityTransformation ]
  !gapprompt@gap>| !gapinput@IteratorCanonical(S);|
  <iterator>
  !gapprompt@gap>| !gapinput@EnumeratorCanonical(S);|
  [ Transformation( [ 1, 3, 2 ] ), IdentityTransformation ]
  !gapprompt@gap>| !gapinput@S := Monoid([Matrix(IsBooleanMat, [[1, 0, 0],|
  !gapprompt@>| !gapinput@                                      [0, 1, 0],|
  !gapprompt@>| !gapinput@                                      [0, 1, 0]])]);|
  <commutative monoid of 3x3 boolean matrices with 1 generator>
  !gapprompt@gap>| !gapinput@it := IteratorCanonical(S);|
  <iterator>
  !gapprompt@gap>| !gapinput@NextIterator(it);|
  Matrix(IsBooleanMat, [[1, 0, 0], [0, 1, 0], [0, 0, 1]])
  !gapprompt@gap>| !gapinput@en := EnumeratorCanonical(S);|
  <enumerator of <commutative monoid of size 2, 3x3 boolean matrices
   with 1 generator>>
  !gapprompt@gap>| !gapinput@en[1];|
  Matrix(IsBooleanMat, [[1, 0, 0], [0, 1, 0], [0, 0, 1]])
  !gapprompt@gap>| !gapinput@Position(en, en[1]);|
  1
  !gapprompt@gap>| !gapinput@Length(en);|
  2
\end{Verbatim}
 }

 

\subsection{\textcolor{Chapter }{PositionCanonical}}
\logpage{[ 11, 1, 2 ]}\nobreak
\hyperdef{L}{X7B4B10AE81602D4E}{}
{\noindent\textcolor{FuncColor}{$\triangleright$\enspace\texttt{PositionCanonical({\mdseries\slshape S, x})\index{PositionCanonical@\texttt{PositionCanonical}}
\label{PositionCanonical}
}\hfill{\scriptsize (operation)}}\\
\textbf{\indent Returns:\ }
A positive integer or \texttt{fail}.



 When the argument \mbox{\texttt{\mdseries\slshape S}} is a semigroup satisfying \texttt{CanUseFroidurePin} (\ref{CanUseFroidurePin}) and \mbox{\texttt{\mdseries\slshape x}} is an element of \mbox{\texttt{\mdseries\slshape S}}, \texttt{PositionCanonical} returns the position of \mbox{\texttt{\mdseries\slshape x}} in \texttt{AsListCanonical(\mbox{\texttt{\mdseries\slshape S}})} or equivalently the position of \mbox{\texttt{\mdseries\slshape x}} in \texttt{EnumeratorCanonical(\mbox{\texttt{\mdseries\slshape S}})}.

 See also \texttt{AsListCanonical} (\ref{AsListCanonical}) and \texttt{EnumeratorCanonical} (\ref{EnumeratorCanonical}). 
\begin{Verbatim}[commandchars=!@|,fontsize=\small,frame=single,label=Example]
  !gapprompt@gap>| !gapinput@S := FullTropicalMaxPlusMonoid(2, 3);|
  <monoid of 2x2 tropical max-plus matrices with 13 generators>
  !gapprompt@gap>| !gapinput@x := Matrix(IsTropicalMaxPlusMatrix, [[1, 3], [2, 1]], 3);|
  Matrix(IsTropicalMaxPlusMatrix, [[1, 3], [2, 1]], 3)
  !gapprompt@gap>| !gapinput@PositionCanonical(S, x);|
  234
  !gapprompt@gap>| !gapinput@EnumeratorCanonical(S)[234] = x;|
  true
\end{Verbatim}
 }

 

\subsection{\textcolor{Chapter }{Enumerate}}
\logpage{[ 11, 1, 3 ]}\nobreak
\hyperdef{L}{X7BCD5342793C7A7E}{}
{\noindent\textcolor{FuncColor}{$\triangleright$\enspace\texttt{Enumerate({\mdseries\slshape S[, limit]})\index{Enumerate@\texttt{Enumerate}}
\label{Enumerate}
}\hfill{\scriptsize (operation)}}\\
\textbf{\indent Returns:\ }
A semigroup (the argument).



 If \mbox{\texttt{\mdseries\slshape S}} is a semigroup with representation \texttt{CanUseFroidurePin} (\ref{CanUseFroidurePin}) and \mbox{\texttt{\mdseries\slshape limit}} is a positive integer, then this operation can be used to enumerate at least \mbox{\texttt{\mdseries\slshape limit}} elements of \mbox{\texttt{\mdseries\slshape S}}, or \texttt{Size(\mbox{\texttt{\mdseries\slshape S}})} elements if this is less than \mbox{\texttt{\mdseries\slshape limit}}, using the Froidure\texttt{\symbol{45}}Pin Algorithm. 

 If the optional second argument \mbox{\texttt{\mdseries\slshape limit}} is not given, then the semigroup is enumerated until all of its elements have
been found. 

 
\begin{Verbatim}[commandchars=!@|,fontsize=\small,frame=single,label=Example]
  !gapprompt@gap>| !gapinput@S := FullTransformationMonoid(7);|
  <full transformation monoid of degree 7>
  !gapprompt@gap>| !gapinput@Enumerate(S, 1000);|
  <full transformation monoid of degree 7>
\end{Verbatim}
 }

 

\subsection{\textcolor{Chapter }{IsEnumerated}}
\logpage{[ 11, 1, 4 ]}\nobreak
\hyperdef{L}{X877FAAA67F834897}{}
{\noindent\textcolor{FuncColor}{$\triangleright$\enspace\texttt{IsEnumerated({\mdseries\slshape S})\index{IsEnumerated@\texttt{IsEnumerated}}
\label{IsEnumerated}
}\hfill{\scriptsize (operation)}}\\
\textbf{\indent Returns:\ }
\texttt{true} or \texttt{false}.



 If \mbox{\texttt{\mdseries\slshape S}} is a semigroup with representation \texttt{CanUseFroidurePin} (\ref{CanUseFroidurePin}), then this operation returns \texttt{true} if the Froidure\texttt{\symbol{45}}Pin Algorithm has been run to completion
(i.e. all of the elements of \mbox{\texttt{\mdseries\slshape S}} have been found) and \texttt{false} if \mbox{\texttt{\mdseries\slshape S}} has not been fully enumerated. }

 }

   
\section{\textcolor{Chapter }{ Cayley graphs }}\logpage{[ 11, 2, 0 ]}
\hyperdef{L}{X789D5E5A8558AA07}{}
{
  

\subsection{\textcolor{Chapter }{RightCayleyDigraph}}
\logpage{[ 11, 2, 1 ]}\nobreak
\hyperdef{L}{X7EA002E27B10CCE0}{}
{\noindent\textcolor{FuncColor}{$\triangleright$\enspace\texttt{RightCayleyDigraph({\mdseries\slshape S})\index{RightCayleyDigraph@\texttt{RightCayleyDigraph}}
\label{RightCayleyDigraph}
}\hfill{\scriptsize (attribute)}}\\
\noindent\textcolor{FuncColor}{$\triangleright$\enspace\texttt{LeftCayleyDigraph({\mdseries\slshape S})\index{LeftCayleyDigraph@\texttt{LeftCayleyDigraph}}
\label{LeftCayleyDigraph}
}\hfill{\scriptsize (attribute)}}\\
\textbf{\indent Returns:\ }
A digraph.



 When the argument \mbox{\texttt{\mdseries\slshape S}} is a semigroup satisfying \texttt{CanUseFroidurePin} (\ref{CanUseFroidurePin}), \texttt{RightCayleyDigraph} returns the right Cayley graph of \mbox{\texttt{\mdseries\slshape S}}, as a \texttt{Digraph} (\textbf{Digraphs: Digraph}) \texttt{digraph} where vertex \texttt{OutNeighbours(digraph)[i][j]} is \texttt{PositionCanonical(\mbox{\texttt{\mdseries\slshape S}}, AsListCanonical(\mbox{\texttt{\mdseries\slshape S}})[i] * GeneratorsOfSemigroup(\mbox{\texttt{\mdseries\slshape S}})[j])}. The digraph returned by \texttt{LeftCayleyDigraph} is defined analogously.

 The digraph returned by this attribute belongs to the category \texttt{IsCayleyDigraph} (\textbf{Digraphs: IsCayleyDigraph}), the semigroup \mbox{\texttt{\mdseries\slshape S}} and the generators used to create the digraph can be recovered from the
digraph using \texttt{SemigroupOfCayleyDigraph} (\textbf{Digraphs: SemigroupOfCayleyDigraph}) and \texttt{GeneratorsOfCayleyDigraph} (\textbf{Digraphs: GeneratorsOfCayleyDigraph}). 
\begin{Verbatim}[commandchars=!@|,fontsize=\small,frame=single,label=Example]
  !gapprompt@gap>| !gapinput@S := FullTransformationMonoid(2);|
  <full transformation monoid of degree 2>
  !gapprompt@gap>| !gapinput@RightCayleyDigraph(S);|
  <immutable multidigraph with 4 vertices, 12 edges>
  !gapprompt@gap>| !gapinput@LeftCayleyDigraph(S);|
  <immutable multidigraph with 4 vertices, 12 edges>
\end{Verbatim}
 }

 }

   
\section{\textcolor{Chapter }{ Random elements of a semigroup }}\logpage{[ 11, 3, 0 ]}
\hyperdef{L}{X824184C785BF12FF}{}
{
  

\subsection{\textcolor{Chapter }{Random (for a semigroup)}}
\logpage{[ 11, 3, 1 ]}\nobreak
\hyperdef{L}{X7BB7FDFE7AFFD672}{}
{\noindent\textcolor{FuncColor}{$\triangleright$\enspace\texttt{Random({\mdseries\slshape S})\index{Random@\texttt{Random}!for a semigroup}
\label{Random:for a semigroup}
}\hfill{\scriptsize (method)}}\\
\textbf{\indent Returns:\ }
A random element.



 This function returns a random element of the semigroup \mbox{\texttt{\mdseries\slshape S}}. If the elements of \mbox{\texttt{\mdseries\slshape S}} have been calculated, then one of these is chosen randomly. Otherwise, if the
data structure for \mbox{\texttt{\mdseries\slshape S}} is known, then a random element of a randomly chosen $\mathcal{R}$\texttt{\symbol{45}}class is returned. If the data structure for \mbox{\texttt{\mdseries\slshape S}} has not been calculated, then a short product (at most \texttt{2 * Length(GeneratorsOfSemigroup(\mbox{\texttt{\mdseries\slshape S}}))}) of generators is returned. }

 }

   
\section{\textcolor{Chapter }{ Properties of elements in a semigroup }}\logpage{[ 11, 4, 0 ]}
\hyperdef{L}{X80EB463F7E5D8920}{}
{
  

\subsection{\textcolor{Chapter }{IndexPeriodOfSemigroupElement}}
\logpage{[ 11, 4, 1 ]}\nobreak
\hyperdef{L}{X869AC4247E2BA4A2}{}
{\noindent\textcolor{FuncColor}{$\triangleright$\enspace\texttt{IndexPeriodOfSemigroupElement({\mdseries\slshape x})\index{IndexPeriodOfSemigroupElement@\texttt{IndexPeriodOfSemigroupElement}}
\label{IndexPeriodOfSemigroupElement}
}\hfill{\scriptsize (operation)}}\\
\textbf{\indent Returns:\ }
A list of two positive integers.



 If \mbox{\texttt{\mdseries\slshape x}} is a semigroup element, then \texttt{IndexPeriodOfSemigroupElement(\mbox{\texttt{\mdseries\slshape x}})} returns the pair \texttt{[m, r]}, where \texttt{m} and \texttt{r} are the least positive integers such that \texttt{\mbox{\texttt{\mdseries\slshape x}}\texttt{\symbol{94}}(m + r) = \mbox{\texttt{\mdseries\slshape x}} \texttt{\symbol{94}} m}. The number \texttt{m} is known as the \emph{index} of \mbox{\texttt{\mdseries\slshape x}}, and the number\texttt{r} is known as the \emph{period} of \mbox{\texttt{\mdseries\slshape x}}. 
\begin{Verbatim}[commandchars=!@|,fontsize=\small,frame=single,label=Example]
  !gapprompt@gap>| !gapinput@x := Transformation([2, 6, 3, 5, 6, 1]);;|
  !gapprompt@gap>| !gapinput@IndexPeriodOfSemigroupElement(x);|
  [ 2, 3 ]
  !gapprompt@gap>| !gapinput@m := IndexPeriodOfSemigroupElement(x)[1];;|
  !gapprompt@gap>| !gapinput@r := IndexPeriodOfSemigroupElement(x)[2];;|
  !gapprompt@gap>| !gapinput@x ^ (m + r) = x ^ m;|
  true
  !gapprompt@gap>| !gapinput@x := PartialPerm([0, 2, 3, 0, 5]);|
  <identity partial perm on [ 2, 3, 5 ]>
  !gapprompt@gap>| !gapinput@IsIdempotent(x);|
  true
  !gapprompt@gap>| !gapinput@IndexPeriodOfSemigroupElement(x);|
  [ 1, 1 ]
\end{Verbatim}
 }

 

\subsection{\textcolor{Chapter }{SmallestIdempotentPower}}
\logpage{[ 11, 4, 2 ]}\nobreak
\hyperdef{L}{X84E1A41F84B70DBB}{}
{\noindent\textcolor{FuncColor}{$\triangleright$\enspace\texttt{SmallestIdempotentPower({\mdseries\slshape x})\index{SmallestIdempotentPower@\texttt{SmallestIdempotentPower}}
\label{SmallestIdempotentPower}
}\hfill{\scriptsize (attribute)}}\\
\textbf{\indent Returns:\ }
A positive integer.



 If \mbox{\texttt{\mdseries\slshape x}} is a semigroup element, then \texttt{SmallestIdempotentPower(\mbox{\texttt{\mdseries\slshape x}})} returns the least positive integer \texttt{n} such that \texttt{\mbox{\texttt{\mdseries\slshape x}}\texttt{\symbol{94}}n} is an idempotent. The smallest idempotent power of \mbox{\texttt{\mdseries\slshape x}} is the least multiple of the period of \mbox{\texttt{\mdseries\slshape x}} that is greater than or equal to the index of \mbox{\texttt{\mdseries\slshape x}}; see \texttt{IndexPeriodOfSemigroupElement} (\ref{IndexPeriodOfSemigroupElement}). 
\begin{Verbatim}[commandchars=!@|,fontsize=\small,frame=single,label=Example]
  !gapprompt@gap>| !gapinput@x := Transformation([4, 1, 4, 5, 1]);|
  Transformation( [ 4, 1, 4, 5, 1 ] )
  !gapprompt@gap>| !gapinput@SmallestIdempotentPower(x);|
  3
  !gapprompt@gap>| !gapinput@ForAll([1 .. 2], i -> not IsIdempotent(x ^ i));|
  true
  !gapprompt@gap>| !gapinput@IsIdempotent(x ^ 3);|
  true
  !gapprompt@gap>| !gapinput@x := Bipartition([[1, 2, -3, -4], [3, -5], [4, -1], [5, -2]]);|
  <block bijection: [ 1, 2, -3, -4 ], [ 3, -5 ], [ 4, -1 ], [ 5, -2 ]>
  !gapprompt@gap>| !gapinput@SmallestIdempotentPower(x);|
  4
  !gapprompt@gap>| !gapinput@ForAll([1 .. 3], i -> not IsIdempotent(x ^ i));|
  true
  !gapprompt@gap>| !gapinput@x := PartialPerm([]);|
  <empty partial perm>
  !gapprompt@gap>| !gapinput@SmallestIdempotentPower(x);|
  1
  !gapprompt@gap>| !gapinput@IsIdempotent(x);|
  true
\end{Verbatim}
 }

 }

   
\section{\textcolor{Chapter }{ Operations for elements in a semigroup }}\logpage{[ 11, 5, 0 ]}
\hyperdef{L}{X7A20EC348515E37B}{}
{
  

\subsection{\textcolor{Chapter }{OneInverseOfSemigroupElement}}
\logpage{[ 11, 5, 1 ]}\nobreak
\hyperdef{L}{X7BDAC0B581B71D0F}{}
{\noindent\textcolor{FuncColor}{$\triangleright$\enspace\texttt{OneInverseOfSemigroupElement({\mdseries\slshape S, x})\index{OneInverseOfSemigroupElement@\texttt{OneInverseOfSemigroupElement}}
\label{OneInverseOfSemigroupElement}
}\hfill{\scriptsize (operation)}}\\
\textbf{\indent Returns:\ }
 One inverse of an element of a semigroup or \texttt{fail}.



 \texttt{OneInverseOfSemigroupElement} returns one inverse of the element \texttt{x} in the semigroup \mbox{\texttt{\mdseries\slshape S}} and returns \texttt{fail} if this element has no inverse in \mbox{\texttt{\mdseries\slshape S}}.

 An \emph{inverse} of an element \mbox{\texttt{\mdseries\slshape x}} in a semigroup \mbox{\texttt{\mdseries\slshape S}} is an element \texttt{y} of \mbox{\texttt{\mdseries\slshape S}} such that \texttt{\mbox{\texttt{\mdseries\slshape x}} * y * \mbox{\texttt{\mdseries\slshape x}} = \mbox{\texttt{\mdseries\slshape x}}} and \texttt{y * \mbox{\texttt{\mdseries\slshape x}} * y = \mbox{\texttt{\mdseries\slshape y}}}. See also \texttt{OnePseudoInverseOfSemigroupElement} (\ref{OnePseudoInverseOfSemigroupElement}). 
\begin{Verbatim}[commandchars=!@|,fontsize=\small,frame=single,label=Example]
  !gapprompt@gap>| !gapinput@S := FullTransformationMonoid(4);|
  <full transformation monoid of degree 4>
  !gapprompt@gap>| !gapinput@s := Transformation([2, 3, 1, 1]);|
  Transformation( [ 2, 3, 1, 1 ] )
  !gapprompt@gap>| !gapinput@OneInverseOfSemigroupElement(S, s);|
  Transformation( [ 3, 1, 2, 2 ] )
  !gapprompt@gap>| !gapinput@e := IdentityTransformation;|
  IdentityTransformation
  !gapprompt@gap>| !gapinput@OneInverseOfSemigroupElement(S, e);|
  IdentityTransformation
  !gapprompt@gap>| !gapinput@F := FreeSemigroup(1);|
  <free semigroup on the generators [ s1 ]>
  !gapprompt@gap>| !gapinput@OneInverseOfSemigroupElement(F, F.1);|
  Error, the semigroup is not finite
\end{Verbatim}
 }

 

\subsection{\textcolor{Chapter }{OnePseudoInverseOfSemigroupElement}}
\logpage{[ 11, 5, 2 ]}\nobreak
\hyperdef{L}{X7EF33B447AB95CCA}{}
{\noindent\textcolor{FuncColor}{$\triangleright$\enspace\texttt{OnePseudoInverseOfSemigroupElement({\mdseries\slshape x})\index{OnePseudoInverseOfSemigroupElement@\texttt{OnePseudoInverseOfSemigroupElement}}
\label{OnePseudoInverseOfSemigroupElement}
}\hfill{\scriptsize (operation)}}\\
\textbf{\indent Returns:\ }
 One pseudo\texttt{\symbol{45}}inverse of an element of a semigroup.



 \texttt{OnePseudoInverseOfSemigroupElement} returns one pseudo\texttt{\symbol{45}}inverse of the element \mbox{\texttt{\mdseries\slshape x}}. This is an element of the same type as \mbox{\texttt{\mdseries\slshape x}} belonging to some semigroup or monoid containing \mbox{\texttt{\mdseries\slshape x}}. 

 An \emph{pseudo\texttt{\symbol{45}}inverse} of an element \mbox{\texttt{\mdseries\slshape x}} is an element \texttt{y} such that \texttt{\mbox{\texttt{\mdseries\slshape x}} * y * \mbox{\texttt{\mdseries\slshape x}} = \mbox{\texttt{\mdseries\slshape x}}}. 

 Methods are installed for this operation for the following types of elements: 
\begin{description}
\item[{transformations}]  \texttt{IsTransformation} (\textbf{Reference: IsTransformation}); 
\item[{partial perms}]  \texttt{IsPartialPerm} (\textbf{Reference: IsPartialPerm}); 
\item[{bipartitions}]  \texttt{IsBipartition} (\ref{IsBipartition}); 
\item[{matrices over finite fields}]  \textsf{GAP} matrices with entries in a finite field; 
\item[{elements of McAlister triple semigroups}]  \texttt{IsMcAlisterTripleSemigroupElement} (\ref{IsMcAlisterTripleSemigroupElement}); 
\item[{elements of regular Rees 0\texttt{\symbol{45}}matrix semigroups over groups}]  \texttt{IsReesZeroMatrixSemigroupElement} (\textbf{Reference: IsReesZeroMatrixSemigroupElement}). 
\end{description}
 See also \texttt{OneInverseOfSemigroupElement} (\ref{OneInverseOfSemigroupElement}). 
\begin{Verbatim}[commandchars=@|A,fontsize=\small,frame=single,label=Example]
  @gapprompt|gap>A @gapinput|s := Transformation([2, 3, 1, 1]);A
  Transformation( [ 2, 3, 1, 1 ] )
  @gapprompt|gap>A @gapinput|OnePseudoInverseOfSemigroupElement(s);A
  Transformation( [ 3, 1, 2, 1 ] )
  @gapprompt|gap>A @gapinput|e := IdentityTransformation;A
  IdentityTransformation
  @gapprompt|gap>A @gapinput|OnePseudoInverseOfSemigroupElement(e);A
  IdentityTransformation
  @gapprompt|gap>A @gapinput|F := FreeSemigroup(1);A
  <free semigroup on the generators [ s1 ]>
  @gapprompt|gap>A @gapinput|OnePseudoInverseOfSemigroupElement(F, F.1);A
  Error, no method found! For debugging hints type ?Recovery from NoMeth\
  odFound
  Error, no 1st choice method found for `OnePseudoInverseOfSemigroupElem\
  ent' on 2 arguments
\end{Verbatim}
 }

 }

   
\section{\textcolor{Chapter }{ Expressing semigroup elements as words in generators }}\logpage{[ 11, 6, 0 ]}
\hyperdef{L}{X81CEB3717E021643}{}
{
  It is possible to express an element of a semigroup as a word in the
generators of that semigroup. This section describes how to accomplish this in \textsf{Semigroups}.

 

\subsection{\textcolor{Chapter }{EvaluateWord}}
\logpage{[ 11, 6, 1 ]}\nobreak
\hyperdef{L}{X799D2F3C866B9AED}{}
{\noindent\textcolor{FuncColor}{$\triangleright$\enspace\texttt{EvaluateWord({\mdseries\slshape gens, w})\index{EvaluateWord@\texttt{EvaluateWord}}
\label{EvaluateWord}
}\hfill{\scriptsize (operation)}}\\
\textbf{\indent Returns:\ }
A semigroup element.



 The argument \mbox{\texttt{\mdseries\slshape gens}} should be a collection of generators of a semigroup and the argument \mbox{\texttt{\mdseries\slshape w}} should be a list of positive integers less than or equal to the length of \mbox{\texttt{\mdseries\slshape gens}}. This operation evaluates the word \mbox{\texttt{\mdseries\slshape w}} in the generators \mbox{\texttt{\mdseries\slshape gens}}. More precisely, \texttt{EvaluateWord(\mbox{\texttt{\mdseries\slshape gens}}, \mbox{\texttt{\mdseries\slshape w}})} returns the equivalent of: 
\begin{Verbatim}[commandchars=!@|,fontsize=\small,frame=single,label=Example]
  Product(List(w, i -> gens[i]));
\end{Verbatim}
 see also \texttt{Factorization} (\ref{Factorization}).

 
\begin{description}
\item[{for elements of a semigroup}]  When \mbox{\texttt{\mdseries\slshape gens}} is a list of elements of a semigroup and \mbox{\texttt{\mdseries\slshape w}} is a list of positive integers less than or equal to the length of \mbox{\texttt{\mdseries\slshape gens}}, this operation returns the product \texttt{gens[w[1]] * gens[w[2]] * .. . * gens[w[n]]} when the length of \mbox{\texttt{\mdseries\slshape w}} is \texttt{n}. 
\item[{for elements of an inverse semigroup}]  When \mbox{\texttt{\mdseries\slshape gens}} is a list of elements with a semigroup inverse and \mbox{\texttt{\mdseries\slshape w}} is a list of non\texttt{\symbol{45}}zero integers whose absolute value does
not exceed the length of \mbox{\texttt{\mdseries\slshape gens}}, this operation returns the product \texttt{gens[AbsInt(w[1])] \texttt{\symbol{94}} SignInt(w[1]) * .. . *
gens[AbsInt(w[n])] \texttt{\symbol{94}} SignInt(w[n])} where \texttt{n} is the length of \mbox{\texttt{\mdseries\slshape w}}. 
\end{description}
 Note that \texttt{EvaluateWord(\mbox{\texttt{\mdseries\slshape gens}}, [])} returns \texttt{One(\mbox{\texttt{\mdseries\slshape gens}})} if \mbox{\texttt{\mdseries\slshape gens}} belongs to the category \texttt{IsMultiplicativeElementWithOne} (\textbf{Reference: IsMultiplicativeElementWithOne}). 
\begin{Verbatim}[commandchars=!@|,fontsize=\small,frame=single,label=Example]
  !gapprompt@gap>| !gapinput@gens := [|
  !gapprompt@>| !gapinput@Transformation([2, 4, 4, 6, 8, 8, 6, 6]),|
  !gapprompt@>| !gapinput@Transformation([2, 7, 4, 1, 4, 6, 5, 2]),|
  !gapprompt@>| !gapinput@Transformation([3, 6, 2, 4, 2, 2, 2, 8]),|
  !gapprompt@>| !gapinput@Transformation([4, 3, 6, 4, 2, 1, 2, 6]),|
  !gapprompt@>| !gapinput@Transformation([4, 5, 1, 3, 8, 5, 8, 2])];;|
  !gapprompt@gap>| !gapinput@S := Semigroup(gens);;|
  !gapprompt@gap>| !gapinput@x := Transformation([1, 4, 6, 1, 7, 2, 7, 6]);;|
  !gapprompt@gap>| !gapinput@word := Factorization(S, x);|
  [ 4, 2 ]
  !gapprompt@gap>| !gapinput@EvaluateWord(gens, word);|
  Transformation( [ 1, 4, 6, 1, 7, 2, 7, 6 ] )
  !gapprompt@gap>| !gapinput@S := SymmetricInverseMonoid(10);;|
  !gapprompt@gap>| !gapinput@x := PartialPerm([2, 6, 7, 0, 0, 9, 0, 1, 0, 5]);|
  [3,7][8,1,2,6,9][10,5]
  !gapprompt@gap>| !gapinput@word := Factorization(S, x);|
  [ -2, -2, -2, -2, -3, -2, -2, -2, -2, -2, 5, 2, 5, 5, 2, 5, 2, 2, 2,
    2, -3, 2, 2, 2, 3, 2, 2, 2, 2, 2, 2, 2, 2, 3, 2, 3, 2 ]
  !gapprompt@gap>| !gapinput@EvaluateWord(GeneratorsOfSemigroup(S), word);|
  [3,7][8,1,2,6,9][10,5]
\end{Verbatim}
 }

 

\subsection{\textcolor{Chapter }{Factorization}}
\logpage{[ 11, 6, 2 ]}\nobreak
\hyperdef{L}{X8357294D7B164106}{}
{\noindent\textcolor{FuncColor}{$\triangleright$\enspace\texttt{Factorization({\mdseries\slshape S, x})\index{Factorization@\texttt{Factorization}}
\label{Factorization}
}\hfill{\scriptsize (operation)}}\\
\textbf{\indent Returns:\ }
A word in the generators.



 
\begin{description}
\item[{for semigroups}]  When \mbox{\texttt{\mdseries\slshape S}} is a semigroup and \mbox{\texttt{\mdseries\slshape x}} belongs to \mbox{\texttt{\mdseries\slshape S}}, \texttt{Factorization} return a word in the generators of \mbox{\texttt{\mdseries\slshape S}} that is equal to \mbox{\texttt{\mdseries\slshape x}}. In this case, a word is a list of positive integers where an entry \texttt{i} corresponds to \texttt{GeneratorsOfSemigroups(S)[i]}. More specifically, 
\begin{Verbatim}[commandchars=!@|,fontsize=\small,frame=single,label=Example]
  EvaluateWord(GeneratorsOfSemigroup(S), Factorization(S, x)) = x;
\end{Verbatim}
 
\item[{for inverse semigroups}]  When \mbox{\texttt{\mdseries\slshape S}} is an inverse semigroup and \mbox{\texttt{\mdseries\slshape x}} belongs to \mbox{\texttt{\mdseries\slshape S}}, \texttt{Factorization} return a word in the generators of \mbox{\texttt{\mdseries\slshape S}} that is equal to \mbox{\texttt{\mdseries\slshape x}}. In this case, a word is a list of non\texttt{\symbol{45}}zero integers where
an entry \texttt{i} corresponds to \texttt{GeneratorsOfSemigroup(S)[i]} and \texttt{\texttt{\symbol{45}}i} corresponds to \texttt{GeneratorsOfSemigroup(S)[i] \texttt{\symbol{94}} \texttt{\symbol{45}}1}. As in the previous case, 
\begin{Verbatim}[commandchars=!@|,fontsize=\small,frame=single,label=Example]
  EvaluateWord(GeneratorsOfSemigroup(S), Factorization(S, x)) = x;
\end{Verbatim}
 
\end{description}
 Note that \texttt{Factorization} does not always return a word of minimum length; see \texttt{MinimalFactorization} (\ref{MinimalFactorization}). 

 See also \texttt{EvaluateWord} (\ref{EvaluateWord}) and \texttt{GeneratorsOfSemigroup} (\textbf{Reference: GeneratorsOfSemigroup}). 
\begin{Verbatim}[commandchars=!@|,fontsize=\small,frame=single,label=Example]
  !gapprompt@gap>| !gapinput@gens := [Transformation([2, 2, 9, 7, 4, 9, 5, 5, 4, 8]),|
  !gapprompt@>| !gapinput@            Transformation([4, 10, 5, 6, 4, 1, 2, 7, 1, 2])];;|
  !gapprompt@gap>| !gapinput@S := Semigroup(gens);;|
  !gapprompt@gap>| !gapinput@x := Transformation([1, 10, 2, 10, 1, 2, 7, 10, 2, 7]);;|
  !gapprompt@gap>| !gapinput@word := Factorization(S, x);|
  [ 2, 2, 1, 2 ]
  !gapprompt@gap>| !gapinput@EvaluateWord(gens, word);|
  Transformation( [ 1, 10, 2, 10, 1, 2, 7, 10, 2, 7 ] )
  !gapprompt@gap>| !gapinput@S := SymmetricInverseMonoid(8);|
  <symmetric inverse monoid of degree 8>
  !gapprompt@gap>| !gapinput@x := PartialPerm([1, 2, 3, 4, 5, 8], [7, 1, 4, 3, 2, 6]);|
  [5,2,1,7][8,6](3,4)
  !gapprompt@gap>| !gapinput@word := Factorization(S, x);|
  [ -2, -2, -2, -2, -2, -2, 2, 4, 4, 2, 3, 2, -3, -2, -2, 3, 2, -3, -2,
    -2, 4, -3, -4, 2, 2, 3, -2, -3, 4, -3, -4, 2, 2, 3, -2, -3, 2, 2,
    3, -2, -3, 2, 2, 3, -2, -3, 4, -3, -4, 3, 2, -3, -2, -2, 3, 2, -3,
    -2, -2, 4, 3, -4, 3, 2, -3, -2, -2, 3, 2, -3, -2, -2, 3, 2, 2, 3,
    2, 2, 2, 2 ]
  !gapprompt@gap>| !gapinput@EvaluateWord(GeneratorsOfSemigroup(S), word);|
  [5,2,1,7][8,6](3,4)
  !gapprompt@gap>| !gapinput@S := DualSymmetricInverseMonoid(6);;|
  !gapprompt@gap>| !gapinput@x := S.1 * S.2 * S.3 * S.2 * S.1;|
  <block bijection: [ 1, 6, -4 ], [ 2, -2, -3 ], [ 3, -5 ], [ 4, -6 ],
   [ 5, -1 ]>
  !gapprompt@gap>| !gapinput@word := Factorization(S, x);|
  [ -2, -2, -2, -2, -2, 4, 2 ]
  !gapprompt@gap>| !gapinput@EvaluateWord(GeneratorsOfSemigroup(S), word);|
  <block bijection: [ 1, 6, -4 ], [ 2, -2, -3 ], [ 3, -5 ], [ 4, -6 ],
   [ 5, -1 ]>
\end{Verbatim}
 }

 

\subsection{\textcolor{Chapter }{MinimalFactorization}}
\logpage{[ 11, 6, 3 ]}\nobreak
\hyperdef{L}{X83A4D71382C5B6C3}{}
{\noindent\textcolor{FuncColor}{$\triangleright$\enspace\texttt{MinimalFactorization({\mdseries\slshape S, x})\index{MinimalFactorization@\texttt{MinimalFactorization}}
\label{MinimalFactorization}
}\hfill{\scriptsize (operation)}}\\
\textbf{\indent Returns:\ }
A minimal word in the generators.



 This operation returns a minimal length word in the generators of the
semigroup \mbox{\texttt{\mdseries\slshape S}} that equals the element \mbox{\texttt{\mdseries\slshape x}}. In this case, a word is a list of positive integers where an entry \texttt{i} corresponds to \texttt{GeneratorsOfSemigroups(\mbox{\texttt{\mdseries\slshape S}})[i]}. More specifically, 
\begin{Verbatim}[commandchars=!@|,fontsize=\small,frame=single,label=Example]
  EvaluateWord(GeneratorsOfSemigroup(S), MinimalFactorization(S, x)) = x;
\end{Verbatim}
 

 \texttt{MinimalFactorization} involves exhaustively enumerating \mbox{\texttt{\mdseries\slshape S}} until the element \mbox{\texttt{\mdseries\slshape x}} is found, and so \texttt{MinimalFactorization} may be less efficient than \texttt{Factorization} (\ref{Factorization}) for some semigroups. 

 Unlike \texttt{Factorization} (\ref{Factorization}) this operation does not distinguish between semigroups and inverse semigroups.
See also \texttt{EvaluateWord} (\ref{EvaluateWord}) and \texttt{GeneratorsOfSemigroup} (\textbf{Reference: GeneratorsOfSemigroup}). 
\begin{Verbatim}[commandchars=!@|,fontsize=\small,frame=single,label=Example]
  !gapprompt@gap>| !gapinput@S := Semigroup(Transformation([2, 2, 9, 7, 4, 9, 5, 5, 4, 8]),|
  !gapprompt@>| !gapinput@                  Transformation([4, 10, 5, 6, 4, 1, 2, 7, 1, 2]));|
  <transformation semigroup of degree 10 with 2 generators>
  !gapprompt@gap>| !gapinput@x := Transformation([8, 8, 2, 2, 9, 2, 8, 8, 9, 9]);|
  Transformation( [ 8, 8, 2, 2, 9, 2, 8, 8, 9, 9 ] )
  !gapprompt@gap>| !gapinput@Factorization(S, x);|
  [ 1, 2, 1, 1, 1, 2, 1, 1, 1, 2, 1, 1, 1, 2, 1 ]
  !gapprompt@gap>| !gapinput@MinimalFactorization(S, x);|
  [ 1, 2, 1, 1, 1, 1, 2, 2, 1 ]
\end{Verbatim}
 }

 

\subsection{\textcolor{Chapter }{NonTrivialFactorization}}
\logpage{[ 11, 6, 4 ]}\nobreak
\hyperdef{L}{X86261F4682DC9842}{}
{\noindent\textcolor{FuncColor}{$\triangleright$\enspace\texttt{NonTrivialFactorization({\mdseries\slshape S, x})\index{NonTrivialFactorization@\texttt{NonTrivialFactorization}}
\label{NonTrivialFactorization}
}\hfill{\scriptsize (operation)}}\\
\textbf{\indent Returns:\ }
A non\texttt{\symbol{45}}trivial word in the generators, or \texttt{fail}.



 When \mbox{\texttt{\mdseries\slshape S}} is a semigroup and \mbox{\texttt{\mdseries\slshape x}} belongs to \mbox{\texttt{\mdseries\slshape S}}, this operation returns a non\texttt{\symbol{45}}trivial word in the
generators of the semigroup \mbox{\texttt{\mdseries\slshape S}} that equals \mbox{\texttt{\mdseries\slshape x}}, if one exists. The definition of a word in the generators is the same as
given in \texttt{Factorization} (\ref{Factorization}) for semigroups and inverse semigroups. A word is
non\texttt{\symbol{45}}trivial if it has length two or more. 

 If no non\texttt{\symbol{45}}trivial word for \mbox{\texttt{\mdseries\slshape x}} exists, then \mbox{\texttt{\mdseries\slshape x}} is an indecomposable element of \mbox{\texttt{\mdseries\slshape S}} and this operation returns \texttt{fail}; see \texttt{IndecomposableElements} (\ref{IndecomposableElements}). 

 When \mbox{\texttt{\mdseries\slshape x}} does not belong to \texttt{GeneratorsOfSemigroup(\mbox{\texttt{\mdseries\slshape S}})}, any factorization of \mbox{\texttt{\mdseries\slshape x}} is non\texttt{\symbol{45}}trivial. In this case, \texttt{NonTrivialFactorization} returns the same word as \texttt{Factorization} (\ref{Factorization}). 

 See also \texttt{EvaluateWord} (\ref{EvaluateWord}) and \texttt{GeneratorsOfSemigroup} (\textbf{Reference: GeneratorsOfSemigroup}). 
\begin{Verbatim}[commandchars=!@|,fontsize=\small,frame=single,label=Example]
  !gapprompt@gap>| !gapinput@x := Transformation([5, 4, 2, 1, 3]);;|
  !gapprompt@gap>| !gapinput@y := Transformation([4, 4, 2, 4, 1]);;|
  !gapprompt@gap>| !gapinput@S := Semigroup([x, y]);|
  <transformation semigroup of degree 5 with 2 generators>
  !gapprompt@gap>| !gapinput@NonTrivialFactorization(S, x * y);|
  [ 1, 2 ]
  !gapprompt@gap>| !gapinput@Factorization(S, x);|
  [ 1 ]
  !gapprompt@gap>| !gapinput@NonTrivialFactorization(S, x);|
  [ 1, 1, 1, 1, 1, 1 ]
  !gapprompt@gap>| !gapinput@Factorization(S, y);|
  [ 2 ]
  !gapprompt@gap>| !gapinput@NonTrivialFactorization(S, y);|
  [ 2, 1, 1, 1, 1, 1 ]
  !gapprompt@gap>| !gapinput@z := PartialPerm([2]);;|
  !gapprompt@gap>| !gapinput@S := Semigroup(z);|
  <commutative partial perm semigroup of rank 1 with 1 generator>
  !gapprompt@gap>| !gapinput@NonTrivialFactorization(S, z);|
  fail
\end{Verbatim}
 }

 }

   
\section{\textcolor{Chapter }{ Generating sets }}\logpage{[ 11, 7, 0 ]}
\hyperdef{L}{X7E4AA1437A6C7B40}{}
{
  

\subsection{\textcolor{Chapter }{Generators}}
\logpage{[ 11, 7, 1 ]}\nobreak
\hyperdef{L}{X7BD5B55C802805B4}{}
{\noindent\textcolor{FuncColor}{$\triangleright$\enspace\texttt{Generators({\mdseries\slshape S})\index{Generators@\texttt{Generators}}
\label{Generators}
}\hfill{\scriptsize (attribute)}}\\
\textbf{\indent Returns:\ }
A list of generators.



 \texttt{Generators} returns a generating set that can be used to define the semigroup \mbox{\texttt{\mdseries\slshape S}}. The generators of a monoid or inverse semigroup \mbox{\texttt{\mdseries\slshape S}}, say, can be defined in several ways, for example, including or excluding the
identity element, including or not the inverses of the generators. \texttt{Generators} uses the definition that returns the least number of generators. If no
generating set for \mbox{\texttt{\mdseries\slshape S}} is known, then \texttt{GeneratorsOfSemigroup} is used by default.

 
\begin{description}
\item[{for a group}] \texttt{Generators(\mbox{\texttt{\mdseries\slshape S}})} is a synonym for \texttt{GeneratorsOfGroup} (\textbf{Reference: GeneratorsOfGroup}). 
\item[{for an ideal of semigroup}] \texttt{Generators(\mbox{\texttt{\mdseries\slshape S}})} is a synonym for \texttt{GeneratorsOfSemigroupIdeal} (\ref{GeneratorsOfSemigroupIdeal}). 
\item[{for a semigroup}] \texttt{Generators(\mbox{\texttt{\mdseries\slshape S}})} is a synonym for \texttt{GeneratorsOfSemigroup} (\textbf{Reference: GeneratorsOfSemigroup}). 
\item[{for a monoid}] \texttt{Generators(\mbox{\texttt{\mdseries\slshape S}})} is a synonym for \texttt{GeneratorsOfMonoid} (\textbf{Reference: GeneratorsOfMonoid}). 
\item[{for an inverse semigroup}] \texttt{Generators(\mbox{\texttt{\mdseries\slshape S}})} is a synonym for \texttt{GeneratorsOfInverseSemigroup} (\textbf{Reference: GeneratorsOfInverseSemigroup}). 
\item[{for an inverse monoid}] \texttt{Generators(\mbox{\texttt{\mdseries\slshape S}})} is a synonym for \texttt{GeneratorsOfInverseMonoid} (\textbf{Reference: GeneratorsOfInverseMonoid}). 
\end{description}
 
\begin{Verbatim}[commandchars=!@|,fontsize=\small,frame=single,label=Example]
  !gapprompt@gap>| !gapinput@M := Monoid([|
  !gapprompt@>| !gapinput@Transformation([1, 4, 6, 2, 5, 3, 7, 8, 9, 9]),|
  !gapprompt@>| !gapinput@Transformation([6, 3, 2, 7, 5, 1, 8, 8, 9, 9])]);;|
  !gapprompt@gap>| !gapinput@GeneratorsOfSemigroup(M);|
  [ IdentityTransformation,
    Transformation( [ 1, 4, 6, 2, 5, 3, 7, 8, 9, 9 ] ),
    Transformation( [ 6, 3, 2, 7, 5, 1, 8, 8, 9, 9 ] ) ]
  !gapprompt@gap>| !gapinput@GeneratorsOfMonoid(M);|
  [ Transformation( [ 1, 4, 6, 2, 5, 3, 7, 8, 9, 9 ] ),
    Transformation( [ 6, 3, 2, 7, 5, 1, 8, 8, 9, 9 ] ) ]
  !gapprompt@gap>| !gapinput@Generators(M);|
  [ Transformation( [ 1, 4, 6, 2, 5, 3, 7, 8, 9, 9 ] ),
    Transformation( [ 6, 3, 2, 7, 5, 1, 8, 8, 9, 9 ] ) ]
  !gapprompt@gap>| !gapinput@S := Semigroup(Generators(M));;|
  !gapprompt@gap>| !gapinput@Generators(S);|
  [ Transformation( [ 1, 4, 6, 2, 5, 3, 7, 8, 9, 9 ] ),
    Transformation( [ 6, 3, 2, 7, 5, 1, 8, 8, 9, 9 ] ) ]
  !gapprompt@gap>| !gapinput@GeneratorsOfSemigroup(S);|
  [ Transformation( [ 1, 4, 6, 2, 5, 3, 7, 8, 9, 9 ] ),
    Transformation( [ 6, 3, 2, 7, 5, 1, 8, 8, 9, 9 ] ) ]
\end{Verbatim}
 }

 

\subsection{\textcolor{Chapter }{SmallGeneratingSet}}
\logpage{[ 11, 7, 2 ]}\nobreak
\hyperdef{L}{X814DBABC878D5232}{}
{\noindent\textcolor{FuncColor}{$\triangleright$\enspace\texttt{SmallGeneratingSet({\mdseries\slshape coll})\index{SmallGeneratingSet@\texttt{SmallGeneratingSet}}
\label{SmallGeneratingSet}
}\hfill{\scriptsize (attribute)}}\\
\noindent\textcolor{FuncColor}{$\triangleright$\enspace\texttt{SmallSemigroupGeneratingSet({\mdseries\slshape coll})\index{SmallSemigroupGeneratingSet@\texttt{SmallSemigroupGeneratingSet}}
\label{SmallSemigroupGeneratingSet}
}\hfill{\scriptsize (attribute)}}\\
\noindent\textcolor{FuncColor}{$\triangleright$\enspace\texttt{SmallMonoidGeneratingSet({\mdseries\slshape coll})\index{SmallMonoidGeneratingSet@\texttt{SmallMonoidGeneratingSet}}
\label{SmallMonoidGeneratingSet}
}\hfill{\scriptsize (attribute)}}\\
\noindent\textcolor{FuncColor}{$\triangleright$\enspace\texttt{SmallInverseSemigroupGeneratingSet({\mdseries\slshape coll})\index{SmallInverseSemigroupGeneratingSet@\texttt{SmallInverseSemigroupGeneratingSet}}
\label{SmallInverseSemigroupGeneratingSet}
}\hfill{\scriptsize (attribute)}}\\
\noindent\textcolor{FuncColor}{$\triangleright$\enspace\texttt{SmallInverseMonoidGeneratingSet({\mdseries\slshape coll})\index{SmallInverseMonoidGeneratingSet@\texttt{SmallInverseMonoidGeneratingSet}}
\label{SmallInverseMonoidGeneratingSet}
}\hfill{\scriptsize (attribute)}}\\
\textbf{\indent Returns:\ }
A small generating set for a semigroup.



 The attributes \texttt{SmallXGeneratingSet} return a relatively small generating subset of the collection of elements \mbox{\texttt{\mdseries\slshape coll}}, which can also be a semigroup. The returned value of \texttt{SmallXGeneratingSet}, where applicable, has the property that 
\begin{Verbatim}[commandchars=!@|,fontsize=\small,frame=single,label=Example]
        X(SmallXGeneratingSet(coll)) = X(coll);
\end{Verbatim}
 where \texttt{X} is any of \texttt{Semigroup} (\textbf{Reference: Semigroup}), \texttt{Monoid} (\textbf{Reference: Monoid}), \texttt{InverseSemigroup} (\textbf{Reference: InverseSemigroup}), or \texttt{InverseMonoid} (\textbf{Reference: InverseMonoid}).

 If the number of generators for \mbox{\texttt{\mdseries\slshape S}} is already relatively small, then these functions will often return the
original generating set. These functions may return different results in
different \textsf{GAP} sessions.

 \texttt{SmallGeneratingSet} returns the smallest of the returned values of \texttt{SmallXGeneratingSet} which is applicable to \mbox{\texttt{\mdseries\slshape coll}}; see \texttt{Generators} (\ref{Generators}).

 As neither irredundancy, nor minimal length are proven, these functions
usually return an answer much more quickly than \texttt{IrredundantGeneratingSubset} (\ref{IrredundantGeneratingSubset}). These functions can be used whenever a small generating set is desired which
does not necessarily needs to be minimal. 
\begin{Verbatim}[commandchars=!@|,fontsize=\small,frame=single,label=Example]
  !gapprompt@gap>| !gapinput@S := Semigroup([|
  !gapprompt@>| !gapinput@Transformation([1, 2, 3, 2, 4]),|
  !gapprompt@>| !gapinput@Transformation([1, 5, 4, 3, 2]),|
  !gapprompt@>| !gapinput@Transformation([2, 1, 4, 2, 2]),|
  !gapprompt@>| !gapinput@Transformation([2, 4, 4, 2, 1]),|
  !gapprompt@>| !gapinput@Transformation([3, 1, 4, 3, 2]),|
  !gapprompt@>| !gapinput@Transformation([3, 2, 3, 4, 1]),|
  !gapprompt@>| !gapinput@Transformation([4, 4, 3, 3, 5]),|
  !gapprompt@>| !gapinput@Transformation([5, 1, 5, 5, 3]),|
  !gapprompt@>| !gapinput@Transformation([5, 4, 3, 5, 2]),|
  !gapprompt@>| !gapinput@Transformation([5, 5, 4, 5, 5])]);;|
  !gapprompt@gap>| !gapinput@SmallGeneratingSet(S);|
  [ Transformation( [ 1, 5, 4, 3, 2 ] ), Transformation( [ 3, 2, 3, 4, 1 ] ),
    Transformation( [ 5, 4, 3, 5, 2 ] ), Transformation( [ 1, 2, 3, 2, 4 ] ),
    Transformation( [ 4, 4, 3, 3, 5 ] ) ]
  !gapprompt@gap>| !gapinput@S := RandomInverseMonoid(IsPartialPermMonoid, 10000, 10);;|
  !gapprompt@gap>| !gapinput@SmallGeneratingSet(S);|
  [ [ 1 .. 10 ] -> [ 3, 2, 4, 5, 6, 1, 7, 10, 9, 8 ],
    [ 1 .. 10 ] -> [ 5, 10, 8, 9, 3, 2, 4, 7, 6, 1 ],
    [ 1, 3, 4, 5, 6, 7, 8, 9, 10 ] -> [ 1, 6, 4, 8, 2, 10, 7, 3, 9 ] ]
  !gapprompt@gap>| !gapinput@M := MathieuGroup(24);;|
  !gapprompt@gap>| !gapinput@mat := List([1 .. 1000], x -> Random(M));;|
  !gapprompt@gap>| !gapinput@Append(mat, [1 .. 1000] * 0);|
  !gapprompt@gap>| !gapinput@mat := List([1 .. 138], x -> List([1 .. 57], x -> Random(mat)));;|
  !gapprompt@gap>| !gapinput@R := ReesZeroMatrixSemigroup(M, mat);;|
  !gapprompt@gap>| !gapinput@U := Semigroup(List([1 .. 200], x -> Random(R)));|
  <subsemigroup of 57x138 Rees 0-matrix semigroup with 100 generators>
  !gapprompt@gap>| !gapinput@Length(SmallGeneratingSet(U));|
  84
  !gapprompt@gap>| !gapinput@S := RandomSemigroup(IsBipartitionSemigroup, 100, 4);|
  <bipartition semigroup of degree 4 with 96 generators>
  !gapprompt@gap>| !gapinput@Length(SmallGeneratingSet(S));|
  13
\end{Verbatim}
 }

 

\subsection{\textcolor{Chapter }{IrredundantGeneratingSubset}}
\logpage{[ 11, 7, 3 ]}\nobreak
\hyperdef{L}{X7F88DA9487720D2B}{}
{\noindent\textcolor{FuncColor}{$\triangleright$\enspace\texttt{IrredundantGeneratingSubset({\mdseries\slshape coll})\index{IrredundantGeneratingSubset@\texttt{IrredundantGeneratingSubset}}
\label{IrredundantGeneratingSubset}
}\hfill{\scriptsize (operation)}}\\
\textbf{\indent Returns:\ }
 A list of irredundant generators. 



 If \mbox{\texttt{\mdseries\slshape coll}} is a collection of elements of a semigroup, then this function returns a
subset \texttt{U} of \mbox{\texttt{\mdseries\slshape coll}} such that no element of \texttt{U} is generated by the other elements of \texttt{U}.

 
\begin{Verbatim}[commandchars=!@|,fontsize=\small,frame=single,label=Example]
  !gapprompt@gap>| !gapinput@S := Semigroup([|
  !gapprompt@>| !gapinput@Transformation([5, 1, 4, 6, 2, 3]),|
  !gapprompt@>| !gapinput@Transformation([1, 2, 3, 4, 5, 6]),|
  !gapprompt@>| !gapinput@Transformation([4, 6, 3, 4, 2, 5]),|
  !gapprompt@>| !gapinput@Transformation([5, 4, 6, 3, 1, 3]),|
  !gapprompt@>| !gapinput@Transformation([2, 2, 6, 5, 4, 3]),|
  !gapprompt@>| !gapinput@Transformation([3, 5, 5, 1, 2, 4]),|
  !gapprompt@>| !gapinput@Transformation([6, 5, 1, 3, 3, 4]),|
  !gapprompt@>| !gapinput@Transformation([1, 3, 4, 3, 2, 1])]);;|
  !gapprompt@gap>| !gapinput@IrredundantGeneratingSubset(S);|
  [ Transformation( [ 1, 3, 4, 3, 2, 1 ] ),
    Transformation( [ 2, 2, 6, 5, 4, 3 ] ),
    Transformation( [ 3, 5, 5, 1, 2, 4 ] ),
    Transformation( [ 5, 1, 4, 6, 2, 3 ] ),
    Transformation( [ 5, 4, 6, 3, 1, 3 ] ),
    Transformation( [ 6, 5, 1, 3, 3, 4 ] ) ]
  !gapprompt@gap>| !gapinput@S := RandomInverseMonoid(IsPartialPermMonoid, 1000, 10);|
  <inverse partial perm monoid of degree 10 with 1000 generators>
  !gapprompt@gap>| !gapinput@SmallGeneratingSet(S);|
  [ [ 1 .. 10 ] -> [ 6, 5, 1, 9, 8, 3, 10, 4, 7, 2 ],
    [ 1 .. 10 ] -> [ 1, 4, 6, 2, 8, 5, 7, 10, 3, 9 ],
    [ 1, 2, 3, 4, 6, 7, 8, 9 ] -> [ 7, 5, 10, 1, 8, 4, 9, 6 ]
    [ 1 .. 9 ] -> [ 4, 3, 5, 7, 10, 9, 1, 6, 8 ] ]
  !gapprompt@gap>| !gapinput@IrredundantGeneratingSubset(last);|
  [ [ 1 .. 9 ] -> [ 4, 3, 5, 7, 10, 9, 1, 6, 8 ],
    [ 1 .. 10 ] -> [ 1, 4, 6, 2, 8, 5, 7, 10, 3, 9 ],
    [ 1 .. 10 ] -> [ 6, 5, 1, 9, 8, 3, 10, 4, 7, 2 ] ]
  !gapprompt@gap>| !gapinput@S := RandomSemigroup(IsBipartitionSemigroup, 1000, 4);|
  <bipartition semigroup of degree 4 with 749 generators>
  !gapprompt@gap>| !gapinput@SmallGeneratingSet(S);|
  [ <bipartition: [ 1, -3 ], [ 2, -2 ], [ 3, -1 ], [ 4, -4 ]>,
    <bipartition: [ 1, 3, -2 ], [ 2, -1, -3 ], [ 4, -4 ]>,
    <bipartition: [ 1, -4 ], [ 2, 4, -1, -3 ], [ 3, -2 ]>,
    <bipartition: [ 1, -1, -3 ], [ 2, -4 ], [ 3, 4, -2 ]>,
    <bipartition: [ 1, -2, -4 ], [ 2 ], [ 3, -3 ], [ 4, -1 ]>,
    <bipartition: [ 1, -2 ], [ 2, -1, -3 ], [ 3, 4, -4 ]>,
    <bipartition: [ 1, 3, -1 ], [ 2, -3 ], [ 4, -2, -4 ]>,
    <bipartition: [ 1, -1 ], [ 2, 4, -4 ], [ 3, -2, -3 ]>,
    <bipartition: [ 1, 3, -1 ], [ 2, -2 ], [ 4, -3, -4 ]>,
    <bipartition: [ 1, 2, -2 ], [ 3, -1, -4 ], [ 4, -3 ]>,
    <bipartition: [ 1, -2, -3 ], [ 2, -4 ], [ 3 ], [ 4, -1 ]>,
    <bipartition: [ 1, -1 ], [ 2, 4, -3 ], [ 3, -2 ], [ -4 ]>,
    <bipartition: [ 1, -3 ], [ 2, -1 ], [ 3, 4, -4 ], [ -2 ]>,
    <bipartition: [ 1, 2, -4 ], [ 3, -1 ], [ 4, -2 ], [ -3 ]>,
    <bipartition: [ 1, -3 ], [ 2, -4 ], [ 3, -1, -2 ], [ 4 ]> ]
  !gapprompt@gap>| !gapinput@IrredundantGeneratingSubset(last);|
  [ <bipartition: [ 1, 2, -4 ], [ 3, -1 ], [ 4, -2 ], [ -3 ]>,
    <bipartition: [ 1, 3, -1 ], [ 2, -2 ], [ 4, -3, -4 ]>,
    <bipartition: [ 1, 3, -2 ], [ 2, -1, -3 ], [ 4, -4 ]>,
    <bipartition: [ 1, -1 ], [ 2, 4, -3 ], [ 3, -2 ], [ -4 ]>,
    <bipartition: [ 1, -3 ], [ 2, -1 ], [ 3, 4, -4 ], [ -2 ]>,
    <bipartition: [ 1, -3 ], [ 2, -2 ], [ 3, -1 ], [ 4, -4 ]>,
    <bipartition: [ 1, -3 ], [ 2, -4 ], [ 3, -1, -2 ], [ 4 ]>,
    <bipartition: [ 1, -2, -3 ], [ 2, -4 ], [ 3 ], [ 4, -1 ]>,
    <bipartition: [ 1, -2, -4 ], [ 2 ], [ 3, -3 ], [ 4, -1 ]> ]
\end{Verbatim}
 }

 

\subsection{\textcolor{Chapter }{MinimalSemigroupGeneratingSet}}
\logpage{[ 11, 7, 4 ]}\nobreak
\hyperdef{L}{X8409DBED7996D495}{}
{\noindent\textcolor{FuncColor}{$\triangleright$\enspace\texttt{MinimalSemigroupGeneratingSet({\mdseries\slshape S})\index{MinimalSemigroupGeneratingSet@\texttt{MinimalSemigroupGeneratingSet}}
\label{MinimalSemigroupGeneratingSet}
}\hfill{\scriptsize (attribute)}}\\
\noindent\textcolor{FuncColor}{$\triangleright$\enspace\texttt{MinimalMonoidGeneratingSet({\mdseries\slshape S})\index{MinimalMonoidGeneratingSet@\texttt{MinimalMonoidGeneratingSet}}
\label{MinimalMonoidGeneratingSet}
}\hfill{\scriptsize (attribute)}}\\
\noindent\textcolor{FuncColor}{$\triangleright$\enspace\texttt{MinimalInverseSemigroupGeneratingSet({\mdseries\slshape S})\index{MinimalInverseSemigroupGeneratingSet@\texttt{Minimal}\-\texttt{Inverse}\-\texttt{Semigroup}\-\texttt{GeneratingSet}}
\label{MinimalInverseSemigroupGeneratingSet}
}\hfill{\scriptsize (attribute)}}\\
\noindent\textcolor{FuncColor}{$\triangleright$\enspace\texttt{MinimalInverseMonoidGeneratingSet({\mdseries\slshape S})\index{MinimalInverseMonoidGeneratingSet@\texttt{MinimalInverseMonoidGeneratingSet}}
\label{MinimalInverseMonoidGeneratingSet}
}\hfill{\scriptsize (attribute)}}\\
\textbf{\indent Returns:\ }
A minimal generating set for a semigroup.



 The attribute \texttt{MinimalXGeneratingSet} returns a minimal generating set for the semigroup \mbox{\texttt{\mdseries\slshape S}}, with respect to length. The returned value of \texttt{MinimalXGeneratingSet}, where applicable, is a minimal\texttt{\symbol{45}}length list of elements of \mbox{\texttt{\mdseries\slshape S}} with the property that 
\begin{Verbatim}[commandchars=!@|,fontsize=\small,frame=single,label=Example]
        X(MinimalXGeneratingSet(S)) = S;
      
\end{Verbatim}
 where \texttt{X} is one of \texttt{Semigroup} (\textbf{Reference: Semigroup}), \texttt{Monoid} (\textbf{Reference: Monoid}), \texttt{InverseSemigroup} (\textbf{Reference: InverseSemigroup}), or \texttt{InverseMonoid} (\textbf{Reference: InverseMonoid}).

 For many types of semigroup, it is not currently possible to find a \texttt{MinimalXGeneratingSet} with the \textsf{Semigroups} package. 

 See also \texttt{SmallGeneratingSet} (\ref{SmallGeneratingSet}) and \texttt{IrredundantGeneratingSubset} (\ref{IrredundantGeneratingSubset}). 
\begin{Verbatim}[commandchars=!@|,fontsize=\small,frame=single,label=Example]
  !gapprompt@gap>| !gapinput@S := MonogenicSemigroup(3, 6);;|
  !gapprompt@gap>| !gapinput@MinimalSemigroupGeneratingSet(S);|
  [ Transformation( [ 2, 3, 4, 5, 6, 1, 6, 7, 8 ] ) ]
  !gapprompt@gap>| !gapinput@S := FullTransformationMonoid(4);;|
  !gapprompt@gap>| !gapinput@MinimalSemigroupGeneratingSet(S);|
  [ Transformation( [ 1, 4, 2, 3 ] ), Transformation( [ 4, 3, 1, 2 ] ),
    Transformation( [ 1, 2, 3, 1 ] ) ]
  !gapprompt@gap>| !gapinput@S := Monoid([|
  !gapprompt@>| !gapinput@ PartialPerm([2, 3, 4, 5, 1, 0, 6, 7]),|
  !gapprompt@>| !gapinput@ PartialPerm([3, 4, 5, 1, 2, 0, 0, 6])]);|
  <partial perm monoid of rank 8 with 2 generators>
  !gapprompt@gap>| !gapinput@IsMonogenicMonoid(S);|
  true
  !gapprompt@gap>| !gapinput@MinimalMonoidGeneratingSet(S);|
  [ [8,7,6](1,2,3,4,5) ]
\end{Verbatim}
 }

 

\subsection{\textcolor{Chapter }{GeneratorsSmallest (for a semigroup)}}
\logpage{[ 11, 7, 5 ]}\nobreak
\hyperdef{L}{X82B02F0887AD1B78}{}
{\noindent\textcolor{FuncColor}{$\triangleright$\enspace\texttt{GeneratorsSmallest({\mdseries\slshape S})\index{GeneratorsSmallest@\texttt{GeneratorsSmallest}!for a semigroup}
\label{GeneratorsSmallest:for a semigroup}
}\hfill{\scriptsize (attribute)}}\\
\textbf{\indent Returns:\ }
A set of elements.



 For a semigroup \mbox{\texttt{\mdseries\slshape S}}, \texttt{GeneratorsSmallest} returns the lexicographically least set of elements \texttt{X} such that \texttt{X} generates \mbox{\texttt{\mdseries\slshape S}} as a semigroup, and such that \texttt{X} is lexicographically ordered and has the property that each \texttt{X[i]} is not generated by \texttt{X[1], X[2], ..., X[i\texttt{\symbol{45}}1]}. 

 It can be difficult to find the set of generators \texttt{X}, and it might contain a substantial proportion of the elements of \mbox{\texttt{\mdseries\slshape S}}. 

 Two semigroups have the same set of elements if and only if their smallest
generating sets are equal. However, due to the complexity of determining the \texttt{GeneratorsSmallest}, this is not the method used by the \textsf{Semigroups} package when comparing semigroups. 
\begin{Verbatim}[commandchars=!@|,fontsize=\small,frame=single,label=Example]
  !gapprompt@gap>| !gapinput@S := Monoid([|
  !gapprompt@>| !gapinput@ Transformation([1, 3, 4, 1]),|
  !gapprompt@>| !gapinput@ Transformation([2, 4, 1, 2]),|
  !gapprompt@>| !gapinput@ Transformation([3, 1, 1, 3]),|
  !gapprompt@>| !gapinput@ Transformation([3, 3, 4, 1])]);|
  <transformation monoid of degree 4 with 4 generators>
  !gapprompt@gap>| !gapinput@GeneratorsSmallest(S);|
  [ Transformation( [ 1, 1, 1, 1 ] ), Transformation( [ 1, 1, 1, 2 ] ),
    Transformation( [ 1, 1, 1, 3 ] ), Transformation( [ 1, 1, 1 ] ),
    Transformation( [ 1, 1, 2, 1 ] ), Transformation( [ 1, 1, 2, 2 ] ),
    Transformation( [ 1, 1, 3, 1 ] ), Transformation( [ 1, 1, 3, 3 ] ),
    Transformation( [ 1, 1 ] ), Transformation( [ 1, 1, 4, 1 ] ),
    Transformation( [ 1, 2, 1, 1 ] ), Transformation( [ 1, 2, 2, 1 ] ),
    IdentityTransformation, Transformation( [ 1, 3, 1, 1 ] ),
    Transformation( [ 1, 3, 4, 1 ] ), Transformation( [ 2, 1, 1, 2 ] ),
    Transformation( [ 2, 2, 2 ] ), Transformation( [ 2, 4, 1, 2 ] ),
    Transformation( [ 3, 3, 3 ] ), Transformation( [ 3, 3, 4, 1 ] ) ]
  !gapprompt@gap>| !gapinput@T := Semigroup(Bipartition([[1, 2, 3], [4, -1], [-2], [-3], [-4]]),|
  !gapprompt@>| !gapinput@                  Bipartition([[1, -3, -4], [2, 3, 4, -2], [-1]]),|
  !gapprompt@>| !gapinput@                  Bipartition([[1, 2, 3, 4, -2], [-1, -4], [-3]]),|
  !gapprompt@>| !gapinput@                  Bipartition([[1, 2, 3, 4], [-1], [-2], [-3, -4]]),|
  !gapprompt@>| !gapinput@                  Bipartition([[1, 2, -1, -2], [3, 4, -3], [-4]]));|
  <bipartition semigroup of degree 4 with 5 generators>
  !gapprompt@gap>| !gapinput@GeneratorsSmallest(T);|
  [ <bipartition: [ 1, 2, 3, 4, -1, -2, -3 ], [ -4 ]>,
    <bipartition: [ 1, 2, 3, 4, -1, -2 ], [ -3 ], [ -4 ]>,
    <bipartition: [ 1, 2, 3, 4, -1 ], [ -2 ], [ -3 ], [ -4 ]>,
    <bipartition: [ 1, 2, 3, 4, -2, -3, -4 ], [ -1 ]>,
    <bipartition: [ 1, 2, 3, 4, -2 ], [ -1, -4 ], [ -3 ]>,
    <bipartition: [ 1, 2, 3, 4, -2 ], [ -1 ], [ -3, -4 ]>,
    <bipartition: [ 1, 2, 3, 4, -3 ], [ -1, -2 ], [ -4 ]>,
    <bipartition: [ 1, 2, 3, 4 ], [ -1, -2, -3 ], [ -4 ]>,
    <bipartition: [ 1, 2, 3, 4, -3, -4 ], [ -1 ], [ -2 ]>,
    <bipartition: [ 1, 2, 3 ], [ 4, -1, -2, -3 ], [ -4 ]>,
    <bipartition: [ 1, 2, -1, -2 ], [ 3, 4, -3 ], [ -4 ]>,
    <bipartition: [ 1, -3 ], [ 2, 3, 4, -1, -2 ], [ -4 ]>,
    <bipartition: [ 1, -3, -4 ], [ 2, 3, 4, -2 ], [ -1 ]> ]
\end{Verbatim}
 }

 

\subsection{\textcolor{Chapter }{IndecomposableElements}}
\logpage{[ 11, 7, 6 ]}\nobreak
\hyperdef{L}{X7B4CD8937858A895}{}
{\noindent\textcolor{FuncColor}{$\triangleright$\enspace\texttt{IndecomposableElements({\mdseries\slshape S})\index{IndecomposableElements@\texttt{IndecomposableElements}}
\label{IndecomposableElements}
}\hfill{\scriptsize (attribute)}}\\
\textbf{\indent Returns:\ }
A list of elements.



 If \mbox{\texttt{\mdseries\slshape S}} is a semigroup, then this attribute returns the set of elements of \mbox{\texttt{\mdseries\slshape S}} that are not decomposable. A element of \mbox{\texttt{\mdseries\slshape S}} is \emph{decomposable} if it can be written as the product of two elements in \mbox{\texttt{\mdseries\slshape S}}. An element of \mbox{\texttt{\mdseries\slshape S}} is \emph{indecomposable} if it is not decomposable. 

 See also \texttt{IsSurjectiveSemigroup} (\ref{IsSurjectiveSemigroup}). 

 Note that any generating set for \mbox{\texttt{\mdseries\slshape S}} contains each indecomposable element of \mbox{\texttt{\mdseries\slshape S}}. Thus \texttt{IndecomposableElements(\mbox{\texttt{\mdseries\slshape S}})} is a subset of \texttt{GeneratorsOfSemigroup(\mbox{\texttt{\mdseries\slshape S}})}. 

 
\begin{Verbatim}[commandchars=!@|,fontsize=\small,frame=single,label=Example]
  !gapprompt@gap>| !gapinput@S := Semigroup([|
  !gapprompt@>| !gapinput@Transformation([1, 1, 2, 3]),|
  !gapprompt@>| !gapinput@Transformation([1, 1, 1, 2])]);|
  <transformation semigroup of degree 4 with 2 generators>
  !gapprompt@gap>| !gapinput@x := IndecomposableElements(S);|
  [ Transformation( [ 1, 1, 2, 3 ] ) ]
  !gapprompt@gap>| !gapinput@IsSubset(GeneratorsOfSemigroup(S), x);|
  true
  !gapprompt@gap>| !gapinput@T := FullTransformationMonoid(10);|
  <full transformation monoid of degree 10>
  !gapprompt@gap>| !gapinput@IndecomposableElements(T);|
  [  ]
  !gapprompt@gap>| !gapinput@B := ZeroSemigroup(IsBipartitionSemigroup, 3);|
  <commutative non-regular bipartition semigroup of size 3, degree 4
   with 2 generators>
  !gapprompt@gap>| !gapinput@IndecomposableElements(B);|
  [ <bipartition: [ 1, 2, 3, -1 ], [ 4, -2 ], [ -3 ], [ -4 ]>,
    <bipartition: [ 1, 2, 4, -1 ], [ 3, -2 ], [ -3 ], [ -4 ]> ]
\end{Verbatim}
 }

 }

   
\section{\textcolor{Chapter }{ Minimal ideals and multiplicative zeros }}\logpage{[ 11, 8, 0 ]}
\hyperdef{L}{X830E18747A0B5BED}{}
{
  In this section we describe the attributes of a semigroup that can be found
using the \textsf{Semigroups} package. 

 

\subsection{\textcolor{Chapter }{MinimalIdeal}}
\logpage{[ 11, 8, 1 ]}\nobreak
\hyperdef{L}{X7BC68589879C3BE9}{}
{\noindent\textcolor{FuncColor}{$\triangleright$\enspace\texttt{MinimalIdeal({\mdseries\slshape S})\index{MinimalIdeal@\texttt{MinimalIdeal}}
\label{MinimalIdeal}
}\hfill{\scriptsize (attribute)}}\\
\textbf{\indent Returns:\ }
 The minimal ideal of a semigroup. 



 The minimal ideal of a semigroup is the least ideal with respect to
containment. 

 It is significantly easier to find the minimal $\mathcal{D}$\texttt{\symbol{45}}class of a semigroup, than to find its $\mathcal{D}$\texttt{\symbol{45}}classes. 

 See also \texttt{RepresentativeOfMinimalIdeal} (\ref{RepresentativeOfMinimalIdeal}), \texttt{PartialOrderOfDClasses} (\ref{PartialOrderOfDClasses}), \texttt{IsGreensLessThanOrEqual} (\textbf{Reference: IsGreensLessThanOrEqual}), and \texttt{MinimalDClass} (\ref{MinimalDClass}). 
\begin{Verbatim}[commandchars=!@|,fontsize=\small,frame=single,label=Example]
  !gapprompt@gap>| !gapinput@S := Semigroup(|
  !gapprompt@>| !gapinput@Transformation([3, 4, 1, 3, 6, 3, 4, 6, 10, 1]),|
  !gapprompt@>| !gapinput@Transformation([8, 2, 3, 8, 4, 1, 3, 4, 9, 7]));;|
  !gapprompt@gap>| !gapinput@MinimalIdeal(S);|
  <simple transformation semigroup ideal of degree 10 with 1 generator>
  !gapprompt@gap>| !gapinput@Elements(MinimalIdeal(S));|
  [ Transformation( [ 1, 1, 1, 1, 1, 1, 1, 1, 1, 1 ] ),
    Transformation( [ 3, 3, 3, 3, 3, 3, 3, 3, 3, 3 ] ),
    Transformation( [ 4, 4, 4, 4, 4, 4, 4, 4, 4, 4 ] ),
    Transformation( [ 6, 6, 6, 6, 6, 6, 6, 6, 6, 6 ] ),
    Transformation( [ 8, 8, 8, 8, 8, 8, 8, 8, 8, 8 ] ) ]
  !gapprompt@gap>| !gapinput@x := Transformation([8, 8, 8, 8, 8, 8, 8, 8, 8, 8]);;|
  !gapprompt@gap>| !gapinput@D := DClass(S, x);;|
  !gapprompt@gap>| !gapinput@ForAll(GreensDClasses(S), x -> IsGreensLessThanOrEqual(D, x));|
  true
  !gapprompt@gap>| !gapinput@MinimalIdeal(POI(10));|
  <partial perm group of rank 0>
  !gapprompt@gap>| !gapinput@MinimalIdeal(BrauerMonoid(6));|
  <simple bipartition *-semigroup ideal of degree 6 with 1 generator>
\end{Verbatim}
 }

 

\subsection{\textcolor{Chapter }{RepresentativeOfMinimalIdeal}}
\logpage{[ 11, 8, 2 ]}\nobreak
\hyperdef{L}{X7CA6744182D07C5B}{}
{\noindent\textcolor{FuncColor}{$\triangleright$\enspace\texttt{RepresentativeOfMinimalIdeal({\mdseries\slshape S})\index{RepresentativeOfMinimalIdeal@\texttt{RepresentativeOfMinimalIdeal}}
\label{RepresentativeOfMinimalIdeal}
}\hfill{\scriptsize (attribute)}}\\
\noindent\textcolor{FuncColor}{$\triangleright$\enspace\texttt{RepresentativeOfMinimalDClass({\mdseries\slshape S})\index{RepresentativeOfMinimalDClass@\texttt{RepresentativeOfMinimalDClass}}
\label{RepresentativeOfMinimalDClass}
}\hfill{\scriptsize (attribute)}}\\
\textbf{\indent Returns:\ }
 An element of the minimal ideal of a semigroup. 



 The minimal ideal of a semigroup is the least ideal with respect to
containment.

 This method returns a representative element of the minimal ideal of \mbox{\texttt{\mdseries\slshape S}} without having to create the minimal ideal itself. In general, beyond being a
member of the minimal ideal, the returned element is not guaranteed to have
any special properties. However, the element will coincide with the zero
element of \mbox{\texttt{\mdseries\slshape S}} if one exists. 

 This method works particularly well if \mbox{\texttt{\mdseries\slshape S}} is a semigroup of transformations or partial permutations.

 See also \texttt{MinimalIdeal} (\ref{MinimalIdeal}) and \texttt{MinimalDClass} (\ref{MinimalDClass}). 
\begin{Verbatim}[commandchars=!@|,fontsize=\small,frame=single,label=Example]
  !gapprompt@gap>| !gapinput@S := SymmetricInverseSemigroup(10);;|
  !gapprompt@gap>| !gapinput@RepresentativeOfMinimalIdeal(S);|
  <empty partial perm>
  !gapprompt@gap>| !gapinput@B := Semigroup([|
  !gapprompt@>| !gapinput@Bipartition([[1, 2], [3, 6, -2], [4, 5, -3, -4], [-1, -6], [-5]]),|
  !gapprompt@>| !gapinput@Bipartition([[1, -1], [2], [3], [4, -3], [5, 6, -5, -6],|
  !gapprompt@>| !gapinput@  [-2, -4]])]);;|
  !gapprompt@gap>| !gapinput@RepresentativeOfMinimalIdeal(B);|
  <bipartition: [ 1, 2 ], [ 3, 6 ], [ 4, 5 ], [ -1, -5, -6 ],
   [ -2, -4 ], [ -3 ]>
  !gapprompt@gap>| !gapinput@S := Semigroup(Transformation([5, 1, 6, 2, 2, 4]),|
  !gapprompt@>| !gapinput@                  Transformation([3, 5, 5, 1, 6, 2]));;|
  !gapprompt@gap>| !gapinput@RepresentativeOfMinimalDClass(S);|
  Transformation( [ 5, 6, 6, 3, 3, 5 ] )
  !gapprompt@gap>| !gapinput@MinimalDClass(S);|
  <Green's D-class: Transformation( [ 5, 6, 6, 3, 3, 5 ] )>
\end{Verbatim}
 }

 

\subsection{\textcolor{Chapter }{MultiplicativeZero}}
\logpage{[ 11, 8, 3 ]}\nobreak
\hyperdef{L}{X7B39F93C8136D642}{}
{\noindent\textcolor{FuncColor}{$\triangleright$\enspace\texttt{MultiplicativeZero({\mdseries\slshape S})\index{MultiplicativeZero@\texttt{MultiplicativeZero}}
\label{MultiplicativeZero}
}\hfill{\scriptsize (attribute)}}\\
\textbf{\indent Returns:\ }
 The zero element of a semigroup. 



 \texttt{MultiplicativeZero} returns the zero element of the semigroup \mbox{\texttt{\mdseries\slshape S}} if it exists and \texttt{fail} if it does not. See also \texttt{MultiplicativeZero} (\textbf{Reference: MultiplicativeZero}). 
\begin{Verbatim}[commandchars=!@|,fontsize=\small,frame=single,label=Example]
  !gapprompt@gap>| !gapinput@S := Semigroup(Transformation([1, 4, 2, 6, 6, 5, 2]),|
  !gapprompt@>| !gapinput@                  Transformation([1, 6, 3, 6, 2, 1, 6]));;|
  !gapprompt@gap>| !gapinput@MultiplicativeZero(S);|
  Transformation( [ 1, 1, 1, 1, 1, 1, 1 ] )
  !gapprompt@gap>| !gapinput@S := Semigroup(Transformation([2, 8, 3, 7, 1, 5, 2, 6]),|
  !gapprompt@>| !gapinput@                  Transformation([3, 5, 7, 2, 5, 6, 3, 8]),|
  !gapprompt@>| !gapinput@                  Transformation([6, 7, 4, 1, 4, 1, 6, 2]),|
  !gapprompt@>| !gapinput@                  Transformation([8, 8, 5, 1, 7, 5, 2, 8]));;|
  !gapprompt@gap>| !gapinput@MultiplicativeZero(S);|
  fail
  !gapprompt@gap>| !gapinput@S := InverseSemigroup(|
  !gapprompt@>| !gapinput@PartialPerm([1, 3, 4], [5, 3, 1]),|
  !gapprompt@>| !gapinput@PartialPerm([1, 2, 3, 4], [4, 3, 1, 2]),|
  !gapprompt@>| !gapinput@PartialPerm([1, 3, 4, 5], [2, 4, 5, 3]));;|
  !gapprompt@gap>| !gapinput@MultiplicativeZero(S);|
  <empty partial perm>
  !gapprompt@gap>| !gapinput@S := PartitionMonoid(6);|
  <regular bipartition *-monoid of size 4213597, degree 6 with 4
   generators>
  !gapprompt@gap>| !gapinput@MultiplicativeZero(S);|
  fail
  !gapprompt@gap>| !gapinput@S := DualSymmetricInverseMonoid(6);|
  <inverse block bijection monoid of degree 6 with 3 generators>
  !gapprompt@gap>| !gapinput@MultiplicativeZero(S);|
  <block bijection: [ 1, 2, 3, 4, 5, 6, -1, -2, -3, -4, -5, -6 ]>
\end{Verbatim}
 }

 

\subsection{\textcolor{Chapter }{UnderlyingSemigroupOfSemigroupWithAdjoinedZero}}
\logpage{[ 11, 8, 4 ]}\nobreak
\hyperdef{L}{X7CD6F5CB83B030B6}{}
{\noindent\textcolor{FuncColor}{$\triangleright$\enspace\texttt{UnderlyingSemigroupOfSemigroupWithAdjoinedZero({\mdseries\slshape S})\index{UnderlyingSemigroupOfSemigroupWithAdjoinedZero@\texttt{Underlying}\-\texttt{Semigroup}\-\texttt{Of}\-\texttt{Semigroup}\-\texttt{With}\-\texttt{Adjoined}\-\texttt{Zero}}
\label{UnderlyingSemigroupOfSemigroupWithAdjoinedZero}
}\hfill{\scriptsize (attribute)}}\\
\textbf{\indent Returns:\ }
 A semigroup, or \texttt{fail}. 



 If \mbox{\texttt{\mdseries\slshape S}} is a semigroup for which the property \texttt{IsSemigroupWithAdjoinedZero} (\ref{IsSemigroupWithAdjoinedZero}) is true, (i.e. \mbox{\texttt{\mdseries\slshape S}} has a \texttt{MultiplicativeZero} (\ref{MultiplicativeZero}) and the set $\mbox{\texttt{\mdseries\slshape S}} \setminus \{ 0 \}$ is a subsemigroup of \mbox{\texttt{\mdseries\slshape S}}), then this method returns the semigroup $\mbox{\texttt{\mdseries\slshape S}} \setminus \{ 0 \}$. 

 Otherwise, if \mbox{\texttt{\mdseries\slshape S}} is a semigroup for which the property \texttt{IsSemigroupWithAdjoinedZero} (\ref{IsSemigroupWithAdjoinedZero}) is \texttt{false}, then this method returns \texttt{fail}. 
\begin{Verbatim}[commandchars=!@|,fontsize=\small,frame=single,label=Example]
  !gapprompt@gap>| !gapinput@S := Semigroup([|
  !gapprompt@>| !gapinput@Transformation([2, 3, 4, 5, 1, 6]),|
  !gapprompt@>| !gapinput@Transformation([2, 1, 3, 4, 5, 6]),|
  !gapprompt@>| !gapinput@Transformation([6, 6, 6, 6, 6, 6])]);|
  <transformation semigroup of degree 6 with 3 generators>
  !gapprompt@gap>| !gapinput@MultiplicativeZero(S);|
  Transformation( [ 6, 6, 6, 6, 6, 6 ] )
  !gapprompt@gap>| !gapinput@G := UnderlyingSemigroupOfSemigroupWithAdjoinedZero(S);|
  <transformation semigroup of degree 5 with 2 generators>
  !gapprompt@gap>| !gapinput@IsGroupAsSemigroup(G);|
  true
  !gapprompt@gap>| !gapinput@IsZeroGroup(S);|
  true
  !gapprompt@gap>| !gapinput@S := SymmetricInverseMonoid(6);;|
  !gapprompt@gap>| !gapinput@MultiplicativeZero(S);|
  <empty partial perm>
  !gapprompt@gap>| !gapinput@G := UnderlyingSemigroupOfSemigroupWithAdjoinedZero(S);|
  fail
\end{Verbatim}
 }

 }

   
\section{\textcolor{Chapter }{ Group of units and identity elements }}\logpage{[ 11, 9, 0 ]}
\hyperdef{L}{X7CAB17667ED5A6E8}{}
{
  

\subsection{\textcolor{Chapter }{GroupOfUnits}}
\logpage{[ 11, 9, 1 ]}\nobreak
\hyperdef{L}{X811AEDD88280C277}{}
{\noindent\textcolor{FuncColor}{$\triangleright$\enspace\texttt{GroupOfUnits({\mdseries\slshape S})\index{GroupOfUnits@\texttt{GroupOfUnits}}
\label{GroupOfUnits}
}\hfill{\scriptsize (attribute)}}\\
\textbf{\indent Returns:\ }
The group of units of a semigroup or \texttt{fail}.



 \texttt{GroupOfUnits} returns the group of units of the semigroup \mbox{\texttt{\mdseries\slshape S}} as a subsemigroup of \mbox{\texttt{\mdseries\slshape S}} if it exists and returns \texttt{fail} if it does not. Use \texttt{IsomorphismPermGroup} (\ref{IsomorphismPermGroup}) if you require a permutation representation of the group of units.

 If a semigroup \mbox{\texttt{\mdseries\slshape S}} has an identity \texttt{e}, then the \emph{group of units} of \mbox{\texttt{\mdseries\slshape S}} is the set of those \texttt{s} in \mbox{\texttt{\mdseries\slshape S}} such that there exists \texttt{t} in \mbox{\texttt{\mdseries\slshape S}} where \texttt{s*t=t*s=e}. Equivalently, the group of units is the $\mathcal{H}$\texttt{\symbol{45}}class of the identity of \mbox{\texttt{\mdseries\slshape S}}.

 See also \texttt{GreensHClassOfElement} (\textbf{Reference: GreensHClassOfElement}), \texttt{IsMonoidAsSemigroup} (\ref{IsMonoidAsSemigroup}), and \texttt{MultiplicativeNeutralElement} (\textbf{Reference: MultiplicativeNeutralElement}). 
\begin{Verbatim}[commandchars=!@|,fontsize=\small,frame=single,label=Example]
  !gapprompt@gap>| !gapinput@S := Semigroup(|
  !gapprompt@>| !gapinput@Transformation([1, 2, 5, 4, 3, 8, 7, 6]),|
  !gapprompt@>| !gapinput@Transformation([1, 6, 3, 4, 7, 2, 5, 8]),|
  !gapprompt@>| !gapinput@Transformation([2, 1, 6, 7, 8, 3, 4, 5]),|
  !gapprompt@>| !gapinput@Transformation([3, 2, 3, 6, 1, 6, 1, 2]),|
  !gapprompt@>| !gapinput@Transformation([5, 2, 3, 6, 3, 4, 7, 4]));;|
  !gapprompt@gap>| !gapinput@Size(S);|
  5304
  !gapprompt@gap>| !gapinput@StructureDescription(GroupOfUnits(S));|
  "C2 x S4"
  !gapprompt@gap>| !gapinput@S := InverseSemigroup(|
  !gapprompt@>| !gapinput@PartialPerm([1, 2, 3, 4, 5, 6, 7, 8, 9, 10],|
  !gapprompt@>| !gapinput@            [2, 4, 5, 3, 6, 7, 10, 9, 8, 1]),|
  !gapprompt@>| !gapinput@PartialPerm([1, 2, 3, 4, 5, 6, 7, 8, 10],|
  !gapprompt@>| !gapinput@            [8, 2, 3, 1, 4, 5, 10, 6, 9]));;|
  !gapprompt@gap>| !gapinput@StructureDescription(GroupOfUnits(S));|
  "C8"
  !gapprompt@gap>| !gapinput@S := InverseSemigroup(|
  !gapprompt@>| !gapinput@PartialPerm([1, 3, 4], [4, 3, 5]),|
  !gapprompt@>| !gapinput@PartialPerm([1, 2, 3, 5], [3, 1, 5, 2]));;|
  !gapprompt@gap>| !gapinput@GroupOfUnits(S);|
  fail
  !gapprompt@gap>| !gapinput@S := Semigroup(|
  !gapprompt@>| !gapinput@Bipartition([[1, 2, 3, -1, -3], [-2]]),|
  !gapprompt@>| !gapinput@Bipartition([[1, -1], [2, 3, -2, -3]]),|
  !gapprompt@>| !gapinput@Bipartition([[1, -2], [2, -3], [3, -1]]),|
  !gapprompt@>| !gapinput@Bipartition([[1], [2, 3, -2], [-1, -3]]));;|
  !gapprompt@gap>| !gapinput@StructureDescription(GroupOfUnits(S));|
  "C3"
\end{Verbatim}
 }

  }

   
\section{\textcolor{Chapter }{ Idempotents }}\logpage{[ 11, 10, 0 ]}
\hyperdef{L}{X7C651C9C78398FFF}{}
{
  

\subsection{\textcolor{Chapter }{Idempotents}}
\logpage{[ 11, 10, 1 ]}\nobreak
\hyperdef{L}{X7C651C9C78398FFF}{}
{\noindent\textcolor{FuncColor}{$\triangleright$\enspace\texttt{Idempotents({\mdseries\slshape obj[, n]})\index{Idempotents@\texttt{Idempotents}}
\label{Idempotents}
}\hfill{\scriptsize (attribute)}}\\
\textbf{\indent Returns:\ }
A list of idempotents.



 The argument \mbox{\texttt{\mdseries\slshape obj}} should be a semigroup, $\mathcal{D}$\texttt{\symbol{45}}class, $\mathcal{H}$\texttt{\symbol{45}}class, $\mathcal{L}$\texttt{\symbol{45}}class, or $\mathcal{R}$\texttt{\symbol{45}}class.

 If the optional second argument \mbox{\texttt{\mdseries\slshape n}} is present and \mbox{\texttt{\mdseries\slshape obj}} is a semigroup, then a list of the idempotents in \mbox{\texttt{\mdseries\slshape obj}} of rank \mbox{\texttt{\mdseries\slshape n}} is returned. If you are only interested in the idempotents of a given rank,
then the second version of the function will probably be faster. However, if
the optional second argument is present, then nothing is stored in \mbox{\texttt{\mdseries\slshape obj}} and so every time the function is called the computation must be repeated. 

 This functions produce essentially the same output as the \textsf{GAP} library function with the same name; see \texttt{Idempotents} (\textbf{Reference: Idempotents}). The main difference is that this function can be applied to a wider class of
objects as described above.

 See also \texttt{IsRegularDClass} (\textbf{Reference: IsRegularDClass}), \texttt{IsRegularGreensClass} (\ref{IsRegularGreensClass}) \texttt{IsGroupHClass} (\textbf{Reference: IsGroupHClass}), \texttt{NrIdempotents} (\ref{NrIdempotents}), and \texttt{GroupHClass} (\ref{GroupHClass}). 
\begin{Verbatim}[commandchars=!@|,fontsize=\small,frame=single,label=Example]
  !gapprompt@gap>| !gapinput@S := Semigroup(Transformation([2, 3, 4, 1]),|
  !gapprompt@>| !gapinput@                  Transformation([3, 3, 1, 1]));;|
  !gapprompt@gap>| !gapinput@Idempotents(S, 1);|
  [  ]
  !gapprompt@gap>| !gapinput@AsSet(Idempotents(S, 2));|
  [ Transformation( [ 1, 1, 3, 3 ] ), Transformation( [ 1, 3, 3, 1 ] ),
    Transformation( [ 2, 2, 4, 4 ] ), Transformation( [ 4, 2, 2, 4 ] ) ]
  !gapprompt@gap>| !gapinput@AsSet(Idempotents(S));|
  [ Transformation( [ 1, 1, 3, 3 ] ), IdentityTransformation,
    Transformation( [ 1, 3, 3, 1 ] ), Transformation( [ 2, 2, 4, 4 ] ),
    Transformation( [ 4, 2, 2, 4 ] ) ]
  !gapprompt@gap>| !gapinput@x := Transformation([2, 2, 4, 4]);;|
  !gapprompt@gap>| !gapinput@R := GreensRClassOfElement(S, x);;|
  !gapprompt@gap>| !gapinput@Idempotents(R);|
  [ Transformation( [ 1, 1, 3, 3 ] ), Transformation( [ 2, 2, 4, 4 ] ) ]
  !gapprompt@gap>| !gapinput@x := Transformation([4, 2, 2, 4]);;|
  !gapprompt@gap>| !gapinput@L := GreensLClassOfElement(S, x);;|
  !gapprompt@gap>| !gapinput@AsSet(Idempotents(L));|
  [ Transformation( [ 2, 2, 4, 4 ] ), Transformation( [ 4, 2, 2, 4 ] ) ]
  !gapprompt@gap>| !gapinput@D := DClassOfLClass(L);;|
  !gapprompt@gap>| !gapinput@AsSet(Idempotents(D));|
  [ Transformation( [ 1, 1, 3, 3 ] ), Transformation( [ 1, 3, 3, 1 ] ),
    Transformation( [ 2, 2, 4, 4 ] ), Transformation( [ 4, 2, 2, 4 ] ) ]
  !gapprompt@gap>| !gapinput@L := GreensLClassOfElement(S, Transformation([3, 1, 1, 3]));;|
  !gapprompt@gap>| !gapinput@AsSet(Idempotents(L));|
  [ Transformation( [ 1, 1, 3, 3 ] ), Transformation( [ 1, 3, 3, 1 ] ) ]
  !gapprompt@gap>| !gapinput@H := GroupHClass(D);|
  <Green's H-class: Transformation( [ 1, 1, 3, 3 ] )>
  !gapprompt@gap>| !gapinput@Idempotents(H);|
  [ Transformation( [ 1, 1, 3, 3 ] ) ]
  !gapprompt@gap>| !gapinput@S := InverseSemigroup(|
  !gapprompt@>| !gapinput@PartialPerm([10, 6, 3, 4, 9, 0, 1]),|
  !gapprompt@>| !gapinput@PartialPerm([6, 10, 7, 4, 8, 2, 9, 1]));;|
  !gapprompt@gap>| !gapinput@Idempotents(S, 1);|
  [ <identity partial perm on [ 4 ]> ]
  !gapprompt@gap>| !gapinput@Idempotents(S, 0);|
  [  ]
\end{Verbatim}
 }

 

\subsection{\textcolor{Chapter }{NrIdempotents}}
\logpage{[ 11, 10, 2 ]}\nobreak
\hyperdef{L}{X7CFC4DB387452320}{}
{\noindent\textcolor{FuncColor}{$\triangleright$\enspace\texttt{NrIdempotents({\mdseries\slshape obj})\index{NrIdempotents@\texttt{NrIdempotents}}
\label{NrIdempotents}
}\hfill{\scriptsize (attribute)}}\\
\textbf{\indent Returns:\ }
 A positive integer. 



 This function returns the number of idempotents in \mbox{\texttt{\mdseries\slshape obj}} where \mbox{\texttt{\mdseries\slshape obj}} can be a semigroup, $\mathcal{D}$\texttt{\symbol{45}}, $\mathcal{L}$\texttt{\symbol{45}}, $\mathcal{H}$\texttt{\symbol{45}}, or $\mathcal{R}$\texttt{\symbol{45}}class. If the actual idempotents are not required, then it
is more efficient to use \texttt{NrIdempotents(obj)} than \texttt{Length(Idempotents(obj))} since the idempotents themselves are not created when \texttt{NrIdempotents} is called.

 See also \texttt{Idempotents} (\textbf{Reference: Idempotents}) and \texttt{Idempotents} (\ref{Idempotents}), \texttt{IsRegularDClass} (\textbf{Reference: IsRegularDClass}), \texttt{IsRegularGreensClass} (\ref{IsRegularGreensClass}) \texttt{IsGroupHClass} (\textbf{Reference: IsGroupHClass}), and \texttt{GroupHClass} (\ref{GroupHClass}). 
\begin{Verbatim}[commandchars=!@|,fontsize=\small,frame=single,label=Example]
  !gapprompt@gap>| !gapinput@S := Semigroup(Transformation([2, 3, 4, 1]),|
  !gapprompt@>| !gapinput@                  Transformation([3, 3, 1, 1]));;|
  !gapprompt@gap>| !gapinput@NrIdempotents(S);|
  5
  !gapprompt@gap>| !gapinput@f := Transformation([2, 2, 4, 4]);;|
  !gapprompt@gap>| !gapinput@R := GreensRClassOfElement(S, f);;|
  !gapprompt@gap>| !gapinput@NrIdempotents(R);|
  2
  !gapprompt@gap>| !gapinput@f := Transformation([4, 2, 2, 4]);;|
  !gapprompt@gap>| !gapinput@L := GreensLClassOfElement(S, f);;|
  !gapprompt@gap>| !gapinput@NrIdempotents(L);|
  2
  !gapprompt@gap>| !gapinput@D := DClassOfLClass(L);;|
  !gapprompt@gap>| !gapinput@NrIdempotents(D);|
  4
  !gapprompt@gap>| !gapinput@L := GreensLClassOfElement(S, Transformation([3, 1, 1, 3]));;|
  !gapprompt@gap>| !gapinput@NrIdempotents(L);|
  2
  !gapprompt@gap>| !gapinput@H := GroupHClass(D);;|
  !gapprompt@gap>| !gapinput@NrIdempotents(H);|
  1
  !gapprompt@gap>| !gapinput@S := InverseSemigroup(|
  !gapprompt@>| !gapinput@PartialPerm([1, 2, 3, 5, 7, 9, 10],|
  !gapprompt@>| !gapinput@             [6, 7, 2, 9, 1, 5, 3]),|
  !gapprompt@>| !gapinput@PartialPerm([1, 2, 3, 5, 6, 7, 9, 10],|
  !gapprompt@>| !gapinput@            [8, 1, 9, 4, 10, 5, 6, 7]));;|
  !gapprompt@gap>| !gapinput@NrIdempotents(S);|
  236
  !gapprompt@gap>| !gapinput@f := PartialPerm([2, 3, 7, 9, 10],|
  !gapprompt@>| !gapinput@                    [7, 2, 1, 5, 3]);;|
  !gapprompt@gap>| !gapinput@D := DClassNC(S, f);;|
  !gapprompt@gap>| !gapinput@NrIdempotents(D);|
  13
\end{Verbatim}
 }

 

\subsection{\textcolor{Chapter }{IdempotentGeneratedSubsemigroup}}
\logpage{[ 11, 10, 3 ]}\nobreak
\hyperdef{L}{X83970D028143B79B}{}
{\noindent\textcolor{FuncColor}{$\triangleright$\enspace\texttt{IdempotentGeneratedSubsemigroup({\mdseries\slshape S})\index{IdempotentGeneratedSubsemigroup@\texttt{IdempotentGeneratedSubsemigroup}}
\label{IdempotentGeneratedSubsemigroup}
}\hfill{\scriptsize (attribute)}}\\
\textbf{\indent Returns:\ }
A semigroup. 



 \texttt{IdempotentGeneratedSubsemigroup} returns the subsemigroup of the semigroup \mbox{\texttt{\mdseries\slshape S}} generated by the idempotents of \mbox{\texttt{\mdseries\slshape S}}.

 See also \texttt{Idempotents} (\ref{Idempotents}) and \texttt{SmallGeneratingSet} (\ref{SmallGeneratingSet}). 
\begin{Verbatim}[commandchars=!@|,fontsize=\small,frame=single,label=Example]
  !gapprompt@gap>| !gapinput@S := Semigroup(Transformation([1, 1]),|
  !gapprompt@>| !gapinput@                  Transformation([2, 1]),|
  !gapprompt@>| !gapinput@                  Transformation([1, 2, 2]),|
  !gapprompt@>| !gapinput@                  Transformation([1, 2, 3, 4, 5, 1]),|
  !gapprompt@>| !gapinput@                  Transformation([1, 2, 3, 4, 5, 5]),|
  !gapprompt@>| !gapinput@                  Transformation([1, 2, 3, 4, 6, 5]),|
  !gapprompt@>| !gapinput@                  Transformation([1, 2, 3, 5, 4]),|
  !gapprompt@>| !gapinput@                  Transformation([1, 2, 3, 7, 4, 5, 7]),|
  !gapprompt@>| !gapinput@                  Transformation([1, 2, 4, 8, 8, 3, 8, 7]),|
  !gapprompt@>| !gapinput@                  Transformation([1, 2, 8, 4, 5, 6, 7, 8]),|
  !gapprompt@>| !gapinput@                  Transformation([7, 7, 7, 4, 5, 6, 1]));;|
  !gapprompt@gap>| !gapinput@IdempotentGeneratedSubsemigroup(S) =|
  !gapprompt@>| !gapinput@Monoid(Transformation([1, 1]),|
  !gapprompt@>| !gapinput@       Transformation([1, 2, 1]),|
  !gapprompt@>| !gapinput@       Transformation([1, 2, 2]),|
  !gapprompt@>| !gapinput@       Transformation([1, 2, 3, 1]),|
  !gapprompt@>| !gapinput@       Transformation([1, 2, 3, 2]),|
  !gapprompt@>| !gapinput@       Transformation([1, 2, 3, 4, 1]),|
  !gapprompt@>| !gapinput@       Transformation([1, 2, 3, 4, 2]),|
  !gapprompt@>| !gapinput@       Transformation([1, 2, 3, 4, 4]),|
  !gapprompt@>| !gapinput@       Transformation([1, 2, 3, 4, 5, 1]),|
  !gapprompt@>| !gapinput@       Transformation([1, 2, 3, 4, 5, 2]),|
  !gapprompt@>| !gapinput@       Transformation([1, 2, 3, 4, 5, 5]),|
  !gapprompt@>| !gapinput@       Transformation([1, 2, 3, 4, 5, 7, 7]),|
  !gapprompt@>| !gapinput@       Transformation([1, 2, 3, 4, 7, 6, 7]),|
  !gapprompt@>| !gapinput@       Transformation([1, 2, 3, 6, 5, 6]),|
  !gapprompt@>| !gapinput@       Transformation([1, 2, 3, 7, 5, 6, 7]),|
  !gapprompt@>| !gapinput@       Transformation([1, 2, 8, 4, 5, 6, 7, 8]),|
  !gapprompt@>| !gapinput@       Transformation([2, 2]));|
  true
  !gapprompt@gap>| !gapinput@S := SymmetricInverseSemigroup(5);|
  <symmetric inverse monoid of degree 5>
  !gapprompt@gap>| !gapinput@IdempotentGeneratedSubsemigroup(S);|
  <inverse partial perm monoid of rank 5 with 5 generators>
  !gapprompt@gap>| !gapinput@S := DualSymmetricInverseSemigroup(5);|
  <inverse block bijection monoid of degree 5 with 3 generators>
  !gapprompt@gap>| !gapinput@IdempotentGeneratedSubsemigroup(S);|
  <inverse block bijection monoid of degree 5 with 10 generators>
  !gapprompt@gap>| !gapinput@IsSemilattice(last);|
  true
\end{Verbatim}
 }

 }

   
\section{\textcolor{Chapter }{ Maximal subsemigroups }}\label{Computing maximal subsemigroups}
\logpage{[ 11, 11, 0 ]}
\hyperdef{L}{X7D490B867CEFCBEF}{}
{
  The \textsf{Semigroups} package provides methods to calculate the maximal subsemigroups of a finite
semigroup, subject to various conditions. A \emph{maximal subsemigroup} of a semigroup is a proper subsemigroup that is contained in no other proper
subsemigroup of the semigroup. 

 When computing the maximal subsemigroups of a regular Rees
(0\texttt{\symbol{45}})matrix semigroup over a group, additional functionality
is available. As described in \cite{Graham1968aa}, a maximal subsemigroup of a finite regular Rees
(0\texttt{\symbol{45}})matrix semigroup over a group is one of 6 possible
types. Using the \textsf{Semigroups} package, it is possible to search for only those maximal subsemigroups of
certain types. 

 A maximal subsemigroup of such a Rees (0\texttt{\symbol{45}})matrix semigroup \texttt{R} over a group \texttt{G} is either: 
\begin{enumerate}
\item \texttt{\texttt{\symbol{123}}0\texttt{\symbol{125}};}
\item formed by removing \texttt{0};
\item formed by removing a column (a non\texttt{\symbol{45}}zero $\mathcal{L}$\texttt{\symbol{45}}class);
\item formed by removing a row (a non\texttt{\symbol{45}}zero $\mathcal{R}$\texttt{\symbol{45}}class);
\item formed by removing a set of both rows and columns;
\item isomorphic to a Rees (0\texttt{\symbol{45}})matrix semigroup of the same
dimensions over a maximal subgroup of \texttt{G} (in particular, the maximal subsemigroup intersects every $\mathcal{H}$\texttt{\symbol{45}}class of \texttt{R}).
\end{enumerate}
 Note that if \texttt{R} is a Rees matrix semigroup then it has no maximal subsemigroups of types 1, 2,
or 5. Only types 3, 4, and 6 are relevant to a Rees matrix semigroup. 

\subsection{\textcolor{Chapter }{MaximalSubsemigroups (for a finite semigroup)}}
\logpage{[ 11, 11, 1 ]}\nobreak
\hyperdef{L}{X860A10E387C19150}{}
{\noindent\textcolor{FuncColor}{$\triangleright$\enspace\texttt{MaximalSubsemigroups({\mdseries\slshape S})\index{MaximalSubsemigroups@\texttt{MaximalSubsemigroups}!for a finite semigroup}
\label{MaximalSubsemigroups:for a finite semigroup}
}\hfill{\scriptsize (attribute)}}\\
\noindent\textcolor{FuncColor}{$\triangleright$\enspace\texttt{MaximalSubsemigroups({\mdseries\slshape S, opts})\index{MaximalSubsemigroups@\texttt{MaximalSubsemigroups}!for a finite semigroup and a record}
\label{MaximalSubsemigroups:for a finite semigroup and a record}
}\hfill{\scriptsize (operation)}}\\
\textbf{\indent Returns:\ }
 The maximal subsemigroups of \mbox{\texttt{\mdseries\slshape S}}. 



 If \mbox{\texttt{\mdseries\slshape S}} is a finite semigroup, then the attribute \texttt{MaximalSubsemigroups} returns a list of the non\texttt{\symbol{45}}empty maximal subsemigroups of \mbox{\texttt{\mdseries\slshape S}}. The methods used by \texttt{MaximalSubsemigroups} are based on \cite{Graham1968aa}, and are described in \cite{Donoven2018aa}. 

 It is computationally expensive to search for the maximal subsemigroups of a
semigroup, and so computations involving \texttt{MaximalSubsemigroups} may be very lengthy. A substantial amount of information on the progress of \texttt{MaximalSubsemigroups} is provided through the info class \texttt{InfoSemigroups} (\ref{InfoSemigroups}), with increasingly detailed information given at levels 1, 2, and 3. 

 The behaviour of \texttt{MaximalSubsemigroups} can be altered via the second argument \mbox{\texttt{\mdseries\slshape opts}}, which should be a record. The optional components of \mbox{\texttt{\mdseries\slshape opts}} are: 
\begin{description}
\item[{gens (a boolean)}]  If \texttt{\mbox{\texttt{\mdseries\slshape opts}}.gens} is \texttt{false} or unspecified, then the maximal subsemigroups themselves are returned and not
just generating sets for these subsemigroups. 

 It can be more computationally expensive to return the generating sets for the
maximal subsemigroups, than to return the maximal subsemigroups themselves. 
\item[{contain (a list)}]  If \texttt{\mbox{\texttt{\mdseries\slshape opts}}.contain} is duplicate\texttt{\symbol{45}}free list of elements of \mbox{\texttt{\mdseries\slshape S}}, then \texttt{MaximalSubsemigroups} will search for the maximal subsemigroups of \mbox{\texttt{\mdseries\slshape S}} which contain those elements. 
\item[{D (a $\mathcal{D}$\texttt{\symbol{45}}class)}]  For a maximal subsemigroup \texttt{M} of a finite semigroup \mbox{\texttt{\mdseries\slshape S}}, there exists a unique $\mathcal{D}$\texttt{\symbol{45}}class which contains the complement of \texttt{M} in \mbox{\texttt{\mdseries\slshape S}}. In other words, the elements of \mbox{\texttt{\mdseries\slshape S}} which \texttt{M} lacks are contained in a unique $\mathcal{D}$\texttt{\symbol{45}}class. 

 If \texttt{\mbox{\texttt{\mdseries\slshape opts}}.D} is a $\mathcal{D}$\texttt{\symbol{45}}class of \mbox{\texttt{\mdseries\slshape S}}, then \texttt{MaximalSubsemigroups} will search exclusively for those maximal subsemigroups of \mbox{\texttt{\mdseries\slshape S}} whose complement is contained in \texttt{\mbox{\texttt{\mdseries\slshape opts}}.D}. 
\item[{types (a list)}]  \emph{This option is relevant only if \mbox{\texttt{\mdseries\slshape S}} is a regular Rees (0\texttt{\symbol{45}})matrix semigroup over a group.} 

 As described at the start of this subsection, \ref{Computing maximal subsemigroups}, a maximal subsemigroup of a regular Rees (0\texttt{\symbol{45}})matrix
semigroup over a group is one of 6 possible types. 

 If \mbox{\texttt{\mdseries\slshape S}} is a regular Rees (0\texttt{\symbol{45}})matrix semigroup over a group and \texttt{\mbox{\texttt{\mdseries\slshape opts}}.types} is a subset of \texttt{[1 .. 6]}, then \texttt{MaximalSubsemigroups} will search for those maximal subsemigroups of \mbox{\texttt{\mdseries\slshape S}} of the types enumerated by \texttt{\mbox{\texttt{\mdseries\slshape opts}}.types}. 

 The default value for this option is \texttt{[1 .. 6]} (i.e. no restriction). 
\end{description}
 
\begin{Verbatim}[commandchars=!@|,fontsize=\small,frame=single,label=Example]
  !gapprompt@gap>| !gapinput@S := FullTransformationSemigroup(3);|
  <full transformation monoid of degree 3>
  !gapprompt@gap>| !gapinput@MaximalSubsemigroups(S);|
  [ <transformation semigroup of degree 3 with 7 generators>,
    <transformation semigroup of degree 3 with 7 generators>,
    <transformation semigroup of degree 3 with 7 generators>,
    <transformation semigroup of degree 3 with 7 generators>,
    <transformation monoid of degree 3 with 5 generators> ]
  !gapprompt@gap>| !gapinput@MaximalSubsemigroups(S,|
  !gapprompt@>| !gapinput@rec(gens := true, D := DClass(S, Transformation([2, 2, 3]))));|
  [ [ Transformation( [ 1, 1, 1 ] ), Transformation( [ 3, 3, 3 ] ),
        Transformation( [ 2, 2, 2 ] ), IdentityTransformation,
        Transformation( [ 2, 3, 1 ] ), Transformation( [ 2, 1 ] ) ] ]
  !gapprompt@gap>| !gapinput@MaximalSubsemigroups(S,|
  !gapprompt@>| !gapinput@rec(contain := [Transformation([2, 3, 1])]));|
  [ <transformation semigroup of degree 3 with 7 generators>,
    <transformation monoid of degree 3 with 5 generators> ]
  !gapprompt@gap>| !gapinput@R := PrincipalFactor(|
  !gapprompt@>| !gapinput@DClass(FullTransformationMonoid(4), Transformation([2, 2])));|
  <Rees 0-matrix semigroup 6x4 over Group([ (2,3,4), (2,4) ])>
  !gapprompt@gap>| !gapinput@MaximalSubsemigroups(R, rec(types := [5],|
  !gapprompt@>| !gapinput@  contain := [RMSElement(R, 1, (), 1),|
  !gapprompt@>| !gapinput@              RMSElement(R, 1, (2, 3), 2)]));|
  [ <subsemigroup of 6x4 Rees 0-matrix semigroup with 10 generators>,
    <subsemigroup of 6x4 Rees 0-matrix semigroup with 10 generators>,
    <subsemigroup of 6x4 Rees 0-matrix semigroup with 10 generators>,
    <subsemigroup of 6x4 Rees 0-matrix semigroup with 10 generators> ]
\end{Verbatim}
 }

 

\subsection{\textcolor{Chapter }{NrMaximalSubsemigroups}}
\logpage{[ 11, 11, 2 ]}\nobreak
\hyperdef{L}{X7A52B4117BCEF379}{}
{\noindent\textcolor{FuncColor}{$\triangleright$\enspace\texttt{NrMaximalSubsemigroups({\mdseries\slshape S})\index{NrMaximalSubsemigroups@\texttt{NrMaximalSubsemigroups}}
\label{NrMaximalSubsemigroups}
}\hfill{\scriptsize (attribute)}}\\
\textbf{\indent Returns:\ }
 The number of maximal subsemigroups of \mbox{\texttt{\mdseries\slshape S}}. 



 If \mbox{\texttt{\mdseries\slshape S}} is a finite semigroup, then \texttt{NrMaximalSubsemigroups} returns the number of non\texttt{\symbol{45}}empty maximal subsemigroups of \mbox{\texttt{\mdseries\slshape S}}. The methods used by \texttt{MaximalSubsemigroups} are based on \cite{Graham1968aa}, and are described in \cite{Donoven2018aa}. 

 It can be significantly faster to find the number of maximal subsemigroups of
a semigroup than to find the maximal subsemigroups themselves. 

 Unless the maximal subsemigroups of \mbox{\texttt{\mdseries\slshape S}} are already known, the command \texttt{NrMaximalSubsemigroups(\mbox{\texttt{\mdseries\slshape S}})} simply calls the command \texttt{MaximalSubsemigroups(\mbox{\texttt{\mdseries\slshape S}}, rec(number := true))}. 

 For more information about searching for maximal subsemigroups of a finite
semigroup in the \textsf{Semigroups} package, and for information about the options available to alter the search,
see \texttt{MaximalSubsemigroups} (\ref{MaximalSubsemigroups:for a finite semigroup and a record}). By supplying the additional option \texttt{\mbox{\texttt{\mdseries\slshape opts}}.number :=} \texttt{true}, the number of maximal subsemigroups will be returned rather than the
subsemigroups themselves. 
\begin{Verbatim}[commandchars=!@|,fontsize=\small,frame=single,label=Example]
  !gapprompt@gap>| !gapinput@S := FullTransformationSemigroup(3);|
  <full transformation monoid of degree 3>
  !gapprompt@gap>| !gapinput@NrMaximalSubsemigroups(S);|
  5
  !gapprompt@gap>| !gapinput@S := RectangularBand(3, 4);;|
  !gapprompt@gap>| !gapinput@NrMaximalSubsemigroups(S);|
  7
  !gapprompt@gap>| !gapinput@R := PrincipalFactor(|
  !gapprompt@>| !gapinput@DClass(FullTransformationMonoid(4), Transformation([2, 2])));|
  <Rees 0-matrix semigroup 6x4 over Group([ (2,3,4), (2,4) ])>
  !gapprompt@gap>| !gapinput@MaximalSubsemigroups(R, rec(number := true, types := [3, 4]));|
  10
\end{Verbatim}
 }

 

\subsection{\textcolor{Chapter }{IsMaximalSubsemigroup}}
\logpage{[ 11, 11, 3 ]}\nobreak
\hyperdef{L}{X82D74C2478A49FD5}{}
{\noindent\textcolor{FuncColor}{$\triangleright$\enspace\texttt{IsMaximalSubsemigroup({\mdseries\slshape S, T})\index{IsMaximalSubsemigroup@\texttt{IsMaximalSubsemigroup}}
\label{IsMaximalSubsemigroup}
}\hfill{\scriptsize (operation)}}\\
\textbf{\indent Returns:\ }
 \texttt{true} or \texttt{false}. 



 If \mbox{\texttt{\mdseries\slshape S}} and \mbox{\texttt{\mdseries\slshape T}} are semigroups, then \texttt{IsMaximalSubsemigroup} returns \texttt{true} if and only if \mbox{\texttt{\mdseries\slshape T}} is a maximal subsemigroup of \mbox{\texttt{\mdseries\slshape S}}. 

 A \emph{maximal subsemigroup} of \mbox{\texttt{\mdseries\slshape S}} is a proper subsemigroup of \mbox{\texttt{\mdseries\slshape S}} which is contained in no other proper subsemigroup of \mbox{\texttt{\mdseries\slshape S}}. 
\begin{Verbatim}[commandchars=!@|,fontsize=\small,frame=single,label=Example]
  !gapprompt@gap>| !gapinput@S := ZeroSemigroup(2);;|
  !gapprompt@gap>| !gapinput@IsMaximalSubsemigroup(S, Semigroup(MultiplicativeZero(S)));|
  true
  !gapprompt@gap>| !gapinput@S := FullTransformationSemigroup(4);|
  <full transformation monoid of degree 4>
  !gapprompt@gap>| !gapinput@T := Semigroup(Transformation([3, 4, 1, 2]),|
  !gapprompt@>| !gapinput@                  Transformation([1, 4, 2, 3]),|
  !gapprompt@>| !gapinput@                  Transformation([2, 1, 1, 3]));|
  <transformation semigroup of degree 4 with 3 generators>
  !gapprompt@gap>| !gapinput@IsMaximalSubsemigroup(S, T);|
  true
  !gapprompt@gap>| !gapinput@R := Semigroup(Transformation([3, 4, 1, 2]),|
  !gapprompt@>| !gapinput@                  Transformation([1, 4, 2, 2]),|
  !gapprompt@>| !gapinput@                  Transformation([2, 1, 1, 3]));|
  <transformation semigroup of degree 4 with 3 generators>
  !gapprompt@gap>| !gapinput@IsMaximalSubsemigroup(S, R);|
  false
\end{Verbatim}
 }

 }

   
\section{\textcolor{Chapter }{ Attributes of transformations and transformation semigroups }}\logpage{[ 11, 12, 0 ]}
\hyperdef{L}{X87696C597F453F4F}{}
{
  

\subsection{\textcolor{Chapter }{ComponentRepsOfTransformationSemigroup}}
\logpage{[ 11, 12, 1 ]}\nobreak
\hyperdef{L}{X8065DBC48722B085}{}
{\noindent\textcolor{FuncColor}{$\triangleright$\enspace\texttt{ComponentRepsOfTransformationSemigroup({\mdseries\slshape S})\index{ComponentRepsOfTransformationSemigroup@\texttt{Component}\-\texttt{Reps}\-\texttt{Of}\-\texttt{Transformation}\-\texttt{Semigroup}}
\label{ComponentRepsOfTransformationSemigroup}
}\hfill{\scriptsize (attribute)}}\\
\textbf{\indent Returns:\ }
The representatives of components of a transformation semigroup.



 This function returns the representatives of the components of the action of
the transformation semigroup \mbox{\texttt{\mdseries\slshape S}} on the set of positive integers not greater than the degree of \mbox{\texttt{\mdseries\slshape S}}. 

 The representatives are the least set of points such that every point can be
reached from some representative under the action of \mbox{\texttt{\mdseries\slshape S}}. 
\begin{Verbatim}[commandchars=!@|,fontsize=\small,frame=single,label=Example]
  !gapprompt@gap>| !gapinput@S := Semigroup(|
  !gapprompt@>| !gapinput@Transformation([11, 11, 9, 6, 4, 1, 4, 1, 6, 7, 12, 5]),|
  !gapprompt@>| !gapinput@Transformation([12, 10, 7, 10, 4, 1, 12, 9, 11, 9, 1, 12]));;|
  !gapprompt@gap>| !gapinput@ComponentRepsOfTransformationSemigroup(S);|
  [ 2, 3, 8 ]
\end{Verbatim}
 }

 

\subsection{\textcolor{Chapter }{ComponentsOfTransformationSemigroup}}
\logpage{[ 11, 12, 2 ]}\nobreak
\hyperdef{L}{X8706A72A7F3EE532}{}
{\noindent\textcolor{FuncColor}{$\triangleright$\enspace\texttt{ComponentsOfTransformationSemigroup({\mdseries\slshape S})\index{ComponentsOfTransformationSemigroup@\texttt{ComponentsOfTransformationSemigroup}}
\label{ComponentsOfTransformationSemigroup}
}\hfill{\scriptsize (attribute)}}\\
\textbf{\indent Returns:\ }
The components of a transformation semigroup.



 This function returns the components of the action of the transformation
semigroup \mbox{\texttt{\mdseries\slshape S}} on the set of positive integers not greater than the degree of \mbox{\texttt{\mdseries\slshape S}}; the components of \mbox{\texttt{\mdseries\slshape S}} partition this set. 
\begin{Verbatim}[commandchars=!@|,fontsize=\small,frame=single,label=Example]
  !gapprompt@gap>| !gapinput@S := Semigroup(|
  !gapprompt@>| !gapinput@Transformation([11, 11, 9, 6, 4, 1, 4, 1, 6, 7, 12, 5]),|
  !gapprompt@>| !gapinput@Transformation([12, 10, 7, 10, 4, 1, 12, 9, 11, 9, 1, 12]));;|
  !gapprompt@gap>| !gapinput@ComponentsOfTransformationSemigroup(S);|
  [ [ 1, 2, 3, 4, 5, 6, 7, 8, 9, 10, 11, 12 ] ]
\end{Verbatim}
 }

 

\subsection{\textcolor{Chapter }{CyclesOfTransformationSemigroup}}
\logpage{[ 11, 12, 3 ]}\nobreak
\hyperdef{L}{X7AA697B186301F54}{}
{\noindent\textcolor{FuncColor}{$\triangleright$\enspace\texttt{CyclesOfTransformationSemigroup({\mdseries\slshape S})\index{CyclesOfTransformationSemigroup@\texttt{CyclesOfTransformationSemigroup}}
\label{CyclesOfTransformationSemigroup}
}\hfill{\scriptsize (attribute)}}\\
\textbf{\indent Returns:\ }
The cycles of a transformation semigroup.



 This function returns the cycles, or strongly connected components, of the
action of the transformation semigroup \mbox{\texttt{\mdseries\slshape S}} on the set of positive integers not greater than the degree of \mbox{\texttt{\mdseries\slshape S}}. 
\begin{Verbatim}[commandchars=!@|,fontsize=\small,frame=single,label=Example]
  !gapprompt@gap>| !gapinput@S := Semigroup(|
  !gapprompt@>| !gapinput@Transformation([11, 11, 9, 6, 4, 1, 4, 1, 6, 7, 12, 5]),|
  !gapprompt@>| !gapinput@Transformation([12, 10, 7, 10, 4, 1, 12, 9, 11, 9, 1, 12]));;|
  !gapprompt@gap>| !gapinput@CyclesOfTransformationSemigroup(S);|
  [ [ 1, 11, 12, 5, 4, 6, 10, 7, 9 ], [ 2 ], [ 3 ], [ 8 ] ]
\end{Verbatim}
 }

 

\subsection{\textcolor{Chapter }{DigraphOfAction (for a transformation semigroup, list, and action)}}
\logpage{[ 11, 12, 4 ]}\nobreak
\hyperdef{L}{X8089CF7182AD1925}{}
{\noindent\textcolor{FuncColor}{$\triangleright$\enspace\texttt{DigraphOfAction({\mdseries\slshape S, list, act})\index{DigraphOfAction@\texttt{DigraphOfAction}!for a transformation semigroup, list, and action}
\label{DigraphOfAction:for a transformation semigroup, list, and action}
}\hfill{\scriptsize (operation)}}\\
\textbf{\indent Returns:\ }
A digraph, or \texttt{fail}.



 If \mbox{\texttt{\mdseries\slshape S}} is a transformation semigroup and \mbox{\texttt{\mdseries\slshape list}} is list such that \mbox{\texttt{\mdseries\slshape S}} acts on the items in \mbox{\texttt{\mdseries\slshape list}} via the function \mbox{\texttt{\mdseries\slshape act}}, then \texttt{DigraphOfAction} returns a digraph representing the action of \mbox{\texttt{\mdseries\slshape S}} on the items in \mbox{\texttt{\mdseries\slshape list}} and any further items output by \texttt{\mbox{\texttt{\mdseries\slshape act}}(\mbox{\texttt{\mdseries\slshape list}}[i], \mbox{\texttt{\mdseries\slshape S}}.j)}.

 If \texttt{\mbox{\texttt{\mdseries\slshape act}}(\mbox{\texttt{\mdseries\slshape list}}[i], \mbox{\texttt{\mdseries\slshape S}}.j)} is the \texttt{k}\texttt{\symbol{45}}th item in \mbox{\texttt{\mdseries\slshape list}}, then in the output digraph there is an edge from the vertex \texttt{i} to the vertex \texttt{k} labelled by \texttt{j}. 

 The values in \mbox{\texttt{\mdseries\slshape list}} and the additional values generated are stored in the vertex labels of the
output digraph; see \texttt{DigraphVertexLabels} (\textbf{Digraphs: DigraphVertexLabels}), and the edge labels are stored in the \texttt{DigraphEdgeLabels} (\textbf{Digraphs: DigraphEdgeLabels}) 

 The digraph returned by \texttt{DigraphOfAction} has no multiple edges; see \texttt{IsMultiDigraph} (\textbf{Digraphs: IsMultiDigraph}). 
\begin{Verbatim}[commandchars=!@|,fontsize=\small,frame=single,label=Example]
  !gapprompt@gap>| !gapinput@S := Semigroup(Transformation([2, 4, 3, 4, 7, 1, 6]),|
  !gapprompt@>| !gapinput@                  Transformation([3, 3, 2, 3, 5, 1, 5]));|
  <transformation semigroup of degree 7 with 2 generators>
  !gapprompt@gap>| !gapinput@list := Concatenation(List([1 .. 7], x -> [x]),|
  !gapprompt@>| !gapinput@                         Combinations([1 .. 7], 2));|
  [ [ 1 ], [ 2 ], [ 3 ], [ 4 ], [ 5 ], [ 6 ], [ 7 ], [ 1, 2 ],
    [ 1, 3 ], [ 1, 4 ], [ 1, 5 ], [ 1, 6 ], [ 1, 7 ], [ 2, 3 ],
    [ 2, 4 ], [ 2, 5 ], [ 2, 6 ], [ 2, 7 ], [ 3, 4 ], [ 3, 5 ],
    [ 3, 6 ], [ 3, 7 ], [ 4, 5 ], [ 4, 6 ], [ 4, 7 ], [ 5, 6 ],
    [ 5, 7 ], [ 6, 7 ] ]
  !gapprompt@gap>| !gapinput@D := DigraphOfAction(S, list, OnSets);|
  <immutable digraph with 28 vertices, 54 edges>
  !gapprompt@gap>| !gapinput@OnSets([2, 5], S.1);|
  [ 4, 7 ]
  !gapprompt@gap>| !gapinput@Position(DigraphVertexLabels(D), [2, 5]);|
  16
  !gapprompt@gap>| !gapinput@DigraphVertexLabel(D, 25);|
  [ 4, 7 ]
  !gapprompt@gap>| !gapinput@DigraphEdgeLabel(D, 16, 25);|
  1
\end{Verbatim}
 }

 

\subsection{\textcolor{Chapter }{DigraphOfActionOnPoints (for a transformation semigroup)}}
\logpage{[ 11, 12, 5 ]}\nobreak
\hyperdef{L}{X7B5ACD5C7E7612A2}{}
{\noindent\textcolor{FuncColor}{$\triangleright$\enspace\texttt{DigraphOfActionOnPoints({\mdseries\slshape S})\index{DigraphOfActionOnPoints@\texttt{DigraphOfActionOnPoints}!for a transformation semigroup}
\label{DigraphOfActionOnPoints:for a transformation semigroup}
}\hfill{\scriptsize (attribute)}}\\
\noindent\textcolor{FuncColor}{$\triangleright$\enspace\texttt{DigraphOfActionOnPoints({\mdseries\slshape S, n})\index{DigraphOfActionOnPoints@\texttt{DigraphOfActionOnPoints}!for a transformation semigroup and an integer}
\label{DigraphOfActionOnPoints:for a transformation semigroup and an integer}
}\hfill{\scriptsize (attribute)}}\\
\textbf{\indent Returns:\ }
A digraph.



 If \mbox{\texttt{\mdseries\slshape S}} is a transformation semigroup and \mbox{\texttt{\mdseries\slshape n}} is a non\texttt{\symbol{45}}negative integer, then \texttt{DigraphOfActionOnPoints(\mbox{\texttt{\mdseries\slshape S}}, \mbox{\texttt{\mdseries\slshape n}})} returns a digraph representing the \texttt{OnPoints} (\textbf{Reference: OnPoints}) action of \mbox{\texttt{\mdseries\slshape S}} on the set \texttt{[1 .. \mbox{\texttt{\mdseries\slshape n}}]}. 

 If the optional argument \mbox{\texttt{\mdseries\slshape n}} is not specified, then by default the degree of \mbox{\texttt{\mdseries\slshape S}} will be chosen for \mbox{\texttt{\mdseries\slshape n}}; see \texttt{DegreeOfTransformationSemigroup} (\textbf{Reference: DegreeOfTransformationSemigroup}). 

 The digraph returned by \texttt{DigraphOfActionOnPoints} has \mbox{\texttt{\mdseries\slshape n}} vertices, where the vertex \texttt{i} corresponds to the point \texttt{i}. For each point \texttt{i} in \texttt{[1 .. \mbox{\texttt{\mdseries\slshape n}}]}, and for each generator \texttt{f} in \texttt{GeneratorsOfSemigroup(\mbox{\texttt{\mdseries\slshape S}})}, there is an edge from the vertex \texttt{i} to the vertex \texttt{i \texttt{\symbol{94}} f}. See \texttt{GeneratorsOfSemigroup} (\textbf{Reference: GeneratorsOfSemigroup}) for further information. 

 
\begin{Verbatim}[commandchars=!@|,fontsize=\small,frame=single,label=Example]
  !gapprompt@gap>| !gapinput@S := Semigroup(Transformation([2, 4, 2, 4, 7, 1, 6]),|
  !gapprompt@>| !gapinput@                  Transformation([3, 3, 2, 3, 5, 1, 5]));|
  <transformation semigroup of degree 7 with 2 generators>
  !gapprompt@gap>| !gapinput@D1 := DigraphOfActionOnPoints(S);|
  <immutable digraph with 7 vertices, 12 edges>
  !gapprompt@gap>| !gapinput@OnPoints(2, S.1);|
  4
  !gapprompt@gap>| !gapinput@D2 := DigraphOfActionOnPoints(S, 4);|
  <immutable digraph with 4 vertices, 7 edges>
  !gapprompt@gap>| !gapinput@D2 = InducedSubdigraph(D1, [1 .. 4]);|
  true
  !gapprompt@gap>| !gapinput@DigraphOfActionOnPoints(S, 5);|
  <immutable digraph with 5 vertices, 8 edges>
\end{Verbatim}
 }

 

\subsection{\textcolor{Chapter }{FixedPointsOfTransformationSemigroup (for a transformation semigroup)}}
\logpage{[ 11, 12, 6 ]}\nobreak
\hyperdef{L}{X7C6D8689819AEEE2}{}
{\noindent\textcolor{FuncColor}{$\triangleright$\enspace\texttt{FixedPointsOfTransformationSemigroup({\mdseries\slshape S})\index{FixedPointsOfTransformationSemigroup@\texttt{Fixed}\-\texttt{Points}\-\texttt{Of}\-\texttt{Transformation}\-\texttt{Semigroup}!for a transformation semigroup}
\label{FixedPointsOfTransformationSemigroup:for a transformation semigroup}
}\hfill{\scriptsize (attribute)}}\\
\textbf{\indent Returns:\ }
A set of positive integers.



 If \mbox{\texttt{\mdseries\slshape S}} is a transformation semigroup, then \texttt{FixedPointsOfTransformationSemigroup(\mbox{\texttt{\mdseries\slshape S}})} returns the set of points \texttt{i} in \texttt{[1 .. DegreeOfTransformationSemigroup(\mbox{\texttt{\mdseries\slshape S}})]} such that \texttt{i \texttt{\symbol{94}} f = i} for all \texttt{f} in \mbox{\texttt{\mdseries\slshape S}}. 

 
\begin{Verbatim}[commandchars=!@|,fontsize=\small,frame=single,label=Example]
  !gapprompt@gap>| !gapinput@f := Transformation([1, 4, 2, 4, 3, 7, 7]);|
  Transformation( [ 1, 4, 2, 4, 3, 7, 7 ] )
  !gapprompt@gap>| !gapinput@S := Semigroup(f);|
  <commutative transformation semigroup of degree 7 with 1 generator>
  !gapprompt@gap>| !gapinput@FixedPointsOfTransformationSemigroup(S);|
  [ 1, 4, 7 ]
\end{Verbatim}
 }

 

\subsection{\textcolor{Chapter }{IsTransitive (for a transformation
    semigroup and a set)}}
\logpage{[ 11, 12, 7 ]}\nobreak
\hyperdef{L}{X83DA161F875F63B1}{}
{\noindent\textcolor{FuncColor}{$\triangleright$\enspace\texttt{IsTransitive({\mdseries\slshape S[, X]})\index{IsTransitive@\texttt{IsTransitive}!for a transformation
    semigroup and a set}
\label{IsTransitive:for a transformation
    semigroup and a set}
}\hfill{\scriptsize (property)}}\\
\noindent\textcolor{FuncColor}{$\triangleright$\enspace\texttt{IsTransitive({\mdseries\slshape S[, n]})\index{IsTransitive@\texttt{IsTransitive}!for a transformation
    semigroup and a pos int}
\label{IsTransitive:for a transformation
    semigroup and a pos int}
}\hfill{\scriptsize (property)}}\\
\textbf{\indent Returns:\ }
\texttt{true} or \texttt{false}.



 A transformation semigroup \mbox{\texttt{\mdseries\slshape S}} is \emph{transitive} or \emph{strongly connected} on the set \mbox{\texttt{\mdseries\slshape X}} if for every \texttt{i, j} in \mbox{\texttt{\mdseries\slshape X}} there is an element \texttt{s} in \mbox{\texttt{\mdseries\slshape S}} such that \texttt{i \texttt{\symbol{94}} s = j}. 

 If the optional second argument is a positive integer \mbox{\texttt{\mdseries\slshape n}}, then \texttt{IsTransitive} returns \texttt{true} if \mbox{\texttt{\mdseries\slshape S}} is transitive on \texttt{[1 .. \mbox{\texttt{\mdseries\slshape n}}]}, and \texttt{false} if it is not. 

 If the optional second argument is not provided, then the degree of \mbox{\texttt{\mdseries\slshape S}} is used by default; see \texttt{DegreeOfTransformationSemigroup} (\textbf{Reference: DegreeOfTransformationSemigroup}). 
\begin{Verbatim}[commandchars=!@|,fontsize=\small,frame=single,label=Example]
  !gapprompt@gap>| !gapinput@S := Semigroup([|
  !gapprompt@>| !gapinput@ Bipartition([|
  !gapprompt@>| !gapinput@   [1, 2], [3, 6, -2], [4, 5, -3, -4], [-1, -6], [-5]]),|
  !gapprompt@>| !gapinput@ Bipartition([|
  !gapprompt@>| !gapinput@   [1, -4], [2, 3, 4, 5], [6], [-1, -6], [-2, -3], [-5]])]);|
  <bipartition semigroup of degree 6 with 2 generators>
  !gapprompt@gap>| !gapinput@AsSemigroup(IsTransformationSemigroup, S);|
  <transformation semigroup of size 11, degree 12 with 2 generators>
  !gapprompt@gap>| !gapinput@IsTransitive(last);|
  false
  !gapprompt@gap>| !gapinput@IsTransitive(AsSemigroup(Group((1, 2, 3))));|
  true
\end{Verbatim}
 }

 

\subsection{\textcolor{Chapter }{SmallestElementSemigroup}}
\logpage{[ 11, 12, 8 ]}\nobreak
\hyperdef{L}{X7C65202187A9C9F5}{}
{\noindent\textcolor{FuncColor}{$\triangleright$\enspace\texttt{SmallestElementSemigroup({\mdseries\slshape S})\index{SmallestElementSemigroup@\texttt{SmallestElementSemigroup}}
\label{SmallestElementSemigroup}
}\hfill{\scriptsize (attribute)}}\\
\noindent\textcolor{FuncColor}{$\triangleright$\enspace\texttt{LargestElementSemigroup({\mdseries\slshape S})\index{LargestElementSemigroup@\texttt{LargestElementSemigroup}}
\label{LargestElementSemigroup}
}\hfill{\scriptsize (attribute)}}\\
\textbf{\indent Returns:\ }
A transformation.



 These attributes return the smallest and largest element of the transformation
semigroup \mbox{\texttt{\mdseries\slshape S}}, respectively. Smallest means the first element in the sorted set of elements
of \mbox{\texttt{\mdseries\slshape S}} and largest means the last element in the set of elements. 

 It is not necessary to find the elements of the semigroup to determine the
smallest or largest element, and this function has considerable better
performance than the equivalent \texttt{Elements(\mbox{\texttt{\mdseries\slshape S}})[1]} and \texttt{Elements(\mbox{\texttt{\mdseries\slshape S}})[Size(\mbox{\texttt{\mdseries\slshape S}})]}. 
\begin{Verbatim}[commandchars=!@|,fontsize=\small,frame=single,label=Example]
  !gapprompt@gap>| !gapinput@S := Monoid(|
  !gapprompt@>| !gapinput@Transformation([1, 4, 11, 11, 7, 2, 6, 2, 5, 5, 10]),|
  !gapprompt@>| !gapinput@Transformation([2, 4, 4, 2, 10, 5, 11, 11, 11, 6, 7]));|
  <transformation monoid of degree 11 with 2 generators>
  !gapprompt@gap>| !gapinput@SmallestElementSemigroup(S);|
  IdentityTransformation
  !gapprompt@gap>| !gapinput@LargestElementSemigroup(S);|
  Transformation( [ 11, 11, 10, 10, 7, 6, 5, 6, 2, 2, 4 ] )
\end{Verbatim}
 }

 

\subsection{\textcolor{Chapter }{CanonicalTransformation}}
\logpage{[ 11, 12, 9 ]}\nobreak
\hyperdef{L}{X84792D3D804413CD}{}
{\noindent\textcolor{FuncColor}{$\triangleright$\enspace\texttt{CanonicalTransformation({\mdseries\slshape trans[, n]})\index{CanonicalTransformation@\texttt{CanonicalTransformation}}
\label{CanonicalTransformation}
}\hfill{\scriptsize (function)}}\\
\textbf{\indent Returns:\ }
 A transformation. 



 If \mbox{\texttt{\mdseries\slshape trans}} is a transformation, and \mbox{\texttt{\mdseries\slshape n}} is a non\texttt{\symbol{45}}negative integer such that the restriction of \mbox{\texttt{\mdseries\slshape trans}} to \texttt{[1 .. n]} defines a transformation of \texttt{[1 .. n]}, then \texttt{CanonicalTransformation} returns a canonical representative of the transformation \mbox{\texttt{\mdseries\slshape trans}} restricted to \texttt{[1 .. n]}. 

 More specifically, let \texttt{C(n)} be a class of transformations of degree \mbox{\texttt{\mdseries\slshape n}} such that \texttt{AsDigraph} returns isomorphic digraphs for every pair of element elements in \texttt{C(n)}. Recall that for a transformation \mbox{\texttt{\mdseries\slshape trans}} and integer \mbox{\texttt{\mdseries\slshape n}} the function \texttt{AsDigraph} returns a digraph with \mbox{\texttt{\mdseries\slshape n}} vertices and an edge with source \texttt{x} and range \texttt{x\texttt{\symbol{94}}trans} for every \texttt{x} in \texttt{[1 .. n]}. See \texttt{AsDigraph} (\textbf{Digraphs: AsDigraph for a binary relation}). Then \texttt{CanonicalTransformation} returns a canonical representative of the class \texttt{C(n)} that contains \mbox{\texttt{\mdseries\slshape trans}}. 
\begin{Verbatim}[commandchars=!@|,fontsize=\small,frame=single,label=Example]
  !gapprompt@gap>| !gapinput@x := Transformation([5, 1, 4, 1, 1]);|
  Transformation( [ 5, 1, 4, 1, 1 ] )
  !gapprompt@gap>| !gapinput@y := Transformation([3, 3, 2, 3, 1]);|
  Transformation( [ 3, 3, 2, 3, 1 ] )
  !gapprompt@gap>| !gapinput@CanonicalTransformation(x);|
  Transformation( [ 3, 5, 2, 2, 2 ] )
  !gapprompt@gap>| !gapinput@CanonicalTransformation(y);|
  Transformation( [ 3, 5, 2, 2, 2 ] )
\end{Verbatim}
 }

 

\subsection{\textcolor{Chapter }{IsConnectedTransformationSemigroup (for a transformation semigroup)}}
\logpage{[ 11, 12, 10 ]}\nobreak
\hyperdef{L}{X82ABE03F80B8CA2B}{}
{\noindent\textcolor{FuncColor}{$\triangleright$\enspace\texttt{IsConnectedTransformationSemigroup({\mdseries\slshape S})\index{IsConnectedTransformationSemigroup@\texttt{IsConnectedTransformationSemigroup}!for a transformation semigroup}
\label{IsConnectedTransformationSemigroup:for a transformation semigroup}
}\hfill{\scriptsize (property)}}\\
\textbf{\indent Returns:\ }
\texttt{true} or \texttt{false}.



 A transformation semigroup \mbox{\texttt{\mdseries\slshape S}} is connected if the digraph returned by the function \texttt{DigraphOfActionOnPoints} is connected. See \texttt{IsConnectedDigraph} (\textbf{Digraphs: IsConnectedDigraph}) and \texttt{DigraphOfActionOnPoints} (\ref{DigraphOfActionOnPoints:for a transformation semigroup}). The function \texttt{IsConnectedTransformationSemigroup} returns \texttt{true} if the semigroup \mbox{\texttt{\mdseries\slshape S}} is connected and \texttt{false} otherwise. 
\begin{Verbatim}[commandchars=!@|,fontsize=\small,frame=single,label=Example]
  !gapprompt@gap>| !gapinput@S := Semigroup([|
  !gapprompt@>| !gapinput@ Transformation([2, 4, 3, 4]),|
  !gapprompt@>| !gapinput@ Transformation([3, 3, 2, 3, 3])]);|
  <transformation semigroup of degree 5 with 2 generators>
  !gapprompt@gap>| !gapinput@IsConnectedTransformationSemigroup(S);|
  true
\end{Verbatim}
 }

 }

   
\section{\textcolor{Chapter }{ Attributes of partial perms and partial perm semigroups }}\logpage{[ 11, 13, 0 ]}
\hyperdef{L}{X8669D0DA7EE1ECE0}{}
{
  

\subsection{\textcolor{Chapter }{ComponentRepsOfPartialPermSemigroup}}
\logpage{[ 11, 13, 1 ]}\nobreak
\hyperdef{L}{X7BC22CB47C7B5EBB}{}
{\noindent\textcolor{FuncColor}{$\triangleright$\enspace\texttt{ComponentRepsOfPartialPermSemigroup({\mdseries\slshape S})\index{ComponentRepsOfPartialPermSemigroup@\texttt{ComponentRepsOfPartialPermSemigroup}}
\label{ComponentRepsOfPartialPermSemigroup}
}\hfill{\scriptsize (attribute)}}\\
\textbf{\indent Returns:\ }
The representatives of components of a partial perm semigroup.



 This function returns the representatives of the components of the action of
the partial perm semigroup \mbox{\texttt{\mdseries\slshape S}} on the set of positive integers where it is defined. 

 The representatives are the least set of points such that every point can be
reached from some representative under the action of \mbox{\texttt{\mdseries\slshape S}}. 
\begin{Verbatim}[commandchars=!@|,fontsize=\small,frame=single,label=Example]
  !gapprompt@gap>| !gapinput@S := Semigroup([|
  !gapprompt@>| !gapinput@PartialPerm([1, 2, 3, 5, 6, 7, 8, 11, 12, 16, 19],|
  !gapprompt@>| !gapinput@            [9, 18, 20, 11, 5, 16, 8, 19, 14, 13, 1]),|
  !gapprompt@>| !gapinput@PartialPerm([1, 2, 3, 4, 5, 6, 8, 9, 10, 12, 14, 16, 18, 19, 20],|
  !gapprompt@>| !gapinput@            [13, 1, 8, 5, 4, 14, 11, 12, 9, 20, 2, 18, 7, 3, 19])]);;|
  !gapprompt@gap>| !gapinput@ComponentRepsOfPartialPermSemigroup(S);|
  [ 6, 10, 15, 17 ]
\end{Verbatim}
 }

 

\subsection{\textcolor{Chapter }{ComponentsOfPartialPermSemigroup}}
\logpage{[ 11, 13, 2 ]}\nobreak
\hyperdef{L}{X8464BC397ACBF2F1}{}
{\noindent\textcolor{FuncColor}{$\triangleright$\enspace\texttt{ComponentsOfPartialPermSemigroup({\mdseries\slshape S})\index{ComponentsOfPartialPermSemigroup@\texttt{ComponentsOfPartialPermSemigroup}}
\label{ComponentsOfPartialPermSemigroup}
}\hfill{\scriptsize (attribute)}}\\
\textbf{\indent Returns:\ }
The components of a partial perm semigroup.



 This function returns the components of the action of the partial perm
semigroup \mbox{\texttt{\mdseries\slshape S}} on the set of positive integers where it is defined; the components of \mbox{\texttt{\mdseries\slshape S}} partition this set. 
\begin{Verbatim}[commandchars=!@|,fontsize=\small,frame=single,label=Example]
  !gapprompt@gap>| !gapinput@S := Semigroup([|
  !gapprompt@>| !gapinput@PartialPerm([1, 2, 3, 5, 6, 7, 8, 11, 12, 16, 19],|
  !gapprompt@>| !gapinput@            [9, 18, 20, 11, 5, 16, 8, 19, 14, 13, 1]),|
  !gapprompt@>| !gapinput@PartialPerm([1, 2, 3, 4, 5, 6, 8, 9, 10, 12, 14, 16, 18, 19, 20],|
  !gapprompt@>| !gapinput@            [13, 1, 8, 5, 4, 14, 11, 12, 9, 20, 2, 18, 7, 3, 19])]);;|
  !gapprompt@gap>| !gapinput@ComponentsOfPartialPermSemigroup(S);|
  [ [ 1, 2, 3, 4, 5, 6, 7, 8, 9, 10, 11, 12, 13, 14, 16, 18, 19, 20 ],
    [ 15 ], [ 17 ] ]
\end{Verbatim}
 }

 

\subsection{\textcolor{Chapter }{CyclesOfPartialPerm}}
\logpage{[ 11, 13, 3 ]}\nobreak
\hyperdef{L}{X832937BB87EB4349}{}
{\noindent\textcolor{FuncColor}{$\triangleright$\enspace\texttt{CyclesOfPartialPerm({\mdseries\slshape x})\index{CyclesOfPartialPerm@\texttt{CyclesOfPartialPerm}}
\label{CyclesOfPartialPerm}
}\hfill{\scriptsize (attribute)}}\\
\textbf{\indent Returns:\ }
The cycles of a partial perm.



 This function returns the cycles, or strongly connected components, of the
action of the partial perm \mbox{\texttt{\mdseries\slshape x}} on the set of positive integers where it is defined. 
\begin{Verbatim}[commandchars=!@|,fontsize=\small,frame=single,label=Example]
  !gapprompt@gap>| !gapinput@x := PartialPerm([3, 1, 4, 2, 5, 0, 0, 6, 0, 7]);|
  [8,6][10,7](1,3,4,2)(5)
  !gapprompt@gap>| !gapinput@CyclesOfPartialPerm(x);|
  [ [ 5 ], [ 1, 3, 4, 2 ] ]
\end{Verbatim}
 }

 

\subsection{\textcolor{Chapter }{CyclesOfPartialPermSemigroup}}
\logpage{[ 11, 13, 4 ]}\nobreak
\hyperdef{L}{X7F7A5E5E8355E230}{}
{\noindent\textcolor{FuncColor}{$\triangleright$\enspace\texttt{CyclesOfPartialPermSemigroup({\mdseries\slshape S})\index{CyclesOfPartialPermSemigroup@\texttt{CyclesOfPartialPermSemigroup}}
\label{CyclesOfPartialPermSemigroup}
}\hfill{\scriptsize (attribute)}}\\
\textbf{\indent Returns:\ }
The cycles of a partial perm semigroup.



 This function returns the cycles, or strongly connected components, of the
action of the partial perm semigroup \mbox{\texttt{\mdseries\slshape S}} on the set of positive integers where it is defined. 
\begin{Verbatim}[commandchars=!@|,fontsize=\small,frame=single,label=Example]
  !gapprompt@gap>| !gapinput@S := Semigroup([|
  !gapprompt@>| !gapinput@PartialPerm([1, 2, 3, 5, 6, 7, 8, 11, 12, 16, 19],|
  !gapprompt@>| !gapinput@            [9, 18, 20, 11, 5, 16, 8, 19, 14, 13, 1]),|
  !gapprompt@>| !gapinput@PartialPerm([1, 2, 3, 4, 5, 6, 8, 9, 10, 12, 14, 16, 18, 19, 20],|
  !gapprompt@>| !gapinput@            [13, 1, 8, 5, 4, 14, 11, 12, 9, 20, 2, 18, 7, 3, 19])]);;|
  !gapprompt@gap>| !gapinput@CyclesOfPartialPermSemigroup(S);|
  [ [ 18, 7, 16 ], [ 1, 9, 12, 14, 2, 20, 19, 3, 8, 11 ], [ 4, 5 ] ]
\end{Verbatim}
 }

 

\subsection{\textcolor{Chapter }{UniformRandomPartialPerm}}
\logpage{[ 11, 13, 5 ]}\nobreak
\hyperdef{L}{X7E7003C07CB2D794}{}
{\noindent\textcolor{FuncColor}{$\triangleright$\enspace\texttt{UniformRandomPartialPerm({\mdseries\slshape n})\index{UniformRandomPartialPerm@\texttt{UniformRandomPartialPerm}}
\label{UniformRandomPartialPerm}
}\hfill{\scriptsize (operation)}}\\
\textbf{\indent Returns:\ }
 A partial perm of degree \mbox{\texttt{\mdseries\slshape n}} chosen uniformly at random among all partial permutations of that degree. 



 This function returns a partial perm of degree \mbox{\texttt{\mdseries\slshape n}} chosen uniformly at random among all partial permutations of that degree. This
function differs from \texttt{RandomPartialPerm} (\textbf{Reference: RandomPartialPerm}) in that the returned partial perm is chosen \emph{uniformly} at random, whereas \texttt{RandomPartialPerm} (\textbf{Reference: RandomPartialPerm}) does not use a uniform distribution. 

 The downside to \texttt{UniformRandomPartialPerm} is that it can use a significant amount of memory, and is rather slower than \texttt{RandomPartialPerm} (\textbf{Reference: RandomPartialPerm}). 
\begin{Verbatim}[commandchars=!@|,fontsize=\small,frame=single,label=Example]
  !gapprompt@gap>| !gapinput@UniformRandomPartialPerm(1000);|
  <partial perm on 971 pts with degree 1000, codegree 1000>
      
\end{Verbatim}
 }

 }

   
\section{\textcolor{Chapter }{ Attributes of Rees (0\texttt{\symbol{45}})matrix semigroups }}\logpage{[ 11, 14, 0 ]}
\hyperdef{L}{X7AF313CF7CBE98D7}{}
{
  

\subsection{\textcolor{Chapter }{RZMSDigraph}}
\logpage{[ 11, 14, 1 ]}\nobreak
\hyperdef{L}{X7EA1B28785B9D38C}{}
{\noindent\textcolor{FuncColor}{$\triangleright$\enspace\texttt{RZMSDigraph({\mdseries\slshape R})\index{RZMSDigraph@\texttt{RZMSDigraph}}
\label{RZMSDigraph}
}\hfill{\scriptsize (attribute)}}\\
\textbf{\indent Returns:\ }
A digraph.



 If \mbox{\texttt{\mdseries\slshape R}} is an $n$ by $m$ Rees 0\texttt{\symbol{45}}matrix semigroup $M^{0}[I, T, \Lambda; P]$ (so that $I = \{1,2,\ldots,n\}$ and $\Lambda = \{1,2,\ldots,m\}$) then \texttt{RZMSDigraph} returns a symmetric bipartite digraph with $n+m$ vertices. An index $i \in I$ corresponds to the vertex $i$ and an index $j \in \Lambda$ corresponds to the vertex $j + n$. 

 Two vertices $v$ and $w$ in \texttt{RZMSDigraph(}\mbox{\texttt{\mdseries\slshape R}}\texttt{)} are adjacent if and only if $v\in I$, $w - n\in \Lambda$, and \texttt{P[w \texttt{\symbol{45}} n][v]} $\neq 0$. 

 This digraph is commonly called the \emph{Graham\texttt{\symbol{45}}Houghton graph} of \mbox{\texttt{\mdseries\slshape R}}. 
\begin{Verbatim}[commandchars=!@|,fontsize=\small,frame=single,label=Example]
  !gapprompt@gap>| !gapinput@R := PrincipalFactor(|
  !gapprompt@>| !gapinput@DClass(FullTransformationMonoid(5),|
  !gapprompt@>| !gapinput@       Transformation([2, 4, 1, 5, 5])));|
  <Rees 0-matrix semigroup 10x5 over Group([ (1,2,3,4), (1,2) ])>
  !gapprompt@gap>| !gapinput@gr := RZMSDigraph(R);|
  <immutable bipartite digraph with bicomponent sizes 10 and 5>
  !gapprompt@gap>| !gapinput@e := DigraphEdges(gr)[1];|
  [ 1, 11 ]
  !gapprompt@gap>| !gapinput@Matrix(R)[e[2] - 10][e[1]] <> 0;|
  true
\end{Verbatim}
 }

 

\subsection{\textcolor{Chapter }{RZMSConnectedComponents}}
\logpage{[ 11, 14, 2 ]}\nobreak
\hyperdef{L}{X79B062917AB34542}{}
{\noindent\textcolor{FuncColor}{$\triangleright$\enspace\texttt{RZMSConnectedComponents({\mdseries\slshape R})\index{RZMSConnectedComponents@\texttt{RZMSConnectedComponents}}
\label{RZMSConnectedComponents}
}\hfill{\scriptsize (attribute)}}\\
\textbf{\indent Returns:\ }
The connected components of a Rees 0\texttt{\symbol{45}}matrix semigroup.



 If \mbox{\texttt{\mdseries\slshape R}} is an $n$ by $m$ Rees 0\texttt{\symbol{45}}matrix semigroup $M^{0}[I, T, \Lambda; P]$ (so that $I = \{1,2,\ldots,n\}$ and $\Lambda = \{1,2,\ldots,m\}$) then \texttt{RZMSConnectedComponents} returns the connected components of \mbox{\texttt{\mdseries\slshape R}}. 

 \emph{Connectedness} is an equivalence relation on the indices of \mbox{\texttt{\mdseries\slshape R}}: the equivalence classes of the relation are called the \emph{connected components} of \mbox{\texttt{\mdseries\slshape R}}, and two indices in $I \cup \Lambda$ are connected if and only if their corresponding vertices in \texttt{RZMSDigraph(}\mbox{\texttt{\mdseries\slshape R}}\texttt{)} are connected (see \texttt{RZMSDigraph} (\ref{RZMSDigraph})). If \mbox{\texttt{\mdseries\slshape R}} has $n$ connected components, then \texttt{RZMSConnectedComponents} will return a list of pairs: 

 \texttt{[ [ }$I_{1}, \Lambda_{1}$\texttt{ ], }$\ldots$\texttt{, [ }$I_{k}, \Lambda_{k}$\texttt{ ] ]} 

 where $I = I_1 \sqcup \cdots \sqcup I_k$, $\Lambda = \Lambda_1 \sqcup \cdots \sqcup \Lambda_k$, and for each $l$ the set $I_{l}\cup\Lambda_{l}$ is a connected component of \mbox{\texttt{\mdseries\slshape R}}. Note that at most one of $I_l$ and $\Lambda_l$ is possibly empty. The ordering of the connected components in the result in
unspecified. 
\begin{Verbatim}[commandchars=!@|,fontsize=\small,frame=single,label=Example]
  !gapprompt@gap>| !gapinput@R := ReesZeroMatrixSemigroup(SymmetricGroup(5),|
  !gapprompt@>| !gapinput@[[(), 0, (1, 3), (4, 5), 0],|
  !gapprompt@>| !gapinput@ [0, (), 0, 0, (1, 3, 4, 5)],|
  !gapprompt@>| !gapinput@ [0, 0, (1, 5)(2, 3), 0, 0],|
  !gapprompt@>| !gapinput@ [0, (2, 3)(1, 4), 0, 0, 0]]);|
  <Rees 0-matrix semigroup 5x4 over Sym( [ 1 .. 5 ] )>
  !gapprompt@gap>| !gapinput@RZMSConnectedComponents(R);|
  [ [ [ 1, 3, 4 ], [ 1, 3 ] ], [ [ 2, 5 ], [ 2, 4 ] ] ]
\end{Verbatim}
 }

 }

 
\section{\textcolor{Chapter }{ Attributes of inverse semigroups }}\logpage{[ 11, 15, 0 ]}
\hyperdef{L}{X822D030682BC1275}{}
{
  

\subsection{\textcolor{Chapter }{NaturalLeqInverseSemigroup}}
\logpage{[ 11, 15, 1 ]}\nobreak
\hyperdef{L}{X7A75A6C486F1DC71}{}
{\noindent\textcolor{FuncColor}{$\triangleright$\enspace\texttt{NaturalLeqInverseSemigroup({\mdseries\slshape S})\index{NaturalLeqInverseSemigroup@\texttt{NaturalLeqInverseSemigroup}}
\label{NaturalLeqInverseSemigroup}
}\hfill{\scriptsize (attribute)}}\\
\textbf{\indent Returns:\ }
 An function. 



 \texttt{NaturalLeqInverseSemigroup} returns a function that, when given two elements \texttt{x, y} of the inverse semigroup \mbox{\texttt{\mdseries\slshape S}}, returns \texttt{true} if \texttt{x} is less than or equal to \texttt{y} in the natural partial order on \mbox{\texttt{\mdseries\slshape S}}. 
\begin{Verbatim}[commandchars=!@|,fontsize=\small,frame=single,label=Example]
  !gapprompt@gap>| !gapinput@S := Monoid(Transformation([1, 3, 4, 4]),|
  !gapprompt@>| !gapinput@               Transformation([1, 4, 2, 4]));|
  <transformation monoid of degree 4 with 2 generators>
  !gapprompt@gap>| !gapinput@IsInverseSemigroup(S);|
  true
  !gapprompt@gap>| !gapinput@Size(S);|
  6
  !gapprompt@gap>| !gapinput@NaturalPartialOrder(S);|
  [ [ 2, 5, 6 ], [ 6 ], [ 6 ], [ 6 ], [ 6 ], [  ] ]
\end{Verbatim}
 }

 

\subsection{\textcolor{Chapter }{JoinIrreducibleDClasses}}
\logpage{[ 11, 15, 2 ]}\nobreak
\hyperdef{L}{X85CDF93C805AF82A}{}
{\noindent\textcolor{FuncColor}{$\triangleright$\enspace\texttt{JoinIrreducibleDClasses({\mdseries\slshape S})\index{JoinIrreducibleDClasses@\texttt{JoinIrreducibleDClasses}}
\label{JoinIrreducibleDClasses}
}\hfill{\scriptsize (attribute)}}\\
\textbf{\indent Returns:\ }
 A list of $\mathcal{D}$\texttt{\symbol{45}}classes. 



 \texttt{JoinIrreducibleDClasses} returns a list of the join irreducible $\mathcal{D}$\texttt{\symbol{45}}classes of the inverse semigroup of partial permutations,
block bijections or partial permutation bipartitions \mbox{\texttt{\mdseries\slshape S}}.

 A \emph{join irreducible $\mathcal{D}$\texttt{\symbol{45}}class} is a $\mathcal{D}$\texttt{\symbol{45}}class containing only join irreducible elements. See \texttt{IsJoinIrreducible} (\ref{IsJoinIrreducible}). If a $\mathcal{D}$\texttt{\symbol{45}}class contains one join irreducible element, then all of
the elements in the $\mathcal{D}$\texttt{\symbol{45}}class are join irreducible. 
\begin{Verbatim}[commandchars=!@|,fontsize=\small,frame=single,label=Example]
  !gapprompt@gap>| !gapinput@S := SymmetricInverseSemigroup(3);|
  <symmetric inverse monoid of degree 3>
  !gapprompt@gap>| !gapinput@JoinIrreducibleDClasses(S);|
  [ <Green's D-class: <identity partial perm on [ 2 ]>> ]
  !gapprompt@gap>| !gapinput@T := InverseSemigroup([|
  !gapprompt@>| !gapinput@  PartialPerm([1, 2, 4, 3]),|
  !gapprompt@>| !gapinput@  PartialPerm([1]),|
  !gapprompt@>| !gapinput@  PartialPerm([0, 2])]);|
  <inverse partial perm semigroup of rank 4 with 3 generators>
  !gapprompt@gap>| !gapinput@JoinIrreducibleDClasses(T);|
  [ <Green's D-class: <identity partial perm on [ 1, 2, 3, 4 ]>>,
    <Green's D-class: <identity partial perm on [ 1 ]>>,
    <Green's D-class: <identity partial perm on [ 2 ]>> ]
  !gapprompt@gap>| !gapinput@D := DualSymmetricInverseSemigroup(3);|
  <inverse block bijection monoid of degree 3 with 3 generators>
  !gapprompt@gap>| !gapinput@JoinIrreducibleDClasses(D);|
  [ <Green's D-class: <block bijection: [ 1, 2, -1, -2 ], [ 3, -3 ]>> ]
\end{Verbatim}
 }

 

\subsection{\textcolor{Chapter }{MajorantClosure}}
\logpage{[ 11, 15, 3 ]}\nobreak
\hyperdef{L}{X801CC67E80898608}{}
{\noindent\textcolor{FuncColor}{$\triangleright$\enspace\texttt{MajorantClosure({\mdseries\slshape S, T})\index{MajorantClosure@\texttt{MajorantClosure}}
\label{MajorantClosure}
}\hfill{\scriptsize (operation)}}\\
\textbf{\indent Returns:\ }
 A majorantly closed list of elements. 



 \texttt{MajorantClosure} returns a majorantly closed subset of an inverse semigroup of partial
permutations, block bijections or partial permutation bipartitions, \mbox{\texttt{\mdseries\slshape S}}, as a list. See \texttt{IsMajorantlyClosed} (\ref{IsMajorantlyClosed}).

 The result contains all elements of \mbox{\texttt{\mdseries\slshape S}} which are greater than or equal to any element of \mbox{\texttt{\mdseries\slshape T}} (with respect to the natural partial order \texttt{NaturalLeqPartialPerm} (\textbf{Reference: NaturalLeqPartialPerm})). In particular, the result is a superset of \mbox{\texttt{\mdseries\slshape T}}.

 Note that \mbox{\texttt{\mdseries\slshape T}} can be a subset of \mbox{\texttt{\mdseries\slshape S}} or a subsemigroup of \mbox{\texttt{\mdseries\slshape S}}. }

 
\begin{Verbatim}[commandchars=!@|,fontsize=\small,frame=single,label=Example]
  !gapprompt@gap>| !gapinput@S := SymmetricInverseSemigroup(4);|
  <symmetric inverse monoid of degree 4>
  !gapprompt@gap>| !gapinput@T := [PartialPerm([1, 0, 3, 0])];|
  [ <identity partial perm on [ 1, 3 ]> ]
  !gapprompt@gap>| !gapinput@U := MajorantClosure(S, T);|
  [ <identity partial perm on [ 1, 3 ]>,
    <identity partial perm on [ 1, 2, 3 ]>, [2,4](1)(3), [4,2](1)(3),
    <identity partial perm on [ 1, 3, 4 ]>,
    <identity partial perm on [ 1, 2, 3, 4 ]>, (1)(2,4)(3) ]
  !gapprompt@gap>| !gapinput@B := InverseSemigroup([|
  !gapprompt@>| !gapinput@ Bipartition([[1, -2], [2, -1], [3, -3], [4, 5, -4, -5]]),|
  !gapprompt@>| !gapinput@ Bipartition([[1, -3], [2, -4], [3, -2], [4, -1], [5, -5]])]);;|
  !gapprompt@gap>| !gapinput@T := [Bipartition([[1, -2], [2, 3, 5, -1, -3, -5], [4, -4]]),|
  !gapprompt@>| !gapinput@ Bipartition([[1, -4], [2, 3, 5, -1, -3, -5], [4, -2]])];;|
  !gapprompt@gap>| !gapinput@IsMajorantlyClosed(B, T);|
  false
  !gapprompt@gap>| !gapinput@MajorantClosure(B, T);|
  [ <block bijection: [ 1, -2 ], [ 2, 3, 5, -1, -3, -5 ], [ 4, -4 ]>,
    <block bijection: [ 1, -4 ], [ 2, 3, 5, -1, -3, -5 ], [ 4, -2 ]>,
    <block bijection: [ 1, -2 ], [ 2, 5, -1, -5 ], [ 3, -3 ], [ 4, -4 ]>
      , <block bijection: [ 1, -2 ], [ 2, -1 ], [ 3, 5, -3, -5 ],
       [ 4, -4 ]>,
    <block bijection: [ 1, -4 ], [ 2, 5, -3, -5 ], [ 3, -1 ], [ 4, -2 ]>
      , <block bijection: [ 1, -4 ], [ 2, -3 ], [ 3, 5, -1, -5 ],
       [ 4, -2 ]>, <block bijection: [ 1, -4 ], [ 2, -3 ], [ 3, -1 ],
       [ 4, -2 ], [ 5, -5 ]> ]
  !gapprompt@gap>| !gapinput@IsMajorantlyClosed(B, last);|
  true
\end{Verbatim}
 

\subsection{\textcolor{Chapter }{Minorants}}
\logpage{[ 11, 15, 4 ]}\nobreak
\hyperdef{L}{X84A3DB79816374DB}{}
{\noindent\textcolor{FuncColor}{$\triangleright$\enspace\texttt{Minorants({\mdseries\slshape S, f})\index{Minorants@\texttt{Minorants}}
\label{Minorants}
}\hfill{\scriptsize (operation)}}\\
\textbf{\indent Returns:\ }
 A list of elements. 



 \texttt{Minorants} takes an element \mbox{\texttt{\mdseries\slshape f}} from an inverse semigroup of partial permutations, block bijections or partial
permutation bipartitions \mbox{\texttt{\mdseries\slshape S}}, and returns a list of the minorants of \mbox{\texttt{\mdseries\slshape f}} in \mbox{\texttt{\mdseries\slshape S}}. 

 A \emph{minorant} of \mbox{\texttt{\mdseries\slshape f}} is an element of \mbox{\texttt{\mdseries\slshape S}} which is strictly less than \mbox{\texttt{\mdseries\slshape f}} in the natural partial order of \mbox{\texttt{\mdseries\slshape S}}. See \texttt{NaturalLeqPartialPerm} (\textbf{Reference: NaturalLeqPartialPerm}). }

 
\begin{Verbatim}[commandchars=!@|,fontsize=\small,frame=single,label=Example]
  !gapprompt@gap>| !gapinput@S := SymmetricInverseSemigroup(3);|
  <symmetric inverse monoid of degree 3>
  !gapprompt@gap>| !gapinput@x := Elements(S)[13];|
  [1,3](2)
  !gapprompt@gap>| !gapinput@Minorants(S, x);|
  [ <empty partial perm>, [1,3], <identity partial perm on [ 2 ]> ]
  !gapprompt@gap>| !gapinput@x := PartialPerm([3, 2, 4, 0]);|
  [1,3,4](2)
  !gapprompt@gap>| !gapinput@S := InverseSemigroup(x);|
  <inverse partial perm semigroup of rank 4 with 1 generator>
  !gapprompt@gap>| !gapinput@Minorants(S, x);|
  [ <identity partial perm on [ 2 ]>, [1,3](2), [3,4](2) ]
\end{Verbatim}
 

\subsection{\textcolor{Chapter }{PrimitiveIdempotents}}
\logpage{[ 11, 15, 5 ]}\nobreak
\hyperdef{L}{X80C0C6C37C4A2ABD}{}
{\noindent\textcolor{FuncColor}{$\triangleright$\enspace\texttt{PrimitiveIdempotents({\mdseries\slshape S})\index{PrimitiveIdempotents@\texttt{PrimitiveIdempotents}}
\label{PrimitiveIdempotents}
}\hfill{\scriptsize (attribute)}}\\
\textbf{\indent Returns:\ }
A list of elements.



 An idempotent in an inverse semigroup \mbox{\texttt{\mdseries\slshape S}} is \emph{primitive} if it is non\texttt{\symbol{45}}zero and minimal with respect to the \texttt{NaturalPartialOrder} (\textbf{Reference: NaturalPartialOrder}) on \mbox{\texttt{\mdseries\slshape S}}. \texttt{PrimitiveIdempotents} returns the list of primitive idempotents in the inverse semigroup \mbox{\texttt{\mdseries\slshape S}}. 
\begin{Verbatim}[commandchars=!@|,fontsize=\small,frame=single,label=Example]
  !gapprompt@gap>| !gapinput@S := InverseMonoid(|
  !gapprompt@>| !gapinput@PartialPerm([1], [4]),|
  !gapprompt@>| !gapinput@PartialPerm([1, 2, 3], [2, 1, 3]),|
  !gapprompt@>| !gapinput@PartialPerm([1, 2, 3], [3, 1, 2]));;|
  !gapprompt@gap>| !gapinput@MultiplicativeZero(S);|
  <empty partial perm>
  !gapprompt@gap>| !gapinput@Set(PrimitiveIdempotents(S));|
  [ <identity partial perm on [ 1 ]>, <identity partial perm on [ 2 ]>,
    <identity partial perm on [ 3 ]>, <identity partial perm on [ 4 ]> ]
  !gapprompt@gap>| !gapinput@S := DualSymmetricInverseMonoid(4);|
  <inverse block bijection monoid of degree 4 with 3 generators>
  !gapprompt@gap>| !gapinput@Set(PrimitiveIdempotents(S));|
  [ <block bijection: [ 1, 2, 3, -1, -2, -3 ], [ 4, -4 ]>,
    <block bijection: [ 1, 2, 4, -1, -2, -4 ], [ 3, -3 ]>,
    <block bijection: [ 1, 2, -1, -2 ], [ 3, 4, -3, -4 ]>,
    <block bijection: [ 1, 3, 4, -1, -3, -4 ], [ 2, -2 ]>,
    <block bijection: [ 1, 3, -1, -3 ], [ 2, 4, -2, -4 ]>,
    <block bijection: [ 1, 4, -1, -4 ], [ 2, 3, -2, -3 ]>,
    <block bijection: [ 1, -1 ], [ 2, 3, 4, -2, -3, -4 ]> ]
\end{Verbatim}
 }

 

\subsection{\textcolor{Chapter }{RightCosetsOfInverseSemigroup}}
\logpage{[ 11, 15, 6 ]}\nobreak
\hyperdef{L}{X7B5D89A585F8B5EA}{}
{\noindent\textcolor{FuncColor}{$\triangleright$\enspace\texttt{RightCosetsOfInverseSemigroup({\mdseries\slshape S, T})\index{RightCosetsOfInverseSemigroup@\texttt{RightCosetsOfInverseSemigroup}}
\label{RightCosetsOfInverseSemigroup}
}\hfill{\scriptsize (operation)}}\\
\textbf{\indent Returns:\ }
 A list of lists of elements. 



 \texttt{RightCosetsOfInverseSemigroup} takes a majorantly closed inverse subsemigroup \mbox{\texttt{\mdseries\slshape T}} of an inverse semigroup of partial permutations, block bijections or partial
permutation bipartitions \mbox{\texttt{\mdseries\slshape S}}. See \texttt{IsMajorantlyClosed} (\ref{IsMajorantlyClosed}). The result is a list of the right cosets of \mbox{\texttt{\mdseries\slshape T}} in \mbox{\texttt{\mdseries\slshape S}}.

 For $s \in S$, the right coset $\overline{Ts}$ is defined if and only if $ss^{-1} \in T$, in which case it is defined to be the majorant closure of the set $Ts$. See \texttt{MajorantClosure} (\ref{MajorantClosure}). Distinct cosets are disjoint but do not necessarily partition \mbox{\texttt{\mdseries\slshape S}}. }

 
\begin{Verbatim}[commandchars=!@|,fontsize=\small,frame=single,label=Example]
  !gapprompt@gap>| !gapinput@S := SymmetricInverseSemigroup(3);|
  <symmetric inverse monoid of degree 3>
  !gapprompt@gap>| !gapinput@T := InverseSemigroup(MajorantClosure(S, [PartialPerm([1])]));|
  <inverse partial perm monoid of rank 3 with 6 generators>
  !gapprompt@gap>| !gapinput@IsMajorantlyClosed(S, T);|
  true
  !gapprompt@gap>| !gapinput@RC := RightCosetsOfInverseSemigroup(S, T);|
  [ [ <identity partial perm on [ 1 ]>,
        <identity partial perm on [ 1, 2 ]>, [2,3](1), [3,2](1),
        <identity partial perm on [ 1, 3 ]>,
        <identity partial perm on [ 1, 2, 3 ]>, (1)(2,3) ],
    [ [1,3], [2,1,3], [1,3](2), (1,3), [1,3,2], (1,3,2), (1,3)(2) ],
    [ [1,2], (1,2), [1,2,3], [3,1,2], [1,2](3), (1,2)(3), (1,2,3) ] ]
\end{Verbatim}
 

\subsection{\textcolor{Chapter }{SameMinorantsSubgroup}}
\logpage{[ 11, 15, 7 ]}\nobreak
\hyperdef{L}{X83298E9E86A343FF}{}
{\noindent\textcolor{FuncColor}{$\triangleright$\enspace\texttt{SameMinorantsSubgroup({\mdseries\slshape H})\index{SameMinorantsSubgroup@\texttt{SameMinorantsSubgroup}}
\label{SameMinorantsSubgroup}
}\hfill{\scriptsize (attribute)}}\\
\textbf{\indent Returns:\ }
 A list of elements of the group $\mathcal{H}$\texttt{\symbol{45}}class \mbox{\texttt{\mdseries\slshape H}}. 



 Given a group $\mathcal{H}$\texttt{\symbol{45}}class \mbox{\texttt{\mdseries\slshape H}} in an inverse semigroup of partial permutations, block bijections or partial
permutation bipartitions \texttt{S}, \texttt{SameMinorantsSubgroup} returns a list of the elements of \mbox{\texttt{\mdseries\slshape H}} which have the same strict minorants as the identity element of \mbox{\texttt{\mdseries\slshape H}}. A \emph{strict minorant} of \texttt{x} in \mbox{\texttt{\mdseries\slshape H}} is an element of \texttt{S} which is less than \texttt{x} (with respect to the natural partial order), but is not equal to \texttt{x}.

 The returned list of elements of \mbox{\texttt{\mdseries\slshape H}} describe a subgroup of \mbox{\texttt{\mdseries\slshape H}}. }

 
\begin{Verbatim}[commandchars=!@|,fontsize=\small,frame=single,label=Example]
  !gapprompt@gap>| !gapinput@S := SymmetricInverseSemigroup(3);|
  <symmetric inverse monoid of degree 3>
  !gapprompt@gap>| !gapinput@H := GroupHClass(DClass(S, PartialPerm([1, 2, 3])));|
  <Green's H-class: <identity partial perm on [ 1, 2, 3 ]>>
  !gapprompt@gap>| !gapinput@Elements(H);|
  [ <identity partial perm on [ 1, 2, 3 ]>, (1)(2,3), (1,2)(3),
    (1,2,3), (1,3,2), (1,3)(2) ]
  !gapprompt@gap>| !gapinput@SameMinorantsSubgroup(H);|
  [ <identity partial perm on [ 1, 2, 3 ]> ]
  !gapprompt@gap>| !gapinput@T := InverseSemigroup(|
  !gapprompt@>| !gapinput@PartialPerm([1, 2, 3, 4], [1, 2, 4, 3]),|
  !gapprompt@>| !gapinput@PartialPerm([1], [1]), PartialPerm([2], [2]));|
  <inverse partial perm semigroup of rank 4 with 3 generators>
  !gapprompt@gap>| !gapinput@Elements(T);|
  [ <empty partial perm>, <identity partial perm on [ 1 ]>,
    <identity partial perm on [ 2 ]>,
    <identity partial perm on [ 1, 2, 3, 4 ]>, (1)(2)(3,4) ]
  !gapprompt@gap>| !gapinput@x := GroupHClass(DClass(T, PartialPerm([1, 2, 3, 4])));|
  <Green's H-class: <identity partial perm on [ 1, 2, 3, 4 ]>>
  !gapprompt@gap>| !gapinput@Elements(x);|
  [ <identity partial perm on [ 1, 2, 3, 4 ]>, (1)(2)(3,4) ]
  !gapprompt@gap>| !gapinput@AsSet(SameMinorantsSubgroup(x));|
  [ <identity partial perm on [ 1, 2, 3, 4 ]>, (1)(2)(3,4) ]
\end{Verbatim}
 

\subsection{\textcolor{Chapter }{SmallerDegreePartialPermRepresentation}}
\logpage{[ 11, 15, 8 ]}\nobreak
\hyperdef{L}{X786D4E397EA4445D}{}
{\noindent\textcolor{FuncColor}{$\triangleright$\enspace\texttt{SmallerDegreePartialPermRepresentation({\mdseries\slshape S})\index{SmallerDegreePartialPermRepresentation@\texttt{Smaller}\-\texttt{Degree}\-\texttt{Partial}\-\texttt{Perm}\-\texttt{Representation}}
\label{SmallerDegreePartialPermRepresentation}
}\hfill{\scriptsize (attribute)}}\\
\textbf{\indent Returns:\ }
 An isomorphism. 



 \texttt{SmallerDegreePartialPermRepresentation} attempts to find an isomorphism from the inverse semigroup \mbox{\texttt{\mdseries\slshape S}} to an inverse semigroup of partial permutations with small degree. If \mbox{\texttt{\mdseries\slshape S}} is already a partial permutation semigroup, and the function cannot reduce the
degree, the identity mapping is returned. 

 There is no guarantee that the smallest possible degree representation is
returned. For more information see \cite{Schein1992aa}. }

 
\begin{Verbatim}[commandchars=!@|,fontsize=\small,frame=single,label=Example]
  !gapprompt@gap>| !gapinput@S := InverseSemigroup(PartialPerm([2, 1, 4, 3, 6, 5, 8, 7]));|
  <partial perm group of rank 8 with 1 generator>
  !gapprompt@gap>| !gapinput@Elements(S);|
  [ <identity partial perm on [ 1, 2, 3, 4, 5, 6, 7, 8 ]>,
    (1,2)(3,4)(5,6)(7,8) ]
  !gapprompt@gap>| !gapinput@iso := SmallerDegreePartialPermRepresentation(S);;|
  !gapprompt@gap>| !gapinput@Source(iso) = S;|
  true
  !gapprompt@gap>| !gapinput@R := Range(iso);|
  <partial perm group of rank 2 with 1 generator>
  !gapprompt@gap>| !gapinput@Elements(R);|
  [ <identity partial perm on [ 1, 2 ]>, (1,2) ]
  !gapprompt@gap>| !gapinput@S := DualSymmetricInverseMonoid(5);;|
  !gapprompt@gap>| !gapinput@T := Range(IsomorphismPartialPermSemigroup(S));|
  <inverse partial perm monoid of size 6721, rank 6721 with 3
   generators>
  !gapprompt@gap>| !gapinput@SmallerDegreePartialPermRepresentation(T);|
  <inverse partial perm monoid of size 6721, rank 6721 with 3
    generators> -> <inverse partial perm monoid of rank 30 with 3
    generators>
\end{Verbatim}
 

\subsection{\textcolor{Chapter }{VagnerPrestonRepresentation}}
\logpage{[ 11, 15, 9 ]}\nobreak
\hyperdef{L}{X7BC49C3487384364}{}
{\noindent\textcolor{FuncColor}{$\triangleright$\enspace\texttt{VagnerPrestonRepresentation({\mdseries\slshape S})\index{VagnerPrestonRepresentation@\texttt{VagnerPrestonRepresentation}}
\label{VagnerPrestonRepresentation}
}\hfill{\scriptsize (attribute)}}\\
\textbf{\indent Returns:\ }
 An isomorphism to an inverse semigroup of partial permutations. 



 \texttt{VagnerPrestonRepresentation} returns an isomorphism from an inverse semigroup \mbox{\texttt{\mdseries\slshape S}} where the elements of \mbox{\texttt{\mdseries\slshape S}} have a unique semigroup inverse accessible via \texttt{Inverse} (\textbf{Reference: Inverse}), to the inverse semigroup of partial permutations \mbox{\texttt{\mdseries\slshape T}} of degree equal to the size of \mbox{\texttt{\mdseries\slshape S}}, which is obtained using the Vagner\texttt{\symbol{45}}Preston Representation
Theorem. 

 More precisely, if $f:S\to T$ is the isomorphism returned by \texttt{VagnerPrestonRepresentation(\mbox{\texttt{\mdseries\slshape S}})} and $x$ is in \mbox{\texttt{\mdseries\slshape S}}, then $f(x)$ is the partial permutation with domain $Sx^{-1}$ and range $Sx^{-1}x$ defined by $f(x): sx^{-1}\mapsto sx^{-1}x$. 

 In many cases, it is possible to find a smaller degree representation than
that provided by \texttt{VagnerPrestonRepresentation} using \texttt{IsomorphismPartialPermSemigroup} (\textbf{Reference: IsomorphismPartialPermSemigroup}) or \texttt{SmallerDegreePartialPermRepresentation} (\ref{SmallerDegreePartialPermRepresentation}). }

 
\begin{Verbatim}[commandchars=!@|,fontsize=\small,frame=single,label=Example]
  !gapprompt@gap>| !gapinput@S := SymmetricInverseSemigroup(2);|
  <symmetric inverse monoid of degree 2>
  !gapprompt@gap>| !gapinput@Size(S);|
  7
  !gapprompt@gap>| !gapinput@iso := VagnerPrestonRepresentation(S);|
  <symmetric inverse monoid of degree 2> ->
  <inverse partial perm monoid of rank 7 with 2 generators>
  !gapprompt@gap>| !gapinput@RespectsMultiplication(iso);|
  true
  !gapprompt@gap>| !gapinput@inv := InverseGeneralMapping(iso);;|
  !gapprompt@gap>| !gapinput@ForAll(S, x -> (x ^ iso) ^ inv = x);|
  true
  !gapprompt@gap>| !gapinput@V := InverseSemigroup(|
  !gapprompt@>| !gapinput@Bipartition([[1, -4], [2, -1], [3, -5], [4], [5], [-2], [-3]]),|
  !gapprompt@>| !gapinput@Bipartition([[1, -5], [2, -1], [3, -3], [4], [5], [-2], [-4]]),|
  !gapprompt@>| !gapinput@Bipartition([[1, -2], [2, -4], [3, -5], [4, -1], [5, -3]]));|
  <inverse bipartition semigroup of degree 5 with 3 generators>
  !gapprompt@gap>| !gapinput@IsInverseSemigroup(V);|
  true
  !gapprompt@gap>| !gapinput@VagnerPrestonRepresentation(V);|
  <inverse bipartition semigroup of size 394, degree 5 with 3
    generators> -> <inverse partial perm semigroup of rank 394 with 5
    generators>
\end{Verbatim}
 

\subsection{\textcolor{Chapter }{CharacterTableOfInverseSemigroup}}
\logpage{[ 11, 15, 10 ]}\nobreak
\hyperdef{L}{X7C83DF9A7973AF6D}{}
{\noindent\textcolor{FuncColor}{$\triangleright$\enspace\texttt{CharacterTableOfInverseSemigroup({\mdseries\slshape S})\index{CharacterTableOfInverseSemigroup@\texttt{CharacterTableOfInverseSemigroup}}
\label{CharacterTableOfInverseSemigroup}
}\hfill{\scriptsize (attribute)}}\\
\textbf{\indent Returns:\ }
 The character table of the inverse semigroup \mbox{\texttt{\mdseries\slshape S}} and a list of conjugacy class representatives of \mbox{\texttt{\mdseries\slshape S}}. 



 Returns a list with two entries: the first entry being the character table of
the inverse semigroup \mbox{\texttt{\mdseries\slshape S}} as a matrix, while the second entry is a list of conjugacy class
representatives of \mbox{\texttt{\mdseries\slshape S}}. 

 The order of the columns in the character table matrix follows the order of
the conjugacy class representatives list. The conjugacy representatives are
grouped by $\mathcal{D}$\texttt{\symbol{45}}class and then sorted by rank. Also, as is typical of
character tables, the rows of the matrix correspond to the irreducible
characters and the columns correspond to the conjugacy classes. 

 This function was contributed by Jhevon Smith and Ben Steinberg. 
\begin{Verbatim}[commandchars=!@|,fontsize=\small,frame=single,label=Example]
  !gapprompt@gap>| !gapinput@S := InverseMonoid([|
  !gapprompt@>| !gapinput@  PartialPerm([1, 2], [3, 1]),|
  !gapprompt@>| !gapinput@ PartialPerm([1, 2, 3], [1, 3, 4]),|
  !gapprompt@>| !gapinput@ PartialPerm([1, 2, 3], [2, 4, 1]),|
  !gapprompt@>| !gapinput@ PartialPerm([1, 3, 4], [3, 4, 1])]);;|
  !gapprompt@gap>| !gapinput@CharacterTableOfInverseSemigroup(S);|
  [ [ [ 1, 0, 0, 0, 0, 0, 0, 0 ], [ 3, 1, 1, 1, 0, 0, 0, 0 ],
        [ 3, 1, E(3), E(3)^2, 0, 0, 0, 0 ],
        [ 3, 1, E(3)^2, E(3), 0, 0, 0, 0 ], [ 6, 3, 0, 0, 1, -1, 0, 0 ],
        [ 6, 3, 0, 0, 1, 1, 0, 0 ], [ 4, 3, 0, 0, 2, 0, 1, 0 ],
        [ 1, 1, 1, 1, 1, 1, 1, 1 ] ],
    [ <identity partial perm on [ 1, 2, 3, 4 ]>,
        <identity partial perm on [ 1, 3, 4 ]>, (1,3,4), (1,4,3),
        <identity partial perm on [ 1, 3 ]>, (1,3),
        <identity partial perm on [ 3 ]>, <empty partial perm> ] ]
  !gapprompt@gap>| !gapinput@S := SymmetricInverseMonoid(4);;|
  !gapprompt@gap>| !gapinput@CharacterTableOfInverseSemigroup(S);|
  [ [ [ 1, -1, 1, 1, -1, 0, 0, 0, 0, 0, 0, 0 ],
        [ 3, -1, 0, -1, 1, 0, 0, 0, 0, 0, 0, 0 ],
        [ 2, 0, -1, 2, 0, 0, 0, 0, 0, 0, 0, 0 ],
        [ 3, 1, 0, -1, -1, 0, 0, 0, 0, 0, 0, 0 ],
        [ 1, 1, 1, 1, 1, 0, 0, 0, 0, 0, 0, 0 ],
        [ 4, -2, 1, 0, 0, 1, -1, 1, 0, 0, 0, 0 ],
        [ 8, 0, -1, 0, 0, 2, 0, -1, 0, 0, 0, 0 ],
        [ 4, 2, 1, 0, 0, 1, 1, 1, 0, 0, 0, 0 ],
        [ 6, 0, 0, -2, 0, 3, -1, 0, 1, -1, 0, 0 ],
        [ 6, 2, 0, 2, 0, 3, 1, 0, 1, 1, 0, 0 ],
        [ 4, 2, 1, 0, 0, 3, 1, 0, 2, 0, 1, 0 ],
        [ 1, 1, 1, 1, 1, 1, 1, 1, 1, 1, 1, 1 ] ],
    [ <identity partial perm on [ 1, 2, 3, 4 ]>, (1)(2)(3,4),
        (1)(2,3,4), (1,2)(3,4), (1,2,3,4),
        <identity partial perm on [ 1, 2, 3 ]>, (1)(2,3), (1,2,3),
        <identity partial perm on [ 2, 3 ]>, (2,3),
        <identity partial perm on [ 1 ]>, <empty partial perm> ] ]
\end{Verbatim}
 }

 

\subsection{\textcolor{Chapter }{EUnitaryInverseCover}}
\logpage{[ 11, 15, 11 ]}\nobreak
\hyperdef{L}{X8383E6747D02D975}{}
{\noindent\textcolor{FuncColor}{$\triangleright$\enspace\texttt{EUnitaryInverseCover({\mdseries\slshape S})\index{EUnitaryInverseCover@\texttt{EUnitaryInverseCover}}
\label{EUnitaryInverseCover}
}\hfill{\scriptsize (attribute)}}\\
\textbf{\indent Returns:\ }
A homomorphism between semigroups.



 If the argument \mbox{\texttt{\mdseries\slshape S}} is an inverse semigroup then this function returns a finite
E\texttt{\symbol{45}}unitary inverse cover of \mbox{\texttt{\mdseries\slshape S}}. A finite E\texttt{\symbol{45}}unitary cover of \mbox{\texttt{\mdseries\slshape S}} is a surjective idempotent separating homomorphism from a finite semigroup
satisfying \texttt{IsEUnitaryInverseSemigroup} (\ref{IsEUnitaryInverseSemigroup}) to \mbox{\texttt{\mdseries\slshape S}}. A semigroup homomorphism is said to be idempotent separating if no two
idempotents are mapped to the same element of the image. 
\begin{Verbatim}[commandchars=!@|,fontsize=\small,frame=single,label=Example]
  !gapprompt@gap>| !gapinput@S := InverseSemigroup([PartialPermNC([1, 2], [2, 1]),|
  !gapprompt@>| !gapinput@PartialPermNC([1], [1])]);|
  <inverse partial perm semigroup of rank 2 with 2 generators>
  !gapprompt@gap>| !gapinput@cov := EUnitaryInverseCover(S);|
  <inverse partial perm semigroup of rank 4 with 2 generators> ->
  <inverse partial perm semigroup of rank 2 with 2 generators>
  !gapprompt@gap>| !gapinput@IsEUnitaryInverseSemigroup(Source(cov));|
  true
  !gapprompt@gap>| !gapinput@S = Range(cov);|
  true
\end{Verbatim}
 }

 }

   
\section{\textcolor{Chapter }{ Nambooripad partial order }}\logpage{[ 11, 16, 0 ]}
\hyperdef{L}{X7AA4CE887EEA661A}{}
{
  

\subsection{\textcolor{Chapter }{NambooripadLeqRegularSemigroup}}
\logpage{[ 11, 16, 1 ]}\nobreak
\hyperdef{L}{X7A7EB0DA8398886E}{}
{\noindent\textcolor{FuncColor}{$\triangleright$\enspace\texttt{NambooripadLeqRegularSemigroup({\mdseries\slshape S})\index{NambooripadLeqRegularSemigroup@\texttt{NambooripadLeqRegularSemigroup}}
\label{NambooripadLeqRegularSemigroup}
}\hfill{\scriptsize (attribute)}}\\
\textbf{\indent Returns:\ }
 A function. 



 \texttt{NambooripadLeqRegularSemigroup} returns a function that, when given two elements \texttt{x, y} of the regular semigroup \mbox{\texttt{\mdseries\slshape S}}, returns \texttt{true} if \texttt{x} is less than or equal to \texttt{y} in the Nambooripad partial order on \mbox{\texttt{\mdseries\slshape S}}. See also \texttt{NambooripadPartialOrder} (\ref{NambooripadPartialOrder}). 
\begin{Verbatim}[commandchars=!@|,fontsize=\small,frame=single,label=Example]
  !gapprompt@gap>| !gapinput@S := BrauerMonoid(3);|
  <regular bipartition *-monoid of degree 3 with 3 generators>
  !gapprompt@gap>| !gapinput@IsRegularSemigroup(S);|
  true
  !gapprompt@gap>| !gapinput@Size(S);|
  15
  !gapprompt@gap>| !gapinput@NambooripadPartialOrder(S);|
  [ [  ], [  ], [  ], [  ], [  ], [  ], [  ], [  ], [  ],
    [ 1, 2, 3, 4, 5, 6, 7, 8, 9 ], [ 1, 2, 3, 4, 5, 6, 7, 8, 9 ],
    [ 1, 2, 3, 4, 5, 6, 7, 8, 9 ], [ 1, 2, 3, 4, 5, 6, 7, 8, 9 ],
    [ 1, 2, 3, 4, 5, 6, 7, 8, 9 ], [ 1, 2, 3, 4, 5, 6, 7, 8, 9 ] ]
  !gapprompt@gap>| !gapinput@NambooripadLeqRegularSemigroup(S)(Elements(S)[3], Elements(S)[9]);|
  false
  !gapprompt@gap>| !gapinput@NambooripadLeqRegularSemigroup(S)(Elements(S)[2], Elements(S)[15]);|
  true
\end{Verbatim}
 }

 

\subsection{\textcolor{Chapter }{NambooripadPartialOrder}}
\logpage{[ 11, 16, 2 ]}\nobreak
\hyperdef{L}{X7928C7D37A9BCBD5}{}
{\noindent\textcolor{FuncColor}{$\triangleright$\enspace\texttt{NambooripadPartialOrder({\mdseries\slshape S})\index{NambooripadPartialOrder@\texttt{NambooripadPartialOrder}}
\label{NambooripadPartialOrder}
}\hfill{\scriptsize (attribute)}}\\
\textbf{\indent Returns:\ }
The Nambooripad partial order on a regular semigroup.



 The \emph{Nambooripad partial order} $\leq$ on a regular semigroup \mbox{\texttt{\mdseries\slshape S}} is defined by \texttt{s}$\leq$\texttt{t} if the principal right ideal of \texttt{S} generated by \texttt{s} is contained in the principal right ideal of \texttt{S} generated by \texttt{t} and there is an idempotent \texttt{e} in the $\mathcal{R}$\texttt{\symbol{45}}class of \texttt{s} such that \texttt{s}$=$\texttt{et}. The Nambooripad partial order coincides with the natural partial order when
considering inverse semigroups \texttt{NaturalPartialOrder} (\textbf{Reference: NaturalPartialOrder}).

 \texttt{NambooripadPartialOrder} returns the Nambooripad partial order on the regular semigroup \mbox{\texttt{\mdseries\slshape S}} as a list of sets of positive integers where entry \texttt{i} in \texttt{NambooripadPartialOrder(\mbox{\texttt{\mdseries\slshape S}})} is the set of positions in \texttt{Elements(\mbox{\texttt{\mdseries\slshape S}})} of elements which are less than \texttt{Elements(\mbox{\texttt{\mdseries\slshape S}})[i]}. See also \texttt{NambooripadLeqRegularSemigroup} (\ref{NambooripadLeqRegularSemigroup}). 
\begin{Verbatim}[commandchars=!@|,fontsize=\small,frame=single,label=Example]
  !gapprompt@gap>| !gapinput@S := BrauerMonoid(3);|
  <regular bipartition *-monoid of degree 3 with 3 generators>
  !gapprompt@gap>| !gapinput@IsRegularSemigroup(S);|
  true
  !gapprompt@gap>| !gapinput@Size(S);|
  15
  !gapprompt@gap>| !gapinput@NambooripadPartialOrder(S);|
  [ [  ], [  ], [  ], [  ], [  ], [  ], [  ], [  ], [  ],
    [ 1, 2, 3, 4, 5, 6, 7, 8, 9 ], [ 1, 2, 3, 4, 5, 6, 7, 8, 9 ],
    [ 1, 2, 3, 4, 5, 6, 7, 8, 9 ], [ 1, 2, 3, 4, 5, 6, 7, 8, 9 ],
    [ 1, 2, 3, 4, 5, 6, 7, 8, 9 ], [ 1, 2, 3, 4, 5, 6, 7, 8, 9 ] ]
  !gapprompt@gap>| !gapinput@NambooripadLeqRegularSemigroup(S)(Elements(S)[3], Elements(S)[9]);|
  false
  !gapprompt@gap>| !gapinput@NambooripadLeqRegularSemigroup(S)(Elements(S)[2], Elements(S)[15]);|
  true
\end{Verbatim}
 }

 }

   
\section{\textcolor{Chapter }{ Monoid character table and Cartan matrix }}\logpage{[ 11, 17, 0 ]}
\hyperdef{L}{X861A6A947F43E8A3}{}
{
  In this section we describe some operations in \textsf{Semigroups} for computing character tables and Cartan matrices of monoids. These functions
are currently implemented for finite monoids over a splitting field of the
monoid in characteristic zero. The case for computing character tables and
Cartan matrices for finite monoids over a splitting field of the monoid in
characteristic zero originates from \cite{Balthazar2025aa} by Balthazar Charles. An example of a splitting field for a finite monoid is
the cyclotomic field with $g$th roots where $g$ is the least common multiple of the exponents of every maximal subgroup of the
finite monoid. In what follows cyclotomic field with $g$th roots can be assumed to be what is meant by a splitting field of the
monoid.

 

\subsection{\textcolor{Chapter }{GeneralizedConjugacyClasses (for a semigroup)}}
\logpage{[ 11, 17, 1 ]}\nobreak
\hyperdef{L}{X8266019D85749451}{}
{\noindent\textcolor{FuncColor}{$\triangleright$\enspace\texttt{GeneralizedConjugacyClasses({\mdseries\slshape S})\index{GeneralizedConjugacyClasses@\texttt{GeneralizedConjugacyClasses}!for a semigroup}
\label{GeneralizedConjugacyClasses:for a semigroup}
}\hfill{\scriptsize (attribute)}}\\
\noindent\textcolor{FuncColor}{$\triangleright$\enspace\texttt{GeneralizedConjugacyClassesRepresentatives({\mdseries\slshape S})\index{GeneralizedConjugacyClassesRepresentatives@\texttt{Generalized}\-\texttt{Conjugacy}\-\texttt{Classes}\-\texttt{Representatives}!for a semigroup}
\label{GeneralizedConjugacyClassesRepresentatives:for a semigroup}
}\hfill{\scriptsize (attribute)}}\\
\textbf{\indent Returns:\ }
A complete list of generalized conjugacy classes.



 \texttt{GeneralizedConjugacyClasses} returns a list of the generalized conjugacy classes of the semigroup \mbox{\texttt{\mdseries\slshape S}}. Let $R$ be a relation on \mbox{\texttt{\mdseries\slshape S}} such that it is the union of the set $\{(mn,nm) | m,n \in S \}$ and the set $\{(m,m^{(w+1)}) | m \in S \}$ where $w$ is \texttt{SmallestIdempotentPower} (\ref{SmallestIdempotentPower}) of \texttt{m}. A generalized conjugacy class is an equivalence class of the smallest
equivalence relation that contains $R$. \texttt{GeneralizedConjugacyClassesRepresentatives} returns a list of the representatives of the generalized conjugacy classes in
the semigroup \mbox{\texttt{\mdseries\slshape S}}. In a group generalized conjugacy classes coincide with conjugacy classes.

 Note that each generalized conjugacy class as currently implemented does not
store a list of elements in the conjugacy class, only the representative with
no way to compute all the elements in the conjugacy class. This means that \texttt{GeneralizedConjugacyClasses} as currently implemented is only as effective as \texttt{GeneralizedConjugacyClassesRepresentatives}. 
\begin{Verbatim}[commandchars=!@|,fontsize=\small,frame=single,label=Example]
  !gapprompt@gap>| !gapinput@S := FullTransformationMonoid(6);;|
  !gapprompt@gap>| !gapinput@GeneralizedConjugacyClassesRepresentatives(S);|
  [ IdentityTransformation, Transformation( [ 1, 2, 3, 4, 6, 5 ] ),
    Transformation( [ 1, 2, 3, 5, 6, 4 ] ),
    Transformation( [ 1, 2, 4, 3, 6, 5 ] ),
    Transformation( [ 1, 2, 4, 5, 6, 3 ] ),
    Transformation( [ 1, 3, 2, 5, 6, 4 ] ),
    Transformation( [ 1, 3, 4, 5, 6, 2 ] ),
    Transformation( [ 2, 1, 4, 3, 6, 5 ] ),
    Transformation( [ 2, 1, 4, 5, 6, 3 ] ),
    Transformation( [ 2, 3, 1, 5, 6, 4 ] ),
    Transformation( [ 2, 3, 4, 5, 6, 1 ] ),
    Transformation( [ 1, 2, 3, 4, 5, 1 ] ),
    Transformation( [ 1, 2, 3, 5, 4, 1 ] ),
    Transformation( [ 1, 2, 4, 5, 3, 1 ] ),
    Transformation( [ 1, 3, 2, 5, 4, 1 ] ),
    Transformation( [ 1, 3, 4, 5, 2, 1 ] ),
    Transformation( [ 2, 1, 4, 5, 3, 2 ] ),
    Transformation( [ 2, 3, 4, 5, 1, 2 ] ),
    Transformation( [ 2, 2, 3, 4, 5, 2 ] ),
    Transformation( [ 2, 2, 3, 5, 4, 2 ] ),
    Transformation( [ 2, 2, 4, 5, 3, 2 ] ),
    Transformation( [ 3, 3, 2, 5, 4, 3 ] ),
    Transformation( [ 3, 3, 4, 5, 2, 3 ] ),
    Transformation( [ 1, 1, 3, 4, 3, 1 ] ),
    Transformation( [ 1, 1, 4, 3, 4, 1 ] ),
    Transformation( [ 3, 3, 4, 1, 4, 3 ] ),
    Transformation( [ 1, 2, 2, 1, 2, 1 ] ),
    Transformation( [ 2, 1, 1, 2, 1, 2 ] ),
    Transformation( [ 1, 1, 1, 1, 1, 1 ] ) ]
\end{Verbatim}
 }

 

\subsection{\textcolor{Chapter }{DClassBicharacter (for a D-classs)}}
\logpage{[ 11, 17, 2 ]}\nobreak
\hyperdef{L}{X87B16B0B8763C20C}{}
{\noindent\textcolor{FuncColor}{$\triangleright$\enspace\texttt{DClassBicharacter({\mdseries\slshape D})\index{DClassBicharacter@\texttt{DClassBicharacter}!for a D-classs}
\label{DClassBicharacter:for a D-classs}
}\hfill{\scriptsize (attribute)}}\\
\textbf{\indent Returns:\ }
A matrix of non\texttt{\symbol{45}}negative integers.



 \texttt{DClassBicharacter} returns a matrix whose \texttt{i}th row and \texttt{j}th column entry is the size of the set $\{m \in D | hmk = m \}$ where \texttt{h} is an element in the \texttt{i}th generalized conjugacy class of the \texttt{Parent} (\textbf{Reference: Parent}) monoid of \mbox{\texttt{\mdseries\slshape D}} and \texttt{k} is an element in the \texttt{j}th generalized conjugacy class of the \texttt{Parent} (\textbf{Reference: Parent}) monoid of \mbox{\texttt{\mdseries\slshape D}}. The order of the generalized conjugacy classes of the \texttt{Parent} (\textbf{Reference: Parent}) monoid of \mbox{\texttt{\mdseries\slshape D}} is determined by the list returned by the attribute \texttt{GeneralizedConjugacyClasses} (\ref{GeneralizedConjugacyClasses:for a semigroup}) of the \texttt{Parent} (\textbf{Reference: Parent}) monoid of \mbox{\texttt{\mdseries\slshape D}}. 
\begin{Verbatim}[commandchars=!@|,fontsize=\small,frame=single,label=Example]
  !gapprompt@gap>| !gapinput@S := FullTransformationMonoid(3);;|
  !gapprompt@gap>| !gapinput@D := DClasses(S)[1];;|
  !gapprompt@gap>| !gapinput@DClassBicharacter(D);|
  [ [ 6, 0, 0, 0, 0, 0 ], [ 0, 2, 0, 0, 0, 0 ], [ 0, 0, 3, 0, 0, 0 ],
    [ 0, 0, 0, 0, 0, 0 ], [ 0, 0, 0, 0, 0, 0 ], [ 0, 0, 0, 0, 0, 0 ] ]
    
\end{Verbatim}
 }

 

\subsection{\textcolor{Chapter }{RegularRepresentationBicharacter (for a semigroup)}}
\logpage{[ 11, 17, 3 ]}\nobreak
\hyperdef{L}{X862B49E0867BDE6E}{}
{\noindent\textcolor{FuncColor}{$\triangleright$\enspace\texttt{RegularRepresentationBicharacter({\mdseries\slshape S})\index{RegularRepresentationBicharacter@\texttt{RegularRepresentationBicharacter}!for a semigroup}
\label{RegularRepresentationBicharacter:for a semigroup}
}\hfill{\scriptsize (attribute)}}\\
\textbf{\indent Returns:\ }
A matrix of non\texttt{\symbol{45}}negative integers.



 \texttt{RegularRepresentationBicharacter} returns a matrix whose \texttt{i}th row and \texttt{j}th column entry is the size of the set $\{m \in \mbox{\texttt{\mdseries\slshape S}} | hmk = m \}$ where $h$ is an element in the \texttt{i}th generalized conjugacy class of \mbox{\texttt{\mdseries\slshape S}} and $k$ is an element in the \texttt{j}th generalized conjugacy class of \mbox{\texttt{\mdseries\slshape S}}. The order of the generalized conjugacy classes of \mbox{\texttt{\mdseries\slshape S}} is determined by the list returned by the attribute \texttt{GeneralizedConjugacyClasses} (\ref{GeneralizedConjugacyClasses:for a semigroup}) of \mbox{\texttt{\mdseries\slshape S}}. 
\begin{Verbatim}[commandchars=!@|,fontsize=\small,frame=single,label=Example]
  !gapprompt@gap>| !gapinput@S := FullTransformationMonoid(3);;|
  !gapprompt@gap>| !gapinput@RegularRepresentationBicharacter(S);|
  [ [ 27, 1, 0, 8, 0, 1 ], [ 9, 3, 0, 4, 0, 1 ], [ 3, 1, 3, 2, 0, 1 ],
    [ 9, 1, 0, 4, 0, 1 ], [ 3, 3, 0, 2, 2, 1 ], [ 3, 1, 0, 2, 0, 1 ] ]
    
\end{Verbatim}
 }

 

\subsection{\textcolor{Chapter }{RClassBicharacterOfGroupHClass (for a group H-class)}}
\logpage{[ 11, 17, 4 ]}\nobreak
\hyperdef{L}{X7CC118E584268C7D}{}
{\noindent\textcolor{FuncColor}{$\triangleright$\enspace\texttt{RClassBicharacterOfGroupHClass({\mdseries\slshape H})\index{RClassBicharacterOfGroupHClass@\texttt{RClassBicharacterOfGroupHClass}!for a group H-class}
\label{RClassBicharacterOfGroupHClass:for a group H-class}
}\hfill{\scriptsize (attribute)}}\\
\textbf{\indent Returns:\ }
A matrix of non\texttt{\symbol{45}}negative integers.



 \texttt{RClassBicharacterOfGroupHClass} returns a matrix whose \texttt{i}th row and \texttt{j}th column entry is the size of the set $\{m \in R | hmk = m \}$ where \texttt{R} is the \texttt{RClassOfHClass} (\ref{RClassOfHClass}) of \mbox{\texttt{\mdseries\slshape H}}. Furthermore, $h$ is an element in the \texttt{i}th generalized conjugacy class of \mbox{\texttt{\mdseries\slshape H}} as a group and $k$ is an element in the \texttt{j}th generalized conjugacy class of the \texttt{Parent} (\textbf{Reference: Parent}) semigroup of \mbox{\texttt{\mdseries\slshape H}}. The order of the generalized conjugacy classes of the \texttt{Parent} (\textbf{Reference: Parent}) semigroup of \mbox{\texttt{\mdseries\slshape H}} is determined by the list returned by the attribute \texttt{GeneralizedConjugacyClasses} (\ref{GeneralizedConjugacyClasses:for a semigroup}) of the \texttt{Parent} (\textbf{Reference: Parent}) semigroup of \mbox{\texttt{\mdseries\slshape H}}. The order of the generalized conjugacy classes of \mbox{\texttt{\mdseries\slshape H}} is determined by the list returned by the attribute \texttt{ConjugacyClasses} (\textbf{Reference: ConjugacyClasses attribute}) of \texttt{OrdinaryCharacterTable} (\textbf{Reference: OrdinaryCharacterTable}) of \texttt{Range} (\textbf{Reference: range}) of \texttt{IsomorphismPermGroup} (\ref{IsomorphismPermGroup}) of \mbox{\texttt{\mdseries\slshape H}}. 
\begin{Verbatim}[commandchars=!@|,fontsize=\small,frame=single,label=Example]
  !gapprompt@gap>| !gapinput@S := FullTransformationMonoid(3);;|
  !gapprompt@gap>| !gapinput@D := RegularDClasses(S)[1];;|
  !gapprompt@gap>| !gapinput@H := GroupHClass(D);;|
  !gapprompt@gap>| !gapinput@RClassBicharacterOfGroupHClass(H);|
  [ [ 6, 0, 0, 0, 0, 0 ], [ 0, 2, 0, 0, 0, 0 ], [ 0, 0, 3, 0, 0, 0 ] ]
\end{Verbatim}
 }

 

\subsection{\textcolor{Chapter }{BlockDiagonalMatrixOfCharacterTables (for a semigroup)}}
\logpage{[ 11, 17, 5 ]}\nobreak
\hyperdef{L}{X796302FF875B56A4}{}
{\noindent\textcolor{FuncColor}{$\triangleright$\enspace\texttt{BlockDiagonalMatrixOfCharacterTables({\mdseries\slshape S})\index{BlockDiagonalMatrixOfCharacterTables@\texttt{Block}\-\texttt{Diagonal}\-\texttt{Matrix}\-\texttt{Of}\-\texttt{Character}\-\texttt{Tables}!for a semigroup}
\label{BlockDiagonalMatrixOfCharacterTables:for a semigroup}
}\hfill{\scriptsize (attribute)}}\\
\textbf{\indent Returns:\ }
A matrix.



 \texttt{BlockDiagonalMatrixOfCharacterTables} returns a block diagonal matrix which has a matrix block for each regular $\mathcal{D}$\texttt{\symbol{45}}class of \mbox{\texttt{\mdseries\slshape S}}. Each diagonal blocks is the character table of one (arbitrary) group $\mathcal{H}$\texttt{\symbol{45}}class in each regular $\mathcal{D}$\texttt{\symbol{45}}class of \mbox{\texttt{\mdseries\slshape S}}. The character tables are determined by \texttt{OrdinaryCharacterTable} (\textbf{Reference: OrdinaryCharacterTable}) of \texttt{Range} (\textbf{Reference: range}) of \texttt{IsomorphismPermGroup} (\ref{IsomorphismPermGroup}) of \texttt{GroupHClass} (\ref{GroupHClass}) of \texttt{RegularDClasses} (\ref{RegularDClasses}) of \mbox{\texttt{\mdseries\slshape S}}. 
\begin{Verbatim}[commandchars=!@|,fontsize=\small,frame=single,label=Example]
  !gapprompt@gap>| !gapinput@S := FullTransformationMonoid(3);;|
  !gapprompt@gap>| !gapinput@BlockDiagonalMatrixOfCharacterTables(S);|
  [ [ 1, -1, 1, 0, 0, 0 ], [ 2, 0, -1, 0, 0, 0 ], [ 1, 1, 1, 0, 0, 0 ],
    [ 0, 0, 0, 1, -1, 0 ], [ 0, 0, 0, 1, 1, 0 ], [ 0, 0, 0, 0, 0, 1 ] ]
    
\end{Verbatim}
 }

 

\subsection{\textcolor{Chapter }{MonoidCharacterTable}}
\logpage{[ 11, 17, 6 ]}\nobreak
\hyperdef{L}{X7915B1EB85DC974A}{}
{\noindent\textcolor{FuncColor}{$\triangleright$\enspace\texttt{MonoidCharacterTable({\mdseries\slshape M[, F]})\index{MonoidCharacterTable@\texttt{MonoidCharacterTable}}
\label{MonoidCharacterTable}
}\hfill{\scriptsize (attribute)}}\\
\textbf{\indent Returns:\ }
The character table object for \mbox{\texttt{\mdseries\slshape M}} over \mbox{\texttt{\mdseries\slshape F}}.



 Called with a finite monoid \mbox{\texttt{\mdseries\slshape M}} and optionally a field \mbox{\texttt{\mdseries\slshape F}}, \texttt{MonoidCharacterTable} returns the character table of the monoid which is defined as the matrix $Trace(X(m))$, where $X$ runs over the simple $FM$\texttt{\symbol{45}}modules and $m$ runs over the generalized conjugacy class representatives of \mbox{\texttt{\mdseries\slshape M}}.

 To get the character table of a monoid to display like the example below, the
irreducible characters need to be computed first using \texttt{Irr} (\ref{Irr:for a monoid character table}).

 If \mbox{\texttt{\mdseries\slshape F}} is not given, then \texttt{MonoidCharacterTable} returns the character table of $M$ over a characteristic zero splitting field of \mbox{\texttt{\mdseries\slshape M}}.

 At the moment, methods are available for the following cases: if \mbox{\texttt{\mdseries\slshape F}} is not given (i.e. it defaults to a splitting field of \mbox{\texttt{\mdseries\slshape M}} over the rationals) and \mbox{\texttt{\mdseries\slshape M}} is a finite monoid.

 For other cases no methods are implemented yet.

 
\begin{Verbatim}[commandchars=!@|,fontsize=\small,frame=single,label=Example]
  !gapprompt@gap>| !gapinput@S := FullTransformationMonoid(3);;|
  !gapprompt@gap>| !gapinput@ct := MonoidCharacterTable(S);;|
  !gapprompt@gap>| !gapinput@Irr(ct);;|
  !gapprompt@gap>| !gapinput@Display(ct);|
      c.1 c.2 c.3 c.4 c.5 c.6
  
  X.1   1  -1   1   .   .   .
  X.2   2   .  -1   .   .   .
  X.3   1   1   1   .   .   .
  X.4   2   .  -1   1  -1   .
  X.5   3   1   .   1   1   .
  X.6   1   1   1   1   1   1
\end{Verbatim}
 }

 

\subsection{\textcolor{Chapter }{Irr (for a monoid character table)}}
\logpage{[ 11, 17, 7 ]}\nobreak
\hyperdef{L}{X8343917B8684DD94}{}
{\noindent\textcolor{FuncColor}{$\triangleright$\enspace\texttt{Irr({\mdseries\slshape T})\index{Irr@\texttt{Irr}!for a monoid character table}
\label{Irr:for a monoid character table}
}\hfill{\scriptsize (attribute)}}\\
\textbf{\indent Returns:\ }
A list of monoid characters.



 \texttt{Irr} returns a list of the characters of the irreducible representations of the \texttt{Parent} (\textbf{Reference: Parent}) monoid of the monoid character table \mbox{\texttt{\mdseries\slshape T}}.  
\begin{Verbatim}[commandchars=!@|,fontsize=\small,frame=single,label=Example]
  !gapprompt@gap>| !gapinput@M := FullBooleanMatMonoid(2);;|
  !gapprompt@gap>| !gapinput@ct := MonoidCharacterTable(M);;|
  !gapprompt@gap>| !gapinput@Irr(ct);|
  [ MonoidCharacter( MonoidCharacterTable( Monoid( [ Matrix(IsBooleanMat\
  , [ [ false, true ], [ true, false ] ]), Matrix(IsBooleanMat, [ [ true\
  , false ], [ true, true ] ]), Matrix(IsBooleanMat, [ [ true, false ], \
  [ false, false ] ]) ] ) ) , [ 1, 1, 1, 1, 1 ] ),
    MonoidCharacter( MonoidCharacterTable( Monoid( [ Matrix(IsBooleanMat\
  , [ [ false, true ], [ true, false ] ]), Matrix(IsBooleanMat, [ [ true\
  , false ], [ true, true ] ]), Matrix(IsBooleanMat, [ [ true, false ], \
  [ false, false ] ]) ] ) ) , [ 0, 1, 2, 3, 1 ] ),
    MonoidCharacter( MonoidCharacterTable( Monoid( [ Matrix(IsBooleanMat\
  , [ [ false, true ], [ true, false ] ]), Matrix(IsBooleanMat, [ [ true\
  , false ], [ true, true ] ]), Matrix(IsBooleanMat, [ [ true, false ], \
  [ false, false ] ]) ] ) ) , [ 0, 0, 1, 2, 0 ] ),
    MonoidCharacter( MonoidCharacterTable( Monoid( [ Matrix(IsBooleanMat\
  , [ [ false, true ], [ true, false ] ]), Matrix(IsBooleanMat, [ [ true\
  , false ], [ true, true ] ]), Matrix(IsBooleanMat, [ [ true, false ], \
  [ false, false ] ]) ] ) ) , [ 0, 0, 0, 1, -1 ] ),
    MonoidCharacter( MonoidCharacterTable( Monoid( [ Matrix(IsBooleanMat\
  , [ [ false, true ], [ true, false ] ]), Matrix(IsBooleanMat, [ [ true\
  , false ], [ true, true ] ]), Matrix(IsBooleanMat, [ [ true, false ], \
  [ false, false ] ]) ] ) ) , [ 0, 0, 0, 1, 1 ] ) ]
\end{Verbatim}
 }

 

\subsection{\textcolor{Chapter }{MonoidCartanMatrix}}
\logpage{[ 11, 17, 8 ]}\nobreak
\hyperdef{L}{X7A7384F679AAABFA}{}
{\noindent\textcolor{FuncColor}{$\triangleright$\enspace\texttt{MonoidCartanMatrix({\mdseries\slshape M[, F]})\index{MonoidCartanMatrix@\texttt{MonoidCartanMatrix}}
\label{MonoidCartanMatrix}
}\hfill{\scriptsize (attribute)}}\\
\textbf{\indent Returns:\ }
An object.



 Called with a finite monoid \mbox{\texttt{\mdseries\slshape M}} and a field \mbox{\texttt{\mdseries\slshape F}}, \texttt{MonoidCartanMatrix} returns the Cartan matrix of the monoid algebra $FM$ which is defined as the matrix $dim Hom(P,Q)/dim End(P / rad(FM))$, where $P$ and $Q$ run over the right indecomposable projective modules of FM.

 To get the Cartan matrix of a monoid to display like the example below, the
projective indecomposable modules need to be computed first using \texttt{Pims} (\ref{Pims:for a monoid cartan matrix}).

 If \mbox{\texttt{\mdseries\slshape M}} is the only argument then \texttt{MonoidCartanMatrix} returns the Cartan matrix of the monoid algebra $FM$, where \mbox{\texttt{\mdseries\slshape F}} is a splitting field of \mbox{\texttt{\mdseries\slshape M}} over the rationals.

 At the moment, methods are available for the following cases: if \mbox{\texttt{\mdseries\slshape F}} is not given (i.e. it defaults to a splitting field of \mbox{\texttt{\mdseries\slshape M}} over the rationals) and \mbox{\texttt{\mdseries\slshape M}} is a finite monoid.

 For other cases no methods are implemented yet.

 
\begin{Verbatim}[commandchars=!@|,fontsize=\small,frame=single,label=Example]
  !gapprompt@gap>| !gapinput@S := FullTransformationMonoid(3);;|
  !gapprompt@gap>| !gapinput@cm := MonoidCartanMatrix(S);;|
  !gapprompt@gap>| !gapinput@Pims(cm);;|
  !gapprompt@gap>| !gapinput@Display(cm);|
      X.1 X.2 X.3 X.4 X.5 X.6
  
  P.1   1   .   .   .   .   .
  P.2   .   1   .   .   .   .
  P.3   1   .   1   1   .   .
  P.4   1   .   .   1   .   .
  P.5   .   .   .   .   1   .
  P.6   .   .   .   1   .   1
\end{Verbatim}
 }

 

\subsection{\textcolor{Chapter }{Pims (for a monoid cartan matrix)}}
\logpage{[ 11, 17, 9 ]}\nobreak
\hyperdef{L}{X7D9AA972875D54C4}{}
{\noindent\textcolor{FuncColor}{$\triangleright$\enspace\texttt{Pims({\mdseries\slshape T})\index{Pims@\texttt{Pims}!for a monoid cartan matrix}
\label{Pims:for a monoid cartan matrix}
}\hfill{\scriptsize (attribute)}}\\
\textbf{\indent Returns:\ }
A list of monoid characters.



 \texttt{Pims} returns a list of the characters of the projective indecomposable modules of
the \texttt{Parent} (\textbf{Reference: Parent}) monoid of the monoid character table \mbox{\texttt{\mdseries\slshape T}}.  
\begin{Verbatim}[commandchars=!@|,fontsize=\small,frame=single,label=Example]
  !gapprompt@gap>| !gapinput@M := FullBooleanMatMonoid(2);;|
  !gapprompt@gap>| !gapinput@cm := MonoidCartanMatrix(M);;|
  !gapprompt@gap>| !gapinput@Pims(cm);|
  [ MonoidCharacter( MonoidCharacterTable( Monoid( [ Matrix(IsBooleanMat\
  , [ [ false, true ], [ true, false ] ]), Matrix(IsBooleanMat, [ [ true\
  , false ], [ true, true ] ]), Matrix(IsBooleanMat, [ [ true, false ], \
  [ false, false ] ]) ] ) ) , Projective Cover Of MonoidCharacter( Monoi\
  dCharacterTable( Monoid( [ Matrix(IsBooleanMat, [ [ false, true ], [ t\
  rue, false ] ]), Matrix(IsBooleanMat, [ [ true, false ], [ true, true \
  ] ]), Matrix(IsBooleanMat, [ [ true, false ], [ false, false ] ]) ] ) \
  ) , [ 1, 1, 1, 1, 1 ] ) ),
    MonoidCharacter( MonoidCharacterTable( Monoid( [ Matrix(IsBooleanMat\
  , [ [ false, true ], [ true, false ] ]), Matrix(IsBooleanMat, [ [ true\
  , false ], [ true, true ] ]), Matrix(IsBooleanMat, [ [ true, false ], \
  [ false, false ] ]) ] ) ) , Projective Cover Of MonoidCharacter( Monoi\
  dCharacterTable( Monoid( [ Matrix(IsBooleanMat, [ [ false, true ], [ t\
  rue, false ] ]), Matrix(IsBooleanMat, [ [ true, false ], [ true, true \
  ] ]), Matrix(IsBooleanMat, [ [ true, false ], [ false, false ] ]) ] ) \
  ) , [ 0, 1, 2, 3, 1 ] ) ),
    MonoidCharacter( MonoidCharacterTable( Monoid( [ Matrix(IsBooleanMat\
  , [ [ false, true ], [ true, false ] ]), Matrix(IsBooleanMat, [ [ true\
  , false ], [ true, true ] ]), Matrix(IsBooleanMat, [ [ true, false ], \
  [ false, false ] ]) ] ) ) , Projective Cover Of MonoidCharacter( Monoi\
  dCharacterTable( Monoid( [ Matrix(IsBooleanMat, [ [ false, true ], [ t\
  rue, false ] ]), Matrix(IsBooleanMat, [ [ true, false ], [ true, true \
  ] ]), Matrix(IsBooleanMat, [ [ true, false ], [ false, false ] ]) ] ) \
  ) , [ 0, 0, 1, 2, 0 ] ) ),
    MonoidCharacter( MonoidCharacterTable( Monoid( [ Matrix(IsBooleanMat\
  , [ [ false, true ], [ true, false ] ]), Matrix(IsBooleanMat, [ [ true\
  , false ], [ true, true ] ]), Matrix(IsBooleanMat, [ [ true, false ], \
  [ false, false ] ]) ] ) ) , Projective Cover Of MonoidCharacter( Monoi\
  dCharacterTable( Monoid( [ Matrix(IsBooleanMat, [ [ false, true ], [ t\
  rue, false ] ]), Matrix(IsBooleanMat, [ [ true, false ], [ true, true \
  ] ]), Matrix(IsBooleanMat, [ [ true, false ], [ false, false ] ]) ] ) \
  ) , [ 0, 0, 0, 1, -1 ] ) ),
    MonoidCharacter( MonoidCharacterTable( Monoid( [ Matrix(IsBooleanMat\
  , [ [ false, true ], [ true, false ] ]), Matrix(IsBooleanMat, [ [ true\
  , false ], [ true, true ] ]), Matrix(IsBooleanMat, [ [ true, false ], \
  [ false, false ] ]) ] ) ) , Projective Cover Of MonoidCharacter( Monoi\
  dCharacterTable( Monoid( [ Matrix(IsBooleanMat, [ [ false, true ], [ t\
  rue, false ] ]), Matrix(IsBooleanMat, [ [ true, false ], [ true, true \
  ] ]), Matrix(IsBooleanMat, [ [ true, false ], [ false, false ] ]) ] ) \
  ) , [ 0, 0, 0, 1, 1 ] ) ) ]
\end{Verbatim}
 }

 }

   }

 
\chapter{\textcolor{Chapter }{ Properties of semigroups }}\label{Properties of semigroups}
\logpage{[ 12, 0, 0 ]}
\hyperdef{L}{X78274024827F306D}{}
{
  In this chapter we describe the methods that are available in \textsf{Semigroups} for determining various properties of a semigroup or monoid.   
\section{\textcolor{Chapter }{ Arbitrary semigroups }}\logpage{[ 12, 1, 0 ]}
\hyperdef{L}{X7D297AEC827F3D4E}{}
{
  In this section we describe the properties of an arbitrary semigroup or monoid
that can be determined using the \textsf{Semigroups} package. 

 

\subsection{\textcolor{Chapter }{IsBand}}
\logpage{[ 12, 1, 1 ]}\nobreak
\hyperdef{L}{X7C8DB14587D1B55A}{}
{\noindent\textcolor{FuncColor}{$\triangleright$\enspace\texttt{IsBand({\mdseries\slshape S})\index{IsBand@\texttt{IsBand}}
\label{IsBand}
}\hfill{\scriptsize (property)}}\\
\textbf{\indent Returns:\ }
\texttt{true} or \texttt{false}.



 \texttt{IsBand} returns \texttt{true} if every element of the semigroup \mbox{\texttt{\mdseries\slshape S}} is an idempotent and \texttt{false} if it is not. An inverse semigroup is band if and only if it is a semilattice;
see \texttt{IsSemilattice} (\ref{IsSemilattice}). 
\begin{Verbatim}[commandchars=!@|,fontsize=\small,frame=single,label=Example]
  !gapprompt@gap>| !gapinput@S := Semigroup(|
  !gapprompt@>| !gapinput@Transformation([1, 1, 1, 4, 4, 4, 7, 7, 7, 1]),|
  !gapprompt@>| !gapinput@Transformation([2, 2, 2, 5, 5, 5, 8, 8, 8, 2]),|
  !gapprompt@>| !gapinput@Transformation([3, 3, 3, 6, 6, 6, 9, 9, 9, 3]),|
  !gapprompt@>| !gapinput@Transformation([1, 1, 1, 4, 4, 4, 7, 7, 7, 4]),|
  !gapprompt@>| !gapinput@Transformation([1, 1, 1, 4, 4, 4, 7, 7, 7, 7]));;|
  !gapprompt@gap>| !gapinput@IsBand(S);|
  true
  !gapprompt@gap>| !gapinput@S := InverseSemigroup(|
  !gapprompt@>| !gapinput@PartialPerm([1, 2, 3, 4, 8, 9], [5, 8, 7, 6, 9, 1]),|
  !gapprompt@>| !gapinput@PartialPerm([1, 3, 4, 7, 8, 9, 10], [2, 3, 8, 7, 10, 6, 1]));;|
  !gapprompt@gap>| !gapinput@IsBand(S);|
  false
  !gapprompt@gap>| !gapinput@IsBand(IdempotentGeneratedSubsemigroup(S));|
  true
  !gapprompt@gap>| !gapinput@S := PartitionMonoid(4);|
  <regular bipartition *-monoid of size 4140, degree 4 with 4
   generators>
  !gapprompt@gap>| !gapinput@M := MinimalIdeal(S);|
  <simple bipartition *-semigroup ideal of degree 4 with 1 generator>
  !gapprompt@gap>| !gapinput@IsBand(M);|
  true
\end{Verbatim}
 }

 

\subsection{\textcolor{Chapter }{IsBlockGroup}}
\logpage{[ 12, 1, 2 ]}\nobreak
\hyperdef{L}{X79659C467C8A7EBD}{}
{\noindent\textcolor{FuncColor}{$\triangleright$\enspace\texttt{IsBlockGroup({\mdseries\slshape S})\index{IsBlockGroup@\texttt{IsBlockGroup}}
\label{IsBlockGroup}
}\hfill{\scriptsize (property)}}\\
\textbf{\indent Returns:\ }
\texttt{true} or \texttt{false}.



 \texttt{IsBlockGroup} returns \texttt{true} if the semigroup \mbox{\texttt{\mdseries\slshape S}} is a block group and \texttt{false} if it is not.

 A semigroup \mbox{\texttt{\mdseries\slshape S}} is a \emph{block group} if every $\mathcal{L}$\texttt{\symbol{45}}class and every $\mathcal{R}$\texttt{\symbol{45}}class of \mbox{\texttt{\mdseries\slshape S}} contains at most one idempotent. Every semigroup of partial permutations is a
block group. 
\begin{Verbatim}[commandchars=!@|,fontsize=\small,frame=single,label=Example]
  !gapprompt@gap>| !gapinput@S := Semigroup(Transformation([5, 6, 7, 3, 1, 4, 2, 8]),|
  !gapprompt@>| !gapinput@                  Transformation([3, 6, 8, 5, 7, 4, 2, 8]));;|
  !gapprompt@gap>| !gapinput@IsBlockGroup(S);|
  true
  !gapprompt@gap>| !gapinput@S := Semigroup(|
  !gapprompt@>| !gapinput@Transformation([2, 1, 10, 4, 5, 9, 7, 4, 8, 4]),|
  !gapprompt@>| !gapinput@Transformation([10, 7, 5, 6, 1, 3, 9, 7, 10, 2]));;|
  !gapprompt@gap>| !gapinput@IsBlockGroup(S);|
  false
  !gapprompt@gap>| !gapinput@S := Semigroup(PartialPerm([1, 2], [5, 4]),|
  !gapprompt@>| !gapinput@                  PartialPerm([1, 2, 3], [1, 2, 5]),|
  !gapprompt@>| !gapinput@                  PartialPerm([1, 2, 3], [2, 1, 5]),|
  !gapprompt@>| !gapinput@                  PartialPerm([1, 3, 4], [3, 1, 2]),|
  !gapprompt@>| !gapinput@                  PartialPerm([1, 3, 4, 5], [5, 4, 3, 2]));;|
  !gapprompt@gap>| !gapinput@T := AsSemigroup(IsBlockBijectionSemigroup, S);|
  <block bijection semigroup of degree 6 with 5 generators>
  !gapprompt@gap>| !gapinput@IsBlockGroup(T);|
  true
  !gapprompt@gap>| !gapinput@IsBlockGroup(AsSemigroup(IsBipartitionSemigroup, S));|
  true
  !gapprompt@gap>| !gapinput@S := Semigroup(|
  !gapprompt@>| !gapinput@Bipartition([[1, -2], [2, -3], [3, -4], [4, -1]]),|
  !gapprompt@>| !gapinput@Bipartition([[1, -2], [2, -1], [3, -3], [4, -4]]),|
  !gapprompt@>| !gapinput@Bipartition([[1, 2, -3], [3, -1, -2], [4, -4]]),|
  !gapprompt@>| !gapinput@Bipartition([[1, -1], [2, -2], [3, -3], [4, -4]]));;|
  !gapprompt@gap>| !gapinput@IsBlockGroup(S);|
  true
\end{Verbatim}
 }

 

\subsection{\textcolor{Chapter }{IsCommutativeSemigroup}}
\logpage{[ 12, 1, 3 ]}\nobreak
\hyperdef{L}{X843EFDA4807FDC31}{}
{\noindent\textcolor{FuncColor}{$\triangleright$\enspace\texttt{IsCommutativeSemigroup({\mdseries\slshape S})\index{IsCommutativeSemigroup@\texttt{IsCommutativeSemigroup}}
\label{IsCommutativeSemigroup}
}\hfill{\scriptsize (property)}}\\
\textbf{\indent Returns:\ }
\texttt{true} or \texttt{false}.



 \texttt{IsCommutativeSemigroup} returns \texttt{true} if the semigroup \mbox{\texttt{\mdseries\slshape S}} is commutative and \texttt{false} if it is not. The function \texttt{IsCommutative} (\textbf{Reference: IsCommutative}) can also be used to test if a semigroup is commutative.

 A semigroup \mbox{\texttt{\mdseries\slshape S}} is \emph{commutative} if \texttt{x * y = y * x} for all \texttt{x, y} in \mbox{\texttt{\mdseries\slshape S}}. 
\begin{Verbatim}[commandchars=!@|,fontsize=\small,frame=single,label=Example]
  !gapprompt@gap>| !gapinput@S := Semigroup(Transformation([2, 4, 5, 3, 7, 8, 6, 9, 1]),|
  !gapprompt@>| !gapinput@                  Transformation([3, 5, 6, 7, 8, 1, 9, 2, 4]));;|
  !gapprompt@gap>| !gapinput@IsCommutativeSemigroup(S);|
  true
  !gapprompt@gap>| !gapinput@IsCommutative(S);|
  true
  !gapprompt@gap>| !gapinput@S := InverseSemigroup(|
  !gapprompt@>| !gapinput@ PartialPerm([1, 2, 3, 4, 5, 6], [2, 5, 1, 3, 9, 6]),|
  !gapprompt@>| !gapinput@ PartialPerm([1, 2, 3, 4, 6, 8], [8, 5, 7, 6, 2, 1]));;|
  !gapprompt@gap>| !gapinput@IsCommutativeSemigroup(S);|
  false
  !gapprompt@gap>| !gapinput@S := Semigroup(|
  !gapprompt@>| !gapinput@Bipartition([[1, 2, 3, 6, 7, -1, -4, -6],|
  !gapprompt@>| !gapinput@             [4, 5, 8, -2, -3, -5, -7, -8]]),|
  !gapprompt@>| !gapinput@ Bipartition([[1, 2, -3, -4], [3, -5], [4, -6], [5, -7],|
  !gapprompt@>| !gapinput@              [6, -8], [7, -1], [8, -2]]));;|
  !gapprompt@gap>| !gapinput@IsCommutativeSemigroup(S);|
  true
\end{Verbatim}
 }

 

\subsection{\textcolor{Chapter }{IsCompletelyRegularSemigroup}}
\logpage{[ 12, 1, 4 ]}\nobreak
\hyperdef{L}{X7AFA23AF819FBF3D}{}
{\noindent\textcolor{FuncColor}{$\triangleright$\enspace\texttt{IsCompletelyRegularSemigroup({\mdseries\slshape S})\index{IsCompletelyRegularSemigroup@\texttt{IsCompletelyRegularSemigroup}}
\label{IsCompletelyRegularSemigroup}
}\hfill{\scriptsize (property)}}\\
\textbf{\indent Returns:\ }
\texttt{true} or \texttt{false}.



 \texttt{IsCompletelyRegularSemigroup} returns \texttt{true} if every element of the semigroup \mbox{\texttt{\mdseries\slshape S}} is contained in a subgroup of \mbox{\texttt{\mdseries\slshape S}}.

 An inverse semigroup is completely regular if and only if it is a Clifford
semigroup; see \texttt{IsCliffordSemigroup} (\ref{IsCliffordSemigroup}). 
\begin{Verbatim}[commandchars=!@|,fontsize=\small,frame=single,label=Example]
  !gapprompt@gap>| !gapinput@S := Semigroup(Transformation([1, 2, 4, 3, 6, 5, 4]),|
  !gapprompt@>| !gapinput@                  Transformation([1, 2, 5, 6, 3, 4, 5]),|
  !gapprompt@>| !gapinput@                  Transformation([2, 1, 2, 2, 2, 2, 2]));;|
  !gapprompt@gap>| !gapinput@IsCompletelyRegularSemigroup(S);|
  true
  !gapprompt@gap>| !gapinput@IsInverseSemigroup(S);|
  true
  !gapprompt@gap>| !gapinput@T := Range(IsomorphismPartialPermSemigroup(S));;|
  !gapprompt@gap>| !gapinput@IsCompletelyRegularSemigroup(T);|
  true
  !gapprompt@gap>| !gapinput@IsCliffordSemigroup(T);|
  true
  !gapprompt@gap>| !gapinput@S := Semigroup(|
  !gapprompt@>| !gapinput@Bipartition([[1, 3, -4], [2, 4, -1, -2], [-3]]),|
  !gapprompt@>| !gapinput@Bipartition([[1, -1], [2, 3, 4, -3], [-2, -4]]));;|
  !gapprompt@gap>| !gapinput@IsCompletelyRegularSemigroup(S);|
  false
\end{Verbatim}
 }

 

\subsection{\textcolor{Chapter }{IsCongruenceFreeSemigroup}}
\logpage{[ 12, 1, 5 ]}\nobreak
\hyperdef{L}{X855088F378D8F5E1}{}
{\noindent\textcolor{FuncColor}{$\triangleright$\enspace\texttt{IsCongruenceFreeSemigroup({\mdseries\slshape S})\index{IsCongruenceFreeSemigroup@\texttt{IsCongruenceFreeSemigroup}}
\label{IsCongruenceFreeSemigroup}
}\hfill{\scriptsize (property)}}\\
\textbf{\indent Returns:\ }
\texttt{true} or \texttt{false}.



 \texttt{IsCongruenceFreeSemigroup} returns \texttt{true} if the semigroup \mbox{\texttt{\mdseries\slshape S}} is a congruence\texttt{\symbol{45}}free semigroup and \texttt{false} if it is not.

 A semigroup \mbox{\texttt{\mdseries\slshape S}} is \emph{congruence\texttt{\symbol{45}}free} if it has no non\texttt{\symbol{45}}trivial proper congruences.

 A semigroup with zero is congruence\texttt{\symbol{45}}free if and only if it
is isomorphic to a regular Rees 0\texttt{\symbol{45}}matrix semigroup \texttt{R} whose underlying semigroup is the trivial group, no two rows of the matrix of \texttt{R} are identical, and no two columns are identical; see Theorem 3.7.1 in \cite{Howie1995aa}.

 A semigroup without zero is congruence\texttt{\symbol{45}}free if and only if
it is a simple group or has order 2; see Theorem 3.7.2 in \cite{Howie1995aa}. 
\begin{Verbatim}[commandchars=!@|,fontsize=\small,frame=single,label=Example]
  !gapprompt@gap>| !gapinput@S := Semigroup(Transformation([4, 2, 3, 3, 4]));;|
  !gapprompt@gap>| !gapinput@IsCongruenceFreeSemigroup(S);|
  true
  !gapprompt@gap>| !gapinput@S := Semigroup(Transformation([2, 2, 4, 4]),|
  !gapprompt@>| !gapinput@                  Transformation([5, 3, 4, 4, 6, 6]));;|
  !gapprompt@gap>| !gapinput@IsCongruenceFreeSemigroup(S);|
  false
\end{Verbatim}
 }

 

\subsection{\textcolor{Chapter }{IsCryptoGroup}}
\logpage{[ 12, 1, 6 ]}\nobreak
\hyperdef{L}{X7D3560117D2014FE}{}
{\noindent\textcolor{FuncColor}{$\triangleright$\enspace\texttt{IsCryptoGroup({\mdseries\slshape S})\index{IsCryptoGroup@\texttt{IsCryptoGroup}}
\label{IsCryptoGroup}
}\hfill{\scriptsize (property)}}\\
\textbf{\indent Returns:\ }
\texttt{true} or \texttt{false}.



 A semigroup \mbox{\texttt{\mdseries\slshape S}} is a \emph{crypto group} if it is completely regular and Green's $\mathcal{H}$\texttt{\symbol{45}}relation is a congruence. \texttt{IsCryptoGroup(\mbox{\texttt{\mdseries\slshape S}})} returns \texttt{true} if the semigroup \mbox{\texttt{\mdseries\slshape S}} is a crypto group, and \texttt{false} if it is not. 

 
\begin{Verbatim}[commandchars=!@|,fontsize=\small,frame=single,label=Example]
  !gapprompt@gap>| !gapinput@G := SymmetricGroup(5);;|
  !gapprompt@gap>| !gapinput@M := [[(1, 2), (2, 4)], [(1, 4), (1, 2, 3, 4, 5)]];;|
  !gapprompt@gap>| !gapinput@R := ReesMatrixSemigroup(G, M);;|
  !gapprompt@gap>| !gapinput@IsCryptoGroup(R);|
  true
  !gapprompt@gap>| !gapinput@S := FullTransformationMonoid(2);;|
  !gapprompt@gap>| !gapinput@IsCryptoGroup(S);|
  false
  !gapprompt@gap>| !gapinput@G1 := AlternatingGroup(4);;|
  !gapprompt@gap>| !gapinput@G2 := SymmetricGroup(3);;|
  !gapprompt@gap>| !gapinput@G3 := AlternatingGroup(5);;|
  !gapprompt@gap>| !gapinput@gr := Digraph([[1, 3], [2, 3], [3]]);;|
  !gapprompt@gap>| !gapinput@sgn := function(x)|
  !gapprompt@>| !gapinput@if SignPerm(x) = 1 then|
  !gapprompt@>| !gapinput@return ();|
  !gapprompt@>| !gapinput@fi;|
  !gapprompt@>| !gapinput@return (1, 2);|
  !gapprompt@>| !gapinput@end;;|
  !gapprompt@gap>| !gapinput@hom13 := GroupHomomorphismByFunction(G1, G3, sgn);;|
  !gapprompt@gap>| !gapinput@hom23 := GroupHomomorphismByFunction(G2, G3, sgn);;|
  !gapprompt@gap>| !gapinput@S := AsSemigroup(IsPartialPermSemigroup,|
  !gapprompt@>| !gapinput@gr,|
  !gapprompt@>| !gapinput@[G1, G2, G3], [[1, 3, hom13], [2, 3, hom23]]);;|
  !gapprompt@gap>| !gapinput@IsCryptoGroup(S);|
  true
\end{Verbatim}
 }

 

\subsection{\textcolor{Chapter }{IsSurjectiveSemigroup}}
\logpage{[ 12, 1, 7 ]}\nobreak
\hyperdef{L}{X7C9560A18348AEE3}{}
{\noindent\textcolor{FuncColor}{$\triangleright$\enspace\texttt{IsSurjectiveSemigroup({\mdseries\slshape S})\index{IsSurjectiveSemigroup@\texttt{IsSurjectiveSemigroup}}
\label{IsSurjectiveSemigroup}
}\hfill{\scriptsize (property)}}\\
\textbf{\indent Returns:\ }
\texttt{true} or \texttt{false}.



 A semigroup is \emph{surjective} if each of its elements can be written as a product of two elements in the
semigroup. \texttt{IsSurjectiveSemigroup(\mbox{\texttt{\mdseries\slshape S}})} returns \texttt{true} if the semigroup \mbox{\texttt{\mdseries\slshape S}} is surjective, and \texttt{false} if it is not. 

 See also \texttt{IndecomposableElements} (\ref{IndecomposableElements}). 

 Note that every monoid, and every regular semigroup, is surjective. 
\begin{Verbatim}[commandchars=!@|,fontsize=\small,frame=single,label=Example]
  !gapprompt@gap>| !gapinput@S := FullTransformationMonoid(100);|
  <full transformation monoid of degree 100>
  !gapprompt@gap>| !gapinput@IsSurjectiveSemigroup(S);|
  true
  !gapprompt@gap>| !gapinput@F := FreeSemigroup(3);;|
  !gapprompt@gap>| !gapinput@P := F / [[F.1, F.2 * F.1], [F.3 ^ 3, F.3]];|
  <fp semigroup with 3 generators and 2 relations of length 10>
  !gapprompt@gap>| !gapinput@IsSurjectiveSemigroup(P);|
  false
  !gapprompt@gap>| !gapinput@I := SingularTransformationMonoid(5);|
  <regular transformation semigroup ideal of degree 5 with 1 generator>
  !gapprompt@gap>| !gapinput@IsSurjectiveSemigroup(I);|
  true
  !gapprompt@gap>| !gapinput@M := MonogenicSemigroup(IsBipartitionSemigroup, 3, 2);|
  <commutative non-regular block bijection semigroup of size 4,
   degree 6 with 1 generator>
  !gapprompt@gap>| !gapinput@IsSurjectiveSemigroup(M);|
  false
\end{Verbatim}
 }

 

\subsection{\textcolor{Chapter }{IsGroupAsSemigroup}}
\logpage{[ 12, 1, 8 ]}\nobreak
\hyperdef{L}{X852F29E8795FA489}{}
{\noindent\textcolor{FuncColor}{$\triangleright$\enspace\texttt{IsGroupAsSemigroup({\mdseries\slshape S})\index{IsGroupAsSemigroup@\texttt{IsGroupAsSemigroup}}
\label{IsGroupAsSemigroup}
}\hfill{\scriptsize (property)}}\\
\textbf{\indent Returns:\ }
\texttt{true} or \texttt{false}.



 \texttt{IsGroupAsSemigroup} returns \texttt{true} if and only if the semigroup \mbox{\texttt{\mdseries\slshape S}} is mathematically a group. 
\begin{Verbatim}[commandchars=!@|,fontsize=\small,frame=single,label=Example]
  !gapprompt@gap>| !gapinput@S := Semigroup(Transformation([2, 4, 5, 3, 7, 8, 6, 9, 1]),|
  !gapprompt@>| !gapinput@                  Transformation([3, 5, 6, 7, 8, 1, 9, 2, 4]));;|
  !gapprompt@gap>| !gapinput@IsGroupAsSemigroup(S);|
  true
  !gapprompt@gap>| !gapinput@G := SymmetricGroup(5);;|
  !gapprompt@gap>| !gapinput@IsGroupAsSemigroup(G);|
  true
  !gapprompt@gap>| !gapinput@S := AsSemigroup(IsPartialPermSemigroup, G);|
  <partial perm group of size 120, rank 5 with 2 generators>
  !gapprompt@gap>| !gapinput@IsGroupAsSemigroup(S);|
  true
  !gapprompt@gap>| !gapinput@G := SymmetricGroup([1, 2, 10]);;|
  !gapprompt@gap>| !gapinput@T := AsSemigroup(IsBlockBijectionSemigroup, G);|
  <inverse block bijection semigroup of size 6, degree 11 with 2
   generators>
  !gapprompt@gap>| !gapinput@IsGroupAsSemigroup(T);|
  true
\end{Verbatim}
 }

 
\subsection{\textcolor{Chapter }{IsIdempotentGenerated}}\logpage{[ 12, 1, 9 ]}
\hyperdef{L}{X835484C481CF3DDD}{}
{
\noindent\textcolor{FuncColor}{$\triangleright$\enspace\texttt{IsIdempotentGenerated({\mdseries\slshape S})\index{IsIdempotentGenerated@\texttt{IsIdempotentGenerated}}
\label{IsIdempotentGenerated}
}\hfill{\scriptsize (property)}}\\
\noindent\textcolor{FuncColor}{$\triangleright$\enspace\texttt{IsSemiband({\mdseries\slshape S})\index{IsSemiband@\texttt{IsSemiband}}
\label{IsSemiband}
}\hfill{\scriptsize (property)}}\\
\textbf{\indent Returns:\ }
 \texttt{true} or \texttt{false}. 



 \texttt{IsIdempotentGenerated} and \texttt{IsSemiband} return \texttt{true} if the semigroup \mbox{\texttt{\mdseries\slshape S}} is generated by its idempotents and \texttt{false} if it is not. See also \texttt{Idempotents} (\ref{Idempotents}) and \texttt{IdempotentGeneratedSubsemigroup} (\ref{IdempotentGeneratedSubsemigroup}). 

 An inverse semigroup is idempotent\texttt{\symbol{45}}generated if and only if
it is a semilattice; see \texttt{IsSemilattice} (\ref{IsSemilattice}).

 The terms semiband and idempotent\texttt{\symbol{45}}generated are synonymous
in this context. 
\begin{Verbatim}[commandchars=!@|,fontsize=\small,frame=single,label=Example]
  !gapprompt@gap>| !gapinput@S := SingularTransformationSemigroup(4);|
  <regular transformation semigroup ideal of degree 4 with 1 generator>
  !gapprompt@gap>| !gapinput@IsIdempotentGenerated(S);|
  true
  !gapprompt@gap>| !gapinput@S := SingularBrauerMonoid(5);|
  <regular bipartition *-semigroup ideal of degree 5 with 1 generator>
  !gapprompt@gap>| !gapinput@IsIdempotentGenerated(S);|
  true
\end{Verbatim}
 }

 

\subsection{\textcolor{Chapter }{IsLeftSimple}}
\logpage{[ 12, 1, 10 ]}\nobreak
\hyperdef{L}{X8206D2B0809952EF}{}
{\noindent\textcolor{FuncColor}{$\triangleright$\enspace\texttt{IsLeftSimple({\mdseries\slshape S})\index{IsLeftSimple@\texttt{IsLeftSimple}}
\label{IsLeftSimple}
}\hfill{\scriptsize (property)}}\\
\noindent\textcolor{FuncColor}{$\triangleright$\enspace\texttt{IsRightSimple({\mdseries\slshape S})\index{IsRightSimple@\texttt{IsRightSimple}}
\label{IsRightSimple}
}\hfill{\scriptsize (property)}}\\
\textbf{\indent Returns:\ }
\texttt{true} or \texttt{false}.



 \texttt{IsLeftSimple} and \texttt{IsRightSimple} returns \texttt{true} if the semigroup \mbox{\texttt{\mdseries\slshape S}} has only one $\mathcal{L}$\texttt{\symbol{45}}class or one $\mathcal{R}$\texttt{\symbol{45}}class, respectively, and returns \texttt{false} if it has more than one. 

 An inverse semigroup is left simple if and only if it is right simple if and
only if it is a group; see \texttt{IsGroupAsSemigroup} (\ref{IsGroupAsSemigroup}). 
\begin{Verbatim}[commandchars=!@|,fontsize=\small,frame=single,label=Example]
  !gapprompt@gap>| !gapinput@S := Semigroup(Transformation([6, 7, 9, 6, 8, 9, 8, 7, 6]),|
  !gapprompt@>| !gapinput@                  Transformation([6, 8, 9, 6, 8, 8, 7, 9, 6]),|
  !gapprompt@>| !gapinput@                  Transformation([6, 8, 9, 7, 8, 8, 7, 9, 6]),|
  !gapprompt@>| !gapinput@                  Transformation([6, 9, 8, 6, 7, 9, 7, 8, 6]),|
  !gapprompt@>| !gapinput@                  Transformation([6, 9, 9, 6, 8, 8, 7, 9, 6]),|
  !gapprompt@>| !gapinput@                  Transformation([6, 9, 9, 7, 8, 8, 6, 9, 7]),|
  !gapprompt@>| !gapinput@                  Transformation([7, 8, 8, 7, 9, 9, 7, 8, 6]),|
  !gapprompt@>| !gapinput@                  Transformation([7, 9, 9, 7, 6, 9, 6, 8, 7]),|
  !gapprompt@>| !gapinput@                  Transformation([8, 7, 6, 9, 8, 6, 8, 7, 9]),|
  !gapprompt@>| !gapinput@                  Transformation([9, 6, 6, 7, 8, 8, 7, 6, 9]),|
  !gapprompt@>| !gapinput@                  Transformation([9, 6, 6, 7, 9, 6, 9, 8, 7]),|
  !gapprompt@>| !gapinput@                  Transformation([9, 6, 7, 9, 6, 6, 9, 7, 8]),|
  !gapprompt@>| !gapinput@                  Transformation([9, 6, 8, 7, 9, 6, 9, 8, 7]),|
  !gapprompt@>| !gapinput@                  Transformation([9, 7, 6, 8, 7, 7, 9, 6, 8]),|
  !gapprompt@>| !gapinput@                  Transformation([9, 7, 7, 8, 9, 6, 9, 7, 8]),|
  !gapprompt@>| !gapinput@                  Transformation([9, 8, 8, 9, 6, 7, 6, 8, 9]));;|
  !gapprompt@gap>| !gapinput@IsRightSimple(S);|
  false
  !gapprompt@gap>| !gapinput@IsLeftSimple(S);|
  true
  !gapprompt@gap>| !gapinput@IsGroupAsSemigroup(S);|
  false
  !gapprompt@gap>| !gapinput@NrRClasses(S);|
  16
  !gapprompt@gap>| !gapinput@S := BrauerMonoid(6);;|
  !gapprompt@gap>| !gapinput@S := Semigroup(RClass(S, Random(MinimalDClass(S))));;|
  !gapprompt@gap>| !gapinput@IsLeftSimple(S);|
  false
  !gapprompt@gap>| !gapinput@IsRightSimple(S);|
  true
\end{Verbatim}
 }

 

\subsection{\textcolor{Chapter }{IsLeftZeroSemigroup}}
\logpage{[ 12, 1, 11 ]}\nobreak
\hyperdef{L}{X7E9261367C8C52C0}{}
{\noindent\textcolor{FuncColor}{$\triangleright$\enspace\texttt{IsLeftZeroSemigroup({\mdseries\slshape S})\index{IsLeftZeroSemigroup@\texttt{IsLeftZeroSemigroup}}
\label{IsLeftZeroSemigroup}
}\hfill{\scriptsize (property)}}\\
\textbf{\indent Returns:\ }
\texttt{true} or \texttt{false}.



 \texttt{IsLeftZeroSemigroup} returns \texttt{true} if the semigroup \mbox{\texttt{\mdseries\slshape S}} is a left zero semigroup and \texttt{false} if it is not. 

 A semigroup is a \emph{left zero semigroup} if \texttt{x*y=x} for all \texttt{x,y}. An inverse semigroup is a left zero semigroup if and only if it is trivial. 
\begin{Verbatim}[commandchars=!@|,fontsize=\small,frame=single,label=Example]
  !gapprompt@gap>| !gapinput@S := Semigroup(Transformation([2, 1, 4, 3, 5]),|
  !gapprompt@>| !gapinput@                  Transformation([3, 2, 3, 1, 1]));;|
  !gapprompt@gap>| !gapinput@IsRightZeroSemigroup(S);|
  false
  !gapprompt@gap>| !gapinput@S := Semigroup(Transformation([1, 2, 3, 3, 1]),|
  !gapprompt@>| !gapinput@                  Transformation([1, 2, 3, 3, 3]));;|
  !gapprompt@gap>| !gapinput@IsLeftZeroSemigroup(S);|
  true
\end{Verbatim}
 }

 

\subsection{\textcolor{Chapter }{IsMonogenicSemigroup}}
\logpage{[ 12, 1, 12 ]}\nobreak
\hyperdef{L}{X79D46BAB7E327AD1}{}
{\noindent\textcolor{FuncColor}{$\triangleright$\enspace\texttt{IsMonogenicSemigroup({\mdseries\slshape S})\index{IsMonogenicSemigroup@\texttt{IsMonogenicSemigroup}}
\label{IsMonogenicSemigroup}
}\hfill{\scriptsize (property)}}\\
\textbf{\indent Returns:\ }
\texttt{true} or \texttt{false}.



 \texttt{IsMonogenicSemigroup} returns \texttt{true} if the semigroup \mbox{\texttt{\mdseries\slshape S}} is monogenic and it returns \texttt{false} if it is not. 

 A semigroup is \emph{monogenic} if it is generated by a single element. See also \texttt{IsMonogenicMonoid} (\ref{IsMonogenicMonoid}), \texttt{IsMonogenicInverseSemigroup} (\ref{IsMonogenicInverseSemigroup}), and \texttt{IsMonogenicInverseMonoid} (\ref{IsMonogenicInverseMonoid}). 
\begin{Verbatim}[commandchars=!@|,fontsize=\small,frame=single,label=Example]
  !gapprompt@gap>| !gapinput@S := Semigroup(|
  !gapprompt@>| !gapinput@Transformation(|
  !gapprompt@>| !gapinput@ [2, 2, 2, 11, 10, 8, 10, 11, 2, 11, 10, 2, 11, 11, 10]),|
  !gapprompt@>| !gapinput@Transformation(|
  !gapprompt@>| !gapinput@ [2, 2, 2, 8, 11, 15, 11, 10, 2, 10, 11, 2, 10, 4, 7]),|
  !gapprompt@>| !gapinput@Transformation(|
  !gapprompt@>| !gapinput@ [2, 2, 2, 11, 10, 8, 10, 11, 2, 11, 10, 2, 11, 11, 10]),|
  !gapprompt@>| !gapinput@Transformation(|
  !gapprompt@>| !gapinput@ [2, 2, 12, 7, 8, 14, 8, 11, 2, 11, 10, 2, 11, 15, 4]));;|
  !gapprompt@gap>| !gapinput@IsMonogenicSemigroup(S);|
  true
  !gapprompt@gap>| !gapinput@S := Semigroup(|
  !gapprompt@>| !gapinput@Bipartition([[1, 2, 3, 4, 5, 6, 7, 8, 9, 10, -2, -5, -7, -9],|
  !gapprompt@>| !gapinput@             [-1, -10], [-3, -4, -6, -8]]),|
  !gapprompt@>| !gapinput@ Bipartition([[1, 4, 7, 8, -2], [2, 3, 5, 10, -5],|
  !gapprompt@>| !gapinput@              [6, 9, -7, -9], [-1, -10], [-3, -4, -6, -8]]));;|
  !gapprompt@gap>| !gapinput@IsMonogenicSemigroup(S);|
  true
  !gapprompt@gap>| !gapinput@S := FullTransformationSemigroup(5);;|
  !gapprompt@gap>| !gapinput@IsMonogenicSemigroup(S);|
  false
\end{Verbatim}
 }

 

\subsection{\textcolor{Chapter }{IsMonogenicMonoid}}
\logpage{[ 12, 1, 13 ]}\nobreak
\hyperdef{L}{X790DC9F4798DBB09}{}
{\noindent\textcolor{FuncColor}{$\triangleright$\enspace\texttt{IsMonogenicMonoid({\mdseries\slshape S})\index{IsMonogenicMonoid@\texttt{IsMonogenicMonoid}}
\label{IsMonogenicMonoid}
}\hfill{\scriptsize (property)}}\\
\textbf{\indent Returns:\ }
\texttt{true} or \texttt{false}.



 \texttt{IsMonogenicMonoid} returns \texttt{true} if the monoid \mbox{\texttt{\mdseries\slshape S}} is a monogenic monoid and it returns \texttt{false} if it is not. 

 A monoid is \emph{monogenic} if it is generated as a monoid by a single element. See also \texttt{IsMonogenicSemigroup} (\ref{IsMonogenicSemigroup}) and \texttt{IsMonogenicInverseMonoid} (\ref{IsMonogenicInverseMonoid}). 
\begin{Verbatim}[commandchars=!@|,fontsize=\small,frame=single,label=Example]
  !gapprompt@gap>| !gapinput@x := PartialPerm([1, 2, 3, 6, 8, 10], [2, 6, 7, 9, 1, 5]);;|
  !gapprompt@gap>| !gapinput@S := Monoid(x, x ^ 2, x ^ 3);;|
  !gapprompt@gap>| !gapinput@IsMonogenicSemigroup(S);|
  false
  !gapprompt@gap>| !gapinput@IsMonogenicMonoid(S);|
  true
  !gapprompt@gap>| !gapinput@S := FullTransformationMonoid(5);;|
  !gapprompt@gap>| !gapinput@IsMonogenicMonoid(S);|
  false
\end{Verbatim}
 }

 

\subsection{\textcolor{Chapter }{IsMonoidAsSemigroup}}
\logpage{[ 12, 1, 14 ]}\nobreak
\hyperdef{L}{X7E4DEECD7CD9886D}{}
{\noindent\textcolor{FuncColor}{$\triangleright$\enspace\texttt{IsMonoidAsSemigroup({\mdseries\slshape S})\index{IsMonoidAsSemigroup@\texttt{IsMonoidAsSemigroup}}
\label{IsMonoidAsSemigroup}
}\hfill{\scriptsize (property)}}\\
\textbf{\indent Returns:\ }
\texttt{true} or \texttt{false}.



 \texttt{IsMonoidAsSemigroup} returns \texttt{true} if and only if the semigroup \mbox{\texttt{\mdseries\slshape S}} is mathematically a monoid, i.e. if and only if it contains a \texttt{MultiplicativeNeutralElement} (\textbf{Reference: MultiplicativeNeutralElement}). 

 It is possible that a semigroup which satisfies \texttt{IsMonoidAsSemigroup} is not in the \textsf{GAP} category \texttt{IsMonoid} (\textbf{Reference: IsMonoid}). This is possible if the \texttt{MultiplicativeNeutralElement} (\textbf{Reference: MultiplicativeNeutralElement}) of \mbox{\texttt{\mdseries\slshape S}} is not equal to the \texttt{One} (\textbf{Reference: One}) of any element in \mbox{\texttt{\mdseries\slshape S}}. Therefore a semigroup satisfying \texttt{IsMonoidAsSemigroup} may not possess the attributes of a monoid (such as, \texttt{GeneratorsOfMonoid} (\textbf{Reference: GeneratorsOfMonoid})).

 See also \texttt{One} (\textbf{Reference: One}), \texttt{IsInverseMonoid} (\textbf{Reference: IsInverseMonoid}) and \texttt{IsomorphismTransformationMonoid} (\textbf{Reference: IsomorphismTransformationMonoid}). 
\begin{Verbatim}[commandchars=!@|,fontsize=\small,frame=single,label=Example]
  !gapprompt@gap>| !gapinput@S := Semigroup(Transformation([1, 4, 6, 2, 5, 3, 7, 8, 9, 9]),|
  !gapprompt@>| !gapinput@                  Transformation([6, 3, 2, 7, 5, 1, 8, 8, 9, 9]));;|
  !gapprompt@gap>| !gapinput@IsMonoidAsSemigroup(S);|
  true
  !gapprompt@gap>| !gapinput@IsMonoid(S);|
  false
  !gapprompt@gap>| !gapinput@MultiplicativeNeutralElement(S);|
  Transformation( [ 1, 2, 3, 4, 5, 6, 7, 8, 9, 9 ] )
  !gapprompt@gap>| !gapinput@T := AsSemigroup(IsBipartitionSemigroup, S);;|
  !gapprompt@gap>| !gapinput@IsMonoidAsSemigroup(T);|
  true
  !gapprompt@gap>| !gapinput@IsMonoid(T);|
  false
  !gapprompt@gap>| !gapinput@One(T);|
  fail
  !gapprompt@gap>| !gapinput@S := Monoid(Transformation([8, 2, 8, 9, 10, 6, 2, 8, 7, 8]),|
  !gapprompt@>| !gapinput@               Transformation([9, 2, 6, 3, 6, 4, 5, 5, 3, 2]));;|
  !gapprompt@gap>| !gapinput@IsMonoidAsSemigroup(S);|
  true
\end{Verbatim}
 }

 

\subsection{\textcolor{Chapter }{IsOrthodoxSemigroup}}
\logpage{[ 12, 1, 15 ]}\nobreak
\hyperdef{L}{X7935C476808C8773}{}
{\noindent\textcolor{FuncColor}{$\triangleright$\enspace\texttt{IsOrthodoxSemigroup({\mdseries\slshape S})\index{IsOrthodoxSemigroup@\texttt{IsOrthodoxSemigroup}}
\label{IsOrthodoxSemigroup}
}\hfill{\scriptsize (property)}}\\
\textbf{\indent Returns:\ }
\texttt{true} or \texttt{false}.



 \texttt{IsOrthodoxSemigroup} returns \texttt{true} if the semigroup \mbox{\texttt{\mdseries\slshape S}} is orthodox and \texttt{false} if it is not.

 A semigroup is \emph{orthodox} if it is regular and its idempotent elements form a subsemigroup. Every
inverse semigroup is also an orthodox semigroup.

 See also \texttt{IsRegularSemigroup} (\ref{IsRegularSemigroup}) and \texttt{IsRegularSemigroup} (\textbf{Reference: IsRegularSemigroup}). 
\begin{Verbatim}[commandchars=!@|,fontsize=\small,frame=single,label=Example]
  !gapprompt@gap>| !gapinput@S := Semigroup(Transformation([1, 1, 1, 4, 5, 4]),|
  !gapprompt@>| !gapinput@                  Transformation([1, 2, 3, 1, 1, 2]),|
  !gapprompt@>| !gapinput@                  Transformation([1, 2, 3, 1, 1, 3]),|
  !gapprompt@>| !gapinput@                  Transformation([5, 5, 5, 5, 5, 5]));;|
  !gapprompt@gap>| !gapinput@IsOrthodoxSemigroup(S);|
  true
  !gapprompt@gap>| !gapinput@S := DualSymmetricInverseMonoid(5);;|
  !gapprompt@gap>| !gapinput@S := Semigroup(GeneratorsOfSemigroup(S));;|
  !gapprompt@gap>| !gapinput@IsOrthodoxSemigroup(S);|
  true
\end{Verbatim}
 }

 

\subsection{\textcolor{Chapter }{IsRectangularBand}}
\logpage{[ 12, 1, 16 ]}\nobreak
\hyperdef{L}{X7E9B674D781B072C}{}
{\noindent\textcolor{FuncColor}{$\triangleright$\enspace\texttt{IsRectangularBand({\mdseries\slshape S})\index{IsRectangularBand@\texttt{IsRectangularBand}}
\label{IsRectangularBand}
}\hfill{\scriptsize (property)}}\\
\textbf{\indent Returns:\ }
\texttt{true} or \texttt{false}.



 \texttt{IsRectangularBand} returns \texttt{true} if the semigroup \mbox{\texttt{\mdseries\slshape S}} is a rectangular band and \texttt{false} if it is not.

 A semigroup \mbox{\texttt{\mdseries\slshape S}} is a \emph{rectangular band} if for all $x, y, z$ in \mbox{\texttt{\mdseries\slshape S}} we have that $x ^ 2 = x$ and $xyz = xz$.

 Equivalently, \mbox{\texttt{\mdseries\slshape S}} is a \emph{rectangular band} if \mbox{\texttt{\mdseries\slshape S}} is isomorphic to a semigroup of the form $I \times \Lambda$ with multiplication $(i, \lambda)(j, \mu) = (i, \mu)$. In this case, \mbox{\texttt{\mdseries\slshape S}} is called an $|I| \times |\Lambda|$ \emph{rectangular band}.

 An inverse semigroup is a rectangular band if and only if it is a group. 
\begin{Verbatim}[commandchars=!@|,fontsize=\small,frame=single,label=Example]
  !gapprompt@gap>| !gapinput@S := Semigroup(|
  !gapprompt@>| !gapinput@Transformation([1, 1, 1, 4, 4, 4, 7, 7, 7, 1]),|
  !gapprompt@>| !gapinput@Transformation([2, 2, 2, 5, 5, 5, 8, 8, 8, 2]),|
  !gapprompt@>| !gapinput@Transformation([3, 3, 3, 6, 6, 6, 9, 9, 9, 3]),|
  !gapprompt@>| !gapinput@Transformation([1, 1, 1, 4, 4, 4, 7, 7, 7, 4]),|
  !gapprompt@>| !gapinput@Transformation([1, 1, 1, 4, 4, 4, 7, 7, 7, 7]));;|
  !gapprompt@gap>| !gapinput@IsRectangularBand(S);|
  true
  !gapprompt@gap>| !gapinput@IsRectangularBand(MinimalIdeal(PartitionMonoid(4)));|
  true
\end{Verbatim}
 }

 

\subsection{\textcolor{Chapter }{IsRectangularGroup}}
\logpage{[ 12, 1, 17 ]}\nobreak
\hyperdef{L}{X80E682BB78547F41}{}
{\noindent\textcolor{FuncColor}{$\triangleright$\enspace\texttt{IsRectangularGroup({\mdseries\slshape S})\index{IsRectangularGroup@\texttt{IsRectangularGroup}}
\label{IsRectangularGroup}
}\hfill{\scriptsize (property)}}\\
\textbf{\indent Returns:\ }
\texttt{true} or \texttt{false}.



 A semigroup is \emph{rectangular group} if it is the direct product of a group and a rectangular band. Or
equivalently, if it is orthodox and simple. 
\begin{Verbatim}[commandchars=!@|,fontsize=\small,frame=single,label=Example]
  !gapprompt@gap>| !gapinput@G := AsSemigroup(IsTransformationSemigroup, MathieuGroup(11));|
  <transformation group of size 7920, degree 11 with 2 generators>
  !gapprompt@gap>| !gapinput@R := RectangularBand(3, 2);|
  <regular transformation semigroup of size 6, degree 6 with 3
   generators>
  !gapprompt@gap>| !gapinput@S := DirectProduct(G, R);;|
  !gapprompt@gap>| !gapinput@IsRectangularGroup(R);|
  true
  !gapprompt@gap>| !gapinput@IsRectangularGroup(G);|
  true
  !gapprompt@gap>| !gapinput@IsRectangularGroup(S);|
  true
  !gapprompt@gap>| !gapinput@IsRectangularGroup(JonesMonoid(3));|
  false
\end{Verbatim}
 }

 

\subsection{\textcolor{Chapter }{IsRegularSemigroup}}
\logpage{[ 12, 1, 18 ]}\nobreak
\hyperdef{L}{X7C4663827C5ACEF1}{}
{\noindent\textcolor{FuncColor}{$\triangleright$\enspace\texttt{IsRegularSemigroup({\mdseries\slshape S})\index{IsRegularSemigroup@\texttt{IsRegularSemigroup}}
\label{IsRegularSemigroup}
}\hfill{\scriptsize (property)}}\\
\textbf{\indent Returns:\ }
\texttt{true} or \texttt{false}.



 \texttt{IsRegularSemigroup} returns \texttt{true} if the semigroup \mbox{\texttt{\mdseries\slshape S}} is regular and \texttt{false} if it is not. 

 A semigroup \texttt{S} is \emph{regular} if for all \texttt{x} in \texttt{S} there exists \texttt{y} in \texttt{S} such that \texttt{x * y * x = x}. Every inverse semigroup is regular, and a semigroup of partial permutations
is regular if and only if it is an inverse semigroup.

 See also \texttt{IsRegularDClass} (\textbf{Reference: IsRegularDClass}), \texttt{IsRegularGreensClass} (\ref{IsRegularGreensClass}), and \texttt{IsRegularSemigroupElement} (\textbf{Reference: IsRegularSemigroupElement}). 
\begin{Verbatim}[commandchars=!@|,fontsize=\small,frame=single,label=Example]
  !gapprompt@gap>| !gapinput@IsRegularSemigroup(FullTransformationSemigroup(5));|
  true
  !gapprompt@gap>| !gapinput@IsRegularSemigroup(JonesMonoid(5));|
  true
\end{Verbatim}
 }

 

\subsection{\textcolor{Chapter }{IsRightZeroSemigroup}}
\logpage{[ 12, 1, 19 ]}\nobreak
\hyperdef{L}{X7CB099958658F979}{}
{\noindent\textcolor{FuncColor}{$\triangleright$\enspace\texttt{IsRightZeroSemigroup({\mdseries\slshape S})\index{IsRightZeroSemigroup@\texttt{IsRightZeroSemigroup}}
\label{IsRightZeroSemigroup}
}\hfill{\scriptsize (property)}}\\
\textbf{\indent Returns:\ }
\texttt{true} or \texttt{false}.



 \texttt{IsRightZeroSemigroup} returns \texttt{true} if the \mbox{\texttt{\mdseries\slshape S}} is a right zero semigroup and \texttt{false} if it is not.

 A semigroup \texttt{S} is a \emph{right zero semigroup} if \texttt{x * y = y} for all \texttt{x, y} in \texttt{S}. An inverse semigroup is a right zero semigroup if and only if it is trivial. 
\begin{Verbatim}[commandchars=!@|,fontsize=\small,frame=single,label=Example]
  !gapprompt@gap>| !gapinput@S := Semigroup(Transformation([2, 1, 4, 3, 5]),|
  !gapprompt@>| !gapinput@                  Transformation([3, 2, 3, 1, 1]));;|
  !gapprompt@gap>| !gapinput@IsRightZeroSemigroup(S);|
  false
  !gapprompt@gap>| !gapinput@S := Semigroup(Transformation([1, 2, 3, 3, 1]),|
  !gapprompt@>| !gapinput@                  Transformation([1, 2, 4, 4, 1]));;|
  !gapprompt@gap>| !gapinput@IsRightZeroSemigroup(S);|
  true
\end{Verbatim}
 }

 
\subsection{\textcolor{Chapter }{IsXTrivial}}\logpage{[ 12, 1, 20 ]}
\hyperdef{L}{X8752642C7F7E512B}{}
{
\noindent\textcolor{FuncColor}{$\triangleright$\enspace\texttt{IsRTrivial({\mdseries\slshape S})\index{IsRTrivial@\texttt{IsRTrivial}}
\label{IsRTrivial}
}\hfill{\scriptsize (property)}}\\
\noindent\textcolor{FuncColor}{$\triangleright$\enspace\texttt{IsLTrivial({\mdseries\slshape S})\index{IsLTrivial@\texttt{IsLTrivial}}
\label{IsLTrivial}
}\hfill{\scriptsize (property)}}\\
\noindent\textcolor{FuncColor}{$\triangleright$\enspace\texttt{IsHTrivial({\mdseries\slshape S})\index{IsHTrivial@\texttt{IsHTrivial}}
\label{IsHTrivial}
}\hfill{\scriptsize (property)}}\\
\noindent\textcolor{FuncColor}{$\triangleright$\enspace\texttt{IsDTrivial({\mdseries\slshape S})\index{IsDTrivial@\texttt{IsDTrivial}}
\label{IsDTrivial}
}\hfill{\scriptsize (property)}}\\
\noindent\textcolor{FuncColor}{$\triangleright$\enspace\texttt{IsAperiodicSemigroup({\mdseries\slshape S})\index{IsAperiodicSemigroup@\texttt{IsAperiodicSemigroup}}
\label{IsAperiodicSemigroup}
}\hfill{\scriptsize (property)}}\\
\noindent\textcolor{FuncColor}{$\triangleright$\enspace\texttt{IsCombinatorialSemigroup({\mdseries\slshape S})\index{IsCombinatorialSemigroup@\texttt{IsCombinatorialSemigroup}}
\label{IsCombinatorialSemigroup}
}\hfill{\scriptsize (property)}}\\
\textbf{\indent Returns:\ }
 \texttt{true} or \texttt{false}. 



 \texttt{IsXTrivial} returns \texttt{true} if Green's $\mathcal{R}$\texttt{\symbol{45}}relation, $\mathcal{L}$\texttt{\symbol{45}}relation, $\mathcal{H}$\texttt{\symbol{45}}relation, $\mathcal{D}$\texttt{\symbol{45}}relation, respectively, on the semigroup \mbox{\texttt{\mdseries\slshape S}} is trivial and \texttt{false} if it is not. These properties can also be applied to a Green's class instead
of a semigroup where applicable. 

 For inverse semigroups, the properties of being $\mathcal{R}$\texttt{\symbol{45}}trivial, $\mathcal{L}$\texttt{\symbol{45}}trivial, $\mathcal{D}$\texttt{\symbol{45}}trivial, and a semilattice are equivalent; see \texttt{IsSemilattice} (\ref{IsSemilattice}). 

 A semigroup is \emph{aperiodic} if its contains no non\texttt{\symbol{45}}trivial subgroups (equivalently, all
of its group $\mathcal{H}$\texttt{\symbol{45}}classes are trivial). A finite semigroup is aperiodic if
and only if it is $\mathcal{H}$\texttt{\symbol{45}}trivial. 

 \emph{Combinatorial} is a synonym for aperiodic in this context. 
\begin{Verbatim}[commandchars=!@|,fontsize=\small,frame=single,label=Example]
  !gapprompt@gap>| !gapinput@S := Semigroup(|
  !gapprompt@>| !gapinput@ Transformation([1, 5, 1, 3, 7, 10, 6, 2, 7, 10]),|
  !gapprompt@>| !gapinput@ Transformation([4, 4, 5, 6, 7, 7, 7, 4, 3, 10]));;|
  !gapprompt@gap>| !gapinput@IsHTrivial(S);|
  true
  !gapprompt@gap>| !gapinput@Size(S);|
  108
  !gapprompt@gap>| !gapinput@IsRTrivial(S);|
  false
  !gapprompt@gap>| !gapinput@IsLTrivial(S);|
  false
\end{Verbatim}
 }

 

\subsection{\textcolor{Chapter }{IsSemigroupWithAdjoinedZero}}
\logpage{[ 12, 1, 21 ]}\nobreak
\hyperdef{L}{X7826DDF8808EC4D9}{}
{\noindent\textcolor{FuncColor}{$\triangleright$\enspace\texttt{IsSemigroupWithAdjoinedZero({\mdseries\slshape S})\index{IsSemigroupWithAdjoinedZero@\texttt{IsSemigroupWithAdjoinedZero}}
\label{IsSemigroupWithAdjoinedZero}
}\hfill{\scriptsize (property)}}\\
\textbf{\indent Returns:\ }
\texttt{true} or \texttt{false}. 



 \texttt{IsSemigroupWithAdjoinedZero} returns \texttt{true} if the semigroup \mbox{\texttt{\mdseries\slshape S}} can be expressed as the disjoint union of subsemigroups $\mbox{\texttt{\mdseries\slshape S}} \setminus \{ 0 \}$ and $\{ 0 \}$ (where $0$ is the \texttt{MultiplicativeZero} (\ref{MultiplicativeZero}) of \mbox{\texttt{\mdseries\slshape S}}). 

 If this is not the case, then either \mbox{\texttt{\mdseries\slshape S}} lacks a multiplicative zero, or the set $\mbox{\texttt{\mdseries\slshape S}} \setminus \{ 0 \}$ is not a subsemigroup of \mbox{\texttt{\mdseries\slshape S}}, and so \texttt{IsSemigroupWithAdjoinedZero} returns \texttt{false}. 
\begin{Verbatim}[commandchars=!@|,fontsize=\small,frame=single,label=Example]
  !gapprompt@gap>| !gapinput@S := Semigroup(Transformation([2, 3, 4, 5, 1, 6]),|
  !gapprompt@>| !gapinput@                  Transformation([2, 1, 3, 4, 5, 6]),|
  !gapprompt@>| !gapinput@                  Transformation([6, 6, 6, 6, 6, 6]));|
  <transformation semigroup of degree 6 with 3 generators>
  !gapprompt@gap>| !gapinput@IsZeroGroup(S);|
  true
  !gapprompt@gap>| !gapinput@IsSemigroupWithAdjoinedZero(S);|
  true
  !gapprompt@gap>| !gapinput@S := FullTransformationMonoid(4);;|
  !gapprompt@gap>| !gapinput@IsSemigroupWithAdjoinedZero(S);|
  false
\end{Verbatim}
 }

 

\subsection{\textcolor{Chapter }{IsSemilattice}}
\logpage{[ 12, 1, 22 ]}\nobreak
\hyperdef{L}{X833D24AE7C900B85}{}
{\noindent\textcolor{FuncColor}{$\triangleright$\enspace\texttt{IsSemilattice({\mdseries\slshape S})\index{IsSemilattice@\texttt{IsSemilattice}}
\label{IsSemilattice}
}\hfill{\scriptsize (property)}}\\
\textbf{\indent Returns:\ }
\texttt{true} or \texttt{false}.



 \texttt{IsSemilattice} returns \texttt{true} if the semigroup \mbox{\texttt{\mdseries\slshape S}} is a semilattice and \texttt{false} if it is not. 

 A semigroup is a \emph{semilattice} if it is commutative and every element is an idempotent. The idempotents of an
inverse semigroup form a semilattice. 
\begin{Verbatim}[commandchars=!@|,fontsize=\small,frame=single,label=Example]
  !gapprompt@gap>| !gapinput@S := Semigroup(Transformation([2, 5, 1, 7, 3, 7, 7]),|
  !gapprompt@>| !gapinput@                  Transformation([3, 6, 5, 7, 2, 1, 7]));;|
  !gapprompt@gap>| !gapinput@Size(S);|
  631
  !gapprompt@gap>| !gapinput@IsInverseSemigroup(S);|
  true
  !gapprompt@gap>| !gapinput@A := Semigroup(Idempotents(S));|
  <transformation semigroup of degree 7 with 32 generators>
  !gapprompt@gap>| !gapinput@IsSemilattice(A);|
  true
  !gapprompt@gap>| !gapinput@S := FactorisableDualSymmetricInverseMonoid(5);;|
  !gapprompt@gap>| !gapinput@S := IdempotentGeneratedSubsemigroup(S);;|
  !gapprompt@gap>| !gapinput@IsSemilattice(S);|
  true
\end{Verbatim}
 }

 
\subsection{\textcolor{Chapter }{IsSimpleSemigroup}}\logpage{[ 12, 1, 23 ]}
\hyperdef{L}{X836F4692839F4874}{}
{
\noindent\textcolor{FuncColor}{$\triangleright$\enspace\texttt{IsSimpleSemigroup({\mdseries\slshape S})\index{IsSimpleSemigroup@\texttt{IsSimpleSemigroup}}
\label{IsSimpleSemigroup}
}\hfill{\scriptsize (property)}}\\
\noindent\textcolor{FuncColor}{$\triangleright$\enspace\texttt{IsCompletelySimpleSemigroup({\mdseries\slshape S})\index{IsCompletelySimpleSemigroup@\texttt{IsCompletelySimpleSemigroup}}
\label{IsCompletelySimpleSemigroup}
}\hfill{\scriptsize (property)}}\\
\textbf{\indent Returns:\ }
\texttt{true} or \texttt{false}.



 \texttt{IsSimpleSemigroup} returns \texttt{true} if the semigroup \mbox{\texttt{\mdseries\slshape S}} is simple and \texttt{false} if it is not.

 A semigroup is \emph{simple} if it has no proper 2\texttt{\symbol{45}}sided ideals. A semigroup is \emph{completely simple} if it is simple and possesses minimal left and right ideals. A finite
semigroup is simple if and only if it is completely simple. An inverse
semigroup is simple if and only if it is a group. 
\begin{Verbatim}[commandchars=!@|,fontsize=\small,frame=single,label=Example]
  !gapprompt@gap>| !gapinput@S := Semigroup(|
  !gapprompt@>| !gapinput@ Transformation([2, 2, 4, 4, 6, 6, 8, 8, 10, 10, 12, 12, 2]),|
  !gapprompt@>| !gapinput@ Transformation([1, 1, 3, 3, 5, 5, 7, 7, 9, 9, 11, 11, 3]),|
  !gapprompt@>| !gapinput@ Transformation([1, 7, 3, 9, 5, 11, 7, 1, 9, 3, 11, 5, 5]),|
  !gapprompt@>| !gapinput@ Transformation([7, 7, 9, 9, 11, 11, 1, 1, 3, 3, 5, 5, 7]));;|
  !gapprompt@gap>| !gapinput@IsSimpleSemigroup(S);|
  true
  !gapprompt@gap>| !gapinput@IsCompletelySimpleSemigroup(S);|
  true
  !gapprompt@gap>| !gapinput@IsSimpleSemigroup(MinimalIdeal(BrauerMonoid(6)));|
  true
  !gapprompt@gap>| !gapinput@R := Range(IsomorphismReesMatrixSemigroup(|
  !gapprompt@>| !gapinput@MinimalIdeal(BrauerMonoid(6))));|
  <Rees matrix semigroup 15x15 over Group(())>
\end{Verbatim}
 }

 

\subsection{\textcolor{Chapter }{IsSynchronizingSemigroup (for a transformation semigroup)}}
\logpage{[ 12, 1, 24 ]}\nobreak
\hyperdef{L}{X7EEC85187D315398}{}
{\noindent\textcolor{FuncColor}{$\triangleright$\enspace\texttt{IsSynchronizingSemigroup({\mdseries\slshape S})\index{IsSynchronizingSemigroup@\texttt{IsSynchronizingSemigroup}!for a transformation semigroup}
\label{IsSynchronizingSemigroup:for a transformation semigroup}
}\hfill{\scriptsize (property)}}\\
\textbf{\indent Returns:\ }
\texttt{true} or \texttt{false}.



 For a positive integer \mbox{\texttt{\mdseries\slshape n}}, \texttt{IsSynchronizingSemigroup} returns \texttt{true} if the semigroup of transformations \mbox{\texttt{\mdseries\slshape S}} contains a transformation with constant value on \texttt{[1 .. \mbox{\texttt{\mdseries\slshape n}}]} where \texttt{n} is the degree of the semigroup. See also \texttt{ConstantTransformation} (\textbf{Reference: ConstantTransformation}). 

 Note that the semigroup consisting of the identity transformation is the
unique transformation semigroup with degree \texttt{0}. In this special case, the function \texttt{IsSynchronizingSemigroup} will return \texttt{false}.

 
\begin{Verbatim}[commandchars=!@|,fontsize=\small,frame=single,label=Example]
  !gapprompt@gap>| !gapinput@S := Semigroup(|
  !gapprompt@>| !gapinput@Transformation([1, 1, 8, 7, 6, 6, 4, 1, 8, 9]),|
  !gapprompt@>| !gapinput@Transformation([5, 8, 7, 6, 10, 8, 7, 6, 9, 7]));;|
  !gapprompt@gap>| !gapinput@IsSynchronizingSemigroup(S);|
  true
  !gapprompt@gap>| !gapinput@S := Semigroup(|
  !gapprompt@>| !gapinput@Transformation([3, 8, 1, 1, 9, 9, 8, 7, 9, 6]),|
  !gapprompt@>| !gapinput@Transformation([7, 6, 8, 7, 5, 6, 8, 7, 8, 9]));;|
  !gapprompt@gap>| !gapinput@IsSynchronizingSemigroup(S);|
  false
  !gapprompt@gap>| !gapinput@Representative(MinimalIdeal(S));|
  Transformation( [ 8, 7, 7, 7, 8, 8, 7, 8, 8, 8 ] )
\end{Verbatim}
 }

 

\subsection{\textcolor{Chapter }{IsUnitRegularMonoid}}
\logpage{[ 12, 1, 25 ]}\nobreak
\hyperdef{L}{X80F9A4B87997839F}{}
{\noindent\textcolor{FuncColor}{$\triangleright$\enspace\texttt{IsUnitRegularMonoid({\mdseries\slshape S})\index{IsUnitRegularMonoid@\texttt{IsUnitRegularMonoid}}
\label{IsUnitRegularMonoid}
}\hfill{\scriptsize (property)}}\\
\textbf{\indent Returns:\ }
\texttt{true} if the semigroup \mbox{\texttt{\mdseries\slshape S}} is unit regular and \texttt{false} if it is not. 



 A monoid is \emph{unit regular} if and only if for every \texttt{{\textgreater}x} in \mbox{\texttt{\mdseries\slshape S}} there exists an element \texttt{y} in the group of units of \mbox{\texttt{\mdseries\slshape S}} such that \texttt{x*y*x=x}. 
\begin{Verbatim}[commandchars=!@|,fontsize=\small,frame=single,label=Example]
  !gapprompt@gap>| !gapinput@IsUnitRegularMonoid(FullTransformationMonoid(3));|
  true
\end{Verbatim}
 }

 

\subsection{\textcolor{Chapter }{IsZeroGroup}}
\logpage{[ 12, 1, 26 ]}\nobreak
\hyperdef{L}{X85F7E5CD86F0643B}{}
{\noindent\textcolor{FuncColor}{$\triangleright$\enspace\texttt{IsZeroGroup({\mdseries\slshape S})\index{IsZeroGroup@\texttt{IsZeroGroup}}
\label{IsZeroGroup}
}\hfill{\scriptsize (property)}}\\
\textbf{\indent Returns:\ }
\texttt{true} or \texttt{false}.



 \texttt{IsZeroGroup} returns \texttt{true} if the semigroup \mbox{\texttt{\mdseries\slshape S}} is a zero group and \texttt{false} if it is not.

 A semigroup \texttt{S} is a \emph{zero group} if there exists an element \texttt{z} in \texttt{S} such that \texttt{S} without \texttt{z} is a group and \texttt{x*z=z*x=z} for all \texttt{x} in \texttt{S}. Every zero group is an inverse semigroup. 
\begin{Verbatim}[commandchars=!@|,fontsize=\small,frame=single,label=Example]
  !gapprompt@gap>| !gapinput@S := Semigroup(Transformation([2, 2, 3, 4, 6, 8, 5, 5, 9]),|
  !gapprompt@>| !gapinput@                  Transformation([3, 3, 8, 2, 5, 6, 4, 4, 9]),|
  !gapprompt@>| !gapinput@                  ConstantTransformation(9, 9));;|
  !gapprompt@gap>| !gapinput@IsZeroGroup(S);|
  true
  !gapprompt@gap>| !gapinput@T := Range(IsomorphismPartialPermSemigroup(S));;|
  !gapprompt@gap>| !gapinput@IsZeroGroup(T);|
  true
  !gapprompt@gap>| !gapinput@IsZeroGroup(JonesMonoid(2));|
  true
\end{Verbatim}
 }

 

\subsection{\textcolor{Chapter }{IsZeroRectangularBand}}
\logpage{[ 12, 1, 27 ]}\nobreak
\hyperdef{L}{X7C6787D07B95B450}{}
{\noindent\textcolor{FuncColor}{$\triangleright$\enspace\texttt{IsZeroRectangularBand({\mdseries\slshape S})\index{IsZeroRectangularBand@\texttt{IsZeroRectangularBand}}
\label{IsZeroRectangularBand}
}\hfill{\scriptsize (property)}}\\
\textbf{\indent Returns:\ }
\texttt{true} or \texttt{false}.



 \texttt{IsZeroRectangularBand} returns \texttt{true} if the semigroup \mbox{\texttt{\mdseries\slshape S}} is a zero rectangular band and \texttt{false} if it is not.

 A semigroup is a \emph{0\texttt{\symbol{45}}rectangular band} if it is 0\texttt{\symbol{45}}simple and $\mathcal{H}$\texttt{\symbol{45}}trivial; see also \texttt{IsZeroSimpleSemigroup} (\ref{IsZeroSimpleSemigroup}) and \texttt{IsHTrivial} (\ref{IsHTrivial}). An inverse semigroup is a 0\texttt{\symbol{45}}rectangular band if and only
if it is a 0\texttt{\symbol{45}}group; see \texttt{IsZeroGroup} (\ref{IsZeroGroup}). 
\begin{Verbatim}[commandchars=!@|,fontsize=\small,frame=single,label=Example]
  !gapprompt@gap>| !gapinput@S := Semigroup(|
  !gapprompt@>| !gapinput@ Transformation([1, 3, 7, 9, 1, 12, 13, 1, 15, 9, 1, 18, 1, 1, 13,|
  !gapprompt@>| !gapinput@                 1, 1, 21, 1, 1, 1, 1, 1, 25, 26, 1]),|
  !gapprompt@>| !gapinput@Transformation([1, 5, 1, 5, 11, 1, 1, 14, 1, 16, 17, 1, 1, 19, 1,|
  !gapprompt@>| !gapinput@                11, 1, 1, 1, 23, 1, 16, 19, 1, 1, 1]),|
  !gapprompt@>| !gapinput@Transformation([1, 4, 8, 1, 10, 1, 8, 1, 1, 1, 10, 1, 8, 10, 1, 1,|
  !gapprompt@>| !gapinput@                20, 1, 22, 1, 8, 1, 1, 1, 1, 1]),|
  !gapprompt@>| !gapinput@Transformation([1, 6, 6, 1, 1, 1, 6, 1, 1, 1, 1, 1, 6, 1, 6, 1, 1,|
  !gapprompt@>| !gapinput@                6, 1, 1, 24, 1, 1, 1, 1, 6]));;|
  !gapprompt@gap>| !gapinput@D := DClass(S,|
  !gapprompt@>| !gapinput@Transformation([1, 8, 1, 1, 8, 1, 1, 1, 1, 1, 8, 1, 1, 8, 1, 1, 1,|
  !gapprompt@>| !gapinput@                1, 1, 1, 1, 1, 1, 1, 1, 1]));;|
  !gapprompt@gap>| !gapinput@IsZeroRectangularBand(Semigroup(D));|
  true
  !gapprompt@gap>| !gapinput@IsZeroRectangularBand(Semigroup(GreensDClasses(S)[1]));|
  false
\end{Verbatim}
 }

 

\subsection{\textcolor{Chapter }{IsZeroSemigroup}}
\logpage{[ 12, 1, 28 ]}\nobreak
\hyperdef{L}{X81A1882181B75CC9}{}
{\noindent\textcolor{FuncColor}{$\triangleright$\enspace\texttt{IsZeroSemigroup({\mdseries\slshape S})\index{IsZeroSemigroup@\texttt{IsZeroSemigroup}}
\label{IsZeroSemigroup}
}\hfill{\scriptsize (property)}}\\
\textbf{\indent Returns:\ }
\texttt{true} or \texttt{false}.



 \texttt{IsZeroSemigroup} returns \texttt{true} if the semigroup \mbox{\texttt{\mdseries\slshape S}} is a zero semigroup and \texttt{false} if it is not.

 A semigroup \texttt{S} is a \emph{zero semigroup} if there exists an element \texttt{z} in \texttt{S} such that \texttt{x*y=z} for all \texttt{x,y} in \texttt{S}. An inverse semigroup is a zero semigroup if and only if it is trivial. 
\begin{Verbatim}[commandchars=!@|,fontsize=\small,frame=single,label=Example]
  !gapprompt@gap>| !gapinput@S := Semigroup(|
  !gapprompt@>| !gapinput@Transformation([4, 7, 6, 3, 1, 5, 3, 6, 5, 9]),|
  !gapprompt@>| !gapinput@Transformation([5, 3, 5, 1, 9, 3, 8, 7, 4, 3]));;|
  !gapprompt@gap>| !gapinput@IsZeroSemigroup(S);|
  false
  !gapprompt@gap>| !gapinput@S := Semigroup(|
  !gapprompt@>| !gapinput@ Transformation([7, 8, 8, 8, 5, 8, 8, 8]),|
  !gapprompt@>| !gapinput@ Transformation([8, 8, 8, 8, 5, 7, 8, 8]),|
  !gapprompt@>| !gapinput@ Transformation([8, 7, 8, 8, 5, 8, 8, 8]),|
  !gapprompt@>| !gapinput@ Transformation([8, 8, 8, 7, 5, 8, 8, 8]),|
  !gapprompt@>| !gapinput@ Transformation([8, 8, 7, 8, 5, 8, 8, 8]));;|
  !gapprompt@gap>| !gapinput@IsZeroSemigroup(S);|
  true
  !gapprompt@gap>| !gapinput@MultiplicativeZero(S);|
  Transformation( [ 8, 8, 8, 8, 5, 8, 8, 8 ] )
\end{Verbatim}
 }

 

\subsection{\textcolor{Chapter }{IsZeroSimpleSemigroup}}
\logpage{[ 12, 1, 29 ]}\nobreak
\hyperdef{L}{X8193A60F839C064E}{}
{\noindent\textcolor{FuncColor}{$\triangleright$\enspace\texttt{IsZeroSimpleSemigroup({\mdseries\slshape S})\index{IsZeroSimpleSemigroup@\texttt{IsZeroSimpleSemigroup}}
\label{IsZeroSimpleSemigroup}
}\hfill{\scriptsize (property)}}\\
\textbf{\indent Returns:\ }
\texttt{true} or \texttt{false}.



 \texttt{IsZeroSimpleSemigroup} returns \texttt{true} if the semigroup \mbox{\texttt{\mdseries\slshape S}} is 0\texttt{\symbol{45}}simple and \texttt{false} if it is not.

 A semigroup is a \emph{0\texttt{\symbol{45}}simple} if it has no two\texttt{\symbol{45}}sided ideals other than itself and the set
containing the zero element; see also \texttt{MultiplicativeZero} (\ref{MultiplicativeZero}). An inverse semigroup is 0\texttt{\symbol{45}}simple if and only if it is a
Brandt semigroup; see \texttt{IsBrandtSemigroup} (\ref{IsBrandtSemigroup}). 
\begin{Verbatim}[commandchars=!@|,fontsize=\small,frame=single,label=Example]
  !gapprompt@gap>| !gapinput@S := Semigroup(|
  !gapprompt@>| !gapinput@ Transformation([1, 17, 17, 17, 17, 17, 17, 17, 17, 17, 5, 17,|
  !gapprompt@>| !gapinput@                 17, 17, 17, 17, 17]),|
  !gapprompt@>| !gapinput@ Transformation([1, 17, 17, 17, 11, 17, 17, 17, 17, 17, 17, 17,|
  !gapprompt@>| !gapinput@                 17, 17, 17, 17, 17]),|
  !gapprompt@>| !gapinput@ Transformation([1, 17, 17, 17, 17, 17, 17, 17, 17, 17, 4, 17,|
  !gapprompt@>| !gapinput@                 17, 17, 17, 17, 17]),|
  !gapprompt@>| !gapinput@ Transformation([1, 17, 17, 5, 17, 17, 17, 17, 17, 17, 17, 17,|
  !gapprompt@>| !gapinput@                 17, 17, 17, 17, 17]));;|
  !gapprompt@gap>| !gapinput@IsZeroSimpleSemigroup(S);|
  true
  !gapprompt@gap>| !gapinput@S := Semigroup(|
  !gapprompt@>| !gapinput@Transformation([2, 3, 4, 5, 1, 8, 7, 6, 2, 7]),|
  !gapprompt@>| !gapinput@Transformation([2, 3, 4, 5, 6, 8, 7, 1, 2, 2]));;|
  !gapprompt@gap>| !gapinput@IsZeroSimpleSemigroup(S);|
  false
\end{Verbatim}
 }

 

\subsection{\textcolor{Chapter }{IsSelfDualSemigroup}}
\logpage{[ 12, 1, 30 ]}\nobreak
\hyperdef{L}{X846FC6247EE31607}{}
{\noindent\textcolor{FuncColor}{$\triangleright$\enspace\texttt{IsSelfDualSemigroup({\mdseries\slshape S})\index{IsSelfDualSemigroup@\texttt{IsSelfDualSemigroup}}
\label{IsSelfDualSemigroup}
}\hfill{\scriptsize (property)}}\\
\textbf{\indent Returns:\ }
\texttt{true} or \texttt{false}.



 Returns \texttt{true} if the semigroup \mbox{\texttt{\mdseries\slshape S}} is self dual and \texttt{false} otherwise. 

 A semigroup is \emph{self dual} if it is isomorphic to its dual, that is, the semigroup \texttt{T} with multiplication \texttt{*} defined by \texttt{x*y=yx} where \texttt{yx} denotes the product in \mbox{\texttt{\mdseries\slshape S}}. 
\begin{Verbatim}[commandchars=!@|,fontsize=\small,frame=single,label=Example]
  !gapprompt@gap>| !gapinput@F := FreeSemigroup("a", "b");|
  <free semigroup on the generators [ a, b ]>
  !gapprompt@gap>| !gapinput@AssignGeneratorVariables(F);|
  !gapprompt@gap>| !gapinput@R := [[a ^ 3, a], [b ^ 2, b], [(a * b) ^ 2, a]];|
  [ [ a^3, a ], [ b^2, b ], [ (a*b)^2, a ] ]
  !gapprompt@gap>| !gapinput@S := F / R;|
  <fp semigroup with 2 generators and 3 relations of length 14>
  !gapprompt@gap>| !gapinput@IsSelfDualSemigroup(S);|
  false
  !gapprompt@gap>| !gapinput@IsSelfDualSemigroup(FreeBand(3));|
  true
  !gapprompt@gap>| !gapinput@S := DualSymmetricInverseMonoid(3);|
  <inverse block bijection monoid of degree 3 with 3 generators>
  !gapprompt@gap>| !gapinput@IsSelfDualSemigroup(S);|
  true
\end{Verbatim}
 }

 }

 
\section{\textcolor{Chapter }{ Inverse semigroups }}\logpage{[ 12, 2, 0 ]}
\hyperdef{L}{X80F2725581B166EE}{}
{
  In this section we describe the properties of an inverse semigroup or monoid
that can be determined using the \textsf{Semigroups} package. 

 

\subsection{\textcolor{Chapter }{IsCliffordSemigroup}}
\logpage{[ 12, 2, 1 ]}\nobreak
\hyperdef{L}{X81DE11987BB81017}{}
{\noindent\textcolor{FuncColor}{$\triangleright$\enspace\texttt{IsCliffordSemigroup({\mdseries\slshape S})\index{IsCliffordSemigroup@\texttt{IsCliffordSemigroup}}
\label{IsCliffordSemigroup}
}\hfill{\scriptsize (property)}}\\
\textbf{\indent Returns:\ }
\texttt{true} or \texttt{false}.



 \texttt{IsCliffordSemigroup} returns \texttt{true} if the semigroup \mbox{\texttt{\mdseries\slshape S}} is regular and its idempotents are central, and \texttt{false} if it is not. 
\begin{Verbatim}[commandchars=!@|,fontsize=\small,frame=single,label=Example]
  !gapprompt@gap>| !gapinput@S := Semigroup(Transformation([1, 2, 4, 5, 6, 3, 7, 8]),|
  !gapprompt@>| !gapinput@                  Transformation([3, 3, 4, 5, 6, 2, 7, 8]),|
  !gapprompt@>| !gapinput@                  Transformation([1, 2, 5, 3, 6, 8, 4, 4]));;|
  !gapprompt@gap>| !gapinput@IsCliffordSemigroup(S);|
  true
  !gapprompt@gap>| !gapinput@T := Range(IsomorphismPartialPermSemigroup(S));;|
  !gapprompt@gap>| !gapinput@IsCliffordSemigroup(S);|
  true
  !gapprompt@gap>| !gapinput@S := DualSymmetricInverseMonoid(5);;|
  !gapprompt@gap>| !gapinput@T := IdempotentGeneratedSubsemigroup(S);;|
  !gapprompt@gap>| !gapinput@IsCliffordSemigroup(T);|
  true
\end{Verbatim}
 }

 

\subsection{\textcolor{Chapter }{IsBrandtSemigroup}}
\logpage{[ 12, 2, 2 ]}\nobreak
\hyperdef{L}{X7EFDBA687DCDA6FA}{}
{\noindent\textcolor{FuncColor}{$\triangleright$\enspace\texttt{IsBrandtSemigroup({\mdseries\slshape S})\index{IsBrandtSemigroup@\texttt{IsBrandtSemigroup}}
\label{IsBrandtSemigroup}
}\hfill{\scriptsize (property)}}\\
\textbf{\indent Returns:\ }
\texttt{true} or \texttt{false}.



 \texttt{IsBrandtSemigroup} return \texttt{true} if the semigroup \mbox{\texttt{\mdseries\slshape S}} is a finite 0\texttt{\symbol{45}}simple inverse semigroup, and \texttt{false} if it is not. See also \texttt{IsZeroSimpleSemigroup} (\ref{IsZeroSimpleSemigroup}) and \texttt{IsInverseSemigroup} (\textbf{Reference: IsInverseSemigroup}). 
\begin{Verbatim}[commandchars=!@|,fontsize=\small,frame=single,label=Example]
  !gapprompt@gap>| !gapinput@S := Semigroup(|
  !gapprompt@>| !gapinput@Transformation([2, 8, 8, 8, 8, 8, 8, 8]),|
  !gapprompt@>| !gapinput@Transformation([5, 8, 8, 8, 8, 8, 8, 8]),|
  !gapprompt@>| !gapinput@Transformation([8, 3, 8, 8, 8, 8, 8, 8]),|
  !gapprompt@>| !gapinput@Transformation([8, 6, 8, 8, 8, 8, 8, 8]),|
  !gapprompt@>| !gapinput@Transformation([8, 8, 1, 8, 8, 8, 8, 8]),|
  !gapprompt@>| !gapinput@Transformation([8, 8, 8, 1, 8, 8, 8, 8]),|
  !gapprompt@>| !gapinput@Transformation([8, 8, 8, 8, 4, 8, 8, 8]),|
  !gapprompt@>| !gapinput@Transformation([8, 8, 8, 8, 8, 7, 8, 8]),|
  !gapprompt@>| !gapinput@Transformation([8, 8, 8, 8, 8, 8, 2, 8]));;|
  !gapprompt@gap>| !gapinput@IsBrandtSemigroup(S);|
  true
  !gapprompt@gap>| !gapinput@T := Range(IsomorphismPartialPermSemigroup(S));;|
  !gapprompt@gap>| !gapinput@IsBrandtSemigroup(T);|
  true
  !gapprompt@gap>| !gapinput@S := DualSymmetricInverseMonoid(4);;|
  !gapprompt@gap>| !gapinput@D := DClass(S,|
  !gapprompt@>| !gapinput@               Bipartition([[1, 2, 3, -1, -2, -3], [4, -4]]));;|
  !gapprompt@gap>| !gapinput@R := InjectionPrincipalFactor(D);;|
  !gapprompt@gap>| !gapinput@S := Semigroup(PreImages(R, GeneratorsOfSemigroup(Range(R))));;|
  !gapprompt@gap>| !gapinput@IsBrandtSemigroup(S);|
  true
\end{Verbatim}
 }

 

\subsection{\textcolor{Chapter }{IsEUnitaryInverseSemigroup}}
\logpage{[ 12, 2, 3 ]}\nobreak
\hyperdef{L}{X843EA0E37C968BBF}{}
{\noindent\textcolor{FuncColor}{$\triangleright$\enspace\texttt{IsEUnitaryInverseSemigroup({\mdseries\slshape S})\index{IsEUnitaryInverseSemigroup@\texttt{IsEUnitaryInverseSemigroup}}
\label{IsEUnitaryInverseSemigroup}
}\hfill{\scriptsize (property)}}\\
\textbf{\indent Returns:\ }
\texttt{true} or \texttt{false}.



 As described in Section 5.9 of \cite{Howie1995aa}, an inverse semigroup \mbox{\texttt{\mdseries\slshape S}} with semilattice of idempotents \mbox{\texttt{\mdseries\slshape E}} is \emph{E\texttt{\symbol{45}}unitary} if for  
\[ s \in S\textrm{ and }e \in E\textrm{: }es \in E \Rightarrow s \in E. \]
  

 Equivalently, \mbox{\texttt{\mdseries\slshape S}} is \emph{E\texttt{\symbol{45}}unitary} if \mbox{\texttt{\mdseries\slshape E}} is closed in the natural partial order (see Proposition 5.9.1 in \cite{Howie1995aa}):  
\[ \textrm{for } s \in S\textrm{ and }e \in E\textrm{: }e \le s \Rightarrow s \in
E. \]
  

 This condition is equivalent to \mbox{\texttt{\mdseries\slshape E}} being majorantly closed in \mbox{\texttt{\mdseries\slshape S}}. See \texttt{IdempotentGeneratedSubsemigroup} (\ref{IdempotentGeneratedSubsemigroup}) and \texttt{IsMajorantlyClosed} (\ref{IsMajorantlyClosed}). Hence an inverse semigroup of partial permutations, block bijections or
partial permutation bipartitions is \emph{E\texttt{\symbol{45}}unitary} if and only if the idempotent semilattice is majorantly closed. 
\begin{Verbatim}[commandchars=!@|,fontsize=\small,frame=single,label=Example]
  !gapprompt@gap>| !gapinput@S := InverseSemigroup(|
  !gapprompt@>| !gapinput@ PartialPerm([1, 2, 3, 4], [2, 3, 1, 6]),|
  !gapprompt@>| !gapinput@ PartialPerm([1, 2, 3, 5], [3, 2, 1, 6]));;|
  !gapprompt@gap>| !gapinput@IsEUnitaryInverseSemigroup(S);|
  true
  !gapprompt@gap>| !gapinput@e := IdempotentGeneratedSubsemigroup(S);;|
  !gapprompt@gap>| !gapinput@ForAll(Difference(S, e), x -> not ForAny(e, y -> y * x in e));|
  true
  !gapprompt@gap>| !gapinput@T := InverseSemigroup([|
  !gapprompt@>| !gapinput@ PartialPerm([1, 3, 4, 6, 8], [2, 5, 10, 7, 9]),|
  !gapprompt@>| !gapinput@ PartialPerm([1, 2, 3, 5, 6, 7, 8], [5, 8, 9, 2, 10, 1, 3]),|
  !gapprompt@>| !gapinput@ PartialPerm([1, 2, 3, 5, 6, 7, 9], [9, 8, 4, 1, 6, 7, 2])]);;|
  !gapprompt@gap>| !gapinput@IsEUnitaryInverseSemigroup(T);|
  false
  !gapprompt@gap>| !gapinput@U := InverseSemigroup([|
  !gapprompt@>| !gapinput@ PartialPerm([1, 2, 3, 4, 5], [2, 3, 4, 5, 1]),|
  !gapprompt@>| !gapinput@ PartialPerm([1, 2, 3, 4, 5], [2, 1, 3, 4, 5])]);;|
  !gapprompt@gap>| !gapinput@IsEUnitaryInverseSemigroup(U);|
  true
  !gapprompt@gap>| !gapinput@IsGroupAsSemigroup(U);|
  true
  !gapprompt@gap>| !gapinput@StructureDescription(U);|
  "S5"
\end{Verbatim}
 }

 

\subsection{\textcolor{Chapter }{IsFInverseSemigroup}}
\logpage{[ 12, 2, 4 ]}\nobreak
\hyperdef{L}{X86F942F48158DAC3}{}
{\noindent\textcolor{FuncColor}{$\triangleright$\enspace\texttt{IsFInverseSemigroup({\mdseries\slshape S})\index{IsFInverseSemigroup@\texttt{IsFInverseSemigroup}}
\label{IsFInverseSemigroup}
}\hfill{\scriptsize (property)}}\\
\textbf{\indent Returns:\ }
\texttt{true} or \texttt{false}.



 This functions determines whether a given semigroup is an
F\texttt{\symbol{45}}inverse semigroup. An F\texttt{\symbol{45}}inverse
semigroup is a semigroup which satisfies \texttt{IsEUnitaryInverseSemigroup} (\ref{IsEUnitaryInverseSemigroup}) as well as being isomorphic to some \texttt{McAlisterTripleSemigroup} (\ref{McAlisterTripleSemigroup}) where the \texttt{McAlisterTripleSemigroupPartialOrder} (\ref{McAlisterTripleSemigroupPartialOrder}) satisfies \texttt{IsJoinSemilatticeDigraph} (\textbf{Digraphs: IsJoinSemilatticeDigraph}). McAlister triple semigroups are a representation of
E\texttt{\symbol{45}}unitary inverse semigroups and more can be read about
them in Chapter \ref{McAlister triple semigroups}. 
\begin{Verbatim}[commandchars=!@|,fontsize=\small,frame=single,label=Example]
  !gapprompt@gap>| !gapinput@S := InverseMonoid([PartialPermNC([1, 2], [1, 2]),|
  !gapprompt@>| !gapinput@PartialPermNC([1, 2, 3], [1, 2, 3]),|
  !gapprompt@>| !gapinput@PartialPermNC([1, 2, 4], [1, 2, 4]),|
  !gapprompt@>| !gapinput@PartialPermNC([1, 2], [2, 1]), PartialPermNC([1, 2, 3], [2, 1, 3]),|
  !gapprompt@>| !gapinput@PartialPermNC([1, 2, 4], [2, 1, 4])]);;|
  !gapprompt@gap>| !gapinput@IsEUnitaryInverseSemigroup(S);|
  true
  !gapprompt@gap>| !gapinput@IsFInverseSemigroup(S);|
  false
  !gapprompt@gap>| !gapinput@IsFInverseSemigroup(IdempotentGeneratedSubsemigroup(S));|
  true
\end{Verbatim}
 }

 

\subsection{\textcolor{Chapter }{IsFInverseMonoid}}
\logpage{[ 12, 2, 5 ]}\nobreak
\hyperdef{L}{X864F1906858BB8CF}{}
{\noindent\textcolor{FuncColor}{$\triangleright$\enspace\texttt{IsFInverseMonoid({\mdseries\slshape S})\index{IsFInverseMonoid@\texttt{IsFInverseMonoid}}
\label{IsFInverseMonoid}
}\hfill{\scriptsize (property)}}\\
\textbf{\indent Returns:\ }
\texttt{true} or \texttt{false}.



 This function determines whether a given semigroup is an
F\texttt{\symbol{45}}inverse monoid. A semigroup is an
F\texttt{\symbol{45}}inverse monoid when it satisfies \texttt{IsMonoid} (\textbf{Reference: IsMonoid}) and \texttt{IsFInverseSemigroup} (\ref{IsFInverseSemigroup}). }

 

\subsection{\textcolor{Chapter }{IsFactorisableInverseMonoid}}
\logpage{[ 12, 2, 6 ]}\nobreak
\hyperdef{L}{X8440E22681BD1EE6}{}
{\noindent\textcolor{FuncColor}{$\triangleright$\enspace\texttt{IsFactorisableInverseMonoid({\mdseries\slshape S})\index{IsFactorisableInverseMonoid@\texttt{IsFactorisableInverseMonoid}}
\label{IsFactorisableInverseMonoid}
}\hfill{\scriptsize (property)}}\\
\textbf{\indent Returns:\ }
\texttt{true} or \texttt{false}.



 An inverse monoid is \emph{factorisable} if every element is the product of an element of the group of units and an
idempotent; see also \texttt{GroupOfUnits} (\ref{GroupOfUnits}) and \texttt{Idempotents} (\ref{Idempotents}). Hence an inverse semigroup of partial permutations is factorisable if and
only if each of its generators is the restriction of some element in the group
of units. 
\begin{Verbatim}[commandchars=!@|,fontsize=\small,frame=single,label=Example]
  !gapprompt@gap>| !gapinput@S := InverseSemigroup(|
  !gapprompt@>| !gapinput@PartialPerm([1, 2, 4], [3, 1, 4]),|
  !gapprompt@>| !gapinput@PartialPerm([1, 2, 3, 5], [4, 1, 5, 2]));;|
  !gapprompt@gap>| !gapinput@IsFactorisableInverseMonoid(S);|
  false
  !gapprompt@gap>| !gapinput@IsFactorisableInverseMonoid(SymmetricInverseSemigroup(5));|
  true
  !gapprompt@gap>| !gapinput@IsFactorisableInverseMonoid(DualSymmetricInverseMonoid(5));|
  false
  !gapprompt@gap>| !gapinput@S := FactorisableDualSymmetricInverseMonoid(5);;|
  !gapprompt@gap>| !gapinput@IsFactorisableInverseMonoid(S);|
  true
\end{Verbatim}
 }

 

\subsection{\textcolor{Chapter }{IsJoinIrreducible}}
\logpage{[ 12, 2, 7 ]}\nobreak
\hyperdef{L}{X817F9F3984FC842C}{}
{\noindent\textcolor{FuncColor}{$\triangleright$\enspace\texttt{IsJoinIrreducible({\mdseries\slshape S, x})\index{IsJoinIrreducible@\texttt{IsJoinIrreducible}}
\label{IsJoinIrreducible}
}\hfill{\scriptsize (operation)}}\\
\textbf{\indent Returns:\ }
 \texttt{true} or \texttt{false}. 



 \texttt{IsJoinIrreducible} determines whether an element \mbox{\texttt{\mdseries\slshape x}} of an inverse semigroup \mbox{\texttt{\mdseries\slshape S}} of partial permutations, block bijections or partial permutation bipartitions
is join irreducible.

 An element \mbox{\texttt{\mdseries\slshape x}} is \emph{join irreducible} when it is not the least upper bound (with respect to the natural partial
order \texttt{NaturalLeqPartialPerm} (\textbf{Reference: NaturalLeqPartialPerm})) of any subset of \mbox{\texttt{\mdseries\slshape S}} not containing \mbox{\texttt{\mdseries\slshape x}}. }

 
\begin{Verbatim}[commandchars=!@|,fontsize=\small,frame=single,label=Example]
  !gapprompt@gap>| !gapinput@S := SymmetricInverseSemigroup(3);|
  <symmetric inverse monoid of degree 3>
  !gapprompt@gap>| !gapinput@x := PartialPerm([1, 2, 3]);|
  <identity partial perm on [ 1, 2, 3 ]>
  !gapprompt@gap>| !gapinput@IsJoinIrreducible(S, x);|
  false
  !gapprompt@gap>| !gapinput@T := InverseSemigroup([|
  !gapprompt@>| !gapinput@ PartialPerm([1, 2, 4, 3]),|
  !gapprompt@>| !gapinput@ PartialPerm([1]),|
  !gapprompt@>| !gapinput@ PartialPerm([0, 2])]);|
  <inverse partial perm semigroup of rank 4 with 3 generators>
  !gapprompt@gap>| !gapinput@y := PartialPerm([1, 2, 3, 4]);|
  <identity partial perm on [ 1, 2, 3, 4 ]>
  !gapprompt@gap>| !gapinput@IsJoinIrreducible(T, y);|
  true
  !gapprompt@gap>| !gapinput@B := InverseSemigroup([|
  !gapprompt@>| !gapinput@ Bipartition([|
  !gapprompt@>| !gapinput@   [1, -5], [2, -2], [3, 5, 6, 7, -1, -4, -6, -7], [4, -3]]),|
  !gapprompt@>| !gapinput@ Bipartition([|
  !gapprompt@>| !gapinput@   [1, -1], [2, -3], [3, -4], [4, 5, 7, -2, -6, -7], [6, -5]]),|
  !gapprompt@>| !gapinput@ Bipartition([|
  !gapprompt@>| !gapinput@   [1, -2], [2, -4], [3, -6], [4, -1], [5, 7, -3, -7], [6, -5]]),|
  !gapprompt@>| !gapinput@ Bipartition([|
  !gapprompt@>| !gapinput@   [1, -5], [2, -1], [3, -6], [4, 5, 7, -2, -4, -7], [6, -3]])]);|
  <inverse block bijection semigroup of degree 7 with 4 generators>
  !gapprompt@gap>| !gapinput@x := Bipartition([|
  !gapprompt@>| !gapinput@  [1, 2, 3, 5, 6, 7, -2, -3, -4, -5, -6, -7], [4, -1]]);|
  <block bijection: [ 1, 2, 3, 5, 6, 7, -2, -3, -4, -5, -6, -7 ],
   [ 4, -1 ]>
  !gapprompt@gap>| !gapinput@IsJoinIrreducible(B, x);|
  true
  !gapprompt@gap>| !gapinput@IsJoinIrreducible(B, B.1);|
  false
\end{Verbatim}
 

\subsection{\textcolor{Chapter }{IsMajorantlyClosed}}
\logpage{[ 12, 2, 8 ]}\nobreak
\hyperdef{L}{X81E6D24F852A7937}{}
{\noindent\textcolor{FuncColor}{$\triangleright$\enspace\texttt{IsMajorantlyClosed({\mdseries\slshape S, T})\index{IsMajorantlyClosed@\texttt{IsMajorantlyClosed}}
\label{IsMajorantlyClosed}
}\hfill{\scriptsize (operation)}}\\
\textbf{\indent Returns:\ }
 \texttt{true} or \texttt{false}. 



 \texttt{IsMajorantlyClosed} determines whether the subset \mbox{\texttt{\mdseries\slshape T}} of the inverse semigroup of partial permutations, block bijections or partial
permutation bipartitions \mbox{\texttt{\mdseries\slshape S}} is majorantly closed in \mbox{\texttt{\mdseries\slshape S}}. See also \texttt{MajorantClosure} (\ref{MajorantClosure}).

 We say that \mbox{\texttt{\mdseries\slshape T}} is \emph{majorantly closed} in \mbox{\texttt{\mdseries\slshape S}} if it contains all elements of \mbox{\texttt{\mdseries\slshape S}} which are greater than or equal to any element of \mbox{\texttt{\mdseries\slshape T}}, with respect to the natural partial order. See \texttt{NaturalLeqPartialPerm} (\textbf{Reference: NaturalLeqPartialPerm}).

 Note that \mbox{\texttt{\mdseries\slshape T}} can be a subset of \mbox{\texttt{\mdseries\slshape S}} or a subsemigroup of \mbox{\texttt{\mdseries\slshape S}}. }

 
\begin{Verbatim}[commandchars=!@|,fontsize=\small,frame=single,label=Example]
  !gapprompt@gap>| !gapinput@S := SymmetricInverseSemigroup(2);|
  <symmetric inverse monoid of degree 2>
  !gapprompt@gap>| !gapinput@T := [Elements(S)[2]];|
  [ <identity partial perm on [ 1 ]> ]
  !gapprompt@gap>| !gapinput@IsMajorantlyClosed(S, T);|
  false
  !gapprompt@gap>| !gapinput@U := [Elements(S)[2], Elements(S)[6]];|
  [ <identity partial perm on [ 1 ]>, <identity partial perm on [ 1, 2 ]
      > ]
  !gapprompt@gap>| !gapinput@IsMajorantlyClosed(S, U);|
  true
  !gapprompt@gap>| !gapinput@D := DualSymmetricInverseSemigroup(3);|
  <inverse block bijection monoid of degree 3 with 3 generators>
  !gapprompt@gap>| !gapinput@x := Bipartition([[1, -2], [2, -3], [3, -1]]);;|
  !gapprompt@gap>| !gapinput@IsMajorantlyClosed(D, [x]);|
  true
  !gapprompt@gap>| !gapinput@y := Bipartition([[1, 2, -1, -2], [3, -3]]);;|
  !gapprompt@gap>| !gapinput@IsMajorantlyClosed(D, [x, y]);|
  false
\end{Verbatim}
 

\subsection{\textcolor{Chapter }{IsMonogenicInverseSemigroup}}
\logpage{[ 12, 2, 9 ]}\nobreak
\hyperdef{L}{X7D2641AD830DEC1C}{}
{\noindent\textcolor{FuncColor}{$\triangleright$\enspace\texttt{IsMonogenicInverseSemigroup({\mdseries\slshape S})\index{IsMonogenicInverseSemigroup@\texttt{IsMonogenicInverseSemigroup}}
\label{IsMonogenicInverseSemigroup}
}\hfill{\scriptsize (property)}}\\
\textbf{\indent Returns:\ }
\texttt{true} or \texttt{false}.



 \texttt{IsMonogenicInverseSemigroup} returns \texttt{true} if the semigroup \mbox{\texttt{\mdseries\slshape S}} is a monogenic inverse semigroup and it returns \texttt{false} if it is not. 

 A inverse semigroup is \emph{monogenic} if it is generated as an inverse semigroup by a single element. See also \texttt{IsMonogenicSemigroup} (\ref{IsMonogenicSemigroup}) and \texttt{IsMonogenicInverseMonoid} (\ref{IsMonogenicInverseMonoid}). 
\begin{Verbatim}[commandchars=!@|,fontsize=\small,frame=single,label=Example]
  !gapprompt@gap>| !gapinput@x := PartialPerm([1, 2, 3, 6, 8, 10], [2, 6, 7, 9, 1, 5]);;|
  !gapprompt@gap>| !gapinput@S := InverseSemigroup(x, x ^ 2, x ^ 3);;|
  !gapprompt@gap>| !gapinput@IsMonogenicSemigroup(S);|
  false
  !gapprompt@gap>| !gapinput@IsMonogenicInverseSemigroup(S);|
  true
  !gapprompt@gap>| !gapinput@x := RandomBlockBijection(100);;|
  !gapprompt@gap>| !gapinput@S := InverseSemigroup(x, x ^ 2, x ^ 20);;|
  !gapprompt@gap>| !gapinput@IsMonogenicInverseSemigroup(S);|
  true
  !gapprompt@gap>| !gapinput@S := SymmetricInverseSemigroup(5);;|
  !gapprompt@gap>| !gapinput@IsMonogenicInverseSemigroup(S);|
  false
\end{Verbatim}
 }

 

\subsection{\textcolor{Chapter }{IsMonogenicInverseMonoid}}
\logpage{[ 12, 2, 10 ]}\nobreak
\hyperdef{L}{X7EDFA6CA86645DBE}{}
{\noindent\textcolor{FuncColor}{$\triangleright$\enspace\texttt{IsMonogenicInverseMonoid({\mdseries\slshape S})\index{IsMonogenicInverseMonoid@\texttt{IsMonogenicInverseMonoid}}
\label{IsMonogenicInverseMonoid}
}\hfill{\scriptsize (property)}}\\
\textbf{\indent Returns:\ }
\texttt{true} or \texttt{false}.



 \texttt{IsMonogenicInverseMonoid} returns \texttt{true} if the monoid \mbox{\texttt{\mdseries\slshape S}} is a monogenic inverse monoid and it returns \texttt{false} if it is not. 

 A inverse monoid is \emph{monogenic} if it is generated as an inverse monoid by a single element. See also \texttt{IsMonogenicInverseSemigroup} (\ref{IsMonogenicInverseSemigroup}) and \texttt{IsMonogenicMonoid} (\ref{IsMonogenicMonoid}). 
\begin{Verbatim}[commandchars=!@|,fontsize=\small,frame=single,label=Example]
  !gapprompt@gap>| !gapinput@x := PartialPerm([1, 2, 3, 6, 8, 10], [2, 6, 7, 9, 1, 5]);;|
  !gapprompt@gap>| !gapinput@S := InverseMonoid(x, x ^ 2, x ^ 3);;|
  !gapprompt@gap>| !gapinput@IsMonogenicMonoid(S);|
  false
  !gapprompt@gap>| !gapinput@IsMonogenicInverseSemigroup(S);|
  false
  !gapprompt@gap>| !gapinput@IsMonogenicInverseMonoid(S);|
  true
  !gapprompt@gap>| !gapinput@x := RandomBlockBijection(100);;|
  !gapprompt@gap>| !gapinput@S := InverseMonoid(x, x ^ 2, x ^ 20);;|
  !gapprompt@gap>| !gapinput@IsMonogenicInverseMonoid(S);|
  true
  !gapprompt@gap>| !gapinput@S := SymmetricInverseMonoid(5);;|
  !gapprompt@gap>| !gapinput@IsMonogenicInverseMonoid(S);|
  false
\end{Verbatim}
 }

 }

   }

 
\chapter{\textcolor{Chapter }{Congruences}}\label{Congruences}
\logpage{[ 13, 0, 0 ]}
\hyperdef{L}{X82BD951079E3C349}{}
{
  Congruences in \textsf{Semigroups} can be described in several different ways:

 
\begin{itemize}
\item  Generating pairs \texttt{\symbol{45}}\texttt{\symbol{45}} the minimal
congruence which contains these pairs 
\item  Rees congruences \texttt{\symbol{45}}\texttt{\symbol{45}} the congruence
specified by a given ideal 
\item  Universal congruences \texttt{\symbol{45}}\texttt{\symbol{45}} the unique
congruence with only one class 
\item  Linked triples \texttt{\symbol{45}}\texttt{\symbol{45}} only for simple or
0\texttt{\symbol{45}}simple semigroups (see below) 
\item  Kernel and trace \texttt{\symbol{45}}\texttt{\symbol{45}} only for inverse
semigroups 
\item  Word graph \texttt{\symbol{45}}\texttt{\symbol{45}} only for congruences
created via \texttt{IteratorOfLeftCongruences} (\ref{IteratorOfLeftCongruences:for a semigroup}) or \texttt{IteratorOfRightCongruences} (\ref{IteratorOfRightCongruences:for a semigroup}) 
\item  Wang pairs \texttt{\symbol{45}}\texttt{\symbol{45}} only for graph inverse
semigroup 
\end{itemize}
 The operation \texttt{SemigroupCongruence} (\ref{SemigroupCongruence}) can be used to create any of these, interpreting the arguments in a smart way.
The usual way of specifying a congruence will be by giving a set of generating
pairs, but a user with an ideal could instead create a Rees congruence or
universal congruence.

 If a congruence is specified by generating pairs on a simple,
0\texttt{\symbol{45}}simple, or inverse semigroup, then the congruence may be
converted automatically to one of the last two items in the above list, to
reduce the complexity of any calculations to be performed. The user need not
manually specify, or even be aware of, the congruence's linked triple or
kernel and trace.

 We can also create left congruences and right congruences, using the \texttt{LeftSemigroupCongruence} (\ref{LeftSemigroupCongruence}) and \texttt{RightSemigroupCongruence} (\ref{RightSemigroupCongruence}) functions.

 Please note that congruence objects made in \textsf{GAP} before loading the \textsf{Semigroups} package may not behave correctly after \textsf{Semigroups} is loaded. If \textsf{Semigroups} is loaded at the beginning of the session, or before any congruence work is
done, then the objects should behave correctly.

 
\section{\textcolor{Chapter }{Semigroup congruence objects}}\logpage{[ 13, 1, 0 ]}
\hyperdef{L}{X784770137D98FEB9}{}
{
 

\subsection{\textcolor{Chapter }{IsSemigroupCongruence}}
\logpage{[ 13, 1, 1 ]}\nobreak
\hyperdef{L}{X78E34B737F0E009F}{}
{\noindent\textcolor{FuncColor}{$\triangleright$\enspace\texttt{IsSemigroupCongruence({\mdseries\slshape obj})\index{IsSemigroupCongruence@\texttt{IsSemigroupCongruence}}
\label{IsSemigroupCongruence}
}\hfill{\scriptsize (property)}}\\


 A semigroup congruence \texttt{cong} is an equivalence relation on a semigroup \texttt{S} which respects left and right multiplication. 

 That is, if $(a,b)$ is a pair in \texttt{cong}, and $x$ is an element of \texttt{S}, then $(ax,bx)$ and $(xa,xb)$ are both in \texttt{cong}. 

 The simplest way of creating a congruence in \textsf{Semigroups} is by a set of \emph{generating pairs}. See \texttt{SemigroupCongruence} (\ref{SemigroupCongruence}). 
\begin{Verbatim}[commandchars=!@|,fontsize=\small,frame=single,label=Example]
  !gapprompt@gap>| !gapinput@S := Semigroup([|
  !gapprompt@>| !gapinput@  Transformation([2, 1, 1, 2, 1]),|
  !gapprompt@>| !gapinput@  Transformation([3, 4, 3, 4, 4]),|
  !gapprompt@>| !gapinput@  Transformation([3, 4, 3, 4, 3]),|
  !gapprompt@>| !gapinput@  Transformation([4, 3, 3, 4, 4])]);;|
  !gapprompt@gap>| !gapinput@pair1 := [Transformation([3, 4, 3, 4, 3]),|
  !gapprompt@>| !gapinput@             Transformation([1, 2, 1, 2, 1])];;|
  !gapprompt@gap>| !gapinput@pair2 := [Transformation([4, 3, 4, 3, 4]),|
  !gapprompt@>| !gapinput@             Transformation([3, 4, 3, 4, 3])];;|
  !gapprompt@gap>| !gapinput@cong := SemigroupCongruence(S, [pair1, pair2]);|
  <semigroup congruence over <simple transformation semigroup of
   degree 5 with 4 generators> with linked triple (2,4,1)>
  !gapprompt@gap>| !gapinput@IsSemigroupCongruence(cong);|
  true
\end{Verbatim}
 }

 

\subsection{\textcolor{Chapter }{IsLeftSemigroupCongruence}}
\logpage{[ 13, 1, 2 ]}\nobreak
\hyperdef{L}{X7E909A78830D42A6}{}
{\noindent\textcolor{FuncColor}{$\triangleright$\enspace\texttt{IsLeftSemigroupCongruence({\mdseries\slshape obj})\index{IsLeftSemigroupCongruence@\texttt{IsLeftSemigroupCongruence}}
\label{IsLeftSemigroupCongruence}
}\hfill{\scriptsize (property)}}\\


 A left semigroup congruence \texttt{cong} is an equivalence relation on a semigroup \texttt{S} which respects left multiplication. 

 That is, if $(a,b)$ is a pair in \texttt{cong}, and $x$ is an element of \texttt{S}, then $(xa,xb)$ is also in \texttt{cong}. 

 The simplest way of creating a left congruence in \textsf{Semigroups} is by a set of \emph{generating pairs}. See \texttt{LeftSemigroupCongruence} (\ref{LeftSemigroupCongruence}). 
\begin{Verbatim}[commandchars=!@|,fontsize=\small,frame=single,label=Example]
  !gapprompt@gap>| !gapinput@S := Semigroup([|
  !gapprompt@>| !gapinput@  Transformation([2, 1, 1, 2, 1]),|
  !gapprompt@>| !gapinput@  Transformation([3, 4, 3, 4, 4]),|
  !gapprompt@>| !gapinput@  Transformation([3, 4, 3, 4, 3]),|
  !gapprompt@>| !gapinput@  Transformation([4, 3, 3, 4, 4])]);;|
  !gapprompt@gap>| !gapinput@pair1 := [Transformation([3, 4, 3, 4, 3]),|
  !gapprompt@>| !gapinput@             Transformation([1, 2, 1, 2, 1])];;|
  !gapprompt@gap>| !gapinput@pair2 := [Transformation([4, 3, 4, 3, 4]),|
  !gapprompt@>| !gapinput@             Transformation([3, 4, 3, 4, 3])];;|
  !gapprompt@gap>| !gapinput@cong := LeftSemigroupCongruence(S, [pair1, pair2]);|
  <left semigroup congruence over <transformation semigroup of degree 5
   with 4 generators> with 2 generating pairs>
  !gapprompt@gap>| !gapinput@IsLeftSemigroupCongruence(cong);|
  true
\end{Verbatim}
 }

 

\subsection{\textcolor{Chapter }{IsRightSemigroupCongruence}}
\logpage{[ 13, 1, 3 ]}\nobreak
\hyperdef{L}{X839EEA797B1CCB67}{}
{\noindent\textcolor{FuncColor}{$\triangleright$\enspace\texttt{IsRightSemigroupCongruence({\mdseries\slshape obj})\index{IsRightSemigroupCongruence@\texttt{IsRightSemigroupCongruence}}
\label{IsRightSemigroupCongruence}
}\hfill{\scriptsize (property)}}\\


 A right semigroup congruence \texttt{cong} is an equivalence relation on a semigroup \texttt{S} which respects right multiplication. 

 That is, if $(a,b)$ is a pair in \texttt{cong}, and $x$ is an element of \texttt{S}, then $(ax,bx)$ is also in \texttt{cong}. 

 The simplest way of creating a right congruence in \textsf{Semigroups} is by a set of \emph{generating pairs}. See \texttt{RightSemigroupCongruence} (\ref{RightSemigroupCongruence}). 
\begin{Verbatim}[commandchars=!@|,fontsize=\small,frame=single,label=Example]
  !gapprompt@gap>| !gapinput@S := Semigroup([|
  !gapprompt@>| !gapinput@  Transformation([2, 1, 1, 2, 1]),|
  !gapprompt@>| !gapinput@  Transformation([3, 4, 3, 4, 4]),|
  !gapprompt@>| !gapinput@  Transformation([3, 4, 3, 4, 3]),|
  !gapprompt@>| !gapinput@  Transformation([4, 3, 3, 4, 4])]);;|
  !gapprompt@gap>| !gapinput@pair1 := [Transformation([3, 4, 3, 4, 3]),|
  !gapprompt@>| !gapinput@             Transformation([1, 2, 1, 2, 1])];;|
  !gapprompt@gap>| !gapinput@pair2 := [Transformation([4, 3, 4, 3, 4]),|
  !gapprompt@>| !gapinput@             Transformation([3, 4, 3, 4, 3])];;|
  !gapprompt@gap>| !gapinput@RightSemigroupCongruence(S, [pair1, pair2]);|
  <right semigroup congruence over <transformation semigroup of
   degree 5 with 4 generators> with 2 generating pairs>
  !gapprompt@gap>| !gapinput@IsRightSemigroupCongruence(cong);|
  true
\end{Verbatim}
 }

 }

 
\section{\textcolor{Chapter }{ Creating congruences }}\logpage{[ 13, 2, 0 ]}
\hyperdef{L}{X7D49787B7B2589B2}{}
{
  

\subsection{\textcolor{Chapter }{SemigroupCongruence}}
\logpage{[ 13, 2, 1 ]}\nobreak
\hyperdef{L}{X85CE56AC84FA5D33}{}
{\noindent\textcolor{FuncColor}{$\triangleright$\enspace\texttt{SemigroupCongruence({\mdseries\slshape S, pairs})\index{SemigroupCongruence@\texttt{SemigroupCongruence}}
\label{SemigroupCongruence}
}\hfill{\scriptsize (function)}}\\
\textbf{\indent Returns:\ }
A semigroup congruence.



 This function returns a semigroup congruence over the semigroup \mbox{\texttt{\mdseries\slshape S}}.

 If \mbox{\texttt{\mdseries\slshape pairs}} is a list of lists of size 2 with elements from \mbox{\texttt{\mdseries\slshape S}}, then this function will return the semigroup congruence defined by these
generating pairs. The individual pairs may instead be given as separate
arguments.

 
\begin{Verbatim}[commandchars=!@|,fontsize=\small,frame=single,label=Example]
  !gapprompt@gap>| !gapinput@S := Semigroup([|
  !gapprompt@>| !gapinput@  Transformation([2, 1, 1, 2, 1]),|
  !gapprompt@>| !gapinput@  Transformation([3, 4, 3, 4, 4]),|
  !gapprompt@>| !gapinput@  Transformation([3, 4, 3, 4, 3]),|
  !gapprompt@>| !gapinput@  Transformation([4, 3, 3, 4, 4])]);;|
  !gapprompt@gap>| !gapinput@pair1 := [Transformation([3, 4, 3, 4, 3]),|
  !gapprompt@>| !gapinput@             Transformation([1, 2, 1, 2, 1])];;|
  !gapprompt@gap>| !gapinput@pair2 := [Transformation([4, 3, 4, 3, 4]),|
  !gapprompt@>| !gapinput@             Transformation([3, 4, 3, 4, 3])];;|
  !gapprompt@gap>| !gapinput@SemigroupCongruence(S, [pair1, pair2]);|
  <semigroup congruence over <simple transformation semigroup of
   degree 5 with 4 generators> with linked triple (2,4,1)>
  !gapprompt@gap>| !gapinput@SemigroupCongruence(S, pair1, pair2);|
  <semigroup congruence over <simple transformation semigroup of
   degree 5 with 4 generators> with linked triple (2,4,1)>
\end{Verbatim}
 }

 

\subsection{\textcolor{Chapter }{LeftSemigroupCongruence}}
\logpage{[ 13, 2, 2 ]}\nobreak
\hyperdef{L}{X8757DB087BE7E55A}{}
{\noindent\textcolor{FuncColor}{$\triangleright$\enspace\texttt{LeftSemigroupCongruence({\mdseries\slshape S, pairs})\index{LeftSemigroupCongruence@\texttt{LeftSemigroupCongruence}}
\label{LeftSemigroupCongruence}
}\hfill{\scriptsize (function)}}\\
\textbf{\indent Returns:\ }
A left semigroup congruence.



 This function returns a left semigroup congruence over the semigroup \mbox{\texttt{\mdseries\slshape S}}.

 If \mbox{\texttt{\mdseries\slshape pairs}} is a list of lists of size 2 with elements from \mbox{\texttt{\mdseries\slshape S}}, then this function will return the least left semigroup congruence on \mbox{\texttt{\mdseries\slshape S}} which contains these generating pairs. The individual pairs may instead be
given as separate arguments.

 
\begin{Verbatim}[commandchars=!@|,fontsize=\small,frame=single,label=Example]
  !gapprompt@gap>| !gapinput@S := Semigroup([|
  !gapprompt@>| !gapinput@  Transformation([2, 1, 1, 2, 1]),|
  !gapprompt@>| !gapinput@  Transformation([3, 4, 3, 4, 4]),|
  !gapprompt@>| !gapinput@  Transformation([3, 4, 3, 4, 3]),|
  !gapprompt@>| !gapinput@  Transformation([4, 3, 3, 4, 4])]);;|
  !gapprompt@gap>| !gapinput@pair1 := [Transformation([3, 4, 3, 4, 3]),|
  !gapprompt@>| !gapinput@             Transformation([1, 2, 1, 2, 1])];;|
  !gapprompt@gap>| !gapinput@pair2 := [Transformation([4, 3, 4, 3, 4]),|
  !gapprompt@>| !gapinput@             Transformation([3, 4, 3, 4, 3])];;|
  !gapprompt@gap>| !gapinput@LeftSemigroupCongruence(S, [pair1, pair2]);|
  <left semigroup congruence over <transformation semigroup of degree 5
   with 4 generators> with 2 generating pairs>
  !gapprompt@gap>| !gapinput@LeftSemigroupCongruence(S, pair1, pair2);|
  <left semigroup congruence over <transformation semigroup of degree 5
   with 4 generators> with 2 generating pairs>
\end{Verbatim}
 }

 

\subsection{\textcolor{Chapter }{RightSemigroupCongruence}}
\logpage{[ 13, 2, 3 ]}\nobreak
\hyperdef{L}{X7D01176B788FEE60}{}
{\noindent\textcolor{FuncColor}{$\triangleright$\enspace\texttt{RightSemigroupCongruence({\mdseries\slshape S, pairs})\index{RightSemigroupCongruence@\texttt{RightSemigroupCongruence}}
\label{RightSemigroupCongruence}
}\hfill{\scriptsize (function)}}\\
\textbf{\indent Returns:\ }
A right semigroup congruence.



 This function returns a right semigroup congruence over the semigroup \mbox{\texttt{\mdseries\slshape S}}.

 If \mbox{\texttt{\mdseries\slshape pairs}} is a list of lists of size 2 with elements from \mbox{\texttt{\mdseries\slshape S}}, then this function will return the least right semigroup congruence on \mbox{\texttt{\mdseries\slshape S}} which contains these generating pairs. The individual pairs may instead be
given as separate arguments.

 
\begin{Verbatim}[commandchars=!@|,fontsize=\small,frame=single,label=Example]
  !gapprompt@gap>| !gapinput@S := Semigroup([|
  !gapprompt@>| !gapinput@  Transformation([2, 1, 1, 2, 1]),|
  !gapprompt@>| !gapinput@  Transformation([3, 4, 3, 4, 4]),|
  !gapprompt@>| !gapinput@  Transformation([3, 4, 3, 4, 3]),|
  !gapprompt@>| !gapinput@  Transformation([4, 3, 3, 4, 4])]);;|
  !gapprompt@gap>| !gapinput@pair1 := [Transformation([3, 4, 3, 4, 3]),|
  !gapprompt@>| !gapinput@             Transformation([1, 2, 1, 2, 1])];;|
  !gapprompt@gap>| !gapinput@pair2 := [Transformation([4, 3, 4, 3, 4]),|
  !gapprompt@>| !gapinput@             Transformation([3, 4, 3, 4, 3])];;|
  !gapprompt@gap>| !gapinput@RightSemigroupCongruence(S, [pair1, pair2]);|
  <right semigroup congruence over <transformation semigroup of
   degree 5 with 4 generators> with 2 generating pairs>
  !gapprompt@gap>| !gapinput@RightSemigroupCongruence(S, pair1, pair2);|
  <right semigroup congruence over <transformation semigroup of
   degree 5 with 4 generators> with 2 generating pairs>
\end{Verbatim}
 }

 }

 
\section{\textcolor{Chapter }{ Congruence classes }}\logpage{[ 13, 3, 0 ]}
\hyperdef{L}{X7D65BB067A762CD6}{}
{
  The main operations and attributes for congruences in the \textsf{GAP} library are: 
\begin{itemize}
\item  \texttt{EquivalenceClasses} (\textbf{Reference: EquivalenceClasses attribute}) 
\item  \texttt{NrEquivalenceClasses} 
\item  \texttt{EquivalenceClassOfElement} (\textbf{Reference: EquivalenceClassOfElement}) 
\end{itemize}
 

\subsection{\textcolor{Chapter }{IsCongruenceClass}}
\logpage{[ 13, 3, 1 ]}\nobreak
\hyperdef{L}{X7B1F67A97E23E6A4}{}
{\noindent\textcolor{FuncColor}{$\triangleright$\enspace\texttt{IsCongruenceClass({\mdseries\slshape obj})\index{IsCongruenceClass@\texttt{IsCongruenceClass}}
\label{IsCongruenceClass}
}\hfill{\scriptsize (category)}}\\


 This category contains any object which is an equivalence class of a semigroup
congruence (see \texttt{IsSemigroupCongruence} (\ref{IsSemigroupCongruence})). An object will only be in this category if the relation is known to be a
semigroup congruence when the congruence class is created. 
\begin{Verbatim}[commandchars=!@|,fontsize=\small,frame=single,label=Example]
  !gapprompt@gap>| !gapinput@S := Monoid([|
  !gapprompt@>| !gapinput@ Transformation([1, 2, 2]), Transformation([3, 1, 3])]);;|
  !gapprompt@gap>| !gapinput@cong := SemigroupCongruence(S, [Transformation([1, 2, 1]),|
  !gapprompt@>| !gapinput@                                   Transformation([2, 1, 2])]);;|
  !gapprompt@gap>| !gapinput@class := EquivalenceClassOfElement(cong,|
  !gapprompt@>| !gapinput@                                      Transformation([3, 1, 1]));|
  <2-sided congruence class of Transformation( [ 3, 1, 1 ] )>
  !gapprompt@gap>| !gapinput@IsCongruenceClass(class);|
  true
\end{Verbatim}
 }

 

\subsection{\textcolor{Chapter }{IsLeftCongruenceClass}}
\logpage{[ 13, 3, 2 ]}\nobreak
\hyperdef{L}{X7C803E8C84E81A0B}{}
{\noindent\textcolor{FuncColor}{$\triangleright$\enspace\texttt{IsLeftCongruenceClass({\mdseries\slshape obj})\index{IsLeftCongruenceClass@\texttt{IsLeftCongruenceClass}}
\label{IsLeftCongruenceClass}
}\hfill{\scriptsize (category)}}\\


 This category contains any object which is an equivalence class of a left
semigroup congruence (see \texttt{IsLeftSemigroupCongruence} (\ref{IsLeftSemigroupCongruence})). An object will only be in this category if the relation is known to be a
left semigroup congruence when the class is created. 
\begin{Verbatim}[commandchars=!@|,fontsize=\small,frame=single,label=Example]
  !gapprompt@gap>| !gapinput@S := Monoid([|
  !gapprompt@>| !gapinput@ Transformation([1, 2, 2]), Transformation([3, 1, 3])]);;|
  !gapprompt@gap>| !gapinput@pairs := [Transformation([1, 2, 1]),|
  !gapprompt@>| !gapinput@             Transformation([2, 1, 2])];;|
  !gapprompt@gap>| !gapinput@cong := LeftSemigroupCongruence(S, pairs);;|
  !gapprompt@gap>| !gapinput@class := EquivalenceClassOfElement(cong,|
  !gapprompt@>| !gapinput@                                      Transformation([3, 1, 1]));|
  <left congruence class of Transformation( [ 3, 1, 1 ] )>
  !gapprompt@gap>| !gapinput@IsLeftCongruenceClass(class);|
  true
\end{Verbatim}
 }

 

\subsection{\textcolor{Chapter }{IsRightCongruenceClass}}
\logpage{[ 13, 3, 3 ]}\nobreak
\hyperdef{L}{X7D2F11C28470BC5A}{}
{\noindent\textcolor{FuncColor}{$\triangleright$\enspace\texttt{IsRightCongruenceClass({\mdseries\slshape obj})\index{IsRightCongruenceClass@\texttt{IsRightCongruenceClass}}
\label{IsRightCongruenceClass}
}\hfill{\scriptsize (category)}}\\


 This category contains any object which is an equivalence class of a right
semigroup congruence (see \texttt{IsRightSemigroupCongruence} (\ref{IsRightSemigroupCongruence})). An object will only be in this category if the relation is known to be a
right semigroup congruence when the class is created. 
\begin{Verbatim}[commandchars=!@|,fontsize=\small,frame=single,label=Example]
  !gapprompt@gap>| !gapinput@S := Monoid([|
  !gapprompt@>| !gapinput@ Transformation([1, 2, 2]), Transformation([3, 1, 3])]);;|
  !gapprompt@gap>| !gapinput@pairs := [Transformation([1, 2, 1]),|
  !gapprompt@>| !gapinput@             Transformation([2, 1, 2])];;|
  !gapprompt@gap>| !gapinput@cong := RightSemigroupCongruence(S, pairs);;|
  !gapprompt@gap>| !gapinput@class := EquivalenceClassOfElement(cong,|
  !gapprompt@>| !gapinput@                                      Transformation([3, 1, 1]));|
  <right congruence class of Transformation( [ 3, 1, 1 ] )>
  !gapprompt@gap>| !gapinput@IsRightCongruenceClass(class);|
  true
\end{Verbatim}
 }

 

\subsection{\textcolor{Chapter }{NonTrivialEquivalenceClasses}}
\logpage{[ 13, 3, 4 ]}\nobreak
\hyperdef{L}{X86C05F31797C1D6D}{}
{\noindent\textcolor{FuncColor}{$\triangleright$\enspace\texttt{NonTrivialEquivalenceClasses({\mdseries\slshape eq})\index{NonTrivialEquivalenceClasses@\texttt{NonTrivialEquivalenceClasses}}
\label{NonTrivialEquivalenceClasses}
}\hfill{\scriptsize (attribute)}}\\
\textbf{\indent Returns:\ }
A list of equivalence classes.



 If \mbox{\texttt{\mdseries\slshape eq}} is an equivalence relation, then this attribute returns a list of all
equivalence classes of \mbox{\texttt{\mdseries\slshape eq}} which contain more than one element. 

 
\begin{Verbatim}[commandchars=!@|,fontsize=\small,frame=single,label=Example]
  !gapprompt@gap>| !gapinput@S := Monoid([Transformation([1, 2, 2]),|
  !gapprompt@>| !gapinput@                Transformation([3, 1, 3])]);;|
  !gapprompt@gap>| !gapinput@cong := SemigroupCongruence(S, [Transformation([1, 2, 1]),|
  !gapprompt@>| !gapinput@                                   Transformation([2, 1, 2])]);;|
  !gapprompt@gap>| !gapinput@classes := NonTrivialEquivalenceClasses(cong);;|
  !gapprompt@gap>| !gapinput@Set(classes);|
  [ <2-sided congruence class of Transformation( [ 1, 2, 2 ] )>, 
    <2-sided congruence class of Transformation( [ 3, 1, 1 ] )>, 
    <2-sided congruence class of Transformation( [ 3, 1, 3 ] )>, 
    <2-sided congruence class of Transformation( [ 2, 1, 2 ] )>, 
    <2-sided congruence class of Transformation( [ 3, 3, 3 ] )> ]
  !gapprompt@gap>| !gapinput@cong := RightSemigroupCongruence(S, [Transformation([1, 2, 1]),|
  !gapprompt@>| !gapinput@                                        Transformation([2, 1, 2])]);;|
  !gapprompt@gap>| !gapinput@classes := NonTrivialEquivalenceClasses(cong);;|
  !gapprompt@gap>| !gapinput@Set(classes);|
  [ <right congruence class of Transformation( [ 3, 1, 3 ] )>,
    <right congruence class of Transformation( [ 2, 1, 2 ] )> ]
\end{Verbatim}
 }

 

\subsection{\textcolor{Chapter }{EquivalenceRelationLookup (for an equivalence relation over a finite semigroup)}}
\logpage{[ 13, 3, 5 ]}\nobreak
\hyperdef{L}{X7DA4BABC7891A7F1}{}
{\noindent\textcolor{FuncColor}{$\triangleright$\enspace\texttt{EquivalenceRelationLookup({\mdseries\slshape equiv})\index{EquivalenceRelationLookup@\texttt{EquivalenceRelationLookup}!for an equivalence relation over a finite semigroup}
\label{EquivalenceRelationLookup:for an equivalence relation over a finite semigroup}
}\hfill{\scriptsize (attribute)}}\\
\textbf{\indent Returns:\ }
A list.



 This attribute describes the equivalence relation \mbox{\texttt{\mdseries\slshape equiv}}, defined over a finite semigroup, as a list of positive integers of length
the size of the finite semigroup over which \mbox{\texttt{\mdseries\slshape equiv}} is defined.

 Each position in the list corresponds to an element of the semigroup (in a
consistent canonical order) and the integer at that position is a unique
identifier for that element's equivalence class under \mbox{\texttt{\mdseries\slshape equiv}}. Two elements of the semigroup on which the equivalence is defined are
related in the equivalence if and only if they have the same number at their
respective positions in the lookup.

 Note that the order in which numbers appear in the list is
non\texttt{\symbol{45}}deterministic, and two equivalence relations describing
the same mathematical relation might therefore have different lookups. Note
also that the maximum value of the list may not be the number of classes of \mbox{\texttt{\mdseries\slshape equiv}}, and that any integer might not be included. However, see \texttt{EquivalenceRelationCanonicalLookup} (\ref{EquivalenceRelationCanonicalLookup:for an equivalence relation over a finite semigroup}). 

 See also \texttt{EquivalenceRelationPartition} (\textbf{Reference: EquivalenceRelationPartition}). 
\begin{Verbatim}[commandchars=!@|,fontsize=\small,frame=single,label=Example]
  !gapprompt@gap>| !gapinput@S := Monoid([|
  !gapprompt@>| !gapinput@ Transformation([1, 2, 2]), Transformation([3, 1, 3])]);;|
  !gapprompt@gap>| !gapinput@cong := SemigroupCongruence(S,|
  !gapprompt@>| !gapinput@[Transformation([1, 2, 1]), Transformation([2, 1, 2])]);;|
  !gapprompt@gap>| !gapinput@lookup := EquivalenceRelationLookup(cong);;|
  !gapprompt@gap>| !gapinput@lookup[3] = lookup[8];|
  true
  !gapprompt@gap>| !gapinput@lookup[2] = lookup[9];|
  false
\end{Verbatim}
 }

 

\subsection{\textcolor{Chapter }{EquivalenceRelationCanonicalLookup (for an equivalence relation over a finite semigroup)}}
\logpage{[ 13, 3, 6 ]}\nobreak
\hyperdef{L}{X8022B7898553F624}{}
{\noindent\textcolor{FuncColor}{$\triangleright$\enspace\texttt{EquivalenceRelationCanonicalLookup({\mdseries\slshape equiv})\index{EquivalenceRelationCanonicalLookup@\texttt{EquivalenceRelationCanonicalLookup}!for an equivalence relation over a finite semigroup}
\label{EquivalenceRelationCanonicalLookup:for an equivalence relation over a finite semigroup}
}\hfill{\scriptsize (attribute)}}\\
\textbf{\indent Returns:\ }
A list.



 This attribute describes the equivalence relation \mbox{\texttt{\mdseries\slshape equiv}}, defined over a finite semigroup, as a list of positive integers of length
the size of the semigroup.

 Each position in the list corresponds to an element of the semigroup (in a
consistent canonical order as defined by \texttt{PositionCanonical} (\ref{PositionCanonical})) and the integer at that position is a unique identifier for that element's
equivalence class under \mbox{\texttt{\mdseries\slshape equiv}}. The value of \texttt{EquivalenceRelationCanonicalLookup} has the property that the first appearance of the value \texttt{i} is strictly later than the first appearance of \texttt{i\texttt{\symbol{45}}1}, and that all entries in the list will be from the range \texttt{[1 .. NrEquivalenceClasses(\mbox{\texttt{\mdseries\slshape equiv}})]}. As such, two equivalence relations on a given semigroup are equal if and
only if their canonical lookups are equal.

 Two elements of the semigroup on which the equivalence relation is defined are
related in the equivalence relation if and only if they have the same number
at their respective positions in the lookup. 

 See also \texttt{EquivalenceRelationLookup} (\ref{EquivalenceRelationLookup:for an equivalence relation over a finite semigroup}) and \texttt{EquivalenceRelationPartition} (\textbf{Reference: EquivalenceRelationPartition}). 
\begin{Verbatim}[commandchars=!@|,fontsize=\small,frame=single,label=Example]
  !gapprompt@gap>| !gapinput@S := Monoid([|
  !gapprompt@>| !gapinput@ Transformation([1, 2, 2]), Transformation([3, 1, 3])]);;|
  !gapprompt@gap>| !gapinput@cong := SemigroupCongruence(S,|
  !gapprompt@>| !gapinput@[Transformation([1, 2, 1]), Transformation([2, 1, 2])]);;|
  !gapprompt@gap>| !gapinput@EquivalenceRelationCanonicalLookup(cong);|
  [ 1, 2, 3, 4, 5, 6, 2, 3, 6, 4, 5, 6 ]
\end{Verbatim}
 }

 

\subsection{\textcolor{Chapter }{EquivalenceRelationCanonicalPartition}}
\logpage{[ 13, 3, 7 ]}\nobreak
\hyperdef{L}{X842D567F79648FEB}{}
{\noindent\textcolor{FuncColor}{$\triangleright$\enspace\texttt{EquivalenceRelationCanonicalPartition({\mdseries\slshape cong})\index{EquivalenceRelationCanonicalPartition@\texttt{Equivalence}\-\texttt{Relation}\-\texttt{Canonical}\-\texttt{Partition}}
\label{EquivalenceRelationCanonicalPartition}
}\hfill{\scriptsize (attribute)}}\\
\textbf{\indent Returns:\ }
A list of lists.



 This attribute returns a list of lists of elements of the underlying set of
the semigroup congruence \mbox{\texttt{\mdseries\slshape cong}}. These lists are precisely the nontrivial equivalence classes of \mbox{\texttt{\mdseries\slshape cong}}. The order in which the classes appear is deterministic, and the order of the
elements inside each class is also deterministic. Hence, two congruence
objects have the same \texttt{EquivalenceRelationCanonicalPartition} if and only if they describe the same relation. 

 See also \texttt{EquivalenceRelationPartition} (\textbf{Reference: EquivalenceRelationPartition}), a similar attribute which does not have canonical ordering, but which is
likely to be faster. 

 
\begin{Verbatim}[commandchars=!@|,fontsize=\small,frame=single,label=Example]
  !gapprompt@gap>| !gapinput@S := Semigroup(Transformation([1, 4, 3, 3]),|
  !gapprompt@>| !gapinput@                  Transformation([2, 4, 3, 3]));;|
  !gapprompt@gap>| !gapinput@cong := SemigroupCongruence(S, [Transformation([1, 4, 3, 3]),|
  !gapprompt@>| !gapinput@                                   Transformation([1, 3, 3, 3])]);;|
  !gapprompt@gap>| !gapinput@EquivalenceRelationCanonicalPartition(cong);|
  [ [ Transformation( [ 1, 4, 3, 3 ] ),
        Transformation( [ 1, 3, 3, 3 ] ) ],
    [ Transformation( [ 4, 3, 3, 3 ] ),
        Transformation( [ 3, 3, 3, 3 ] ) ] ]
\end{Verbatim}
 }

 

\subsection{\textcolor{Chapter }{OnLeftCongruenceClasses}}
\logpage{[ 13, 3, 8 ]}\nobreak
\hyperdef{L}{X85400841879E41A5}{}
{\noindent\textcolor{FuncColor}{$\triangleright$\enspace\texttt{OnLeftCongruenceClasses({\mdseries\slshape class, elm})\index{OnLeftCongruenceClasses@\texttt{OnLeftCongruenceClasses}}
\label{OnLeftCongruenceClasses}
}\hfill{\scriptsize (operation)}}\\
\textbf{\indent Returns:\ }
A left congruence class.



 If \mbox{\texttt{\mdseries\slshape class}} is an equivalence class of the left semigroup congruence \texttt{cong} on the semigroup \texttt{S}, and \mbox{\texttt{\mdseries\slshape elm}} is an element of \texttt{S}, then this operation returns the equivalence class of \texttt{cong} containing the element \texttt{\mbox{\texttt{\mdseries\slshape elm}} * x}, where \texttt{x} is any element of \mbox{\texttt{\mdseries\slshape class}}. The result is well\texttt{\symbol{45}}defined by the definition of a left
congruence. 

 See \texttt{IsLeftSemigroupCongruence} (\ref{IsLeftSemigroupCongruence}) and \texttt{IsLeftCongruenceClass} (\ref{IsLeftCongruenceClass}). 

 
\begin{Verbatim}[commandchars=!@|,fontsize=\small,frame=single,label=Example]
  !gapprompt@gap>| !gapinput@S := Semigroup([|
  !gapprompt@>| !gapinput@  Transformation([2, 1, 1, 2, 1]),|
  !gapprompt@>| !gapinput@  Transformation([3, 4, 3, 4, 4]),|
  !gapprompt@>| !gapinput@  Transformation([3, 4, 3, 4, 3]),|
  !gapprompt@>| !gapinput@  Transformation([4, 3, 3, 4, 4])]);;|
  !gapprompt@gap>| !gapinput@pair1 := [Transformation([3, 4, 3, 4, 3]),|
  !gapprompt@>| !gapinput@             Transformation([1, 2, 1, 2, 1])];;|
  !gapprompt@gap>| !gapinput@pair2 := [Transformation([4, 3, 4, 3, 4]),|
  !gapprompt@>| !gapinput@             Transformation([3, 4, 3, 4, 3])];;|
  !gapprompt@gap>| !gapinput@cong := LeftSemigroupCongruence(S, [pair1, pair2]);|
  <left semigroup congruence over <transformation semigroup of degree 5
   with 4 generators> with 2 generating pairs>
  !gapprompt@gap>| !gapinput@x := Transformation([3, 4, 3, 4, 3]);;|
  !gapprompt@gap>| !gapinput@class := EquivalenceClassOfElement(cong, x);|
  <left congruence class of Transformation( [ 3, 4, 3, 4, 3 ] )>
  !gapprompt@gap>| !gapinput@elm := Transformation([1, 2, 2, 1, 2]);;|
  !gapprompt@gap>| !gapinput@OnLeftCongruenceClasses(class, elm);|
  <left congruence class of Transformation( [ 3, 4, 4, 3, 4 ] )>
\end{Verbatim}
 }

 

\subsection{\textcolor{Chapter }{OnRightCongruenceClasses}}
\logpage{[ 13, 3, 9 ]}\nobreak
\hyperdef{L}{X7D66F8607B4F0D8F}{}
{\noindent\textcolor{FuncColor}{$\triangleright$\enspace\texttt{OnRightCongruenceClasses({\mdseries\slshape class, elm})\index{OnRightCongruenceClasses@\texttt{OnRightCongruenceClasses}}
\label{OnRightCongruenceClasses}
}\hfill{\scriptsize (operation)}}\\
\textbf{\indent Returns:\ }
A right congruence class.



 If \mbox{\texttt{\mdseries\slshape class}} is an equivalence class of the right semigroup congruence \texttt{cong} on the semigroup \texttt{S}, and \mbox{\texttt{\mdseries\slshape elm}} is an element of \texttt{S}, then this operation returns the equivalence class of \texttt{cong} containing the element \texttt{x * \mbox{\texttt{\mdseries\slshape elm}}}, where \texttt{x} is any element of \mbox{\texttt{\mdseries\slshape class}}. The result is well\texttt{\symbol{45}}defined by the definition of a right
congruence. 

 See \texttt{IsRightSemigroupCongruence} (\ref{IsRightSemigroupCongruence}) and \texttt{IsRightCongruenceClass} (\ref{IsRightCongruenceClass}). 

 
\begin{Verbatim}[commandchars=!@|,fontsize=\small,frame=single,label=Example]
  !gapprompt@gap>| !gapinput@S := Semigroup([|
  !gapprompt@>| !gapinput@  Transformation([2, 1, 1, 2, 1]),|
  !gapprompt@>| !gapinput@  Transformation([3, 4, 3, 4, 4]),|
  !gapprompt@>| !gapinput@  Transformation([3, 4, 3, 4, 3]),|
  !gapprompt@>| !gapinput@  Transformation([4, 3, 3, 4, 4])]);;|
  !gapprompt@gap>| !gapinput@pair1 := [Transformation([3, 4, 3, 4, 3]),|
  !gapprompt@>| !gapinput@             Transformation([1, 2, 1, 2, 1])];;|
  !gapprompt@gap>| !gapinput@pair2 := [Transformation([4, 3, 4, 3, 4]),|
  !gapprompt@>| !gapinput@             Transformation([3, 4, 3, 4, 3])];;|
  !gapprompt@gap>| !gapinput@cong := RightSemigroupCongruence(S, [pair1, pair2]);|
  <right semigroup congruence over <transformation semigroup of
   degree 5 with 4 generators> with 2 generating pairs>
  !gapprompt@gap>| !gapinput@x := Transformation([3, 4, 3, 4, 3]);;|
  !gapprompt@gap>| !gapinput@class := EquivalenceClassOfElement(cong, x);|
  <right congruence class of Transformation( [ 3, 4, 3, 4, 3 ] )>
  !gapprompt@gap>| !gapinput@elm := Transformation([1, 2, 2, 1, 2]);;|
  !gapprompt@gap>| !gapinput@OnRightCongruenceClasses(class, elm);|
  <right congruence class of Transformation( [ 2, 1, 2, 1, 2 ] )>
\end{Verbatim}
 }

 }

 
\section{\textcolor{Chapter }{ Finding the congruences of a semigroup }}\logpage{[ 13, 4, 0 ]}
\hyperdef{L}{X806DEBC07E6D8FCA}{}
{
  

\subsection{\textcolor{Chapter }{CongruencesOfSemigroup (for a semigroup)}}
\logpage{[ 13, 4, 1 ]}\nobreak
\hyperdef{L}{X7E8D5BA27CB5A4DA}{}
{\noindent\textcolor{FuncColor}{$\triangleright$\enspace\texttt{CongruencesOfSemigroup({\mdseries\slshape S})\index{CongruencesOfSemigroup@\texttt{CongruencesOfSemigroup}!for a semigroup}
\label{CongruencesOfSemigroup:for a semigroup}
}\hfill{\scriptsize (attribute)}}\\
\noindent\textcolor{FuncColor}{$\triangleright$\enspace\texttt{LeftCongruencesOfSemigroup({\mdseries\slshape S})\index{LeftCongruencesOfSemigroup@\texttt{LeftCongruencesOfSemigroup}!for a semigroup}
\label{LeftCongruencesOfSemigroup:for a semigroup}
}\hfill{\scriptsize (attribute)}}\\
\noindent\textcolor{FuncColor}{$\triangleright$\enspace\texttt{RightCongruencesOfSemigroup({\mdseries\slshape S})\index{RightCongruencesOfSemigroup@\texttt{RightCongruencesOfSemigroup}!for a semigroup}
\label{RightCongruencesOfSemigroup:for a semigroup}
}\hfill{\scriptsize (attribute)}}\\
\noindent\textcolor{FuncColor}{$\triangleright$\enspace\texttt{CongruencesOfSemigroup({\mdseries\slshape S, restriction})\index{CongruencesOfSemigroup@\texttt{CongruencesOfSemigroup}!for a semigroup and a multiplicative element collection}
\label{CongruencesOfSemigroup:for a semigroup and a multiplicative element collection}
}\hfill{\scriptsize (operation)}}\\
\noindent\textcolor{FuncColor}{$\triangleright$\enspace\texttt{LeftCongruencesOfSemigroup({\mdseries\slshape S, restriction})\index{LeftCongruencesOfSemigroup@\texttt{LeftCongruencesOfSemigroup}!for a semigroup and a multiplicative element collection}
\label{LeftCongruencesOfSemigroup:for a semigroup and a multiplicative element collection}
}\hfill{\scriptsize (operation)}}\\
\noindent\textcolor{FuncColor}{$\triangleright$\enspace\texttt{RightCongruencesOfSemigroup({\mdseries\slshape S, restriction})\index{RightCongruencesOfSemigroup@\texttt{RightCongruencesOfSemigroup}!for a semigroup and a multiplicative element collection}
\label{RightCongruencesOfSemigroup:for a semigroup and a multiplicative element collection}
}\hfill{\scriptsize (operation)}}\\
\textbf{\indent Returns:\ }
The congruences of a semigroup.



 This attribute gives a list of the left, right, or 2\texttt{\symbol{45}}sided
congruences of the semigroup \mbox{\texttt{\mdseries\slshape S}}. 

 If \mbox{\texttt{\mdseries\slshape restriction}} is specified and is a collection of elements from \mbox{\texttt{\mdseries\slshape S}}, then the result will only include congruences generated by pairs of elements
from \mbox{\texttt{\mdseries\slshape restriction}}. Otherwise, all congruences will be calculated.

 See also \texttt{LatticeOfCongruences} (\ref{LatticeOfCongruences:for a semigroup}). 

 
\begin{Verbatim}[commandchars=!@|,fontsize=\small,frame=single,label=Example]
  !gapprompt@gap>| !gapinput@S := ReesZeroMatrixSemigroup(SymmetricGroup(3),|
  !gapprompt@>| !gapinput@                                [[(), (1, 3, 2)], [(1, 2), 0]]);;|
  !gapprompt@gap>| !gapinput@congs := CongruencesOfSemigroup(S);;|
  !gapprompt@gap>| !gapinput@Length(congs);|
  4
  !gapprompt@gap>| !gapinput@Set(congs, NrEquivalenceClasses);|
  [ 1, 5, 9, 25 ]
  !gapprompt@gap>| !gapinput@pos := Position(congs, UniversalSemigroupCongruence(S));;|
  !gapprompt@gap>| !gapinput@congs[pos];|
  <universal semigroup congruence over
  <Rees 0-matrix semigroup 2x2 over Sym( [ 1 .. 3 ] )>>
\end{Verbatim}
 }

 

\subsection{\textcolor{Chapter }{MinimalCongruencesOfSemigroup (for a semigroup)}}
\logpage{[ 13, 4, 2 ]}\nobreak
\hyperdef{L}{X7838738987B2DB41}{}
{\noindent\textcolor{FuncColor}{$\triangleright$\enspace\texttt{MinimalCongruencesOfSemigroup({\mdseries\slshape S})\index{MinimalCongruencesOfSemigroup@\texttt{MinimalCongruencesOfSemigroup}!for a semigroup}
\label{MinimalCongruencesOfSemigroup:for a semigroup}
}\hfill{\scriptsize (attribute)}}\\
\noindent\textcolor{FuncColor}{$\triangleright$\enspace\texttt{MinimalLeftCongruencesOfSemigroup({\mdseries\slshape S})\index{MinimalLeftCongruencesOfSemigroup@\texttt{MinimalLeftCongruencesOfSemigroup}!for a semigroup}
\label{MinimalLeftCongruencesOfSemigroup:for a semigroup}
}\hfill{\scriptsize (attribute)}}\\
\noindent\textcolor{FuncColor}{$\triangleright$\enspace\texttt{MinimalRightCongruencesOfSemigroup({\mdseries\slshape S})\index{MinimalRightCongruencesOfSemigroup@\texttt{MinimalRightCongruencesOfSemigroup}!for a semigroup}
\label{MinimalRightCongruencesOfSemigroup:for a semigroup}
}\hfill{\scriptsize (attribute)}}\\
\noindent\textcolor{FuncColor}{$\triangleright$\enspace\texttt{MinimalCongruencesOfSemigroup({\mdseries\slshape S, restriction})\index{MinimalCongruencesOfSemigroup@\texttt{MinimalCongruencesOfSemigroup}!for a semigroup and a multiplicative element collection}
\label{MinimalCongruencesOfSemigroup:for a semigroup and a multiplicative element collection}
}\hfill{\scriptsize (operation)}}\\
\noindent\textcolor{FuncColor}{$\triangleright$\enspace\texttt{MinimalLeftCongruencesOfSemigroup({\mdseries\slshape S, restriction})\index{MinimalLeftCongruencesOfSemigroup@\texttt{MinimalLeftCongruencesOfSemigroup}!for a semigroup and a multiplicative element collection}
\label{MinimalLeftCongruencesOfSemigroup:for a semigroup and a multiplicative element collection}
}\hfill{\scriptsize (operation)}}\\
\noindent\textcolor{FuncColor}{$\triangleright$\enspace\texttt{MinimalRightCongruencesOfSemigroup({\mdseries\slshape S, restriction})\index{MinimalRightCongruencesOfSemigroup@\texttt{MinimalRightCongruencesOfSemigroup}!for a semigroup and a multiplicative element collection}
\label{MinimalRightCongruencesOfSemigroup:for a semigroup and a multiplicative element collection}
}\hfill{\scriptsize (operation)}}\\
\textbf{\indent Returns:\ }
The congruences of a semigroup.



 If \mbox{\texttt{\mdseries\slshape S}} is a semigroup, then the attribute \texttt{MinimalCongruencesOfSemigroup} gives a list of all the congruences on \mbox{\texttt{\mdseries\slshape S}} which are \emph{minimal}. A congruence is minimal iff it is non\texttt{\symbol{45}}trivial and
contains no other congruences as subrelations (apart from the trivial
congruence). 

 \texttt{MinimalLeftCongruencesOfSemigroup} and \texttt{MinimalRightCongruencesOfSemigroup} do the same thing, but for left congruences and right congruences
respectively. Note that any congruence is also a left congruence, but that a
minimal congruence may not be a minimal left congruence. 

 If \mbox{\texttt{\mdseries\slshape restriction}} is specified and is a collection of elements from \mbox{\texttt{\mdseries\slshape S}}, then the result will only include congruences generated by pairs of elements
from \mbox{\texttt{\mdseries\slshape restriction}}. Otherwise, all congruences will be calculated.

 See also \texttt{CongruencesOfSemigroup} (\ref{CongruencesOfSemigroup:for a semigroup}) and \texttt{PrincipalCongruencesOfSemigroup} (\ref{PrincipalCongruencesOfSemigroup:for a semigroup}). 

 
\begin{Verbatim}[commandchars=!@|,fontsize=\small,frame=single,label=Example]
  !gapprompt@gap>| !gapinput@S := Semigroup(Transformation([1, 3, 2]),|
  !gapprompt@>| !gapinput@                  Transformation([3, 1, 3]));;|
  !gapprompt@gap>| !gapinput@min := MinimalCongruencesOfSemigroup(S);|
  [ <2-sided semigroup congruence over <transformation semigroup
       of size 13, degree 3 with 2 generators> with 1 generating pairs>
   ]
  !gapprompt@gap>| !gapinput@minl := MinimalLeftCongruencesOfSemigroup(S);|
  [ <left semigroup congruence over <transformation semigroup
       of size 13, degree 3 with 2 generators> with 1 generating pairs>,
    <left semigroup congruence over <transformation semigroup
       of size 13, degree 3 with 2 generators> with 1 generating pairs>,
    <left semigroup congruence over <transformation semigroup
       of size 13, degree 3 with 2 generators> with 1 generating pairs>
   ]
\end{Verbatim}
 }

 

\subsection{\textcolor{Chapter }{PrincipalCongruencesOfSemigroup (for a semigroup)}}
\logpage{[ 13, 4, 3 ]}\nobreak
\hyperdef{L}{X7986F3597F9DF7AF}{}
{\noindent\textcolor{FuncColor}{$\triangleright$\enspace\texttt{PrincipalCongruencesOfSemigroup({\mdseries\slshape S})\index{PrincipalCongruencesOfSemigroup@\texttt{PrincipalCongruencesOfSemigroup}!for a semigroup}
\label{PrincipalCongruencesOfSemigroup:for a semigroup}
}\hfill{\scriptsize (attribute)}}\\
\noindent\textcolor{FuncColor}{$\triangleright$\enspace\texttt{PrincipalLeftCongruencesOfSemigroup({\mdseries\slshape S})\index{PrincipalLeftCongruencesOfSemigroup@\texttt{PrincipalLeftCongruencesOfSemigroup}!for a semigroup}
\label{PrincipalLeftCongruencesOfSemigroup:for a semigroup}
}\hfill{\scriptsize (attribute)}}\\
\noindent\textcolor{FuncColor}{$\triangleright$\enspace\texttt{PrincipalRightCongruencesOfSemigroup({\mdseries\slshape S})\index{PrincipalRightCongruencesOfSemigroup@\texttt{Principal}\-\texttt{Right}\-\texttt{Congruences}\-\texttt{Of}\-\texttt{Semigroup}!for a semigroup}
\label{PrincipalRightCongruencesOfSemigroup:for a semigroup}
}\hfill{\scriptsize (attribute)}}\\
\noindent\textcolor{FuncColor}{$\triangleright$\enspace\texttt{PrincipalCongruencesOfSemigroup({\mdseries\slshape S, restriction})\index{PrincipalCongruencesOfSemigroup@\texttt{PrincipalCongruencesOfSemigroup}!for a semigroup and a multiplicative element collection}
\label{PrincipalCongruencesOfSemigroup:for a semigroup and a multiplicative element collection}
}\hfill{\scriptsize (operation)}}\\
\noindent\textcolor{FuncColor}{$\triangleright$\enspace\texttt{PrincipalLeftCongruencesOfSemigroup({\mdseries\slshape S, restriction})\index{PrincipalLeftCongruencesOfSemigroup@\texttt{PrincipalLeftCongruencesOfSemigroup}!for a semigroup and a multiplicative element collection}
\label{PrincipalLeftCongruencesOfSemigroup:for a semigroup and a multiplicative element collection}
}\hfill{\scriptsize (operation)}}\\
\noindent\textcolor{FuncColor}{$\triangleright$\enspace\texttt{PrincipalRightCongruencesOfSemigroup({\mdseries\slshape S, restriction})\index{PrincipalRightCongruencesOfSemigroup@\texttt{Principal}\-\texttt{Right}\-\texttt{Congruences}\-\texttt{Of}\-\texttt{Semigroup}!for a semigroup and a multiplicative element collection}
\label{PrincipalRightCongruencesOfSemigroup:for a semigroup and a multiplicative element collection}
}\hfill{\scriptsize (operation)}}\\
\textbf{\indent Returns:\ }
A list.



 If \mbox{\texttt{\mdseries\slshape S}} is a semigroup, then the attribute \texttt{PrincipalCongruencesOfSemigroup} gives a list of all the congruences on \mbox{\texttt{\mdseries\slshape S}} which are \emph{principal}. A congruence is principal if and only if it is
non\texttt{\symbol{45}}trivial and can be defined by a single generating pair. 

 \texttt{PrincipalLeftCongruencesOfSemigroup} and \texttt{PrincipalRightCongruencesOfSemigroup} do the same thing, but for left congruences and right congruences
respectively. Note that any congruence is a left congruence and a right
congruence, but that a principal congruence may not be a principal left
congruence or a principal right congruence. 

 If \mbox{\texttt{\mdseries\slshape restriction}} is specified and is a collection of elements from \mbox{\texttt{\mdseries\slshape S}}, then the result will only include congruences generated by pairs of elements
from \mbox{\texttt{\mdseries\slshape restriction}}. Otherwise, all congruences will be calculated.

 See also \texttt{CongruencesOfSemigroup} (\ref{CongruencesOfSemigroup:for a semigroup}) and \texttt{MinimalCongruencesOfSemigroup} (\ref{MinimalCongruencesOfSemigroup:for a semigroup}). 

 
\begin{Verbatim}[commandchars=!@|,fontsize=\small,frame=single,label=Example]
  !gapprompt@gap>| !gapinput@S := Semigroup(Transformation([1, 3, 2]),|
  !gapprompt@>| !gapinput@                  Transformation([3, 1, 3]));;|
  !gapprompt@gap>| !gapinput@congs := PrincipalCongruencesOfSemigroup(S);|
  [ <universal semigroup congruence over <transformation semigroup
       of size 13, degree 3 with 2 generators>>,
    <2-sided semigroup congruence over <transformation semigroup
       of size 13, degree 3 with 2 generators> with 1 generating pairs>,
    <2-sided semigroup congruence over <transformation semigroup
       of size 13, degree 3 with 2 generators> with 1 generating pairs>,
    <2-sided semigroup congruence over <transformation semigroup
       of size 13, degree 3 with 2 generators> with 1 generating pairs>,
    <2-sided semigroup congruence over <transformation semigroup
       of size 13, degree 3 with 2 generators> with 1 generating pairs>
   ]
\end{Verbatim}
 }

 

\subsection{\textcolor{Chapter }{IsCongruencePoset}}
\logpage{[ 13, 4, 4 ]}\nobreak
\hyperdef{L}{X8195D6F47EE52806}{}
{\noindent\textcolor{FuncColor}{$\triangleright$\enspace\texttt{IsCongruencePoset({\mdseries\slshape poset})\index{IsCongruencePoset@\texttt{IsCongruencePoset}}
\label{IsCongruencePoset}
}\hfill{\scriptsize (Category)}}\\
\noindent\textcolor{FuncColor}{$\triangleright$\enspace\texttt{IsCayleyDigraphOfCongruences({\mdseries\slshape poset})\index{IsCayleyDigraphOfCongruences@\texttt{IsCayleyDigraphOfCongruences}}
\label{IsCayleyDigraphOfCongruences}
}\hfill{\scriptsize (Category)}}\\
\textbf{\indent Returns:\ }
\texttt{true} or \texttt{false}.



 This category contains all congruence posets. A \emph{congruence poset} is a partially ordered set of congruences over a specific semigroup, where the
ordering is defined by containment according to \texttt{IsSubrelation} (\ref{IsSubrelation}): given two congruences \texttt{cong1} and \texttt{cong2}, we say that \texttt{cong1} {\textless} \texttt{cong2} if and only if \texttt{cong1} is a subrelation (a refinement) of \texttt{cong2}. The congruences in a congruence poset can be left, right, or
two\texttt{\symbol{45}}sided. 

 A congruence poset is a digraph (see \texttt{IsDigraph} (\textbf{Digraphs: IsDigraph})) with a vertex for each congruence, and an edge from vertex \texttt{i} to vertex \texttt{j} only if the congruence numbered \texttt{i} is a subrelation of the congruence numbered \texttt{j}. To avoid using an unnecessarily large amount of memory in some cases, a
congruence poset does not necessarily belong to \texttt{IsPartialOrderDigraph} (\textbf{Digraphs: IsPartialOrderDigraph}). In other words, although every congruence poset represents a partial order
it is not necessarily the case that there is an edge from vertex \texttt{i} to vertex \texttt{j} if and only if the congruence numbered \texttt{i} is a subrelation of the congruence numbered \texttt{j}. 

 The list of congruences can be obtained using \texttt{CongruencesOfPoset} (\ref{CongruencesOfPoset}); and the underlying semigroup of the poset can be obtained using \texttt{UnderlyingSemigroupOfCongruencePoset} (\ref{UnderlyingSemigroupOfCongruencePoset}).

 Congruence posets can be created using any of: 
\begin{itemize}
\item  \texttt{PosetOfCongruences} (\ref{PosetOfCongruences}), 
\item  \texttt{JoinSemilatticeOfCongruences} (\ref{JoinSemilatticeOfCongruences}) 
\item  \texttt{LatticeOfCongruences} (\ref{LatticeOfCongruences:for a semigroup}), \texttt{LatticeOfLeftCongruences} (\ref{LatticeOfLeftCongruences:for a semigroup}), or \texttt{LatticeOfRightCongruences} (\ref{LatticeOfRightCongruences:for a semigroup}) 
\item  \texttt{CayleyDigraphOfCongruences} (\ref{CayleyDigraphOfCongruences:for a semigroup}), \texttt{CayleyDigraphOfLeftCongruences} (\ref{CayleyDigraphOfLeftCongruences:for a semigroup}), or \texttt{CayleyDigraphOfRightCongruences} (\ref{CayleyDigraphOfRightCongruences:for a semigroup}). 
\end{itemize}
 \texttt{IsCayleyDigraphOfCongruences} only applies to the output of \texttt{JoinSemilatticeOfCongruences} (\ref{JoinSemilatticeOfCongruences}), \texttt{CayleyDigraphOfCongruences} (\ref{CayleyDigraphOfCongruences:for a semigroup}), \texttt{CayleyDigraphOfLeftCongruences} (\ref{CayleyDigraphOfLeftCongruences:for a semigroup}), and \texttt{CayleyDigraphOfRightCongruences} (\ref{CayleyDigraphOfRightCongruences:for a semigroup}). The congruences used as the generating set for these operations can be
obtained using \texttt{GeneratingCongruencesOfJoinSemilattice} (\ref{GeneratingCongruencesOfJoinSemilattice}). 
\begin{Verbatim}[commandchars=!@|,fontsize=\small,frame=single,label=Example]
  !gapprompt@gap>| !gapinput@S := SymmetricInverseMonoid(2);;|
  !gapprompt@gap>| !gapinput@poset := LatticeOfCongruences(S);|
  <lattice of 4 two-sided congruences over
   <symmetric inverse monoid of degree 2>>
  !gapprompt@gap>| !gapinput@IsCongruencePoset(poset);|
  true
  !gapprompt@gap>| !gapinput@IsDigraph(poset);|
  true
  !gapprompt@gap>| !gapinput@IsIsomorphicDigraph(poset,|
  !gapprompt@>| !gapinput@Digraph([[1, 2, 3, 4], [2], [2, 3], [2, 3, 4]]));|
  true
  !gapprompt@gap>| !gapinput@T := FullTransformationMonoid(3);;|
  !gapprompt@gap>| !gapinput@congs := PrincipalCongruencesOfSemigroup(T);;|
  !gapprompt@gap>| !gapinput@poset := JoinSemilatticeOfCongruences(PosetOfCongruences(congs));|
  <lattice of 6 two-sided congruences over
   <full transformation monoid of degree 3>>
  !gapprompt@gap>| !gapinput@IsCayleyDigraphOfCongruences(poset);|
  false
  !gapprompt@gap>| !gapinput@IsCongruencePoset(poset);|
  true
  !gapprompt@gap>| !gapinput@DigraphNrVertices(poset);|
  6
  !gapprompt@gap>| !gapinput@poset := CayleyDigraphOfCongruences(T);|
  <poset of 7 two-sided congruences over
   <full transformation monoid of degree 3>>
  !gapprompt@gap>| !gapinput@IsCayleyDigraphOfCongruences(poset);|
  true
\end{Verbatim}
 }

 

\subsection{\textcolor{Chapter }{LatticeOfCongruences (for a semigroup)}}
\logpage{[ 13, 4, 5 ]}\nobreak
\hyperdef{L}{X86C9C5BA7FE93F4C}{}
{\noindent\textcolor{FuncColor}{$\triangleright$\enspace\texttt{LatticeOfCongruences({\mdseries\slshape S})\index{LatticeOfCongruences@\texttt{LatticeOfCongruences}!for a semigroup}
\label{LatticeOfCongruences:for a semigroup}
}\hfill{\scriptsize (attribute)}}\\
\noindent\textcolor{FuncColor}{$\triangleright$\enspace\texttt{LatticeOfLeftCongruences({\mdseries\slshape S})\index{LatticeOfLeftCongruences@\texttt{LatticeOfLeftCongruences}!for a semigroup}
\label{LatticeOfLeftCongruences:for a semigroup}
}\hfill{\scriptsize (attribute)}}\\
\noindent\textcolor{FuncColor}{$\triangleright$\enspace\texttt{LatticeOfRightCongruences({\mdseries\slshape S})\index{LatticeOfRightCongruences@\texttt{LatticeOfRightCongruences}!for a semigroup}
\label{LatticeOfRightCongruences:for a semigroup}
}\hfill{\scriptsize (attribute)}}\\
\noindent\textcolor{FuncColor}{$\triangleright$\enspace\texttt{LatticeOfCongruences({\mdseries\slshape S, restriction})\index{LatticeOfCongruences@\texttt{LatticeOfCongruences}!for a semigroup and a list or collection}
\label{LatticeOfCongruences:for a semigroup and a list or collection}
}\hfill{\scriptsize (operation)}}\\
\noindent\textcolor{FuncColor}{$\triangleright$\enspace\texttt{LatticeOfLeftCongruences({\mdseries\slshape S, restriction})\index{LatticeOfLeftCongruences@\texttt{LatticeOfLeftCongruences}!for a semigroup and a list or collection}
\label{LatticeOfLeftCongruences:for a semigroup and a list or collection}
}\hfill{\scriptsize (operation)}}\\
\noindent\textcolor{FuncColor}{$\triangleright$\enspace\texttt{LatticeOfRightCongruences({\mdseries\slshape S, restriction})\index{LatticeOfRightCongruences@\texttt{LatticeOfRightCongruences}!for a semigroup and a list or collection}
\label{LatticeOfRightCongruences:for a semigroup and a list or collection}
}\hfill{\scriptsize (operation)}}\\
\textbf{\indent Returns:\ }
A lattice digraph.



 If \mbox{\texttt{\mdseries\slshape S}} is a semigroup, then \texttt{LatticeOfCongruences} returns a congruence poset object containing all the congruences of \mbox{\texttt{\mdseries\slshape S}} and information about how they are contained in each other. See \texttt{IsCongruencePoset} (\ref{IsCongruencePoset}) for more details. 

 \texttt{LatticeOfLeftCongruences} and \texttt{LatticeOfRightCongruences} do the same thing for left and right congruences, respectively. 

 If \mbox{\texttt{\mdseries\slshape restriction}} is specified and is a collection of elements from \mbox{\texttt{\mdseries\slshape S}}, then the result will only include congruences generated by pairs of elements
from \mbox{\texttt{\mdseries\slshape restriction}}. Otherwise, all congruences will be calculated.

 See \texttt{CongruencesOfSemigroup} (\ref{CongruencesOfSemigroup:for a semigroup}). 

 
\begin{Verbatim}[commandchars=!@|,fontsize=\small,frame=single,label=Example]
  !gapprompt@gap>| !gapinput@S := OrderEndomorphisms(2);;|
  !gapprompt@gap>| !gapinput@LatticeOfCongruences(S);|
  <lattice of 3 two-sided congruences over <regular transformation
   monoid of size 3, degree 2 with 2 generators>>
  !gapprompt@gap>| !gapinput@LatticeOfLeftCongruences(S);|
  <lattice of 3 left congruences over <regular transformation monoid
   of size 3, degree 2 with 2 generators>>
  !gapprompt@gap>| !gapinput@LatticeOfRightCongruences(S);|
  <lattice of 5 right congruences over <regular transformation monoid
   of size 3, degree 2 with 2 generators>>
  !gapprompt@gap>| !gapinput@IsIsomorphicDigraph(LatticeOfRightCongruences(S),|
  !gapprompt@>| !gapinput@Digraph([[1, 2, 3, 4, 5], [2], [2, 3], [2, 4], [2, 5]]));|
  true
  !gapprompt@gap>| !gapinput@S := FullTransformationMonoid(4);;|
  !gapprompt@gap>| !gapinput@restriction := [Transformation([1, 1, 1, 1]),|
  !gapprompt@>| !gapinput@                   Transformation([1, 1, 1, 2]),|
  !gapprompt@>| !gapinput@                   Transformation([1, 1, 1, 3])];;|
  !gapprompt@gap>| !gapinput@latt := LatticeOfCongruences(S, Combinations(restriction, 2));|
  <lattice of 2 two-sided congruences over
   <full transformation monoid of degree 4>>
\end{Verbatim}
 }

 

\subsection{\textcolor{Chapter }{CayleyDigraphOfCongruences (for a semigroup)}}
\logpage{[ 13, 4, 6 ]}\nobreak
\hyperdef{L}{X784CFDE37A0B4F84}{}
{\noindent\textcolor{FuncColor}{$\triangleright$\enspace\texttt{CayleyDigraphOfCongruences({\mdseries\slshape S})\index{CayleyDigraphOfCongruences@\texttt{CayleyDigraphOfCongruences}!for a semigroup}
\label{CayleyDigraphOfCongruences:for a semigroup}
}\hfill{\scriptsize (attribute)}}\\
\noindent\textcolor{FuncColor}{$\triangleright$\enspace\texttt{CayleyDigraphOfLeftCongruences({\mdseries\slshape S})\index{CayleyDigraphOfLeftCongruences@\texttt{CayleyDigraphOfLeftCongruences}!for a semigroup}
\label{CayleyDigraphOfLeftCongruences:for a semigroup}
}\hfill{\scriptsize (attribute)}}\\
\noindent\textcolor{FuncColor}{$\triangleright$\enspace\texttt{CayleyDigraphOfRightCongruences({\mdseries\slshape S})\index{CayleyDigraphOfRightCongruences@\texttt{CayleyDigraphOfRightCongruences}!for a semigroup}
\label{CayleyDigraphOfRightCongruences:for a semigroup}
}\hfill{\scriptsize (attribute)}}\\
\noindent\textcolor{FuncColor}{$\triangleright$\enspace\texttt{CayleyDigraphOfCongruences({\mdseries\slshape S, restriction})\index{CayleyDigraphOfCongruences@\texttt{CayleyDigraphOfCongruences}!for a semigroup and a list or collection}
\label{CayleyDigraphOfCongruences:for a semigroup and a list or collection}
}\hfill{\scriptsize (operation)}}\\
\noindent\textcolor{FuncColor}{$\triangleright$\enspace\texttt{CayleyDigraphOfLeftCongruences({\mdseries\slshape S, restriction})\index{CayleyDigraphOfLeftCongruences@\texttt{CayleyDigraphOfLeftCongruences}!for a semigroup and a list or collection}
\label{CayleyDigraphOfLeftCongruences:for a semigroup and a list or collection}
}\hfill{\scriptsize (operation)}}\\
\noindent\textcolor{FuncColor}{$\triangleright$\enspace\texttt{CayleyDigraphOfRightCongruences({\mdseries\slshape S, restriction})\index{CayleyDigraphOfRightCongruences@\texttt{CayleyDigraphOfRightCongruences}!for a semigroup and a list or collection}
\label{CayleyDigraphOfRightCongruences:for a semigroup and a list or collection}
}\hfill{\scriptsize (operation)}}\\
\textbf{\indent Returns:\ }
A digraph.



 If \mbox{\texttt{\mdseries\slshape S}} is a semigroup, then \texttt{CayleyDigraphOfCongruences} returns the right Cayley graph of the semilattice of congruences of \mbox{\texttt{\mdseries\slshape S}} with respect to the generating set consisting of the principal congruences
congruence poset. See \texttt{IsCayleyDigraphOfCongruences} (\ref{IsCayleyDigraphOfCongruences}) for more details. 

 \texttt{CayleyDigraphOfLeftCongruences} and \texttt{CayleyDigraphOfRightCongruences} do the same thing for left and right congruences, respectively. 

 If \mbox{\texttt{\mdseries\slshape restriction}} is specified and is a collection of elements from \mbox{\texttt{\mdseries\slshape S}}, then the result will only include congruences generated by pairs of elements
from \mbox{\texttt{\mdseries\slshape restriction}}. Otherwise, all congruences will be calculated.

 Note that \texttt{LatticeOfCongruences} (\ref{LatticeOfCongruences:for a semigroup}), and its analogues for right and left congruences, return the reflexive
transitive closure of the digraph returned by this function (with any multiple
edges removed). If there are a large number of congruences, then it might be
the case that forming the reflexive transitive closure takes a significant
amount of time, and so it might be desirable to use this function instead. 

 See \texttt{CongruencesOfSemigroup} (\ref{CongruencesOfSemigroup:for a semigroup}). 

 
\begin{Verbatim}[commandchars=!@|,fontsize=\small,frame=single,label=Example]
  !gapprompt@gap>| !gapinput@S := OrderEndomorphisms(2);;|
  !gapprompt@gap>| !gapinput@CayleyDigraphOfCongruences(S);|
  <poset of 3 two-sided congruences over <regular transformation monoid
   of size 3, degree 2 with 2 generators>>
  !gapprompt@gap>| !gapinput@CayleyDigraphOfLeftCongruences(S);|
  <poset of 3 left congruences over <regular transformation monoid
   of size 3, degree 2 with 2 generators>>
  !gapprompt@gap>| !gapinput@CayleyDigraphOfRightCongruences(S);|
  <poset of 5 right congruences over <regular transformation monoid
   of size 3, degree 2 with 2 generators>>
  !gapprompt@gap>| !gapinput@IsIsomorphicDigraph(CayleyDigraphOfRightCongruences(S),|
  !gapprompt@>| !gapinput@Digraph([[2, 3, 4], [2, 5, 5], [5, 3, 5], [5, 5, 4], [5, 5, 5]]));|
  true
  !gapprompt@gap>| !gapinput@S := FullTransformationMonoid(4);;|
  !gapprompt@gap>| !gapinput@restriction := [Transformation([1, 1, 1, 1]),|
  !gapprompt@>| !gapinput@                   Transformation([1, 1, 1, 2]),|
  !gapprompt@>| !gapinput@                   Transformation([1, 1, 1, 3])];;|
  !gapprompt@gap>| !gapinput@CayleyDigraphOfCongruences(S, Combinations(restriction, 2));|
  <poset of 2 two-sided congruences over
   <full transformation monoid of degree 4>>
\end{Verbatim}
 }

 

\subsection{\textcolor{Chapter }{PosetOfPrincipalCongruences (for a semigroup)}}
\logpage{[ 13, 4, 7 ]}\nobreak
\hyperdef{L}{X83ACF2C0789F621B}{}
{\noindent\textcolor{FuncColor}{$\triangleright$\enspace\texttt{PosetOfPrincipalCongruences({\mdseries\slshape S})\index{PosetOfPrincipalCongruences@\texttt{PosetOfPrincipalCongruences}!for a semigroup}
\label{PosetOfPrincipalCongruences:for a semigroup}
}\hfill{\scriptsize (attribute)}}\\
\noindent\textcolor{FuncColor}{$\triangleright$\enspace\texttt{PosetOfPrincipalLeftCongruences({\mdseries\slshape S})\index{PosetOfPrincipalLeftCongruences@\texttt{PosetOfPrincipalLeftCongruences}!for a semigroup}
\label{PosetOfPrincipalLeftCongruences:for a semigroup}
}\hfill{\scriptsize (attribute)}}\\
\noindent\textcolor{FuncColor}{$\triangleright$\enspace\texttt{PosetOfPrincipalRightCongruences({\mdseries\slshape S})\index{PosetOfPrincipalRightCongruences@\texttt{PosetOfPrincipalRightCongruences}!for a semigroup}
\label{PosetOfPrincipalRightCongruences:for a semigroup}
}\hfill{\scriptsize (attribute)}}\\
\noindent\textcolor{FuncColor}{$\triangleright$\enspace\texttt{PosetOfPrincipalCongruences({\mdseries\slshape S, restriction})\index{PosetOfPrincipalCongruences@\texttt{PosetOfPrincipalCongruences}!for a semigroup and a multiplicative element collection}
\label{PosetOfPrincipalCongruences:for a semigroup and a multiplicative element collection}
}\hfill{\scriptsize (operation)}}\\
\noindent\textcolor{FuncColor}{$\triangleright$\enspace\texttt{PosetOfPrincipalLeftCongruences({\mdseries\slshape S, restriction})\index{PosetOfPrincipalLeftCongruences@\texttt{PosetOfPrincipalLeftCongruences}!for a semigroup and a multiplicative element collection}
\label{PosetOfPrincipalLeftCongruences:for a semigroup and a multiplicative element collection}
}\hfill{\scriptsize (operation)}}\\
\noindent\textcolor{FuncColor}{$\triangleright$\enspace\texttt{PosetOfPrincipalRightCongruences({\mdseries\slshape S, restriction})\index{PosetOfPrincipalRightCongruences@\texttt{PosetOfPrincipalRightCongruences}!for a semigroup and a multiplicative element collection}
\label{PosetOfPrincipalRightCongruences:for a semigroup and a multiplicative element collection}
}\hfill{\scriptsize (operation)}}\\
\textbf{\indent Returns:\ }
A congruence poset.



 If \mbox{\texttt{\mdseries\slshape S}} is a semigroup, then \texttt{PosetOfPrincipalCongruences} returns a congruence poset object which contains all the principal congruences
of \mbox{\texttt{\mdseries\slshape S}}, ordered by containment according to \texttt{IsSubrelation} (\ref{IsSubrelation}). A congruence is \emph{principal} if it can be defined by a single generating pair. \texttt{PosetOfPrincipalLeftCongruences} and \texttt{PosetOfPrincipalRightCongruences} do the same thing for left and right congruences respectively. 

 If \mbox{\texttt{\mdseries\slshape restriction}} is specified and is a collection of elements from \mbox{\texttt{\mdseries\slshape S}}, then the result will only include principal congruences generated by pairs
of elements from \mbox{\texttt{\mdseries\slshape restriction}}. Otherwise, all principal congruences will be calculated.

 See also \texttt{LatticeOfCongruences} (\ref{LatticeOfCongruences:for a semigroup}) and \texttt{PrincipalCongruencesOfSemigroup} (\ref{PrincipalCongruencesOfSemigroup:for a semigroup}). 

 
\begin{Verbatim}[commandchars=!@|,fontsize=\small,frame=single,label=Example]
  !gapprompt@gap>| !gapinput@S := Semigroup(Transformation([1, 3, 1]),|
  !gapprompt@>| !gapinput@                  Transformation([2, 3, 3]));;|
  !gapprompt@gap>| !gapinput@PosetOfPrincipalLeftCongruences(S);|
  <poset of 12 left congruences over <transformation semigroup
   of size 11, degree 3 with 2 generators>>
  !gapprompt@gap>| !gapinput@PosetOfPrincipalCongruences(S);|
  <lattice of 3 two-sided congruences over <transformation semigroup
   of size 11, degree 3 with 2 generators>>
  !gapprompt@gap>| !gapinput@restriction := [Transformation([3, 2, 3]),|
  !gapprompt@>| !gapinput@                   Transformation([3, 1, 3]),|
  !gapprompt@>| !gapinput@                   Transformation([2, 2, 2])];;|
  !gapprompt@gap>| !gapinput@poset := PosetOfPrincipalRightCongruences(S,|
  !gapprompt@>| !gapinput@Combinations(restriction, 2));|
  <poset of 3 right congruences over <transformation semigroup
   of size 11, degree 3 with 2 generators>>
\end{Verbatim}
 }

 

\subsection{\textcolor{Chapter }{CongruencesOfPoset}}
\logpage{[ 13, 4, 8 ]}\nobreak
\hyperdef{L}{X7B2E2CEE8626FBC3}{}
{\noindent\textcolor{FuncColor}{$\triangleright$\enspace\texttt{CongruencesOfPoset({\mdseries\slshape poset})\index{CongruencesOfPoset@\texttt{CongruencesOfPoset}}
\label{CongruencesOfPoset}
}\hfill{\scriptsize (attribute)}}\\
\textbf{\indent Returns:\ }
A list.



 If \mbox{\texttt{\mdseries\slshape poset}} is a congruence poset object, then this attribute returns a list of all the
congruence objects in the poset (these may be left, right, or
two\texttt{\symbol{45}}sided). The order of this list corresponds to the order
of the entries in the poset. 

 See also \texttt{LatticeOfCongruences} (\ref{LatticeOfCongruences:for a semigroup}) and \texttt{CongruencesOfSemigroup} (\ref{CongruencesOfSemigroup:for a semigroup}). 

 
\begin{Verbatim}[commandchars=!@|,fontsize=\small,frame=single,label=Example]
  !gapprompt@gap>| !gapinput@S := OrderEndomorphisms(2);;|
  !gapprompt@gap>| !gapinput@latt := LatticeOfRightCongruences(S);|
  <lattice of 5 right congruences over <regular transformation monoid
   of size 3, degree 2 with 2 generators>>
  !gapprompt@gap>| !gapinput@CongruencesOfPoset(latt);|
  [ <2-sided semigroup congruence over <regular transformation monoid
       of size 3, degree 2 with 2 generators> with 0 generating pairs>,
    <right semigroup congruence over <regular transformation monoid
       of size 3, degree 2 with 2 generators> with 1 generating pairs>,
    <right semigroup congruence over <regular transformation monoid
       of size 3, degree 2 with 2 generators> with 1 generating pairs>,
    <right semigroup congruence over <regular transformation monoid
       of size 3, degree 2 with 2 generators> with 1 generating pairs>,
    <right semigroup congruence over <regular transformation monoid
       of size 3, degree 2 with 2 generators> with 2 generating pairs> ]
\end{Verbatim}
 }

 

\subsection{\textcolor{Chapter }{UnderlyingSemigroupOfCongruencePoset}}
\logpage{[ 13, 4, 9 ]}\nobreak
\hyperdef{L}{X830F42B582A2FAA0}{}
{\noindent\textcolor{FuncColor}{$\triangleright$\enspace\texttt{UnderlyingSemigroupOfCongruencePoset({\mdseries\slshape poset})\index{UnderlyingSemigroupOfCongruencePoset@\texttt{Underlying}\-\texttt{Semigroup}\-\texttt{Of}\-\texttt{Congruence}\-\texttt{Poset}}
\label{UnderlyingSemigroupOfCongruencePoset}
}\hfill{\scriptsize (attribute)}}\\
\textbf{\indent Returns:\ }
A semigroup.



 If \mbox{\texttt{\mdseries\slshape poset}} is a congruence poset object, then this attribute returns the semigroup on
which all its congruences are defined. 
\begin{Verbatim}[commandchars=!@|,fontsize=\small,frame=single,label=Example]
  !gapprompt@gap>| !gapinput@S := OrderEndomorphisms(2);|
  <regular transformation monoid of degree 2 with 2 generators>
  !gapprompt@gap>| !gapinput@latt := LatticeOfRightCongruences(S);|
  <lattice of 5 right congruences over <regular transformation monoid
   of size 3, degree 2 with 2 generators>>
  !gapprompt@gap>| !gapinput@UnderlyingSemigroupOfCongruencePoset(latt) = S;|
  true
\end{Verbatim}
 }

 

\subsection{\textcolor{Chapter }{PosetOfCongruences}}
\logpage{[ 13, 4, 10 ]}\nobreak
\hyperdef{L}{X78A91138818A4FAE}{}
{\noindent\textcolor{FuncColor}{$\triangleright$\enspace\texttt{PosetOfCongruences({\mdseries\slshape coll})\index{PosetOfCongruences@\texttt{PosetOfCongruences}}
\label{PosetOfCongruences}
}\hfill{\scriptsize (operation)}}\\
\textbf{\indent Returns:\ }
A congruence poset.



 If \mbox{\texttt{\mdseries\slshape coll}} is a list or collection of semigroup congruences (which may be left, right, or
two\texttt{\symbol{45}}sided) then this operation returns the congruence poset
formed by these congruences partially ordered by containment. 

 This operation does not create any new congruences or take any joins. See also \texttt{JoinSemilatticeOfCongruences} (\ref{JoinSemilatticeOfCongruences}), \texttt{IsCongruencePoset} (\ref{IsCongruencePoset}), and \texttt{LatticeOfCongruences} (\ref{LatticeOfCongruences:for a semigroup}). 

 
\begin{Verbatim}[commandchars=!@|,fontsize=\small,frame=single,label=Example]
  !gapprompt@gap>| !gapinput@S := OrderEndomorphisms(2);;|
  !gapprompt@gap>| !gapinput@pair1 := [Transformation([1, 1]), IdentityTransformation];;|
  !gapprompt@gap>| !gapinput@pair2 := [IdentityTransformation, Transformation([2, 2])];;|
  !gapprompt@gap>| !gapinput@coll := [RightSemigroupCongruence(S, pair1),|
  !gapprompt@>| !gapinput@            RightSemigroupCongruence(S, pair2),|
  !gapprompt@>| !gapinput@            RightSemigroupCongruence(S, [])];;|
  !gapprompt@gap>| !gapinput@poset := PosetOfCongruences(coll);|
  <poset of 3 right congruences over <regular transformation monoid
   of size 3, degree 2 with 2 generators>>
  !gapprompt@gap>| !gapinput@OutNeighbours(poset);|
  [ [ 1 ], [ 2 ], [ 1, 2, 3 ] ]
\end{Verbatim}
 }

 

\subsection{\textcolor{Chapter }{JoinSemilatticeOfCongruences}}
\logpage{[ 13, 4, 11 ]}\nobreak
\hyperdef{L}{X87CF25A178B7F1AF}{}
{\noindent\textcolor{FuncColor}{$\triangleright$\enspace\texttt{JoinSemilatticeOfCongruences({\mdseries\slshape poset})\index{JoinSemilatticeOfCongruences@\texttt{JoinSemilatticeOfCongruences}}
\label{JoinSemilatticeOfCongruences}
}\hfill{\scriptsize (attribute)}}\\
\textbf{\indent Returns:\ }
A congruence poset.



 If \mbox{\texttt{\mdseries\slshape poset}} is a congruence poset (i.e. it satisfies \texttt{IsCongruencePoset} (\ref{IsCongruencePoset})), then this function returns the congruence poset formed by these congruences
partially ordered by containment, along with all their joins. This includes
the empty join which equals the trivial congruence. 

 The digraph returned by this function represents the Cayley graph of the
semilattice generated by \texttt{CongruencesOfPoset} (\ref{CongruencesOfPoset}) with identity adjoined. The reflexive transitive closure of this digraph is a
join semillatice in the sense of \texttt{IsJoinSemilatticeDigraph} (\textbf{Digraphs: IsJoinSemilatticeDigraph}). 

 See also \texttt{IsCongruencePoset} (\ref{IsCongruencePoset}) and \texttt{PosetOfCongruences} (\ref{PosetOfCongruences}). 
\begin{Verbatim}[commandchars=!@|,fontsize=\small,frame=single,label=Example]
  !gapprompt@gap>| !gapinput@S := SymmetricInverseMonoid(2);;|
  !gapprompt@gap>| !gapinput@pair1 := [PartialPerm([1], [1]), PartialPerm([2], [1])];;|
  !gapprompt@gap>| !gapinput@pair2 := [PartialPerm([1], [1]), PartialPerm([1, 2], [1, 2])];;|
  !gapprompt@gap>| !gapinput@pair3 := [PartialPerm([1, 2], [1, 2]),|
  !gapprompt@>| !gapinput@             PartialPerm([1, 2], [2, 1])];;|
  !gapprompt@gap>| !gapinput@coll := [RightSemigroupCongruence(S, pair1),|
  !gapprompt@>| !gapinput@            RightSemigroupCongruence(S, pair2),|
  !gapprompt@>| !gapinput@            RightSemigroupCongruence(S, pair3)];;|
  !gapprompt@gap>| !gapinput@D := JoinSemilatticeOfCongruences(PosetOfCongruences(coll));|
  <poset of 4 right congruences over
   <symmetric inverse monoid of degree 2>>
  !gapprompt@gap>| !gapinput@IsJoinSemilatticeDigraph(DigraphReflexiveTransitiveClosure(D));|
  true
\end{Verbatim}
 }

 

\subsection{\textcolor{Chapter }{GeneratingCongruencesOfJoinSemilattice}}
\logpage{[ 13, 4, 12 ]}\nobreak
\hyperdef{L}{X7ECE04727B6A58A3}{}
{\noindent\textcolor{FuncColor}{$\triangleright$\enspace\texttt{GeneratingCongruencesOfJoinSemilattice({\mdseries\slshape poset})\index{GeneratingCongruencesOfJoinSemilattice@\texttt{Generating}\-\texttt{Congruences}\-\texttt{Of}\-\texttt{Join}\-\texttt{Semilattice}}
\label{GeneratingCongruencesOfJoinSemilattice}
}\hfill{\scriptsize (attribute)}}\\
\textbf{\indent Returns:\ }
A list of congruences.



 If \mbox{\texttt{\mdseries\slshape poset}} satisfies \texttt{IsCayleyDigraphOfCongruences} (\ref{IsCayleyDigraphOfCongruences}), then this attribute holds the generating set for the semilattice of
congruences (where the operation is join). 
\begin{Verbatim}[commandchars=!@|,fontsize=\small,frame=single,label=Example]
  !gapprompt@gap>| !gapinput@S := OrderEndomorphisms(3);;|
  !gapprompt@gap>| !gapinput@D := CayleyDigraphOfCongruences(S);|
  <poset of 4 two-sided congruences over <regular transformation monoid
   of size 10, degree 3 with 3 generators>>
  !gapprompt@gap>| !gapinput@GeneratingCongruencesOfJoinSemilattice(D);|
  [ <universal semigroup congruence over <regular transformation monoid
       of size 10, degree 3 with 3 generators>>,
    <2-sided semigroup congruence over <regular transformation monoid
       of size 10, degree 3 with 3 generators> with 1 generating pairs>,
    <2-sided semigroup congruence over <regular transformation monoid
       of size 10, degree 3 with 3 generators> with 1 generating pairs>
   ]
\end{Verbatim}
 }

 

\subsection{\textcolor{Chapter }{MinimalCongruences (for a list or collection)}}
\logpage{[ 13, 4, 13 ]}\nobreak
\hyperdef{L}{X780E2B3F8509CE32}{}
{\noindent\textcolor{FuncColor}{$\triangleright$\enspace\texttt{MinimalCongruences({\mdseries\slshape coll})\index{MinimalCongruences@\texttt{MinimalCongruences}!for a list or collection}
\label{MinimalCongruences:for a list or collection}
}\hfill{\scriptsize (attribute)}}\\
\noindent\textcolor{FuncColor}{$\triangleright$\enspace\texttt{MinimalCongruences({\mdseries\slshape poset})\index{MinimalCongruences@\texttt{MinimalCongruences}!for a congruence poset}
\label{MinimalCongruences:for a congruence poset}
}\hfill{\scriptsize (attribute)}}\\
\textbf{\indent Returns:\ }
A list.



 If \mbox{\texttt{\mdseries\slshape coll}} is a list or collection of semigroup congruences (which may be left, right, or
two\texttt{\symbol{45}}sided) then this attribute returns a list of all the
congruences from \mbox{\texttt{\mdseries\slshape coll}} which do not contain any of the others as subrelations. 

 Alternatively, a congruence poset \mbox{\texttt{\mdseries\slshape poset}} can be specified; in this case, the congruences contained in \mbox{\texttt{\mdseries\slshape poset}} will be used in place of \mbox{\texttt{\mdseries\slshape coll}}, and information already known about their containments will be used. 

 This function should not be confused with \texttt{MinimalCongruencesOfSemigroup} (\ref{MinimalCongruencesOfSemigroup:for a semigroup}). See also \texttt{IsCongruencePoset} (\ref{IsCongruencePoset}) and \texttt{PosetOfCongruences} (\ref{PosetOfCongruences}). 

 
\begin{Verbatim}[commandchars=!@|,fontsize=\small,frame=single,label=Example]
  !gapprompt@gap>| !gapinput@S := SymmetricInverseMonoid(2);;|
  !gapprompt@gap>| !gapinput@pair1 := [PartialPerm([1], [1]), PartialPerm([2], [1])];;|
  !gapprompt@gap>| !gapinput@pair2 := [PartialPerm([1], [1]), PartialPerm([1, 2], [1, 2])];;|
  !gapprompt@gap>| !gapinput@pair3 := [PartialPerm([1, 2], [1, 2]),|
  !gapprompt@>| !gapinput@             PartialPerm([1, 2], [2, 1])];;|
  !gapprompt@gap>| !gapinput@coll := [RightSemigroupCongruence(S, pair1),|
  !gapprompt@>| !gapinput@            RightSemigroupCongruence(S, pair2),|
  !gapprompt@>| !gapinput@            RightSemigroupCongruence(S, pair3)];;|
  !gapprompt@gap>| !gapinput@MinimalCongruences(PosetOfCongruences(coll));|
  [ <right semigroup congruence over <symmetric inverse monoid of degree\
   2> with 1 generating pairs>,
    <right semigroup congruence over <symmetric inverse monoid of degree\
   2> with 1 generating pairs> ]
  !gapprompt@gap>| !gapinput@poset := LatticeOfCongruences(S);|
  <lattice of 4 two-sided congruences over
   <symmetric inverse monoid of degree 2>>
  !gapprompt@gap>| !gapinput@MinimalCongruences(poset);|
  [ <2-sided semigroup congruence over <symmetric inverse monoid of degr\
  ee 2> with 0 generating pairs> ]
\end{Verbatim}
 }

 

\subsection{\textcolor{Chapter }{NumberOfRightCongruences (for a semigroup, positive integer, and list or collection)}}
\logpage{[ 13, 4, 14 ]}\nobreak
\hyperdef{L}{X7AE16F237E862934}{}
{\noindent\textcolor{FuncColor}{$\triangleright$\enspace\texttt{NumberOfRightCongruences({\mdseries\slshape S, n, extra})\index{NumberOfRightCongruences@\texttt{NumberOfRightCongruences}!for a semigroup, positive integer, and list or collection}
\label{NumberOfRightCongruences:for a semigroup, positive integer, and list or collection}
}\hfill{\scriptsize (operation)}}\\
\noindent\textcolor{FuncColor}{$\triangleright$\enspace\texttt{NumberOfLeftCongruences({\mdseries\slshape S, n, extra})\index{NumberOfLeftCongruences@\texttt{NumberOfLeftCongruences}!for a semigroup, positive integer, and list or collection}
\label{NumberOfLeftCongruences:for a semigroup, positive integer, and list or collection}
}\hfill{\scriptsize (operation)}}\\
\noindent\textcolor{FuncColor}{$\triangleright$\enspace\texttt{NumberOfRightCongruences({\mdseries\slshape S, n})\index{NumberOfRightCongruences@\texttt{NumberOfRightCongruences}!for a semigroup, and a positive integer}
\label{NumberOfRightCongruences:for a semigroup, and a positive integer}
}\hfill{\scriptsize (operation)}}\\
\noindent\textcolor{FuncColor}{$\triangleright$\enspace\texttt{NumberOfLeftCongruences({\mdseries\slshape S, n})\index{NumberOfLeftCongruences@\texttt{NumberOfLeftCongruences}!for a semigroup, and a positive integer}
\label{NumberOfLeftCongruences:for a semigroup, and a positive integer}
}\hfill{\scriptsize (operation)}}\\
\noindent\textcolor{FuncColor}{$\triangleright$\enspace\texttt{NumberOfRightCongruences({\mdseries\slshape S})\index{NumberOfRightCongruences@\texttt{NumberOfRightCongruences}!for a semigroup}
\label{NumberOfRightCongruences:for a semigroup}
}\hfill{\scriptsize (attribute)}}\\
\noindent\textcolor{FuncColor}{$\triangleright$\enspace\texttt{NumberOfLeftCongruences({\mdseries\slshape S})\index{NumberOfLeftCongruences@\texttt{NumberOfLeftCongruences}!for a semigroup}
\label{NumberOfLeftCongruences:for a semigroup}
}\hfill{\scriptsize (operation)}}\\
\textbf{\indent Returns:\ }
A non\texttt{\symbol{45}}negative integer.



 \texttt{NumberOfRightCongruences} returns the number of right congruences of the semigroup \mbox{\texttt{\mdseries\slshape S}} with at most \mbox{\texttt{\mdseries\slshape n}} classes that contain the pairs in \mbox{\texttt{\mdseries\slshape extra}}; \texttt{NumberOfLeftCongruences} is defined dually for left congruences rather than right congruences. 

 If the optional third argument \mbox{\texttt{\mdseries\slshape extra}} is not present, then \texttt{NumberOfRightCongruences} returns the number of right congruences of \mbox{\texttt{\mdseries\slshape S}} with at most \mbox{\texttt{\mdseries\slshape n}} classes. 

 If the optional second argument \mbox{\texttt{\mdseries\slshape n}} is not present, then \texttt{NumberOfRightCongruences} returns the number of right congruences of \mbox{\texttt{\mdseries\slshape S}}. 

 Note that the 2 and 3 argument variants of this function can be applied to
infinite semigroups, but the 1 argument variant cannot. 

 If the lattice of right or left congruences of \mbox{\texttt{\mdseries\slshape S}} is known, then that is used by \texttt{NumberOfRightCongruences}. If this lattice is not known, then Sim's low index congruence algorithm is
used.

 See \texttt{IteratorOfRightCongruences} (\ref{IteratorOfRightCongruences:for a semigroup}) to actually obtain the congruences counted by this function. 
\begin{Verbatim}[commandchars=!@|,fontsize=\small,frame=single,label=Example]
  !gapprompt@gap>| !gapinput@S := PartitionMonoid(2);|
  <regular bipartition *-monoid of size 15, degree 2 with 3 generators>
  !gapprompt@gap>| !gapinput@NumberOfRightCongruences(S, 10);|
  86
  !gapprompt@gap>| !gapinput@NumberOfLeftCongruences(S, 10);|
  86
  !gapprompt@gap>| !gapinput@NumberOfRightCongruences(S, Size(S), [[S.1, S.2], [S.1, S.3]]);|
  1
  !gapprompt@gap>| !gapinput@NumberOfLeftCongruences(S, Size(S), [[S.1, S.2], [S.1, S.3]]);|
  1
\end{Verbatim}
 }

 

\subsection{\textcolor{Chapter }{IteratorOfRightCongruences (for a semigroup, positive integer, and list or collection)}}
\logpage{[ 13, 4, 15 ]}\nobreak
\hyperdef{L}{X807A5FCC87661FA4}{}
{\noindent\textcolor{FuncColor}{$\triangleright$\enspace\texttt{IteratorOfRightCongruences({\mdseries\slshape S, n, extra})\index{IteratorOfRightCongruences@\texttt{IteratorOfRightCongruences}!for a semigroup, positive integer, and list or collection}
\label{IteratorOfRightCongruences:for a semigroup, positive integer, and list or collection}
}\hfill{\scriptsize (operation)}}\\
\noindent\textcolor{FuncColor}{$\triangleright$\enspace\texttt{IteratorOfLeftCongruences({\mdseries\slshape S, n, extra})\index{IteratorOfLeftCongruences@\texttt{IteratorOfLeftCongruences}!for a semigroup, positive integer, and list or collection}
\label{IteratorOfLeftCongruences:for a semigroup, positive integer, and list or collection}
}\hfill{\scriptsize (operation)}}\\
\noindent\textcolor{FuncColor}{$\triangleright$\enspace\texttt{IteratorOfRightCongruences({\mdseries\slshape S, n})\index{IteratorOfRightCongruences@\texttt{IteratorOfRightCongruences}!for a semigroup, and a positive integer}
\label{IteratorOfRightCongruences:for a semigroup, and a positive integer}
}\hfill{\scriptsize (operation)}}\\
\noindent\textcolor{FuncColor}{$\triangleright$\enspace\texttt{IteratorOfLeftCongruences({\mdseries\slshape S, n})\index{IteratorOfLeftCongruences@\texttt{IteratorOfLeftCongruences}!for a semigroup, and a positive integer}
\label{IteratorOfLeftCongruences:for a semigroup, and a positive integer}
}\hfill{\scriptsize (operation)}}\\
\noindent\textcolor{FuncColor}{$\triangleright$\enspace\texttt{IteratorOfRightCongruences({\mdseries\slshape S})\index{IteratorOfRightCongruences@\texttt{IteratorOfRightCongruences}!for a semigroup}
\label{IteratorOfRightCongruences:for a semigroup}
}\hfill{\scriptsize (attribute)}}\\
\noindent\textcolor{FuncColor}{$\triangleright$\enspace\texttt{IteratorOfLeftCongruences({\mdseries\slshape S})\index{IteratorOfLeftCongruences@\texttt{IteratorOfLeftCongruences}!for a semigroup}
\label{IteratorOfLeftCongruences:for a semigroup}
}\hfill{\scriptsize (operation)}}\\
\textbf{\indent Returns:\ }
An iterator.



 \texttt{IteratorOfRightCongruences} returns an iterator where calling \texttt{NextIterator} (\textbf{Reference: NextIterator}) returns the next right congruence of the semigroup \mbox{\texttt{\mdseries\slshape S}} with at most \mbox{\texttt{\mdseries\slshape n}} classes that contain the pairs in \mbox{\texttt{\mdseries\slshape extra}}; \texttt{IteratorOfLeftCongruences} is defined dually for left congruences rather than right congruences. 

 If the optional third argument \mbox{\texttt{\mdseries\slshape extra}} is not present, then \texttt{IteratorOfRightCongruences} uses an empty list by default. 

 If the optional second argument \mbox{\texttt{\mdseries\slshape n}} is not present, then \texttt{IteratorOfRightCongruences} uses \texttt{Size(\mbox{\texttt{\mdseries\slshape S}})} by default. 

 Note that the 2 and 3 argument variants of this function can be applied to
infinite semigroups, but the 1 argument variant cannot. 

 If the lattice of right or left congruences of \mbox{\texttt{\mdseries\slshape S}} is known, then that is used by \texttt{IteratorOfRightCongruences}. If this lattice is not known, then Sim's low index congruence algorithm is
used. 
\begin{Verbatim}[commandchars=!@|,fontsize=\small,frame=single,label=Example]
  !gapprompt@gap>| !gapinput@F := FreeMonoidAndAssignGeneratorVars("a", "b");|
  <free monoid on the generators [ a, b ]>
  !gapprompt@gap>| !gapinput@R := [[a ^ 3, a], [b ^ 2, b], [(a * b) ^ 2, a]];|
  [ [ a^3, a ], [ b^2, b ], [ (a*b)^2, a ] ]
  !gapprompt@gap>| !gapinput@S := F / R;|
  <fp monoid with 2 generators and 3 relations of length 14>
  !gapprompt@gap>| !gapinput@NumberOfRightCongruences(S);|
  6
  !gapprompt@gap>| !gapinput@it := IteratorOfRightCongruences(S);|
  <iterator>
  !gapprompt@gap>| !gapinput@OutNeighbours(WordGraph(NextIterator(it)));|
  [ [ 1, 1 ] ]
  !gapprompt@gap>| !gapinput@OutNeighbours(WordGraph(NextIterator(it)));|
  [ [ 2, 1 ], [ 2, 2 ] ]
  !gapprompt@gap>| !gapinput@OutNeighbours(WordGraph(NextIterator(it)));|
  [ [ 2, 2 ], [ 2, 2 ] ]
  !gapprompt@gap>| !gapinput@OutNeighbours(WordGraph(NextIterator(it)));|
  [ [ 2, 3 ], [ 2, 2 ], [ 2, 3 ] ]
  !gapprompt@gap>| !gapinput@OutNeighbours(WordGraph(NextIterator(it)));|
  [ [ 2, 3 ], [ 2, 2 ], [ 3, 3 ] ]
  !gapprompt@gap>| !gapinput@OutNeighbours(WordGraph(NextIterator(it)));|
  [ [ 2, 3 ], [ 2, 2 ], [ 4, 3 ], [ 4, 4 ] ]
  !gapprompt@gap>| !gapinput@NextIterator(it);|
  fail
\end{Verbatim}
 }

 }

 
\section{\textcolor{Chapter }{ Comparing congruences }}\logpage{[ 13, 5, 0 ]}
\hyperdef{L}{X857D750579B87DBF}{}
{
  

\subsection{\textcolor{Chapter }{IsSubrelation}}
\logpage{[ 13, 5, 1 ]}\nobreak
\hyperdef{L}{X85075F1D878512F5}{}
{\noindent\textcolor{FuncColor}{$\triangleright$\enspace\texttt{IsSubrelation({\mdseries\slshape cong1, cong2})\index{IsSubrelation@\texttt{IsSubrelation}}
\label{IsSubrelation}
}\hfill{\scriptsize (operation)}}\\
\textbf{\indent Returns:\ }
True or false.



 If \mbox{\texttt{\mdseries\slshape cong1}} and \mbox{\texttt{\mdseries\slshape cong2}} are congruences over the same semigroup, then this operation returns whether \mbox{\texttt{\mdseries\slshape cong2}} is a refinement of \mbox{\texttt{\mdseries\slshape cong1}}, i.e. whether every pair in \mbox{\texttt{\mdseries\slshape cong2}} is contained in \mbox{\texttt{\mdseries\slshape cong1}}. 

 
\begin{Verbatim}[commandchars=!@|,fontsize=\small,frame=single,label=Example]
  !gapprompt@gap>| !gapinput@S := ReesZeroMatrixSemigroup(SymmetricGroup(3),|
  !gapprompt@>| !gapinput@[[(), (1, 3, 2)], [(1, 2), 0]]);;|
  !gapprompt@gap>| !gapinput@cong1 := SemigroupCongruence(S, [RMSElement(S, 1, (1, 2, 3), 1),|
  !gapprompt@>| !gapinput@                                    RMSElement(S, 1, (), 1)]);;|
  !gapprompt@gap>| !gapinput@cong2 := SemigroupCongruence(S, []);;|
  !gapprompt@gap>| !gapinput@IsSubrelation(cong1, cong2);|
  true
  !gapprompt@gap>| !gapinput@IsSubrelation(cong2, cong1);|
  false
\end{Verbatim}
 }

 

\subsection{\textcolor{Chapter }{IsSuperrelation}}
\logpage{[ 13, 5, 2 ]}\nobreak
\hyperdef{L}{X83878AED7A8E75BE}{}
{\noindent\textcolor{FuncColor}{$\triangleright$\enspace\texttt{IsSuperrelation({\mdseries\slshape cong1, cong2})\index{IsSuperrelation@\texttt{IsSuperrelation}}
\label{IsSuperrelation}
}\hfill{\scriptsize (operation)}}\\
\textbf{\indent Returns:\ }
True or false.



 If \mbox{\texttt{\mdseries\slshape cong1}} and \mbox{\texttt{\mdseries\slshape cong2}} are congruences over the same semigroup, then this operation returns whether \mbox{\texttt{\mdseries\slshape cong1}} is a refinement of \mbox{\texttt{\mdseries\slshape cong2}}, i.e. whether every pair in \mbox{\texttt{\mdseries\slshape cong1}} is contained in \mbox{\texttt{\mdseries\slshape cong2}}. 

 See \texttt{IsSubrelation} (\ref{IsSubrelation}). 

 
\begin{Verbatim}[commandchars=!@|,fontsize=\small,frame=single,label=Example]
  !gapprompt@gap>| !gapinput@S := ReesZeroMatrixSemigroup(SymmetricGroup(3),|
  !gapprompt@>| !gapinput@[[(), (1, 3, 2)], [(1, 2), 0]]);;|
  !gapprompt@gap>| !gapinput@cong1 := SemigroupCongruence(S, [RMSElement(S, 1, (1, 2, 3), 1),|
  !gapprompt@>| !gapinput@                                    RMSElement(S, 1, (), 1)]);;|
  !gapprompt@gap>| !gapinput@cong2 := SemigroupCongruence(S, []);;|
  !gapprompt@gap>| !gapinput@IsSuperrelation(cong1, cong2);|
  false
  !gapprompt@gap>| !gapinput@IsSuperrelation(cong2, cong1);|
  true
\end{Verbatim}
 }

 

\subsection{\textcolor{Chapter }{MeetSemigroupCongruences}}
\logpage{[ 13, 5, 3 ]}\nobreak
\hyperdef{L}{X7952A5A5789C6F60}{}
{\noindent\textcolor{FuncColor}{$\triangleright$\enspace\texttt{MeetSemigroupCongruences({\mdseries\slshape c1, c2})\index{MeetSemigroupCongruences@\texttt{MeetSemigroupCongruences}}
\label{MeetSemigroupCongruences}
}\hfill{\scriptsize (operation)}}\\
\noindent\textcolor{FuncColor}{$\triangleright$\enspace\texttt{MeetLeftSemigroupCongruences({\mdseries\slshape c1, c2})\index{MeetLeftSemigroupCongruences@\texttt{MeetLeftSemigroupCongruences}}
\label{MeetLeftSemigroupCongruences}
}\hfill{\scriptsize (operation)}}\\
\noindent\textcolor{FuncColor}{$\triangleright$\enspace\texttt{MeetRightSemigroupCongruences({\mdseries\slshape c1, c2})\index{MeetRightSemigroupCongruences@\texttt{MeetRightSemigroupCongruences}}
\label{MeetRightSemigroupCongruences}
}\hfill{\scriptsize (operation)}}\\
\textbf{\indent Returns:\ }
A semigroup congruence.



 This operation returns the \emph{meet} of the two semigroup congruences \mbox{\texttt{\mdseries\slshape c1}} and \mbox{\texttt{\mdseries\slshape c2}} \texttt{\symbol{45}}\texttt{\symbol{45}} that is, the largest semigroup
congruence contained in both \mbox{\texttt{\mdseries\slshape c1}} and \mbox{\texttt{\mdseries\slshape c2}}.

 
\begin{Verbatim}[commandchars=!@|,fontsize=\small,frame=single,label=Example]
  !gapprompt@gap>| !gapinput@S := ReesZeroMatrixSemigroup(SymmetricGroup(3),|
  !gapprompt@>| !gapinput@[[(), (1, 3, 2)], [(1, 2), 0]]);;|
  !gapprompt@gap>| !gapinput@cong1 := SemigroupCongruence(S, [RMSElement(S, 1, (1, 2, 3), 1),|
  !gapprompt@>| !gapinput@                                    RMSElement(S, 1, (), 1)]);;|
  !gapprompt@gap>| !gapinput@cong2 := SemigroupCongruence(S, []);;|
  !gapprompt@gap>| !gapinput@MeetSemigroupCongruences(cong1, cong2);|
  <semigroup congruence over <Rees 0-matrix semigroup 2x2 over
    Sym( [ 1 .. 3 ] )> with linked triple (1,2,2)>
\end{Verbatim}
 }

 

\subsection{\textcolor{Chapter }{JoinSemigroupCongruences}}
\logpage{[ 13, 5, 4 ]}\nobreak
\hyperdef{L}{X8262D5207DBF3C5B}{}
{\noindent\textcolor{FuncColor}{$\triangleright$\enspace\texttt{JoinSemigroupCongruences({\mdseries\slshape c1, c2})\index{JoinSemigroupCongruences@\texttt{JoinSemigroupCongruences}}
\label{JoinSemigroupCongruences}
}\hfill{\scriptsize (operation)}}\\
\noindent\textcolor{FuncColor}{$\triangleright$\enspace\texttt{JoinLeftSemigroupCongruences({\mdseries\slshape c1, c2})\index{JoinLeftSemigroupCongruences@\texttt{JoinLeftSemigroupCongruences}}
\label{JoinLeftSemigroupCongruences}
}\hfill{\scriptsize (operation)}}\\
\noindent\textcolor{FuncColor}{$\triangleright$\enspace\texttt{JoinRightSemigroupCongruences({\mdseries\slshape c1, c2})\index{JoinRightSemigroupCongruences@\texttt{JoinRightSemigroupCongruences}}
\label{JoinRightSemigroupCongruences}
}\hfill{\scriptsize (operation)}}\\
\textbf{\indent Returns:\ }
A semigroup congruence.



 This operation returns the \emph{join} of the two semigroup congruences \mbox{\texttt{\mdseries\slshape c1}} and \mbox{\texttt{\mdseries\slshape c2}} \texttt{\symbol{45}}\texttt{\symbol{45}} that is, the smallest semigroup
congruence containing all the relations in both \mbox{\texttt{\mdseries\slshape c1}} and \mbox{\texttt{\mdseries\slshape c2}}.

 
\begin{Verbatim}[commandchars=!@|,fontsize=\small,frame=single,label=Example]
  !gapprompt@gap>| !gapinput@S := ReesZeroMatrixSemigroup(SymmetricGroup(3),|
  !gapprompt@>| !gapinput@[[(), (1, 3, 2)], [(1, 2), 0]]);;|
  !gapprompt@gap>| !gapinput@cong1 := SemigroupCongruence(S, [RMSElement(S, 1, (1, 2, 3), 1),|
  !gapprompt@>| !gapinput@                                    RMSElement(S, 1, (), 1)]);;|
  !gapprompt@gap>| !gapinput@cong2 := SemigroupCongruence(S, []);;|
  !gapprompt@gap>| !gapinput@JoinSemigroupCongruences(cong1, cong2);|
  <semigroup congruence over <Rees 0-matrix semigroup 2x2 over
    Sym( [ 1 .. 3 ] )> with linked triple (3,2,2)>
\end{Verbatim}
 }

 }

 
\section{\textcolor{Chapter }{Congruences on Rees matrix semigroups}}\logpage{[ 13, 6, 0 ]}
\hyperdef{L}{X7A6478D1831DD787}{}
{
  This section describes the implementation of congruences of simple and
0\texttt{\symbol{45}}simple semigroups in the \textsf{Semigroups} package, and the functions associated with them. This code and this part of
the manual were written by Michael Young. Most of the theorems used in this
chapter are from Section 3.5 of \cite{Howie1995aa}.

 By the Rees Theorem, any 0\texttt{\symbol{45}}simple semigroup $S$ is isomorphic to a \emph{Rees 0\texttt{\symbol{45}}matrix semigroup} (see  (\textbf{Reference: Rees Matrix Semigroups})) over a group, with a regular sandwich matrix.  That is, 
\[S \cong \mathcal{M} ^ 0[G; I, \Lambda; P], \]
 where $G$ is a group, $\Lambda$ and $I$ are non\texttt{\symbol{45}}empty sets, and $P$ is regular in the sense that it has no rows or columns consisting solely of
zeroes.  

 The congruences of a Rees 0\texttt{\symbol{45}}matrix semigroup are in
1\texttt{\symbol{45}}1 correspondence with the \emph{linked triple}, which is a triple of the form \texttt{[N, S, T]} where: 
\begin{itemize}
\item  \texttt{N} is a normal subgroup of the underlying group \texttt{G}, 
\item  \texttt{S} is an equivalence relation on the columns of \texttt{P}, 
\item  \texttt{T} is an equivalence relation on the rows of \texttt{P}, 
\end{itemize}
 satisfying the following conditions: 
\begin{itemize}
\item  a pair of \texttt{S}\texttt{\symbol{45}}related columns must contain zeroes in precisely the same
rows, 
\item  a pair of \texttt{T}\texttt{\symbol{45}}related rows must contain zeroes in precisely the same
columns, 
\item  if \texttt{i} and \texttt{j} are \texttt{S}\texttt{\symbol{45}}related, \texttt{k} and \texttt{l} are \texttt{T}\texttt{\symbol{45}}related and the matrix entries $p_{k, i}, p_{k, j}, p_{l, i}, p_{l, j} \neq 0$, then $q_{k, l, i, j} \in N$, where 
\[q_{k, l, i, j} = p_{k, i} p_{l, i} ^ {-1} p_{l, j} p_{k, j} ^ {-1}.\]
 
\end{itemize}
 By Theorem 3.5.9 in \cite{Howie1995aa}, for any finite 0\texttt{\symbol{45}}simple Rees 0\texttt{\symbol{45}}matrix
semigroup, there is a bijection between its non\texttt{\symbol{45}}universal
congruences and its linked triples. In this way, we can internally represent
any congruence of such a semigroup by storing its associated linked triple
instead of a set of generating pairs, and thus perform many calculations on it
more efficiently.

 If a congruence is defined by a linked triple \texttt{(N, S, T)}, then a single class of that congruence can be defined by a triple \texttt{(Nx, i / S, k / S)}, where \texttt{Nx} is a right coset of \texttt{N}, \texttt{i / S} is the equivalence class of \texttt{i} in \texttt{S}, and \texttt{k / S} is the equivalence class of \texttt{k} in \texttt{T}. Thus we can internally represent any class of such a congruence as a triple
simply consisting of a right coset and two positive integers.

 An analogous condition exists for finite simple Rees matrix semigroups without
zero.

 

\subsection{\textcolor{Chapter }{IsRMSCongruenceByLinkedTriple}}
\logpage{[ 13, 6, 1 ]}\nobreak
\hyperdef{L}{X7F4AFD7F7E163022}{}
{\noindent\textcolor{FuncColor}{$\triangleright$\enspace\texttt{IsRMSCongruenceByLinkedTriple({\mdseries\slshape obj})\index{IsRMSCongruenceByLinkedTriple@\texttt{IsRMSCongruenceByLinkedTriple}}
\label{IsRMSCongruenceByLinkedTriple}
}\hfill{\scriptsize (category)}}\\
\noindent\textcolor{FuncColor}{$\triangleright$\enspace\texttt{IsRZMSCongruenceByLinkedTriple({\mdseries\slshape obj})\index{IsRZMSCongruenceByLinkedTriple@\texttt{IsRZMSCongruenceByLinkedTriple}}
\label{IsRZMSCongruenceByLinkedTriple}
}\hfill{\scriptsize (category)}}\\
\textbf{\indent Returns:\ }
\texttt{true} or \texttt{false}.



 These categories describe a type of semigroup congruence over a Rees matrix or
0\texttt{\symbol{45}}matrix semigroup. Externally, an object of this type may
be used in the same way as any other object in the category \texttt{IsSemigroupCongruence} (\textbf{Reference: IsSemigroupCongruence}) but it is represented internally by its \emph{linked triple}, and certain functions may take advantage of this information to reduce
computation times.

 An object of this type may be constructed with \texttt{RMSCongruenceByLinkedTriple} or \texttt{RZMSCongruenceByLinkedTriple}, or this representation may be selected automatically by \texttt{SemigroupCongruence} (\ref{SemigroupCongruence}). 
\begin{Verbatim}[commandchars=!@|,fontsize=\small,frame=single,label=Example]
  !gapprompt@gap>| !gapinput@G := Group([(1, 4, 5), (1, 5, 3, 4)]);;|
  !gapprompt@gap>| !gapinput@mat := [[0, 0, (1, 4, 5), 0, 0, (1, 4, 3, 5)],|
  !gapprompt@>| !gapinput@           [0, (), 0, 0, (3, 5), 0],|
  !gapprompt@>| !gapinput@           [(), 0, 0, (3, 5), 0, 0]];;|
  !gapprompt@gap>| !gapinput@S := ReesZeroMatrixSemigroup(G, mat);;|
  !gapprompt@gap>| !gapinput@N := Group([(1, 4)(3, 5), (1, 5)(3, 4)]);;|
  !gapprompt@gap>| !gapinput@colBlocks := [[1], [2, 5], [3, 6], [4]];;|
  !gapprompt@gap>| !gapinput@rowBlocks := [[1], [2], [3]];;|
  !gapprompt@gap>| !gapinput@cong := RZMSCongruenceByLinkedTriple(S, N, colBlocks, rowBlocks);;|
  !gapprompt@gap>| !gapinput@IsRZMSCongruenceByLinkedTriple(cong);|
  true
\end{Verbatim}
 }

 

\subsection{\textcolor{Chapter }{RMSCongruenceByLinkedTriple}}
\logpage{[ 13, 6, 2 ]}\nobreak
\hyperdef{L}{X87A475E4847D3C96}{}
{\noindent\textcolor{FuncColor}{$\triangleright$\enspace\texttt{RMSCongruenceByLinkedTriple({\mdseries\slshape S, N, colBlocks, rowBlocks})\index{RMSCongruenceByLinkedTriple@\texttt{RMSCongruenceByLinkedTriple}}
\label{RMSCongruenceByLinkedTriple}
}\hfill{\scriptsize (function)}}\\
\noindent\textcolor{FuncColor}{$\triangleright$\enspace\texttt{RZMSCongruenceByLinkedTriple({\mdseries\slshape S, N, colBlocks, rowBlocks})\index{RZMSCongruenceByLinkedTriple@\texttt{RZMSCongruenceByLinkedTriple}}
\label{RZMSCongruenceByLinkedTriple}
}\hfill{\scriptsize (function)}}\\
\textbf{\indent Returns:\ }
 A Rees matrix or 0\texttt{\symbol{45}}matrix semigroup congruence by linked
triple. 



 This function returns a semigroup congruence over the Rees matrix or
0\texttt{\symbol{45}}matrix semigroup \mbox{\texttt{\mdseries\slshape S}} corresponding to the linked triple (\mbox{\texttt{\mdseries\slshape N}}, \mbox{\texttt{\mdseries\slshape colBlocks}}, \mbox{\texttt{\mdseries\slshape rowBlocks}}). The argument \mbox{\texttt{\mdseries\slshape N}} should be a normal subgroup of the underlying semigroup of \mbox{\texttt{\mdseries\slshape S}}; \mbox{\texttt{\mdseries\slshape colBlocks}} should be a partition of the columns of the matrix of \mbox{\texttt{\mdseries\slshape S}}; and \mbox{\texttt{\mdseries\slshape rowBlocks}} should be a partition of the rows of the matrix of \mbox{\texttt{\mdseries\slshape S}}. For example, if the matrix has 5 rows, then a possibility for \mbox{\texttt{\mdseries\slshape rowBlocks}} might be \texttt{[[1, 3], [2, 5], [4]]}.

 If the arguments describe a valid linked triple on \mbox{\texttt{\mdseries\slshape S}}, then an object in the category \texttt{IsRZMSCongruenceByLinkedTriple} is returned. This object can be used like any other semigroup congruence in \textsf{GAP}.

 If the arguments describe a triple which is not \emph{linked} in the sense described above, then this function returns an error. 
\begin{Verbatim}[commandchars=!@|,fontsize=\small,frame=single,label=Example]
  !gapprompt@gap>| !gapinput@G := Group([(1, 4, 5), (1, 5, 3, 4)]);;|
  !gapprompt@gap>| !gapinput@mat := [[0, 0, (1, 4, 5), 0, 0, (1, 4, 3, 5)],|
  !gapprompt@>| !gapinput@           [0, (), 0, 0, (3, 5), 0],|
  !gapprompt@>| !gapinput@           [(), 0, 0, (3, 5), 0, 0]];;|
  !gapprompt@gap>| !gapinput@S := ReesZeroMatrixSemigroup(G, mat);;|
  !gapprompt@gap>| !gapinput@N := Group([(1, 4)(3, 5), (1, 5)(3, 4)]);;|
  !gapprompt@gap>| !gapinput@colBlocks := [[1], [2, 5], [3, 6], [4]];;|
  !gapprompt@gap>| !gapinput@rowBlocks := [[1], [2], [3]];;|
  !gapprompt@gap>| !gapinput@cong := RZMSCongruenceByLinkedTriple(S, N, colBlocks, rowBlocks);;|
\end{Verbatim}
 }

 

\subsection{\textcolor{Chapter }{IsRMSCongruenceClassByLinkedTriple}}
\logpage{[ 13, 6, 3 ]}\nobreak
\hyperdef{L}{X79E4CF7B79B25AE3}{}
{\noindent\textcolor{FuncColor}{$\triangleright$\enspace\texttt{IsRMSCongruenceClassByLinkedTriple({\mdseries\slshape obj})\index{IsRMSCongruenceClassByLinkedTriple@\texttt{IsRMSCongruenceClassByLinkedTriple}}
\label{IsRMSCongruenceClassByLinkedTriple}
}\hfill{\scriptsize (category)}}\\
\noindent\textcolor{FuncColor}{$\triangleright$\enspace\texttt{IsRZMSCongruenceClassByLinkedTriple({\mdseries\slshape obj})\index{IsRZMSCongruenceClassByLinkedTriple@\texttt{IsRZMSCongruenceClassByLinkedTriple}}
\label{IsRZMSCongruenceClassByLinkedTriple}
}\hfill{\scriptsize (category)}}\\
\textbf{\indent Returns:\ }
\texttt{true} or \texttt{false}.



 These categories contain the congruence classes of a semigroup congruence of
the categories \texttt{IsRMSCongruenceByLinkedTriple} (\ref{IsRMSCongruenceByLinkedTriple}) and \texttt{IsRZMSCongruenceByLinkedTriple} (\ref{IsRZMSCongruenceByLinkedTriple}) respectively. 

 An object of one of these types may be used in the same way as any other
object in the category \texttt{IsCongruenceClass} (\ref{IsCongruenceClass}), but the class is represented internally by information related to the
congruence's \emph{linked triple}, and certain functions may take advantage of this information to reduce
computation times. 

 
\begin{Verbatim}[commandchars=!@|,fontsize=\small,frame=single,label=Example]
  !gapprompt@gap>| !gapinput@G := Group([(1, 4, 5), (1, 5, 3, 4)]);;|
  !gapprompt@gap>| !gapinput@mat := [[0, 0, (1, 4, 5), 0, 0, (1, 4, 3, 5)],|
  !gapprompt@>| !gapinput@           [0, (), 0, 0, (3, 5), 0],|
  !gapprompt@>| !gapinput@           [(), 0, 0, (3, 5), 0, 0]];;|
  !gapprompt@gap>| !gapinput@S := ReesZeroMatrixSemigroup(G, mat);;|
  !gapprompt@gap>| !gapinput@N := Group([(1, 4)(3, 5), (1, 5)(3, 4)]);;|
  !gapprompt@gap>| !gapinput@colBlocks := [[1], [2, 5], [3, 6], [4]];;|
  !gapprompt@gap>| !gapinput@rowBlocks := [[1], [2], [3]];;|
  !gapprompt@gap>| !gapinput@cong := RZMSCongruenceByLinkedTriple(S, N, colBlocks, rowBlocks);;|
  !gapprompt@gap>| !gapinput@classes := EquivalenceClasses(cong);;|
  !gapprompt@gap>| !gapinput@IsRZMSCongruenceClassByLinkedTriple(classes[1]);|
  true
\end{Verbatim}
 }

 

\subsection{\textcolor{Chapter }{RMSCongruenceClassByLinkedTriple}}
\logpage{[ 13, 6, 4 ]}\nobreak
\hyperdef{L}{X7E9AB940868FCC9D}{}
{\noindent\textcolor{FuncColor}{$\triangleright$\enspace\texttt{RMSCongruenceClassByLinkedTriple({\mdseries\slshape cong, nCoset, colClass, rowClass})\index{RMSCongruenceClassByLinkedTriple@\texttt{RMSCongruenceClassByLinkedTriple}}
\label{RMSCongruenceClassByLinkedTriple}
}\hfill{\scriptsize (operation)}}\\
\noindent\textcolor{FuncColor}{$\triangleright$\enspace\texttt{RZMSCongruenceClassByLinkedTriple({\mdseries\slshape cong, nCoset, colClass, rowClass})\index{RZMSCongruenceClassByLinkedTriple@\texttt{RZMSCongruenceClassByLinkedTriple}}
\label{RZMSCongruenceClassByLinkedTriple}
}\hfill{\scriptsize (operation)}}\\
\textbf{\indent Returns:\ }
 A Rees matrix or 0\texttt{\symbol{45}}matrix semigroup congruence class by
linked triple. 



 This operation returns one congruence class of the congruence \mbox{\texttt{\mdseries\slshape cong}}, as defined by the other three parameters.

 The argument \mbox{\texttt{\mdseries\slshape cong}} must be a Rees matrix or 0\texttt{\symbol{45}}matrix semigroup congruence by
linked triple. If the linked triple consists of the three parameters \texttt{N}, \texttt{colBlocks} and \texttt{rowBlocks}, then \mbox{\texttt{\mdseries\slshape nCoset}} must be a right coset of \texttt{N}, \mbox{\texttt{\mdseries\slshape colClass}} must be a positive integer corresponding to a position in the list \texttt{colBlocks}, and \mbox{\texttt{\mdseries\slshape rowClass}} must be a positive integer corresponding to a position in the list \texttt{rowBlocks}.

 If the arguments are valid, an \texttt{IsRMSCongruenceClassByLinkedTriple} or \texttt{IsRZMSCongruenceClassByLinkedTriple} object is returned, which can be used like any other equivalence class in \textsf{GAP}. Otherwise, an error is returned. 
\begin{Verbatim}[commandchars=!@|,fontsize=\small,frame=single,label=Example]
  !gapprompt@gap>| !gapinput@G := Group([(1, 4, 5), (1, 5, 3, 4)]);;|
  !gapprompt@gap>| !gapinput@mat := [[0, 0, (1, 4, 5), 0, 0, (1, 4, 3, 5)],|
  !gapprompt@>| !gapinput@           [0, (), 0, 0, (3, 5), 0],|
  !gapprompt@>| !gapinput@           [(), 0, 0, (3, 5), 0, 0]];;|
  !gapprompt@gap>| !gapinput@S := ReesZeroMatrixSemigroup(G, mat);;|
  !gapprompt@gap>| !gapinput@N := Group([(1, 4)(3, 5), (1, 5)(3, 4)]);;|
  !gapprompt@gap>| !gapinput@colBlocks := [[1], [2, 5], [3, 6], [4]];;|
  !gapprompt@gap>| !gapinput@rowBlocks := [[1], [2], [3]];;|
  !gapprompt@gap>| !gapinput@cong := RZMSCongruenceByLinkedTriple(S, N, colBlocks, rowBlocks);;|
  !gapprompt@gap>| !gapinput@class := RZMSCongruenceClassByLinkedTriple(cong,|
  !gapprompt@>| !gapinput@RightCoset(N, (1, 5)), 2, 3);|
  <2-sided congruence class of (2,(3,4),3)>
\end{Verbatim}
 }

 

\subsection{\textcolor{Chapter }{IsLinkedTriple}}
\logpage{[ 13, 6, 5 ]}\nobreak
\hyperdef{L}{X7B19CACF7A37ADBC}{}
{\noindent\textcolor{FuncColor}{$\triangleright$\enspace\texttt{IsLinkedTriple({\mdseries\slshape S, N, colBlocks, rowBlocks})\index{IsLinkedTriple@\texttt{IsLinkedTriple}}
\label{IsLinkedTriple}
}\hfill{\scriptsize (operation)}}\\
\textbf{\indent Returns:\ }
\texttt{true} or \texttt{false}.



 This operation returns true if and only if the arguments (\mbox{\texttt{\mdseries\slshape N}}, \mbox{\texttt{\mdseries\slshape colBlocks}}, \mbox{\texttt{\mdseries\slshape rowBlocks}}) describe a linked triple of the Rees matrix or 0\texttt{\symbol{45}}matrix
semigroup \mbox{\texttt{\mdseries\slshape S}}, as described above. 
\begin{Verbatim}[commandchars=!@|,fontsize=\small,frame=single,label=Example]
  !gapprompt@gap>| !gapinput@G := Group([(1, 4, 5), (1, 5, 3, 4)]);;|
  !gapprompt@gap>| !gapinput@mat := [[0, 0, (1, 4, 5), 0, 0, (1, 4, 3, 5)],|
  !gapprompt@>| !gapinput@           [0, (), 0, 0, (3, 5), 0],|
  !gapprompt@>| !gapinput@           [(), 0, 0, (3, 5), 0, 0]];;|
  !gapprompt@gap>| !gapinput@S := ReesZeroMatrixSemigroup(G, mat);;|
  !gapprompt@gap>| !gapinput@N := Group([(1, 4)(3, 5), (1, 5)(3, 4)]);;|
  !gapprompt@gap>| !gapinput@colBlocks := [[1], [2, 5], [3, 6], [4]];;|
  !gapprompt@gap>| !gapinput@rowBlocks := [[1], [2], [3]];;|
  !gapprompt@gap>| !gapinput@IsLinkedTriple(S, N, colBlocks, rowBlocks);|
  true
\end{Verbatim}
 }

 

\subsection{\textcolor{Chapter }{AsSemigroupCongruenceByGeneratingPairs}}
\logpage{[ 13, 6, 6 ]}\nobreak
\hyperdef{L}{X7DB7E32E865AD95D}{}
{\noindent\textcolor{FuncColor}{$\triangleright$\enspace\texttt{AsSemigroupCongruenceByGeneratingPairs({\mdseries\slshape cong})\index{AsSemigroupCongruenceByGeneratingPairs@\texttt{AsSemigroup}\-\texttt{Congruence}\-\texttt{By}\-\texttt{Generating}\-\texttt{Pairs}}
\label{AsSemigroupCongruenceByGeneratingPairs}
}\hfill{\scriptsize (operation)}}\\
\textbf{\indent Returns:\ }
A semigroup congruence.



 This operation takes \mbox{\texttt{\mdseries\slshape cong}}, a semigroup congruence, and returns the same congruence relation, but
described by \textsf{GAP}'s default method of defining semigroup congruences: a set of generating pairs
for the congruence. 
\begin{Verbatim}[commandchars=!@|,fontsize=\small,frame=single,label=Example]
  !gapprompt@gap>| !gapinput@S := ReesZeroMatrixSemigroup(SymmetricGroup(3),|
  !gapprompt@>| !gapinput@                                [[(), (1, 3, 2)], [(1, 2), 0]]);;|
  !gapprompt@gap>| !gapinput@cong := CongruencesOfSemigroup(S)[3];;|
  !gapprompt@gap>| !gapinput@AsSemigroupCongruenceByGeneratingPairs(cong);|
  <semigroup congruence over <Rees 0-matrix semigroup 2x2 over
    Sym( [ 1 .. 3 ] )> with linked triple (3,2,2)>
\end{Verbatim}
 }

 }

 
\section{\textcolor{Chapter }{ Congruences on inverse semigroups }}\label{Congruences on inverse semigroups}
\logpage{[ 13, 7, 0 ]}
\hyperdef{L}{X7BFDC38178940AE6}{}
{
  This section describes the implementation of congruences of inverse semigroups
in the \textsf{Semigroups} package, and the functions associated with them. This code and this part of
the manual were written by Michael Young. Most of the theorems used in this
chapter are from Section 5.3 of \cite{Howie1995aa}.

 The congruences of an inverse semigroup are in 1\texttt{\symbol{45}}1
correspondence with its \emph{congruence pairs}. A congruence pair is a pair \texttt{(N, t)} such that: 
\begin{itemize}
\item  \texttt{N} is a normal subsemigroup of \texttt{S} \texttt{\symbol{45}}\texttt{\symbol{45}} that is, a
self\texttt{\symbol{45}}conjugate subsemigroup which contains all the
idempotents of \texttt{S}, 
\item  \texttt{t} is a normal congruence on \texttt{E}, the subsemigroup of all idempotents in \texttt{S} \texttt{\symbol{45}}\texttt{\symbol{45}} that is, a congruence on \texttt{E} such that if $(e, f)$ is a pair in \texttt{t}, then the pair $(a ^ {-1} e a, a ^ {-1} f a)$ is also in \texttt{t}, 
\end{itemize}
 satisfying the following conditions: 
\begin{itemize}
\item  If $ae \in N$ and $(e, a ^ {-1} a) \in t$, then $a \in N$, 
\item  If $a \in N$, then $(aa ^ {-1} , a ^ {-1} a) \in t$. 
\end{itemize}
 By Theorem 5.3.3 in \cite{Howie1995aa}, for any inverse semigroup, there is a bijection between its congruences and
its congruence pairs. In this way, we can internally represent any congruence
of such a semigroup by storing its associated congruence pair instead of a set
of generating pairs, and thus perform many calculations on it more
efficiently.

 If we have a congruence \texttt{C} with congruence pair \texttt{(N, t)}, it turns out that \texttt{N} is its \emph{kernel} (that is, the set of all elements congruent to an idempotent) and that \texttt{t} is its \emph{trace} (that is, the restriction of \texttt{C} to the idempotents). Hence, we refer to a congruence stored in this format as
a "congruence by kernel and trace".

 See \texttt{cong{\textunderscore}by{\textunderscore}ker{\textunderscore}trace{\textunderscore}threshold} in Section \ref{Options when creating semigroups} for details on when this method is used. 

 

\subsection{\textcolor{Chapter }{IsInverseSemigroupCongruenceByKernelTrace}}
\logpage{[ 13, 7, 1 ]}\nobreak
\hyperdef{L}{X8546E48E85A2A7E8}{}
{\noindent\textcolor{FuncColor}{$\triangleright$\enspace\texttt{IsInverseSemigroupCongruenceByKernelTrace({\mdseries\slshape cong})\index{IsInverseSemigroupCongruenceByKernelTrace@\texttt{IsInverse}\-\texttt{Semigroup}\-\texttt{Congruence}\-\texttt{By}\-\texttt{Kernel}\-\texttt{Trace}}
\label{IsInverseSemigroupCongruenceByKernelTrace}
}\hfill{\scriptsize (Category)}}\\
\textbf{\indent Returns:\ }
 \texttt{true} or \texttt{false}.



 This category contains any inverse semigroup congruence \mbox{\texttt{\mdseries\slshape cong}} which is represented internally by its kernel and trace. The \texttt{SemigroupCongruence} (\ref{SemigroupCongruence}) function may create an object of this category if its first argument \mbox{\texttt{\mdseries\slshape S}} is an inverse semigroup and has sufficiently large size. It can be treated
like any other semigroup congruence object. 

 See \cite{Howie1995aa} Section 5.3 for more details. See also \texttt{InverseSemigroupCongruenceByKernelTrace} (\ref{InverseSemigroupCongruenceByKernelTrace}). 
\begin{Verbatim}[commandchars=!@|,fontsize=\small,frame=single,label=Example]
  !gapprompt@gap>| !gapinput@S := InverseSemigroup([|
  !gapprompt@>| !gapinput@ PartialPerm([4, 3, 1, 2]),|
  !gapprompt@>| !gapinput@ PartialPerm([1, 4, 2, 0, 3])],|
  !gapprompt@>| !gapinput@ rec(cong_by_ker_trace_threshold := 0));;|
  !gapprompt@gap>| !gapinput@cong := SemigroupCongruence(S, []);|
  <semigroup congruence over <inverse partial perm semigroup
   of size 351, rank 5 with 2 generators> with congruence pair (24,24)>
  !gapprompt@gap>| !gapinput@IsInverseSemigroupCongruenceByKernelTrace(cong);|
  true
\end{Verbatim}
 }

 

\subsection{\textcolor{Chapter }{InverseSemigroupCongruenceByKernelTrace}}
\logpage{[ 13, 7, 2 ]}\nobreak
\hyperdef{L}{X7A588B737CEEB104}{}
{\noindent\textcolor{FuncColor}{$\triangleright$\enspace\texttt{InverseSemigroupCongruenceByKernelTrace({\mdseries\slshape S, kernel, traceBlocks})\index{InverseSemigroupCongruenceByKernelTrace@\texttt{Inverse}\-\texttt{Semigroup}\-\texttt{Congruence}\-\texttt{By}\-\texttt{Kernel}\-\texttt{Trace}}
\label{InverseSemigroupCongruenceByKernelTrace}
}\hfill{\scriptsize (function)}}\\
\textbf{\indent Returns:\ }
 An inverse semigroup congruence by kernel and trace. 



 If \mbox{\texttt{\mdseries\slshape S}} is an inverse semigroup, \mbox{\texttt{\mdseries\slshape kernel}} is a subsemigroup of \mbox{\texttt{\mdseries\slshape S}}, \mbox{\texttt{\mdseries\slshape traceBlocks}} is a list of lists describing a congruence on the idempotents of \mbox{\texttt{\mdseries\slshape S}}, and $(\mbox{\texttt{\mdseries\slshape kernel}}, \mbox{\texttt{\mdseries\slshape trace}})$ describes a valid congruence pair for \mbox{\texttt{\mdseries\slshape S}} (see \cite{Howie1995aa} Section 5.3) then this function returns the semigroup congruence defined by
that congruence pair. 

 See also \texttt{KernelOfSemigroupCongruence} (\ref{KernelOfSemigroupCongruence}) and \texttt{TraceOfSemigroupCongruence} (\ref{TraceOfSemigroupCongruence}). 
\begin{Verbatim}[commandchars=!@|,fontsize=\small,frame=single,label=Example]
  !gapprompt@gap>| !gapinput@S := InverseSemigroup([|
  !gapprompt@>| !gapinput@  PartialPerm([2, 3]), PartialPerm([2, 0, 3])]);;|
  !gapprompt@gap>| !gapinput@kernel := InverseSemigroup([|
  !gapprompt@>| !gapinput@  PartialPerm([1, 0, 3]), PartialPerm([0, 2, 3]),|
  !gapprompt@>| !gapinput@  PartialPerm([1, 2]), PartialPerm([3]),|
  !gapprompt@>| !gapinput@  PartialPerm([2])]);;|
  !gapprompt@gap>| !gapinput@trace := [|
  !gapprompt@>| !gapinput@ [PartialPerm([0, 2, 3])],|
  !gapprompt@>| !gapinput@ [PartialPerm([1, 2])],|
  !gapprompt@>| !gapinput@ [PartialPerm([1, 0, 3])],|
  !gapprompt@>| !gapinput@ [PartialPerm([0, 0, 3]), PartialPerm([0, 2]),|
  !gapprompt@>| !gapinput@  PartialPerm([1]), PartialPerm([], [])]];;|
  !gapprompt@gap>| !gapinput@cong := InverseSemigroupCongruenceByKernelTrace(S, kernel, trace);|
  <semigroup congruence over <inverse partial perm semigroup of rank 3
   with 2 generators> with congruence pair (13,4)>
\end{Verbatim}
 }

 

\subsection{\textcolor{Chapter }{AsInverseSemigroupCongruenceByKernelTrace}}
\logpage{[ 13, 7, 3 ]}\nobreak
\hyperdef{L}{X87916D4E854F1F6B}{}
{\noindent\textcolor{FuncColor}{$\triangleright$\enspace\texttt{AsInverseSemigroupCongruenceByKernelTrace({\mdseries\slshape cong})\index{AsInverseSemigroupCongruenceByKernelTrace@\texttt{AsInverse}\-\texttt{Semigroup}\-\texttt{Congruence}\-\texttt{By}\-\texttt{Kernel}\-\texttt{Trace}}
\label{AsInverseSemigroupCongruenceByKernelTrace}
}\hfill{\scriptsize (attribute)}}\\
\textbf{\indent Returns:\ }
 An inverse semigroup congruence by kernel and trace. 



 If \mbox{\texttt{\mdseries\slshape cong}} is a semigroup congruence over an inverse semigroup, then this attribute
returns an object which describes the same congruence, but with an internal
representation defined by that congruence's kernel and trace. 

 See \cite{Howie1995aa} section 5.3 for more details. 
\begin{Verbatim}[commandchars=!@|,fontsize=\small,frame=single,label=Example]
  !gapprompt@gap>| !gapinput@I := InverseSemigroup([|
  !gapprompt@>| !gapinput@ PartialPerm([2, 3]), PartialPerm([2, 0, 3])]);;|
  !gapprompt@gap>| !gapinput@cong := SemigroupCongruenceByGeneratingPairs(I,|
  !gapprompt@>| !gapinput@[[PartialPerm([0, 1, 3]), PartialPerm([0, 1])],|
  !gapprompt@>| !gapinput@ [PartialPerm([]), PartialPerm([1, 2])]]);|
  <2-sided semigroup congruence over <inverse partial perm semigroup of
   rank 3 with 2 generators> with 2 generating pairs>
  !gapprompt@gap>| !gapinput@cong2 := AsInverseSemigroupCongruenceByKernelTrace(cong);|
  <semigroup congruence over <inverse partial perm semigroup
   of size 19, rank 3 with 2 generators> with congruence pair (19,1)>
\end{Verbatim}
 }

 

\subsection{\textcolor{Chapter }{KernelOfSemigroupCongruence}}
\logpage{[ 13, 7, 4 ]}\nobreak
\hyperdef{L}{X7D521AFF7876CBC7}{}
{\noindent\textcolor{FuncColor}{$\triangleright$\enspace\texttt{KernelOfSemigroupCongruence({\mdseries\slshape cong})\index{KernelOfSemigroupCongruence@\texttt{KernelOfSemigroupCongruence}}
\label{KernelOfSemigroupCongruence}
}\hfill{\scriptsize (attribute)}}\\
\textbf{\indent Returns:\ }
 An inverse semigroup. 



 If \mbox{\texttt{\mdseries\slshape cong}} is a congruence over a semigroup with inverse op, then this attribute returns
the \emph{kernel} of that congruence; that is, the inverse subsemigroup consisting of all
elements which are related to an idempotent by \mbox{\texttt{\mdseries\slshape cong}}. 

 
\begin{Verbatim}[commandchars=!@|,fontsize=\small,frame=single,label=Example]
  !gapprompt@gap>| !gapinput@I := InverseSemigroup([|
  !gapprompt@>| !gapinput@ PartialPerm([2, 3]), PartialPerm([2, 0, 3])]);;|
  !gapprompt@gap>| !gapinput@cong := SemigroupCongruence(I,|
  !gapprompt@>| !gapinput@[[PartialPerm([0, 1, 3]), PartialPerm([0, 1])],|
  !gapprompt@>| !gapinput@ [PartialPerm([]), PartialPerm([1, 2])]]);|
  <2-sided semigroup congruence over <inverse partial perm semigroup
   of size 19, rank 3 with 2 generators> with 2 generating pairs>
  !gapprompt@gap>| !gapinput@KernelOfSemigroupCongruence(cong);|
  <inverse partial perm semigroup of size 19, rank 3 with 5 generators>
\end{Verbatim}
 }

 

\subsection{\textcolor{Chapter }{TraceOfSemigroupCongruence}}
\logpage{[ 13, 7, 5 ]}\nobreak
\hyperdef{L}{X80972A5B82D88F89}{}
{\noindent\textcolor{FuncColor}{$\triangleright$\enspace\texttt{TraceOfSemigroupCongruence({\mdseries\slshape cong})\index{TraceOfSemigroupCongruence@\texttt{TraceOfSemigroupCongruence}}
\label{TraceOfSemigroupCongruence}
}\hfill{\scriptsize (attribute)}}\\
\textbf{\indent Returns:\ }
A list of lists.



 If \mbox{\texttt{\mdseries\slshape cong}} is an inverse semigroup congruence by kernel and trace, then this attribute
returns the restriction of \mbox{\texttt{\mdseries\slshape cong}} to the idempotents of the semigroup. This is in block form: each idempotent
will appear in precisely one list, and two idempotents will be in the same
list if and only if they are related by \mbox{\texttt{\mdseries\slshape cong}}. 
\begin{Verbatim}[commandchars=!@|,fontsize=\small,frame=single,label=Example]
  !gapprompt@gap>| !gapinput@I := InverseSemigroup([|
  !gapprompt@>| !gapinput@ PartialPerm([2, 3]), PartialPerm([2, 0, 3])]);;|
  !gapprompt@gap>| !gapinput@cong := SemigroupCongruence(I,|
  !gapprompt@>| !gapinput@[[PartialPerm([0, 1, 3]), PartialPerm([0, 1])],|
  !gapprompt@>| !gapinput@ [PartialPerm([]), PartialPerm([1, 2])]]);|
  <2-sided semigroup congruence over <inverse partial perm semigroup
   of size 19, rank 3 with 2 generators> with 2 generating pairs>
  !gapprompt@gap>| !gapinput@TraceOfSemigroupCongruence(cong);|
  [ [ <empty partial perm>, <identity partial perm on [ 1 ]>,
        <identity partial perm on [ 2 ]>,
        <identity partial perm on [ 1, 2 ]>,
        <identity partial perm on [ 3 ]>,
        <identity partial perm on [ 2, 3 ]>,
        <identity partial perm on [ 1, 3 ]> ] ]
\end{Verbatim}
 }

 

\subsection{\textcolor{Chapter }{IsInverseSemigroupCongruenceClassByKernelTrace}}
\logpage{[ 13, 7, 6 ]}\nobreak
\hyperdef{L}{X8049A0E780A7A8D9}{}
{\noindent\textcolor{FuncColor}{$\triangleright$\enspace\texttt{IsInverseSemigroupCongruenceClassByKernelTrace({\mdseries\slshape obj})\index{IsInverseSemigroupCongruenceClassByKernelTrace@\texttt{IsInverse}\-\texttt{Semigroup}\-\texttt{Congruence}\-\texttt{Class}\-\texttt{By}\-\texttt{Kernel}\-\texttt{Trace}}
\label{IsInverseSemigroupCongruenceClassByKernelTrace}
}\hfill{\scriptsize (Category)}}\\
\textbf{\indent Returns:\ }
 \texttt{true} or \texttt{false}. 



 This category contains any congruence class which belongs to a congruence
which is represented internally by its kernel and trace. See \texttt{InverseSemigroupCongruenceByKernelTrace} (\ref{InverseSemigroupCongruenceByKernelTrace}). 

 See \cite{Howie1995aa} Section 5.3 for more details. 
\begin{Verbatim}[commandchars=!@|,fontsize=\small,frame=single,label=Example]
  !gapprompt@gap>| !gapinput@I := InverseSemigroup([|
  !gapprompt@>| !gapinput@ PartialPerm([2, 3]), PartialPerm([2, 0, 3])],|
  !gapprompt@>| !gapinput@rec(cong_by_ker_trace_threshold := 0));;|
  !gapprompt@gap>| !gapinput@cong := SemigroupCongruence(I,|
  !gapprompt@>| !gapinput@[[PartialPerm([0, 1, 3]), PartialPerm([0, 1])],|
  !gapprompt@>| !gapinput@ [PartialPerm([]), PartialPerm([1, 2])]]);;|
  !gapprompt@gap>| !gapinput@class := EquivalenceClassOfElement(cong,|
  !gapprompt@>| !gapinput@                                     PartialPerm([1, 2], [2, 3]));;|
  !gapprompt@gap>| !gapinput@IsInverseSemigroupCongruenceClassByKernelTrace(class);|
  true
\end{Verbatim}
 }

 

\subsection{\textcolor{Chapter }{MinimumGroupCongruence}}
\logpage{[ 13, 7, 7 ]}\nobreak
\hyperdef{L}{X857495647F9A9579}{}
{\noindent\textcolor{FuncColor}{$\triangleright$\enspace\texttt{MinimumGroupCongruence({\mdseries\slshape S})\index{MinimumGroupCongruence@\texttt{MinimumGroupCongruence}}
\label{MinimumGroupCongruence}
}\hfill{\scriptsize (attribute)}}\\
\textbf{\indent Returns:\ }
 An inverse semigroup congruence by kernel and trace. 



 If \mbox{\texttt{\mdseries\slshape S}} is an inverse semigroup, then this function returns the least congruence on \mbox{\texttt{\mdseries\slshape S}} whose quotient is a group. 

 
\begin{Verbatim}[commandchars=!@|,fontsize=\small,frame=single,label=Example]
  !gapprompt@gap>| !gapinput@S := InverseSemigroup([|
  !gapprompt@>| !gapinput@  PartialPerm([5, 2, 0, 0, 1, 4]),|
  !gapprompt@>| !gapinput@  PartialPerm([1, 4, 6, 3, 5, 0, 2])]);;|
  !gapprompt@gap>| !gapinput@cong := MinimumGroupCongruence(S);|
  <semigroup congruence over <inverse partial perm semigroup
   of size 101, rank 7 with 2 generators> with congruence pair (59,1)>
  !gapprompt@gap>| !gapinput@IsGroupAsSemigroup(S / cong);|
  true
\end{Verbatim}
 }

 }

 
\section{\textcolor{Chapter }{ Congruences on graph inverse semigroups }}\logpage{[ 13, 8, 0 ]}
\hyperdef{L}{X8036DAA287C71CAC}{}
{
  

\subsection{\textcolor{Chapter }{IsCongruenceByWangPair}}
\logpage{[ 13, 8, 1 ]}\nobreak
\hyperdef{L}{X7AEB7DA27E76145B}{}
{\noindent\textcolor{FuncColor}{$\triangleright$\enspace\texttt{IsCongruenceByWangPair({\mdseries\slshape cong})\index{IsCongruenceByWangPair@\texttt{IsCongruenceByWangPair}}
\label{IsCongruenceByWangPair}
}\hfill{\scriptsize (property)}}\\


 A congruence by Wang pair \texttt{cong} is a congruence of a graph inverse semigroup \texttt{S} which is expressed in terms of two sets \mbox{\texttt{\mdseries\slshape H}} and \mbox{\texttt{\mdseries\slshape W}} of vertices of the corresponding graph of \mbox{\texttt{\mdseries\slshape S}}. The set \mbox{\texttt{\mdseries\slshape H}} must be a hereditary subset (closed under reachability) and all vertices in \mbox{\texttt{\mdseries\slshape W}} must have all but one of their out\texttt{\symbol{45}}neighbours in \mbox{\texttt{\mdseries\slshape H}}. For more information on Wang pairs see \cite{Wang2019aa} and \cite{Anagnostopoulou-Merkouri2021aa}.

 
\begin{Verbatim}[commandchars=!@|,fontsize=\small,frame=single,label=Example]
  !gapprompt@gap>| !gapinput@D := Digraph([[3, 4], [3, 4], [4], []]);|
  <immutable digraph with 4 vertices, 5 edges>
  !gapprompt@gap>| !gapinput@S := GraphInverseSemigroup(D);|
  <finite graph inverse semigroup with 4 vertices, 5 edges>
  !gapprompt@gap>| !gapinput@cong := CongruenceByWangPair(S, [3, 4], []);|
  <graph inverse semigroup congruence with H = [ 3, 4 ] and W = [  ]>
  !gapprompt@gap>| !gapinput@IsCongruenceByWangPair(cong);|
  true
  !gapprompt@gap>| !gapinput@cong := CongruenceByWangPair(S, [4], [2]);|
  <graph inverse semigroup congruence with H = [ 4 ] and W = [ 2 ]>
  !gapprompt@gap>| !gapinput@IsCongruenceByWangPair(cong);|
  true
  !gapprompt@gap>| !gapinput@e_1 := S.1;|
  e_1
  !gapprompt@gap>| !gapinput@e_3 := S.3;|
  e_3
  !gapprompt@gap>| !gapinput@cong := SemigroupCongruence(S, [[e_1, e_3]]);|
  <2-sided semigroup congruence over <finite graph inverse semigroup wit\
  h 4 vertices, 5 edges> with 1 generating pairs>
  !gapprompt@gap>| !gapinput@IsCongruenceByWangPair(cong);|
  false
\end{Verbatim}
 }

 

\subsection{\textcolor{Chapter }{CongruenceByWangPair}}
\logpage{[ 13, 8, 2 ]}\nobreak
\hyperdef{L}{X7F30D10F7BEEEBB9}{}
{\noindent\textcolor{FuncColor}{$\triangleright$\enspace\texttt{CongruenceByWangPair({\mdseries\slshape S, H, W})\index{CongruenceByWangPair@\texttt{CongruenceByWangPair}}
\label{CongruenceByWangPair}
}\hfill{\scriptsize (function)}}\\
\textbf{\indent Returns:\ }
A semigroup congruence.



 This function returns a semigroup congruence over the graph inverse semigroup \mbox{\texttt{\mdseries\slshape S}} in the form of a Wang pair.

 If \mbox{\texttt{\mdseries\slshape S}} is a finite graph inverse semigroup \mbox{\texttt{\mdseries\slshape H}} and \mbox{\texttt{\mdseries\slshape W}} are two lists of vertices in the graph of \mbox{\texttt{\mdseries\slshape S}} representing a valid hereditary subset and a W\texttt{\symbol{45}}set
respectively, then this function will return the semigroup congruence defined
by this Wang pair. For the definition of Wang pair \texttt{IsCongruenceByWangPair} (\ref{IsCongruenceByWangPair}).

 
\begin{Verbatim}[commandchars=!@|,fontsize=\small,frame=single,label=Example]
  !gapprompt@gap>| !gapinput@D := Digraph([[3, 4], [3, 4], [4], []]);|
  <immutable digraph with 4 vertices, 5 edges>
  !gapprompt@gap>| !gapinput@S := GraphInverseSemigroup(D);|
  <finite graph inverse semigroup with 4 vertices, 5 edges>
  !gapprompt@gap>| !gapinput@cong := CongruenceByWangPair(S, [3, 4], []);|
  <graph inverse semigroup congruence with H = [ 3, 4 ] and W = [  ]>
  !gapprompt@gap>| !gapinput@cong := CongruenceByWangPair(S, [4], [2]);|
  <graph inverse semigroup congruence with H = [ 4 ] and W = [ 2 ]>
  !gapprompt@gap>| !gapinput@cong := CongruenceByWangPair(S, [3, 4], []);|
  <graph inverse semigroup congruence with H = [ 3, 4 ] and W = [  ]>
\end{Verbatim}
 }

 

\subsection{\textcolor{Chapter }{AsCongruenceByWangPair}}
\logpage{[ 13, 8, 3 ]}\nobreak
\hyperdef{L}{X817F4FC27E9BACD8}{}
{\noindent\textcolor{FuncColor}{$\triangleright$\enspace\texttt{AsCongruenceByWangPair({\mdseries\slshape cong})\index{AsCongruenceByWangPair@\texttt{AsCongruenceByWangPair}}
\label{AsCongruenceByWangPair}
}\hfill{\scriptsize (operation)}}\\
\textbf{\indent Returns:\ }
A congruence by Wang pair.



 This operation takes \mbox{\texttt{\mdseries\slshape cong}}, a finite graph inverse semigroup congruence, and returns an object
representing the same congruence, but described as a congruence by Wang pairs:
a pair of sets \mbox{\texttt{\mdseries\slshape H}} and \mbox{\texttt{\mdseries\slshape W}} of the corresponding graph of \mbox{\texttt{\mdseries\slshape S}} that are a hereditary subset and a W\texttt{\symbol{45}}set of the graph of \mbox{\texttt{\mdseries\slshape S}} respectively. For more information about Wang pairs see \cite{Wang2019aa} and \cite{Anagnostopoulou-Merkouri2021aa}. 

 
\begin{Verbatim}[commandchars=!@|,fontsize=\small,frame=single,label=Example]
  !gapprompt@gap>| !gapinput@D := Digraph([[2, 3], [3], [4], []]);|
  <immutable digraph with 4 vertices, 4 edges>
  !gapprompt@gap>| !gapinput@S := GraphInverseSemigroup(D);|
  <finite graph inverse semigroup with 4 vertices, 4 edges>
  !gapprompt@gap>| !gapinput@CongruenceByWangPair(S, [4], [2]);|
  <graph inverse semigroup congruence with H = [ 4 ] and W = [ 2 ]>
  !gapprompt@gap>| !gapinput@cong := AsSemigroupCongruenceByGeneratingPairs(last);|
  <2-sided semigroup congruence over <finite graph inverse semigroup wit\
  h 4 vertices, 4 edges> with 2 generating pairs>
  !gapprompt@gap>| !gapinput@AsCongruenceByWangPair(cong);|
  <graph inverse semigroup congruence with H = [ 4 ] and W = [ 2 ]>
  !gapprompt@gap>| !gapinput@CongruenceByWangPair(S, [3, 4], [1]);|
  <graph inverse semigroup congruence with H = [ 3, 4 ] and W = [ 1 ]>
  !gapprompt@gap>| !gapinput@cong := AsSemigroupCongruenceByGeneratingPairs(last);|
  <2-sided semigroup congruence over <finite graph inverse semigroup wit\
  h 4 vertices, 4 edges> with 3 generating pairs>
  !gapprompt@gap>| !gapinput@AsCongruenceByWangPair(cong);|
  <graph inverse semigroup congruence with H = [ 3, 4 ] and W = [ 1 ]>
\end{Verbatim}
 }

 

\subsection{\textcolor{Chapter }{GeneratingCongruencesOfLattice}}
\logpage{[ 13, 8, 4 ]}\nobreak
\hyperdef{L}{X858AE13379B5C380}{}
{\noindent\textcolor{FuncColor}{$\triangleright$\enspace\texttt{GeneratingCongruencesOfLattice({\mdseries\slshape S})\index{GeneratingCongruencesOfLattice@\texttt{GeneratingCongruencesOfLattice}}
\label{GeneratingCongruencesOfLattice}
}\hfill{\scriptsize (attribute)}}\\
\textbf{\indent Returns:\ }
A semigroup.



 This attribute takes a finite graph inverse semigroup \mbox{\texttt{\mdseries\slshape S}} and returns a minimal generating set for the lattice of congruences of \mbox{\texttt{\mdseries\slshape S}}, as described in \cite{Anagnostopoulou-Merkouri2021aa}. This operation works only if the corresponding digraph of the graph inverse
semigroup is simple. If there are multiple edges, an error is returned.

 
\begin{Verbatim}[commandchars=!@|,fontsize=\small,frame=single,label=Example]
  !gapprompt@gap>| !gapinput@D := Digraph([[2, 3], [3], [4], []]);|
  <immutable digraph with 4 vertices, 4 edges>
  !gapprompt@gap>| !gapinput@S := GraphInverseSemigroup(D);|
  <finite graph inverse semigroup with 4 vertices, 4 edges>
  !gapprompt@gap>| !gapinput@CongruenceByWangPair(S, [4], [2]);|
  <graph inverse semigroup congruence with H = [ 4 ] and W = [ 2 ]>
  !gapprompt@gap>| !gapinput@cong := AsSemigroupCongruenceByGeneratingPairs(last);|
  <2-sided semigroup congruence over <finite graph inverse semigroup wit\
  h 4 vertices, 4 edges> with 2 generating pairs>
  !gapprompt@gap>| !gapinput@AsCongruenceByWangPair(cong);|
  <graph inverse semigroup congruence with H = [ 4 ] and W = [ 2 ]>
  !gapprompt@gap>| !gapinput@CongruenceByWangPair(S, [3, 4], [1]);|
  <graph inverse semigroup congruence with H = [ 3, 4 ] and W = [ 1 ]>
  !gapprompt@gap>| !gapinput@cong := AsSemigroupCongruenceByGeneratingPairs(last);|
  <2-sided semigroup congruence over <finite graph inverse semigroup wit\
  h 4 vertices, 4 edges> with 3 generating pairs>
  !gapprompt@gap>| !gapinput@AsCongruenceByWangPair(cong);|
  <graph inverse semigroup congruence with H = [ 3, 4 ] and W = [ 1 ]>
\end{Verbatim}
 }

 }

 
\section{\textcolor{Chapter }{Rees congruences}}\logpage{[ 13, 9, 0 ]}
\hyperdef{L}{X7CE483078769A4D6}{}
{
  A Rees congruence is defined by a semigroup ideal. It is a congruence on a
semigroup \texttt{S} which has one congruence class equal to a semigroup ideal \texttt{I} of \texttt{S}, and every other congruence class being a singleton. 

 

\subsection{\textcolor{Chapter }{SemigroupIdealOfReesCongruence}}
\logpage{[ 13, 9, 1 ]}\nobreak
\hyperdef{L}{X7BC02E08840E0AF4}{}
{\noindent\textcolor{FuncColor}{$\triangleright$\enspace\texttt{SemigroupIdealOfReesCongruence({\mdseries\slshape cong})\index{SemigroupIdealOfReesCongruence@\texttt{SemigroupIdealOfReesCongruence}}
\label{SemigroupIdealOfReesCongruence}
}\hfill{\scriptsize (attribute)}}\\
\textbf{\indent Returns:\ }
A semigroup ideal.



 If \mbox{\texttt{\mdseries\slshape cong}} is a rees congruence (see \texttt{IsReesCongruence} (\textbf{Reference: IsReesCongruence})) then this attribute returns the two\texttt{\symbol{45}}sided ideal that was
used to define it, i.e.\texttt{\symbol{126}}the ideal of elements in the only
non\texttt{\symbol{45}}trivial congruence class of \mbox{\texttt{\mdseries\slshape cong}}. 
\begin{Verbatim}[commandchars=!@|,fontsize=\small,frame=single,label=Example]
  !gapprompt@gap>| !gapinput@S := Semigroup([|
  !gapprompt@>| !gapinput@Transformation([2, 3, 4, 3, 1, 1]),|
  !gapprompt@>| !gapinput@Transformation([6, 4, 4, 4, 6, 1])]);;|
  !gapprompt@gap>| !gapinput@I := SemigroupIdeal(S,|
  !gapprompt@>| !gapinput@Transformation([4, 4, 4, 4, 4, 2]),|
  !gapprompt@>| !gapinput@Transformation([3, 3, 3, 3, 3, 2]));;|
  !gapprompt@gap>| !gapinput@cong := ReesCongruenceOfSemigroupIdeal(I);;|
  !gapprompt@gap>| !gapinput@SemigroupIdealOfReesCongruence(cong);|
  <non-regular transformation semigroup ideal of degree 6 with
    2 generators>
\end{Verbatim}
 }

 

\subsection{\textcolor{Chapter }{IsReesCongruenceClass}}
\logpage{[ 13, 9, 2 ]}\nobreak
\hyperdef{L}{X7E15F66A8029C64A}{}
{\noindent\textcolor{FuncColor}{$\triangleright$\enspace\texttt{IsReesCongruenceClass({\mdseries\slshape obj})\index{IsReesCongruenceClass@\texttt{IsReesCongruenceClass}}
\label{IsReesCongruenceClass}
}\hfill{\scriptsize (category)}}\\
\textbf{\indent Returns:\ }
\texttt{true} or \texttt{false}.



 This category describes a congruence class of a Rees congruence. A congruence
class of a Rees congruence either contains all the elements of an ideal, or is
a singleton (see \texttt{IsReesCongruence} (\textbf{Reference: IsReesCongruence})).

 An object of this type may be used in the same way as any other congruence
class object. 
\begin{Verbatim}[commandchars=!@|,fontsize=\small,frame=single,label=Example]
  !gapprompt@gap>| !gapinput@S := Semigroup(|
  !gapprompt@>| !gapinput@Transformation([2, 3, 4, 3, 1, 1]),|
  !gapprompt@>| !gapinput@Transformation([6, 4, 4, 4, 6, 1]));;|
  !gapprompt@gap>| !gapinput@I := SemigroupIdeal(S,|
  !gapprompt@>| !gapinput@Transformation([4, 4, 4, 4, 4, 2]),|
  !gapprompt@>| !gapinput@Transformation([3, 3, 3, 3, 3, 2]));;|
  !gapprompt@gap>| !gapinput@cong := ReesCongruenceOfSemigroupIdeal(I);;|
  !gapprompt@gap>| !gapinput@classes := EquivalenceClasses(cong);;|
  !gapprompt@gap>| !gapinput@IsReesCongruenceClass(classes[1]);|
  true
\end{Verbatim}
 }

 }

 
\section{\textcolor{Chapter }{Universal and trivial congruences}}\logpage{[ 13, 10, 0 ]}
\hyperdef{L}{X7C6B4A2980BE9B03}{}
{
  The linked triples of a completely 0\texttt{\symbol{45}}simple Rees
0\texttt{\symbol{45}}matrix semigroup describe only its
non\texttt{\symbol{45}}universal congruences. In any one of these, the zero
element of the semigroup is related only to itself. However, for any semigroup $S$ the universal relation $S \times  S$ is a congruence; called the \emph{universal congruence}. The universal congruence on a semigroup has its own unique representation.

 Since many things we want to calculate about congruences are trivial in the
case of the universal congruence, this package contains a category
specifically designed for it, \texttt{IsUniversalSemigroupCongruence}. We also define \texttt{IsUniversalSemigroupCongruenceClass}, which represents the single congruence class of the universal congruence.

 

\subsection{\textcolor{Chapter }{IsUniversalSemigroupCongruence}}
\logpage{[ 13, 10, 1 ]}\nobreak
\hyperdef{L}{X8751EF557A81A2B1}{}
{\noindent\textcolor{FuncColor}{$\triangleright$\enspace\texttt{IsUniversalSemigroupCongruence({\mdseries\slshape obj})\index{IsUniversalSemigroupCongruence@\texttt{IsUniversalSemigroupCongruence}}
\label{IsUniversalSemigroupCongruence}
}\hfill{\scriptsize (property)}}\\
\textbf{\indent Returns:\ }
\texttt{true} or \texttt{false}.



 This property describes a type of semigroup congruence, which must refer to
the \emph{universal semigroup congruence} $S \times S$. Externally, an object of this type may be used in the same way as any other
object in the category \texttt{IsSemigroupCongruence} (\textbf{Reference: IsSemigroupCongruence}).

 An object of this type may be constructed with \texttt{UniversalSemigroupCongruence} or this representation may be selected automatically as an alternative to an \texttt{IsRZMSCongruenceByLinkedTriple} object (since the universal congruence cannot be represented by a linked
triple). 
\begin{Verbatim}[commandchars=!@|,fontsize=\small,frame=single,label=Example]
  !gapprompt@gap>| !gapinput@S := Semigroup([Transformation([3, 2, 3])]);;|
  !gapprompt@gap>| !gapinput@U := UniversalSemigroupCongruence(S);;|
  !gapprompt@gap>| !gapinput@IsUniversalSemigroupCongruence(U);|
  true
\end{Verbatim}
 }

 

\subsection{\textcolor{Chapter }{IsUniversalSemigroupCongruenceClass}}
\logpage{[ 13, 10, 2 ]}\nobreak
\hyperdef{L}{X8646253C86AFFA29}{}
{\noindent\textcolor{FuncColor}{$\triangleright$\enspace\texttt{IsUniversalSemigroupCongruenceClass({\mdseries\slshape obj})\index{IsUniversalSemigroupCongruenceClass@\texttt{IsUniversalSemigroupCongruenceClass}}
\label{IsUniversalSemigroupCongruenceClass}
}\hfill{\scriptsize (category)}}\\
\textbf{\indent Returns:\ }
\texttt{true} or \texttt{false}.



 This category describes a class of the universal semigroup congruence (see \texttt{IsUniversalSemigroupCongruence} (\ref{IsUniversalSemigroupCongruence})). A universal semigroup congruence by definition has precisely one congruence
class, which contains all of the elements of the semigroup in question. 
\begin{Verbatim}[commandchars=!@|,fontsize=\small,frame=single,label=Example]
  !gapprompt@gap>| !gapinput@S := Semigroup([Transformation([3, 2, 3])]);;|
  !gapprompt@gap>| !gapinput@U := UniversalSemigroupCongruence(S);;|
  !gapprompt@gap>| !gapinput@classes := EquivalenceClasses(U);;|
  !gapprompt@gap>| !gapinput@IsUniversalSemigroupCongruenceClass(classes[1]);|
  true
\end{Verbatim}
 }

 

\subsection{\textcolor{Chapter }{UniversalSemigroupCongruence}}
\logpage{[ 13, 10, 3 ]}\nobreak
\hyperdef{L}{X7B99C6A47F3D375F}{}
{\noindent\textcolor{FuncColor}{$\triangleright$\enspace\texttt{UniversalSemigroupCongruence({\mdseries\slshape S})\index{UniversalSemigroupCongruence@\texttt{UniversalSemigroupCongruence}}
\label{UniversalSemigroupCongruence}
}\hfill{\scriptsize (operation)}}\\
\textbf{\indent Returns:\ }
A universal semigroup congruence.



 This operation returns the universal semigroup congruence for the semigroup \mbox{\texttt{\mdseries\slshape S}}. It can be used in the same way as any other semigroup congruence object. 
\begin{Verbatim}[commandchars=!@|,fontsize=\small,frame=single,label=Example]
  !gapprompt@gap>| !gapinput@S := ReesZeroMatrixSemigroup(SymmetricGroup(3),|
  !gapprompt@>| !gapinput@[[(), (1, 3, 2)], [(1, 2), 0]]);;|
  !gapprompt@gap>| !gapinput@UniversalSemigroupCongruence(S);|
  <universal semigroup congruence over
  <Rees 0-matrix semigroup 2x2 over Sym( [ 1 .. 3 ] )>>
\end{Verbatim}
 }

 

\subsection{\textcolor{Chapter }{TrivialCongruence}}
\logpage{[ 13, 10, 4 ]}\nobreak
\hyperdef{L}{X7A9B483B7B4E0E27}{}
{\noindent\textcolor{FuncColor}{$\triangleright$\enspace\texttt{TrivialCongruence({\mdseries\slshape S})\index{TrivialCongruence@\texttt{TrivialCongruence}}
\label{TrivialCongruence}
}\hfill{\scriptsize (attribute)}}\\
\textbf{\indent Returns:\ }
A trivial semigroup congruence.



 This operation returns the trivial semigroup congruence for the semigroup \mbox{\texttt{\mdseries\slshape S}}. It can be used in the same way as any other semigroup congruence object. 
\begin{Verbatim}[commandchars=!@|,fontsize=\small,frame=single,label=Example]
  !gapprompt@gap>| !gapinput@S := ReesZeroMatrixSemigroup(SymmetricGroup(3),|
  !gapprompt@>| !gapinput@[[(), (1, 3, 2)], [(1, 2), 0]]);;|
  !gapprompt@gap>| !gapinput@TrivialCongruence(S);|
  <semigroup congruence over <Rees 0-matrix semigroup 2x2 over
    Sym( [ 1 .. 3 ] )> with linked triple (1,2,2)>
  !gapprompt@gap>| !gapinput@S := PartitionMonoid(2);|
  <regular bipartition *-monoid of size 15, degree 2 with 3 generators>
  !gapprompt@gap>| !gapinput@TrivialCongruence(S);|
  <2-sided semigroup congruence over <regular bipartition *-monoid
   of size 15, degree 2 with 3 generators> with 0 generating pairs>
\end{Verbatim}
 }

 }

 }

 
\chapter{\textcolor{Chapter }{ Semigroup homomorphisms }}\label{Semigroup homomorphisms}
\logpage{[ 14, 0, 0 ]}
\hyperdef{L}{X861935DB81A478C2}{}
{
  In this chapter we describe the various ways to define a homomorphism from a
semigroup to another semigroup. 

   
\section{\textcolor{Chapter }{ Homomorphisms of arbitrary semigroups }}\logpage{[ 14, 1, 0 ]}
\hyperdef{L}{X7F1FDA9C7C25799A}{}
{
  

\subsection{\textcolor{Chapter }{SemigroupHomomorphismByImages (for two semigroups and two lists)}}
\logpage{[ 14, 1, 1 ]}\nobreak
\hyperdef{L}{X817596438369885B}{}
{\noindent\textcolor{FuncColor}{$\triangleright$\enspace\texttt{SemigroupHomomorphismByImages({\mdseries\slshape S, T, gens, imgs})\index{SemigroupHomomorphismByImages@\texttt{SemigroupHomomorphismByImages}!for two semigroups and two lists}
\label{SemigroupHomomorphismByImages:for two semigroups and two lists}
}\hfill{\scriptsize (operation)}}\\
\noindent\textcolor{FuncColor}{$\triangleright$\enspace\texttt{SemigroupHomomorphismByImages({\mdseries\slshape S, T, imgs})\index{SemigroupHomomorphismByImages@\texttt{SemigroupHomomorphismByImages}!for two semigroups and a list}
\label{SemigroupHomomorphismByImages:for two semigroups and a list}
}\hfill{\scriptsize (operation)}}\\
\noindent\textcolor{FuncColor}{$\triangleright$\enspace\texttt{SemigroupHomomorphismByImages({\mdseries\slshape S, T})\index{SemigroupHomomorphismByImages@\texttt{SemigroupHomomorphismByImages}!for two semigroups}
\label{SemigroupHomomorphismByImages:for two semigroups}
}\hfill{\scriptsize (operation)}}\\
\noindent\textcolor{FuncColor}{$\triangleright$\enspace\texttt{SemigroupHomomorphismByImages({\mdseries\slshape S, gens, imgs})\index{SemigroupHomomorphismByImages@\texttt{SemigroupHomomorphismByImages}!for a semigroup and two lists}
\label{SemigroupHomomorphismByImages:for a semigroup and two lists}
}\hfill{\scriptsize (operation)}}\\
\textbf{\indent Returns:\ }
 A semigroup homomorphism, or \texttt{fail}. 



 \texttt{SemigroupHomomorphismByImages} attempts to construct a homomorphism from the semigroup \mbox{\texttt{\mdseries\slshape S}} to the semigroup \mbox{\texttt{\mdseries\slshape T}} by mapping the \texttt{i}\texttt{\symbol{45}}th element of \mbox{\texttt{\mdseries\slshape gens}} to the \texttt{i}\texttt{\symbol{45}}th element of \mbox{\texttt{\mdseries\slshape imgs}}. If this mapping corresponds to a homomorphism, the homomorphism is returned,
and if not, then \texttt{fail} is returned. Similarly, if \mbox{\texttt{\mdseries\slshape gens}} does not generate \mbox{\texttt{\mdseries\slshape S}}, \texttt{fail} is returned.

 If omitted, the arguments \mbox{\texttt{\mdseries\slshape gens}} and \mbox{\texttt{\mdseries\slshape imgs}} default to the generators of \mbox{\texttt{\mdseries\slshape S}} and \mbox{\texttt{\mdseries\slshape T}} respectively. See \texttt{GeneratorsOfSemigroup} (\textbf{Reference: GeneratorsOfSemigroup}).

 If \mbox{\texttt{\mdseries\slshape T}} is not given, then it defaults to the semigroup generated by \mbox{\texttt{\mdseries\slshape imgs}}, resulting in the mapping being surjective. 
\begin{Verbatim}[commandchars=!@|,fontsize=\small,frame=single,label=Example]
  !gapprompt@gap>| !gapinput@S := FullTransformationMonoid(3);;|
  !gapprompt@gap>| !gapinput@gens := GeneratorsOfSemigroup(S);;|
  !gapprompt@gap>| !gapinput@J := FullTransformationMonoid(4);;|
  !gapprompt@gap>| !gapinput@imgs := ListWithIdenticalEntries(4,|
  !gapprompt@>| !gapinput@ConstantTransformation(3, 1));;|
  !gapprompt@gap>| !gapinput@hom := SemigroupHomomorphismByImages(S, J, gens, imgs);|
  <full transformation monoid of degree 3> ->
  <full transformation monoid of degree 4>
\end{Verbatim}
 }

 

\subsection{\textcolor{Chapter }{SemigroupHomomorphismByFunctionNC}}
\logpage{[ 14, 1, 2 ]}\nobreak
\hyperdef{L}{X7DAA6AD985C22AD6}{}
{\noindent\textcolor{FuncColor}{$\triangleright$\enspace\texttt{SemigroupHomomorphismByFunctionNC({\mdseries\slshape S, T, fun})\index{SemigroupHomomorphismByFunctionNC@\texttt{SemigroupHomomorphismByFunctionNC}}
\label{SemigroupHomomorphismByFunctionNC}
}\hfill{\scriptsize (operation)}}\\
\noindent\textcolor{FuncColor}{$\triangleright$\enspace\texttt{SemigroupHomomorphismByFunction({\mdseries\slshape S, T, fun})\index{SemigroupHomomorphismByFunction@\texttt{SemigroupHomomorphismByFunction}}
\label{SemigroupHomomorphismByFunction}
}\hfill{\scriptsize (operation)}}\\
\textbf{\indent Returns:\ }
 A semigroup homomorphism or \texttt{fail}. 



 \texttt{SemigroupHomomorphismByFunctionNC} returns a semigroup homomorphism with source \mbox{\texttt{\mdseries\slshape S}} and range \mbox{\texttt{\mdseries\slshape T}}, such that each element \texttt{s} in \mbox{\texttt{\mdseries\slshape S}} is mapped to the element \mbox{\texttt{\mdseries\slshape fun}}\texttt{(s)}, where \mbox{\texttt{\mdseries\slshape fun}} is a \textsf{GAP} function.

 The function \texttt{SemigroupHomomorphismByFunctionNC} performs no checks on whether the function actually gives a homomorphism, and
so it is possible for this operation to return a mapping from \mbox{\texttt{\mdseries\slshape S}} to \mbox{\texttt{\mdseries\slshape T}} that is not a homomorphism.

 The function \texttt{SemigroupHomomorphismByFunction} checks that the mapping from \mbox{\texttt{\mdseries\slshape S}} to \mbox{\texttt{\mdseries\slshape T}} defined by \mbox{\texttt{\mdseries\slshape fun}} satisfies \texttt{RespectsMultiplication} (\textbf{Reference: RespectsMultiplication}), which can be expensive. If \texttt{RespectsMultiplication} (\textbf{Reference: RespectsMultiplication}) does not hold, then \texttt{fail} is returned. 
\begin{Verbatim}[commandchars=!@|,fontsize=\small,frame=single,label=Example]
  !gapprompt@gap>| !gapinput@g := Semigroup([(1, 2, 3, 4), (1, 2)]);;|
  !gapprompt@gap>| !gapinput@h := Semigroup([(1, 2, 3), (1, 2)]);;|
  !gapprompt@gap>| !gapinput@hom := SemigroupHomomorphismByFunction(g, h,|
  !gapprompt@>| !gapinput@function(x)|
  !gapprompt@>| !gapinput@if SignPerm(x) = -1 then return (1, 2);|
  !gapprompt@>| !gapinput@else return ();|
  !gapprompt@>| !gapinput@fi; end);|
  <semigroup of size 24, with 2 generators> ->
  <semigroup of size 6, with 2 generators>
\end{Verbatim}
 }

 The following methods relate to semigroup homomorphisms by images or by
function: 
\begin{itemize}
\item  \texttt{Range} (\textbf{Reference: range}), 
\item  \texttt{Image} (\textbf{Reference: Image}), 
\item  \texttt{Images} (\textbf{Reference: Images}), 
\item  \texttt{ImageElm} (\textbf{Reference: ImageElm}), 
\item  \texttt{PreImage} (\textbf{Reference: PreImage}), 
\item  \texttt{PreImages} (\textbf{Reference: PreImages}), 
\item  \texttt{PreImagesRepresentative} (\textbf{Reference: PreImagesRepresentative}), 
\item  \texttt{PreImagesRange} (\textbf{Reference: PreImagesRange}), 
\item  \texttt{PreImagesElm} (\textbf{Reference: PreImagesElm}), 
\item  \texttt{PreImagesSet} (\textbf{Reference: PreImagesSet}), 
\item  \texttt{IsSurjective} (\textbf{Reference: IsSurjective}), 
\item  \texttt{IsInjective} (\textbf{Reference: IsInjective}), 
\item  \texttt{IsBijective} (\textbf{Reference: IsBijective}), 
\item  \texttt{Source} (\textbf{Reference: Source}), 
\item  \texttt{Range} (\textbf{Reference: range}), 
\item  \texttt{ImagesSource} (\textbf{Reference: ImagesSource}), 
\item  \texttt{KernelOfSemigroupHomomorphism} (\ref{KernelOfSemigroupHomomorphism}). 
\end{itemize}
 

\subsection{\textcolor{Chapter }{IsSemigroupHomomorphismByImages}}
\logpage{[ 14, 1, 3 ]}\nobreak
\hyperdef{L}{X7C76C6E5780D4A57}{}
{\noindent\textcolor{FuncColor}{$\triangleright$\enspace\texttt{IsSemigroupHomomorphismByImages({\mdseries\slshape hom})\index{IsSemigroupHomomorphismByImages@\texttt{IsSemigroupHomomorphismByImages}}
\label{IsSemigroupHomomorphismByImages}
}\hfill{\scriptsize (filter)}}\\
\textbf{\indent Returns:\ }
 \texttt{true} or \texttt{false}. 



 \texttt{IsSemigroupHomomorphismByImages} returns \texttt{true} if \mbox{\texttt{\mdseries\slshape hom}} is a semigroup homomorphism by images and \texttt{false} if it is not. A semigroup homomorphism is a mapping from a semigroup \texttt{S} to a semigroup \texttt{T} that respects multiplication. This representation describes semigroup
homomorphisms internally by the generators of \texttt{S} and their images in \texttt{T}. See \texttt{SemigroupHomomorphismByImages} (\ref{SemigroupHomomorphismByImages:for two semigroups and two lists}). 
\begin{Verbatim}[commandchars=!@|,fontsize=\small,frame=single,label=Example]
  !gapprompt@gap>| !gapinput@S := FullTransformationMonoid(3);;|
  !gapprompt@gap>| !gapinput@gens := GeneratorsOfSemigroup(S);;|
  !gapprompt@gap>| !gapinput@T := FullTransformationMonoid(4);;|
  !gapprompt@gap>| !gapinput@imgs := ListWithIdenticalEntries(4, ConstantTransformation(3, 1));;|
  !gapprompt@gap>| !gapinput@hom := SemigroupHomomorphismByImages(S, T, gens, imgs);|
  <full transformation monoid of degree 3> ->
  <full transformation monoid of degree 4>
  !gapprompt@gap>| !gapinput@IsSemigroupHomomorphismByImages(hom);|
  true
\end{Verbatim}
 }

 

\subsection{\textcolor{Chapter }{IsSemigroupHomomorphismByFunction}}
\logpage{[ 14, 1, 4 ]}\nobreak
\hyperdef{L}{X7F9CF9457E84BAE2}{}
{\noindent\textcolor{FuncColor}{$\triangleright$\enspace\texttt{IsSemigroupHomomorphismByFunction({\mdseries\slshape hom})\index{IsSemigroupHomomorphismByFunction@\texttt{IsSemigroupHomomorphismByFunction}}
\label{IsSemigroupHomomorphismByFunction}
}\hfill{\scriptsize (filter)}}\\
\textbf{\indent Returns:\ }
 \texttt{true} or \texttt{false}. 



 \texttt{IsSemigroupHomomorphismByFunction} returns \texttt{true} if \mbox{\texttt{\mdseries\slshape hom}} was created using \texttt{SemigroupHomomorphismByFunction} (\ref{SemigroupHomomorphismByFunction}) and \texttt{false} if it was not. Note that this filter may return \texttt{true} even if the underlying \textsf{GAP} function does not define a homomorphism. A semigroup homomorphism is a mapping
from a semigroup \texttt{S} to a semigroup \texttt{T} that respects multiplication. This representation describes semigroup
homomorphisms internally using a \textsf{GAP} function mapping elements of \texttt{S} to their images in \texttt{T}. 
\begin{Verbatim}[commandchars=!@|,fontsize=\small,frame=single,label=Example]
  !gapprompt@gap>| !gapinput@S := Semigroup([(1, 2, 3, 4), (1, 2)]);;|
  !gapprompt@gap>| !gapinput@T := Semigroup([(1, 2, 3), (1, 2)]);;|
  !gapprompt@gap>| !gapinput@hom := SemigroupHomomorphismByFunction(S, T,|
  !gapprompt@>| !gapinput@function(x) if SignPerm(x) = -1 then return (1, 2);|
  !gapprompt@>| !gapinput@else return ();fi; end);|
  <semigroup of size 24, with 2 generators> ->
  <semigroup of size 6, with 2 generators>
  !gapprompt@gap>| !gapinput@IsSemigroupHomomorphismByFunction(hom);|
  true
\end{Verbatim}
 }

 

\subsection{\textcolor{Chapter }{AsSemigroupHomomorphismByImages (for a semigroup homomorphism by function)}}
\logpage{[ 14, 1, 5 ]}\nobreak
\hyperdef{L}{X7CEBDC767CC184B6}{}
{\noindent\textcolor{FuncColor}{$\triangleright$\enspace\texttt{AsSemigroupHomomorphismByImages({\mdseries\slshape hom})\index{AsSemigroupHomomorphismByImages@\texttt{AsSemigroupHomomorphismByImages}!for a semigroup homomorphism by function}
\label{AsSemigroupHomomorphismByImages:for a semigroup homomorphism by function}
}\hfill{\scriptsize (operation)}}\\
\textbf{\indent Returns:\ }
 A semigroup homomorphism, or \texttt{fail}. 



 \texttt{AsSemigroupHomomorphismByImages} takes \mbox{\texttt{\mdseries\slshape hom}}, a semigroup homomorphism, and returns the same mapping but represented
internally using the generators of \texttt{Source(\mbox{\texttt{\mdseries\slshape hom}})} and their images in \texttt{Range(\mbox{\texttt{\mdseries\slshape hom}})}. If \mbox{\texttt{\mdseries\slshape hom}} not a semigroup homomorphism, then \texttt{fail} is returned. For example, this could happen if \mbox{\texttt{\mdseries\slshape hom}} was created using \texttt{SemigroupIsomorphismByFunction} (\ref{SemigroupIsomorphismByFunction}) and a function which does not give a homomorphism.

 
\begin{Verbatim}[commandchars=!@|,fontsize=\small,frame=single,label=Example]
  !gapprompt@gap>| !gapinput@S := Semigroup([(1, 2, 3, 4), (1, 2)]);;|
  !gapprompt@gap>| !gapinput@T := Semigroup([(1, 2, 3), (1, 2)]);;|
  !gapprompt@gap>| !gapinput@hom := SemigroupHomomorphismByFunction(S, T,|
  !gapprompt@>| !gapinput@function(x) if SignPerm(x) = -1 then return (1, 2);|
  !gapprompt@>| !gapinput@else return (); fi; end);|
  <semigroup of size 24, with 2 generators> ->
  <semigroup of size 6, with 2 generators>
\end{Verbatim}
 }

 

\subsection{\textcolor{Chapter }{AsSemigroupHomomorphismByFunction (for a semigroup homomorphism by images)}}
\logpage{[ 14, 1, 6 ]}\nobreak
\hyperdef{L}{X7973F31986CF0DD4}{}
{\noindent\textcolor{FuncColor}{$\triangleright$\enspace\texttt{AsSemigroupHomomorphismByFunction({\mdseries\slshape hom})\index{AsSemigroupHomomorphismByFunction@\texttt{AsSemigroupHomomorphismByFunction}!for a semigroup homomorphism by images}
\label{AsSemigroupHomomorphismByFunction:for a semigroup homomorphism by images}
}\hfill{\scriptsize (operation)}}\\
\textbf{\indent Returns:\ }
 A semigroup homomorphism. 



 \texttt{AsSemigroupHomomorphismByFunction} takes \mbox{\texttt{\mdseries\slshape hom}}, a semigroup homomorphism, and returns the same mapping but described by a \textsf{GAP} function mapping elements of \texttt{Source(\mbox{\texttt{\mdseries\slshape hom}})} to their images in \texttt{Range(\mbox{\texttt{\mdseries\slshape hom}})}.

 
\begin{Verbatim}[commandchars=!@|,fontsize=\small,frame=single,label=Example]
  !gapprompt@gap>| !gapinput@T := TrivialSemigroup();;|
  !gapprompt@gap>| !gapinput@S := GLM(2, 2);;|
  !gapprompt@gap>| !gapinput@gens := GeneratorsOfSemigroup(S);;|
  !gapprompt@gap>| !gapinput@imgs := ListX(gens, x -> IdentityTransformation);;|
  !gapprompt@gap>| !gapinput@hom := SemigroupHomomorphismByImages(S, T, gens, imgs);;|
  !gapprompt@gap>| !gapinput@hom := AsSemigroupHomomorphismByFunction(hom);|
  <general linear monoid 2x2 over GF(2)> ->
  <trivial transformation group of degree 0 with 1 generator>
\end{Verbatim}
 }

 

\subsection{\textcolor{Chapter }{KernelOfSemigroupHomomorphism}}
\logpage{[ 14, 1, 7 ]}\nobreak
\hyperdef{L}{X86BCE2207E55FC9F}{}
{\noindent\textcolor{FuncColor}{$\triangleright$\enspace\texttt{KernelOfSemigroupHomomorphism({\mdseries\slshape hom})\index{KernelOfSemigroupHomomorphism@\texttt{KernelOfSemigroupHomomorphism}}
\label{KernelOfSemigroupHomomorphism}
}\hfill{\scriptsize (attribute)}}\\
\textbf{\indent Returns:\ }
 A semigroup congruence. 



 \texttt{KernelOfSemigroupHomomorphism} returns the kernel of the semigroup homomorphism \mbox{\texttt{\mdseries\slshape hom}}. The kernel of a semigroup homomorphism \mbox{\texttt{\mdseries\slshape hom}} is a semigroup congruence relating pairs of elements in \texttt{Source(\mbox{\texttt{\mdseries\slshape hom}})} mapping to the same element under \mbox{\texttt{\mdseries\slshape hom}}. 
\begin{Verbatim}[commandchars=!@|,fontsize=\small,frame=single,label=Example]
  !gapprompt@gap>| !gapinput@S := Semigroup([Transformation([2, 1, 5, 1, 5]),|
  !gapprompt@>| !gapinput@      Transformation([1, 1, 1, 5, 3]),|
  !gapprompt@>| !gapinput@      Transformation([2, 5, 3, 5, 3])]);;|
  !gapprompt@gap>| !gapinput@congs := CongruencesOfSemigroup(S);;|
  !gapprompt@gap>| !gapinput@cong := congs[3];;|
  !gapprompt@gap>| !gapinput@T := S / cong;;|
  !gapprompt@gap>| !gapinput@gens := GeneratorsOfSemigroup(S);;|
  !gapprompt@gap>| !gapinput@images := List(gens, gen -> EquivalenceClassOfElement(cong, gen));;|
  !gapprompt@gap>| !gapinput@hom1 := SemigroupHomomorphismByImages(S, T, gens, images);;|
  !gapprompt@gap>| !gapinput@cong = KernelOfSemigroupHomomorphism(hom1);|
  true
\end{Verbatim}
 }

 }

   
\section{\textcolor{Chapter }{ Isomorphisms of arbitrary semigroups }}\logpage{[ 14, 2, 0 ]}
\hyperdef{L}{X7A8945817BD44943}{}
{
  

\subsection{\textcolor{Chapter }{IsIsomorphicSemigroup}}
\logpage{[ 14, 2, 1 ]}\nobreak
\hyperdef{L}{X7A6D59247F15935E}{}
{\noindent\textcolor{FuncColor}{$\triangleright$\enspace\texttt{IsIsomorphicSemigroup({\mdseries\slshape S, T})\index{IsIsomorphicSemigroup@\texttt{IsIsomorphicSemigroup}}
\label{IsIsomorphicSemigroup}
}\hfill{\scriptsize (operation)}}\\
\textbf{\indent Returns:\ }
 \texttt{true} or \texttt{false}. 



 If \mbox{\texttt{\mdseries\slshape S}} and \mbox{\texttt{\mdseries\slshape T}} are semigroups, then this operation attempts to determine whether \mbox{\texttt{\mdseries\slshape S}} and \mbox{\texttt{\mdseries\slshape T}} are isomorphic semigroups by using the operation \texttt{IsomorphismSemigroups} (\ref{IsomorphismSemigroups}). If \texttt{IsomorphismSemigroups(\mbox{\texttt{\mdseries\slshape S}}, \mbox{\texttt{\mdseries\slshape T}})} returns an isomorphism, then \texttt{IsIsomorphicSemigroup(\mbox{\texttt{\mdseries\slshape S}}, \mbox{\texttt{\mdseries\slshape T}})} returns \texttt{true}, while if \texttt{IsomorphismSemigroups(\mbox{\texttt{\mdseries\slshape S}}, \mbox{\texttt{\mdseries\slshape T}})} returns \texttt{fail}, then \texttt{IsIsomorphicSemigroup(\mbox{\texttt{\mdseries\slshape S}}, \mbox{\texttt{\mdseries\slshape T}})} returns \texttt{false}.

 Note that in some cases, at present, there is no method for determining
whether \mbox{\texttt{\mdseries\slshape S}} is isomorphic to \mbox{\texttt{\mdseries\slshape T}}, even if it is obvious to the user whether or not \mbox{\texttt{\mdseries\slshape S}} and \mbox{\texttt{\mdseries\slshape T}} are isomorphic. There are plans to improve this in the future. 

 
\begin{Verbatim}[commandchars=!@|,fontsize=\small,frame=single,label=Example]
  !gapprompt@gap>| !gapinput@S := Semigroup(PartialPerm([1, 2, 4], [1, 3, 5]),|
  !gapprompt@>| !gapinput@                  PartialPerm([1, 3, 5], [1, 2, 4]));;|
  !gapprompt@gap>| !gapinput@T := AsSemigroup(IsTransformationSemigroup, S);;|
  !gapprompt@gap>| !gapinput@IsIsomorphicSemigroup(S, T);|
  true
  !gapprompt@gap>| !gapinput@IsIsomorphicSemigroup(FullTransformationMonoid(4),|
  !gapprompt@>| !gapinput@PartitionMonoid(4));|
  false
\end{Verbatim}
 }

 

\subsection{\textcolor{Chapter }{SmallestMultiplicationTable}}
\logpage{[ 14, 2, 2 ]}\nobreak
\hyperdef{L}{X7DE212DF7DF0A4E9}{}
{\noindent\textcolor{FuncColor}{$\triangleright$\enspace\texttt{SmallestMultiplicationTable({\mdseries\slshape S})\index{SmallestMultiplicationTable@\texttt{SmallestMultiplicationTable}}
\label{SmallestMultiplicationTable}
}\hfill{\scriptsize (attribute)}}\\
\textbf{\indent Returns:\ }
 The lex\texttt{\symbol{45}}least multiplication table of a semigroup. 



 This function returns the lex\texttt{\symbol{45}}least multiplication table of
a semigroup isomorphic to the semigroup \mbox{\texttt{\mdseries\slshape S}}. \texttt{SmallestMultiplicationTable} returns the lex\texttt{\symbol{45}}least multiplication of any semigroup
isomorphic to \mbox{\texttt{\mdseries\slshape S}}. Due to the high complexity of computing the smallest multiplication table of
a semigroup, this function only performs well for semigroups with at most
approximately 50 elements.

 \texttt{SmallestMultiplicationTable} is based on the function \texttt{IdSmallSemigroup} (\textbf{Smallsemi: IdSmallSemigroup}) by Andreas Distler.

 From Version 3.3.0 of \textsf{Semigroups} this attribute is computed using \texttt{MinimalImage} (\textbf{images: MinimalImage}) from the the \href{https://gap-packages.github.io/images/} {images}  package. See also: \texttt{CanonicalMultiplicationTable} (\ref{CanonicalMultiplicationTable}). 
\begin{Verbatim}[commandchars=!@|,fontsize=\small,frame=single,label=Example]
  !gapprompt@gap>| !gapinput@S := Semigroup(|
  !gapprompt@>| !gapinput@Bipartition([[1, 2, 3, -1, -3], [-2]]),|
  !gapprompt@>| !gapinput@Bipartition([[1, 2, 3, -1], [-2], [-3]]),|
  !gapprompt@>| !gapinput@Bipartition([[1, 2, 3], [-1], [-2, -3]]),|
  !gapprompt@>| !gapinput@Bipartition([[1, 2, -1], [3, -2], [-3]]));;|
  !gapprompt@gap>| !gapinput@Size(S);|
  8
  !gapprompt@gap>| !gapinput@SmallestMultiplicationTable(S);|
  [ [ 1, 1, 3, 4, 5, 6, 7, 8 ], [ 1, 1, 3, 4, 5, 6, 7, 8 ],
    [ 1, 1, 3, 4, 5, 6, 7, 8 ], [ 1, 3, 3, 4, 5, 6, 7, 8 ],
    [ 5, 5, 6, 7, 5, 6, 7, 8 ], [ 5, 5, 6, 7, 5, 6, 7, 8 ],
    [ 5, 6, 6, 7, 5, 6, 7, 8 ], [ 5, 6, 6, 7, 5, 6, 7, 8 ] ]
\end{Verbatim}
 }

 

\subsection{\textcolor{Chapter }{CanonicalMultiplicationTable}}
\logpage{[ 14, 2, 3 ]}\nobreak
\hyperdef{L}{X7FFEEFF484039A42}{}
{\noindent\textcolor{FuncColor}{$\triangleright$\enspace\texttt{CanonicalMultiplicationTable({\mdseries\slshape S})\index{CanonicalMultiplicationTable@\texttt{CanonicalMultiplicationTable}}
\label{CanonicalMultiplicationTable}
}\hfill{\scriptsize (attribute)}}\\
\textbf{\indent Returns:\ }
 A canonical multiplication table (up to isomorphism) of a semigroup. 



 This function returns a multiplication table of a semigroup isomorphic to the
semigroup \mbox{\texttt{\mdseries\slshape S}}. \texttt{CanonicalMultiplicationTable} returns a multiplication that is canonical, in the sense that if two
semigroups \texttt{S} and \texttt{T} are isomorphic, then the return values of \texttt{CanonicalMultiplicationTable} are equal. 

 \texttt{CanonicalMultiplicationTable} uses the machinery for canonical labelling of vertex coloured digraphs in \href{http://www.tcs.tkk.fi/Software/bliss/} {bliss}  via \texttt{BlissCanonicalLabelling} (\textbf{Digraphs: BlissCanonicalLabelling for a digraph and a list}).

 The multiplication table returned by this function is the result of \texttt{OnMultiplicationTable(MultiplicationTable(\mbox{\texttt{\mdseries\slshape S}}), CanonicalMultiplicationTablePerm(\mbox{\texttt{\mdseries\slshape S}}));}

 Note that the performance of \texttt{CanonicalMultiplicationTable} is vastly superior to that of \texttt{SmallestMultiplicationTable}.

 See also: \texttt{CanonicalMultiplicationTablePerm} (\ref{CanonicalMultiplicationTablePerm}) and \texttt{OnMultiplicationTable} (\ref{OnMultiplicationTable}). 
\begin{Verbatim}[commandchars=!@|,fontsize=\small,frame=single,label=Example]
  !gapprompt@gap>| !gapinput@S := Semigroup(|
  !gapprompt@>| !gapinput@Bipartition([[1, 2, 3, -1, -3], [-2]]),|
  !gapprompt@>| !gapinput@Bipartition([[1, 2, 3, -1], [-2], [-3]]),|
  !gapprompt@>| !gapinput@Bipartition([[1, 2, 3], [-1], [-2, -3]]),|
  !gapprompt@>| !gapinput@Bipartition([[1, 2, -1], [3, -2], [-3]]));;|
  !gapprompt@gap>| !gapinput@Size(S);|
  8
  !gapprompt@gap>| !gapinput@CanonicalMultiplicationTable(S);|
  [ [ 1, 2, 2, 8, 1, 2, 7, 8 ], [ 1, 2, 2, 8, 1, 2, 7, 8 ],
    [ 1, 2, 6, 4, 5, 6, 7, 8 ], [ 1, 2, 5, 4, 5, 6, 7, 8 ],
    [ 1, 2, 6, 4, 5, 6, 7, 8 ], [ 1, 2, 6, 4, 5, 6, 7, 8 ],
    [ 1, 2, 1, 8, 1, 2, 7, 8 ], [ 1, 2, 1, 8, 1, 2, 7, 8 ] ]
\end{Verbatim}
 }

 

\subsection{\textcolor{Chapter }{CanonicalMultiplicationTablePerm}}
\logpage{[ 14, 2, 4 ]}\nobreak
\hyperdef{L}{X869533A7819EC2F8}{}
{\noindent\textcolor{FuncColor}{$\triangleright$\enspace\texttt{CanonicalMultiplicationTablePerm({\mdseries\slshape S})\index{CanonicalMultiplicationTablePerm@\texttt{CanonicalMultiplicationTablePerm}}
\label{CanonicalMultiplicationTablePerm}
}\hfill{\scriptsize (attribute)}}\\
\textbf{\indent Returns:\ }
 A permutation. 



 This function returns a permutation \texttt{p} such that \texttt{OnMultiplicationTable(MultiplicationTable(\mbox{\texttt{\mdseries\slshape S}}), p);} equals \texttt{CanonicalMultiplicationTable(\mbox{\texttt{\mdseries\slshape S}})}. 

 See \texttt{CanonicalMultiplicationTable} (\ref{CanonicalMultiplicationTable}) for more details.

 \texttt{CanonicalMultiplicationTablePerm} uses the machinery for canonical labelling of vertex coloured digraphs in \href{http://www.tcs.tkk.fi/Software/bliss/} {bliss}  via \texttt{BlissCanonicalLabelling} (\textbf{Digraphs: BlissCanonicalLabelling for a digraph and a list}).

 
\begin{Verbatim}[commandchars=!@|,fontsize=\small,frame=single,label=Example]
  !gapprompt@gap>| !gapinput@S := Semigroup(|
  !gapprompt@>| !gapinput@Bipartition([[1, 2, 3, -1, -3], [-2]]),|
  !gapprompt@>| !gapinput@Bipartition([[1, 2, 3, -1], [-2], [-3]]),|
  !gapprompt@>| !gapinput@Bipartition([[1, 2, 3], [-1], [-2, -3]]),|
  !gapprompt@>| !gapinput@Bipartition([[1, 2, -1], [3, -2], [-3]]));;|
  !gapprompt@gap>| !gapinput@Size(S);|
  8
  !gapprompt@gap>| !gapinput@CanonicalMultiplicationTablePerm(S);|
  (1,5,8,3,6,7,2,4)
\end{Verbatim}
 }

 

\subsection{\textcolor{Chapter }{OnMultiplicationTable}}
\logpage{[ 14, 2, 5 ]}\nobreak
\hyperdef{L}{X83BC6B998479BD27}{}
{\noindent\textcolor{FuncColor}{$\triangleright$\enspace\texttt{OnMultiplicationTable({\mdseries\slshape table, p})\index{OnMultiplicationTable@\texttt{OnMultiplicationTable}}
\label{OnMultiplicationTable}
}\hfill{\scriptsize (operation)}}\\
\noindent\textcolor{FuncColor}{$\triangleright$\enspace\texttt{PermuteMultiplicationTable({\mdseries\slshape output, table, p})\index{PermuteMultiplicationTable@\texttt{PermuteMultiplicationTable}}
\label{PermuteMultiplicationTable}
}\hfill{\scriptsize (operation)}}\\
\noindent\textcolor{FuncColor}{$\triangleright$\enspace\texttt{PermuteMultiplicationTableNC({\mdseries\slshape output, table, p})\index{PermuteMultiplicationTableNC@\texttt{PermuteMultiplicationTableNC}}
\label{PermuteMultiplicationTableNC}
}\hfill{\scriptsize (operation)}}\\


 If \mbox{\texttt{\mdseries\slshape table}} is a multiplication table of a semigroup and \mbox{\texttt{\mdseries\slshape p}} is a permutation of \texttt{[1 .. Length(\mbox{\texttt{\mdseries\slshape table}})]}, then \texttt{OnMultiplicationTable} returns the multiplication table obtained from \mbox{\texttt{\mdseries\slshape table}} where the elements \texttt{[1 .. Length(\mbox{\texttt{\mdseries\slshape table}})]} are relabelled according to \mbox{\texttt{\mdseries\slshape p}}.

 \texttt{PermuteMultiplicationTable} works similarly but takes an additional argument \mbox{\texttt{\mdseries\slshape output}}, which should be a list of mutable lists of the same dimensions as \mbox{\texttt{\mdseries\slshape table}}. The argument \mbox{\texttt{\mdseries\slshape output}} is modified in\texttt{\symbol{45}}place to store the multiplication table that
would be given by \texttt{OnMultiplicationTable(\mbox{\texttt{\mdseries\slshape table}}, \mbox{\texttt{\mdseries\slshape p}})}, while the argument \mbox{\texttt{\mdseries\slshape table}} remains unchanged. This approach is preferred over \texttt{OnMultiplicationTable} when the user does not need to create a separate result variable.

 \texttt{PermuteMultiplicationTableNC} is the operation called by \texttt{PermuteMultiplicationTable}. It assumes that the arguments \mbox{\texttt{\mdseries\slshape output}}, \mbox{\texttt{\mdseries\slshape table}} and \mbox{\texttt{\mdseries\slshape p}} are well\texttt{\symbol{45}}formed and can therefore lead to unexpected
results. 

 
\begin{Verbatim}[commandchars=!@|,fontsize=\small,frame=single,label=Example]
  !gapprompt@gap>| !gapinput@table := [[1, 1, 3, 4, 5, 6, 7, 8],|
  !gapprompt@>| !gapinput@             [1, 1, 3, 4, 5, 6, 7, 8],|
  !gapprompt@>| !gapinput@             [1, 1, 3, 4, 5, 6, 7, 8],|
  !gapprompt@>| !gapinput@             [1, 3, 3, 4, 5, 6, 7, 8],|
  !gapprompt@>| !gapinput@             [5, 5, 6, 7, 5, 6, 7, 8],|
  !gapprompt@>| !gapinput@             [5, 5, 6, 7, 5, 6, 7, 8],|
  !gapprompt@>| !gapinput@             [5, 6, 6, 7, 5, 6, 7, 8],|
  !gapprompt@>| !gapinput@             [5, 6, 6, 7, 5, 6, 7, 8]];;|
  !gapprompt@gap>| !gapinput@p := (1, 2, 3, 4)(10, 11, 12);;|
  !gapprompt@gap>| !gapinput@OnMultiplicationTable(table, p);|
  [ [ 1, 2, 4, 4, 5, 6, 7, 8 ], [ 1, 2, 2, 4, 5, 6, 7, 8 ],
    [ 1, 2, 2, 4, 5, 6, 7, 8 ], [ 1, 2, 2, 4, 5, 6, 7, 8 ],
    [ 7, 5, 5, 6, 5, 6, 7, 8 ], [ 7, 5, 5, 6, 5, 6, 7, 8 ],
    [ 7, 5, 6, 6, 5, 6, 7, 8 ], [ 7, 5, 6, 6, 5, 6, 7, 8 ] ]
  !gapprompt@gap>| !gapinput@temp := List([1 .. Size(table)], x -> [1 .. Size(table)]);;|
  !gapprompt@gap>| !gapinput@PermuteMultiplicationTable(temp, table, p);;|
  !gapprompt@gap>| !gapinput@temp = OnMultiplicationTable(table, p);|
  true
    
\end{Verbatim}
 }

 

\subsection{\textcolor{Chapter }{IsomorphismSemigroups}}
\logpage{[ 14, 2, 6 ]}\nobreak
\hyperdef{L}{X8248C522825E2684}{}
{\noindent\textcolor{FuncColor}{$\triangleright$\enspace\texttt{IsomorphismSemigroups({\mdseries\slshape S, T})\index{IsomorphismSemigroups@\texttt{IsomorphismSemigroups}}
\label{IsomorphismSemigroups}
}\hfill{\scriptsize (operation)}}\\
\textbf{\indent Returns:\ }
 An isomorphism, or \texttt{fail}. 



 This operation attempts to find an isomorphism from the semigroup \mbox{\texttt{\mdseries\slshape S}} to the semigroup \mbox{\texttt{\mdseries\slshape T}}. If it finds one, then it is returned, and if not, then \texttt{fail} is returned. 

 \texttt{IsomorphismSemigroups} uses the machinery for finding isomorphisms between vertex coloured digraphs
in \href{http://www.tcs.tkk.fi/Software/bliss/} {bliss}  via \texttt{IsomorphismDigraphs} (\textbf{Digraphs: IsomorphismDigraphs for digraphs and homogeneous lists}) using digraphs constructed from the multiplication tables of \mbox{\texttt{\mdseries\slshape S}} and \mbox{\texttt{\mdseries\slshape T}}. 

 Note that finding an isomorphism between two semigroups is difficult, and may
not be possible for semigroups whose size exceeds a few hundred elements. On
the other hand, \texttt{IsomorphismSemigroups} may be able deduce that \mbox{\texttt{\mdseries\slshape S}} and \mbox{\texttt{\mdseries\slshape T}} are not isomorphic by finding that some of their
semigroup\texttt{\symbol{45}}theoretic properties differ. 
\begin{Verbatim}[commandchars=!@|,fontsize=\small,frame=single,label=Example]
  !gapprompt@gap>| !gapinput@S := RectangularBand(IsTransformationSemigroup, 4, 5);|
  <regular transformation semigroup of size 20, degree 9 with 5
   generators>
  !gapprompt@gap>| !gapinput@T := RectangularBand(IsBipartitionSemigroup, 4, 5);|
  <regular bipartition semigroup of size 20, degree 3 with 5 generators>
  !gapprompt@gap>| !gapinput@IsomorphismSemigroups(S, T) <> fail;|
  true
  !gapprompt@gap>| !gapinput@D := DClass(FullTransformationMonoid(5),|
  !gapprompt@>| !gapinput@               Transformation([1, 2, 3, 4, 1]));;|
  !gapprompt@gap>| !gapinput@S := PrincipalFactor(D);;|
  !gapprompt@gap>| !gapinput@StructureDescription(UnderlyingSemigroup(S));|
  "S4"
  !gapprompt@gap>| !gapinput@S;|
  <Rees 0-matrix semigroup 10x5 over S4>
  !gapprompt@gap>| !gapinput@D := DClass(PartitionMonoid(5),|
  !gapprompt@>| !gapinput@Bipartition([[1], [2, -2], [3, -3], [4, -4], [5, -5], [-1]]));;|
  !gapprompt@gap>| !gapinput@T := PrincipalFactor(D);;|
  !gapprompt@gap>| !gapinput@StructureDescription(UnderlyingSemigroup(T));|
  "S4"
  !gapprompt@gap>| !gapinput@T;|
  <Rees 0-matrix semigroup 15x15 over S4>
  !gapprompt@gap>| !gapinput@IsomorphismSemigroups(S, T);|
  fail
  !gapprompt@gap>| !gapinput@I := SemigroupIdeal(FullTransformationMonoid(5),|
  !gapprompt@>| !gapinput@                       Transformation([1, 1, 2, 3, 4]));;|
  !gapprompt@gap>| !gapinput@T := PrincipalFactor(DClass(I, I.1));;|
  !gapprompt@gap>| !gapinput@StructureDescription(UnderlyingSemigroup(T));|
  "S4"
  !gapprompt@gap>| !gapinput@T;|
  <Rees 0-matrix semigroup 10x5 over S4>
  !gapprompt@gap>| !gapinput@IsomorphismSemigroups(S, T) <> fail;|
  true
\end{Verbatim}
 }

 

\subsection{\textcolor{Chapter }{AutomorphismGroup (for a semigroup)}}
\logpage{[ 14, 2, 7 ]}\nobreak
\hyperdef{L}{X79BFF4E77A8090EF}{}
{\noindent\textcolor{FuncColor}{$\triangleright$\enspace\texttt{AutomorphismGroup({\mdseries\slshape S})\index{AutomorphismGroup@\texttt{AutomorphismGroup}!for a semigroup}
\label{AutomorphismGroup:for a semigroup}
}\hfill{\scriptsize (operation)}}\\
\textbf{\indent Returns:\ }
 A group. 



 This operation returns the group of automorphisms of the semigroup \mbox{\texttt{\mdseries\slshape S}}. \texttt{AutomorphismGroup} uses \href{http://www.tcs.tkk.fi/Software/bliss/} {bliss}  via \texttt{AutomorphismGroup} (\textbf{Digraphs: AutomorphismGroup for a digraph and a homogeneous list}) using a vertex coloured digraph constructed from the multiplication table of \mbox{\texttt{\mdseries\slshape S}}. Consequently, this method is only really feasible for semigroups whose size
does not exceed a few hundred elements. 

 
\begin{Verbatim}[commandchars=!@|,fontsize=\small,frame=single,label=Example]
  !gapprompt@gap>| !gapinput@S := RectangularBand(IsTransformationSemigroup, 4, 5);|
  <regular transformation semigroup of size 20, degree 9 with 5
   generators>
  !gapprompt@gap>| !gapinput@StructureDescription(AutomorphismGroup(S));|
  "S4 x S5"
\end{Verbatim}
 }

 

\subsection{\textcolor{Chapter }{SemigroupIsomorphismByImages (for two semigroups and two lists)}}
\logpage{[ 14, 2, 8 ]}\nobreak
\hyperdef{L}{X80FE565183A9410D}{}
{\noindent\textcolor{FuncColor}{$\triangleright$\enspace\texttt{SemigroupIsomorphismByImages({\mdseries\slshape S, T, gens, imgs})\index{SemigroupIsomorphismByImages@\texttt{SemigroupIsomorphismByImages}!for two semigroups and two lists}
\label{SemigroupIsomorphismByImages:for two semigroups and two lists}
}\hfill{\scriptsize (operation)}}\\
\noindent\textcolor{FuncColor}{$\triangleright$\enspace\texttt{SemigroupIsomorphismByImages({\mdseries\slshape S, T, imgs})\index{SemigroupIsomorphismByImages@\texttt{SemigroupIsomorphismByImages}!for two semigroups and a list}
\label{SemigroupIsomorphismByImages:for two semigroups and a list}
}\hfill{\scriptsize (operation)}}\\
\noindent\textcolor{FuncColor}{$\triangleright$\enspace\texttt{SemigroupIsomorphismByImages({\mdseries\slshape S, T})\index{SemigroupIsomorphismByImages@\texttt{SemigroupIsomorphismByImages}!for two semigroups}
\label{SemigroupIsomorphismByImages:for two semigroups}
}\hfill{\scriptsize (operation)}}\\
\noindent\textcolor{FuncColor}{$\triangleright$\enspace\texttt{SemigroupIsomorphismByImages({\mdseries\slshape S, gens, imgs})\index{SemigroupIsomorphismByImages@\texttt{SemigroupIsomorphismByImages}!for a semigroup and two list}
\label{SemigroupIsomorphismByImages:for a semigroup and two list}
}\hfill{\scriptsize (operation)}}\\
\textbf{\indent Returns:\ }
 A semigroup isomorphism, or \texttt{fail}. 



 \texttt{SemigroupIsomorphismByImages} attempts to construct a isomorphism from the semigroup \mbox{\texttt{\mdseries\slshape S}} to the semigroup \mbox{\texttt{\mdseries\slshape T}}, by mapping the \texttt{i}\texttt{\symbol{45}}th element of \mbox{\texttt{\mdseries\slshape gens}} to the \texttt{i}\texttt{\symbol{45}}th element of \mbox{\texttt{\mdseries\slshape imgs}}. If this mapping corresponds to an isomorphism, the isomorphism is returned,
and if not, then \texttt{fail} is returned. An isomorphism is a bijective homomorphism. See also \texttt{SemigroupHomomorphismByImages} (\ref{SemigroupHomomorphismByImages:for two semigroups and two lists}). 
\begin{Verbatim}[commandchars=!@|,fontsize=\small,frame=single,label=Example]
  !gapprompt@gap>| !gapinput@S := Semigroup([|
  !gapprompt@>| !gapinput@ Matrix(IsNTPMatrix, [[0, 1, 2], [4, 3, 0], [0, 2, 0]], 9, 4),|
  !gapprompt@>| !gapinput@ Matrix(IsNTPMatrix, [[1, 1, 0], [4, 1, 1], [0, 0, 0]], 9, 4)]);;|
  !gapprompt@gap>| !gapinput@T := AsSemigroup(IsTransformationSemigroup, S);;|
  !gapprompt@gap>| !gapinput@iso := SemigroupIsomorphismByImages(S, T);|
  <semigroup of size 46, 3x3 ntp matrices with 2 generators> ->
  <transformation semigroup of size 46, degree 47 with 2 generators>
\end{Verbatim}
 }

 

\subsection{\textcolor{Chapter }{SemigroupIsomorphismByFunctionNC}}
\logpage{[ 14, 2, 9 ]}\nobreak
\hyperdef{L}{X7B44408D8309C3DC}{}
{\noindent\textcolor{FuncColor}{$\triangleright$\enspace\texttt{SemigroupIsomorphismByFunctionNC({\mdseries\slshape S, T, fun, invFun})\index{SemigroupIsomorphismByFunctionNC@\texttt{SemigroupIsomorphismByFunctionNC}}
\label{SemigroupIsomorphismByFunctionNC}
}\hfill{\scriptsize (operation)}}\\
\noindent\textcolor{FuncColor}{$\triangleright$\enspace\texttt{SemigroupIsomorphismByFunction({\mdseries\slshape S, T, fun, invFun})\index{SemigroupIsomorphismByFunction@\texttt{SemigroupIsomorphismByFunction}}
\label{SemigroupIsomorphismByFunction}
}\hfill{\scriptsize (operation)}}\\
\textbf{\indent Returns:\ }
 A semigroup isomorphism or \texttt{fail}. 



 \texttt{SemigroupIsomorphismByFunctionNC} returns a semigroup isomorphism with source \mbox{\texttt{\mdseries\slshape S}} and range \mbox{\texttt{\mdseries\slshape T}}, such that each element \texttt{s} in \mbox{\texttt{\mdseries\slshape S}} is mapped to the element \mbox{\texttt{\mdseries\slshape fun}}\texttt{(s)}, where \mbox{\texttt{\mdseries\slshape fun}} is a \textsf{GAP} function, and \mbox{\texttt{\mdseries\slshape invFun}} its inverse, mapping \mbox{\texttt{\mdseries\slshape fun}}\texttt{(s)} back to \texttt{s}. 

 The function \texttt{SemigroupIsomorphismByFunctionNC} performs no checks on whether the function actually gives an isomorphism, and
so it is possible for this operation to return a mapping from \mbox{\texttt{\mdseries\slshape S}} to \mbox{\texttt{\mdseries\slshape T}} that is not a homomorphism, or not a bijection, or where the return value of \texttt{InverseGeneralMapping} (\textbf{Reference: InverseGeneralMapping}) is not the inverse of the returned function.

 The function \texttt{SemigroupIsomorphismByFunction} checks that: the mapping from \mbox{\texttt{\mdseries\slshape S}} to \mbox{\texttt{\mdseries\slshape T}} defined by \mbox{\texttt{\mdseries\slshape fun}} satisfies \texttt{RespectsMultiplication} (\textbf{Reference: RespectsMultiplication}); that the function from \mbox{\texttt{\mdseries\slshape T}} to \mbox{\texttt{\mdseries\slshape S}} defined by \mbox{\texttt{\mdseries\slshape invFun}} satisfies \texttt{RespectsMultiplication} (\textbf{Reference: RespectsMultiplication}); and that these functions are mutual inverses. This can be expensive. If any
of these checks fails, then \texttt{fail} is returned. 
\begin{Verbatim}[commandchars=!@|,fontsize=\small,frame=single,label=Example]
  !gapprompt@gap>| !gapinput@S := MonogenicSemigroup(IsTransformationSemigroup, 3, 2);;|
  !gapprompt@gap>| !gapinput@T := MonogenicSemigroup(IsBipartitionSemigroup, 3, 2);;|
  !gapprompt@gap>| !gapinput@map := x -> T.1 ^ Length(Factorization(S, x));;|
  !gapprompt@gap>| !gapinput@inv := x -> S.1 ^ Length(Factorization(T, x));;|
  !gapprompt@gap>| !gapinput@iso := SemigroupIsomorphismByFunction(S, T, map, inv);|
  <commutative non-regular transformation semigroup of size 4, degree 5
    with 1 generator> -> <commutative non-regular block bijection
    semigroup of size 4, degree 6 with 1 generator>
\end{Verbatim}
 }

 

\subsection{\textcolor{Chapter }{IsSemigroupIsomorphismByFunction}}
\logpage{[ 14, 2, 10 ]}\nobreak
\hyperdef{L}{X7EFDBD2C7A4FB6AF}{}
{\noindent\textcolor{FuncColor}{$\triangleright$\enspace\texttt{IsSemigroupIsomorphismByFunction({\mdseries\slshape iso})\index{IsSemigroupIsomorphismByFunction@\texttt{IsSemigroupIsomorphismByFunction}}
\label{IsSemigroupIsomorphismByFunction}
}\hfill{\scriptsize (filter)}}\\
\textbf{\indent Returns:\ }
 \texttt{true} or \texttt{false}. 



 \texttt{IsSemigroupIsomorphismByFunction} returns \texttt{true} if \mbox{\texttt{\mdseries\slshape hom}} satisfies \texttt{IsSemigroupHomomorphismByFunction} (\ref{IsSemigroupHomomorphismByFunction}) and \texttt{IsBijective} (\textbf{Reference: IsBijective}), and \texttt{false} if does not. Note that this filter may return \texttt{true} even if the underlying \textsf{GAP} function does not define a homomorphism. A semigroup isomorphism is a mapping
from a semigroup \texttt{S} to a semigroup \texttt{T} that respects multiplication. This representation describes semigroup
isomorphisms internally by using a \textsf{GAP} function mapping elements of \texttt{S} to their images in \texttt{T}. See \texttt{SemigroupIsomorphismByFunction} (\ref{SemigroupIsomorphismByFunction}).

 
\begin{Verbatim}[commandchars=!@|,fontsize=\small,frame=single,label=Example]
  !gapprompt@gap>| !gapinput@S := MonogenicSemigroup(IsTransformationSemigroup, 3, 2);;|
  !gapprompt@gap>| !gapinput@T := MonogenicSemigroup(IsBipartitionSemigroup, 3, 2);;|
  !gapprompt@gap>| !gapinput@map := x -> T.1 ^ Length(Factorization(S, x));;|
  !gapprompt@gap>| !gapinput@inv := x -> S.1 ^ Length(Factorization(T, x));;|
  !gapprompt@gap>| !gapinput@iso := SemigroupIsomorphismByFunction(S, T, map, inv);|
  <commutative non-regular transformation semigroup of size 4, degree 5
    with 1 generator> -> <commutative non-regular block bijection
    semigroup of size 4, degree 6 with 1 generator>
  !gapprompt@gap>| !gapinput@IsSemigroupIsomorphismByFunction(iso);|
  true
\end{Verbatim}
 }

 

\subsection{\textcolor{Chapter }{AsSemigroupIsomorphismByFunction (for a semigroup homomorphism by images)}}
\logpage{[ 14, 2, 11 ]}\nobreak
\hyperdef{L}{X86C4FC857AF125BD}{}
{\noindent\textcolor{FuncColor}{$\triangleright$\enspace\texttt{AsSemigroupIsomorphismByFunction({\mdseries\slshape hom})\index{AsSemigroupIsomorphismByFunction@\texttt{AsSemigroupIsomorphismByFunction}!for a semigroup homomorphism by images}
\label{AsSemigroupIsomorphismByFunction:for a semigroup homomorphism by images}
}\hfill{\scriptsize (operation)}}\\
\textbf{\indent Returns:\ }
 A semigroup isomorphism, or \texttt{fail}. 



 \texttt{AsSemigroupIsomorphismByFunction} takes a semigroup homomorphism \mbox{\texttt{\mdseries\slshape hom}} and returns a semigroup isomorphism represented using \textsf{GAP} functions for the isomorphism and its inverse. If \mbox{\texttt{\mdseries\slshape hom}} is not bijective, then \texttt{fail} is returned. 
\begin{Verbatim}[commandchars=!@|,fontsize=\small,frame=single,label=Example]
  !gapprompt@gap>| !gapinput@S := FullTransformationMonoid(3);;|
  !gapprompt@gap>| !gapinput@gens := GeneratorsOfSemigroup(S);;|
  !gapprompt@gap>| !gapinput@imgs := ListWithIdenticalEntries(4, ConstantTransformation(3, 1));;|
  !gapprompt@gap>| !gapinput@hom := SemigroupHomomorphismByImages(S, S, gens, gens);;|
  !gapprompt@gap>| !gapinput@AsSemigroupIsomorphismByFunction(hom);|
  <full transformation monoid of degree 3> ->
  <full transformation monoid of degree 3>
\end{Verbatim}
 }

 

\subsection{\textcolor{Chapter }{SmallerDegreeTransformationRepresentation}}
\logpage{[ 14, 2, 12 ]}\nobreak
\hyperdef{L}{X794E5DA4872989E4}{}
{\noindent\textcolor{FuncColor}{$\triangleright$\enspace\texttt{SmallerDegreeTransformationRepresentation({\mdseries\slshape S})\index{SmallerDegreeTransformationRepresentation@\texttt{Smaller}\-\texttt{Degree}\-\texttt{Transformation}\-\texttt{Representation}}
\label{SmallerDegreeTransformationRepresentation}
}\hfill{\scriptsize (attribute)}}\\
\textbf{\indent Returns:\ }
An isomorphism to a transformation semigroup.



 This function attempts to find a small degree transformation representation of
the semigroup \mbox{\texttt{\mdseries\slshape S}}. The implementation attempts to find a right congruence of \mbox{\texttt{\mdseries\slshape S}} that \mbox{\texttt{\mdseries\slshape S}} acts on (the equivalence classes of) faithfully. 

 If \mbox{\texttt{\mdseries\slshape S}} is not a finitely presented semigroup, then the returned isomorphism is the
composition of an isomorphism to a finitely presented semigroup and an
isomorphism from that finitely presented semigroup to a transformation
semigroup. 

 The runtime of this function depends on the presentation for \mbox{\texttt{\mdseries\slshape S}} that is either given explicitly or computed by the \textsf{Semigroups} package, but it is difficult to predict what properties of the presentation
lead to a shorter runtime. This is unlikely to terminate in a reasonable
amount of time for semigroups with more than approx. \texttt{10000} elements, but might also not terminate quickly for smaller semigroups
depending on the presentation used. 
\begin{Verbatim}[commandchars=!@|,fontsize=\small,frame=single,label=Example]
  !gapprompt@gap>| !gapinput@S := BrauerMonoid(3);|
  <regular bipartition *-monoid of degree 3 with 3 generators>
  !gapprompt@gap>| !gapinput@IsomorphismTransformationSemigroup(S);|
  <regular bipartition *-monoid of size 15, degree 3 with 3 generators>
  -> <transformation monoid of size 15, degree 15 with 3 generators>
  !gapprompt@gap>| !gapinput@SmallerDegreeTransformationRepresentation(S);|
  CompositionMapping(
  <fp semigroup with 4 generators and 20 relations of length 81> ->
  <transformation monoid of degree 7 with 3 generators>,
  <regular bipartition *-monoid of size 15, degree 3 with 3 generators>
  -> <fp semigroup with 4 generators and 20 relations of length 81> )
  !gapprompt@gap>| !gapinput@S := JonesMonoid(5);|
  <regular bipartition *-monoid of degree 5 with 4 generators>
  !gapprompt@gap>| !gapinput@Size(S);|
  42
  !gapprompt@gap>| !gapinput@SmallerDegreeTransformationRepresentation(S);|
  CompositionMapping(
  <fp semigroup with 5 generators and 28 relations of length 120> ->
  <transformation monoid of degree 10 with 4 generators>,
  <regular bipartition *-monoid of size 42, degree 5 with 4 generators>
  -> <fp semigroup with 5 generators and 28 relations of length 120> )
\end{Verbatim}
 }

 

\subsection{\textcolor{Chapter }{MinimalFaithfulTransformationDegree}}
\logpage{[ 14, 2, 13 ]}\nobreak
\hyperdef{L}{X867264587CFD0013}{}
{\noindent\textcolor{FuncColor}{$\triangleright$\enspace\texttt{MinimalFaithfulTransformationDegree({\mdseries\slshape S})\index{MinimalFaithfulTransformationDegree@\texttt{MinimalFaithfulTransformationDegree}}
\label{MinimalFaithfulTransformationDegree}
}\hfill{\scriptsize (attribute)}}\\
\textbf{\indent Returns:\ }
A positive integer.



 This function returns the minimal degree of a faithful transformation
representation of the semigroup \mbox{\texttt{\mdseries\slshape S}}. This is currently only implemented for a very small number of types of
semigroups. 

  
\begin{Verbatim}[commandchars=!@|,fontsize=\small,frame=single,label=Example]
  !gapprompt@gap>| !gapinput@S := RightZeroSemigroup(10);|
  <transformation semigroup of degree 7 with 10 generators>
  !gapprompt@gap>| !gapinput@MinimalFaithfulTransformationDegree(S);|
  7
\end{Verbatim}
 }

 }

   
\section{\textcolor{Chapter }{ Isomorphisms of Rees (0\texttt{\symbol{45}})matrix semigroups }}\logpage{[ 14, 3, 0 ]}
\hyperdef{L}{X80DE3DB0782D9358}{}
{
  An isomorphism between two regular finite Rees (0\texttt{\symbol{45}})matrix
semigroups whose underlying semigroups are groups can be described by a triple
defined in terms of the matrices and underlying groups of the semigroups. For
a full description of the theory involved, see Section 3.4 of \cite{Howie1995aa}. 

 An isomorphism described in this way can be constructed using \texttt{RMSIsoByTriple} (\ref{RMSIsoByTriple}) or \texttt{RZMSIsoByTriple} (\ref{RZMSIsoByTriple}), and will satisfy the filter \texttt{IsRMSIsoByTriple} (\ref{IsRMSIsoByTriple}) or \texttt{IsRZMSIsoByTriple} (\ref{IsRZMSIsoByTriple}). 

 

\subsection{\textcolor{Chapter }{IsRMSIsoByTriple}}
\logpage{[ 14, 3, 1 ]}\nobreak
\hyperdef{L}{X82FCB1E585429FEA}{}
{\noindent\textcolor{FuncColor}{$\triangleright$\enspace\texttt{IsRMSIsoByTriple\index{IsRMSIsoByTriple@\texttt{IsRMSIsoByTriple}}
\label{IsRMSIsoByTriple}
}\hfill{\scriptsize (Category)}}\\
\noindent\textcolor{FuncColor}{$\triangleright$\enspace\texttt{IsRZMSIsoByTriple\index{IsRZMSIsoByTriple@\texttt{IsRZMSIsoByTriple}}
\label{IsRZMSIsoByTriple}
}\hfill{\scriptsize (Category)}}\\


 The isomorphisms between finite Rees matrix or 0\texttt{\symbol{45}}matrix
semigroups \texttt{S} and \texttt{T} over groups \texttt{G} and \texttt{H}, respectively, specified by a triple consisting of: 
\begin{enumerate}
\item  an isomorphism of the underlying graph of \texttt{S} to the underlying graph of of \texttt{T} 
\item  an isomorphism from \texttt{G} to \texttt{H} 
\item  a function from \texttt{Rows(S)} union \texttt{Columns(S)} to \texttt{H} 
\end{enumerate}
 belong to the categories \texttt{IsRMSIsoByTriple} and \texttt{IsRZMSIsoByTriple}. Basic operators for such isomorphism are given in \ref{Operators for isomorphisms of Rees (0-)matrix semigroups}, and basic operations are: \texttt{Range} (\textbf{Reference: range}), \texttt{Source} (\textbf{Reference: Source}), \texttt{ELM{\textunderscore}LIST} (\ref{ELMuScoreLIST:for IsRMSIsoByTriple}), \texttt{CompositionMapping} (\textbf{Reference: CompositionMapping}), \texttt{ImagesElm} (\ref{ImagesElm:for IsRMSIsoByTriple}), \texttt{ImagesRepresentative} (\ref{ImagesRepresentative:for IsRMSIsoByTriple}), \texttt{InverseGeneralMapping} (\textbf{Reference: InverseGeneralMapping}), \texttt{PreImagesRepresentative} (\textbf{Reference: PreImagesRepresentative}), \texttt{IsOne} (\textbf{Reference: IsOne}). }

 

\subsection{\textcolor{Chapter }{RMSIsoByTriple}}
\logpage{[ 14, 3, 2 ]}\nobreak
\hyperdef{L}{X82B0BDCD7CBDCC2E}{}
{\noindent\textcolor{FuncColor}{$\triangleright$\enspace\texttt{RMSIsoByTriple({\mdseries\slshape R1, R2, triple})\index{RMSIsoByTriple@\texttt{RMSIsoByTriple}}
\label{RMSIsoByTriple}
}\hfill{\scriptsize (operation)}}\\
\noindent\textcolor{FuncColor}{$\triangleright$\enspace\texttt{RZMSIsoByTriple({\mdseries\slshape R1, R2, triple})\index{RZMSIsoByTriple@\texttt{RZMSIsoByTriple}}
\label{RZMSIsoByTriple}
}\hfill{\scriptsize (operation)}}\\
\textbf{\indent Returns:\ }
An isomorphism.



 If \mbox{\texttt{\mdseries\slshape R1}} and \mbox{\texttt{\mdseries\slshape R2}} are isomorphic regular Rees 0\texttt{\symbol{45}}matrix semigroups whose
underlying semigroups are groups then \texttt{RZMSIsoByTriple} returns the isomorphism between \mbox{\texttt{\mdseries\slshape R1}} and \mbox{\texttt{\mdseries\slshape R2}} defined by \mbox{\texttt{\mdseries\slshape triple}}, which should be a list consisting of the following: 
\begin{itemize}
\item  \texttt{\mbox{\texttt{\mdseries\slshape triple}}[1]} should be a permutation describing an isomorphism from the graph of \mbox{\texttt{\mdseries\slshape R1}} to the graph of \mbox{\texttt{\mdseries\slshape R2}}, i.e. it should satisfy \texttt{OnDigraphs(RZMSDigraph(\mbox{\texttt{\mdseries\slshape R1}}), \mbox{\texttt{\mdseries\slshape triple}}[1]) = RZMSDigraph(\mbox{\texttt{\mdseries\slshape R2}})}. 
\item  \texttt{\mbox{\texttt{\mdseries\slshape triple}}[2]} should be an isomorphism from the underlying group of \mbox{\texttt{\mdseries\slshape R1}} to the underlying group of \mbox{\texttt{\mdseries\slshape R2}} (see \texttt{UnderlyingSemigroup} (\textbf{Reference: UnderlyingSemigroup for a Rees 0\texttt{\symbol{45}}matrix semigroup})). 
\item  \texttt{\mbox{\texttt{\mdseries\slshape triple}}[3]} should be a list of elements from the underlying group of \mbox{\texttt{\mdseries\slshape R2}}. If the \texttt{Matrix} (\textbf{Reference: Matrix}) of \mbox{\texttt{\mdseries\slshape R1}} has $m$ columns and $n$ rows, then the list should have length $m + n$, where the first $m$ entries should correspond to the columns of \mbox{\texttt{\mdseries\slshape R1}}'s matrix, and the last $n$ entries should correspond to the rows. These column and row entries should
correspond to the $u_i$ and $v_\lambda$ elements in Theorem 3.4.1 of \cite{Howie1995aa}. 
\end{itemize}
 If \mbox{\texttt{\mdseries\slshape triple}} describes a valid isomorphism from \mbox{\texttt{\mdseries\slshape R1}} to \mbox{\texttt{\mdseries\slshape R2}} then this will return an object in the category \texttt{IsRZMSIsoByTriple} (\ref{IsRZMSIsoByTriple}); otherwise an error will be returned. 

 If \mbox{\texttt{\mdseries\slshape R1}} and \mbox{\texttt{\mdseries\slshape R2}} are instead Rees matrix semigroups (without zero) then \texttt{RMSIsoByTriple} should be used instead. This operation is used in the same way, but it should
be noted that since an RMS's graph is a complete bipartite graph, \texttt{\mbox{\texttt{\mdseries\slshape triple}}[1]} can be any permutation on \texttt{[1 .. m + n]}, so long as no point in \texttt{[1 .. m]} is mapped to a point in \texttt{[m + 1 .. m + n]}. 

 
\begin{Verbatim}[commandchars=!@|,fontsize=\small,frame=single,label=Example]
  !gapprompt@gap>| !gapinput@g := SymmetricGroup(3);;|
  !gapprompt@gap>| !gapinput@mat := [[0, 0, (1, 3)], [(1, 2, 3), (), (2, 3)], [0, 0, ()]];;|
  !gapprompt@gap>| !gapinput@R := ReesZeroMatrixSemigroup(g, mat);;|
  !gapprompt@gap>| !gapinput@id := IdentityMapping(g);;|
  !gapprompt@gap>| !gapinput@g_elms_list := [(), (), (), (), (), ()];;|
  !gapprompt@gap>| !gapinput@RZMSIsoByTriple(R, R, [(), id, g_elms_list]);|
  ((), IdentityMapping( SymmetricGroup( [ 1 .. 3 ] ) ),
  [ (), (), (), (), (), () ])
\end{Verbatim}
 }

 

\subsection{\textcolor{Chapter }{ELM{\textunderscore}LIST (for IsRMSIsoByTriple)}}
\logpage{[ 14, 3, 3 ]}\nobreak
\hyperdef{L}{X81C4DE427D4A3D6C}{}
{\noindent\textcolor{FuncColor}{$\triangleright$\enspace\texttt{ELM{\textunderscore}LIST({\mdseries\slshape map, pos})\index{ELM{\textunderscore}LIST@\texttt{ELM{\textunderscore}LIST}!for IsRMSIsoByTriple}
\label{ELMuScoreLIST:for IsRMSIsoByTriple}
}\hfill{\scriptsize (operation)}}\\
\textbf{\indent Returns:\ }
 A component of an isomorphism of Rees (0\texttt{\symbol{45}})matrix semigroups
by triple. 



 \texttt{ELM{\textunderscore}LIST(\mbox{\texttt{\mdseries\slshape map}}, \mbox{\texttt{\mdseries\slshape i}})} returns the \texttt{i}th component of the Rees (0\texttt{\symbol{45}})matrix semigroup isomorphism
by triple \mbox{\texttt{\mdseries\slshape map}} when \texttt{i = 1, 2, 3}. 

 The components of an isomorphism of Rees (0\texttt{\symbol{45}})matrix
semigroups by triple are: 
\begin{enumerate}
\item  An isomorphism of the underlying graphs of the source and range of \mbox{\texttt{\mdseries\slshape map}}, respectively. 
\item  An isomorphism of the underlying groups of the source and range of \mbox{\texttt{\mdseries\slshape map}}, respectively. 
\item  An function from the union of the rows and columns of the source of \mbox{\texttt{\mdseries\slshape map}} to the underlying group of the range of \mbox{\texttt{\mdseries\slshape map}}. 
\end{enumerate}
 }

 

\subsection{\textcolor{Chapter }{CompositionMapping2 (for IsRMSIsoByTriple)}}
\logpage{[ 14, 3, 4 ]}\nobreak
\hyperdef{L}{X7A02528F8721F378}{}
{\noindent\textcolor{FuncColor}{$\triangleright$\enspace\texttt{CompositionMapping2({\mdseries\slshape map1, map2})\index{CompositionMapping2@\texttt{CompositionMapping2}!for IsRMSIsoByTriple}
\label{CompositionMapping2:for IsRMSIsoByTriple}
}\hfill{\scriptsize (operation)}}\\
\noindent\textcolor{FuncColor}{$\triangleright$\enspace\texttt{CompositionMapping2({\mdseries\slshape map1, map2})\index{CompositionMapping2@\texttt{CompositionMapping2}!for IsRZMSIsoByTriple}
\label{CompositionMapping2:for IsRZMSIsoByTriple}
}\hfill{\scriptsize (operation)}}\\
\textbf{\indent Returns:\ }
 A Rees (0\texttt{\symbol{45}})matrix semigroup by triple. 



 If \mbox{\texttt{\mdseries\slshape map1}} and \mbox{\texttt{\mdseries\slshape map2}} are isomorphisms of Rees matrix or 0\texttt{\symbol{45}}matrix semigroups
specified by triples and the range of \mbox{\texttt{\mdseries\slshape map2}} is contained in the source of \mbox{\texttt{\mdseries\slshape map1}}, then \texttt{CompositionMapping2(\mbox{\texttt{\mdseries\slshape map1}}, \mbox{\texttt{\mdseries\slshape map2}})} returns the isomorphism from \texttt{Source(\mbox{\texttt{\mdseries\slshape map2}})} to \texttt{Range(\mbox{\texttt{\mdseries\slshape map1}})} specified by the triple with components: 
\begin{enumerate}
\item  \texttt{\mbox{\texttt{\mdseries\slshape map1}}[1] * \mbox{\texttt{\mdseries\slshape map2}}[1]} 
\item  \texttt{\mbox{\texttt{\mdseries\slshape map1}}[2] * \mbox{\texttt{\mdseries\slshape map2}}[2]} 
\item  the componentwise product of \texttt{\mbox{\texttt{\mdseries\slshape map1}}[1] * \mbox{\texttt{\mdseries\slshape map2}}[3]} and \texttt{\mbox{\texttt{\mdseries\slshape map1}}[3] * \mbox{\texttt{\mdseries\slshape map2}}[2]}. 
\end{enumerate}
 
\begin{Verbatim}[commandchars=!@|,fontsize=\small,frame=single,label=Example]
  !gapprompt@gap>| !gapinput@R := ReesZeroMatrixSemigroup(Group([(1, 2, 3, 4)]),|
  !gapprompt@>| !gapinput@[[(1, 3)(2, 4), (1, 4, 3, 2), (), (1, 2, 3, 4), (1, 3)(2, 4), 0],|
  !gapprompt@>| !gapinput@ [(1, 4, 3, 2), 0, (), (1, 4, 3, 2), (1, 2, 3, 4), (1, 2, 3, 4)],|
  !gapprompt@>| !gapinput@ [(), (), (1, 4, 3, 2), (1, 2, 3, 4), 0, (1, 2, 3, 4)],|
  !gapprompt@>| !gapinput@ [(1, 2, 3, 4), (1, 4, 3, 2), (1, 2, 3, 4), 0, (), (1, 2, 3, 4)],|
  !gapprompt@>| !gapinput@ [(1, 3)(2, 4), (1, 2, 3, 4), 0, (), (1, 4, 3, 2), (1, 2, 3, 4)],|
  !gapprompt@>| !gapinput@ [0, (1, 2, 3, 4), (1, 2, 3, 4), (1, 2, 3, 4), (1, 2, 3, 4), ()]]);|
  <Rees 0-matrix semigroup 6x6 over Group([ (1,2,3,4) ])>
  !gapprompt@gap>| !gapinput@G := AutomorphismGroup(R);;|
  !gapprompt@gap>| !gapinput@G.2;|
  ((), IdentityMapping( Group( [ (1,2,3,4) ] ) ),
  [ (), (), (), (), (), (), (), (), (), (), (), () ])
  !gapprompt@gap>| !gapinput@G.3;|
  (( 2, 4)( 3, 5)( 8,10)( 9,11), GroupHomomorphismByImages( Group(
  [ (1,2,3,4) ] ), Group( [ (1,2,3,4) ] ), [ (1,2,3,4) ],
  [ (1,2,3,4) ] ), [ (), (1,3)(2,4), (1,3)(2,4), (1,3)(2,4),
    (1,3)(2,4), (1,3)(2,4), (), (1,3)(2,4), (1,3)(2,4), (1,3)(2,4),
    (1,3)(2,4), (1,3)(2,4) ])
  !gapprompt@gap>| !gapinput@CompositionMapping2(G.2, G.3);|
  (( 2, 4)( 3, 5)( 8,10)( 9,11), GroupHomomorphismByImages( Group(
  [ (1,2,3,4) ] ), Group( [ (1,2,3,4) ] ), [ (1,2,3,4) ],
  [ (1,2,3,4) ] ), [ (), (1,3)(2,4), (1,3)(2,4), (1,3)(2,4),
    (1,3)(2,4), (1,3)(2,4), (), (1,3)(2,4), (1,3)(2,4), (1,3)(2,4),
    (1,3)(2,4), (1,3)(2,4) ])
\end{Verbatim}
 }

 

\subsection{\textcolor{Chapter }{ImagesElm (for IsRMSIsoByTriple)}}
\logpage{[ 14, 3, 5 ]}\nobreak
\hyperdef{L}{X7F159C1179C93C11}{}
{\noindent\textcolor{FuncColor}{$\triangleright$\enspace\texttt{ImagesElm({\mdseries\slshape map, pt})\index{ImagesElm@\texttt{ImagesElm}!for IsRMSIsoByTriple}
\label{ImagesElm:for IsRMSIsoByTriple}
}\hfill{\scriptsize (operation)}}\\
\noindent\textcolor{FuncColor}{$\triangleright$\enspace\texttt{ImagesRepresentative({\mdseries\slshape map, pt})\index{ImagesRepresentative@\texttt{ImagesRepresentative}!for IsRMSIsoByTriple}
\label{ImagesRepresentative:for IsRMSIsoByTriple}
}\hfill{\scriptsize (operation)}}\\
\textbf{\indent Returns:\ }
 An element of a Rees (0\texttt{\symbol{45}})matrix semigroup or a list
containing such an element. 



 If \mbox{\texttt{\mdseries\slshape map}} is an isomorphism of Rees matrix or 0\texttt{\symbol{45}}matrix semigroups
specified by a triple and \mbox{\texttt{\mdseries\slshape pt}} is an element of the source of \mbox{\texttt{\mdseries\slshape map}}, then \texttt{ImagesRepresentative(\mbox{\texttt{\mdseries\slshape map}}, \mbox{\texttt{\mdseries\slshape pt}}) = \mbox{\texttt{\mdseries\slshape pt}} \texttt{\symbol{94}} \mbox{\texttt{\mdseries\slshape map}}} returns the image of \mbox{\texttt{\mdseries\slshape pt}} under \mbox{\texttt{\mdseries\slshape map}}. 

 The image of \mbox{\texttt{\mdseries\slshape pt}} under \mbox{\texttt{\mdseries\slshape map}} of \texttt{Range(\mbox{\texttt{\mdseries\slshape map}})} is the element with components: 
\begin{enumerate}
\item  \texttt{\mbox{\texttt{\mdseries\slshape pt}}[1] \texttt{\symbol{94}} \mbox{\texttt{\mdseries\slshape map}}[1]} 
\item  \texttt{(\mbox{\texttt{\mdseries\slshape pt}}[1] \texttt{\symbol{94}} \mbox{\texttt{\mdseries\slshape map}}[3]) * (\mbox{\texttt{\mdseries\slshape pt}}[2] \texttt{\symbol{94}} \mbox{\texttt{\mdseries\slshape map}}[2]) * (\mbox{\texttt{\mdseries\slshape pt}}[3] \texttt{\symbol{94}} \mbox{\texttt{\mdseries\slshape map}}[3]) \texttt{\symbol{94}} \texttt{\symbol{45}}1} 
\item  \texttt{\mbox{\texttt{\mdseries\slshape pt}}[3] \texttt{\symbol{94}} \mbox{\texttt{\mdseries\slshape map}}[1]}. 
\end{enumerate}
 \texttt{ImagesElm(\mbox{\texttt{\mdseries\slshape map}}, \mbox{\texttt{\mdseries\slshape pt}})} simply returns \texttt{[ImagesRepresentative(\mbox{\texttt{\mdseries\slshape map}}, \mbox{\texttt{\mdseries\slshape pt}})]}. 
\begin{Verbatim}[commandchars=!@|,fontsize=\small,frame=single,label=Example]
  !gapprompt@gap>| !gapinput@R := ReesZeroMatrixSemigroup(Group([(2, 8), (2, 8, 6)]),|
  !gapprompt@>| !gapinput@[[0, (2, 8), 0, 0, 0, (2, 8, 6)],|
  !gapprompt@>| !gapinput@ [(), 0, (2, 8, 6), (2, 6), (2, 6, 8), 0],|
  !gapprompt@>| !gapinput@ [(2, 8, 6), 0, (2, 6, 8), (2, 8), (), 0],|
  !gapprompt@>| !gapinput@ [(2, 8, 6), 0, (2, 6, 8), (2, 8), (), 0],|
  !gapprompt@>| !gapinput@ [0, (2, 8, 6), 0, 0, 0, (2, 8)],|
  !gapprompt@>| !gapinput@ [(2, 8, 6), 0, (2, 6, 8), (2, 8), (), 0]]);|
  <Rees 0-matrix semigroup 6x6 over Group([ (2,8), (2,8,6) ])>
  !gapprompt@gap>| !gapinput@map := RZMSIsoByTriple(R, R,|
  !gapprompt@>| !gapinput@[(), IdentityMapping(Group([(2, 8), (2, 8, 6)])),|
  !gapprompt@>| !gapinput@[(), (2, 6, 8), (), (), (), (2, 8, 6),|
  !gapprompt@>| !gapinput@ (2, 8, 6), (), (), (), (2, 6, 8), ()]]);;|
  !gapprompt@gap>| !gapinput@ImagesElm(map, RMSElement(R, 1, (2, 8), 2));|
  [ (1,(2,8),2) ]
\end{Verbatim}
 }

 

\subsection{\textcolor{Chapter }{CanonicalReesZeroMatrixSemigroup}}
\logpage{[ 14, 3, 6 ]}\nobreak
\hyperdef{L}{X8765885F784557B9}{}
{\noindent\textcolor{FuncColor}{$\triangleright$\enspace\texttt{CanonicalReesZeroMatrixSemigroup({\mdseries\slshape S})\index{CanonicalReesZeroMatrixSemigroup@\texttt{CanonicalReesZeroMatrixSemigroup}}
\label{CanonicalReesZeroMatrixSemigroup}
}\hfill{\scriptsize (attribute)}}\\
\noindent\textcolor{FuncColor}{$\triangleright$\enspace\texttt{CanonicalReesMatrixSemigroup({\mdseries\slshape S})\index{CanonicalReesMatrixSemigroup@\texttt{CanonicalReesMatrixSemigroup}}
\label{CanonicalReesMatrixSemigroup}
}\hfill{\scriptsize (attribute)}}\\
\textbf{\indent Returns:\ }
 A Rees zero matrix semigroup. 



 If \mbox{\texttt{\mdseries\slshape S}} is a Rees 0\texttt{\symbol{45}}matrix semigroup then \texttt{CanonicalReesZeroMatrixSemigroup} returns an isomorphic Rees 0\texttt{\symbol{45}}matrix semigroup \texttt{T} with the same \texttt{UnderlyingSemigroup} (\textbf{Reference: UnderlyingSemigroup for a Rees 0\texttt{\symbol{45}}matrix semigroup}) as \mbox{\texttt{\mdseries\slshape S}} but the \texttt{Matrix} (\textbf{Reference: Matrix}) of \texttt{T} has been canonicalized. The output \texttt{T} is canonical in the sense that for any two inputs which are isomorphic Rees
zero matrix semigroups the output of this function is the same.

 \texttt{CanonicalReesMatrixSemigroup} works the same but for Rees matrix semigroups. 
\begin{Verbatim}[commandchars=!@|,fontsize=\small,frame=single,label=Example]
  !gapprompt@gap>| !gapinput@S := ReesZeroMatrixSemigroup(SymmetricGroup(3),|
  !gapprompt@>| !gapinput@[[(), (1, 3, 2)], [(), ()]]);;|
  !gapprompt@gap>| !gapinput@T := CanonicalReesZeroMatrixSemigroup(S);|
  <Rees 0-matrix semigroup 2x2 over Sym( [ 1 .. 3 ] )>
  !gapprompt@gap>| !gapinput@Matrix(S);|
  [ [ (), (1,3,2) ], [ (), () ] ]
  !gapprompt@gap>| !gapinput@Matrix(T);|
  [ [ (), () ], [ (), (1,2,3) ] ]
\end{Verbatim}
 }

  
\subsection{\textcolor{Chapter }{ Operators for isomorphisms of Rees (0\texttt{\symbol{45}})matrix semigroups }}\label{Operators for isomorphisms of Rees (0-)matrix semigroups}
\logpage{[ 14, 3, 7 ]}
\hyperdef{L}{X7ED8BF227F4229E2}{}
{
  
\begin{description}
\item[{\texttt{\mbox{\texttt{\mdseries\slshape map}}[\mbox{\texttt{\mdseries\slshape i}}]}}]  \index{ELM_LIST@\texttt{ELM{\textunderscore}LIST} (for Rees (0\texttt{\symbol{45}})matrix semigroup isomorphisms by triples)} \texttt{\mbox{\texttt{\mdseries\slshape map}}[i]} returns the \mbox{\texttt{\mdseries\slshape i}}th component of the Rees (0\texttt{\symbol{45}})matrix semigroup isomorphism
by triple \mbox{\texttt{\mdseries\slshape map}} when \texttt{\mbox{\texttt{\mdseries\slshape i}} = 1, 2, 3}; see \texttt{ELM{\textunderscore}LIST} (\ref{ELMuScoreLIST:for IsRMSIsoByTriple}). 
\item[{\texttt{\mbox{\texttt{\mdseries\slshape map1}} * \mbox{\texttt{\mdseries\slshape map2}}}}]  \index{*@\texttt{ * } (for Rees (0\texttt{\symbol{45}})matrix semigroup isomorphisms by triples)} returns the composition of \mbox{\texttt{\mdseries\slshape map2}} and \mbox{\texttt{\mdseries\slshape map1}}; see \texttt{CompositionMapping2} (\ref{CompositionMapping2:for IsRMSIsoByTriple}). 
\item[{\texttt{\mbox{\texttt{\mdseries\slshape map1}} {\textless} \mbox{\texttt{\mdseries\slshape map2}}}}]  \index{<@\texttt{{\textless}} (for Rees (0\texttt{\symbol{45}})matrix semigroup isomorphisms by triples)} returns \texttt{true} if \mbox{\texttt{\mdseries\slshape map1}} is lexicographically less than \mbox{\texttt{\mdseries\slshape map2}}. 
\item[{\texttt{\mbox{\texttt{\mdseries\slshape map1}} = \mbox{\texttt{\mdseries\slshape map2}}}}]  \index{=@\texttt{ = } (for Rees (0\texttt{\symbol{45}})matrix semigroup isomorphisms by triples)} returns \texttt{true} if the Rees (0\texttt{\symbol{45}})matrix semigroup isomorphisms by triple \mbox{\texttt{\mdseries\slshape map1}} and \mbox{\texttt{\mdseries\slshape map2}} have equal source and range, and are equal as functions, and \texttt{false} otherwise. 

 It is possible for \mbox{\texttt{\mdseries\slshape map1}} and \mbox{\texttt{\mdseries\slshape map2}} to be equal but to have distinct components. 
\item[{\texttt{\mbox{\texttt{\mdseries\slshape pt}} \texttt{\symbol{94}} \mbox{\texttt{\mdseries\slshape map}}}}]  \index{^@\texttt{ \texttt{\symbol{94}} } (for Rees (0\texttt{\symbol{45}})matrix semigroup isomorphisms by triples)} returns the image of the element \mbox{\texttt{\mdseries\slshape pt}} of the source of \mbox{\texttt{\mdseries\slshape map}} under the isomorphism \mbox{\texttt{\mdseries\slshape map}}; see \texttt{ImagesElm} (\ref{ImagesElm:for IsRMSIsoByTriple}). 
\end{description}
 }

 }

   }

 
\chapter{\textcolor{Chapter }{ Finitely presented semigroups and Tietze transformations }}\label{Finitely presented semigroups and Tietze transformations}
\logpage{[ 15, 0, 0 ]}
\hyperdef{L}{X7F11EF307D4F409B}{}
{
  In this chapter we describe the functions implemented in \textsf{Semigroups} that extend the features available in \textsf{GAP} for dealing with finitely presented semigroups and monoids. 

 Section \ref{Changing representation for words and strings} (written by Maria Tsalakou and Murray Whyte) and Section \ref{Helper functions} (written by Luke Elliott) demonstrate a number of helper functions that allow
the user to convert between different representations of words and relations. 

 In the later sections, written and implemented by Ben Spiers and Tom
Conti\texttt{\symbol{45}}Leslie, we describe how to change the relations of a
finitely presented semigroup either manually or automatically using Tietze
transformations (which is abbreviated to \textsc{Stz}. 

 
\section{\textcolor{Chapter }{ Changing representation for words and strings }}\label{Changing representation for words and strings}
\logpage{[ 15, 1, 0 ]}
\hyperdef{L}{X87CDF7DF7E47F8FB}{}
{
  This section contains various methods for dealing with words, which for these
purposes are lists of positive integers. 

\subsection{\textcolor{Chapter }{WordToString (for a string and a list)}}
\logpage{[ 15, 1, 1 ]}\nobreak
\hyperdef{L}{X7CACE2997C1D243A}{}
{\noindent\textcolor{FuncColor}{$\triangleright$\enspace\texttt{WordToString({\mdseries\slshape A, w})\index{WordToString@\texttt{WordToString}!for a string and a list}
\label{WordToString:for a string and a list}
}\hfill{\scriptsize (operation)}}\\
\textbf{\indent Returns:\ }
 A string. 



 Returns the word \mbox{\texttt{\mdseries\slshape w}}, in the form of a string of letters of the alphabet \mbox{\texttt{\mdseries\slshape A}}. The alphabet is given as a string containing its members. 
\begin{Verbatim}[commandchars=!@|,fontsize=\small,frame=single,label=Example]
  !gapprompt@gap>| !gapinput@WordToString("abcd", [4, 2, 3, 1, 1, 4, 2, 3]);|
  "dbcaadbc"
\end{Verbatim}
 }

 

\subsection{\textcolor{Chapter }{RandomWord (for two integers)}}
\logpage{[ 15, 1, 2 ]}\nobreak
\hyperdef{L}{X7C81FB0F7BF3D4A3}{}
{\noindent\textcolor{FuncColor}{$\triangleright$\enspace\texttt{RandomWord({\mdseries\slshape l, n})\index{RandomWord@\texttt{RandomWord}!for two integers}
\label{RandomWord:for two integers}
}\hfill{\scriptsize (operation)}}\\
\textbf{\indent Returns:\ }
 A word. 



 Returns a random word of length \mbox{\texttt{\mdseries\slshape l}} over \mbox{\texttt{\mdseries\slshape n}} letters. 
\begin{Verbatim}[commandchars=!@|,fontsize=\small,frame=single,label=Example]
  !gapprompt@gap>| !gapinput@RandomWord(8, 5);|
  [ 2, 4, 3, 4, 5, 3, 3, 2 ]
  !gapprompt@gap>| !gapinput@RandomWord(8, 5);|
  [ 3, 3, 5, 5, 5, 4, 4, 5 ]
  !gapprompt@gap>| !gapinput@RandomWord(8, 4);|
  [ 1, 4, 1, 1, 3, 3, 4, 4 ]
\end{Verbatim}
 }

 

\subsection{\textcolor{Chapter }{StandardiseWord}}
\logpage{[ 15, 1, 3 ]}\nobreak
\hyperdef{L}{X7B338B917D72FBED}{}
{\noindent\textcolor{FuncColor}{$\triangleright$\enspace\texttt{StandardiseWord({\mdseries\slshape w})\index{StandardiseWord@\texttt{StandardiseWord}}
\label{StandardiseWord}
}\hfill{\scriptsize (operation)}}\\
\noindent\textcolor{FuncColor}{$\triangleright$\enspace\texttt{StandardizeWord({\mdseries\slshape w})\index{StandardizeWord@\texttt{StandardizeWord}}
\label{StandardizeWord}
}\hfill{\scriptsize (operation)}}\\
\textbf{\indent Returns:\ }
 A list of positive integers. 



 This function takes a word \mbox{\texttt{\mdseries\slshape w}}, consisting of \texttt{n} distinct positive integers and returns a word \texttt{s} where the characters of \texttt{s} correspond to those of \mbox{\texttt{\mdseries\slshape w}} in order of first appearance.

 The word \mbox{\texttt{\mdseries\slshape w}} is changed in\texttt{\symbol{45}}place into word \texttt{s}. 
\begin{Verbatim}[commandchars=!@|,fontsize=\small,frame=single,label=Example]
  !gapprompt@gap>| !gapinput@w := [3, 1, 2];|
  [ 3, 1, 2 ]
  !gapprompt@gap>| !gapinput@StandardiseWord(w);|
  [ 1, 2, 3 ]
  !gapprompt@gap>| !gapinput@w;|
  [ 1, 2, 3 ]
  !gapprompt@gap>| !gapinput@w := [4, 2, 10, 2];|
  [ 4, 2, 10, 2 ]
  !gapprompt@gap>| !gapinput@StandardiseWord(w);|
  [ 1, 2, 3, 2 ]
\end{Verbatim}
 }

 

\subsection{\textcolor{Chapter }{StringToWord (for a string)}}
\logpage{[ 15, 1, 4 ]}\nobreak
\hyperdef{L}{X7B07470F86FDC7BA}{}
{\noindent\textcolor{FuncColor}{$\triangleright$\enspace\texttt{StringToWord({\mdseries\slshape s})\index{StringToWord@\texttt{StringToWord}!for a string}
\label{StringToWord:for a string}
}\hfill{\scriptsize (operation)}}\\
\textbf{\indent Returns:\ }
 A list of positive integers. 



 This function takes a string \mbox{\texttt{\mdseries\slshape s}}, consisting of \texttt{n} distinct positive integers and returns a word \texttt{w} (i.e. a list of positive integers) over the alphabet \texttt{[1 .. n]}. The positive integers of \texttt{w} correspond to the characters of \mbox{\texttt{\mdseries\slshape s}}, in order of first appearance. 
\begin{Verbatim}[commandchars=!@|,fontsize=\small,frame=single,label=Example]
  !gapprompt@gap>| !gapinput@w := "abac";|
  "abac"
  !gapprompt@gap>| !gapinput@StringToWord(w);|
  [ 1, 2, 1, 3 ]
  !gapprompt@gap>| !gapinput@w := "ccala";|
  "ccala"
  !gapprompt@gap>| !gapinput@StringToWord(w);|
  [ 1, 1, 2, 3, 2 ]
  !gapprompt@gap>| !gapinput@w := "a1b5";|
  "a1b5"
  !gapprompt@gap>| !gapinput@StringToWord(w);|
  [ 1, 2, 3, 4 ]
\end{Verbatim}
 }

 }

 
\section{\textcolor{Chapter }{ Helper functions }}\label{Helper functions}
\logpage{[ 15, 2, 0 ]}
\hyperdef{L}{X7BD4785D8488BAD5}{}
{
  This section describes operations implemented in \textsf{Semigroups} that are designed to interact with standard \textsf{GAP} methods for creating finitely presented semigroups and monoids (see  (\textbf{Reference: Finitely Presented Semigroups and Monoids})). 

\subsection{\textcolor{Chapter }{ParseRelations}}
\logpage{[ 15, 2, 1 ]}\nobreak
\hyperdef{L}{X7C2FCCA487DFC84C}{}
{\noindent\textcolor{FuncColor}{$\triangleright$\enspace\texttt{ParseRelations({\mdseries\slshape gens, rels})\index{ParseRelations@\texttt{ParseRelations}}
\label{ParseRelations}
}\hfill{\scriptsize (operation)}}\\
\textbf{\indent Returns:\ }
A list of pairs of semigroup or monoid elements.



 \texttt{ParseRelations} converts a string describing relations for a semigroup or monoid to the list
of pairs of semigroup or monoid elements it represents. Any white space given
is ignored. The output list is then compatible with other \textsf{GAP} functions. In the below examples we see free semigroups and monoids being
directly quotiented by the output of the \texttt{ParseRelations} function. 

 The argument \mbox{\texttt{\mdseries\slshape gens}} must be a list of generators for a free semigroup, each being a single
alphabet letter (in upper or lower case). The argument \mbox{\texttt{\mdseries\slshape rels}} must be string that lists the equalities desired.

 To take a quotient of a free monoid, it is necessary to use \texttt{GeneratorsOfMonoid} (\textbf{Reference: GeneratorsOfMonoid}) as the 1st argument to \texttt{ParseRelations} and the identity may appear as \texttt{1} in the specified relations. 
\begin{Verbatim}[commandchars=!@|,fontsize=\small,frame=single,label=Example]
  !gapprompt@gap>| !gapinput@f := FreeSemigroup("x", "y", "z");;|
  !gapprompt@gap>| !gapinput@AssignGeneratorVariables(f);|
  !gapprompt@gap>| !gapinput@ParseRelations([x, y, z], "  x=(y^2z) ^2x, y=xxx , z=y^3");|
  [ [ x, (y^2*z)^2*x ], [ y, x^3 ], [ z, y^3 ] ]
  !gapprompt@gap>| !gapinput@r := ParseRelations([x, y, z], "  x=(y^2z)^2x, y=xxx=z , z=y^3");|
  [ [ x, (y^2*z)^2*x ], [ y, x^3 ], [ x^3, z ], [ z, y^3 ] ]
  !gapprompt@gap>| !gapinput@f / r;|
  <fp semigroup with 3 generators and 4 relations of length 23>
  !gapprompt@gap>| !gapinput@f2 := FreeSemigroup("a");|
  <free semigroup on the generators [ a ]>
  !gapprompt@gap>| !gapinput@f2 / ParseRelations(GeneratorsOfSemigroup(f2), "a = a^2");|
  <fp semigroup with 1 generator and 1 relation of length 4>
  !gapprompt@gap>| !gapinput@FreeMonoidAndAssignGeneratorVars("a", "b")|
  !gapprompt@>| !gapinput@/ ParseRelations([a, b], "a^2=1,b^3=1,(ab)^3=1");|
  <fp monoid with 2 generators and 3 relations of length 13>
\end{Verbatim}
 }

 

\subsection{\textcolor{Chapter }{ElementOfFpSemigroup}}
\logpage{[ 15, 2, 2 ]}\nobreak
\hyperdef{L}{X847012347856C55E}{}
{\noindent\textcolor{FuncColor}{$\triangleright$\enspace\texttt{ElementOfFpSemigroup({\mdseries\slshape S, word})\index{ElementOfFpSemigroup@\texttt{ElementOfFpSemigroup}}
\label{ElementOfFpSemigroup}
}\hfill{\scriptsize (operation)}}\\
\textbf{\indent Returns:\ }
An element of the fp semigroup \mbox{\texttt{\mdseries\slshape S}}.



 When \mbox{\texttt{\mdseries\slshape S}} is a finitely presented semigroup and \mbox{\texttt{\mdseries\slshape word}} is an associative word in the associated free semigroup (see \texttt{IsAssocWord} (\textbf{Reference: IsAssocWord})), this returns the fp semigroup element with representative \mbox{\texttt{\mdseries\slshape word}}. 

 This function is just a short form of the \textsf{GAP} library implementation of \texttt{ElementOfFpSemigroup} (\textbf{Reference: ElementOfFpSemigroup}) which does not require retrieving an element family. 
\begin{Verbatim}[commandchars=!@|,fontsize=\small,frame=single,label=Example]
  !gapprompt@gap>| !gapinput@f := FreeSemigroup("x", "y");;|
  !gapprompt@gap>| !gapinput@AssignGeneratorVariables(f);|
  !gapprompt@gap>| !gapinput@s := f / [[x * x, x], [y * y, y]];|
  <fp semigroup with 2 generators and 2 relations of length 8>
  !gapprompt@gap>| !gapinput@a := ElementOfFpSemigroup(s, x * y);|
  x*y
  !gapprompt@gap>| !gapinput@b := ElementOfFpSemigroup(s, x * y * y);|
  x*y^2
  !gapprompt@gap>| !gapinput@a in s;|
  true
  !gapprompt@gap>| !gapinput@a = b;|
  true
\end{Verbatim}
 }

 

\subsection{\textcolor{Chapter }{ElementOfFpMonoid}}
\logpage{[ 15, 2, 3 ]}\nobreak
\hyperdef{L}{X82B7A51B7FE90486}{}
{\noindent\textcolor{FuncColor}{$\triangleright$\enspace\texttt{ElementOfFpMonoid({\mdseries\slshape M, word})\index{ElementOfFpMonoid@\texttt{ElementOfFpMonoid}}
\label{ElementOfFpMonoid}
}\hfill{\scriptsize (operation)}}\\
\textbf{\indent Returns:\ }
An element of the fp monoid \mbox{\texttt{\mdseries\slshape M}}.



 When \mbox{\texttt{\mdseries\slshape M}} is a finitely presented monoid and \mbox{\texttt{\mdseries\slshape word}} is an associative word in the associated free monoid (see \texttt{IsAssocWord} (\textbf{Reference: IsAssocWord})), this returns the fp monoid element with representative \mbox{\texttt{\mdseries\slshape word}}. 

 This is analogous to \texttt{ElementOfFpSemigroup} (\ref{ElementOfFpSemigroup}). }

 

\subsection{\textcolor{Chapter }{FreeMonoidAndAssignGeneratorVars}}
\logpage{[ 15, 2, 4 ]}\nobreak
\hyperdef{L}{X7C3837FA83BE9CD9}{}
{\noindent\textcolor{FuncColor}{$\triangleright$\enspace\texttt{FreeMonoidAndAssignGeneratorVars({\mdseries\slshape arg...})\index{FreeMonoidAndAssignGeneratorVars@\texttt{FreeMonoidAndAssignGeneratorVars}}
\label{FreeMonoidAndAssignGeneratorVars}
}\hfill{\scriptsize (function)}}\\
\noindent\textcolor{FuncColor}{$\triangleright$\enspace\texttt{FreeSemigroupAndAssignGeneratorVars({\mdseries\slshape arg...})\index{FreeSemigroupAndAssignGeneratorVars@\texttt{FreeSemigroupAndAssignGeneratorVars}}
\label{FreeSemigroupAndAssignGeneratorVars}
}\hfill{\scriptsize (function)}}\\
\textbf{\indent Returns:\ }
A free semigroup or monoid.



 \texttt{FreeMonoidAndAssignGeneratorVars} is synonym with: 
\begin{Verbatim}[commandchars=!@|,fontsize=\small,frame=single,label=Example]
        FreeMonoid(arg...);
        AssignGeneratorVariables(last);
\end{Verbatim}
 These functions exist so that the \texttt{String} method for a finitely presented semigroup or monoid to be valid \textsf{GAP} input which can be used to reconstruct the semigroup or monoid. 
\begin{Verbatim}[commandchars=!@|,fontsize=\small,frame=single,label=Example]
  !gapprompt@gap>| !gapinput@F := FreeSemigroupAndAssignGeneratorVars("x", "y");;|
  !gapprompt@gap>| !gapinput@IsBound(x);|
  true
  !gapprompt@gap>| !gapinput@IsBound(y);|
  true
\end{Verbatim}
 }

 

\subsection{\textcolor{Chapter }{IsSubsemigroupOfFpMonoid}}
\logpage{[ 15, 2, 5 ]}\nobreak
\hyperdef{L}{X7FF4A1CF79799314}{}
{\noindent\textcolor{FuncColor}{$\triangleright$\enspace\texttt{IsSubsemigroupOfFpMonoid({\mdseries\slshape S})\index{IsSubsemigroupOfFpMonoid@\texttt{IsSubsemigroupOfFpMonoid}}
\label{IsSubsemigroupOfFpMonoid}
}\hfill{\scriptsize (property)}}\\
\textbf{\indent Returns:\ }
\texttt{true} or \texttt{false}.



 This property is \texttt{true} if the object \mbox{\texttt{\mdseries\slshape S}} is a subsemigroup of an fp monoid, and \texttt{false} otherwise. This property is just a synonym for \texttt{IsSemigroup} (\textbf{Reference: IsSemigroup}) and \texttt{IsElementOfFpMonoidCollection}. 
\begin{Verbatim}[commandchars=!@|,fontsize=\small,frame=single,label=Example]
  !gapprompt@gap>| !gapinput@F := FreeSemigroup("a", "b");|
  <free semigroup on the generators [ a, b ]>
  !gapprompt@gap>| !gapinput@AssignGeneratorVariables(F);|
  !gapprompt@gap>| !gapinput@R := [[a ^ 3, a], [b ^ 2, b], [(a * b) ^ 2, a]];|
  [ [ a^3, a ], [ b^2, b ], [ (a*b)^2, a ] ]
  !gapprompt@gap>| !gapinput@S := F / R;|
  <fp semigroup with 2 generators and 3 relations of length 14>
  !gapprompt@gap>| !gapinput@IsSubsemigroupOfFpMonoid(S);|
  false
  !gapprompt@gap>| !gapinput@map := EmbeddingFpMonoid(S);|
  <fp semigroup with 2 generators and 3 relations of length 14> ->
  <fp monoid with 2 generators and 3 relations of length 14>
  !gapprompt@gap>| !gapinput@IsSubsemigroupOfFpMonoid(Image(map));|
  true
  !gapprompt@gap>| !gapinput@IsSubsemigroupOfFpMonoid(Range(map));|
  true
\end{Verbatim}
 }

 }

 
\section{\textcolor{Chapter }{ Creating Tietze transformation objects }}\label{Creating Tietze transformation objects}
\logpage{[ 15, 3, 0 ]}
\hyperdef{L}{X8379F50C83FA5088}{}
{
  It is possible to use \textsf{GAP} to create finitely presented semigroups without the \textsf{Semigroups} package, by creating a free semigroup, then quotienting by a list of
relations. This is described in the reference manual ( (\textbf{Reference: Finitely Presented Semigroups and Monoids})). 

 However, finitely presented semigroups do not allow for their relations to be
simplified, so in the following sections, we describe how to create and modify
the semigroup Tietze (\texttt{IsStzPresentation} (\ref{IsStzPresentation})) object associated with an fp semigroup. This object can be automatically
simplified, or the user can manually apply Tietze transformations to add or
remove relations or generators in the presentation. 

 This object is analogous to \texttt{PresentationFpGroup} (\textbf{Reference: PresentationFpGroup}) implemented for fp groups in the main \textsf{GAP} distribution ( (\textbf{Reference: Presentations and Tietze Transformations})), but its features are semigroup\texttt{\symbol{45}}specific. Most of the
functions used to create, view and manipulate semigroup Tietze objects are
prefixed with \textsc{Stz}. 

\subsection{\textcolor{Chapter }{StzPresentation}}
\logpage{[ 15, 3, 1 ]}\nobreak
\hyperdef{L}{X7B505DA083E2EC0C}{}
{\noindent\textcolor{FuncColor}{$\triangleright$\enspace\texttt{StzPresentation({\mdseries\slshape S})\index{StzPresentation@\texttt{StzPresentation}}
\label{StzPresentation}
}\hfill{\scriptsize (operation)}}\\
\textbf{\indent Returns:\ }
 A semigroup Tietze (Stz) object. 



 If \mbox{\texttt{\mdseries\slshape s}} is an fp semigroup (\texttt{IsFpSemigroup} (\textbf{Reference: IsFpSemigroup})), then this function returns a modifiable object representing the generators
and relations of \mbox{\texttt{\mdseries\slshape s}}. 
\begin{Verbatim}[commandchars=!@|,fontsize=\small,frame=single,label=Example]
  !gapprompt@gap>| !gapinput@F := FreeSemigroup("a", "b", "c");;|
  !gapprompt@gap>| !gapinput@AssignGeneratorVariables(F);;|
  !gapprompt@gap>| !gapinput@S := F / [[a * b, c], [b * c, a], [c * a, b]];|
  <fp semigroup with 3 generators and 3 relations of length 12>
  !gapprompt@gap>| !gapinput@stz := StzPresentation(S);|
  <fp semigroup presentation with 3 generators and 3 relations
   with length 12>
\end{Verbatim}
 }

 

\subsection{\textcolor{Chapter }{IsStzPresentation}}
\logpage{[ 15, 3, 2 ]}\nobreak
\hyperdef{L}{X7B86C70F84BCF8BD}{}
{\noindent\textcolor{FuncColor}{$\triangleright$\enspace\texttt{IsStzPresentation({\mdseries\slshape stz})\index{IsStzPresentation@\texttt{IsStzPresentation}}
\label{IsStzPresentation}
}\hfill{\scriptsize (filter)}}\\
\textbf{\indent Returns:\ }
 \texttt{true} or \texttt{false}. 



 Every semigroup Tietze object is an element of the category \texttt{IsStzPresentation}. Internally, each Stz object contains a list of generators (each represented
as a string) and a list of relations (each represented as a pair of \texttt{LetterRep} words, see \texttt{LetterRepAssocWord} (\textbf{Reference: LetterRepAssocWord})). These generator and relation lists can be modified using Tietze
transformations (\ref{Changing Tietze transformation objects}). 

 When a \texttt{IsStzPresentation} object \mbox{\texttt{\mdseries\slshape stz}} is created from an fp semigroup \texttt{s} using \texttt{stz := StzPresentation(s)}, the generators and relations of \mbox{\texttt{\mdseries\slshape stz}} are initially equal to the generators and relations of \texttt{s}. However, as the Stz object \mbox{\texttt{\mdseries\slshape stz}} is modified, these lists may change, and their current state can be viewed
using \texttt{GeneratorsOfStzPresentation} (\ref{GeneratorsOfStzPresentation}) and \texttt{RelationsOfStzPresentation} (\ref{RelationsOfStzPresentation}). 
\begin{Verbatim}[commandchars=!@|,fontsize=\small,frame=single,label=Example]
  !gapprompt@gap>| !gapinput@F := FreeSemigroup("a", "b", "c");;|
  !gapprompt@gap>| !gapinput@AssignGeneratorVariables(F);;|
  !gapprompt@gap>| !gapinput@S := F / [[a * b, c], [b * c, a], [c * a, b]];|
  <fp semigroup with 3 generators and 3 relations of length 12>
  !gapprompt@gap>| !gapinput@stz := StzPresentation(S);|
  <fp semigroup presentation with 3 generators and 3 relations
   with length 12>
  !gapprompt@gap>| !gapinput@IsStzPresentation(stz);|
  true
\end{Verbatim}
 }

 

\subsection{\textcolor{Chapter }{GeneratorsOfStzPresentation}}
\logpage{[ 15, 3, 3 ]}\nobreak
\hyperdef{L}{X7F399C5982227D31}{}
{\noindent\textcolor{FuncColor}{$\triangleright$\enspace\texttt{GeneratorsOfStzPresentation({\mdseries\slshape stz})\index{GeneratorsOfStzPresentation@\texttt{GeneratorsOfStzPresentation}}
\label{GeneratorsOfStzPresentation}
}\hfill{\scriptsize (attribute)}}\\
\textbf{\indent Returns:\ }
 A list of strings. 



 If \mbox{\texttt{\mdseries\slshape stz}} is an \texttt{StzPresentation} (\ref{StzPresentation}) object, then \texttt{GeneratorsOfStzPresentation} will return (as strings) the generators of the fp semigroup that the
presentation was created from. In the \texttt{StzPresentation} (\ref{StzPresentation}) object, it is only necessary to know how many generators there are, but for
the purposes of representing generators and relations of the presentation
object and building a new fp semigroup from the object, the strings
representing the generators are stored. 

 As Tietze transformations are performed on \mbox{\texttt{\mdseries\slshape stz}}, the generators will change, but the labels will remain as close to the
original labels as possible, so that if a generator in the fp semigroup
obtained from the presentation is the same as a generator in the original fp
semigroup, then they should have the same label. 
\begin{Verbatim}[commandchars=!@|,fontsize=\small,frame=single,label=Example]
  !gapprompt@gap>| !gapinput@F := FreeSemigroup("a", "b", "c");|
  <free semigroup on the generators [ a, b, c ]>
  !gapprompt@gap>| !gapinput@T := F / [[F.1, F.2 ^ 5 * F.3],|
  !gapprompt@>| !gapinput@             [F.2 ^ 6, F.2 ^ 3]];|
  <fp semigroup with 3 generators and 2 relations of length 19>
  !gapprompt@gap>| !gapinput@stz := StzPresentation(T);|
  <fp semigroup presentation with 3 generators and 2 relations
   with length 19>
  !gapprompt@gap>| !gapinput@GeneratorsOfStzPresentation(stz);|
  [ "a", "b", "c" ]
\end{Verbatim}
 }

 

\subsection{\textcolor{Chapter }{RelationsOfStzPresentation}}
\logpage{[ 15, 3, 4 ]}\nobreak
\hyperdef{L}{X7DCFA5D8834C31BF}{}
{\noindent\textcolor{FuncColor}{$\triangleright$\enspace\texttt{RelationsOfStzPresentation({\mdseries\slshape stz})\index{RelationsOfStzPresentation@\texttt{RelationsOfStzPresentation}}
\label{RelationsOfStzPresentation}
}\hfill{\scriptsize (attribute)}}\\
\textbf{\indent Returns:\ }
 A list of pairs of words in \texttt{LetterRep} (\texttt{LetterRepAssocWord} (\textbf{Reference: LetterRepAssocWord})) form. 



 If \mbox{\texttt{\mdseries\slshape stz}} is an \texttt{StzPresentation} (\ref{StzPresentation}) object, then \texttt{RelationsOfStzPresentation} will return in letter rep form the current relations of the presentation
object. When the presentation object is first created, these will be the \texttt{LetterRep} forms of the relations of the fp semigroup object that is used to create \mbox{\texttt{\mdseries\slshape stz}}. As Tietze transformations are performed on the presentation object, the
relations returned by this function will change to reflect the
transformations. 
\begin{Verbatim}[commandchars=!@|,fontsize=\small,frame=single,label=Example]
  !gapprompt@gap>| !gapinput@F := FreeSemigroup("a", "b", "c");|
  <free semigroup on the generators [ a, b, c ]>
  !gapprompt@gap>| !gapinput@T := F / [[F.1, F.2 ^ 5 * F.3],|
  !gapprompt@>| !gapinput@             [F.2 ^ 6, F.2 ^ 3]];|
  <fp semigroup with 3 generators and 2 relations of length 19>
  !gapprompt@gap>| !gapinput@stz := StzPresentation(T);|
  <fp semigroup presentation with 3 generators and 2 relations
   with length 19>
  !gapprompt@gap>| !gapinput@RelationsOfStzPresentation(stz);|
  [ [ [ 1 ], [ 2, 2, 2, 2, 2, 3 ] ],
    [ [ 2, 2, 2, 2, 2, 2 ], [ 2, 2, 2 ] ] ]
\end{Verbatim}
 }

 

\subsection{\textcolor{Chapter }{UnreducedFpSemigroup (for a presentation)}}
\logpage{[ 15, 3, 5 ]}\nobreak
\hyperdef{L}{X79D7439679E96F28}{}
{\noindent\textcolor{FuncColor}{$\triangleright$\enspace\texttt{UnreducedFpSemigroup({\mdseries\slshape stz})\index{UnreducedFpSemigroup@\texttt{UnreducedFpSemigroup}!for a presentation}
\label{UnreducedFpSemigroup:for a presentation}
}\hfill{\scriptsize (attribute)}}\\
\textbf{\indent Returns:\ }
 An fp semigroup. 



 If \mbox{\texttt{\mdseries\slshape stz}} is an \texttt{StzPresentation} (\ref{StzPresentation}) object, then \texttt{UnreducedFpSemigroup} will return the fp semigroup that was used to create \mbox{\texttt{\mdseries\slshape stz}} using \texttt{StzPresentation} (\ref{StzPresentation}). 
\begin{Verbatim}[commandchars=!@|,fontsize=\small,frame=single,label=Example]
  !gapprompt@gap>| !gapinput@F := FreeSemigroup("a", "b", "c");|
  <free semigroup on the generators [ a, b, c ]>
  !gapprompt@gap>| !gapinput@T := F / [[F.1, F.2 ^ 5 * F.3],|
  !gapprompt@>| !gapinput@             [F.2 ^ 6, F.2 ^ 3]];|
  <fp semigroup with 3 generators and 2 relations of length 19>
  !gapprompt@gap>| !gapinput@stz := StzPresentation(T);|
  <fp semigroup presentation with 3 generators and 2 relations
   with length 19>
  !gapprompt@gap>| !gapinput@UnreducedFpSemigroup(stz) = T;|
  true
\end{Verbatim}
 }

 

\subsection{\textcolor{Chapter }{Length}}
\logpage{[ 15, 3, 6 ]}\nobreak
\hyperdef{L}{X780769238600AFD1}{}
{\noindent\textcolor{FuncColor}{$\triangleright$\enspace\texttt{Length({\mdseries\slshape stz})\index{Length@\texttt{Length}}
\label{Length}
}\hfill{\scriptsize (operation)}}\\
\textbf{\indent Returns:\ }
 A non\texttt{\symbol{45}}negative integer. 



 If \mbox{\texttt{\mdseries\slshape stz}} is an \texttt{StzPresentation} (\ref{StzPresentation}) object, then the \texttt{Length} of the object is defined as the number of generators plus the lengths of each
word in each relation of \mbox{\texttt{\mdseries\slshape stz}}. 
\begin{Verbatim}[commandchars=!@|,fontsize=\small,frame=single,label=Example]
  !gapprompt@gap>| !gapinput@F := FreeSemigroup("a", "b", "c");|
  <free semigroup on the generators [ a, b, c ]>
  !gapprompt@gap>| !gapinput@T := F / [[F.1, F.2 ^ 5 * F.3],|
  !gapprompt@>| !gapinput@             [F.2 ^ 6, F.2 ^ 3], [F.2 ^ 2, F.2]];|
  <fp semigroup with 3 generators and 3 relations of length 22>
  !gapprompt@gap>| !gapinput@stz := StzPresentation(T);|
  <fp semigroup presentation with 3 generators and 3 relations
   with length 22>
  !gapprompt@gap>| !gapinput@Length(stz);|
  22
\end{Verbatim}
 }

 }

 
\section{\textcolor{Chapter }{ Printing Tietze transformation objects }}\label{Printing Tietze transformation objects}
\logpage{[ 15, 4, 0 ]}
\hyperdef{L}{X84E3986C7EA62A06}{}
{
  Since the relations are stored as flat lists of numbers, there are several
methods installed to print the presentations in more
user\texttt{\symbol{45}}friendly forms. 

 All printing methods in this section are displayed as information (\texttt{Info} (\textbf{Reference: Info})) in the class \texttt{InfoFpSemigroup} at level 1. Setting \texttt{SetInfoLevevl(InfoFpSemigroup, 0)} will suppress the messages, while any higher number will display them. 

 

\subsection{\textcolor{Chapter }{StzPrintRelations}}
\logpage{[ 15, 4, 1 ]}\nobreak
\hyperdef{L}{X80F18CED8572886B}{}
{\noindent\textcolor{FuncColor}{$\triangleright$\enspace\texttt{StzPrintRelations({\mdseries\slshape stz[, list]})\index{StzPrintRelations@\texttt{StzPrintRelations}}
\label{StzPrintRelations}
}\hfill{\scriptsize (operation)}}\\


 If \mbox{\texttt{\mdseries\slshape stz}} is an \texttt{StzPresentation} (\ref{StzPresentation}) object and \mbox{\texttt{\mdseries\slshape list}} is a list of positive integers, then \texttt{StzPrintRelations} prints for each \mbox{\texttt{\mdseries\slshape i}} in \mbox{\texttt{\mdseries\slshape list}} the \mbox{\texttt{\mdseries\slshape i}}th relation to the console in terms of the stored labels for the generators
(that is, as words over the alphabet consisting of the generators of \mbox{\texttt{\mdseries\slshape stz}}).

 If \mbox{\texttt{\mdseries\slshape list}} is not specified then \texttt{StzPrintRelations} prints all relations of \mbox{\texttt{\mdseries\slshape stz}} in order. 
\begin{Verbatim}[commandchars=!@|,fontsize=\small,frame=single,label=Example]
  !gapprompt@gap>| !gapinput@F := FreeSemigroup("a", "b", "c");|
  <free semigroup on the generators [ a, b, c ]>
  !gapprompt@gap>| !gapinput@T := F / [[F.1, F.2 ^ 5 * F.3],|
  !gapprompt@>| !gapinput@             [F.2 ^ 6, F.2 ^ 3], [F.2 ^ 2, F.2]];|
  <fp semigroup with 3 generators and 3 relations of length 22>
  !gapprompt@gap>| !gapinput@stz := StzPresentation(T);|
  <fp semigroup presentation with 3 generators and 3 relations
   with length 22>
  !gapprompt@gap>| !gapinput@StzPrintRelations(stz, [2, 3]);|
  #I  2. b^6 = b^3
  #I  3. b^2 = b
  !gapprompt@gap>| !gapinput@StzPrintRelations(stz);|
  #I  1. a = b^5*c
  #I  2. b^6 = b^3
  #I  3. b^2 = b
\end{Verbatim}
 }

 

\subsection{\textcolor{Chapter }{StzPrintRelation}}
\logpage{[ 15, 4, 2 ]}\nobreak
\hyperdef{L}{X81837E037DE5ACBF}{}
{\noindent\textcolor{FuncColor}{$\triangleright$\enspace\texttt{StzPrintRelation({\mdseries\slshape stz, int})\index{StzPrintRelation@\texttt{StzPrintRelation}}
\label{StzPrintRelation}
}\hfill{\scriptsize (operation)}}\\


 If \mbox{\texttt{\mdseries\slshape stz}} is an \texttt{StzPresentation} (\ref{StzPresentation}) object, then \texttt{StzPrintRelation} calls \texttt{StzPrintRelations} with parameters \mbox{\texttt{\mdseries\slshape stz}} and \texttt{[\mbox{\texttt{\mdseries\slshape int}}]}. 
\begin{Verbatim}[commandchars=!@|,fontsize=\small,frame=single,label=Example]
  !gapprompt@gap>| !gapinput@F := FreeSemigroup("a", "b", "c");|
  <free semigroup on the generators [ a, b, c ]>
  !gapprompt@gap>| !gapinput@T := F / [[F.1, F.2 ^ 5 * F.3],|
  !gapprompt@>| !gapinput@             [F.2 ^ 6, F.2 ^ 3], [F.2 ^ 2, F.2]];|
  <fp semigroup with 3 generators and 3 relations of length 22>
  !gapprompt@gap>| !gapinput@stz := StzPresentation(T);|
  <fp semigroup presentation with 3 generators and 3 relations
   with length 22>
  !gapprompt@gap>| !gapinput@StzPrintRelation(stz, 2);|
  #I  2. b^6 = b^3
\end{Verbatim}
 }

 

\subsection{\textcolor{Chapter }{StzPrintGenerators}}
\logpage{[ 15, 4, 3 ]}\nobreak
\hyperdef{L}{X846B86477CBA3A2F}{}
{\noindent\textcolor{FuncColor}{$\triangleright$\enspace\texttt{StzPrintGenerators({\mdseries\slshape stz[, list]})\index{StzPrintGenerators@\texttt{StzPrintGenerators}}
\label{StzPrintGenerators}
}\hfill{\scriptsize (operation)}}\\


 If \mbox{\texttt{\mdseries\slshape stz}} is an \texttt{StzPresentation} (\ref{StzPresentation}) object and \mbox{\texttt{\mdseries\slshape list}} is a list of positive integers, then \texttt{StzPrintGenerators} for each \mbox{\texttt{\mdseries\slshape i}} in \mbox{\texttt{\mdseries\slshape list}} the \mbox{\texttt{\mdseries\slshape i}}th generator and the number of occurrences of that generator in the relations
is printed to the screen.

 If \mbox{\texttt{\mdseries\slshape list}} is not specified then \texttt{StzPrintGenerators} prints all generators of \mbox{\texttt{\mdseries\slshape stz}} in order. 
\begin{Verbatim}[commandchars=!@|,fontsize=\small,frame=single,label=Example]
  !gapprompt@gap>| !gapinput@F := FreeSemigroup("a", "b", "c");|
  <free semigroup on the generators [ a, b, c ]>
  !gapprompt@gap>| !gapinput@T := F / [[F.1, F.2 ^ 5 * F.3],|
  !gapprompt@>| !gapinput@             [F.2 ^ 6, F.2 ^ 3], [F.2 ^ 2, F.2]];|
  <fp semigroup with 3 generators and 3 relations of length 22>
  !gapprompt@gap>| !gapinput@stz := StzPresentation(T);|
  <fp semigroup presentation with 3 generators and 3 relations
   with length 22>
  !gapprompt@gap>| !gapinput@StzPrintGenerators(stz, [1, 2]);|
  #I  1.  a  1 occurrences
  #I  2.  b  17 occurrences
  !gapprompt@gap>| !gapinput@StzPrintGenerators(stz);|
  #I  1.  a  1 occurrences
  #I  2.  b  17 occurrences
  #I  3.  c  1 occurrences
\end{Verbatim}
 }

 

\subsection{\textcolor{Chapter }{StzPrintPresentation}}
\logpage{[ 15, 4, 4 ]}\nobreak
\hyperdef{L}{X7D5096807EC939E1}{}
{\noindent\textcolor{FuncColor}{$\triangleright$\enspace\texttt{StzPrintPresentation({\mdseries\slshape stz})\index{StzPrintPresentation@\texttt{StzPrintPresentation}}
\label{StzPrintPresentation}
}\hfill{\scriptsize (operation)}}\\


 If \mbox{\texttt{\mdseries\slshape stz}} is an \texttt{StzPresentation} (\ref{StzPresentation}) object, then \texttt{StzPrintPresentation} prints a comprehensive overview of \mbox{\texttt{\mdseries\slshape stz}}, including the generators and number of occurrences of each generator in the
relations, the relations as words over the generators, and the forward and
backward maps that indicate how the unreduced semigroup maps to the semigroup
currently described by \mbox{\texttt{\mdseries\slshape stz}}. 
\begin{Verbatim}[commandchars=!@|,fontsize=\small,frame=single,label=Example]
  !gapprompt@gap>| !gapinput@F := FreeSemigroup("a", "b", "c");|
  <free semigroup on the generators [ a, b, c ]>
  !gapprompt@gap>| !gapinput@T := F / [[F.1, F.2 ^ 5 * F.3],|
  !gapprompt@>| !gapinput@             [F.2 ^ 6, F.2 ^ 3], [F.2 ^ 2, F.2]];|
  <fp semigroup with 3 generators and 3 relations of length 22>
  !gapprompt@gap>| !gapinput@stz := StzPresentation(T);|
  <fp semigroup presentation with 3 generators and 3 relations
   with length 22>
  !gapprompt@gap>| !gapinput@StzPrintPresentation(stz);|
  #I  Current generators:
  #I  1.  a  1 occurrences
  #I  2.  b  17 occurrences
  #I  3.  c  1 occurrences
  #I
  #I  Current relations:
  #I  1. a = b^5*c
  #I  2. b^6 = b^3
  #I  3. b^2 = b
  #I
  #I  There are 3 generators and 3 relations of total length 22.
  #I
  #I  Generators of original fp semigroup expressed as
  #I  combinations of generators in current presentation:
  #I  1. a = a
  #I  2. b = b
  #I  3. c = c
  #I
  #I  Generators of current presentation expressed as
  #I  combinations of generators of original fp semigroup:
  #I  1. a = a
  #I  2. b = b
  #I  3. c = c
\end{Verbatim}
 }

 }

 
\section{\textcolor{Chapter }{ Changing Tietze transformation objects }}\label{Changing Tietze transformation objects}
\logpage{[ 15, 5, 0 ]}
\hyperdef{L}{X7EAF8F7A7F7A2A78}{}
{
  Fundamentally, there are four different changes that can be made to a
presentation without changing the algebraic structure of the fp semigroup that
can be derived from it. These four changes are called Tietze transformations,
and they are primarily implemented in this section as operations on an \texttt{StzPresentation} object that will throw errors if the conditions have not been met to perform
the Tietze transformation. 

 However, the checks required in order to ensure that a Tietze transformation
is valid sometimes require verifying equality of two words in an fp semigroup
(for example, to ensure that a relation we are adding to the list of relations
can be derived from the relations already present). Since these checks
sometimes do not terminate, a second implementation of Tietze transformations
assumes good faith and does not perform any checks to see whether the
requested Tietze transformation actually maintains the structure of the
semigroup. This latter type should be used at the user's discretion. If only
the first type are used, the presentation will always give a semigroup
isomorphic to the one used to create the object, but if instead one is not
changing the presentation with the intention of maintaining algebraic
structure, these no\texttt{\symbol{45}}check functions are available for use. 

 The four Tietze transformations on a presentation are adding a relation,
removing a relation, adding a generator, and removing a generator, with
particular conditions on what can be added/removed in order to maintain
structure. More details on each transformation and its arguments and
conditions is given in each entry below. In addition to the four elementary
transformations, there is an additional function \texttt{StzSubstituteRelation} which applies multiple Tietze transformations in sequence. 

 

\subsection{\textcolor{Chapter }{StzAddRelation}}
\logpage{[ 15, 5, 1 ]}\nobreak
\hyperdef{L}{X7FCA079C7C4D398F}{}
{\noindent\textcolor{FuncColor}{$\triangleright$\enspace\texttt{StzAddRelation({\mdseries\slshape stz, pair})\index{StzAddRelation@\texttt{StzAddRelation}}
\label{StzAddRelation}
}\hfill{\scriptsize (operation)}}\\


 If \mbox{\texttt{\mdseries\slshape stz}} is an \texttt{StzPresentation} (\ref{StzPresentation}) object and \mbox{\texttt{\mdseries\slshape pair}} is a list containing two \texttt{LetterRep} words over the generators of \mbox{\texttt{\mdseries\slshape stz}}, then \texttt{StzAddRelation} will perform a Tietze transformation of the first type and add a new relation
to \mbox{\texttt{\mdseries\slshape stz}}. This only happens if the new relation that would be formed from \mbox{\texttt{\mdseries\slshape pair}} can be constructed from the other existing relations; that is, if we can
perform elementary operations using the existing relations of \mbox{\texttt{\mdseries\slshape stz}} to convert \texttt{\mbox{\texttt{\mdseries\slshape pair}}[1]} into \texttt{\mbox{\texttt{\mdseries\slshape pair}}[2]}. 

 If, instead, \mbox{\texttt{\mdseries\slshape pair}} is a list containing two elements of the fp semigroup \texttt{S} that was used to create \texttt{stz}, and the two words are equal in that semigroup, then this function will add
the \texttt{LetterRep} of these words as a new relation to \mbox{\texttt{\mdseries\slshape stz}}. 
\begin{Verbatim}[commandchars=!@|,fontsize=\small,frame=single,label=Example]
  !gapprompt@gap>| !gapinput@F := FreeSemigroup("a", "b", "c");|
  <free semigroup on the generators [ a, b, c ]>
  !gapprompt@gap>| !gapinput@T := F / [[F.1, F.2 ^ 5 * F.3],|
  !gapprompt@>| !gapinput@             [F.2 ^ 6, F.2 ^ 3]];|
  <fp semigroup with 3 generators and 2 relations of length 19>
  !gapprompt@gap>| !gapinput@stz := StzPresentation(T);|
  <fp semigroup presentation with 3 generators and 2 relations
   with length 19>
  !gapprompt@gap>| !gapinput@pair := [[2, 2, 2, 2, 2, 2, 2, 2, 2], [2, 2, 2]];;|
  !gapprompt@gap>| !gapinput@StzAddRelation(stz, pair);|
  !gapprompt@gap>| !gapinput@RelationsOfStzPresentation(stz);|
  [ [ [ 1 ], [ 2, 2, 2, 2, 2, 3 ] ],
    [ [ 2, 2, 2, 2, 2, 2 ], [ 2, 2, 2 ] ],
    [ [ 2, 2, 2, 2, 2, 2, 2, 2, 2 ], [ 2, 2, 2 ] ] ]
  !gapprompt@gap>| !gapinput@pair2 := [[1, 1], [3]];|
  [ [ 1, 1 ], [ 3 ] ]
  !gapprompt@gap>| !gapinput@StzAddRelation(stz, pair2);|
  Error, StzAddRelation: second argument <pair> must list two
  words that are equal in the presentation <stz>
\end{Verbatim}
 }

 

\subsection{\textcolor{Chapter }{StzRemoveRelation}}
\logpage{[ 15, 5, 2 ]}\nobreak
\hyperdef{L}{X802456968490EE7C}{}
{\noindent\textcolor{FuncColor}{$\triangleright$\enspace\texttt{StzRemoveRelation({\mdseries\slshape stz, index})\index{StzRemoveRelation@\texttt{StzRemoveRelation}}
\label{StzRemoveRelation}
}\hfill{\scriptsize (operation)}}\\


 If \mbox{\texttt{\mdseries\slshape stz}} is an \texttt{StzPresentation} (\ref{StzPresentation}) object and \mbox{\texttt{\mdseries\slshape index}} is a positive integer less than or equal to the number of relations of \mbox{\texttt{\mdseries\slshape stz}}, then \texttt{StzRemoveRelation} will perform a Tietze transformation of the second type and remove the \mbox{\texttt{\mdseries\slshape index}}th relation of \mbox{\texttt{\mdseries\slshape stz}} if that relation is such that one side of it can be obtained from the other by
a sequence of elementary operations using only the other relations of \mbox{\texttt{\mdseries\slshape stz}}. 
\begin{Verbatim}[commandchars=!@|,fontsize=\small,frame=single,label=Example]
  !gapprompt@gap>| !gapinput@F := FreeSemigroup("a", "b", "c");|
  <free semigroup on the generators [ a, b, c ]>
  !gapprompt@gap>| !gapinput@T := F / [[F.1, F.2 ^ 5 * F.3],|
  !gapprompt@>| !gapinput@             [F.2 ^ 6, F.2 ^ 3], [F.2 ^ 2, F.2]];|
  <fp semigroup with 3 generators and 3 relations of length 22>
  !gapprompt@gap>| !gapinput@stz := StzPresentation(T);|
  <fp semigroup presentation with 3 generators and 3 relations
   with length 22>
  !gapprompt@gap>| !gapinput@StzRemoveRelation(stz, 2);|
  !gapprompt@gap>| !gapinput@RelationsOfStzPresentation(stz);|
  [ [ [ 1 ], [ 2, 2, 2, 2, 2, 3 ] ], [ [ 2, 2 ], [ 2 ] ] ]
  !gapprompt@gap>| !gapinput@StzRemoveRelation(stz, 2);|
  Error, StzRemoveRelation: second argument <index> must point to
  a relation that is redundant in the presentation <stz>
\end{Verbatim}
 }

 

\subsection{\textcolor{Chapter }{StzAddGenerator}}
\logpage{[ 15, 5, 3 ]}\nobreak
\hyperdef{L}{X7F2E7FAF7F899DB9}{}
{\noindent\textcolor{FuncColor}{$\triangleright$\enspace\texttt{StzAddGenerator({\mdseries\slshape stz, word[, name]})\index{StzAddGenerator@\texttt{StzAddGenerator}}
\label{StzAddGenerator}
}\hfill{\scriptsize (operation)}}\\


 If \mbox{\texttt{\mdseries\slshape stz}} is an \texttt{StzPresentation} (\ref{StzPresentation}) object and \mbox{\texttt{\mdseries\slshape word}} is a \texttt{LetterRep} word over the generators of \mbox{\texttt{\mdseries\slshape stz}}, then \texttt{StzAddGenerator} will perform a Tietze transformation which adds a new generator to \mbox{\texttt{\mdseries\slshape stz}} and a new relation of the form \texttt{newgenerator = \mbox{\texttt{\mdseries\slshape word}}}. 

 If, instead, \mbox{\texttt{\mdseries\slshape word}} is a word over the fp semigroup \mbox{\texttt{\mdseries\slshape S}} that was used to create \mbox{\texttt{\mdseries\slshape stz}}, then this function will add the new generator and a new relation with the
new generator equal to the \texttt{LetterRep} of this word as a new relation to \mbox{\texttt{\mdseries\slshape stz}}. 

 A new name for the generator is chosen and added automatically based on the
names of the existing generators to \texttt{GeneratorsOfStzPresentation} (\ref{GeneratorsOfStzPresentation}) if the argument \mbox{\texttt{\mdseries\slshape name}} is not included. If it is, and if \mbox{\texttt{\mdseries\slshape name}} is a string that is not equal to an existing generator, then the string added
to the list of generators will be \mbox{\texttt{\mdseries\slshape name}} instead. 
\begin{Verbatim}[commandchars=!@|,fontsize=\small,frame=single,label=Example]
  !gapprompt@gap>| !gapinput@F := FreeSemigroup("a", "b", "c");|
  <free semigroup on the generators [ a, b, c ]>
  !gapprompt@gap>| !gapinput@T := F / [[F.1, F.2 ^ 5 * F.3],|
  !gapprompt@>| !gapinput@             [F.2 ^ 6, F.2 ^ 3], [F.2 ^ 2, F.2]];|
  <fp semigroup with 3 generators and 3 relations of length 22>
  !gapprompt@gap>| !gapinput@stz := StzPresentation(T);|
  <fp semigroup presentation with 3 generators and 3 relations
   with length 22>
  !gapprompt@gap>| !gapinput@StzAddGenerator(stz, [2, 2, 2]);|
  !gapprompt@gap>| !gapinput@RelationsOfStzPresentation(stz);|
  [ [ [ 1 ], [ 2, 2, 2, 2, 2, 3 ] ],
    [ [ 2, 2, 2, 2, 2, 2 ], [ 2, 2, 2 ] ], [ [ 2, 2 ], [ 2 ] ],
    [ [ 2, 2, 2 ], [ 4 ] ] ]
\end{Verbatim}
 }

 

\subsection{\textcolor{Chapter }{StzRemoveGenerator}}
\logpage{[ 15, 5, 4 ]}\nobreak
\hyperdef{L}{X85D1304A8461AB59}{}
{\noindent\textcolor{FuncColor}{$\triangleright$\enspace\texttt{StzRemoveGenerator({\mdseries\slshape stz, gen/genname[, index]})\index{StzRemoveGenerator@\texttt{StzRemoveGenerator}}
\label{StzRemoveGenerator}
}\hfill{\scriptsize (operation)}}\\


 If \mbox{\texttt{\mdseries\slshape stz}} is an \texttt{StzPresentation} (\ref{StzPresentation}) object and \mbox{\texttt{\mdseries\slshape gen}} is a positive integer less than or equal to the number of generators of \mbox{\texttt{\mdseries\slshape stz}}, then \texttt{StzRemoveGenerator} will perform a Tietze transformation which removes the \mbox{\texttt{\mdseries\slshape gen}}th generator of \mbox{\texttt{\mdseries\slshape stz}}.

 The argument \mbox{\texttt{\mdseries\slshape stz}} must contain a relation of the form \texttt{\mbox{\texttt{\mdseries\slshape gen}} = word} or \texttt{word = \mbox{\texttt{\mdseries\slshape gen}}}, where \texttt{word} contains no occurrences of the generator \mbox{\texttt{\mdseries\slshape gen}} being removed. The generator is then removed from the presentation by
replacing every instance of \mbox{\texttt{\mdseries\slshape gen}} with a copy of \texttt{word}.

 If the second argument is a string \mbox{\texttt{\mdseries\slshape genname}} rather than a positive integer \mbox{\texttt{\mdseries\slshape gen}}, then the function searches the generators of \mbox{\texttt{\mdseries\slshape stz}} for a generator with the same name and attempts to remove the generator if the
same conditions as above are met. 

 If the argument \mbox{\texttt{\mdseries\slshape index}} is included and is a positive integer less than or equal to the number of
relations, then rather than searching the relations for the first to satisfy
the necessary conditions, the function checks the \mbox{\texttt{\mdseries\slshape index}}th relation to see if it satisfies those conditions, and applies the Tietze
transformation by removing this relation. 
\begin{Verbatim}[commandchars=!@|,fontsize=\small,frame=single,label=Example]
  !gapprompt@gap>| !gapinput@F := FreeSemigroup("a", "b", "c");|
  <free semigroup on the generators [ a, b, c ]>
  !gapprompt@gap>| !gapinput@T := F / [[F.1, F.2 ^ 5 * F.3],|
  !gapprompt@>| !gapinput@             [F.2 ^ 6, F.2 ^ 3], [F.2 ^ 2, F.2]];|
  <fp semigroup with 3 generators and 3 relations of length 22>
  !gapprompt@gap>| !gapinput@stz := StzPresentation(T);|
  <fp semigroup presentation with 3 generators and 3 relations
   with length 22>
  !gapprompt@gap>| !gapinput@StzRemoveGenerator(stz, 1);|
  !gapprompt@gap>| !gapinput@RelationsOfStzPresentation(stz);|
  [ [ [ 1, 1, 1, 1, 1, 1 ], [ 1, 1, 1 ] ], [ [ 1, 1 ], [ 1 ] ] ]
\end{Verbatim}
 }

 

\subsection{\textcolor{Chapter }{StzSubstituteRelation}}
\logpage{[ 15, 5, 5 ]}\nobreak
\hyperdef{L}{X85A89C227B8B1A96}{}
{\noindent\textcolor{FuncColor}{$\triangleright$\enspace\texttt{StzSubstituteRelation({\mdseries\slshape stz, index, side})\index{StzSubstituteRelation@\texttt{StzSubstituteRelation}}
\label{StzSubstituteRelation}
}\hfill{\scriptsize (operation)}}\\


 If \mbox{\texttt{\mdseries\slshape stz}} is an \texttt{StzPresentation} (\ref{StzPresentation}) object and \mbox{\texttt{\mdseries\slshape index}} is a positive integer less than or equal to the number of relations of \mbox{\texttt{\mdseries\slshape stz}} and \mbox{\texttt{\mdseries\slshape side}} is either \texttt{1} or \texttt{2}, then \texttt{StzRemoveGenerator} will perform a sequence of Tietze transformations in order to replace, for the \mbox{\texttt{\mdseries\slshape index}}th relation (say \texttt{[u, v]}), to replace all instances of the \mbox{\texttt{\mdseries\slshape side}}th word of the relation in all other relations by the other side (so, for \texttt{\mbox{\texttt{\mdseries\slshape side}}=1}, all instances of \texttt{u} in all other relations of \mbox{\texttt{\mdseries\slshape stz}} are replaced by \texttt{v}). This requires two Tietze transformations per relation containing \texttt{u}, one to add a new redundant relation with each \texttt{u} replaced by \texttt{v}, and another to remove the original (now redundant) relation. 
\begin{Verbatim}[commandchars=!@|,fontsize=\small,frame=single,label=Example]
  !gapprompt@gap>| !gapinput@F := FreeSemigroup("a", "b", "c");|
  <free semigroup on the generators [ a, b, c ]>
  !gapprompt@gap>| !gapinput@T := F / [[F.1, F.2 ^ 5 * F.3],|
  !gapprompt@>| !gapinput@             [F.2 ^ 6, F.2 ^ 3], [F.2 ^ 2, F.2]];|
  <fp semigroup with 3 generators and 3 relations of length 22>
  !gapprompt@gap>| !gapinput@stz := StzPresentation(T);|
  <fp semigroup presentation with 3 generators and 3 relations
   with length 22>
  !gapprompt@gap>| !gapinput@StzSubstituteRelation(stz, 3, 1);|
  !gapprompt@gap>| !gapinput@RelationsOfStzPresentation(stz);|
  [ [ [ 1 ], [ 2, 2, 2, 3 ] ], [ [ 2, 2, 2 ], [ 2, 2 ] ],
    [ [ 2, 2 ], [ 2 ] ] ]
  !gapprompt@gap>| !gapinput@StzSubstituteRelation(stz, 3, 1);|
  !gapprompt@gap>| !gapinput@RelationsOfStzPresentation(stz);|
  [ [ [ 1 ], [ 2, 2, 3 ] ], [ [ 2, 2 ], [ 2 ] ], [ [ 2, 2 ], [ 2 ] ] ]
\end{Verbatim}
 }

 }

 
\section{\textcolor{Chapter }{ Converting a Tietze transformation object into a fp semigroup }}\label{Converting a Tietze transformation object into a fp semigroup}
\logpage{[ 15, 6, 0 ]}
\hyperdef{L}{X7CBD15BC86CC2080}{}
{
  Semigroup Tietze transformation objects (\texttt{IsStzPresentation} (\ref{IsStzPresentation})) are not actual fp semigroups in the sense of \texttt{IsFpSemigroup} (\textbf{Reference: IsFpSemigroup}). This is because their generators and relations can be modified (see section \ref{Changing Tietze transformation objects}). However, an \texttt{StzPresentation} (\ref{StzPresentation}) can be converted back into an actual finitely presented semigroup using the
methods described in this section.

 The intended use of semigroup Tietze objects is as follows: given an fp
semigroup \texttt{S}, create a modifiable presentation using \texttt{stz := StzPresentation(S)}, apply Tietze transformations to it (perhaps in order to simplify the
presentation), then generate a new fp semigroup \texttt{T} given by \texttt{stz} which is isomorphic to \texttt{S}, but has a simpler presentation. Once \texttt{T} is obtained, it may be of interest to map elements of \texttt{S} over to \texttt{T}, where they may be represented by different combinations of generators. The
isomorphism achieving this is described in this section (see \texttt{StzIsomorphism} (\ref{StzIsomorphism})). 

\subsection{\textcolor{Chapter }{TietzeForwardMap}}
\logpage{[ 15, 6, 1 ]}\nobreak
\hyperdef{L}{X7C46AF3A867F273A}{}
{\noindent\textcolor{FuncColor}{$\triangleright$\enspace\texttt{TietzeForwardMap({\mdseries\slshape stz})\index{TietzeForwardMap@\texttt{TietzeForwardMap}}
\label{TietzeForwardMap}
}\hfill{\scriptsize (attribute)}}\\
\textbf{\indent Returns:\ }
 A list of lists of positive integers. 



 If \mbox{\texttt{\mdseries\slshape stz}} is an \texttt{StzPresentation} (\ref{StzPresentation}) object, then \texttt{TietzeForwardMap} returns a list of lists of positive integers. There is an element of this list
for every generator of the unreduced semigroup of the presentation (see \texttt{UnreducedFpSemigroup} (\ref{UnreducedFpSemigroup:for a presentation})) that indicates the word (in \texttt{LetterRep} form) in the semigroup object currently defined by the presentation that the
generator maps to.

 This mapping is updated as the presentation object is transformed. It begins
as a list of the form \texttt{[[1], [2], [3], . . ., [n]]} where \texttt{n} is the number of generators of the unreduced semigroup. 
\begin{Verbatim}[commandchars=!@|,fontsize=\small,frame=single,label=Example]
  !gapprompt@gap>| !gapinput@F := FreeSemigroup("a", "b", "c");|
  <free semigroup on the generators [ a, b, c ]>
  !gapprompt@gap>| !gapinput@S := F / [[F.1, F.3 ^ 3], [F.2 ^ 2, F.3 ^ 2]];|
  <fp semigroup with 3 generators and 2 relations of length 11>
  !gapprompt@gap>| !gapinput@stz := StzPresentation(S);|
  <fp semigroup presentation with 3 generators and 2 relations
   with length 11>
  !gapprompt@gap>| !gapinput@StzRemoveGenerator(stz, 1);|
  !gapprompt@gap>| !gapinput@TietzeForwardMap(stz);|
  [ [ 2, 2, 2 ], [ 1 ], [ 2 ] ]
\end{Verbatim}
 }

 

\subsection{\textcolor{Chapter }{TietzeBackwardMap}}
\logpage{[ 15, 6, 2 ]}\nobreak
\hyperdef{L}{X7F6E26348503DD53}{}
{\noindent\textcolor{FuncColor}{$\triangleright$\enspace\texttt{TietzeBackwardMap({\mdseries\slshape stz})\index{TietzeBackwardMap@\texttt{TietzeBackwardMap}}
\label{TietzeBackwardMap}
}\hfill{\scriptsize (attribute)}}\\
\textbf{\indent Returns:\ }
 A list of lists of positive integers. 



 If \mbox{\texttt{\mdseries\slshape stz}} is an \texttt{StzPresentation} (\ref{StzPresentation}) object, then \texttt{TietzeBackwardMap} returns a list of lists of positive integers. There is an element of this list
for every generator of the semigroup that the presentation currently defines
that indicates the word (in \texttt{LetterRep} form) in the unreduced semigroup of the presentation (see \texttt{UnreducedFpSemigroup} (\ref{UnreducedFpSemigroup:for a presentation})) that the generator maps to.

 This mapping is updated as the presentation object is transformed. It begins
as a list of the form \texttt{[[1], [2], [3], . . ., [n]]} where \texttt{n} is the number of generators of the unreduced semigroup. 
\begin{Verbatim}[commandchars=!@|,fontsize=\small,frame=single,label=Example]
  !gapprompt@gap>| !gapinput@F := FreeSemigroup("a", "b", "c");|
  <free semigroup on the generators [ a, b, c ]>
  !gapprompt@gap>| !gapinput@S := F / [[F.1, F.3 ^ 3], [F.2 ^ 2, F.3 ^ 2]];|
  <fp semigroup with 3 generators and 2 relations of length 11>
  !gapprompt@gap>| !gapinput@stz := StzPresentation(S);|
  <fp semigroup presentation with 3 generators and 2 relations
   with length 11>
  !gapprompt@gap>| !gapinput@StzRemoveGenerator(stz, 1);|
  !gapprompt@gap>| !gapinput@TietzeBackwardMap(stz);|
  [ [ 2 ], [ 3 ] ]
\end{Verbatim}
 }

 

\subsection{\textcolor{Chapter }{StzIsomorphism}}
\logpage{[ 15, 6, 3 ]}\nobreak
\hyperdef{L}{X80142D917E368E46}{}
{\noindent\textcolor{FuncColor}{$\triangleright$\enspace\texttt{StzIsomorphism({\mdseries\slshape stz})\index{StzIsomorphism@\texttt{StzIsomorphism}}
\label{StzIsomorphism}
}\hfill{\scriptsize (operation)}}\\
\textbf{\indent Returns:\ }
 A mapping object. 



 If \mbox{\texttt{\mdseries\slshape stz}} is an \texttt{StzPresentation} (\ref{StzPresentation}) object, then \texttt{StzIsomorphism} returns a mapping object that maps the unreduced semigroup of the presentation
(see \texttt{UnreducedFpSemigroup} (\ref{UnreducedFpSemigroup:for a presentation})) to an FpSemigroup object that is defined by the generators and relations of
the semigroup presentation at the moment this function is ran.

 If a \texttt{StzIsomorphism} is generated from \mbox{\texttt{\mdseries\slshape stz}}, and the presentation \mbox{\texttt{\mdseries\slshape stz}} is further modified afterwards (for example by applying more Tietze
transformations or \texttt{StzSimplifyOnce} (\ref{StzSimplifyOnce}) to \mbox{\texttt{\mdseries\slshape stz}}), then running \texttt{StzIsomorphism(\mbox{\texttt{\mdseries\slshape stz}})} a second time will produce a different result consistent with the \emph{new} generators and relations of \mbox{\texttt{\mdseries\slshape stz}}.

 This mapping is built from the \texttt{TietzeForwardMap} (\ref{TietzeForwardMap}) and \texttt{TietzeBackwardMap} (\ref{TietzeBackwardMap}) attributes from the presentation object, since if we know how to map the
generators of the respective semigroups, then we know how to map any element
of that semigroup.

 This function is the primary way to obtain the simplified semigroup from the
presentation object, by applying \texttt{Range} to the mapping that this function returns. 
\begin{Verbatim}[commandchars=!@|,fontsize=\small,frame=single,label=Example]
  !gapprompt@gap>| !gapinput@F := FreeSemigroup("a", "b", "c");|
  <free semigroup on the generators [ a, b, c ]>
  !gapprompt@gap>| !gapinput@S := F / [[F.1, F.3 ^ 3], [F.2 ^ 2, F.3 ^ 2]];|
  <fp semigroup with 3 generators and 2 relations of length 11>
  !gapprompt@gap>| !gapinput@stz := StzPresentation(S);|
  <fp semigroup presentation with 3 generators and 2 relations
   with length 11>
  !gapprompt@gap>| !gapinput@StzRemoveGenerator(stz, "a");|
  !gapprompt@gap>| !gapinput@map := StzIsomorphism(stz);|
  <fp semigroup with 3 generators and 2 relations of length 11> ->
  <fp semigroup with 2 generators and 1 relation of length 6>
  !gapprompt@gap>| !gapinput@S.1 ^ map;|
  c^3
\end{Verbatim}
 }

 }

 
\section{\textcolor{Chapter }{ Automatically simplifying a Tietze transformation object }}\label{Automatically simplifying a Tietze transformation object}
\logpage{[ 15, 7, 0 ]}
\hyperdef{L}{X8549F1C87E7BD29A}{}
{
  

 It is possible to create a presentation object from an fp semigroup object and
attempt to manually reduce it by applying Tietze transformations. However,
there may be many different reductions that can be applied, so \texttt{StzSimplifyOnce} can be used to automatically check a number of different possible reductions
and apply the best one. Then, \texttt{StzSimplifyPresentation } repeatedly applies StzSimplifyOnce to the presentation object until it fails
to reduce the presentation any further. The metric with respect to which the \texttt{IsStzPresentation} object is reduced is \texttt{Length} (\ref{Length}). 

\subsection{\textcolor{Chapter }{StzSimplifyOnce}}
\logpage{[ 15, 7, 1 ]}\nobreak
\hyperdef{L}{X818403DA85EA82B8}{}
{\noindent\textcolor{FuncColor}{$\triangleright$\enspace\texttt{StzSimplifyOnce({\mdseries\slshape stz})\index{StzSimplifyOnce@\texttt{StzSimplifyOnce}}
\label{StzSimplifyOnce}
}\hfill{\scriptsize (operation)}}\\
\textbf{\indent Returns:\ }
 \texttt{true} or \texttt{false}. 



 If \mbox{\texttt{\mdseries\slshape stz}} is an \texttt{StzPresentation} (\ref{StzPresentation}) object, then \texttt{StzSimplifyOnce} will check the possible reductions in length for a number of different
possible Tietze transformations, and apply the choice which gives the least
length. If a valid transformation was found then the function returns \texttt{true}, and if no transformation was performed because none would lead to a
reduction in length, then the function returns \texttt{false}.

 There are four different possible transformations that \texttt{StzSimplifyOnce} may apply. The function searches for redundant generators and checks if
removing them would reduce the overall length of the presentation, it checks
whether substituting one side of each relation throughout the rest of the
relations would reduce the length, it checks whether there are any trivial
relations (of the form \mbox{\texttt{\mdseries\slshape w = w}} for some word \mbox{\texttt{\mdseries\slshape w}}) or any duplicated relations (relations which are formed from precisely the
same words as another relation), and it checks whether any frequently
occurring subwords in the relations can be replaced with a new generator to
reduce the length. For more details, see \ref{Changing Tietze transformation objects}

 At \texttt{InfoLevel} 2 (which is the default value, and which can be set using \texttt{SetInfoLevel(InfoFpSemigroup, 2)}), the precise transformations performed are printed to the screen. 
\begin{Verbatim}[commandchars=!@|,fontsize=\small,frame=single,label=Example]
  !gapprompt@gap>| !gapinput@F := FreeSemigroup("a", "b", "c");|
  <free semigroup on the generators [ a, b, c ]>
  !gapprompt@gap>| !gapinput@T := F / [[F.1, F.2 ^ 5 * F.3],|
  !gapprompt@>| !gapinput@             [F.2 ^ 6, F.2 ^ 3]];|
  <fp semigroup with 3 generators and 2 relations of length 19>
  !gapprompt@gap>| !gapinput@stz := StzPresentation(T);|
  <fp semigroup presentation with 3 generators and 2 relations
   with length 19>
  !gapprompt@gap>| !gapinput@StzSimplifyOnce(stz);|
  #I  <Removing redundant generator a using relation : 1. a = b^5*c>
  true
  !gapprompt@gap>| !gapinput@stz;|
  <fp semigroup presentation with 2 generators and 1 relation
   with length 11>
\end{Verbatim}
 }

 

\subsection{\textcolor{Chapter }{StzSimplifyPresentation}}
\logpage{[ 15, 7, 2 ]}\nobreak
\hyperdef{L}{X87B27F5A87103502}{}
{\noindent\textcolor{FuncColor}{$\triangleright$\enspace\texttt{StzSimplifyPresentation({\mdseries\slshape stz})\index{StzSimplifyPresentation@\texttt{StzSimplifyPresentation}}
\label{StzSimplifyPresentation}
}\hfill{\scriptsize (operation)}}\\


 If \mbox{\texttt{\mdseries\slshape stz}} is an \texttt{StzPresentation} object, then \texttt{StzSimplifyPresentation} will repeatedly apply the best of a few possible reductions to \mbox{\texttt{\mdseries\slshape stz}} until it can no longer reduce the length of the presentation.

 
\begin{Verbatim}[commandchars=!@|,fontsize=\small,frame=single,label=Example]
  !gapprompt@gap>| !gapinput@F := FreeSemigroup("a", "b", "c");|
  <free semigroup on the generators [ a, b, c ]>
  !gapprompt@gap>| !gapinput@T := F / [[F.1, F.2 ^ 5 * F.3],|
  !gapprompt@>| !gapinput@             [F.2 ^ 6, F.2 ^ 3]];|
  <fp semigroup on the generators [ a, b, c ]>
  !gapprompt@gap>| !gapinput@stz := StzPresentation(T);|
  <fp semigroup presentation with 3 generators and 2 relations
   with length 19>
  !gapprompt@gap>| !gapinput@StzSimplifyPresentation(stz);|
  !gapprompt@gap>| !gapinput@RelationsOfStzPresentation(stz);|
  [ [ [ 1, 1, 1 ], [ 3 ] ], [ [ 3, 3 ], [ 3 ] ] ]
\end{Verbatim}
 }

 }

 
\section{\textcolor{Chapter }{ Automatically simplifying an fp semigroup }}\label{Automatically simplifying an fp semigroup}
\logpage{[ 15, 8, 0 ]}
\hyperdef{L}{X817332E27D0406A7}{}
{
  

 It may be the case that, rather than working with a Tietze transformation
object, we want to start with an fp semigroup object and obtain the most
simplified version of that fp semigroup that \texttt{StzSimplifyPresentation} can produce. In this case, \texttt{SimplifyFpSemigroup} can be applied to obtain a mapping from its argument to a reduced fp
semigroup. If the mapping is not of interest, \texttt{SimplifiedFpSemigroup} can be used to directly obtain a new fp semigroup isomorphic to the first with
reduced relations and generators (the mapping is stored as an attribute of the
output). With these functions, the user never has to consider precisely what
Tietze transformations to perform, and never has to worry about using the \texttt{StzPresentation} object or its associated operations. They can start with an fp semigroup and
obtain a simplified fp semigroup. 

\subsection{\textcolor{Chapter }{SimplifyFpSemigroup}}
\logpage{[ 15, 8, 1 ]}\nobreak
\hyperdef{L}{X7F9ED6117BA94F01}{}
{\noindent\textcolor{FuncColor}{$\triangleright$\enspace\texttt{SimplifyFpSemigroup({\mdseries\slshape S})\index{SimplifyFpSemigroup@\texttt{SimplifyFpSemigroup}}
\label{SimplifyFpSemigroup}
}\hfill{\scriptsize (operation)}}\\
\textbf{\indent Returns:\ }
 A mapping object. 



 If \mbox{\texttt{\mdseries\slshape S}} is a finitely presented semigroup, then \texttt{SimplifyFpSemigroup} will return a mapping object which will map \mbox{\texttt{\mdseries\slshape S}} to a finitely presented semigroup which has had its presentation simplified. \texttt{SimplifyFpSemigroup} creates an \texttt{StzPresentation} object \mbox{\texttt{\mdseries\slshape stz}} from \mbox{\texttt{\mdseries\slshape S}}, which is then reduced using Tietze transformations until the presentation
cannot be reduced in length any further.

 \texttt{SimplifyFpSemigroup} applies the function \texttt{StzSimplifyPresentation} (\ref{StzSimplifyPresentation}) to \mbox{\texttt{\mdseries\slshape stz}}, which repeatedly checks whether a number of different possible
transformations will cause a reduction in length, and if so applies the best
one. This loop continues until no transformations cause any reductions, in
which case the mapping is returned. The newly reduced FpSemigroup can be
accessed either by taking the range of the mapping or calling \texttt{SimplifiedFpSemigroup}, which first runs \texttt{SimplifyFpSemigroup} and then returns the range of the mapping with the mapping held as an
attribute.

 For more information on how the mapping is created and used, go to \ref{Converting a Tietze transformation object into a fp semigroup}. 
\begin{Verbatim}[commandchars=!@|,fontsize=\small,frame=single,label=Example]
  !gapprompt@gap>| !gapinput@F := FreeSemigroup("a", "b", "c");|
  <free semigroup on the generators [ a, b, c ]>
  !gapprompt@gap>| !gapinput@T := F / [[F.1, F.2 ^ 5 * F.3],|
  !gapprompt@>| !gapinput@             [F.2 ^ 6, F.2 ^ 3]];|
  <fp semigroup on the generators [ a, b, c ]>
  !gapprompt@gap>| !gapinput@map := SimplifyFpSemigroup(T);;|
  #I  Applying StzSimplifyPresentation...
  #I  StzSimplifyPresentation is verbose by default. Use SetInfoLevel(InfoFpSemigroup, 1) to hide
  #I  output while maintaining ability to use StzPrintRelations, StzPrintGenerators, etc.
  #I  Current: <fp semigroup presentation with 3 generators and 2 relations with length 19>
  #I  <Removing redundant generator a using relation : 1. a = b^5*c>
  #I  Current: <fp semigroup presentation with 2 generators and 1 relation with length 11>
  #I  <Creating new generator to replace instances of word: b^3>
  #I  Current: <fp semigroup presentation with 3 generators and 2 relations with length 10>
  !gapprompt@gap>| !gapinput@IsMapping(map);|
  true
  !gapprompt@gap>| !gapinput@T.1;|
  a
  !gapprompt@gap>| !gapinput@T.1 ^ map;|
  b^5*c
  !gapprompt@gap>| !gapinput@RelationsOfFpSemigroup(Range(map));|
  [ [ b^3, d ], [ d^2, d ] ]
\end{Verbatim}
 }

 

\subsection{\textcolor{Chapter }{SimplifiedFpSemigroup}}
\logpage{[ 15, 8, 2 ]}\nobreak
\hyperdef{L}{X7B32AC2B7D5C74D6}{}
{\noindent\textcolor{FuncColor}{$\triangleright$\enspace\texttt{SimplifiedFpSemigroup({\mdseries\slshape S})\index{SimplifiedFpSemigroup@\texttt{SimplifiedFpSemigroup}}
\label{SimplifiedFpSemigroup}
}\hfill{\scriptsize (operation)}}\\
\textbf{\indent Returns:\ }
 A finitely presented semigroup. 



 If \mbox{\texttt{\mdseries\slshape S}} is an fp semigroup object (\texttt{IsFpSemigroup} (\textbf{Reference: IsFpSemigroup})), then \texttt{SimplifiedFpSemigroup} will return an FpSemigroup object \mbox{\texttt{\mdseries\slshape T}} which is isomorphic to \mbox{\texttt{\mdseries\slshape S}} which has been reduced to minimise its length. \texttt{SimplifiedFpSemigroup} applies \texttt{SimplifyFpSemigroup} (\ref{SimplifyFpSemigroup}) and assigns the \texttt{Range} of the isomorphism object which is returned to \mbox{\texttt{\mdseries\slshape T}}, adding the isomorphism to \mbox{\texttt{\mdseries\slshape T}} as an attribute. In this way, while \mbox{\texttt{\mdseries\slshape T}} is a completely new FpSemigroup object, words in \mbox{\texttt{\mdseries\slshape S}} can be mapped to \mbox{\texttt{\mdseries\slshape T}} using the map obtained from the attribute \texttt{FpTietzeIsomorphism} (\ref{FpTietzeIsomorphism}).

 For more information on the mapping between the semigroups and how it is
created, see \ref{Converting a Tietze transformation object into a fp semigroup}. 
\begin{Verbatim}[commandchars=!@|,fontsize=\small,frame=single,label=Example]
  !gapprompt@gap>| !gapinput@F := FreeSemigroup("a", "b", "c");;|
  !gapprompt@gap>| !gapinput@S := F / [[F.1 ^ 4, F.1], [F.1, F.1 ^ 44], [F.1 ^ 8, F.2 * F.3]];;|
  !gapprompt@gap>| !gapinput@T := SimplifiedFpSemigroup(S);;|
  #I  Applying StzSimplifyPresentation...
  #I  StzSimplifyPresentation is verbose by default. Use SetInfoLevel(InfoFpSemigroup, 1) to hide
  #I  output while maintaining ability to use StzPrintRelations, StzPrintGenerators, etc.
  #I  Current: <fp semigroup presentation with 3 generators and 3 relations with length 63>
  #I  <Replacing all instances in other relations of relation: 1. a^4 = a>
  #I  Current: <fp semigroup presentation with 3 generators and 3 relations with length 24>
  #I  <Replacing all instances in other relations of relation: 1. a^4 = a>
  #I  Current: <fp semigroup presentation with 3 generators and 3 relations with length 18>
  #I  <Replacing all instances in other relations of relation: 1. a^4 = a>
  #I  Current: <fp semigroup presentation with 3 generators and 3 relations with length 15>
  #I  <Replacing all instances in other relations of relation: 3. a = a^2>
  #I  Current: <fp semigroup presentation with 3 generators and 3 relations with length 12>
  #I  <Removing duplicate relation: 1. a = a^2>
  #I  Current: <fp semigroup presentation with 3 generators and 2 relations with length 9>
  #I  <Removing redundant generator a using relation : 2. a = b*c>
  #I  Current: <fp semigroup presentation with 2 generators and 1 relation with length 8>
  !gapprompt@gap>| !gapinput@map := FpTietzeIsomorphism(T);;|
  !gapprompt@gap>| !gapinput@S.1 ^ map;|
  b*c
  !gapprompt@gap>| !gapinput@S.1 ^ map = T.1 * T.2;|
  true
  !gapprompt@gap>| !gapinput@invmap := InverseGeneralMapping(map);;|
  !gapprompt@gap>| !gapinput@T.1 ^ invmap = S.2;|
  true
  !gapprompt@gap>| !gapinput@T.1 = S.2;|
  false
\end{Verbatim}
 }

 

\subsection{\textcolor{Chapter }{UnreducedFpSemigroup (for a semigroup)}}
\logpage{[ 15, 8, 3 ]}\nobreak
\hyperdef{L}{X876620BC7BD37EF6}{}
{\noindent\textcolor{FuncColor}{$\triangleright$\enspace\texttt{UnreducedFpSemigroup({\mdseries\slshape S})\index{UnreducedFpSemigroup@\texttt{UnreducedFpSemigroup}!for a semigroup}
\label{UnreducedFpSemigroup:for a semigroup}
}\hfill{\scriptsize (attribute)}}\\
\textbf{\indent Returns:\ }
 \mbox{\texttt{\mdseries\slshape T}}, an fp semigroup object. 



 If \mbox{\texttt{\mdseries\slshape S}} is an fp semigroup object that has been obtained through calling \texttt{SimplifiedFpSemigroup} (\ref{SimplifiedFpSemigroup}) on some fp semigroup \mbox{\texttt{\mdseries\slshape T}} then \texttt{UnreducedFpSemigroup} returns the original semigroup object before simplification. These are
unrelated semigroup objects, except that \mbox{\texttt{\mdseries\slshape S}} will have a \texttt{FpTietzeIsomorphism} (\ref{FpTietzeIsomorphism}) attribute that returns an isomorphic mapping from \mbox{\texttt{\mdseries\slshape T}} to \mbox{\texttt{\mdseries\slshape S}}. 

 If \texttt{SimplifyFpSemigroup} (\ref{SimplifyFpSemigroup}) has been called on an fp semigroup \mbox{\texttt{\mdseries\slshape T}}, then \texttt{UnreducedFpSemigroup} can be used on the \texttt{Range} of the resultant mapping to obtain the domain. 
\begin{Verbatim}[commandchars=!@|,fontsize=\small,frame=single,label=Example]
  !gapprompt@gap>| !gapinput@F := FreeSemigroup("a", "b", "c");|
  <free semigroup on the generators [ a, b, c ]>
  !gapprompt@gap>| !gapinput@T := F / [[F.1, F.2 ^ 5 * F.3],|
  !gapprompt@>| !gapinput@             [F.2 ^ 6, F.2 ^ 3]];|
  <fp semigroup on the generators [ a, b, c ]>
  !gapprompt@gap>| !gapinput@S := SimplifiedFpSemigroup(T);|
  #I  Applying StzSimplifyPresentation...
  #I  StzSimplifyPresentation is verbose by default. Use SetInfoLevel(InfoFpSemigroup, 1) to hide
  #I  output while maintaining ability to use StzPrintRelations, StzPrintGenerators, etc.
  #I  Current: <fp semigroup presentation with 3 generators and 2 relations with length 19>
  #I  <Removing redundant generator a using relation : 1. a = b^5*c>
  #I  Current: <fp semigroup presentation with 2 generators and 1 relation with length 11>
  #I  <Creating new generator to replace instances of word: b^3>
  #I  Current: <fp semigroup presentation with 3 generators and 2 relations with length 10>
  <fp semigroup on the generators [ b, c, d ]>
  !gapprompt@gap>| !gapinput@UnreducedFpSemigroup(S) = T;|
  true
\end{Verbatim}
 }

 

\subsection{\textcolor{Chapter }{FpTietzeIsomorphism}}
\logpage{[ 15, 8, 4 ]}\nobreak
\hyperdef{L}{X80C4E1757D4F3CE5}{}
{\noindent\textcolor{FuncColor}{$\triangleright$\enspace\texttt{FpTietzeIsomorphism({\mdseries\slshape S})\index{FpTietzeIsomorphism@\texttt{FpTietzeIsomorphism}}
\label{FpTietzeIsomorphism}
}\hfill{\scriptsize (attribute)}}\\
\textbf{\indent Returns:\ }
 A mapping object. 



 If \mbox{\texttt{\mdseries\slshape S}} is an fp semigroup object that has been obtained through calling \texttt{SimplifiedFpSemigroup} (\ref{SimplifiedFpSemigroup}) on some fp semigroup \mbox{\texttt{\mdseries\slshape T}}, then \texttt{FpTietzeIsomorphism} returns an isomorphism from \mbox{\texttt{\mdseries\slshape T}} to \mbox{\texttt{\mdseries\slshape S}}. Simplification produces an fp semigroup isomorphic to the original fp
semigroup, and these two fp semigroup objects can interact with each other
through the mapping given by this function. 
\begin{Verbatim}[commandchars=!@|,fontsize=\small,frame=single,label=Example]
  !gapprompt@gap>| !gapinput@F := FreeSemigroup("a", "b", "c");|
  <free semigroup on the generators [ a, b, c ]>
  !gapprompt@gap>| !gapinput@T := F / [[F.1, F.2 ^ 5 * F.3],|
  !gapprompt@>| !gapinput@             [F.2 ^ 6, F.2 ^ 3]];|
  <fp semigroup on the generators [ a, b, c ]>
  !gapprompt@gap>| !gapinput@S := SimplifiedFpSemigroup(T);|
  #I  Applying StzSimplifyPresentation...
  #I  StzSimplifyPresentation is verbose by default. Use SetInfoLevel(InfoFpSemigroup, 1) to hide
  #I  output while maintaining ability to use StzPrintRelations, StzPrintGenerators, etc.
  #I  Current: <fp semigroup presentation with 3 generators and 2 relations with length 19>
  #I  <Removing redundant generator a using relation : 1. a = b^5*c>
  #I  Current: <fp semigroup presentation with 2 generators and 1 relation with length 11>
  #I  <Creating new generator to replace instances of word: b^3>
  #I  Current: <fp semigroup presentation with 3 generators and 2 relations with length 10>
  <fp semigroup on the generators [ b, c, d ]>
  !gapprompt@gap>| !gapinput@T.2;|
  b
  !gapprompt@gap>| !gapinput@S.1;|
  b
  !gapprompt@gap>| !gapinput@T.2 = S.1;|
  false
  !gapprompt@gap>| !gapinput@map := FpTietzeIsomorphism(S);;|
  !gapprompt@gap>| !gapinput@T.2 ^ map = S.1;|
  true
\end{Verbatim}
 }

 }

 }

 
\chapter{\textcolor{Chapter }{ Visualising semigroups and elements }}\label{Visualising semigroups and elements}
\logpage{[ 16, 0, 0 ]}
\hyperdef{L}{X80E82C6785300A86}{}
{
  There are two operations \texttt{TikzString} (\ref{TikzString}) and \texttt{DotString} (\ref{DotString}) in \textsf{Semigroups} for creating {\LaTeX} and \texttt{dot} (also known as GraphViz) format pictures of the Green's class structure of a
semigroup and for visualising certain types of elements of a semigroup. There
is also a function \texttt{Splash} (\textbf{Digraphs: Splash}) for automatically processing the output of \texttt{TikzString} (\ref{TikzString}) and \texttt{DotString} (\ref{DotString}) and opening the resulting pdf file.   
\section{\textcolor{Chapter }{ \texttt{dot} pictures }}\logpage{[ 16, 1, 0 ]}
\hyperdef{L}{X82E16CAB874A1D84}{}
{
  In this section, we describe the operations in \textsf{Semigroups} for creating pictures in \texttt{dot} format. 

 The operations described in this section return strings, which can be written
to a file using the function \texttt{FileString} (\textbf{GAPDoc: FileString}) or viewed using \texttt{Splash} (\textbf{Digraphs: Splash}). 

\subsection{\textcolor{Chapter }{DotString}}
\logpage{[ 16, 1, 1 ]}\nobreak
\hyperdef{L}{X7F51F3CD7E13D199}{}
{\noindent\textcolor{FuncColor}{$\triangleright$\enspace\texttt{DotString({\mdseries\slshape S[, options]})\index{DotString@\texttt{DotString}}
\label{DotString}
}\hfill{\scriptsize (operation)}}\\
\textbf{\indent Returns:\ }
A string.



 If the argument \mbox{\texttt{\mdseries\slshape S}} is a semigroup, and the optional second argument \mbox{\texttt{\mdseries\slshape options}} is a record, then this operation produces a graphical representation of the
partial order of the $\mathcal{D}$\texttt{\symbol{45}}classes of the semigroup \mbox{\texttt{\mdseries\slshape S}} together with the eggbox diagram of each $\mathcal{D}$\texttt{\symbol{45}}class. The output is in \texttt{dot} format (also known as \texttt{GraphViz}) format. For details about this file format, and information about how to
display or edit this format see \href{https://www.graphviz.org} {\texttt{https://www.graphviz.org}}. 

 The string returned by \texttt{DotString} can be written to a file using the command \texttt{FileString} (\textbf{GAPDoc: FileString}). 

 The $\mathcal{D}$\texttt{\symbol{45}}classes are shown as eggbox diagrams with $\mathcal{L}$\texttt{\symbol{45}}classes as rows and $\mathcal{R}$\texttt{\symbol{45}}classes as columns; group $\mathcal{H}$\texttt{\symbol{45}}classes are shaded gray and contain an asterisk. The $\mathcal{L}$\texttt{\symbol{45}}classes and $\mathcal{R}$\texttt{\symbol{45}}classes within a $\mathcal{D}$\texttt{\symbol{45}}class are arranged to correspond to the normalization of
the principal factor given by \texttt{NormalizedPrincipalFactor} (\ref{NormalizedPrincipalFactor}). The $\mathcal{D}$\texttt{\symbol{45}}classes are numbered according to their index in \texttt{GreensDClasses(\mbox{\texttt{\mdseries\slshape S}})}, so that an \texttt{i} appears next to the eggbox diagram of \texttt{GreensDClasses(\mbox{\texttt{\mdseries\slshape S}})[i]}. A line from one $\mathcal{D}$\texttt{\symbol{45}}class to another indicates that the higher $\mathcal{D}$\texttt{\symbol{45}}class is greater than the lower one in the $\mathcal{D}$\texttt{\symbol{45}}order on \mbox{\texttt{\mdseries\slshape S}}. 

 If the optional second argument \mbox{\texttt{\mdseries\slshape options}} is present, it can be used to specify some options for output. 
\begin{description}
\item[{number}]  if \texttt{\mbox{\texttt{\mdseries\slshape options}}.number} is \texttt{false}, then the $\mathcal{D}$\texttt{\symbol{45}}classes in the diagram are not numbered according to their
index in the list of $\mathcal{D}$\texttt{\symbol{45}}classes of \mbox{\texttt{\mdseries\slshape S}}. The default value for this option is \texttt{true}. 
\item[{maximal}]  if \texttt{\mbox{\texttt{\mdseries\slshape options}}.maximal} is \texttt{true}, then the structure description of the group $\mathcal{H}$\texttt{\symbol{45}}classes is displayed; see \texttt{StructureDescription} (\textbf{Reference: StructureDescription}). Setting this attribute to \texttt{true} can adversely affect the performance of \texttt{DotString}. The default value for this option is \texttt{false}. 
\item[{normal}]  if \texttt{\mbox{\texttt{\mdseries\slshape options}}.normal} is \texttt{false}, then the $\mathcal{L}$\texttt{\symbol{45}} and $\mathcal{R}$\texttt{\symbol{45}}classes within each $\mathcal{D}$\texttt{\symbol{45}}class arranged to correspond to \texttt{PrincipalFactor} (\ref{PrincipalFactor}). If \texttt{\mbox{\texttt{\mdseries\slshape options}}.normal} is \texttt{true}, they are instead arranged to correspond to \texttt{NormalizedPrincipalFactor} (\ref{NormalizedPrincipalFactor}). Setting this attribute to \texttt{false} may improve the performance of \texttt{DotString} as it avoids the computation of \texttt{InjectionNormalizedPrincipalFactor} (\ref{InjectionNormalizedPrincipalFactor}). The default value for this option is \texttt{true}. 
\end{description}
 
\begin{Verbatim}[commandchars=!@|,fontsize=\small,frame=single,label=Example]
  !gapprompt@gap>| !gapinput@S := FullTransformationMonoid(3);|
  <full transformation monoid of degree 3>
  !gapprompt@gap>| !gapinput@DotString(S);|
  "//dot\ndigraph  DClasses {\nnode [shape=plaintext]\nedge [color=blac\
  k,arrowhead=none]\n1 [shape=box style=invisible label=<\n<TABLE BORDE\
  R=\"0\" CELLBORDER=\"1\" CELLPADDING=\"10\" CELLSPACING=\"0\" PORT=\"\
  1\">\n<TR BORDER=\"0\"><TD COLSPAN=\"1\" BORDER = \"0\" > 1</TD></TR>\
  <TR><TD BGCOLOR=\"gray\">*</TD></TR>\n</TABLE>>];\n2 [shape=box style\
  =invisible label=<\n<TABLE BORDER=\"0\" CELLBORDER=\"1\" CELLPADDING=\
  \"10\" CELLSPACING=\"0\" PORT=\"2\">\n<TR BORDER=\"0\"><TD COLSPAN=\"\
  3\" BORDER = \"0\" > 2</TD></TR><TR><TD BGCOLOR=\"gray\">*</TD><TD BG\
  COLOR=\"gray\">*</TD><TD BGCOLOR=\"white\"></TD></TR>\n<TR><TD BGCOLO\
  R=\"gray\">*</TD><TD BGCOLOR=\"white\"></TD><TD BGCOLOR=\"gray\">*</T\
  D></TR>\n<TR><TD BGCOLOR=\"white\"></TD><TD BGCOLOR=\"gray\">*</TD><T\
  D BGCOLOR=\"gray\">*</TD></TR>\n</TABLE>>];\n3 [shape=box style=invis\
  ible label=<\n<TABLE BORDER=\"0\" CELLBORDER=\"1\" CELLPADDING=\"10\"\
   CELLSPACING=\"0\" PORT=\"3\">\n<TR BORDER=\"0\"><TD COLSPAN=\"1\" BO\
  RDER = \"0\" > 3</TD></TR><TR><TD BGCOLOR=\"gray\">*</TD></TR>\n<TR><\
  TD BGCOLOR=\"gray\">*</TD></TR>\n<TR><TD BGCOLOR=\"gray\">*</TD></TR>\
  \n</TABLE>>];\n1 -> 2\n2 -> 3\n }"
  !gapprompt@gap>| !gapinput@FileString("t3.dot", DotString(S));|
  1040
\end{Verbatim}
 }

 

\subsection{\textcolor{Chapter }{DotString (for a Cayley digraph)}}
\logpage{[ 16, 1, 2 ]}\nobreak
\hyperdef{L}{X853B81B385E2CF36}{}
{\noindent\textcolor{FuncColor}{$\triangleright$\enspace\texttt{DotString({\mdseries\slshape digraph})\index{DotString@\texttt{DotString}!for a Cayley digraph}
\label{DotString:for a Cayley digraph}
}\hfill{\scriptsize (operation)}}\\
\textbf{\indent Returns:\ }
A string.



 If \mbox{\texttt{\mdseries\slshape digraph}} is a \texttt{Digraph} (\textbf{Digraphs: Digraph}) in the category \texttt{IsCayleyDigraph} (\textbf{Digraphs: IsCayleyDigraph}), then \texttt{DotString} returns a graphical representation of \mbox{\texttt{\mdseries\slshape digraph}}. The output is in \texttt{dot} format (also known as \texttt{GraphViz}) format. For details about this file format, and information about how to
display or edit this format see \href{https://www.graphviz.org} {\texttt{https://www.graphviz.org}}. 

 The string returned by \texttt{DotString} can be written to a file using the command \texttt{FileString} (\textbf{GAPDoc: FileString}).

 See also \texttt{DotLeftCayleyDigraph} (\ref{DotLeftCayleyDigraph}) and \texttt{TikzLeftCayleyDigraph} (\ref{TikzLeftCayleyDigraph}). }

 

\subsection{\textcolor{Chapter }{DotSemilatticeOfIdempotents}}
\logpage{[ 16, 1, 3 ]}\nobreak
\hyperdef{L}{X7C22E8D17D6C23EA}{}
{\noindent\textcolor{FuncColor}{$\triangleright$\enspace\texttt{DotSemilatticeOfIdempotents({\mdseries\slshape S})\index{DotSemilatticeOfIdempotents@\texttt{DotSemilatticeOfIdempotents}}
\label{DotSemilatticeOfIdempotents}
}\hfill{\scriptsize (attribute)}}\\
\textbf{\indent Returns:\ }
A string.



 This function produces a graphical representation of the semilattice of the
idempotents of an inverse semigroup \mbox{\texttt{\mdseries\slshape S}} where the elements of \mbox{\texttt{\mdseries\slshape S}} have a unique semigroup inverse accessible via \texttt{Inverse} (\textbf{Reference: Inverse}). The idempotents are grouped by the $\mathcal{D}$\texttt{\symbol{45}}class they belong to. 

 The output is in \texttt{dot} format (also known as \texttt{GraphViz}) format. For details about this file format, and information about how to
display or edit this format see \href{https://www.graphviz.org} {\texttt{https://www.graphviz.org}}. 

 
\begin{Verbatim}[commandchars=!@|,fontsize=\small,frame=single,label=Example]
  !gapprompt@gap>| !gapinput@S := DualSymmetricInverseMonoid(4);|
  <inverse block bijection monoid of degree 4 with 3 generators>
  !gapprompt@gap>| !gapinput@DotSemilatticeOfIdempotents(S);|
  "//dot\ngraph graphname {\n  node [shape=point]\nranksep=2;\nsubgraph \
  cluster_1{\n15 \n}\nsubgraph cluster_2{\n5 11 14 12 13 8 \n}\nsubgraph\
   cluster_3{\n2 10 6 3 4 9 7 \n}\nsubgraph cluster_4{\n1 \n}\n2 -- 1\n3\
   -- 1\n4 -- 1\n5 -- 2\n5 -- 3\n5 -- 4\n6 -- 1\n7 -- 1\n8 -- 2\n8 -- 6\
  \n8 -- 7\n9 -- 1\n10 -- 1\n11 -- 2\n11 -- 9\n11 -- 10\n12 -- 3\n12 -- \
  6\n12 -- 9\n13 -- 3\n13 -- 7\n13 -- 10\n14 -- 4\n14 -- 6\n14 -- 10\n15\
   -- 5\n15 -- 8\n15 -- 11\n15 -- 12\n15 -- 13\n15 -- 14\n }"
\end{Verbatim}
 }

 

\subsection{\textcolor{Chapter }{DotLeftCayleyDigraph}}
\logpage{[ 16, 1, 4 ]}\nobreak
\hyperdef{L}{X7E38369D7E8BEA4C}{}
{\noindent\textcolor{FuncColor}{$\triangleright$\enspace\texttt{DotLeftCayleyDigraph({\mdseries\slshape S})\index{DotLeftCayleyDigraph@\texttt{DotLeftCayleyDigraph}}
\label{DotLeftCayleyDigraph}
}\hfill{\scriptsize (operation)}}\\
\noindent\textcolor{FuncColor}{$\triangleright$\enspace\texttt{DotRightCayleyDigraph({\mdseries\slshape S})\index{DotRightCayleyDigraph@\texttt{DotRightCayleyDigraph}}
\label{DotRightCayleyDigraph}
}\hfill{\scriptsize (operation)}}\\
\textbf{\indent Returns:\ }
A string.



 If \mbox{\texttt{\mdseries\slshape S}} is a semigroup satisfying \texttt{CanUseFroidurePin} (\ref{CanUseFroidurePin}), then \texttt{DotLeftCayleyDigraph} is simply short for \texttt{DotString(LeftCayleyDigraph(\mbox{\texttt{\mdseries\slshape S}}))}.

 \texttt{DotRightCayleyDigraph} can be used to produce a dot string for the right Cayley graph of \mbox{\texttt{\mdseries\slshape S}}.

 See \texttt{DotString} (\ref{DotString}) for more details, and see also \texttt{TikzLeftCayleyDigraph} (\ref{TikzLeftCayleyDigraph}). 
\begin{Verbatim}[commandchars=!@|,fontsize=\small,frame=single,label=Example]
  !gapprompt@gap>| !gapinput@DotLeftCayleyDigraph(Semigroup(IdentityTransformation));|
  "//dot\ndigraph hgn{\nnode [shape=circle]\n1 [label=\"a\"]\n1 -> 1\n}\
  \n"
  !gapprompt@gap>| !gapinput@DotRightCayleyDigraph(Semigroup(IdentityTransformation));|
  "//dot\ndigraph hgn{\nnode [shape=circle]\n1 [label=\"a\"]\n1 -> 1\n}\
  \n"
\end{Verbatim}
 }

 }

   
\section{\textcolor{Chapter }{ \texttt{tex} output }}\logpage{[ 16, 2, 0 ]}
\hyperdef{L}{X83152CA78114E2BD}{}
{
  In this section, we describe the operations in \textsf{Semigroups} for creating {\LaTeX} representations of \textsf{GAP} objects. For pictures of \textsf{GAP} objects please see Section \ref{tikz pictures}. 

\subsection{\textcolor{Chapter }{TexString}}
\logpage{[ 16, 2, 1 ]}\nobreak
\hyperdef{L}{X80FBF08180F7F181}{}
{\noindent\textcolor{FuncColor}{$\triangleright$\enspace\texttt{TexString({\mdseries\slshape f[, n]})\index{TexString@\texttt{TexString}}
\label{TexString}
}\hfill{\scriptsize (operation)}}\\
\textbf{\indent Returns:\ }
A string.



 This function produces a string containing LaTeX code for the transformation \mbox{\texttt{\mdseries\slshape f}}. If the optional parameter \mbox{\texttt{\mdseries\slshape n}} is used, then this is taken to be the degree of the transformation \mbox{\texttt{\mdseries\slshape f}}, if the parameter \mbox{\texttt{\mdseries\slshape n}} is not given, then \texttt{DegreeOfTransformation} (\textbf{Reference: DegreeOfTransformation}) is used by default. If \mbox{\texttt{\mdseries\slshape n}} is less than the degree of \mbox{\texttt{\mdseries\slshape f}}, then an error is given. 

 
\begin{Verbatim}[commandchars=!@|,fontsize=\small,frame=single,label=Example]
  !gapprompt@gap>| !gapinput@TexString(Transformation([6, 2, 4, 3, 6, 4]));|
  "\\begin{pmatrix}\n  1 & 2 & 3 & 4 & 5 & 6 \\\\\n  6 & 2 & 4 & 3 & 6 &\
   4\n\\end{pmatrix}"
  !gapprompt@gap>| !gapinput@TexString(Transformation([1, 2, 1, 3]), 5);|
  "\\begin{pmatrix}\n  1 & 2 & 3 & 4 & 5 \\\\\n  1 & 2 & 1 & 3 & 5\n\\en\
  d{pmatrix}"
\end{Verbatim}
 }

 }

   
\section{\textcolor{Chapter }{ \texttt{tikz} pictures }}\label{tikz pictures}
\logpage{[ 16, 3, 0 ]}
\hyperdef{L}{X7BDDE0FE80D09887}{}
{
  In this section, we describe the operations in \textsf{Semigroups} for creating pictures in \texttt{tikz} format. 

 The functions described in this section return a string, which can be written
to a file using the function \texttt{FileString} (\textbf{GAPDoc: FileString}) or viewed using \texttt{Splash} (\textbf{Digraphs: Splash}). 

\subsection{\textcolor{Chapter }{TikzString}}
\logpage{[ 16, 3, 1 ]}\nobreak
\hyperdef{L}{X7F0971F678B4FC66}{}
{\noindent\textcolor{FuncColor}{$\triangleright$\enspace\texttt{TikzString({\mdseries\slshape obj[, options]})\index{TikzString@\texttt{TikzString}}
\label{TikzString}
}\hfill{\scriptsize (operation)}}\\
\textbf{\indent Returns:\ }
A string.



 This function produces a graphical representation of the object \mbox{\texttt{\mdseries\slshape obj}} using the \texttt{tikz} package for {\LaTeX}. More precisely, this operation outputs a string containing a minimal {\LaTeX} document which can be compiled using {\LaTeX} to produce a picture of \mbox{\texttt{\mdseries\slshape obj}}. 

 Currently the following types of objects are supported: 
\begin{description}
\item[{blocks}]  If \mbox{\texttt{\mdseries\slshape obj}} is the left or right blocks of a bipartition, then \texttt{TikzString} returns a graphical representation of these blocks; see Section \ref{section-blocks}. 
\item[{bipartitions}]  If \mbox{\texttt{\mdseries\slshape obj}} is a bipartition, then \texttt{TikzString} returns a graphical representation of \mbox{\texttt{\mdseries\slshape obj}}. 

 If the optional second argument \mbox{\texttt{\mdseries\slshape options}} is a record with the component \texttt{colors} set to \texttt{true}, then the blocks of \mbox{\texttt{\mdseries\slshape f}} will be colored using the standard \texttt{tikz} colors. Due to the limited number of colors available in \texttt{tikz} this option only works when the degree of \mbox{\texttt{\mdseries\slshape obj}} is less than 20. See Chapter \ref{Bipartitions and blocks} for more details about bipartitions. 
\item[{pbrs}]  If \mbox{\texttt{\mdseries\slshape obj}} is a \texttt{PBR} (\ref{PBR}), then \texttt{TikzString} returns a graphical representation \mbox{\texttt{\mdseries\slshape obj}}; see Chapter \ref{Partitioned binary relations (PBRs)}. 
\item[{Cayley graphs}]  If \mbox{\texttt{\mdseries\slshape obj}} is a \texttt{Digraph} (\textbf{Digraphs: Digraph}) in the category \texttt{IsCayleyDigraph} (\textbf{Digraphs: IsCayleyDigraph}), then \texttt{TikzString} returns a picture of \mbox{\texttt{\mdseries\slshape obj}}. No attempt is made whatsoever to produce a sensible picture of the digraph \mbox{\texttt{\mdseries\slshape obj}}, in fact, the vertices are all given the same coordinates. Human intervention
is required to produce a meaningful picture from the value returned by this
method. It is intended to make the task of drawing such a Cayley graph more
straightforward by providing everything except the final layout of the graph.
Please use \texttt{DotString} (\ref{DotString}) if you want an automatically laid out diagram of the digraph \mbox{\texttt{\mdseries\slshape obj}}. 
\end{description}
 
\begin{Verbatim}[commandchars=!@|,fontsize=\small,frame=single,label=Example]
  !gapprompt@gap>| !gapinput@x := Bipartition([[1, 4, -2, -3], [2, 3, 5, -5], [-1, -4]]);;|
  !gapprompt@gap>| !gapinput@TikzString(RightBlocks(x));|
  "%tikz\n\\documentclass{minimal}\n\\usepackage{tikz}\n\\begin{documen\
  t}\n\\begin{tikzpicture}\n  \\draw[ultra thick](5,2)circle(.115);\n  \
  \\draw(1.8,5) node [top] {{$1$}};\n  \\fill(4,2)circle(.125);\n  \\dr\
  aw(1.8,4) node [top] {{$2$}};\n  \\fill(3,2)circle(.125);\n  \\draw(1\
  .8,3) node [top] {{$3$}};\n  \\draw[ultra thick](2,2)circle(.115);\n \
   \\draw(1.8,2) node [top] {{$4$}};\n  \\fill(1,2)circle(.125);\n  \\d\
  raw(1.8,1) node [top] {{$5$}};\n\n  \\draw (5,2.125) .. controls (5,2\
  .8) and (2,2.8) .. (2,2.125);\n  \\draw (4,2.125) .. controls (4,2.6)\
   and (3,2.6) .. (3,2.125);\n\\end{tikzpicture}\n\n\\end{document}"
  !gapprompt@gap>| !gapinput@x := Bipartition([[1, 5], [2, 4, -3, -5], [3, -1, -2], [-4]]);;|
  !gapprompt@gap>| !gapinput@TikzString(x);|
  "%tikz\n\\documentclass{minimal}\n\\usepackage{tikz}\n\\begin{documen\
  t}\n\\begin{tikzpicture}\n\n  %block #1\n  %vertices and labels\n  \\\
  fill(1,2)circle(.125);\n  \\draw(0.95, 2.2) node [above] {{ $1$}};\n \
   \\fill(5,2)circle(.125);\n  \\draw(4.95, 2.2) node [above] {{ $5$}};\
  \n\n  %lines\n  \\draw(1,1.875) .. controls (1,1.1) and (5,1.1) .. (5\
  ,1.875);\n\n  %block #2\n  %vertices and labels\n  \\fill(2,2)circle(\
  .125);\n  \\draw(1.95, 2.2) node [above] {{ $2$}};\n  \\fill(4,2)circ\
  le(.125);\n  \\draw(3.95, 2.2) node [above] {{ $4$}};\n  \\fill(3,0)c\
  ircle(.125);\n  \\draw(3, -0.2) node [below] {{ $-3$}};\n  \\fill(5,0\
  )circle(.125);\n  \\draw(5, -0.2) node [below] {{ $-5$}};\n\n  %lines\
  \n  \\draw(2,1.875) .. controls (2,1.3) and (4,1.3) .. (4,1.875);\n  \
  \\draw(3,0.125) .. controls (3,0.7) and (5,0.7) .. (5,0.125);\n  \\dr\
  aw(2,2)--(3,0);\n\n  %block #3\n  %vertices and labels\n  \\fill(3,2)\
  circle(.125);\n  \\draw(2.95, 2.2) node [above] {{ $3$}};\n  \\fill(1\
  ,0)circle(.125);\n  \\draw(1, -0.2) node [below] {{ $-1$}};\n  \\fill\
  (2,0)circle(.125);\n  \\draw(2, -0.2) node [below] {{ $-2$}};\n\n  %l\
  ines\n  \\draw(1,0.125) .. controls (1,0.6) and (2,0.6) .. (2,0.125);\
  \n  \\draw(3,2)--(2,0);\n\n  %block #4\n  %vertices and labels\n  \\f\
  ill(4,0)circle(.125);\n  \\draw(4, -0.2) node [below] {{ $-4$}};\n\n \
  %lines\n\\end{tikzpicture}\n\n\\end{document}"
  !gapprompt@gap>| !gapinput@TikzString(UniversalPBR(2));|
  "%latex\n\\documentclass{minimal}\n\\usepackage{tikz}\n\\begin{docume\
  nt}\n\\usetikzlibrary{arrows}\n\\usetikzlibrary{arrows.meta}\n\\newco\
  mmand{\\arc}{\\draw[semithick, -{>[width = 1.5mm, length = 2.5mm]}]}\
  \n\\begin{tikzpicture}[\n  vertex/.style={circle, draw, fill=black, i\
  nner sep =0.04cm},\n  ghost/.style={circle, draw = none, inner sep = \
  0.14cm},\n  botloop/.style={min distance = 8mm, out = -70, in = -110}\
  ,\n  toploop/.style={min distance = 8mm, out = 70, in = 110}]\n\n  % \
  vertices and labels\n  \\foreach \\i in {1,...,2} {\n    \\node [vert\
  ex] at (\\i/1.5, 3) {};\n    \\node [ghost] (\\i) at (\\i/1.5, 3) {};\
  \n  }\n\n  \\foreach \\i in {1,...,2} {\n    \\node [vertex] at (\\i/\
  1.5, 0) {};\n    \\node [ghost] (-\\i) at (\\i/1.5, 0) {};\n  }\n\n  \
  % arcs from vertex 1\n  \\arc (1) to (-2);\n  \\arc (1) to (-1);\n  \
  \\arc (1) edge [toploop] (1);\n  \\arc (1) .. controls (1.06666666666\
  66667, 2.125) and (0.93333333333333324, 2.125) .. (2);\n\n  % arcs fr\
  om vertex -1\n  \\arc (-1) .. controls (1.0666666666666667, 0.875) an\
  d (0.93333333333333324, 0.875) .. (-2);\n  \\arc (-1) edge [botloop] \
  (-1);\n  \\arc (-1) to (1);\n  \\arc (-1) to (2);\n\n  % arcs from ve\
  rtex 2\n  \\arc (2) to (-2);\n  \\arc (2) to (-1);\n  \\arc (2) .. co\
  ntrols (0.93333333333333324, 2.125) and (1.0666666666666667, 2.125) .\
  . (1);\n  \\arc (2) edge [toploop] (2);\n\n  % arcs from vertex -2\n \
   \\arc (-2) edge [botloop] (-2);\n  \\arc (-2) .. controls (0.9333333\
  3333333324, 0.875) and (1.0666666666666667, 0.875) .. (-1);\n  \\arc \
  (-2) to (1);\n  \\arc (-2) to (2);\n\n\\end{tikzpicture}\n\\end{docum\
  ent}"
\end{Verbatim}
 }

 

\subsection{\textcolor{Chapter }{TikzLeftCayleyDigraph}}
\logpage{[ 16, 3, 2 ]}\nobreak
\hyperdef{L}{X81EF7F777EEF69C1}{}
{\noindent\textcolor{FuncColor}{$\triangleright$\enspace\texttt{TikzLeftCayleyDigraph({\mdseries\slshape S})\index{TikzLeftCayleyDigraph@\texttt{TikzLeftCayleyDigraph}}
\label{TikzLeftCayleyDigraph}
}\hfill{\scriptsize (operation)}}\\
\noindent\textcolor{FuncColor}{$\triangleright$\enspace\texttt{TikzRightCayleyDigraph({\mdseries\slshape S})\index{TikzRightCayleyDigraph@\texttt{TikzRightCayleyDigraph}}
\label{TikzRightCayleyDigraph}
}\hfill{\scriptsize (operation)}}\\
\textbf{\indent Returns:\ }
A string.



 If \mbox{\texttt{\mdseries\slshape S}} is a semigroup satisfying \texttt{CanUseFroidurePin} (\ref{CanUseFroidurePin}), then \texttt{TikzLeftCayleyDigraph} is simply short for \texttt{TikzString(LeftCayleyDigraph(\mbox{\texttt{\mdseries\slshape S}}))}.

 \texttt{TikzRightCayleyDigraph} can be used to produce a tikz string for the right Cayley graph of \mbox{\texttt{\mdseries\slshape S}}.

 See \texttt{TikzString} (\ref{TikzString}) for more details, and see also \texttt{DotLeftCayleyDigraph} (\ref{DotLeftCayleyDigraph}). 
\begin{Verbatim}[commandchars=!@|,fontsize=\small,frame=single,label=Example]
  !gapprompt@gap>| !gapinput@TikzLeftCayleyDigraph(Semigroup(IdentityTransformation));|
  "\\begin{tikzpicture}[scale=1, auto, \n  vertex/.style={c\
  ircle, draw, thick, fill=white, minimum size=0.65cm},\n  \
  edge/.style={arrows={-angle 90}, thick},\n  loop/.style={\
  min distance=5mm,looseness=5,arrows={-angle 90},thick}]\n\
  \n  % Vertices . . .\n  \\node [vertex] (a) at (0, 0) {};\
  \n  \\node at (0, 0) {$a$};\n\n  % Edges . . .\n  \\path[\
  ->] (a) edge [loop]\n           node {$a$} (a);\n\\end{ti\
  kzpicture}"
  !gapprompt@gap>| !gapinput@TikzRightCayleyDigraph(Semigroup(IdentityTransformation));|
  "\\begin{tikzpicture}[scale=1, auto, \n  vertex/.style={c\
  ircle, draw, thick, fill=white, minimum size=0.65cm},\n  \
  edge/.style={arrows={-angle 90}, thick},\n  loop/.style={\
  min distance=5mm,looseness=5,arrows={-angle 90},thick}]\n\
  \n  % Vertices . . .\n  \\node [vertex] (a) at (0, 0) {};\
  \n  \\node at (0, 0) {$a$};\n\n  % Edges . . .\n  \\path[\
  ->] (a) edge [loop]\n           node {$a$} (a);\n\\end{ti\
  kzpicture}"
\end{Verbatim}
 }

 }

 }

 
\chapter{\textcolor{Chapter }{ IO }}\label{IO}
\logpage{[ 17, 0, 0 ]}
\hyperdef{L}{X80CDCB927B3E5BB9}{}
{
  
\section{\textcolor{Chapter }{Reading and writing elements to a file}}\logpage{[ 17, 1, 0 ]}
\hyperdef{L}{X7CE72BB17F2D49F8}{}
{
 The functions \texttt{ReadGenerators} (\ref{ReadGenerators}) and \texttt{WriteGenerators} (\ref{WriteGenerators}) can be used to read or write, respectively, elements of a semigroup to a file. 

\subsection{\textcolor{Chapter }{ReadGenerators}}
\logpage{[ 17, 1, 1 ]}\nobreak
\hyperdef{L}{X8728096E8427EDE8}{}
{\noindent\textcolor{FuncColor}{$\triangleright$\enspace\texttt{ReadGenerators({\mdseries\slshape filename[, nr]})\index{ReadGenerators@\texttt{ReadGenerators}}
\label{ReadGenerators}
}\hfill{\scriptsize (function)}}\\
\textbf{\indent Returns:\ }
A list of lists of semigroup elements.



 If \mbox{\texttt{\mdseries\slshape filename}} is an \textsf{IO} package file object or is the name of a file created using \texttt{WriteGenerators} (\ref{WriteGenerators}), then \texttt{ReadGenerators} returns the contents of this file as a list of lists of elements of a
semigroup. 

 If the optional second argument \mbox{\texttt{\mdseries\slshape nr}} is present, then \texttt{ReadGenerators} returns the elements stored in the \mbox{\texttt{\mdseries\slshape nr}}th line of \mbox{\texttt{\mdseries\slshape filename}}. 
\begin{Verbatim}[commandchars=!@|,fontsize=\small,frame=single,label=Example]
  !gapprompt@gap>| !gapinput@file := Concatenation(SEMIGROUPS.PackageDir,|
  !gapprompt@>| !gapinput@"/data/tst/testdata");;|
  !gapprompt@gap>| !gapinput@ReadGenerators(file, 13);|
  [ <identity partial perm on [ 2, 3, 4, 5, 6 ]>,
    <identity partial perm on [ 2, 3, 5, 6 ]>, [1,2](5)(6) ]
\end{Verbatim}
 }

 

\subsection{\textcolor{Chapter }{WriteGenerators}}
\logpage{[ 17, 1, 2 ]}\nobreak
\hyperdef{L}{X78041E8F87EFDE62}{}
{\noindent\textcolor{FuncColor}{$\triangleright$\enspace\texttt{WriteGenerators({\mdseries\slshape filename, list[, append][, function]})\index{WriteGenerators@\texttt{WriteGenerators}}
\label{WriteGenerators}
}\hfill{\scriptsize (function)}}\\
\textbf{\indent Returns:\ }
\texttt{IO{\textunderscore}OK} or \texttt{IO{\textunderscore}ERROR}.



 This function provides a method for writing collections of elements of a
semigroup to a file. The resulting file can be further compressed using \texttt{gzip} or \texttt{xz}.

 The argument \mbox{\texttt{\mdseries\slshape list}} should be a list of lists of elements, or semigroups. 

 The argument \mbox{\texttt{\mdseries\slshape filename}} should be a string containing the name of a file or an \textsf{IO} package file object where the entries in \mbox{\texttt{\mdseries\slshape list}} will be written; see \texttt{IO{\textunderscore}File} (\textbf{IO: IO{\textunderscore}File mode}) and \texttt{IO{\textunderscore}CompressedFile} (\textbf{IO: IO{\textunderscore}CompressedFile}).

 If the optional third argument \mbox{\texttt{\mdseries\slshape append}} is not present or is given and equals \texttt{"w"}, then the previous content of the file is deleted and overwritten. If the
third argument is \texttt{"a"}, then \texttt{list} is appended to the file. 

 If any element of \mbox{\texttt{\mdseries\slshape list}} is a semigroup, then the generators of that semigroup are written to \mbox{\texttt{\mdseries\slshape filename}}. More specifically, the list returned by \texttt{GeneratorsOfSemigroup} (\textbf{Reference: GeneratorsOfSemigroup}) is written to the file. 

 This function returns \texttt{IO{\textunderscore}OK} (\textbf{IO: IO{\textunderscore}OK}) if everything went well or \texttt{IO{\textunderscore}ERROR} (\textbf{IO: IO{\textunderscore}Error}) if something went wrong.

 The file produced by \texttt{WriteGenerators} can be read using \texttt{ReadGenerators} (\ref{ReadGenerators}). 

 From Version 3.0.0 onwards the \textsf{Semigroups} package used the \href{https://gap-packages.github.io/io} {IO}  package pickling functionality; see  (\textbf{IO: Pickling and unpickling}) for more details. This approach is used because it is more general and more
robust than the methods used by earlier versions of \textsf{Semigroups}, although the performance is somewhat worse, and the resulting files are
somewhat larger. 

 }

 

\subsection{\textcolor{Chapter }{IteratorFromGeneratorsFile}}
\logpage{[ 17, 1, 3 ]}\nobreak
\hyperdef{L}{X8711D6E280F87E67}{}
{\noindent\textcolor{FuncColor}{$\triangleright$\enspace\texttt{IteratorFromGeneratorsFile({\mdseries\slshape filename})\index{IteratorFromGeneratorsFile@\texttt{IteratorFromGeneratorsFile}}
\label{IteratorFromGeneratorsFile}
}\hfill{\scriptsize (function)}}\\
\textbf{\indent Returns:\ }
An iterator.



 If \mbox{\texttt{\mdseries\slshape filename}} is a file or a string containing the name of a file created using \texttt{WriteGenerators} (\ref{WriteGenerators}), then \texttt{IteratorFromGeneratorsFile} returns an iterator \texttt{iter} such that \texttt{NextIterator(iter)} returns the next collection of generators stored in the file \mbox{\texttt{\mdseries\slshape filename}}. 

 This function is a convenient way of, for example, looping over a collection
of generators in a file without loading every object in the file into memory.
This might be useful if the file contains more information than there is
available memory.

 If you want to get an iterator for a file written using \texttt{WriteGenerators} from a version of \textsf{Semigroups} before version 3.0.0, then you can use \texttt{IteratorFromOldGeneratorsFile}. }

 }

 
\section{\textcolor{Chapter }{Reading and writing multiplication tables to a file}}\logpage{[ 17, 2, 0 ]}
\hyperdef{L}{X7AB8E281795A4964}{}
{
  The functions \texttt{ReadMultiplicationTable} (\ref{ReadMultiplicationTable}) and \texttt{WriteMultiplicationTable} (\ref{WriteMultiplicationTable}) can be used to read or write, respectively, multiplication tables to a file. 

\subsection{\textcolor{Chapter }{ReadMultiplicationTable}}
\logpage{[ 17, 2, 1 ]}\nobreak
\hyperdef{L}{X805058C07F9373B4}{}
{\noindent\textcolor{FuncColor}{$\triangleright$\enspace\texttt{ReadMultiplicationTable({\mdseries\slshape filename[, nr]})\index{ReadMultiplicationTable@\texttt{ReadMultiplicationTable}}
\label{ReadMultiplicationTable}
}\hfill{\scriptsize (function)}}\\
\textbf{\indent Returns:\ }
A list of multiplication tables.



 If \mbox{\texttt{\mdseries\slshape filename}} is a file or is the name of a file created using \texttt{WriteMultiplicationTable} (\ref{WriteMultiplicationTable}), then \texttt{ReadMultiplicationTable} returns the contents of this file as a list of multiplication tables.

 If the optional second argument \mbox{\texttt{\mdseries\slshape nr}} is present, then \texttt{ReadMultiplicationTable} returns the multiplication table stored in the \mbox{\texttt{\mdseries\slshape nr}}th line of \mbox{\texttt{\mdseries\slshape filename}}.  
\begin{Verbatim}[commandchars=!@|,fontsize=\small,frame=single,label=Example]
  !gapprompt@gap>| !gapinput@file := Concatenation(SEMIGROUPS.PackageDir,|
  !gapprompt@>| !gapinput@"/data/tst/tables.gz");;|
  !gapprompt@gap>| !gapinput@tab := ReadMultiplicationTable(file, 12);|
  [ [ 1, 1, 3, 4, 5, 6, 7, 8, 9, 6 ], [ 1, 2, 3, 4, 5, 6, 7, 8, 9, 10 ],
    [ 3, 3, 1, 5, 4, 7, 6, 9, 8, 7 ], [ 4, 4, 9, 6, 3, 8, 5, 1, 7, 8 ],
    [ 5, 5, 8, 7, 1, 9, 4, 3, 6, 9 ], [ 6, 6, 7, 8, 9, 1, 3, 4, 5, 1 ],
    [ 7, 7, 6, 9, 8, 3, 1, 5, 4, 3 ], [ 8, 8, 5, 1, 7, 4, 9, 6, 3, 4 ],
    [ 9, 9, 4, 3, 6, 5, 8, 7, 1, 5 ], [ 6, 10, 7, 8, 9, 1, 3, 4, 5, 2 ]
   ]
\end{Verbatim}
 }

 

\subsection{\textcolor{Chapter }{WriteMultiplicationTable}}
\logpage{[ 17, 2, 2 ]}\nobreak
\hyperdef{L}{X868B824B7E24FA96}{}
{\noindent\textcolor{FuncColor}{$\triangleright$\enspace\texttt{WriteMultiplicationTable({\mdseries\slshape filename, list[, append]})\index{WriteMultiplicationTable@\texttt{WriteMultiplicationTable}}
\label{WriteMultiplicationTable}
}\hfill{\scriptsize (function)}}\\
\textbf{\indent Returns:\ }
\texttt{IO{\textunderscore}OK} or \texttt{IO{\textunderscore}ERROR}.



 This function provides a method for writing collections of multiplication
tables to a file. The resulting file can be further compressed using \texttt{gzip} or \texttt{xz}. This function applies to square arrays with a maximum of 255 rows where the
entries are integers from \texttt{[1, 2, .., n]} (where \texttt{n} is the number of rows in the array.

 The argument \mbox{\texttt{\mdseries\slshape list}} should be a list of multiplication tables.

 The argument \mbox{\texttt{\mdseries\slshape filename}} should be a file or a string containing the name of a file where the entries
in \mbox{\texttt{\mdseries\slshape list}} will be written or an \textsf{IO} package file object; see \texttt{IO{\textunderscore}File} (\textbf{IO: IO{\textunderscore}File mode}) and \texttt{IO{\textunderscore}CompressedFile} (\textbf{IO: IO{\textunderscore}CompressedFile}).

 If the optional third argument \mbox{\texttt{\mdseries\slshape append}} is not present or is given and equals \texttt{"w"}, then the previous content of the file is deleted and overwritten. If the
third argument is given and equals \texttt{"a"} then \texttt{list} is appended to the file. This function returns \texttt{IO{\textunderscore}OK} (\textbf{IO: IO{\textunderscore}OK}) if everything went well or \texttt{IO{\textunderscore}ERROR} (\textbf{IO: IO{\textunderscore}Error}) if something went wrong.

 The multiplication tables saved in \mbox{\texttt{\mdseries\slshape filename}} can be recovered from the file using \texttt{ReadMultiplicationTable} (\ref{ReadMultiplicationTable}). }

 

\subsection{\textcolor{Chapter }{IteratorFromMultiplicationTableFile}}
\logpage{[ 17, 2, 3 ]}\nobreak
\hyperdef{L}{X85708F5B7FBE3549}{}
{\noindent\textcolor{FuncColor}{$\triangleright$\enspace\texttt{IteratorFromMultiplicationTableFile({\mdseries\slshape filename})\index{IteratorFromMultiplicationTableFile@\texttt{IteratorFromMultiplicationTableFile}}
\label{IteratorFromMultiplicationTableFile}
}\hfill{\scriptsize (function)}}\\
\textbf{\indent Returns:\ }
An iterator.



 If \mbox{\texttt{\mdseries\slshape filename}} is a file or a string containing the name of a file created using \texttt{WriteMultiplicationTable} (\ref{WriteMultiplicationTable}), then \texttt{IteratorFromMultiplicationTableFile} returns an iterator \texttt{iter} such that \texttt{NextIterator(iter)} returns the next multiplication table stored in the file \mbox{\texttt{\mdseries\slshape filename}}.

 This function is a convenient way of, for example, looping over a collection
of multiplication tables in a file without loading every object in the file
into memory. This might be useful if the file contains more information than
there is available memory.

 }

 }

 }

     
\chapter{\textcolor{Chapter }{Translations}}\label{Chapter_Translations}
\logpage{[ 18, 0, 0 ]}
\hyperdef{L}{X7EC01B437CC2B2C9}{}
{
  In this chapter we describe the functionality in \textsf{Semigroups} for working with translations of semigroups. The notation used (as well as a
number of results) is based on \cite{Petrich1970aa}. Translations are of interest mainly due to their role in ideal extensions. A
description of this role can also be found in \cite{Petrich1970aa}. The implementation of translations in this package is only applicable to
finite semigroups satisfying \texttt{CanUseFroidurePin} (\ref{CanUseFroidurePin}). 

 For a semigroup $S$, a transformation $\lambda$ of $S$ (written on the left) is a left translation if for all $s, t$ in $S$, $\lambda (s)t = \lambda (st) $. A right translation $\rho$ (written on the right) is defined dually. 

 The set $L$ of left translations of $S$ is a semigroup under composition of functions, as is the set $R$ of right translations. The associativity of $S$ guarantees that left [right] multiplication by any element $s$ of $S$ represents a left [right] translation; this is the \emph{inner} left [right] translation $\lambda_s$ [$\rho_s$]. The inner left [right] translations form an ideal in $L$ [$R$]. 

 A left translation $\lambda$ and right translation $\rho$ are \emph{linked} if for all \texttt{s, t} in $S$, $s\lambda(t) = (s)\rho t$. A pair of linked translations is called a \emph{bitranslation}. The set of all bitranslations forms a semigroup $H$ called the \emph{translational hull} of $S$ where the operation is componentwise. If the components are inner translations
corresponding to multiplication by the same element, then the bitranslation is \emph{inner}. The inner bitranslations form an ideal of the translational hull. 

  Translations of a normalized Rees matrix semigroup $T$ (see \texttt{RMSNormalization} (\ref{RMSNormalization})) over a group $G$ can be represented through certain tuples, which can be computed very
efficiently compared to arbitrary translations. For left translations these
tuples consist of a function from the row indices of $T$ to G and a transformation on the row indices of $T$; the same is true of right translations and columns. More specifically, given
a normalised Rees matrix semigroup $S$ over a group $G$ with sandwich matrix $P$, rows $I$ and columns $J$, the left translations are characterised by pairs $(\theta, \chi)$ where $\theta$ is a function from $I$ to $G$ and $\chi$ is a transformation of $I$. The left translation $\lambda$ defined by $(\theta, \chi)$ acts on $S$ via 

 
\[ \lambda((i, a, \mu)) = ((i)\chi, (i)\theta \cdot a, \mu), \]
 

 where $i \in I$, $a \in G$, and $\mu \in J$ Dually, right translations $\rho$ are characterised by pairs $(\omega, \psi)$ where $\omega$ is a function from $J$ to $G$ and $\psi$ is a transformation of $J$, with action given by 

 
\[ ((i, a, \mu))\rho = (i, a \cdot (\mu)\psi, (\mu)\psi). \]
 

 Similarly, bitranslations $(\lambda, \rho)$ of $S$ can be characterised by triples $(g, \chi, \psi)$ such that $g \in G$, $\chi$ and $\psi$ are transformations of $I, J$ respectively, and 

 
\[ p_{\mu, (i)\chi} \cdot g \cdot p_{(1)\psi, i} = p_{\mu, (1)\chi} \cdot g \cdot
p_{(mu)\psi, i}. \]
 

 The action of $\lambda$ on $S$ is then given by 

 
\[ \lambda((i, a, \mu)) = ((i)\chi, b \cdot p_{(1)\psi, i} \cdot a, \mu), \]
 

 and of $\rho$ on $S$ by 

 
\[ ((i, a, \mu))\rho = (i, a \cdot p_{\mu, (1)\chi} \cdot b, (\mu)\psi). \]
 

 Further details may be found in \cite{Clifford1977aa}. 

 
\section{\textcolor{Chapter }{Methods for translations}}\label{Chapter_Translations_Section_Methods_for_translations}
\logpage{[ 18, 1, 0 ]}
\hyperdef{L}{X864C64877E5714AC}{}
{
  
\subsection{\textcolor{Chapter }{IsXTranslation}}\label{IsXTranslation}
\logpage{[ 18, 1, 1 ]}
\hyperdef{L}{X849F15607B774B90}{}
{
\noindent\textcolor{FuncColor}{$\triangleright$\enspace\texttt{IsSemigroupTranslation({\mdseries\slshape arg})\index{IsSemigroupTranslation@\texttt{IsSemigroupTranslation}!for IsAssociativeElement and IsMultiplicativeElementWithOne}
\label{IsSemigroupTranslation:for IsAssociativeElement and IsMultiplicativeElementWithOne}
}\hfill{\scriptsize (filter)}}\\
\noindent\textcolor{FuncColor}{$\triangleright$\enspace\texttt{IsLeftTranslation({\mdseries\slshape arg})\index{IsLeftTranslation@\texttt{IsLeftTranslation}!for IsSemigroupTranslation}
\label{IsLeftTranslation:for IsSemigroupTranslation}
}\hfill{\scriptsize (filter)}}\\
\noindent\textcolor{FuncColor}{$\triangleright$\enspace\texttt{IsRightTranslation({\mdseries\slshape arg})\index{IsRightTranslation@\texttt{IsRightTranslation}!for IsSemigroupTranslation}
\label{IsRightTranslation:for IsSemigroupTranslation}
}\hfill{\scriptsize (filter)}}\\
\textbf{\indent Returns:\ }
\texttt{true} or \texttt{false} 



 All, and only, left [right] translations belong to \texttt{IsLeftTranslation} [\texttt{IsRightTranslation}]. These are both subcategories of \texttt{IsSemigroupTranslation}, which itself is a subcategory of \texttt{IsAssociativeElement}. 
\begin{Verbatim}[commandchars=!@|,fontsize=\small,frame=single,label=Example]
  !gapprompt@gap>| !gapinput@S := RectangularBand(3, 4);;|
  !gapprompt@gap>| !gapinput@l := Representative(LeftTranslations(S));;|
  !gapprompt@gap>| !gapinput@IsSemigroupTranslation(l);|
  true
  !gapprompt@gap>| !gapinput@IsLeftTranslation(l);|
  true
  !gapprompt@gap>| !gapinput@IsRightTranslation(l);|
  false
  !gapprompt@gap>| !gapinput@l = One(LeftTranslations(S));|
  true
  !gapprompt@gap>| !gapinput@l * l = l;|
  true
\end{Verbatim}
 

 }

 

\subsection{\textcolor{Chapter }{IsBitranslation (for IsAssociativeElement and IsMultiplicativeElementWithOne)}}
\logpage{[ 18, 1, 2 ]}\nobreak
\hyperdef{L}{X7F6689E885982816}{}
{\noindent\textcolor{FuncColor}{$\triangleright$\enspace\texttt{IsBitranslation({\mdseries\slshape arg})\index{IsBitranslation@\texttt{IsBitranslation}!for IsAssociativeElement and IsMultiplicativeElementWithOne}
\label{IsBitranslation:for IsAssociativeElement and IsMultiplicativeElementWithOne}
}\hfill{\scriptsize (filter)}}\\
\textbf{\indent Returns:\ }
\texttt{true} or \texttt{false} 



 All, and only, bitranslations belong to \texttt{IsBitranslation}. This is a subcategory of \texttt{IsAssociativeElement} (\textbf{Reference: IsAssociativeElement}). 
\begin{Verbatim}[commandchars=!@|,fontsize=\small,frame=single,label=Example]
  !gapprompt@gap>| !gapinput@G := Group((1, 2), (3, 4));;|
  !gapprompt@gap>| !gapinput@A := AsList(G);;|
  !gapprompt@gap>| !gapinput@mat := [[A[1], 0, A[1]],|
  !gapprompt@>| !gapinput@[A[2], A[2], A[4]]];;|
  !gapprompt@gap>| !gapinput@S := ReesZeroMatrixSemigroup(G, mat);;|
  !gapprompt@gap>| !gapinput@L := LeftTranslations(S);;|
  !gapprompt@gap>| !gapinput@R := RightTranslations(S);;|
  !gapprompt@gap>| !gapinput@l := OneOp(Representative(L));;|
  !gapprompt@gap>| !gapinput@r := OneOp(Representative(R));;|
  !gapprompt@gap>| !gapinput@H := TranslationalHull(S);;|
  !gapprompt@gap>| !gapinput@x := Bitranslation(H, l, r);;|
  !gapprompt@gap>| !gapinput@IsBitranslation(x);|
  true
  !gapprompt@gap>| !gapinput@IsSemigroupTranslation(x);|
  false
\end{Verbatim}
 }

 
\subsection{\textcolor{Chapter }{IsXTranslationCollection}}\label{IsXTranslationCollection}
\logpage{[ 18, 1, 3 ]}
\hyperdef{L}{X7F536B1B85978B63}{}
{
\noindent\textcolor{FuncColor}{$\triangleright$\enspace\texttt{IsSemigroupTranslationCollection({\mdseries\slshape obj})\index{IsSemigroupTranslationCollection@\texttt{IsSemigroupTranslationCollection}}
\label{IsSemigroupTranslationCollection}
}\hfill{\scriptsize (filter)}}\\
\noindent\textcolor{FuncColor}{$\triangleright$\enspace\texttt{IsLeftTranslationCollection({\mdseries\slshape obj})\index{IsLeftTranslationCollection@\texttt{IsLeftTranslationCollection}}
\label{IsLeftTranslationCollection}
}\hfill{\scriptsize (filter)}}\\
\noindent\textcolor{FuncColor}{$\triangleright$\enspace\texttt{IsRightTranslationCollection({\mdseries\slshape obj})\index{IsRightTranslationCollection@\texttt{IsRightTranslationCollection}}
\label{IsRightTranslationCollection}
}\hfill{\scriptsize (filter)}}\\
\noindent\textcolor{FuncColor}{$\triangleright$\enspace\texttt{IsBitranslationCollection({\mdseries\slshape obj})\index{IsBitranslationCollection@\texttt{IsBitranslationCollection}}
\label{IsBitranslationCollection}
}\hfill{\scriptsize (filter)}}\\
\textbf{\indent Returns:\ }
\texttt{true} or \texttt{false} 



 Every collection of left\texttt{\symbol{45}}, right\texttt{\symbol{45}}, or
bi\texttt{\symbol{45}}translations belongs to the respective category \texttt{IsXTranslationCollection} (\ref{IsXTranslationCollection}). 

 }

 
\subsection{\textcolor{Chapter }{XPartOfBitranslation}}\label{XPartOfBitranslation}
\logpage{[ 18, 1, 4 ]}
\hyperdef{L}{X7D52D17E7A28CE0E}{}
{
\noindent\textcolor{FuncColor}{$\triangleright$\enspace\texttt{LeftPartOfBitranslation({\mdseries\slshape h})\index{LeftPartOfBitranslation@\texttt{LeftPartOfBitranslation}}
\label{LeftPartOfBitranslation}
}\hfill{\scriptsize (function)}}\\
\noindent\textcolor{FuncColor}{$\triangleright$\enspace\texttt{RightPartOfBitranslation({\mdseries\slshape arg})\index{RightPartOfBitranslation@\texttt{RightPartOfBitranslation}}
\label{RightPartOfBitranslation}
}\hfill{\scriptsize (function)}}\\
\textbf{\indent Returns:\ }
a left or right translation 



 For a Bitranslation \mbox{\texttt{\mdseries\slshape h}} consisting of a linked pair $(l, r)$, of left and right translations, \texttt{LeftPartOfBitranslation(\mbox{\texttt{\mdseries\slshape b}})} returns the left translation \texttt{l}, and \texttt{RightPartOfBitranslation(\mbox{\texttt{\mdseries\slshape b}})} returns the right translation \texttt{r}. 

 }

 
\subsection{\textcolor{Chapter }{XTranslation}}\label{XTranslation}
\logpage{[ 18, 1, 5 ]}
\hyperdef{L}{X7ACCBAB57E910910}{}
{
\noindent\textcolor{FuncColor}{$\triangleright$\enspace\texttt{LeftTranslation({\mdseries\slshape T, x[, y]})\index{LeftTranslation@\texttt{LeftTranslation}!for IsLeftTranslationsSemigroup, IsGeneralMapping}
\label{LeftTranslation:for IsLeftTranslationsSemigroup, IsGeneralMapping}
}\hfill{\scriptsize (operation)}}\\
\noindent\textcolor{FuncColor}{$\triangleright$\enspace\texttt{RightTranslation({\mdseries\slshape arg1, arg2})\index{RightTranslation@\texttt{RightTranslation}!for IsRightTranslationsSemigroup, IsGeneralMapping}
\label{RightTranslation:for IsRightTranslationsSemigroup, IsGeneralMapping}
}\hfill{\scriptsize (operation)}}\\
\textbf{\indent Returns:\ }
a left or right translation 



 For the semigroup \mbox{\texttt{\mdseries\slshape T}} of left or right translations of a semigroup $ S$ and \mbox{\texttt{\mdseries\slshape x}} one of: 
\begin{itemize}
\item  a mapping on the underlying semigroup; note that in this case only the values
of the mapping on the \texttt{UnderlyingRepresentatives} of \mbox{\texttt{\mdseries\slshape T}} are checked and used, so mappings which do not define translations can be used
to create translations if they are valid on that subset of S; 
\item  a list of indices representing the images of the \texttt{UnderlyingRepresentatives} of \mbox{\texttt{\mdseries\slshape T}}, where the ordering is that of \texttt{PositionCanonical} (\ref{PositionCanonical}) on \mbox{\texttt{\mdseries\slshape S}}; 
\item  (for \texttt{LeftTranslation}) a list of length \texttt{Length(Rows(S))} containing elements of \texttt{UnderlyingSemigroup(S)}; in this case \mbox{\texttt{\mdseries\slshape S}} must be a normalised Rees matrix semigroup and \texttt{y} must be a Transformation of \texttt{Rows(S)}; 
\item  (for \texttt{RightTranslation}) a list of length \texttt{Length(Columns(S))} containing elements of \texttt{UnderlyingSemigroup(S)}; in this case \mbox{\texttt{\mdseries\slshape S}} must be a normalised Rees matrix semigroup and \texttt{y} must be a Transformation of \texttt{Columns(S)}; 
\end{itemize}
 \texttt{LeftTranslation} and \texttt{RightTranslation} return the corresponding translations. 
\begin{Verbatim}[commandchars=!@|,fontsize=\small,frame=single,label=Example]
  !gapprompt@gap>| !gapinput@S := RectangularBand(3, 4);;|
  !gapprompt@gap>| !gapinput@L := LeftTranslations(S);;|
  !gapprompt@gap>| !gapinput@s := AsList(S)[1];;|
  !gapprompt@gap>| !gapinput@map := MappingByFunction(S, S, x -> s * x);;|
  !gapprompt@gap>| !gapinput@l := LeftTranslation(L, map);|
  <left translation on <regular transformation semigroup of size 12, 
   degree 8 with 4 generators>>
  !gapprompt@gap>| !gapinput@s ^ l;|
  Transformation( [ 1, 2, 1, 1, 5, 5, 5, 5 ] )
\end{Verbatim}
 

 }

 

\subsection{\textcolor{Chapter }{Bitranslation (for IsBitranslationsSemigroup, IsLeftTranslation, IsRightTranslation)}}
\logpage{[ 18, 1, 6 ]}\nobreak
\hyperdef{L}{X8664424983C3281F}{}
{\noindent\textcolor{FuncColor}{$\triangleright$\enspace\texttt{Bitranslation({\mdseries\slshape H, l, r})\index{Bitranslation@\texttt{Bitranslation}!for IsBitranslationsSemigroup, IsLeftTranslation, IsRightTranslation}
\label{Bitranslation:for IsBitranslationsSemigroup, IsLeftTranslation, IsRightTranslation}
}\hfill{\scriptsize (operation)}}\\
\textbf{\indent Returns:\ }
a bitranslation 



 If \mbox{\texttt{\mdseries\slshape H}} is a translational hull over a semigroup $S$, and \mbox{\texttt{\mdseries\slshape l}} and \mbox{\texttt{\mdseries\slshape r}} are linked left and right translations respectively over $S$, then this function returns the bitranslation $(\mbox{\texttt{\mdseries\slshape l}}, \mbox{\texttt{\mdseries\slshape r}})$. If \mbox{\texttt{\mdseries\slshape l}} and \mbox{\texttt{\mdseries\slshape r}} are not linked, then an error is produced. 
\begin{Verbatim}[commandchars=!@|,fontsize=\small,frame=single,label=Example]
  !gapprompt@gap>| !gapinput@G := Group((1, 2), (3, 4));;|
  !gapprompt@gap>| !gapinput@A := AsList(G);;|
  !gapprompt@gap>| !gapinput@mat := [[A[1], 0],|
  !gapprompt@>| !gapinput@[A[2], A[2]]];;|
  !gapprompt@gap>| !gapinput@S := ReesZeroMatrixSemigroup(G, mat);;|
  !gapprompt@gap>| !gapinput@L := LeftTranslations(S);;|
  !gapprompt@gap>| !gapinput@R := RightTranslations(S);;|
  !gapprompt@gap>| !gapinput@l := LeftTranslation(L, MappingByFunction(S, S, s -> S.1 * s));;|
  !gapprompt@gap>| !gapinput@r := RightTranslation(R, MappingByFunction(S, S, s -> s * S.1));;|
  !gapprompt@gap>| !gapinput@H := TranslationalHull(S);;|
  !gapprompt@gap>| !gapinput@x := Bitranslation(H, l, r);|
  <bitranslation on <regular semigroup of size 17, with 4 generators>>
\end{Verbatim}
 }

 
\subsection{\textcolor{Chapter }{UnderlyingSemigroup}}\label{UnderlyingSemigroup}
\logpage{[ 18, 1, 7 ]}
\hyperdef{L}{X7B5BB0BA8683A021}{}
{
\noindent\textcolor{FuncColor}{$\triangleright$\enspace\texttt{UnderlyingSemigroup({\mdseries\slshape S})\index{UnderlyingSemigroup@\texttt{UnderlyingSemigroup}!for IsTranslationsSemigroup}
\label{UnderlyingSemigroup:for IsTranslationsSemigroup}
}\hfill{\scriptsize (attribute)}}\\
\noindent\textcolor{FuncColor}{$\triangleright$\enspace\texttt{UnderlyingSemigroup({\mdseries\slshape arg})\index{UnderlyingSemigroup@\texttt{UnderlyingSemigroup}!for IsBitranslationsSemigroup}
\label{UnderlyingSemigroup:for IsBitranslationsSemigroup}
}\hfill{\scriptsize (attribute)}}\\
\textbf{\indent Returns:\ }
a semigroup 



 Given a semigroup of translations or bitranslations, returns the semigroup on
which these translations act. 

 }

 
\subsection{\textcolor{Chapter }{XTranslationsSemigroupOfFamily}}\label{XTranslationsSemigroupOfFamily}
\logpage{[ 18, 1, 8 ]}
\hyperdef{L}{X857C28C8790A35F6}{}
{
\noindent\textcolor{FuncColor}{$\triangleright$\enspace\texttt{LeftTranslationsSemigroupOfFamily({\mdseries\slshape fam})\index{LeftTranslationsSemigroupOfFamily@\texttt{LeftTranslationsSemigroupOfFamily}!for IsFamily}
\label{LeftTranslationsSemigroupOfFamily:for IsFamily}
}\hfill{\scriptsize (attribute)}}\\
\noindent\textcolor{FuncColor}{$\triangleright$\enspace\texttt{RightTranslationsSemigroupOfFamily({\mdseries\slshape arg})\index{RightTranslationsSemigroupOfFamily@\texttt{RightTranslationsSemigroupOfFamily}!for IsFamily}
\label{RightTranslationsSemigroupOfFamily:for IsFamily}
}\hfill{\scriptsize (attribute)}}\\
\noindent\textcolor{FuncColor}{$\triangleright$\enspace\texttt{TranslationalHullOfFamily({\mdseries\slshape arg})\index{TranslationalHullOfFamily@\texttt{TranslationalHullOfFamily}!for IsFamily}
\label{TranslationalHullOfFamily:for IsFamily}
}\hfill{\scriptsize (attribute)}}\\
\textbf{\indent Returns:\ }
 the semigroup of left or right translations, or the translational hull 



 Given a family \mbox{\texttt{\mdseries\slshape fam}} of left\texttt{\symbol{45}}, right\texttt{\symbol{45}} or
bi\texttt{\symbol{45}}translations, returns the translations semigroup or
translational hull to which they belong. 
\begin{Verbatim}[commandchars=!@|,fontsize=\small,frame=single,label=Example]
  !gapprompt@gap>| !gapinput@S := RectangularBand(3, 3);;|
  !gapprompt@gap>| !gapinput@L := LeftTranslations(S);;|
  !gapprompt@gap>| !gapinput@l := Representative(L);;|
  !gapprompt@gap>| !gapinput@LeftTranslationsSemigroupOfFamily(FamilyObj(l)) = L;|
  true
  !gapprompt@gap>| !gapinput@H := TranslationalHull(S);;|
  !gapprompt@gap>| !gapinput@h := Representative(H);;|
  !gapprompt@gap>| !gapinput@TranslationalHullOfFamily(FamilyObj(h)) = H;|
  true
\end{Verbatim}
 

 }

 
\subsection{\textcolor{Chapter }{TypeXTranslationSemigroupElements}}\label{TypeXTranslationSemigroupElements}
\logpage{[ 18, 1, 9 ]}
\hyperdef{L}{X7AA681A283A22C28}{}
{
\noindent\textcolor{FuncColor}{$\triangleright$\enspace\texttt{TypeLeftTranslationsSemigroupElements({\mdseries\slshape arg})\index{TypeLeftTranslationsSemigroupElements@\texttt{Type}\-\texttt{Left}\-\texttt{Translations}\-\texttt{Semigroup}\-\texttt{Elements}!for IsLeftTranslationsSemigroup}
\label{TypeLeftTranslationsSemigroupElements:for IsLeftTranslationsSemigroup}
}\hfill{\scriptsize (attribute)}}\\
\noindent\textcolor{FuncColor}{$\triangleright$\enspace\texttt{TypeRightTranslationsSemigroupElements({\mdseries\slshape arg})\index{TypeRightTranslationsSemigroupElements@\texttt{Type}\-\texttt{Right}\-\texttt{Translations}\-\texttt{Semigroup}\-\texttt{Elements}!for IsRightTranslationsSemigroup}
\label{TypeRightTranslationsSemigroupElements:for IsRightTranslationsSemigroup}
}\hfill{\scriptsize (attribute)}}\\
\noindent\textcolor{FuncColor}{$\triangleright$\enspace\texttt{TypeBitranslations({\mdseries\slshape arg})\index{TypeBitranslations@\texttt{TypeBitranslations}!for IsBitranslationsSemigroup}
\label{TypeBitranslations:for IsBitranslationsSemigroup}
}\hfill{\scriptsize (attribute)}}\\
\textbf{\indent Returns:\ }
a type 



 Given a (bi)translations semigroup, returns the type of the elements that it
contains. 

 }

 
\subsection{\textcolor{Chapter }{XTranslations}}\label{XTranslations}
\logpage{[ 18, 1, 10 ]}
\hyperdef{L}{X7D5CC8A48371410D}{}
{
\noindent\textcolor{FuncColor}{$\triangleright$\enspace\texttt{LeftTranslations({\mdseries\slshape S})\index{LeftTranslations@\texttt{LeftTranslations}!for IsSemigroup and CanUseFroidurePin and IsFinite}
\label{LeftTranslations:for IsSemigroup and CanUseFroidurePin and IsFinite}
}\hfill{\scriptsize (attribute)}}\\
\noindent\textcolor{FuncColor}{$\triangleright$\enspace\texttt{RightTranslations({\mdseries\slshape arg})\index{RightTranslations@\texttt{RightTranslations}!for IsSemigroup and CanUseFroidurePin and IsFinite}
\label{RightTranslations:for IsSemigroup and CanUseFroidurePin and IsFinite}
}\hfill{\scriptsize (attribute)}}\\
\textbf{\indent Returns:\ }
the semigroup of left or right translations 



 Given a finite semigroup \mbox{\texttt{\mdseries\slshape S}} satisfying \texttt{CanUseFroidurePin} (\ref{CanUseFroidurePin}), returns the semigroup of all left or right translations of \mbox{\texttt{\mdseries\slshape S}}. 
\begin{Verbatim}[commandchars=!@|,fontsize=\small,frame=single,label=Example]
  !gapprompt@gap>| !gapinput@S := Semigroup([Transformation([1, 4, 3, 3, 6, 5]),|
  !gapprompt@>| !gapinput@Transformation([3, 4, 1, 1, 4, 2])]);;|
  !gapprompt@gap>| !gapinput@L := LeftTranslations(S);|
  <the semigroup of left translations of <transformation semigroup of 
   degree 6 with 2 generators>>
  !gapprompt@gap>| !gapinput@Size(L);|
  361
\end{Verbatim}
 

 }

 

\subsection{\textcolor{Chapter }{TranslationalHull (for IsSemigroup and CanUseFroidurePin and IsFinite)}}
\logpage{[ 18, 1, 11 ]}\nobreak
\hyperdef{L}{X8231194386449BD4}{}
{\noindent\textcolor{FuncColor}{$\triangleright$\enspace\texttt{TranslationalHull({\mdseries\slshape S})\index{TranslationalHull@\texttt{TranslationalHull}!for IsSemigroup and CanUseFroidurePin and IsFinite}
\label{TranslationalHull:for IsSemigroup and CanUseFroidurePin and IsFinite}
}\hfill{\scriptsize (attribute)}}\\
\textbf{\indent Returns:\ }
the translational hull 



 Given a finite semigroup \mbox{\texttt{\mdseries\slshape S}} satisfying \texttt{CanUseFroidurePin} (\ref{CanUseFroidurePin}), returns the translational hull of \mbox{\texttt{\mdseries\slshape S}}. 
\begin{Verbatim}[commandchars=!@|,fontsize=\small,frame=single,label=Example]
  !gapprompt@gap>| !gapinput@S := Semigroup([Transformation([1, 4, 3, 3, 6, 5]),|
  !gapprompt@>| !gapinput@Transformation([3, 4, 1, 1, 4, 2])]);;|
  !gapprompt@gap>| !gapinput@H := TranslationalHull(S);|
  <translational hull over <transformation semigroup of degree 6 with 2 
   generators>>
  !gapprompt@gap>| !gapinput@Size(H);|
  38
\end{Verbatim}
 }

 
\subsection{\textcolor{Chapter }{NrXTranslations}}\label{NrXTranslations}
\logpage{[ 18, 1, 12 ]}
\hyperdef{L}{X7C826FBA78739FA4}{}
{
\noindent\textcolor{FuncColor}{$\triangleright$\enspace\texttt{NrLeftTranslations({\mdseries\slshape S})\index{NrLeftTranslations@\texttt{NrLeftTranslations}!for IsSemigroup and CanUseFroidurePin and IsFinite}
\label{NrLeftTranslations:for IsSemigroup and CanUseFroidurePin and IsFinite}
}\hfill{\scriptsize (attribute)}}\\
\noindent\textcolor{FuncColor}{$\triangleright$\enspace\texttt{NrRightTranslations({\mdseries\slshape arg})\index{NrRightTranslations@\texttt{NrRightTranslations}!for IsSemigroup and CanUseFroidurePin and IsFinite}
\label{NrRightTranslations:for IsSemigroup and CanUseFroidurePin and IsFinite}
}\hfill{\scriptsize (attribute)}}\\
\noindent\textcolor{FuncColor}{$\triangleright$\enspace\texttt{NrBitranslations({\mdseries\slshape arg})\index{NrBitranslations@\texttt{NrBitranslations}!for IsSemigroup and CanUseFroidurePin and IsFinite}
\label{NrBitranslations:for IsSemigroup and CanUseFroidurePin and IsFinite}
}\hfill{\scriptsize (attribute)}}\\
\textbf{\indent Returns:\ }
the number of left\texttt{\symbol{45}}, right\texttt{\symbol{45}}, or
bi\texttt{\symbol{45}}translations 



 Given a finite semigroup \mbox{\texttt{\mdseries\slshape S}} satisfying \texttt{CanUseFroidurePin} (\ref{CanUseFroidurePin}), returns the number of left\texttt{\symbol{45}}, right\texttt{\symbol{45}},
or bi\texttt{\symbol{45}}translations of \mbox{\texttt{\mdseries\slshape S}}. This is typically more efficient than calling \texttt{Size(LeftTranslations(\mbox{\texttt{\mdseries\slshape S}}))}, as the [bi]translations may not be produced. 
\begin{Verbatim}[commandchars=!@|,fontsize=\small,frame=single,label=Example]
  !gapprompt@gap>| !gapinput@S := Semigroup([Transformation([1, 4, 3, 3, 6, 5, 2]),|
  !gapprompt@>| !gapinput@Transformation([3, 4, 1, 1, 4, 2])]);;|
  !gapprompt@gap>| !gapinput@NrLeftTranslations(S);|
  1444
  !gapprompt@gap>| !gapinput@NrRightTranslations(S);|
  398
  !gapprompt@gap>| !gapinput@NrBitranslations(S);|
  69
\end{Verbatim}
 

 }

 
\subsection{\textcolor{Chapter }{InnerXTranslations}}\label{InnerXTranslations}
\logpage{[ 18, 1, 13 ]}
\hyperdef{L}{X7E9306DF79587A33}{}
{
\noindent\textcolor{FuncColor}{$\triangleright$\enspace\texttt{InnerLeftTranslations({\mdseries\slshape S})\index{InnerLeftTranslations@\texttt{InnerLeftTranslations}!for IsSemigroup and CanUseFroidurePin and IsFinite}
\label{InnerLeftTranslations:for IsSemigroup and CanUseFroidurePin and IsFinite}
}\hfill{\scriptsize (attribute)}}\\
\noindent\textcolor{FuncColor}{$\triangleright$\enspace\texttt{InnerRightTranslations({\mdseries\slshape arg})\index{InnerRightTranslations@\texttt{InnerRightTranslations}!for IsSemigroup and CanUseFroidurePin and IsFinite}
\label{InnerRightTranslations:for IsSemigroup and CanUseFroidurePin and IsFinite}
}\hfill{\scriptsize (attribute)}}\\
\textbf{\indent Returns:\ }
the monoid of inner left or right translations 



 For a finite semigroup \mbox{\texttt{\mdseries\slshape S}} satisfying \texttt{CanUseFroidurePin} (\ref{CanUseFroidurePin}), \texttt{InnerLeftTranslations}(\mbox{\texttt{\mdseries\slshape S}}) returns the inner left translations of S (i.e. the translations defined by
left multiplication by a fixed element of \mbox{\texttt{\mdseries\slshape S}}), and \texttt{InnerRightTranslations}(\mbox{\texttt{\mdseries\slshape S}}) returns the inner right translations of \mbox{\texttt{\mdseries\slshape S}} (i.e. the translations defined by right multiplication by a fixed element of \mbox{\texttt{\mdseries\slshape S}}). 
\begin{Verbatim}[commandchars=!@|,fontsize=\small,frame=single,label=Example]
  !gapprompt@gap>| !gapinput@S := Semigroup([Transformation([1, 4, 3, 3, 6, 5]),|
  !gapprompt@>| !gapinput@Transformation([3, 4, 1, 1, 4, 2])]);;|
  !gapprompt@gap>| !gapinput@I := InnerLeftTranslations(S);|
  <left translations semigroup over <transformation semigroup 
   of size 22, degree 6 with 2 generators>>
  !gapprompt@gap>| !gapinput@Size(I) <= Size(S);|
  true
\end{Verbatim}
 

 }

 

\subsection{\textcolor{Chapter }{InnerTranslationalHull (for IsSemigroup and CanUseFroidurePin and IsFinite)}}
\logpage{[ 18, 1, 14 ]}\nobreak
\hyperdef{L}{X7C109DF080E72F68}{}
{\noindent\textcolor{FuncColor}{$\triangleright$\enspace\texttt{InnerTranslationalHull({\mdseries\slshape S})\index{InnerTranslationalHull@\texttt{InnerTranslationalHull}!for IsSemigroup and CanUseFroidurePin and IsFinite}
\label{InnerTranslationalHull:for IsSemigroup and CanUseFroidurePin and IsFinite}
}\hfill{\scriptsize (attribute)}}\\
\textbf{\indent Returns:\ }
the inner translational hull 



 Given a finite semigroup \mbox{\texttt{\mdseries\slshape S}} satisfying \texttt{CanUseFroidurePin} (\ref{CanUseFroidurePin}), returns the inner translational hull of \mbox{\texttt{\mdseries\slshape S}}, i.e. the bitranslations whose left and right translation components are
inner translations defined by left and right multiplication by the same fixed
element of \mbox{\texttt{\mdseries\slshape S}}. 
\begin{Verbatim}[commandchars=!@|,fontsize=\small,frame=single,label=Example]
  !gapprompt@gap>| !gapinput@S := Semigroup([Transformation([1, 4, 3, 3, 6, 5]),|
  !gapprompt@>| !gapinput@Transformation([3, 4, 1, 1, 4, 2])]);;|
  !gapprompt@gap>| !gapinput@I := InnerTranslationalHull(S);|
  <semigroup of bitranslations over <transformation semigroup 
   of size 22, degree 6 with 2 generators>>
  !gapprompt@gap>| !gapinput@L := LeftTranslations(S);;|
  !gapprompt@gap>| !gapinput@R := RightTranslations(S);;|
  !gapprompt@gap>| !gapinput@H := TranslationalHull(S);;|
  !gapprompt@gap>| !gapinput@inners := [];;|
  !gapprompt@gap>| !gapinput@for s in S do|
  !gapprompt@>| !gapinput@l := LeftTranslation(L, MappingByFunction(S, S, x -> s * x));|
  !gapprompt@>| !gapinput@r := RightTranslation(R, MappingByFunction(S, S, x -> x * s));|
  !gapprompt@>| !gapinput@AddSet(inners, Bitranslation(H, l, r));|
  !gapprompt@>| !gapinput@od;|
  !gapprompt@gap>| !gapinput@Set(I) = inners;|
  true
\end{Verbatim}
 }

 

\subsection{\textcolor{Chapter }{UnderlyingRepresentatives (for IsTranslationsSemigroup)}}
\logpage{[ 18, 1, 15 ]}\nobreak
\hyperdef{L}{X87E30C387BE79FB8}{}
{\noindent\textcolor{FuncColor}{$\triangleright$\enspace\texttt{UnderlyingRepresentatives({\mdseries\slshape T})\index{UnderlyingRepresentatives@\texttt{UnderlyingRepresentatives}!for IsTranslationsSemigroup}
\label{UnderlyingRepresentatives:for IsTranslationsSemigroup}
}\hfill{\scriptsize (attribute)}}\\
\textbf{\indent Returns:\ }
a set of representatives 



 For efficiency, we typically store translations on a semigroup $S$ as their actions on a small subset of $S$. For left translations, this is a set of representatives of the maximal $\mathcal{R}$\texttt{\symbol{45}}classes of $S$; dually for right translations we use representatives of the maximal $\mathcal{L}$\texttt{\symbol{45}}classes. You can use \texttt{UnderlyingRepresentatives} to access these representatives. 
\begin{Verbatim}[commandchars=!@|,fontsize=\small,frame=single,label=Example]
  !gapprompt@gap>| !gapinput@G := Range(IsomorphismPermGroup(SmallGroup(12, 1)));;|
  !gapprompt@gap>| !gapinput@mat := [[G.1, G.2], [G.1, G.1],|
  !gapprompt@>| !gapinput@[G.2, G.3], [G.1 * G.2, G.1 * G.3]];;|
  !gapprompt@gap>| !gapinput@S := ReesMatrixSemigroup(G, mat);;|
  !gapprompt@gap>| !gapinput@L := LeftTranslations(S);;|
  !gapprompt@gap>| !gapinput@R := RightTranslations(S);;|
  !gapprompt@gap>| !gapinput@UnderlyingRepresentatives(L);|
  [ (1,(),1), (2,(),2) ]
  !gapprompt@gap>| !gapinput@UnderlyingRepresentatives(R);|
  [ (1,(),1), (2,(),2), (1,(),3), (1,(),4) ]
\end{Verbatim}
 }

 

\subsection{\textcolor{Chapter }{ImageSetOfTranslation (for IsSemigroupTranslation)}}
\logpage{[ 18, 1, 16 ]}\nobreak
\hyperdef{L}{X7E81252986BB72BB}{}
{\noindent\textcolor{FuncColor}{$\triangleright$\enspace\texttt{ImageSetOfTranslation({\mdseries\slshape t})\index{ImageSetOfTranslation@\texttt{ImageSetOfTranslation}!for IsSemigroupTranslation}
\label{ImageSetOfTranslation:for IsSemigroupTranslation}
}\hfill{\scriptsize (operation)}}\\
\textbf{\indent Returns:\ }
a set of elements 



 Given a left or right translation \mbox{\texttt{\mdseries\slshape t}} on a semigroup $S$, returns the set of elements of $S$ lying in the image of \mbox{\texttt{\mdseries\slshape t}}. 
\begin{Verbatim}[commandchars=!@|,fontsize=\small,frame=single,label=Example]
  !gapprompt@gap>| !gapinput@S := Semigroup([Transformation([1, 3, 3, 4]),|
  !gapprompt@>| !gapinput@Transformation([3, 4, 1, 2])]);;|
  !gapprompt@gap>| !gapinput@t := Set(LeftTranslations(S))[4];|
  <left translation on <transformation semigroup of size 8, degree 4 
   with 2 generators>>
  !gapprompt@gap>| !gapinput@ImageSetOfTranslation(t);|
  [ Transformation( [ 1, 2, 3, 1 ] ), Transformation( [ 1, 3, 3, 1 ] ), 
    Transformation( [ 3, 1, 1, 3 ] ), Transformation( [ 3, 4, 1, 3 ] ) ]
\end{Verbatim}
 }

 }

 }

 \def\bibname{References\logpage{[ "Bib", 0, 0 ]}
\hyperdef{L}{X7A6F98FD85F02BFE}{}
}

\bibliographystyle{alpha}
\bibliography{bibliography.xml}

\addcontentsline{toc}{chapter}{References}

\def\indexname{Index\logpage{[ "Ind", 0, 0 ]}
\hyperdef{L}{X83A0356F839C696F}{}
}

\cleardoublepage
\phantomsection
\addcontentsline{toc}{chapter}{Index}


\printindex

\immediate\write\pagenrlog{["Ind", 0, 0], \arabic{page},}
\newpage
\immediate\write\pagenrlog{["End"], \arabic{page}];}
\immediate\closeout\pagenrlog
\end{document}
